% ============================================================
%  R. Nielsen, "Evangelietroen og den moderne Bevidsthed.
%  Forelæsninger over Jesu Liv. Første Deel"
%  [Gospel Faith and the Modern Consciousness. Lectures on the Life of Jesus.
%  First Part]. Kjøbenhavn: C. A. Reitzel, 1849. X + 530 pp.
%  TRANSLATION (English)
%  Source of truth: transcription.tex (Danish), COMPLETE; sampled clean
%  2026-09-26 (12 pp.; see ACCURACY-LEDGER.tsv).
%  Mirrors transcription.tex 1:1; every printed page carries \opage{N} at the
%  corresponding point in the English, with its "% --- p. N ---" line.
%  Translated 26-27 Sep 2026 in 22 batches of about 24 printed pages
%  (TRANSLATION-PLAYBOOK.md). German, Latin and Greek quotations are kept as
%  printed, misprints included; printer's defects in the Danish (turned
%  sorts such as "uok", "Paastaud"; "Verden" for "Vreden", p. 426) are
%  translated by their intended sense; the nine Rettelser, applied in the
%  Danish, are applied here too. Scripture is translated from Nielsen's
%  Danish (Authorized Version phrasing where it is faithful to it); bold
%  (Fraktur) and letterspaced scripture keep \textbf / \emph.
%  Terms: Evangelietroen = the faith of the Gospel / Gospel faith; den moderne
%  Bevidsthed = the modern consciousness; den Troende = the Believer;
%  Fritænkeren = the Free Thinker; Middelveien = the Middle Way; Under =
%  wonder, Undergjerning = wonder-work, Mirakel = miracle (Nielsen keeps the
%  two apart: p. 40); Aabenbaring = revelation, Aabenbarelse = manifestation;
%  Forargelse = offense; Anstød = stumbling-block; Tilegnelse = appropriation;
%  Inderlighed = inwardness; Forgjængeren = the Forerunner; forklare =
%  explain / transfigure (Nielsen's pun, pp. 345, 435, cannot be kept).
% ============================================================
\documentclass[12pt,a4paper]{book}
\usepackage[utf8]{inputenc}
\usepackage[T1]{fontenc}
\usepackage{libertinus}
\usepackage{libertinust1math}
\usepackage{textalpha}   % Greek typed directly
\usepackage{graphicx}    % for \reflectbox (the old Danish »ɔ:« id-est mark)
\DeclareUnicodeCharacter{0254}{\reflectbox{c}}  % ɔ (U+0254) = reversed c
\usepackage[english]{babel}
\usepackage[protrusion=true,expansion=true]{microtype}
\usepackage{setspace}
\usepackage[margin=1.2in]{geometry}
\usepackage{enumitem}
\usepackage{fancyhdr}
\setlength{\headheight}{28pt}
\addtolength{\topmargin}{-13.5pt}

% Footnotes as the 1849 printing sets them: *) rather than 1, 2, 3.
% A constant \thefootnote rather than \fnsymbol with footmisc[perpage], both of
% which were tried on the 1850 volume and are wrong:
%   - \fnsymbol is math mode, so it sets U+2217 ASTERISK OPERATOR, not a text
%     asterisk;
%   - [perpage] resets on the TYPESET page, not the book's printed page, so a
%     second note landing on the same typeset page came out as a dagger.
% If a printed page of this book is found to carry two notes with distinct
% marks, log it and revisit this line — do not silently renumber.
\renewcommand{\thefootnote}{*)}

\usepackage{hyperref}

\hypersetup{
  pdftitle={Gospel Faith and the Modern Consciousness},
  pdfauthor={R. Nielsen},
  hidelinks
}

\pagestyle{fancy}
\fancyhf{}
\fancyhead[LE,RO]{\thepage}
\fancyhead[RE]{\textit{Gospel Faith and the Modern Consciousness}}
\fancyhead[LO]{\textit{\leftmark}}

\setstretch{1.3}

\title{\textbf{Gospel Faith and the Modern Consciousness}\\[6pt]
  \large Lectures on the Life of Jesus\\[4pt]
  \large First Part}
\author{R.\ Nielsen}
\date{Copenhagen: C.\ A.\ Reitzel, 1849}

% ---- original-page markers (paginate.py) ----
\usepackage{marginnote}
\newcommand{\opage}[1]{\marginnote{\footnotesize\textit{[#1]}}}
\newcommand{\apage}[1]{\marginnote{\footnotesize\textit{[#1]}}}
\begin{document}
\maketitle
\thispagestyle{empty}
\newpage

\tableofcontents
\newpage

% --- p. [I] ---   (PDF 6)
\begin{center}
  {\LARGE Gospel Faith}\\[10pt]
  and\\[10pt]
  {\LARGE the Modern Consciousness.}\\[12pt]
  \rule{0.28\textwidth}{0.4pt}\\[14pt]
  {\large Lectures on the Life of Jesus}\\[10pt]
  by\\[10pt]
  {\large \textit{Lic.\ theol.}\ R.\ Nielsen,}\\[2pt]
  {\footnotesize Professor of Philosophy.}\\[12pt]
  \rule{0.28\textwidth}{0.4pt}\\[16pt]
  {\large First Part.}\\[16pt]
  \rule{0.62\textwidth}{0.4pt}\\[16pt]
  {\large Copenhagen.}\\[4pt]
  On commission with University Bookseller C.\ A.\ Reitzel.\\[4pt]
  Printed by E.\ C.\ Løser.\\[6pt]
  {\large 1849.}
\end{center}

\newpage

% --- p. [III] ---   (PDF 8)
\begin{center}
  {\large To Her Majesty}\\[18pt]
  {\LARGE Queen Dowager Caroline Amalia,}\\[20pt]
  whose comfort is a faith\\
  that believes the Gospel,\\[20pt]
  is dedicated\\[14pt]
  \emph{this book}\\[8pt]
  {\footnotesize on}\\[10pt]
  {\large Gospel Faith}\\[22pt]
  \emph{most humbly}\\[8pt]
  by\\[8pt]
  {\large the Author.}
\end{center}

\newpage

\chapter*{Preface}
\addcontentsline{toc}{chapter}{Preface}
\label{ch:fortale}

% --- p. V ---
Whether and to what extent the reasoning criticism of the Bible can resolve
to declare our four canonical Gospels genuine or spurious, consistent or
inconsistent, reliable or unreliable, is a question for the Schools, whose
answer had surely best be left to the learned.

It is a fact, however, that there do now exist these four Gospels, in which
Jesus of Nazareth is presented to us as the Son of God and of Man, as the
Messiah and Savior of the world.

That the authors of these Gospels must themselves have believed in the
Messiah they describe, there is surely no reason to doubt either; but the
question is how it can now be conceived that any human being who has truly
come to share in the goods of culture and become familiar with the higher
cultivation of the present age will still be able to regard it as the
highest thing to share faith with the four Evangelists.
% --- p. VI ---
The task to whose solution the present work presumably makes some
contribution can then be briefly set forth in the vital question: \emph{What
is it to believe in the Christ of the Gospels?}

A threefold answer is possible: \emph{It is to appropriate the revealed
truth}; --- \emph{It is to let oneself be beguiled by a pious fancy}; ---
\emph{It is to console oneself with a curious mixture of truth and fancy.}

The first answer befits --- \emph{the faith of the Gospel}; the second ---
\emph{the modern consciousness}; the third points to a mediating view that
would prefer to keep to the Middle Way.

The faith of the Gospel lives in --- \emph{the Believer}; the modern
consciousness in --- \emph{the free thinker, the Free Thinker}; the doctrine
of the Middle Way has considerable adherents among the theologians of the
School.

The main idea, to let the spokesmen for the views here designated agree to
go through the four Gospels together, each speaking in his own character,
was already made the basis of the lectures on the Life of Jesus that the
author delivered in the University building in the winter of 1848, before a
respectable
% --- p. VII ---
number of listeners of both sexes; but a fuller presentation could be
communicated only in writing.

In a more developed form, then, these lectures are commended to the
attention of the reading public with the wish, perhaps a rather bold one,
that the written word may now find the same kind reception that previously
fell to the lot of the spoken.

In order, however, not to arouse any untimely expectation, or to occasion
any error that would waste reviewers' time, the author will herewith take
the liberty of drawing the attention of the wider public to the fact that
the eager dispute conducted in these proceedings for and against the
credibility of the Evangelists does not really have quite such an outcome
that one can say with certainty whether it is ``the Free Thinker'' or ``the
Believer'' who carries off the victory in the end.

On the other hand, the attentive reader who is interested in the fundamental
question to such a degree that he provisionally makes do with the
presentation and, if not overlooking its defects, at least excuses them,
will easily convince himself that the dispute is not without result. For
this too is surely a result, that the opposing views become clear, that one
sees what dwells in the disputants. The dispute is then like a
% --- p. VIII ---
trial; the trial takes its course; the case is submitted for judgment; the
reader's conscience is the jury that pronounces the verdict.

In order, however, not to conceal his own personal conviction either, and
thereby turn the whole into a kind of philosophical thought-experiment, the
author has in several places partly spoken out quite plainly, partly also
taken the floor on behalf of the disputants; but, as it became more and more
evident to him that neither of the parties could possibly refute the other,
he has not ventured to undertake any final refutation either.

The faith of the Gospel cannot be refuted; it has its own supernatural
principle: the Believer endures to the end.

The modern consciousness cannot be refuted; it too has its own principle,
and does not fear its consequences: the free thinker does not tire, his
powers grow as the dispute proceeds.

The only one of the ideal powers whose refutation could be contemplated here
would then have to be the doctrine of the Middle Way, and all the more so as
it here certainly finds itself in a rather tight position, midway between
the old simple faith of the Gospel and the
% --- p. IX ---
negative consciousness of modern knowledge. But this too is impossible: the
doctrine of the Middle Way really cannot be refuted. True, it is not very
rigid in its principles, but in return it is all the more many-sided; true,
it does not have the immediate superiority of strength, but it has all the
more tenacity; true, it cannot kindle any notable enthusiasm, but it can
work with all the greater composure and tactical calculation.

How could the doctrine of the Middle Way possibly be refuted, when one
considers that not only is moderation a cardinal virtue, but mediocrity
itself is a power in the world? How could the doctrine of the Middle Way be
refuted, when one considers that it can, after all, appeal partly to the old
and partly to the new, partly to faith and partly to knowledge, partly to
criticism and partly to dogmatics? How could the doctrine of the Middle Way
be refuted, when one considers that whoever took upon himself so enormous a
labor might in the end come to refute virtually the whole of the present
theology of the School? --- \textit{Favete lingvis}: no presumption!

But although the author knows within himself that he neither can nor will
refute the theology of the
% --- p. X ---
School, he is naturally by no means thereby assured that the theology of the
School, should it once find it worth the trouble, could not refute him.
Still, it is not said that it will, even if it easily could.

Yet whether the fundamental idea set forth in this work should occasion a
quiet disapproval, or a loud refutation, or both, or --- nothing at all; the
undersigned will console himself that he has honestly and sincerely labored
toward the answer to his vital question, without in any way wishing to
offend either the deeper learning or the higher science.

\medskip
\noindent\hspace*{2em}{\small Copenhagen, May 1, 1849.}

\medskip
\noindent\hfill\textbf{R.\ Nielsen.}\hspace*{3em}

\chapter*{Contents}
\addcontentsline{toc}{chapter}{Contents}
\label{ch:indhold}

\begin{center}
  {\large First Division.}\\[6pt]
  \emph{The Fore-Gospel of the New Testament:}\\[4pt]
  \textbf{\large The Coming of Jesus into the World.}\\[6pt]
  {\footnotesize I came forth from the Father and am come into the world (Joh. 16, 28).}
\end{center}

\begingroup
\setlength{\parindent}{0pt}
\makebox[2.6em][l]{I.}The Angel Gabriel and the Priest Zacharias\dotfill Page 1--42\par
\makebox[2.6em][l]{II.}The Angel Gabriel and the Virgin Mary\dotfill --- 42--80\par
\makebox[2.6em][l]{III.}The Holy Family\dotfill --- 81--183\par
\hspace*{2.6em}\makebox[1.8em][l]{1.}Mary's Visit to Elisabeth --- The Birth of John\\
\hspace*{4.4em}the Baptist\dotfill ---100--128\par
\hspace*{2.6em}\makebox[1.8em][l]{2.}The Birth of Jesus --- The Shepherds in Bethlehem\dotfill ---128--166\par
\hspace*{2.6em}\makebox[1.8em][l]{3.}The Gospel Forewarnings\dotfill ---166--183\par
\endgroup

\begin{center}
  {\large Second Division.}\\[6pt]
  \emph{The Sacred History:}\\[4pt]
  \textbf{\large The Life of Jesus in the World.}\\[6pt]
  {\footnotesize While I am in the world, I am the light of the world (Joh. 9, 5).}\\[8pt]
  \textit{A.}\ The Redeemer.
\end{center}

\begingroup
\setlength{\parindent}{0pt}
\makebox[2.6em][l]{I.}The Forerunner\dotfill Page 184--201\par
\makebox[2.6em][l]{II.}Christ among His Own\dotfill --- 201--227\par
\makebox[2.6em][l]{III.}Christ Is a Sign That Is Spoken Against\dotfill --- 228--289\par
\makebox[2.6em][l]{IV.}Christ Reveals the Thoughts of Hearts\dotfill --- 289--334\par
\makebox[2.6em][l]{V.}Christ and the Spirit of Darkness\dotfill --- 334--393\par
\makebox[2.6em][l]{VI.}Christ Is Mighty in Word and Deed\dotfill --- 393--458\par
\makebox[2.6em][l]{VII.}Jesus Christ Is the Savior of the World\dotfill --- 458--530\par
\endgroup

\chapter*{Errata}
\addcontentsline{toc}{chapter}{Errata}
\label{ch:rettelser}

\begingroup
\setlength{\parindent}{0pt}
\begin{tabular}{@{}l@{~}r@{~~}l@{~}r@{~~}l@{}}
Page & 97  & line & 8  & from top: Undergjerningens read: Undergjerningernes\\
---  & 124 & ---  & 11 & f. t.: at han tør r.: at han her før\\
---  & 165 & ---  & 3  & from bottom: Loø r.: Lov\\
---  & 193 & ---  & 11 & f. t.: Himmerige r.: Himmeriges Rige\\
---  & 322 & ---  & 10 & f. t.: \textit{le peuble} r.: \textit{le peuple}\\
---  & 346 & ---  & 9  & f. b.: hvor i henlægge r.: hvor I henlægge\\
---  & 359 & ---  & 1  & f. t.: Taabeligheden r.: Taaligheden\\
---  & 368 & ---  & 13 & f. t.: fortjene r.: fortjener\\
---  & 406 & ---  & 9  & f. t.: fordrage r.: foredrage\\
\end{tabular}
\endgroup

\chapter*{First Division\\[4pt]
  \large The Fore-Gospel of the New Testament: The Coming of Jesus into the World}
\addcontentsline{toc}{chapter}{First Division. The Fore-Gospel of the New Testament: The Coming of Jesus into the World}
\markboth{First Division}{First Division}
\label{ch:afsnit1}

% --- p. 1 ---
\opage{1} %
\begin{center}
  {\footnotesize I came forth from the Father and am come into the world (Joh.\ 16, 28).}\\[10pt]
  \rule{0.22\textwidth}{0.4pt}
\end{center}

\section*{I.\\[4pt] The Angel Gabriel and the Priest Zacharias.}
\addcontentsline{toc}{section}{I. The Angel Gabriel and the Priest Zacharias}

\begin{center}
  \rule{0.14\textwidth}{0.4pt}
\end{center}

\begingroup\large
Behold, says Jehovah to the people of Israel, I send you Elias the prophet
before the great and dreadful day of the Lord come. And he shall turn the
heart of the fathers to the children, and the heart of the children to their
fathers, lest I come and smite the earth with a curse.
\par\endgroup

Thus run the last words of the last prophet in the Old Testament
(Malach.\ 4, 5--6), and with the fulfillment of this prophecy the Gospel in
the New Testament begins. (Luk.\ 1, 5--25).

\begingroup\large
In the days of King Herod there lived in Judaea a priest of the course of
Abia, named Zacharias, whose wife
% --- p. 2 ---
\opage{2}Elisabeth was of the daughters of Aaron. They were both righteous
before God, and walked blameless in all the commandments and ordinances of
the Lord. But they had no child; for Elisabeth was barren, and they were
both well stricken in years. But it came to pass, while Zacharias executed
the priest's office before God in the order of his course, and went into the
holy place to kindle the incense offering, that the angel Gabriel appeared
to him. And he saw the angel standing on the right side of the altar of
incense, and he was terrified, and fear fell upon him. But the angel said to
him: Fear not, Zacharias! for thy prayer is heard, and thy wife Elisabeth
shall bear thee a son, and thou shalt call his name John. He shall go forth
in the spirit and power of Elias, to turn the hearts of the fathers to the
children, and the disobedient to the understanding of the just, to make
ready for the Lord a people well prepared. And Zacharias said to the angel:
Whereby shall I know this? for I am old, and my wife also is well stricken
in years. And the angel answered and said to him: I am Gabriel, that stand
before God, and I am sent to speak to thee, and to proclaim this to thee for
joy. --- But because Zacharias doubted, the angel struck him dumb, and in
this punishment gave him a sign that the promise should be fulfilled. And
all the people assembled, who were in the court praying at the hour of
incense, waited for Zacharias, and marveled that he tarried so long; and
when at last he came out, the assembly perceived that he had seen a vision
in the sanctuary. And when the days of his ministration were ended,
Zacharias departed to his own house, and that was fulfilled which the angel
had proclaimed to him.
\par\endgroup

This, in sum, is the content of the narrative
% --- p. 3 ---
\opage{3}with which the Gospel of the Savior of the world is introduced. The
preparation ends, that the consummation may come; that which is in part
ceases, that the perfect may take its beginning. An extraordinary, a
wonderful beginning! The extraordinary shows itself not merely in the pious
priest's receiving a portent that the fullness of time is near; it is of
decisive significance that this portent comes from heaven, that it is God's
angel who brings it to him, that it is the prophet's prophecy that sounds
and resounds in the glad tidings. The wonderful will not let itself be
confined to the annunciation's actually coming to fulfillment; we are
expressly taught that this fulfillment is itself a wonder, a divine work, by
which the work of redemption is prepared. The Evangelist proclaims to us
that it is the God of Israel who now once more visits his people, the God of
the prophets who fulfills his promises, the God of salvation who remembers
his holy covenant according to the oath he swore to Abraham.

If only we might now assume that the fulfillment of the prophecy really took
place as the narrative says, if only we could be assured that this God of
Israel is the only true God, that this God of the prophets really is able to
fulfill his promises, and that this God of salvation has not, in the
forgetfulness of the ages, forgotten himself and the oath he swore to
Abraham; yes, were it only granted us --- to believe: then we could still
seek refuge in the sanctuary of revelation, could still, as at the beginning
of the Gospel, hear the angel's voice: ``I am sent to speak to thee, and to
proclaim this to thee for joy.''

``But,'' say the enlightened of the present day, ``all this has no true
reality at all. The story of the annunciation dissolves into sheer illusion
when the narrative is seen through the spectacles of criticism.'' Might they
perhaps be right? Has critical science now really come so far that it can,
without further ado, begin its
% --- p. 4 ---
\opage{4}dissolution of the Gospel history with the assertion that
revelation and wonders belong among impossible things? Such an assertion
would surely stand in great need of a proof comprehensible to simple people.
But such a proof can hardly be carried out unless the faith of the Gospel
first comes to speak freely and unhindered, and that not off to one side as
a partisan matter, but in the very course of the proceedings as a party to
the case, according to that golden rule: \textit{audiatur et altera pars.}
To attack the narrative before us with an unproved assertion is surely not
only unseemly, unjust and unscientific, but seems even to be so uncritical
that a particular criticism of the narrative itself thereby becomes wholly
superfluous. What can be made of the simple story of Zacharias and
Elisabeth once the angel is sent away, whether it can then perhaps be
explained as a superstitious distortion of an entirely natural event, or
better still explained away as a myth, does not concern us at all. If the
angelic manifestation is lost, then the story is lost to us. The critical
inquiry therefore turns exclusively on the main question whether the account
before us, both as a whole and in its individual features, is in complete
agreement with the fundamental character of revelation, or not. Only when
the task is posed in this way can a scientific discussion be opened between
doubt and faith.

Let the critical doubter, then, come forward with his objections; these
objections will gain in scientific weight precisely in proportion as the
doubter can otherwise relate sympathetically to what he calls the
fundamental idea of Christianity.

``Along the merely rational way,'' says the doubter, ``it is certainly
impossible to attain any decisive religious conviction. But if, on the
other hand, it is already hard enough to have to believe in the reality of
wonder in general, the
% --- p. 5 ---
\opage{5}difficulty becomes far more keenly felt when one is to resolve to
accept this or that individual, determinate wonder in the form of a
historical event. Every time I take up the N.\ T. and prepare myself to go
through the Gospel history, I impress upon myself again and again the
principle that, as surely as there is a real revelation, one must begin from
a determinate, supernatural starting point, and that divine communications,
extraordinary phenomena and the like must therefore be in their proper place
here. But when I then set about familiarizing myself with the very first
account in Luke, why then --- my understanding stands still. The more I
exert myself toward an appropriation, the harder I am repelled; the longer I
stare at the picture, the more clearly I see that the Evangelist, who
explains everything for Theophilus, explains nothing for me. In the first
place, I find no sufficient reason why the Baptist's birth should require so
extraordinary an arrangement on the part of heaven itself. It must certainly
be granted that John the Baptist, as preacher of repentance, as Christ's
Forerunner, as prophetic mediator between the old and the new covenant, must
have been a man of eminent spiritual gifts. But in order to be endowed with
divine power of spirit, is it then absolutely necessary that one's birth be
conditioned by a miracle? On this assumption Abraham and Moses and all the
prophets, indeed the apostles themselves, would have to be able to show that
their birth was owed to a wonder-work. But even if one could overlook this
difficulty, why should John's birth be announced beforehand to the parents
by so extraordinary a messenger? What entitles Zacharias and Elisabeth to so
extraordinary a favor? This awkwardness strikes the eye still more forcibly
when we consider the angelic manifestation described here a little more
closely. Only when the human spirit is laboring in the very highest
% --- p. 6 ---
\opage{6} tension can supernatural visions be thought of: and only when this
tension has been brought about by a great, decisive world crisis, only when
the knot of entanglements is drawn so tight that none but God himself can
untie it, only then may we presumably take the immediate intervention of the
Governor of the world in the course of things for more than a pious
illusion. When Jerusalem was besieged by the Romans, when the Jewish people,
driven to the uttermost by despair, rage and religious fanaticism, fought its
last fight; when the outrage against Jehovah's sanctuary was so revolting
that the fathers must have turned in their graves, and the misery so
boundless that a mother devoured her own child, then just such a grand
turning point, one that challenged heaven, seemed to have arrived. Had the
Evangelist depicted a similar scene in our narrative, a scene of sacrifice
in the temple already kindled by the flames, why then I could understand
that an old priest saw visions in the sanctuary, and the people in the court
believed in the manifestation. But this manifestation lacks the point of the
moment, the involuntary inspiration of the idea. This meeting between
Zacharias and the angel is kept in so dry, one might almost say so
bourgeois a style, that the miraculous arrangement almost looks like a
little private surprise in the midst of everyday life. The priest has, to be
sure, entered the holy place; he is also somewhat frightened the moment he
sees the vision; but he nevertheless collects himself at once, as soon as the
angel begins to speak. And what sort of impression, then, does this angelic
greeting, with the trusty declarations that follow it, produce? `Don't be
afraid, Zacharias! I am neither an evil spirit nor a ghost; my name is
Gabriel. I bring from heaven the message that God Jehovah has now at last
heard your prayer, and resolved to grant you a son, who shall be called John
and shall work like the prophet Elias.' And then the cold-bloodedness with
which Zacharias,
% --- p. 7 ---
\opage{7} once the first fright has settled, considers the difficulty of the
matter on his own and his wife's behalf, and to be on the safe side asks for
a further sign, as if the angel's express declaration were not enough for
him. Holy God, should such a thing be possible? Could there be found
anywhere outside the preserves of the imagination such an intercourse
between the heavenly and the earthly, so familiar, oh, what shall I call it,
so comradely a relation between the majesty of eternity and the
insignificance of finite man? No, no! it is impossible! I cannot, I dare not
believe this narrative; for I cannot and dare not console myself with a
self-deception. So it is this little story, then, that must rise up against
me every time and set itself like a barrier between me and the Gospel; for I
know it: in rejecting it, I reject --- `God's revelation.' So it is the
annunciation of this child's birth that must shut me out of the paradise of
childlikeness; for I know it: in breaking with its angel, I have at the same
stroke passed sentence on the sacred history.''

It is futile to try to meet these and similar objections with general
scientific reasons, since what is at issue here is not a relative difference
between a higher and a lower degree of knowledge, but the absolute
difference between faith and unbelief. It is only a modern hypocrisy when
many nowadays pretend that they would gladly believe in the wonders of the
Gospel, if only their higher enlightenment, their freer cultivation, their
pure scientific spirit would permit them. He who does not know within
himself that the opposition between faith and non-faith is so far from being
dependent on any science whatever that this opposition itself, rather,
grounds, permeates and governs our essential knowledge, may perhaps well
boast of much enlightenment and many-sided cultivation, but is still not so
thoroughly cultivated that
% --- p. 8 ---
\opage{8} he truly sees what it is that is really at issue here. It is
certainly impossible for the doubter to grasp the convincing power of faith;
but it is by no means impossible for the believer to go through the
reflections of doubt. With this self-collection we now return once more to
our Gospel, give ourselves over to contemplation, and strive to interpret
the letter of the narrative in the spirit of faith.

How mysterious, and yet how manifest, is this beginning! History lingers,
and thought rests as at a turning point. The spirits of the age have settled
to rest; the outward world has grown so still: it is in the days of Emperor
Augustus, it is in the days of King Herod. The temple still stands on Mount
Zion: the people of Israel still have a home in the land of the patriarchs.
No outward enemy dares harass Judaea, for the crafty Herod is the Emperor's
friend; no cry of revolt dares sound in the streets of Jerusalem, for the
cruel Herod --- punishes with death. The day's work in the houses, the
sacrificial service in the temple, all goes its quiet, calm course according
to custom and usage: no one expects, no one demands the extraordinary.
Zacharias too enters the holy place according to the usage of the
priesthood; it had now fallen to him that day to ``burn incense'' in the
temple. The multitude of the people is gathered in the court, according to
custom and usage; it was fitting, after all, to pray at ``the hour of
incense.'' No one expects, no one demands the wondrous. Then comes the
manifestation. Sudden as lightning the angel is seen in the sanctuary, and
the vision strikes Zacharias with terror, and the priest tarries in the holy
place, and the congregation must wait --- against custom and usage. And
those waiting wonder why he tarries in there; but no one suspects, no one
demands the supernatural. Behold, then the old priest steps out into the
court, dumb and enigmatic, and with so strange a look; only now does the
assembly perceive that something
% --- p. 9 ---
\opage{9}wondrous must have befallen him, that he must have seen ``a vision in
the temple.''

Such is the beginning of the Gospel. Let worldly wisdom take offense at this
beginning; it is only natural that one demands another starting point when
one has another God. A deity whose essence is immediately one with the
forces of nature and the spirit of history cannot begin its manifestation so
decisively and yet so calmly; for with violent pains, amid strong
convulsions, must the world-spirit give birth to its new ideas. The
ambiguous world-god makes no decision under an Emperor Augustus, and works no wonders under a King Herod. Fermenting in dark possibilities, he must
first await some desperate catastrophe or other, such as the destruction of
Jerusalem or the ruin of Rome. But the God of the Gospel, with his eternal
counsels, has no need at all of such finite awakenings. Or should the God of
redemption forget man's spiritual need because the race dozes, and the
nations live on in security, saying: ``Here is peace, and no danger''? Or
should the God of eternity and of providence perhaps be unable to see 70
years ahead in time, and to complete the cycle of revelation before the
sanctuary of the old covenant falls in ruins and the chosen people is
scattered over the earth?

It is entirely natural that the non-believer must take offense at this
beginning; for a world-god cannot begin with such clear, such ideal facts.
The world-god must begin with obscure stirrings, immature attempts,
ambiguous starts. Only gradually, as the new work of the spirit clarifies
itself out of its chaos and displays its grand figures in the more definite
outlines of execution, does poetizing fantasy enter the service of memory
and glorify the raw beginning with legends and myths. But the God of
revelation begins with the wonder, and the wonder, which itself
% --- p. 10 ---
\opage{10} is the unity of the ideal and the actual as holy fact, wants no
fantastic embellishment of uncertain occurrences, but a simple and truthful
account of the supernatural beginning.

With contempt the dissolving criticism looks down on the pious ones who,
according to our Gospel, are favored with the angel's message. With imagined
superiority one asks what entitles Zacharias and Elisabeth to such a favor.
In the eyes of the world this objection is as apt as it is bitter; for if
the persons concerned really are as lowly and insignificant as criticism
assumes, then things go on in the Gospel as in the fairy tale and the
legend, where a god who takes the notion to surprise by a miracle is no
respecter of persons even in this sense, that he takes up with the first
comer. But is it not highly uncritical to confuse the brilliant greatness
of genius and energy, which can astonish a world, with that quiet, humble
greatness of obedience and self-denial which alone pleases God? In the eyes
of world history Zacharias and Elisabeth are certainly insignificant; it
knows nothing of their deeds, it scarcely even heeds their names. A
psychological observer, who otherwise with unwearied diligence tracks down
things old and new so that a cultivated reading public may always be
supplied with a suitable selection of interesting silhouettes for pleasant
entertainment, will no doubt pass our old couple by with indifference; for
Zacharias's house is a quiet house, without piquant scenes and
nerve-shattering episodes. If, on the other hand, we make do with the
strength of soul that, without exactly causing any noisy stir outwardly, is
exercised and tested in all inwardness before the All-knowing, then the
Evangelist has not left Zacharias and his wife without witness. He reports
of them, after all, that they were both righteous before God, and walked
blameless in all the commandments and
% --- p. 11 ---
\opage{11} ordinances of the Lord. True enough, this testimony, set down in
few words, looks rather like a passing remark, a fleeting hint, than an
enthusiastic eulogy. But this comes partly from the fact that it cannot at
all be thought that Zacharias and Elisabeth should by their good works have
merited a manifestation that is owed solely to divine grace; partly also
from the fact that language is able to express in few and pithy words what
requires a whole life's sustained exertion when it is to be expressed in
deed. It is, speaking in earnest, not so easy to walk blameless in all the
commandments and ordinances of the Lord. He who was capable of this might
well be called great in Israel. And yet one person might perhaps seem to
himself great in obedience without therefore being tested in much
self-denial (for the God of Israel blesses and rewards the obedient); and
another might perhaps seem to himself strong in faith without yet being
tried in much patience, since his fear of God never lacked its temporal
reward. But when the expected reward fails to come, when the gift that is
first to give the present life its right and full significance, the gift
that the God of the patriarchs has promised to all who find favor in his
eyes, is denied the believer; when day by day it becomes more and more
improbable that this gift will ever fall to the lot of the one who hopes,
then strength of soul is needed to preserve faith and hope. It is this
strength of soul that we honor and admire in Zacharias and Elisabeth. ---
He who disdains to enter into the religious and civil conditions of the
Hebrew may, in free and self-complacent cultivation, well find fault that a
pair of old spouses should be so anxious for an heir that at last they hit
upon the thought that God himself might perhaps come to their aid by a
miracle. But it is nevertheless true that the idea of the continuation of
family life in numerous and happy descendants was so grown together with the
patriarchal relation to God that
% --- p. 12 ---
\opage{12} the Israelite who, without propagating his name to posterity, sank
down into the shadow-realm of oblivion could be regarded as rejected by
Jehovah and cut off from his people.

So we can well imagine how Zacharias and Elisabeth must have made their
covenant in the childlike faith and the youthful hope that the God of
Abraham, Isaac and Jacob would grant his blessing to them too, how in this
faith they must have promised each other to walk blameless before God in all
the commandments and ordinances of the Lord. But --- youth vanished, and the
glad expectation was not fulfilled. Neighbors and kinsfolk received the
blessing; but Zacharias and Elisabeth had to stand aside. If want then moved
in with secret sorrow to dwell with them, if at times they had to look at
each other in sadness, so poor, alas, so forsaken, because Jehovah would not
look upon them, would not visit them: yet humble faith was greater than
quiet grief, hope stronger than fear, and the days of expectation not yet
over. The promise stands in letters of gold: \emph{The righteous shall
flourish like a palm tree, he shall grow like a cedar in Lebanon.} ``Oh, that
one might be found righteous before God, and then --- have one's share in the
promise! Oh, that one might walk blameless in all the commandments and
ordinances of the Lord, and then find favor in his eyes, and --- be adorned
with the blessing!'' --- Thus our pious pair comforted each other with the
word of the promise, and reckoned the multitude of the days by the
fulfilling of the law and the number of the years by the service of
righteousness, that they might both be ``righteous before God and walk
blameless in all his commandments and ordinances.'' But the days vanished,
and the years went by, and --- the expectation was not fulfilled. The palm
tree flourished, the cedar grew, the vine bore its fruit; but --- the
expectation was not fulfilled. Pious neighbors, happy kinsfolk lifted up
their voices and praised Jehovah's mercy from
% --- p. 13 ---
\opage{13} generation to generation; only Zacharias and Elisabeth had to be
silent and wait. Pious neighbors, happy kinsfolk saw the blessing of the
fathers upon children and children's children when the feast day gathered
them to the solemnity of sacrifice and the joy of the banquet; but Zacharias
and his wife remained alone with their offering, alone with their
thanksgiving, alone in their house. Old age came and found Zacharias
childless, and Elisabeth got a byname as one cast off by the Lord, and she
felt her reproach among men; but --- they remained alone. So it was, then, a
trial in which they were to be proved, and Zacharias and Elisabeth
understood that it was a trial; for they did not envy their happy kinsfolk,
and they did not murmur against Jehovah, and they did not lose hope in
despondency; but --- they renewed the covenant of their youth, saying in
their hearts: ``We will still walk blameless in all the commandments and
ordinances of the Lord; for this is surely a trial. Why we, we alone, should
thus stand aside among our kin and bear the contempt of men, as though we
were despised by the God whom we would serve in all righteousness, ah, as
though we were rejected by him who yet rejects only the ungodly! Why we, we
alone, should go down into the grave with gray hairs like our fathers, full
of days like our fathers, and yet not receive our fathers' blessing, but
leave our dwelling desolate, as though our life had been but a shadow, yes,
give up our names to oblivion, ah, as though we were forgotten by him who
yet can never forget his pious ones! Why, O God, why? But let us be silent
and wait and be still before God; for who hath directed the Spirit of
Jehovah, and what counsellor hath taught him! Let us rather take our own
thoughts to task, that we may not take the Lord to task; for who knows what
the Almighty will yet do for us? Yes, great and wondrous things the Lord
did of old for his chosen ones, as for Abraham and Sarah, Manoah and his
wife, Elkanah and Hannah;
% --- p. 14 ---
\opage{14} yes, the Lord is a God of knowledge; should his works not be
right!'' This, then, is the turning point of the trial. They who had so long
believed with hope now at last believed against hope, yet hoping in the God
who quickens the dead and calls those things which are not as though they
were. Behold, this is Abraham's faith, Abraham's hope! With this faith and
this hope Zacharias meets the angel in the holy place.

If, however, a wordy presentation set out to prove to modern criticism how
worthy and well prepared Zacharias might thus have been to receive the
angel's message, such an attempt would not only be rather fruitless in and
of itself, but would even betray a questionable misunderstanding of the
Gospel's purpose. If such a thing could be proved, who was nearer to
furnishing the proof than the Evangelist himself? But it is folly to demand
proof for that which by its nature admits of no proof. However urgent the
presentation might be, it would not help. Though the most secret powers of
language were stirred, yes, though one spoke with the tongues of angels in
praise of Zacharias and Elisabeth, what would it avail, when negative
criticism will hear only its own language, adopted once and for all, and the
modern consciousness knows within itself that it cannot learn the tongue of
angels? Therefore the Evangelist has not engaged in any attempt in this
direction either; but, well knowing that the Gospel is only for the simple
in faith, he keeps it brief and says of our old couple: ``But they were both
righteous before God, and walked blameless in all the commandments and
ordinances of the Lord.'' More is not needed; so then --- to the angelic
manifestation.

As criticism breaks into the sanctuary and takes aim at Gabriel, a
remarkable dispute arises. For criticism has found out that, since modern
science does not acknowledge angels,
% --- p. 15 ---
\opage{15} these do not exist at all, and that, if angels do not exist, then an
angel with the name Gabriel does not exist either; further, criticism has
found out that, if angels do not exist, then they cannot manifest themselves
either, and that, if angels cannot manifest themselves, then an angel named
Gabriel cannot have manifested himself to Zacharias either. Therefore it is
now also said, in the language of modern science, that Gabriel cannot stand
before the tribunal of criticism, and that one has for that reason found
oneself justified in assigning him a place among the mythical figures. So
this is what criticism has now found out, and it is all quite correct, and
there is nothing whatever to object to in the critical chain of reasoning
except this one thing, that its presupposition is plucked out of thin air.
An actual proof of the impossibility of the existence of angels modern
science will doubtless refrain from furnishing, and when faith merely knows
the possibility to be safe, it demands no more of the whole of worldly
wisdom. Against the assertion that it is Gabriel who cannot stand before the
tribunal of criticism, we can then, on faith's behalf, set up the assertion
that it is really, perhaps, criticism, with all its critics, that cannot hold
its own against Gabriel, and therefore hastens, the sooner the better, to
consign him --- to the clutches of myth. Science naturally does nothing
without deliberation. If, then, science has after a sharp interrogation
banished the angel from the regions of actuality as an ambiguous and
suspicious person, one may also rest assured that science has had its good
and sufficient reasons for doing so. Everything has been done on the part of
the interpreters to bring this strange being to reason and to explain the
story in a reasonable and natural way; but all the efforts of art have so far
been in vain. The matter is this: if the angel is to be something, he must of
course be either bodily or bodiless. But now the Evangelist has presented
him in such a way that he is both the one and the other, and yet
% --- p. 16 ---
\opage{16} neither the one nor the other; and what rational person can then
make head or tail of him? If the natural art of interpretation wants to hold
him fast as an outward figure, and let him answer to his name Gabriel, i.e.\
\emph{man of God}, then this man of God must naturally be some highly
peculiar man or other who has sneaked into the holy place to surprise the
priest. But this notion falls through at once. The sanctuary was closed,
after all, and one then looks in vain for a back door through which the man
could have slipped in, and through which he could have slipped out again.
If it will not succeed to hold the angel fast in bodily externality, he
must naturally be conceived as a bodiless spirit. But if he is a bodiless
spirit, then he cannot be seen with bodily eyes either: ``so he must be a
figment of the priest's own overheated imagination.'' But against this
conjecture the narrative protests most vigorously. The interpreter who would
grasp the vision grasps at air, and behold! there the angel stands,
\emph{there}! ``on the right side of the altar of incense,'' and looks at
the terrified Zacharias, and speaks so audibly, so authoritatively, so
fully, so calmly, clearly and soberly, that no man of God can speak more
distinctly, and no ecstatic conjure himself into such a vision. Were it now
quite settled and proved once and for all that modern science exhausts the
truth, and that nothing is permitted to exist except that whose mode of
existence science comprehends: why then, indeed, it would be over with the
angel, and over with the wonder, and over with faith. But suppose the
reverse is the case; suppose that the highest insight to which all our
comprehending leads is the insight that the incomprehensible is laid as the
ground of the comprehensible, that the wonder of revelation neither can nor
shall be comprehended; then faith is proved right, and with that the angel
too is proved right; for the angel is the spokesman of the wonder:
{\large I am Gabriel, that stand before God, and am sent to
% --- p. 17 ---
\opage{17} speak to thee, and to proclaim this to thee for joy.} So it will
then be best to let Gabriel be what he is, a supernatural being, an angel of
annunciation in the sanctuary whose veil criticism is unable to lift.

But even a believing eye is not at every moment equally prepared to meet the
angel of annunciation. The calm, historical narrative, which here
communicates the most wondrous thing of all in the form of immediate fact,
and without any circumlocution inserts a supernatural link into the series
of natural events, has something infinitely repellent for awakening
thought. Had the Evangelist at least, before expressly mentioning the
vision, sent ahead an introductory remark, for example that Zacharias, when
he entered the holy place, was in the spirit, or: saw heaven open, --- or
granted us some similar hint; then it would perhaps not be so difficult for
the interpreter to find an explanation. To shed a little more light on this
point, we shall fetch a couple of examples from the prophets. Isaiah
describes his calling (Jes.\ 6, 1--8); what a breakthrough, what a vision!
Upon his high and lofty throne sits Jehovah; the skirts of his robe fill the
temple. The seraphim who stand above him cover their faces, and cry one to
another: ``Holy, holy, holy is the Lord Zebaoth, the whole earth is full of
his glory.'' The doorposts tremble at the voice of him who cries, and the
house of God is filled with smoke. Terrified, the seer cries out: ``Woe is
me, I must perish; for I am a man of unclean lips, and I dwell in the midst
of a people that have unclean lips, and yet mine eyes have seen the King,
the Lord Zebaoth.'' Then the seraph fetches a live coal from the altar, and
touches the prophet's lips with it, ``that his iniquity may depart from him,
and his sin be purged.'' Truly, this consecration works on us with the whole
immediate force of revelation; it is
% --- p. 18 ---
\opage{18} as if we felt the holy fire that glows through the prophet's lip.
But the vision is at the same time presented to us as a pure vision of the
spirit, and just because it is a vision of the spirit, the fact admits of an
explanation. The figure of Jehovah transforms itself into an image, and the
image takes on an allegorical tinge; the wonder is absorbed into the idea,
and the idea satisfies reflection. This holds in still higher degree when
the apocalyptic images turn on the crises of world history. When a Daniel,
in the visions of the night, sees ``the Ancient of Days'' surrounded by the
many thousands of angels, when he sees the judgment set and the books
opened, and one like a Son of Man coming in the clouds of heaven to ``the
Ancient of Days,'' and receiving from him the kingdom and the power and the
glory: then these nocturnal visions have so spiritual, so airy a form that
the weight of finitude vanishes, and the images pass by in great but
fleeting outlines. It is the spirits of history that rush about the
prophet's head; it is the Ruler of all who summons the nations before his
throne, it is ``the Son of Man'' who holds the Day of Judgment. Indeed, so
striking is the resemblance between these apocalyptic images and the
ingenious idea-forms of the philosophy of history that a believing observer
very easily confuses the supernatural signs from the providence of
revelation with the natural determinations in the general concept of the
governance of the world, and ends as a philosophical spectator. Reason has
nothing against man's standing in a general relation to God through the
idea of the world; but this does not satisfy faith. Faith is rooted in the
conviction that the saving God enters into an immediate personal relation
with fallen man by means of the wonder. If it is now settled that the
wonder is the right and proper mark of revelation, then it is also settled
that the manifestation which lays down the whole content of the
determinations of spirit in the immediate fact of the wonder is to be
regarded as the norm and standard for
% --- p. 19 ---
\opage{19} judging those manifestations whose overflowing fullness of ideas
sates reflection and softens the objectionable. With this we have gained the
standpoint from which the angel of annunciation is to be regarded. For it is
so far from being the case that those prophetic visions should have greater
spiritual validity than the angelic manifestation of our Gospel, that the
appropriation of the latter is rather for us the proper guarantee that we do
not misunderstand the former. It must certainly be granted that the angel's
coming is here kept in so naive and external a form of actuality that he who
cannot notice the objectionable at all must either have the perfect faith of
blessed simplicity, or be ensnared in superstition, or be a dull and
thoughtless reader. But this hardness of the fact, against which reflection
stumbles, is precisely the touchstone of appropriation. It is the
stumbling-block that awakens faith. If I cannot believe the wonder of the
annunciation, namely because it is a wonder, then, for all my philosophizing
over the prophetic visions of the future, I can never believe the prophets
themselves either; that is to say: I do not have faith at all. The believer,
on the other hand, who understands the prophets as they, properly speaking,
want to be understood, also understands the angel as Zacharias understood
him; that is, in --- ``fear and trembling.'' The believer who grasps the
spirit of revelation will therefore not dispute the fact, but fight for the
fact, until the light dawns and the angel wins the victory. The believer who
grasps the spirit of revelation will not sit down to brood over the
bodiliness of the bodiless one, lest he lose his understanding; will not
brazenly force his way into the holy place and stare at the glory of the
divine one, lest he lose his sight; but he will in all humility draw near to
the veil, and incline his ear toward the sanctuary, and listen to the angel's
voice as it sounds in the sanctuary: in this voice he will then
% --- p. 20 ---
\opage{20} recognize the prophetic voice; for the angel is in accord with the
prophets.

The modern consciousness can doubtless not deny its liberal view here.
``Should heaven be so gracious; should there really be, yonder behind the
veil, an angel who could still whisper a word of comfort to a grieving soul,
then it would truly be a sin if some loud-voiced person came too near him
and disturbed the consolation. Everyone must have his freedom in this
respect. If faith wants to remain standing in all simplicity, and preach the
incomprehensible: very well, let it stand! But it must put up with the world
walking away from it; for the world simply will not remain standing; the
world goes with the times, and time goes, and goes constantly forward. The
world goes toward the light, and enlightenment goes forward, after all,
constantly forward. The vital questions of the present therefore presuppose
quite other ideals than the angels in the days of King Herod, and require
quite other exemplars than Zacharias and Elisabeth.'' --- But if the present
is really so liberal-minded, then our Gospel is certainly liberal-minded
too, and to that extent the dispute should soon be able to end. The angel is
not obtrusive; he does not step out of the sanctuary; he has no wish, after
all, to stand in the way of any world-ideal whatever. No, Gabriel does not
go about wearying anyone with his presence, or wrangling with anyone about
his authority: he comes as an offer, when he is expressly sent; but he is
sent only to those who need such a sending. Believe him, he has no time at
all to go forward with the times: he must in all simplicity remain standing;
for --- he stands before God. Believe him, he is truly not envious of the
progress of enlightenment; he stands before God, and he who stands before
God, should he lack enlightenment?

With our believing old couple too, who in relation to the angelic
% --- p. 21 ---
\opage{21} manifestation are placed here as evangelical exemplars, the
present will very easily be done. For these pious spouses are not obtrusive
either; they make no pretensions at all, but keep to themselves in all
simplicity. Elisabeth is not so vain that she should now go and change her
dress in order to be modern: she withdraws, after all, in all quietness and
seeks concealment in her own house. Zacharias does not have the ambition to
want to lead the movements of the age; the wonder that struck him with
astonishment stopped the flight of his thoughts: he had to fall dumb. ---
Without being burdened by the wonder of the Gospel, the world can thus, as
is fitting, work on at the solution of its own tasks, and try all the means
of temporality for the attainment of temporal ends. A question as to where
historical progress finally leads, however, one may well permit oneself.
Should the deeper meaning of existence infallibly be this, that every
thought of the eternal ought to be abolished as soon as possible, so that the
whole striving of the spirit may concentrate itself all the more forcefully
on the temporal, then enlightenment will at the very first opportunity be
completely finished with revelation, and worldly prudence easily settle
accounts with both God and his angels. But when this cannot be the meaning,
when temporality still goes on hungering in the midst of its abundance, and
does not itself possess that with which it can still its hunger; when the
world must continually change color, now from indignation, now from shame,
that the expected solution of its tasks, against everyone's expectation,
dissolved into a disappointment: then the movement of culture cannot simply
go on running out into infinite finitude either; but at every decisive
moment progress must halt at the ``melancholy'' question: is there then a
God, is there an eternity? Hence it comes that the world, for all its
unruliness, can never really bring itself to an open break with the higher
governance. Only
% --- p. 22 ---
\opage{22} in a fit of demonic despair does enlightenment now and then let it
be known that, in the rapid course of its progress, it has sped so far past
``God and devil, heaven and hell'' that even the most cowardly need no
longer fear these bugbears; but it is, as was said, only in a fit of
effervescent humor that the foolhardy spokesmen of enlightenment speak so
beyond themselves. The world itself grows ever more sober-minded as it
grows more cultivated; why, then, should a cultivated world be so
ill-advised as to break with God in heaven? As mankind makes progress, it
also gains experience, and, it goes without saying, experience makes one
wise. Worldly prudence has already come so far that it does not dare, in
dead earnest, to tear itself loose from the Gospel. Why indeed should it do
so? Such a hazard is unhealthy, and an inward unhealthiness could easily
cause restless dreams. If the matter is now seen from this side, then it is
by no means the evangelical exemplars that press themselves upon
enlightenment; it is, on the contrary, enlightenment itself that must again
and again turn round and take up the Gospel. At first glance it then looks
as though enlightenment were the superior power that was to draw the
evangelical content along into the higher development, purify the gold of
its base admixtures, and thereby make manifest what revelation really is.
But the more the reason of the world exerts itself in this direction, the
clearer it becomes that the superiority is on the side of the Gospel, that it
is revelation itself which, more deeply understood, is to enlighten us about
what enlightenment is in relation to faith. If, now, the faith of the
Gospel and enlightenment are, under the conditions of the age, as
inseparable as they are incompatible, then there is nothing else for it but
that they must incessantly wrestle with each other. In this wrestling it
becomes apparent that faith not only remains standing in all simplicity,
% --- p. 23 ---
\opage{23} but also, once touched by reflection, sets strong thoughts in
motion in order to assert its simplicity, and thereby transforms itself into
the most critical of all the tasks of criticism. Among the many practical,
unavoidable, pressing questions of the age, the question of the essence and
distinguishing mark of the faith of the Gospel will always make itself heard
as a vital question. If, then, the Gospel before us is to be rightly
appropriated, if the appropriation is to be secured against the
misunderstanding of enlightenment, and the believer is to give himself an
account of what it really is that he believes when he believes the Gospel,
then the holy exemplars must set themselves in motion, then Elisabeth and
Zacharias, indeed the angel himself, must work with critical force and help
us to find the explanation.

Enlightenment readily approves that Elisabeth confided her sorrow to God,
since she believed in God. But if in her heart she cherished the hope that
the Ruler of all would make an extraordinary arrangement to wipe away her
tears and grant her desire, then an enlightened understanding can by no
means approve of this. ``I, this single human being, am not entitled to
assume that the God who governs the countless worlds according to eternal,
unchangeable laws should concern himself with my petty sorrows and take care
of my private affairs. Only in the same degree as I become conscious that
the purpose of my life has significance for the great whole, only in that
same degree may I also believe my lot to be taken up into the special
governance of Providence. Let us for a moment turn from Elisabeth to an
exemplar that stands nearer to the present. A lovable princess shares the
throne with the ruler who, with almost superhuman power, has founded a new, a
glorious realm. The throne demands an heir, the ruler a son; but hope ---
fails the princess. Were such a woman not religious from the first, she
would have to become so by the force of circumstances. So happy,
% --- p. 24 ---
\opage{24} and yet so unhappy! Yes, she is indeed the happy one: no
wonder-flower is so wondrous as her wondrous happiness; no diadem can
sparkle like her glance when it rests in quiet rapture upon Him, the hero
whose praise is the jubilation of the nations; upon Him, the ruler, before
whom the mighty bow in admiration. But --- `the throne demands an heir, the
ruler a son.' She feels it in her anguish, each time she looks at her adored
lord --- that his eye stares; she understands it in secret torment, when she
truly regards Him, the ruler, that a shadow of despondency darkens his
countenance. So she is then, among all the unhappy, the most unhappy of all,
if hope fails. This sorrow a woman cannot bear alone, and none of all those
who would carry her in their arms can bear it for her: then she turns in the
anguish of her soul to God, that for his mercy's sake he might hear her sigh
and be moved by her tears. Were there now an angel of heaven who would
promise this oppressed heart the fulfillment of a single, ah, of only a
single wish: then surely we have here one who prays, who in unspeakable
happiness would offer up her whole life in quiet thanksgiving. Were there
now a mercy for one who must otherwise despair over her lot: then surely
here was one in anguish who on a single throw of the dice could win
everything, oh yes, but also lose everything in the strange game of fortune.
Were there now a governance that for the individual human being's sake would
bend the law of necessity, when the individual had become entangled in the
threads of great events: then surely we would have here a single being whose
destiny reaches far beyond the narrow bounds of a family life, and proves
decisive for the destiny of a realm, for the weal and woe of millions. But it
is an awkward thing, this immediate faith in God's special governance.
Perhaps the expectation is fulfilled in the end: and then thousands will
congratulate the happy one, and ten thousand imagine that they are praising
the governing
% --- p. 25 ---
\opage{25} Providence. Perhaps the hope shall fail: and then the unhappy one, cast off
by the ruler, forsaken by the world, and forgotten by men, must try to
save her soul in the understanding that the governance goes its own ways,
and demands the individual's will as a sacrifice.'' --- In this modern
allusion it must now be easy to find out the real meaning. The old
Elisabeth, after all, holds no high station, and is on the whole too
insignificant to expect the extraordinary from God (for that the expected
son should become so great a man, a prophet like Elias, she could not
know, of course, before the angel had revealed it); ``that she
nevertheless so faithfully cherishes the dear hope, and that for so many
long years, may well be a touching proof of the patriarchal simplicity of
that age, but superstitious it is all the same; of the difficulties in
which immediate faith in Providence entangles the thinking human being
this daughter of Aaron has as yet no inkling, and consequently she cannot
be set up as an exemplar of faith for the present day either.''

Here we must now keep a watchful eye on criticism, lest it criticize us
into a bottomless misunderstanding. If enlightenment will be so resolute
as to reject faith, and with it all religious truth, then it is in order
that it should without further ado dissolve and reject the exemplar; but
so long as there is still something in the world that is called
religion, the Gospel will assert its exemplars, and the daughter of Aaron
will continue to the very last to bear witness to faith. Without faith a
human being can at most grasp himself as a link in the rational order of
the world, and must resign himself to the small part's being sacrificed
for ``the great whole.'' Self-preservation fights a hard fight with
necessity; but the law is inflexible: here there is no rescue for the
individual. Called to work for universal ends, the gifted man will gladly
give himself up to the power of the idea, in case it might perhaps save
him out of corruptibility; but the pure idea is a faithless friend: when
the individual believes himself nearest of all to his
% --- p. 26 ---
\opage{26} high goal, the idea casts the finite instrument away, and plunges
into the course of the world, as the stream into the sea. ``Is there then
no salvation at all for the self of man?'' Yes, surely there is
salvation, as truly as God lives, a salvation in faith. In faith
everything is changed, the world's relation to man, and man's relation to
all things; for man changes his position on earth when he finds his God
in heaven. What the law of nature is not able to do, to rescue the
individual from destruction, since it consists only in the vicissitude of
things; what the power of the idea is not able to do, to tear the single
human being out of the grip of the elemental forces, since in its
peculiarity it is after all only a relative power among the many powers
of existence: behold, that the God of the believers is able to do; for
God is indeed the Almighty. How it can otherwise be conceived that,
despite the countless powers of finitude, there can be an Almighty who
has all things in his power, of that faith knows nothing; it does not
engage in any comprehending, and offers no explanatory paraphrase, but
says to every human being who would save his soul, short and to the
point: this you shall believe, that God is the Almighty. If this, then, is
the distinguishing mark of faith, that it so faithfully preserves the
strong, immediate impression of God's omnipotence; if the Believer simply
cannot get higher than to the Almighty, but must find it blessed to be
under his keeping, and to dwell under the shadow of his wings; then
Elisabeth is still an exemplar of faith.

At this point we can perhaps best enter into closer negotiation with the
enlightened. --- In the world a many-sided distinction is made, and
rightly so, between high and low, great and small, strong and weak,
significant and insignificant human beings. Of the various endeavors and
aims it likewise holds that some are unimportant, others important, and
the one far more important than the other. But should the same
distinction now also, in the same sense, be able to
% --- p. 27 ---
\opage{27} hold for God, the Almighty? Should any worldly end really be so
grandiose that it could impress God? or should anyone need to seat
himself upon a throne in order to attract the attention of Providence?
To attract the attention of a public it is certainly advantageous to
occupy a higher place in society; but with God's attention it is quite
another matter, since the Almighty is at the same time the All-knowing.
What is the whole world to God? --- A nothing, as truly as God is
almighty. (That this lies in the concept of omnipotence the enlightened
must surely know best themselves, since it is precisely they who search
out the concepts.) Either, then, there is no governance at all, because
everything, even the very greatest, is too little for Providence, a
nothing for the Almighty; or else the governing will must be assumed to
work in the least as in the greatest, in the life of the individual as
in the destiny of millions. To forestall any misinterpretation, it must be
expressly maintained that this vanishing of worldly differences by no
means abolishes the concept of the special governance, as if everything
were at bottom of equal validity, that is to say: were in itself
indifferent in the eyes of Providence; since faith on the contrary
presupposes that the Believer stands in such a relation of inwardness to
God that his life acquires precisely thereby a quite peculiar and
particular significance. The Believer now lives at the same time in
finitude; and with this it is given that the hope of finite goods can
also be taken up into true faith. But insofar as the Believer sets his
hope on any temporal good whatsoever, he does not ask at all whether this
good belongs, in the worldly sense, among the significant or the
insignificant things; but he asks himself: ``does the good you desire now
really have such a significance for your own peculiar existence that you
may in truth refer its attainment to God himself, pray to God for it,
and, if it falls to your
% --- p. 28 ---
\opage{28} lot, receive it as a gift of God?'' --- On this condition the temporal
may be united with the eternal; for the attainment or non-attainment of
the finite good then becomes a matter between man and God; which can also
be expressed thus: man commits his cause to God. However much in the world
may depend on circumstances and occurrences, on fortune and misfortune,
there is nevertheless one thing that is withdrawn from the play of
contingencies, and that is the cause which a human being in faith has
committed to God. If it is now quite certain that enlightenment, with all
its long strides of progress, can never get any further with the
governing Providence, then Elisabeth's simplicity is indeed the highest
wisdom. This daughter of Aaron must surely understand what it means to
commit one's cause to God; she has, after all, had to sacrifice a whole
life to learning it. She grows old, and hears the contentious judgments of
men about prosperity and adversity, fortune and misfortune, how they now
hope the best in all security, now again despair and fear the worst; but
her thoughts are not confused, her pious mind is not darkened: she grows
old in equanimity and young again in gentleness; for she has committed
her cause to God. She becomes old, and hears the eager strife of men
about punishment and guilt, reward and merit, how So-and-so got his
deserved reward because he did what is evil; but So-and-so found no
reward with God, though he thought himself so righteous: she knows that
she is not among the fortunate; but there is no bitterness in her soul.
She knows that she cannot boast before men of her reward with God; but
she has no gnawing doubt: she is silent and believes, and bears her shame.
So she has now grown old on the way of election, and yet --- young;
transformed, and yet --- unchanged: the expectation of youth bears its
fruit in the hope of old age. There is an inward peace, a
% --- p. 29 ---
\opage{29} quiet transfiguration over the aged woman: she has committed her cause to
God.

What has been said here of Elisabeth holds in like manner of Zacharias.
The angel's proclamation begins with these words: ``Fear not, Zacharias,
thy prayer is heard!'' This greeting bears express witness that the pious
priest, even in his advanced age (he says himself, V.\ 18: I am old, and
my wife also is well stricken in years), has still not forgotten to pray
the same prayer that he prayed when he stood in the strength of his
manhood; for it is surely not conceivable that the angel should remind
him of a prayer he had long since forgotten. (The divine tolerates no
forgetfulness; if God hears forgotten prayers, it becomes a punishment to
a human being, and not a joy.) To that extent, then, Zacharias has not
heeded the order of the laws of nature, but on this point alone has asked
after God's will, committed the cause to God, and precisely thereby has
come to overstep the boundary of probabilities. If enlightenment is not
satisfied with Elisabeth, it can be still less so with Zacharias: ``this
simple, narrow man with his insignificant family worries has, after all,
no claim either on the sympathy of the present day.'' Of course, the
historical consciousness demands quite another satisfaction. It is
therefore probably best that we (so as not entirely to overlook the
demands of the present day) consider once more that modern exemplar to
which we were referred above. ``The world-genius has thus come into his
own, the idolized hero has won the diadem: for his own sake, for his
throne's, for the millions' sake he then turns in a quiet thought to the
governance, that it might now secure his realm and grant him the hoped-for
son. But it goes without saying that the law of the world must here, as
everywhere, naturally remain in force; fortune and misfortune, fate and
its whims, are taken into the reckoning. If probability stands against
the wish, the ruler must try to find another
% --- p. 30 ---
\opage{30} way out: either he then gives up the expectation and takes other
measures for the security of the realm, or else he must, before it is too
late, dictate a divorce and enter into a new union, so that fortune might
yet comply with him in everything, and probability support the hope.''
Whether the ruler acts wisely or unwisely in this, moreover, is not the
question here; but certain it is that calculations of probability of this
kind have nothing to do with the faith of the Gospel. On faith's behalf
Zacharias accordingly goes another way: he does not give up his hope,
though time passes; he does not seek divorce, though Elisabeth grows old;
but he holds fast to the Almighty, commits the cause to God, and ---
prays.

Perhaps it would sound like a not unwitty mockery if a modern critic now
hit upon the notion that we were here setting up ``a quite apt exemplar
for all those who by prayer and conjuration seek to obtain what in this
sinful world they cannot obtain by the use of natural powers, either
because they lack all power whatsoever, or because they are too
helpless to use the powers our Lord has given them.'' But since it is not
for the mockers and those given to laughter, but for those who in
earnest wish to be taught what faith is able to do, that the Gospel is
proclaimed: let us consider more closely the question whether it does not
lie precisely in the nature of the case that the Believer first comes to
clarity about himself when he begins to break with probability. So long
as the present affords us satisfaction without foreboding any danger for
the future, so long as probability on all sides, as it were, fences in
and secures our whole existence, it is no great matter to trust in God;
it is roughly the same as trusting in the natural order of things. But as
soon as probability merely turns round and shows its dark side, the test
can be passed at once. --- It is dangerous to go into a plague-house;
% --- p. 31 ---
\opage{31} there is every probability that one will be infected. The courageous man
can nevertheless enter and do his duty, though he lacks religious faith;
the Believer can likewise enter and do his duty, though he lacks natural
courage. The former does not care about probability, because he trusts
in himself. The latter, on the other hand, breaks with probability,
because he trusts in God. The courageous man has the fullest recognition
of the significance of probability; he knows that here, more than in any
other place, he is in the greatest danger; but he defies the danger; for
he can die, after all, and with that the matter is ended. The Believer,
on the other hand, cannot defy the danger; so he bids defiance to
probability, and gives himself up to God's will. If God wills that he
shall go forth from there unharmed, then no contagion can touch him; if
God wills that he shall go through the sickness, then He gives him
strength to suffer; if God wills that he shall die, then it is for his
good. But if it is now a distinguishing mark of faith to close the eye to
the finite conditions of nature and to defy the probabilities, then
Zacharias is indeed an exemplar of faith.

Enlightenment, which cannot do without an almighty God, has in place of
the faith of the Gospel gradually had another faith introduced, which is
supposed to fare better with the growing knowledge of the laws of nature,
namely the so-called faith of reason. According to the faith of reason,
God is certainly due the honor of being the Creator of nature and the
Lawgiver of the world from the beginning. But since the time the world
came into being (and it is surely well over 6000 years since everything
was arranged for the best), God has had in a manner to suspend his
omnipotence and let the laws of nature govern themselves; and why?
because he is far too reasonable to want to disturb the order once
introduced. If the concept of governance is now held fast in pure,
indeterminate generality, then the faith of reason seems in all
essentials to agree with the evangelical faith. On both sides it is
granted that
% --- p. 32 ---
\opage{32} God is almighty, that nothing happens without his will, that man ought
therefore to acknowledge the governance of Providence in all things.
Nonetheless there is in reality an infinite difference here. This
difference shows itself most clearly of all in the determination of the
probable and the possible. Let two human beings agree to refer the
dispensations of life to God, but in such a way that the first
understands it according to the faith of reason, while the second holds
to the Gospel: we shall then soon learn that there lies a world of
difference between them. With the mutual assurance that man's life is in
God's hand, they set out on their journey. On the way they fall into
distress at sea, and both seize the same plank of rescue. When the first
now looks about him on all sides and sees nothing but the wild sea, when
the probability of going under stands to the probability of being saved
as a hundred to one, he bursts out: ``Rescue is impossible, we must
perish!'' But the second answers: ``My dear friend! why should rescue be
impossible here? do remember that man's life is in God's hand!'' So they
are both saved, and join in thanking the governance. But in the foreign
land to which they have now come, they are condemned by barbarians to be
cast into a burning furnace. Then the first laments in the bitterness of
his soul: ``See, now all rescue is impossible! Better it would have been
for us had we remained on the wild sea than that we should come to this
barbarous land to suffer so ignominious a death.'' But the second
answers: ``Yes indeed, rescue is possible; even if we must go through the
fire, rescue is still possible.'' This answer will then be foolish talk
to the man of reason; he will be inwardly incensed at so foolish a game
with possibilities: ``Here, where flight is impossible, where mercy is
impossible, where it is as impossible that the death sentence be changed
as it is utterly impossible that the burning fire should change its
nature: how can anyone
% --- p. 33 ---
\opage{33} here want to play with possibility and make mockery of fate?'' But the
second shall answer him and say: ``My dear friend, do not forget what you
yourself said when we set out on this journey: man's life is in God's
hand.'' So they are both saved, and are agreed upon \dots\ yes, what is
it, really, that they are agreed upon? Both are no doubt agreed upon
using their natural understanding insofar as it is a matter of observing
the natural circumstances and calculating the probabilities. In this the
one need not fall short of the other at all. But in acknowledging the
governance they can hardly be agreed. The first looks only at the law of
nature, and is so taken with the probable that in every decisive case he
forgets God; the second too has, no doubt, a worldly eye for the natural
conditions; but in the moment of decision he concentrates, so to speak,
his whole power of sight in the eye with which he looks toward God, in
such a way that he thereby comes to overlook probability. The former has
a dim and hazy concept of a divine omnipotence, which he denies in deed
every time he ought actually to apply it; the latter has so living an
impression of the omnipresence of the Almighty that none of the powers of
finitude is able to overwhelm him. In other words: the first really has
no faith at all; the second, on the other hand, is a Believer.

This difference shows itself most clearly of all when the task is to
expect a finite good that is to be received as a gift of God. Though they
both cherish the same wish, and agree, each for himself, to commit the
matter to God, they cannot for long go together: their ways part. The
faith of reason goes no further than probability and reasonable prospects
permit. ``If the fulfillment of the wish is a flat impossibility; if any
reasonable person can see that nothing will come of it, then there is,
after all, nothing more to be done about the matter. Instead of troubling
God with unreasonable
% --- p. 34 ---
\opage{34} prayers, a reasonable man then does much better to foster another wish,
and look about for other attainable goods.'' The faith of the Gospel, on
the other hand, cannot permit so rich a variety in wishes and hopes. Once
the Believer has understood the significance of his wish and committed
the matter to God, he has by that very act bound himself by the wish, and
must now, despite all the changes of weather in circumstances and
probabilities, remain standing and wait for what God will do. By means of
the wish man learns to understand his God as he learns to understand
himself; and that, after all, is also what he prays for. But would it
not now be presumptuous if one who prays were to begin thus: ``Almighty
God! I have hitherto hoped that thou wouldst fulfill my dearest wish; but
now, since circumstances have changed, and since time has advanced so far
that thou \emph{canst} not possibly fulfill it, I will now let that
matter go, and pray thee another prayer, which thou canst perhaps still
fulfill.'' That such a prayer is a blasphemy can surely not be doubted.
For if the wish was so rash that the one praying must at once have been
able to see its unreasonableness: how then could a human being make so
bold as to come before God with such a wish? But if it is only worldly
improbabilities that stand in the way of its fulfillment: how then dare
anyone venture to speak to the Almighty of the impossibility of
fulfilling it? No, it is something quite different that matters here. The
one who prays holds himself assured that God both can and will hear him,
as surely as the hearing is for his good; but with this it is at the same
time given that the not-hearing may also be for his good. This is
precisely what makes the task difficult. Which of the two cases may be
taking place here, time can surely never teach the Believer; but he can
and shall learn it in time, as he learns to wait upon God and to
understand himself in relation
% --- p. 35 ---
\opage{35} to the Almighty. If such a waiting now belongs among the essential
requirements of faith, then we can well understand what held Zacharias
to the fulfillment of the one wish, so that in a certain sense he did not
notice at all that circumstances changed, that the years passed, and old
age came; but he kept on with the same wish and the same prayer, even
when a whole generation had passed him by, even when he had to say to
himself: ``I am old, and my wife also is well stricken in years.''

It is not exactly as easy as sneering criticism pretends to commit one's
cause to God and to wait for the Almighty. No, it is much easier to move
with the times. For he who moves with the times is carried by
circumstances and driven by probabilities. But he who is to wait upon God
cannot, in the finite sense, even stand still (for then finitude would
sweep him along with it); but in order to remain at his post and wait, he
must incessantly wrestle with the pressing circumstances and work against
the current of probabilities. It is so reassuring to walk the broad road
of probabilities: if I do not find a foothold here, then I find it there;
even if I should fall, the harm would hardly be greater than that I could
get up again: it is not so far to the ground. On the other hand it is very
fearful to walk along a narrow path over an abyss; the slightest misstep
can bring death. The probability of falling is so great that I dare not
heed it at all, lest I grow dizzy; I dare not look to either side at all,
but must look straight ahead without interruption, and am in great need
of someone who could go before me and say: look at me!

The faith of reason, which keeps to the wide-stretching surface of
probabilities, can easily stride forward with sure steps, for the broad
earth is, after all, firm enough to bear a human being. But precisely
because it bears so well, and everything goes as of itself according to
the laws of gravity
% --- p. 36 ---
\opage{36} and equilibrium, omnipotence is here almost superfluous. The faith of the
Gospel, on the other hand, never has more room to spare in finitude than
a human being needs each time just to set his foot on. Therefore the
Believer must beware of looking too deeply into the possibilities and of
turning too often toward the probabilities; in all caution he must go his
way straight ahead, and keep looking at the Almighty, lest he grow dizzy.
It is, after all, only a very small part of the earth's surface that a
human being needs for a foothold. The place on which the foot rests the
eye cannot see; the place the eye sees the foot has not yet trodden. He
who walks must nevertheless have his eyes with him, and look ahead to see
whether the place he means to tread on can also bear him. Essentially
speaking, the eye can surely not foresee with certainty how reliable the
place is; but then one feels one's way and relies on one's judgment: the
natural gait goes according to probability. If, on the other hand,
someone would form a notion of what it might be like to walk straight
against probability, then it is indeed very difficult to describe an
inward movement by an outward image: we will nevertheless make the
following attempt. Suppose that a human being were suddenly changed in
such a way that the earth became transparent to his eye, still more
transparent than the still, clear lake that in its depths mirrors the
clouds and the blue sky. Suppose that the same human being were further
transformed in such a way that the transparent earth would only grow firm
under his foot and bear him when he stepped out boldly with resolute
mind, but would on the other hand dissolve and give way, fleeting and
faithless as the pure air, every time he groped for a foothold in anxious
wavering. Then we would have here a wanderer who could not take a single
step without coming into conflict with probability. Yes, it is a strange
position. See, he stands on the firm earth like all the
% --- p. 37 ---
\opage{37} others; and yet it seems to him as though he stood on a column of air
between heaven and heaven. The earth can really bear him, just as well as
all the others, if only he steps out boldly; and yet it looks in his eyes
as though with a single step he would plunge deep down into --- the
clouds below, yes, through the opening of the clouds down into the blue
infinite deep. But just as his strength is leaving him, and everything is
about to go black before his eyes, a soft voice sounds: ``Fear not, only
step out boldly, set your foot here; but if you waver, then woe to you,
for then you are a dead man!'' So he takes every step straight against
probability; but it is a slow gait; it is as though his breath would stop
every time he moves a foot.

The others can easily step out boldly; for to them the dark earth has no
transparency, and gives no mirror image of the deep heaven; the others
can easily go briskly forward; for to them there is no danger that the
firm earth should dissolve into air: wherever they look they everywhere
find a sure foothold; indeed all the others can easily hurry far past our
wanderer; for it is, in the finite sense, far easier to take a hundred
steps according to probability without the express assistance of
omnipotence than a single step against the law of the probable in
unconditional trust in the Almighty. But if it is nonetheless the
unavoidable requirement of the faith of the Gospel that a human being go
his way even against all probability, then Zacharias is an exemplar of
faith; for just such a gait did the old priest walk when he went into
the holy place to stand before God's angel.

Enlightenment, however, claims to know better about the probable, and
therefore throws itself with all its objections upon Gabriel and his
wondrous annunciation. ``What in the present narrative seems to justify
the pious couple's unconditional faith, and gives their expectation so
edifying a tinge, is after all at
% --- p. 38 ---
\opage{38} bottom nothing other than the angelic manifestation itself, with the
highly gratifying outcome that follows upon it. According to this
direction it can then hardly be so burdensome for a Believer to engage in
conflict with the natural surroundings, since in case of need he can
always count on an angel's assistance. Against such assistance we have,
as far as the evangelical doctrine of Providence itself is concerned,
nothing to object. On the contrary, the angels are indeed `ministering
spirits, sent forth to minister to those who shall inherit salvation.'
Since, moreover, everything is so wisely arranged that the number of the
angels who stand before God's throne is ten thousand times ten thousand
and far more, while the number of true believers on earth must always be
estimated at `a very little flock': there will in every age always be far
more angels than there are believers, and the `ministry' can consequently
never become burdensome to the heavenly ones. If this is now so, then
faith is undeniably the best life insurance any human being can choose,
and a believing soul can easily bid defiance to the probabilities. But
what on the other hand makes the matter somewhat questionable is the
peculiar circumstance that in our days, when even the least remarkable
thing seldom escapes general attention, one notices not the very least
either of the angels or of their miracles. If it happens once that a
human being steps forward and appeals to an immediate manifestation or a
new miracle, then all reasonable people are of course at once agreed in
charging him with enthusiasm. It is not only the frivolous who mock him,
and the unbelieving who are scandalized at him; but even the orthodox
shake their heads, and look at one another with most dubious faces: Good
Lord, what is one to think? Two cases are possible; but whichever of
these one may assume, it nonetheless goes against the faith in miracles.
Either God has, since those days, altogether changed his
% --- p. 39 ---
\opage{39} governance, and given the angels quite another ministry, and then it
naturally follows of itself that the faith in miracles is no longer
justified; or else the miracles must have ceased, and the angels
themselves have withdrawn, because the right faith has vanished from the
earth. But should it not now be a great proof against the divine origin
of the faith in miracles that it can vanish so completely from the
earth?'' --- So much for the modern enlightenment.

In the name of revelation the faith of the Gospel is secured against every
presumptuous cleverness. The Evangelist has not concealed from us that
Zacharias was terrified, and fear fell upon him, when he saw the angel
standing on the right side of the altar of incense. Indeed, so little was
the pious man prepared for the wonder proclaimed to him that even in the
moment of surprise he sins against the Holy One.
{\large Behold, thou shalt become dumb, and not be able to speak, until the
day that this comes to pass, because thou believedst not my words, which
shall be fulfilled in their season (Luc.\ 1, 20).} Thus the angel punishes
the doubter, and by this punishment Zacharias must surely be acquitted of
the superstitious frivolity that thinks a miracle will happen at every
moment. Had the prayer not been fulfilled in this way, would the Believer
then perhaps have forfeited his hope and regarded his life as wasted, and
given up his faith at the last? Oh no, for that the aged man was far too
inwardly assured of the truth of his confidence and of God's
indescribable goodness. With this faith the courageous patriarch would
have gone to his rest, with this faith he would have thanked and praised
the Almighty in his last hour, yes, with this blessed faith he would then
have blessed his wife and smiled in death, when the angels came to carry
him into Abraham's bosom. But, of course, then the fulfillment would have
been a different one, then death would have come before the angels; now,
on the contrary, the angel came before death, and surprised him with the
glad
% --- p. 40 ---
\opage{40} tidings, and terrified him with his speech. It is quite correct, what
enlightenment asserts, that faith presupposes the wonder; for it draws its
nourishment from it, and is itself a wonder of the soul. In this sense the
faith of the Gospel can rightly be called a faith in miracles. It is
likewise certain that a human being who lives sincerely in this faith
must in many respects appear somewhat strange to the world. Why should he
not seem strange who has his real life in a wonder! But from this it by
no means follows that the Believer should brood on arbitrary miracles, and
roam through existence in the insipid fancy that somewhere or other he
might perhaps come upon an angel or a new manifestation. If anyone
impatiently craves new manifestations, then such craving is so far from
proving the strength of his faith that it even proves the opposite, since
it bears witness, after all, that he has not yet learned to appropriate
the revelation once given. To that extent the faith of the Gospel can very
well consist with the law of nature, and will consist with it until the
end of the world. Or should the world-spirit perhaps have reason to lodge
a complaint about the disturbing miracles of revelation? See, kingdoms
arise --- without miracles; they pass away again --- without any miracles
at all. Thrones are raised up --- without miracles; they collapse again
--- without any miracles at all. So it goes in the present age; but so it
went, after all, in the days of the Emperor Augustus and King Herod as
well. The world is on the whole strongly built, and mankind has great
powers: even if ten kingdoms fell ten times, the race would be strong
enough to build them all ten up again the eleventh time. Only of one
kingdom does it hold that the whole world does not have strength enough
to build it up from the ground; of one kingdom it holds that, if it could
ever collapse, all the powers of temporality would be too weak to raise it
from its fall, and that is the kingdom of God, the kingdom of truth, of
% --- p. 41 ---
\opage{41} righteousness and of holiness. When this kingdom is to be founded on
earth, then a wonder must happen, a great wonder, then the powers of
heaven must stir, and ---
\emph{the angels are sent forth to minister to those who shall inherit
salvation.} Therefore Gabriel's voice sounds in the holy place: {\large I
am sent to speak to thee, and to proclaim this to thee for joy.}

Great changes, remarkable advances, can take place in the history of the
nations without any miracles at all. The spirit of the age has workers
enough: where one ends, another begins; what one promises, another will
no doubt fulfill; what one prepares, another can perhaps complete. Thus
the world goes forward, after all. Forerunners are born in droves; they
awaken, they lead, they teach the nations, and point toward ``the great
future'' and ``the great goal.'' And the nations rise up, and --- feel
themselves called, and heroes arise, and feel themselves called; and the
movement begins, and the upheaval takes place, but without any miracles at
all. And there is fighting and conquering by land and by sea, but ---
without any miracles at all; and there is fighting and conquering with
word and with sword, but --- without any miracles at all. And the
awakening ends, and the victory disappoints, and the forerunners
contradict and accuse one another, and the heroes combat and overthrow one
another, and the nations quarrel and lose faith in one another, and the
jubilation falls silent, and the movement halts, and the whole
``wonder-work'' is again like an empty, toilsome nothing, --- without any
miracles at all. All this the world-spirit may venture; for the motherly
nature of the race has powers enough, and natural powers are squandered.
But there is nevertheless one thing that motherly nature is not able to
do, however lavishly it squandered its most glorious gifts on ever so
many gifted men: it cannot equip any man so strongly that by the power of
his word he should be able to awaken the sleeping conscience of men, and
lead a degenerate generation back to obedience under God's law. If such a
man is to be equipped, the power must be given him from above. But
% --- p. 42 ---
\opage{42} with the power from on high wisdom keeps thrift: the supernatural powers
are not squandered. When, therefore, the coming of the Son of God was to be
prepared, there was only one man who was called to carry out the great
work of preparation. ``Only one man''? Yes, but this one man was also
sufficiently equipped from the beginning; care was taken that his calling
and mission should be rightly grasped and understood from the beginning;
therefore the angel of the annunciation is sent forth, therefore Gabriel's
voice sounds in the sanctuary, as the human voice of the heavenly
messenger:
{\large He shall be great in the sight of the Lord; wine and strong drink
shall he not drink, and he shall be filled with the Holy Spirit even from
his mother's womb, and turn many of the children of Israel to the Lord
their God. And he shall go before him in the spirit and power of Elias, to
turn the hearts of the fathers to the children and the disobedient to the
mind of the just, to make ready for the Lord a people well prepared.}

\begin{center}
  \rule{0.38\textwidth}{0.4pt}
\end{center}

\section*{II.\\[4pt] The Angel Gabriel and the Virgin Mary.}
\addcontentsline{toc}{section}{II. The Angel Gabriel and the Virgin Mary}

\begin{center}
  {\footnotesize (Luk.\ 1, 26--38.\ \ Matth.\ 1, 18--25.)}\\[10pt]
  \rule{0.12\textwidth}{0.4pt}
\end{center}

\begingroup\large
In the sixth month of Elisabeth's time the angel Gabriel was sent from God
to the Virgin Mary in Nazareth with the wondrous tidings that she was the
highly favored one, the blessed among women, who should bear the Savior of
the world. And the virgin was terrified at his speech, and
% --- p. 43 ---
\opage{43} marveled at his greeting. But when the angel with divine voice proclaimed
to her: The Holy Spirit shall come upon thee, and the power of the Highest
shall overshadow thee; therefore also that holy thing which shall be born
of thee shall be called the Son of God: then Mary gave herself up as the
handmaid of the Lord, and the angel departed from her. --- But the Virgin
Mary was betrothed to a man named Joseph, of the house of David, and
Joseph misjudged his betrothed wife. But since he was a considerate, if
also a strict and just, man, he would not put her to open shame, but
resolved to part from her secretly. But while he thought on these things,
behold, the angel of the Lord appeared to him in a dream, and said: Joseph,
thou son of David! fear not to take unto thee Mary thy wife; for that which
is conceived in her is of the Holy Spirit. And she shall bear a son, and
thou shalt call his name Jesus; for he shall save his people from their
sins. And Joseph awoke from his dream, and denied himself, and did as the
angel of the Lord had commanded him.
\par\endgroup

The inward connection between this and the preceding narrative is evident
of itself. The same angel who revealed himself in the Temple is also sent
forth to Nazareth; the time when he appears to the virgin is fixed by the
point in time when Zacharias received his manifestation; indeed the
wonder that took place with Elisabeth in her old age is even expounded
outright to Mary as a sign that \emph{with God nothing shall be
impossible}. He who feels himself addressed by the first Gospel account
might then perhaps think to himself that this story of the annunciation
was without further ado vindicated and explained in the former, so that
one could easily believe this one once one had made oneself familiar
% --- p. 44 ---
\opage{44} with the former. Yet so credulous a presupposition might easily be
subject to a somewhat serious misunderstanding. We have, to be sure, in
the preceding development already removed several difficulties that
would otherwise stand in the way here of a believing appropriation of
this Gospel; we have already made intelligible to ourselves faith's
self-confirmation against the assaults of offense and the subtlety of
doubt; but what shows itself in the religious life holds also in relation
to sacred history: the struggle of faith never ceases. The God of the
wonder is always the same; the purpose of revelation is essentially
always the same; but in every new and extraordinary event the wonderful
takes on a new, surprising form, and doubt, which not only attacks the
Believer from without as an outward enemy, but also steals forward in his
own soul as an inward, secret enemy, will ever seize every new occasion to
insist on the unreasonableness of the wonder, and by an unexpected turn to
take the secure by surprise. Although sacred history is in the very
highest sense the history of spirit, is, to speak in the language of
philosophy, an organic whole, a connected chain of eternal ideas: it is
nevertheless so far from being the case that a subsequent miracle could
be derived from a preceding one, that faith's appropriation of every
extraordinary event must rather begin afresh, in order to fight with doubt
over \emph{everything and nothing}. Each time a single Gospel account
comforts, edifies and gladdens us, it is, after all, as though a light
dawned over the whole Gospel; for we are in faith, and possess the whole
in the single part. But conversely, when a single narrative repels and
offends us, then it is also as though the whole revelation were to
dissolve into nothing: we lose the whole with the single part; for we
lose --- faith. This now finds a particular application to the relation
between the two annunciations. Precisely because the wonder in ``the
Annunciation of Mary'' is greater
% --- p. 45 ---
\opage{45} and, so to speak, more wondrous than in the preceding narrative,
precisely for that reason doubt also sharpens its sting, so that the
penetrating appropriation requires all the greater exertion. To convince
ourselves of this we now need only listen for a moment to the mockery of
the Free Thinkers, and let the negative criticism have its say.

``The narrative of the angel's visit to Mary is,'' to speak in the spirit
of the Free Thinker, ``undeniably one of the most guileless narratives in
the N. T. How strong in faith those souls must be who, without suspecting
the difficulty, can put together such dissimilar notions as: \emph{God in
heaven and the province of Galilee, the archangel Gabriel and the market
town of Nazareth, the Holy Spirit and the Virgin Mary}. (Such astonishing
combinations of the heavenly and the earthly are not, to be sure, always a
proof of excessive strength in faith; they can sometimes also be explained
by an excessive weakness in thinking.) The names \emph{Galilee, Nazareth,
Mary}, etc.\ have now, in the course of centuries, acquired the privilege
of being placed in immediate connection with God and his angels. But such
proper names can easily be thought of as changed without the matter
itself being in the least distorted. It is only a question of choosing
such designations as call back to us the original impression, in that they
involuntarily stamp the thing and the person with the true everyday mark
of reality. Just make, then, a quite simple change of names. Choose only
the names of the province and the town from your nearest surroundings;
give the virgin a more earthly, less aesthetic name: then read aloud with
priestly pathos the first two verses of your Gospel, and you shall
perceive that it is not so easy after all to bring the supernatural into
so outward and immediate a contact with the natural.''

``But,'' one says with Gabriel, ``with God nothing
% --- p. 46 ---
\opage{46} shall be impossible. Faith is prepared for everything, even the most
improbable. Very well then! Since the old faith in miracles holds up so
well that even in our 19th century one dares publicly to put forward the
doctrine that Mary `conceived' her son of the Holy Spirit, we must then
ask for somewhat closer information concerning this virgin and her
mythical relation to her betrothed. The Evangelists Luke and Matthew are
rather sparing of words as far as these delicate points are concerned; it
happens quite fortunately, then, that we have more detailed traditions
from roughly the same time and, to judge by the view of the Church
Fathers, in roughly the same spirit as our two canonical ones. Let us
then carefully compare the two ancient source-writings, namely: the Gospel
of the wondrous birth of Mary (\textit{Ev.\ de nativitate Mariæ}) and the
Protevangelium of the Apostle James; perhaps through these endeavors a
clearer light might dawn over the true, miraculous connection of the
matter.''

``The Virgin Mary, who had the extraordinary destination of becoming
mother to the Son of God, naturally cannot have been an ordinary virgin.
One is therefore entitled to expect that she herself must have come into
the world in accordance with a special divine arrangement, and from
childhood on have led a life consecrated to God and his angels. This
conjecture is then also confirmed when we merely consult both the small
Gospels cited. Joachim and Anna, an old, very God-fearing couple, lived in
Nazareth, and were at length misjudged and scorned on account of their
long, childless marriage. But in this sorrowful situation of theirs they
now enjoyed the great grace that the angel of the Lord revealed himself to
them both, each separately, and promised them a daughter who should be
called Mary, and be chosen to bear the Savior of the world. When Mary was
6 months old,
% --- p. 47 ---
\opage{47} her mother once set her down on the ground; but the child walked only 7
steps, and sought her way back again to her mother's bosom. This Anna
then very wisely interpreted as a sign that her God-consecrated daughter
would not tread upon the earth before she had been presented in the
Temple of the Lord. From her 3rd to her 14th year Mary then lived in the
Temple, where every day she enjoyed divine manifestation; for the angels
visited her, indeed she even received her food from an angel's hand.
--- This must now be enough for us as far as the first point is
concerned.''

``Next, on faith's behalf, we wished to know somewhat more closely the
conditions under which the betrothal took place; since the manner in
which the angel, according to the canonical accounts, transforms the
actual union into a merely apparent one, exalts Mary into a \textit{Noli
me tangere} and degrades Joseph to a titular husband with duties without
rights, might otherwise easily offend those weaker in faith as an
arbitrary procedure little worthy of God. In this respect too the two
writings cited render good service. For it is evident both from the
Gospel of the birth of Mary and from the corresponding Protevangelium that
it was never the intention either of Joseph or of Mary to enter into any
ordinary conjugal union with each other. Only after the high priest
Zacharias had inquired of God in the Holy of Holies did Mary, in
accordance with the divine oracle (Jes.\ 11, 1), have to consent to an
apparent spouse, or at least a guardian, being chosen for her from among
the men of the house of David. It was then decided that the marriageable
men of the Davidic family, or, more precisely according to the
Protevangelium, all the widowers among the people, should bring their
rods to the altar, so that by a sign (4 Moseb.\ 17) it might be shown who
was the worthy one. Joseph was at that time already an old man, and at
first would not come
% --- p. 48 ---
\opage{48} forward with his rod, for which reason no sign appeared either. The high
priest then had to inquire of God again, and when Joseph's rod had been
handed over, the sign took place, in that a flower shot forth from the
rod, and a dove that came down from heaven settled on the tip of it. ---
See, this is an explanation as edifying as it is credible. If one is to
believe the great wonder, why not the lesser ones too? With God, as
Gabriel says, nothing shall be impossible. Faith, after all, is prepared
for everything, even for the most improbable.''

It is by no means so easy for the highly learned theologians to take up
the gauntlet that the negative criticism, with a bold, challenging look,
throws down to them on the battlefield of science. To take refuge in the
artifices of natural interpretation does not help; for when the angel is
transformed into a natural lover, then the offense cannot be averted that
Mary, either deliberately or in the delusion of enthusiasm, deceived
Joseph, and that Christ, the Sinless One, who was to save the world from
sin, was thus himself conceived in sin and born in iniquity. But if one is
now outraged, on behalf of religiousness and morality, at so base a
maltreatment of the venerable tradition, without yet being able to bring
oneself to accept the wonder as our canonical Gospels present it, then
there is only one way out left, namely, to declare the whole narrative a
myth. --- That the dissolving criticism especially selects this narrative
as a striking example of the myth's being at home in the N. T. as well is
quite in order; but even such theologians as, with all their scientific
criticism, still seek to preserve what they call the kernel of faith have,
in an age that so easily confuses the Holy Spirit of revelation with the
pure world-idea, and believing appropriation with the knowledge of
philosophy, not been able to resist the temptation. For while the wonder,
as the only one of its kind, has all analogy against it,
the myth has
% --- p. 49 ---
\opage{49} on the other hand, in the present case, a rich multitude of ancient analogies on
its side. Heroes such as Heracles, Romulus and Alexander, sages such as Zoroaster,
Pythagoras and Plato, are of course represented in mythical poetry as virgin-born
sons of gods. The concept of myth is, moreover, so pliant, ambiguous and hovering
that it permits the most diverse shadings: ideal and reality, fable and fact can
mingle here, and enter into combinations in very unequal measure. The mythical
form is really the garb in which the fantast decks out the idea, in order to
introduce it with more-than-real splendor into the meager world of reality. Thus
the mythical idea comes to glitter like the kind of beauty that connoisseurs, in
the language of the salons, call a ``distance dazzler''; thus the idea reposes,
voluptuously swelling, in the forms of myth, and reflection, which would
comprehend the fusion of the eternal with the temporal, now takes refuge first
where it finds itself least hampered by the confines of actual life. But when
youthful reflection has for a time remained enamored of the mythical phenomenon,
it begins in all sobriety to knock against the sharper edges of reality, and so
awakens at last with the critical recognition that the ideal glory of myth is
only an illusion.

``But,'' objects a theological school, ``it is not at all necessary to let
oneself be tossed about between such extremes, so that one either, with the
old-fashioned believers, holds fast to the wonder in the immediate literal form
of the account, or else, with the modern enemies of revelation, denies all
belief in a divine co-operation because the idea of this co-operation has been
handed down to us in the mythical garb of legend. Theological sense will know
how to find the Middle Way, so that something of the essential content of the
narrative is preserved for faith, while the part that rests more directly on the
form is, on the other hand, sacrificed to scientific criticism.''
% --- p. 50 ---
\opage{50} So, a proposal for a reasonable and decent balance between faith and
criticism! Could it be possible? Could a shrewd prudence really succeed in
bringing about so lenient an agreement between the contending powers? Well then!
Why should we be so foolish as to flee the Middle Way, if it really is the way
that leads to the goal? Far be it from us to wish to drown out the voice of
those who walk the Middle Way, if only this voice is the voice of truth. Far be
it from us to wish to deny the doctrine of the Middle Way its due respect. On
the contrary! We admire the shrewdness, we honor the prudence, if it really
accomplishes such great things; we would gladly listen to an ingenious proposal.
``The method is this: Let it be assumed that Jesus, in consequence of the
wondrous co-operation of the divine Spirit, was born into the world without sin:
by this dogmatic presupposition one satisfies Christian faith; let it next be
conceded that Joseph, in his lawful marriage with Mary, is the natural father of
Jesus, and that, as a consequence, the angelic manifestation, with all that
belongs to it, is due solely to legendary poetry: by this concession one
satisfies criticism.'' --- If only one little unpleasantness did not creep in!
``What unpleasantness?'' This one, namely: that both faith and criticism feel so
offended by the proposed division that, in the most extreme disagreement, they
agree to protest against the settlement. ``What do you mean? Should faith, the
true Christian faith, really suffer from such a craving for miracles that, with respect to this wonder, it would permit no alteration at all in the individual
features of the account?'' By no means! True faith does not crave miracles (the
craving for miracles is to be sought in superstition). True faith does not
presume to prescribe to God the way of redemption; it does not dispute that the wonder could, after all, also have taken place in another manner. True faith
does not deny that Mark and John have, each for his part, handed down to us
% --- p. 51 ---
\opage{51} a Gospel complete in its kind; true faith knows within itself that it would
not perish even if it pleased God to take this Gospel of the Annunciation away,
if only in his grace he let the words of the confession of faith remain. True
faith therefore does not attempt, either, to derive the Gospel of the
Annunciation as a necessary consequence of the necessary ground of the idea of
revelation. No, true faith is not so arbitrary, so presumptuous and self-willed.
When faith nevertheless protests against the attempt at mediation just cited,
its objections proceed from a quite different ground. The narrative that in the
present case was to fall as a sacrifice to criticism has, after all, been handed
down to us as an integral member of the canonical Gospel cycle; the same law
that holds for the interpretation of the Gospels in general must then surely
hold also for the interpretation of this Gospel in particular.

But the inviolable fundamental law of the faith of the Gospel is this: that no
wonder can be the object of a believing apprehension unless it is presented to
us in a determinate form of reality; for the wonder is a fact, and every fact
must always be represented in a certain way. If I now learn from criticism that
our present Gospel account is a failed and falsified representation of the idea
of Jesus' sinless birth, then naturally I must give up faith in this Gospel. My
thought may perhaps occupy itself with the idea; but faith, having lost the
evangelical representation, gropes blindly among loose conjectures: \emph{thus}?
or \emph{thus}? or perhaps \emph{quite otherwise}? So as not to be
misunderstood, we must expressly remark here that the determinate \emph{How}
that faith asks about is by no means a How of comprehension, but a \emph{How} of
representation. Indeed, if I could comprehend how the wonder necessarily had to
have taken place, why then I would not need to be so particular about the
evangelical mode of
% --- p. 52 ---
\opage{52} representation; but precisely because faith does not have the concept, it
must necessarily hold to the representation. He who has neither any insight into
the essence of the thing nor any representation of its outward appearance, what
does he have? Loose thoughts, unstable opinions, empty conjectures. The critical
interpreter who strikes out on the recommended Middle Way may turn and twist as
he likes; still he gets into difficulty with faith. If he says: ``We regard the
evangelical account as failed because the wonder in question cannot be
designated by any representation whatsoever''; then this implies that the
wonder itself is incredible, since it contradicts every possible supposition.
The evangelist then meant to give us a representation of that which by its
nature conflicts with all representation; he meant to make the impossible
possible, having set himself an altogether unevangelical task. But against such
an assertion faith must protest most decidedly; for a Gospel writer who so
misunderstands his calling that he begins with the solution of impossible tasks
can surely, on the whole, be regarded as little worthy of trust. If, on the
other hand, the critic replies that the solution of the task is by no means
impossible, only that the evangelist, for easily explicable reasons, let himself
be misled by the mythical legend; if for that reason he undertakes to install
for us a more adequate form of representation in place of the evangelical one:
then --- yes, what then does the critic do? \emph{He dictates a new Gospel!} If
we now cast away all evasive ambiguities, the new Gospel reads, short and
definite, thus: ``\emph{Jesus is indeed a natural fruit of the marriage of
Joseph and Mary; but by means of the co-operation of the divine Spirit he
nevertheless came into the world, a sinless child of sinful parents.}'' One must
truly be very sure of one's case to venture so decisive a step. Faith is surely
within its rights when it asks here, first, what
% --- p. 53 ---
\opage{53} necessity compels the interpreter to reject the old Gospel, and next: what
warrant there is for the reliability of the new? The latter point is easily
settled; the new Gospel is a self-made Gospel, a conjecture that stands today
and falls tomorrow. The first question, on the other hand, requires a more
careful discussion. It is objected against the old Gospel that it does not,
after all, explain the sinless birth, since it leaves the human mother behind as
a sinful link. But surely this difficulty cannot make it necessary to invent a
new Gospel which, instead of removing that link, introduces two sinful links?

``It is not, then,'' say our theologians of the Middle Way, ``for that reason
either that we reject the old form of the Gospel; rather, it is done in
accordance with weighty critical considerations. What objections, for example,
might not be drawn, on the one hand, from the view of the circumstances of
Jesus' birth that prevailed both among his contemporaries and among those
nearest to him (Luc.\ 4, 22; Matth.\ 13, 55; Joh.\ 6, 42 and other passages),
and, on the other hand, from the complete silence with which
this point is passed over both in the evangelical and in the apostolic
demonstration of Jesus' divine-human nature? ---''
But what if these objections, all together, rest on sheer misunderstandings?
What if a believing apprehension of sacred history now shows that it lay in the
nature of the case that those contemporaries had to regard Jesus as Joseph's
son, since Joseph, after all, in accordance with the angel's admonition, took
Mary as his wife; and that those nearest to Jesus could, as little as the other
contemporaries, know anything of a secret which for the persons concerned was so
difficult to communicate that Mary could not even confide it to Joseph? What if,
finally, this wonder is not adduced as proof of Jesus' divine nature, namely,
not in the Gospels,
% --- p. 54 ---
\opage{54} because it is not by telling of wonders but by himself performing wonder-works
that Jesus awakens the recognition of his divine mission; and not in the
apostolic discourses, because these have a different form from the evangelical
narratives: if now, we say, this is how matters stand, has not the faith of the
Gospel reason to complain of such a proceeding?

``Yet what does it help that faith complains, when criticism nevertheless has
truth and science on its side?'' --- Yes! Truth is a sacred power, and the faith
that strives against truth must fall by its own impotence. But to find the truth
one must know the way. General truths force themselves upon the thinker with
necessity; but the truth of the Gospel is a gift of God, which the believer is
to appropriate in freedom. If the test of necessity is now confused with free
appropriation, the confusion is boundless. It is this confusion that betrays the
scientific weakness of the doctrine of the Middle Way. The half-criticism
begins, namely, by setting apart everything that it thinks it can regard as
indifferent and superfluous in relation to the general idea of revelation.
Therefore it sets aside every wonder that cannot be derived with necessity from
the idea. Thus the angelic manifestation to Joseph in a dream is rejected as
unmotivated, since Mary, after all, had had an angel's visit in the waking state;
indeed, even this angel's visit is now, according to the half-criticism's ``new
Gospel,'' just as superfluous as the former. Such economical calculations may
look quite fine at first glance; but let it be well noted that the separating
criticism of the Middle Way theory has thereby become a dissolving criticism.
Once the dissolution has got under way, its scientific momentum cannot be halted
by saying: stop! and meeting it with a remnant of faith. If I have once permitted
myself the arbitrariness of dismissing the angel of the Annunciation, then
% --- p. 55 ---
\opage{55} it is only a new arbitrariness if, in the name of faith, I would begin in my
own fashion to maintain the sinless birth. For the idea as such nothing is
necessary except the idea itself. Let humanity according to its idea be
sinless: the idea of sinlessness I can then perhaps demand with necessity; but
where now is the theologian of the Middle Way who will prove to me that this
idea would dissolve into nothing if Jesus of Nazareth had not been born
sinless? Faith in Jesus' sinlessness stands and falls with childlike faith in
sacred history; but no historical fact can be proved by virtue of pure
necessity. Negative criticism is therefore within its good, scientific right
when, with unbending consistency, it completes the dissolution which the
doctrine of the Middle Way long ago began.

We have followed the investigation through its various instances; we have heard
enough of the mockery of unbelief, of the circumlocutions of doubt and the
critical embarrassments of half-faith, enough to convince us that there is only
one decisive step: either to choose the denial of faith or to remain in faith.

The Believer, who has thus secured himself against every scientific attack, now
turns away from the din of battle, to rest in the Gospel with the quiet joy of
inwardness. The mocker's arrogance shall not wound his feeling, nor the critic's
phantasms dazzle his eye. Only the soft echo of contradiction will make itself
heard yet once more in the reflection in which he strives to become intimate,
truly intimate, with the wonder of the Gospel.

``\emph{In the beginning was the Word, and the Word was with God, and the Word
was God.} And God spoke the word of omnipotence, and behold! a world came to be
through the Word; he spoke the word of election, and behold! a people of God
came to be in the world through the Word.
% --- p. 56 ---
\opage{56} Abraham heard the speaking God and wandered away from the mute gods; Moses
heard the speaking God, and --- Israel became the people of God's own
possession through the Word. But the heart of the people in the house of Israel
was toward the Lord like the carnal heart of a woman, which knows no constancy,
but in weakness and strength, in frailty and defiance, is moved unsteadily
hither and thither. Then the holy word was heard as the bitter word of wrath
that rebukes the refractory, as the crushing word of punishment that chastises
the faithless. And again the heart of the people was, in the elect, like the
heart of a loving woman who repents her guilt, humbles herself in obedience, and
bows in faith and devotion. Then the holy word was heard as a word of mercy that
raises up the fallen, as a word of grace that comforts the mourner, as a word of
love that fills the soul with joy. It was a parable, indeed a dark saying, that
the people was wedded to Jehovah as the woman is wedded to the man, her lord. It
was a discipline of holiness that was to save the elect; a riddle of love that
was to be solved at the last. Never did any word in the world sound so heavy, so
stern, as the word of the Law to Jehovah's people; never was any riddle in the
world solved as the riddle of the covenant in Israel was solved, when the woman
`found favor with God.' And the riddle was solved, and the parable became a
truth, and the Law a Gospel in the angel's greeting:
{\large Hail to thee, thou Favored One; the Lord is with thee, thou blessed among
women.''}%
\footnote{``But the virgin Mary was betrothed to a man named Joseph, of the
house of David, and Joseph misjudged his betrothed wife.'' How the relation
between Mary and Joseph is to be understood more precisely will depend, in the
main, on how one explains to oneself Mary's election. When the Gospel legend, in
keeping with the ascetic tendency and luxuriant belief in miracles of the
ancient Church, represents Mary as a born saint, as the playmate of the angels,
our evangelists can by no means be said to bear the blame for such a
misunderstanding. Luke mentions the betrothal as an actual fact, and Matthew,
who presupposes that the wonder has already taken place with Mary, lays
particular weight on Joseph's natural surprise. That the Free Thinker helps
himself to the apocryphal fabrications can therefore surprise no one who knows
where he is determined to go. The fundamental difference between the canonical
and the apocryphal account, however, by no means lies in the former's admitting
fewer, and the latter more, miraculous occurrences; what is decisive here is
precisely this, that the miracle stories of the legend betray their spurious
origin insofar as their multiplicity serves only to deck out ``the Favored One''
as a fantastic being. With the same right with which the ascetically pious
fancy has annihilated the natural significance of the betrothal by supernatural
preparations, and distorted the original impression of Mary's image by a false
gilding, with the same right a modern fancy might now also take it into its
head to surround the betrothed pair with an erotic nimbus, and to set Joseph and
Mary aglow with the strong passion of the god of love. Instead of the utterly
unsensuous bond between a ninety-year-old graybeard and a young convent maiden
who sits waiting for an extraordinary angel's visit, we should then get a
romantic bond of love, sensuously radiant, like that which entwines a Romeo and
Juliet, an Axel and Valborg; only with the difference that here it is not a
family feud, nor any Black Friar, but an archangel who, as a
\textit{Deus ex machina}, cuts the knot. If the latter poetic invention cannot
be justified with the least semblance of right by the original text of the
Gospel, then the former can just as little; and there is thus no reason to
misjudge the evangelical account out of anxious regard for the extravagances of
the legend. If, then, we keep, as is fitting, to the canonical tradition, the
betrothal retains its full moral validity, and then the supernatural
communication by which Joseph's misunderstanding is removed will also prove to
be grounded in the nature of the case. That Mary gave Joseph her hand clearly
proves that, before the Annunciation occurred, she had no notion that she
herself was ``the Favored One.'' This circumstance, however, is so far from
weakening our conviction of her spiritual preparation that it serves rather to
set the simple faith of her pure heart in a still clearer light. It is
undeniably one of the most beautiful proofs of true humility that she who, in
her deep awareness of sin, her unconditional trust in God's grace and her
infinite devotion to his will, fulfills the inward work of faith, has herself
no inkling of her preeminence over others. Had Mary thought to herself: I dare
not become betrothed to Joseph; I must be prepared for something far higher;
for, should the Lord need a handmaid to accomplish the extraordinary, there will
scarcely be found anyone worthier than I; then such self-complacent reflection
would surely have been enough to make her wholly unworthy.

If it accords both with the letter and with the spirit of the Gospel that the
betrothed were really on the point of consummating their marriage when the
Annunciation occurred, then Mary's silence concerning the wonder that had taken
place becomes intelligible to every believer. Concerning her closer relation to
Joseph, and its significance in the civil and moral world, the angel had, after
all, revealed nothing to her. Should she then break off the union? Should she,
as mother of ``the Son of the Most High,'' stand alone in the world, and expose
him and herself to the mockery of men? Would Joseph understand her if she
confided her secret to him? What should she demand of him? Would he consent to
her demand, even if she really could find out what she should ask of him? That
the critics have hitherto not been properly willing to heed these and similar
difficulties is explicable enough. For when, by virtue of the modern
consciousness, one knows in advance that the wonder is only an imagining, one
passes over such trifles. It is not worth the trouble, after all, to enter
vividly into so absurd a narrative. On the other hand, it furthers the
scientific victory of criticism far more easily if one pretends that an angelic
manifestation was an everyday occurrence; for thereby one opens for oneself the
most convenient way of escaping from the wonder.

Should it, on the other hand, not be altogether unworthy of science to enter
properly into its object, should modern criticism really be able to bring itself
to take a look into the situation of the betrothed at the time, then one will
surely blush to put Mary's silence down, with Strauß, to --- ``the principle of
laziness.''

If the manifestation to Joseph is already sufficiently motivated by regard for
Mary's condition, the communication receives besides its distinctive stamp from
the fact that the angel appears to him in a dream. This form of communication
corresponds entirely to its purpose. Joseph's natural grief over the loss of
Mary sets the man in the right light for us. He feels wounded, and this hurt is
all the more bitter as he at the same time feels what he loses in losing Mary.

He harbors suspicion against her innocence, for she is, after all, as one
fallen; but he stifles this suspicion again: she has, after all, hitherto been
to him a pattern of the purest womanliness. Therefore he dare not even follow
the prompting of jealousy to shame her publicly, but resolves to put her away
privily. To no human being can he confide his sorrow, for she is dear to him
above all; if he must give up faith in her, he must give up faith in humanity.
No friend would understand his words, even if he attempted an explanation; for
there is confusion and contradiction in his own thoughts. Still less can he
approach Mary herself: in what language should he speak to her? Reproach her
with her fall? No! she cannot be guilty! Demand that she justify herself? No!
She cannot possibly justify herself; she must, after all, be guilty! So only
one way out remains to him, namely: to hide his grief, and ``put her away
privily.'' One truly need not exalt the carpenter Joseph of the house of David
into a knight of love of modern reflection in order to explain to oneself his
secret grief. Only grant the man what, according to the Gospel, is surely due to
him at the least, a noble heart, a sound understanding, and the struggle within
him is understood of itself. If Joseph with all his brooding cannot himself
solve the riddle that has been set him by the wonder, then it is by no means
superfluous, but wholly in keeping with the intended purpose, that an angel of
the Lord solves it for him in a dream. Joseph is thus prepared for a
manifestation. His heart demands nothing less than what the angel of the dream
proclaims to him: he demands the wonder, and therefore he also believes in the
wonder. This faith of Joseph's is, however, infinitely different from the loose
superstition of the interpreter of dreams; for it
at once proves its truth in a great moral resolution: he awakes from his dream
to do as the angel of the Lord has commanded him. Thereby Joseph becomes an
evangelical exemplar for all those who with pure and manly self-denial fulfill
the holy law of duty, even when the fulfillment of duty, if it becomes known,
brings upon them only the mockery and misjudgment of the world.}

% --- p. 57 ---
\opage{57} %
``\emph{In the beginning was the Word, and the Word was with God, and the Word
was God.} And God spoke the word of omnipotence, and behold! a world came to be
through the Word; he spoke the word of election, and behold! Israel became the
people of God's own possession through the Word. --- And because the Spirit, the
Holy Spirit, came upon Israel by the Word, and the power of the Most High
overshadowed God's elect through the Word, therefore the people of Israel was,
as in a figure,
% --- p. 58 ---
\opage{58} called the Son of the Most High according to the word of the Lord:
\emph{Israel is my firstborn son} (2 Mos.\ 4, 22). And yet Israel's heart was
hard, and he walked after his own thoughts in a way that was not good. And
Jehovah punished his elect, yea, the Lord chastened his people as a father
chastens his refractory son. But because
% --- p. 59 ---
\opage{59} the Spirit, the Holy Spirit, was nevertheless upon Israel by the Word, and
the power of the Most High overshadowed God's elect through the Word, the most
excellent who were born of the Lord's people were hailed as \emph{children of
the Most High} (Ps.\ 82, 6). And who, then, was more excellent among the elect
than the shepherd from Bethlehem, the singer-king with the
% --- p. 60 ---
\opage{60} %
golden harp, the man after God's own heart? So the word of the Lord came to
David: \emph{Thou art my Son; this day have I begotten thee} (2 Ps.\ 7). --- And
Jehovah God was with David, and David's name was great, `like the name of the
great ones that are on the earth,' and he was the Son of the Most High through
the Spirit. But he was
% --- p. 61 ---
\opage{61} also a son of sin and iniquity through the flesh. And the man after God's own
heart must shudder at his iniquity as no other; the Lord's elect feels God's
wrath as no other; the singer-king with the golden harp is sorrowful, after all,
as no other: he is so pale with grief, `he is
% --- p. 62 ---
\opage{62} so weary with sighing in the long night, he has watered his couch with
tears.' It was then only as in a figure, in a parable, that David was called
\emph{the Son of the Most High}. He received a promise of eternity in the Word,
but the Spirit chastised him in the flesh; --- `he went the way of all the
earth.' But because the Spirit was nevertheless upon Israel, and the power of the
Most High overshadowed God's elect, the word of truth worked in the parable.
Nathan stood before David's throne, and spoke with prophetic voice of the Son of
the Most High who should build the Lord a house, and have the throne of his
kingdom established for ever. \emph{When thy days are fulfilled, and thou restest
with thy fathers, then will I raise up thy seed after thee, which shall proceed
out of thy body, and I will establish his kingdom} (2 Sam.\ 7, 12). --- And who,
now, was more excellent among the elect than Solomon with the treasures of
fortune and the glorious gifts of the spirit? He asked for understanding, and
the Lord God gave him a wise and understanding heart, yea, gave him even that
which he had not asked, namely both riches and honor. All Israel heard King
Solomon's judgment; and they feared before the king's face, for they saw that
the wisdom of God was in him to do judgment. And Solomon became greater than all
the kings of the earth in riches and wisdom. The queen of Sheba came to
Jerusalem to prove him with dark sayings; and there was not a word hidden from
the king which he did not declare to her. The people praised him in their heart
when he turned his face about and blessed all the congregation of Israel; yea,
the people praised his name, for he was a prince of peace: Judah and Israel
dwelt safely, every man under his vine and under his fig tree, from Dan even to
Beersheba, all the days of Solomon. And yet it was only as in a figure, in a
parable, that this son of David could be called \emph{the Son of the Most High}.
The son of wisdom took wisdom in vain, and for the wisdom of
% --- p. 63 ---
\opage{63}%
vanity all things are full of labor, so that no one can utter it: the eye is not
satisfied with seeing, nor the ear filled with hearing; and the Preacher sighs so
heavily: \textbf{Vanity of vanities. All is vanity!} \emph{I the Preacher was
king over Israel in Jerusalem; I saw all the works that are done under the sun,
and behold! all was vanity and vexation of spirit.} No, alas, no! Solomon in all
his glory was not the true Son of God. The power of the Spirit failed him, and
his servants mocked him; the sense of the flesh beguiled him, and his wives
turned his heart. --- Israel is then a riddle to himself: the Spirit is mighty;
but the sense of the flesh is enmity against the Spirit of God, and the Word
cannot become flesh. See Ephraim and Judah, how they hate each other with a
deadly hatred. False spirits and false prophets, foreign enemies and foreign
gods ensnare Israel, consume his strength, and threaten him with fall and ruin.
Where now is \emph{the Son of the Most High}, who was to sit on David's throne
and reign over the house of Jacob for ever? Alas! On the throne sits an Ahaz,
in dread of Pekah and Rezin. It is told to the house of David that the Syrians
rely confidently on Ephraim; and the king's heart trembles, and the heart of his
people, `as the trees of the forest tremble before the wind.' --- But because
the Spirit, the Holy Spirit, was nevertheless upon Israel, and the power of the
Most High overshadowed God's elect, the Word was mighty through the Spirit, and
the Lord's prophets powerful by the strength of the Word; for the Spirit was in
the words of the prophets, and the Word was like `a hammer that cleaves
mountains.' --- The prophet goes to meet the king; Isaiah speaks to Ahaz; thus a
man of faith speaks, in the hour of danger, to the disheartened doubter:
\emph{Take heed, be still; fear not, neither let thy heart be faint} (Jesaj.\ 7,
4). --- Jehovah's people is in the hand of the Almighty. The king in Damascus
cannot
% --- p. 64 ---
\opage{64} subdue us, Ephraim and the son of Remaliah cannot destroy us; only when
Jehovah hisses for the fly that is at the border of the rivers of Egypt, and for
the bee that is in the land of Assur, does ruin come upon us. Assyria cannot
save us; the arm of a man cannot protect us against the wrath of the Holy One.
We are in God's hand, we are humbled under his might, we are raised up by the
strength of his arm, we shall be saved by the wonder of omnipotence. ---
\emph{Therefore Jehovah will give you a sign:} \textbf{Behold, the virgin shall
conceive, and bear a son, and thou shalt call his name Immanuel} (Jes.\ 7, 14).
But the doubter on the throne did not understand it, his servants did not grasp
it. He who was to be as \emph{the Son of the Most High} burned his sons in the
fire after the abominations of the heathen, and plundered the house of the Lord
to buy his salvation from the king of Assyria. And yet --- the prophet does not
dream, he never loses sight of the conditions of the present; but in the midst
of the tribulations of the time a vision of the future passes before his eye,
and he rejoices in the vision, as one who gropes in the night rejoices when the
day dawns; yea, he rejoices in the vision, for the salvation of the people, the
true salvation, appears to him in the vision: \emph{The people that walked in
darkness,} he exclaims, \emph{behold a great light; the light riseth for them
that dwell in the land of the shadow of death.} Isaiah does not dream, he does
not lose sight of the conditions of the present; but out of the vision of the
present a vision of eternity clears itself, and the Son of the Most High stands
before him in the vision: \textbf{A child is born unto us, a son is given unto
us. And the government shall be upon his shoulder; and his name shall be called
Wonder, Counselor, mighty God, Father of eternity, and Prince of Peace} (Jes.\
9, 6).

What a wondrous people this chosen people is, with its kings and prophets!
Behold Israel, then! Far off in
% --- p. 65 ---
\opage{65} history he comes up, a solitary figure, in the mist of myths. His path
crosses the ways of the nations round about: should he then mingle with the
strangers? No! he hates them, they despise him: he would crush all the heathen
if he could. Behold Israel, behold him closely! He is no misty image in gray
antiquity, but a real figure with sharp, marked features. There is light in his
eye when he looks at reality; there is understanding in his brain when he judges
the course of things and calculates conditions and ends. You do not dazzle him
with empty words: he knows the strict law of the natural, he feels the limit of
reality; he demands fact. And yet, how strange! He, precisely he, is a believer,
indeed the only believer in the midst of all the nations, because he is the
only one who truly knows the difference between Creator and creature. He does
not dream himself into faith. No! he is compelled to faith: for only in faith is
his salvation. The God of omnipotence saves him, the God of wonder looses his
chains and bears him on eagles' wings away from the land of bondage. Yes, then
he must surely believe; for he stands, after all, before Mount Sinai, surrounded
by signs and wondrous works. Behold Israel, behold him closely! He is no
marrowless shadow, but a vigorous nature with flesh and blood and a soul full of
lusts and passions. He sees the idol: desire awakens, enjoyment beckons. He sees
the golden calf, its luster enchants him: the eye sparkles, the hand trembles,
the lips worship, for he is almost wild with craving. Then Jehovah's voice
thunders, and the sinner starts up in terror. The altar of the idol falls, the
song of joy falls silent, the dance ends. The disobedient one must now suffer
his punishment, must deny his desire, repent his transgression, and curse the
false gods. Behold Israel when the arm of the Mighty One rests heavily upon him,
when he fights the fight of despair with the bitter enemies, when
% --- p. 66 ---
\opage{66} he cries to heaven for vengeance, but the Lord does not answer him, because
the day of wrath has come: then his courage sinks, then his strength fails; then
the enemies come upon him, and tear him to pieces, and scatter his limbs by the
rivers of Babylon. Then he is as one dead; who would now believe that he could
rise again? But then behold him again, when the Spirit comes upon him, when the
power of the Most High overshadows him! It is a wonder of God. `The dead bones'
gather themselves with flesh and blood; the scattered limbs are joined together,
the dead man breathes, Israel stands up, finds again the way to the land of
promise, and praises Jehovah on Mount Zion. The chosen people is a
world-historical riddle. Of all the other spiritually gifted peoples it holds
true that they begin with mythical nature gods, with glittering dreams of youth,
then pass over gradually from the immediate representations of mythical faith to
a reflective apprehension of the world-idea, and thereby to doubt and unbelief.
Israel, on the other hand, has from the beginning the contradiction within him
that he is at one and the same time a believer and yet a doubter, a zealous
worshiper of God and yet a mocker. He therefore does not advance from faith to
doubt, as from the stage of minority to that of majority, in order to end in a
dissolution; but he goes continually from faith through unbelief back to faith.
Therefore at the end of his course he stands essentially on the same stage as at
the beginning. With all the trials the chosen people has passed through, all the
promises it has received, all the prophecies it has heard, and all the
supplementary provisions by which the Law was gradually augmented, it
comprehends itself and its God as little at the end as at the beginning. Only
this does the believing Israelite understand: that he is to fear God and keep
his commandments, that the people is to be saved by the wonder, and is to await
the time and hour when the Savior, the Son of the Most High, comes.''

% --- p. 67 ---
\opage{67} %
``If I now look upon all this with an earnest and searching eye, my soul gathers
itself in the one decision: that \emph{either} the people of election is a
strange caricature in history's own worldly drama, the wonder an illusion, the
Spirit's dream of the omnipotence of the Word a bold poetic fiction, and the
riddle of the world a natural nodal point in humanity's own essence (and then
the thinker is right when he dissolves all supernatural facts as natural symbols
in the cosmic development of the idea, and then the negative critic is perfectly
right when he in turn dissolves the idea's mythical selfhood into sentimental
fancies, and the defender of half-faith cannot, in the language of science,
answer him one in a thousand); --- \emph{or} else the riddle of the world must be
a mystery of God, the people of election the people of the wonder, and the
visions of the prophets the Spirit's forewarnings of the manifestation of the
eternal Word in the flesh: and then the faith of the Gospel is unconditionally
right when it turns with displeasure away from all obliging half-explanations,
and reduces all conceited objections to silence, in order to hear in peace how
the angel of the Annunciation solves the riddle for the virgin in Nazareth. The
men of half-faith, who here sense something supernatural and yet look anxiously
about for the natural conditions, contradict themselves and know not what they
do: only the angel of the Annunciation, the angel who `stands in the presence of
God' (Luc.\ 1, 19.), knows the solution of the riddle. The angel knows that, as
earth and dust was the natural condition for the creation of the first Adam, so
now the pure soul of a believing woman is the sufficient natural condition for
the conception of the second Adam; the angel knows that the sign of salvation is
`the virgin's son Immanuel'; yes, the angel who stands before God's face, the
angel of the Annunciation, knows that the Word from heaven becomes flesh in
Israel when the Spirit descends with the Creator's power into the virginal heart,
and the God of love accomplishes his wonder.
% --- p. 68 ---
\opage{68} That is why the voice is so clear, the tone so solemn, so gently moved, and
yet so full of authority, when it sounds to Mary: \textbf{The Holy Spirit shall
come upon thee, and the power of the Most High shall overshadow thee; therefore
also the holy thing which shall be born of thee shall be called the Son of
God.''}

``\emph{In the beginning was the Word, and the Word was with God, and the Word
was God.} And the God of the Word spoke to the elect, and Israel heard God's
word, and understood that it was a word of faith, which was to be kept in the
heart and proved in deed. The word of the world, too, will be heard in the
world; the spirit of nature, too, will be glorified in humanity according to its
archetypes. If it is a contest for the prize of beauty, of the wisdom of
thought, of dominion over the world: then Canaan is not the promised land; then
Israel can measure himself neither with the Greek nor the Roman nor the strong
ones of the North. But if the question is of the glory of the Eternal, of the
wisdom of godliness and the struggle of faith for the kingdom of holiness on
earth, then Israel has the Word. If there is to be talk of a human mind that
offers up its flesh and blood, yea, its natural understanding, for the
fulfillment of a divine command; if there is to be talk of the trials in which a
man is tried when he fights the fight of decision with God and his own heart;
if the discourse is to show how a man of the spirit can defy a whole world with
a divine promise, and in the midst of the perils of reality build all his hope
on a wonder of God: then Israel stands alone beyond comparison, then Jehovah's
son has the Word, and all the voices of the nations must be silent when he
speaks. And yet the people of Israel, with the revealed Word, is a riddle to
itself: it hears, and yet does not understand; it sees, and yet does not grasp
--- the mystery of faith. That faith has a power of eternity in it through the
Word must be evident to all who consider its effects in the realm of forces.
But when the carnal
% --- p. 69 ---
\opage{69} mind is kindled by the Spirit, then passion surges, then the real difference
between obedience and disobedience must come forth in visible decision: the Law
demands the outward deed. Only on this condition can the people understand the
greatness of soul of the believer. The struggle of self-denial, the pain of
repentance, the incomprehensible elastic force of self-sacrifice: all must come
to light in such facts that it can become an object of admiration, and pass
over into the history of the people. But that the life of faith can be hidden
in pure inwardness, and yet lift a man above the corruptibility of the outward
world, that the perfect deed can be a meek soul's unconditional devotion to
God's will, that the power of the Most High can be active in the incorruptible
being of a quiet spirit, this would never become truly clear to Jehovah's
people. In this sense it holds that Israel's strong faith is more
outward-historical than inward-psychical, more heroic-masculine than
tender-feminine. Therefore the people of the Law can indeed grasp the beginning
of its faith, but not its consummation. That the election began with the man
who was made perfect in obedience when he prepared to offer God his only son,
this the chosen people could understand, and therefore Israel praised the father
of faith; but that the consummation of faith should come about through the
woman who, in simplicity of heart, gave herself up to the will of the Holy One,
in order to bear God's only-begotten Son, that surpassed sense and
understanding.''

``Not as though Israel would put woman in the shade: no, far from it! Jehovah's
eye rests upon her too; he hears her sigh when she prays humbly; he comforts her
when she speaks out of the `abundance of her complaint'; he blesses her when she
`pours out her heart before his face.' Where faith exercises the highest power,
woman may well, in the moment of decision, step forward and measure herself with
man. Israel knows that the spirit of strength can
% --- p. 70 ---
\opage{70} in the hour of danger rest upon a woman; Israel knows that at times, when the
bold lost courage, when the strong became weak, it has come to pass that the
Lord made the weak strong, and saved his people by a woman's hand. Only say a
word of Deborah under the palm tree, and you shall see how the Hebrew's face
shines: `Deborah's counsel became Barak's victory, the wife of Lapidoth was a
mother in Israel.' Name Judith's name, it kindles a spark in the Jew's eye:
`Judith drew the sword of vengeance in the Assyrian's tent, and saved the honor
of Jerusalem.' Yes, all the world can understand that the people of Israel must
sing the praise of Deborah, must extol an Esther and a Judith.''

``How mighty is Deborah in her lofty, prophetic calm! --- The people is in
bondage under Jabin's hand; for twenty years the ignominy has lasted. Only with
the prophetess on Mount Ephraim do the oppressed find a refuge: she, a woman, is
the only one who judges in Israel. Sisera threatens, and the people cries in
anguish to the Lord, `alas, Jabin's captain has 900 chariots of iron'; but
Deborah is not afraid. Calmly she sits under her palm tree, and judges the
children of Israel. The Spirit stirs in her; she sees, after all, the people's
distress, she knows the enemy's design. And yet the woman sits on her judgment
seat, with the word of the Law in her heart, as firm and calm as the seasoned
helmsman when the sea roars. Then the hour of deliverance strikes: she calls the
son of Abinoam. She marks out the hero's way for him, she raises his courage,
she follows him into battle. There is a certainty in her look, an assurance in
her word, a superiority in her power of action, that only faith can give to one
whom the Lord chooses to save his people.''

``Great and glorious is Judith's strength; she is the glory of her race: all the
passions of the people's faith burn in the Jewess's heart. The Assyrian's soul
was captured by her loveliness, the Persian trembled at
% --- p. 71 ---
\opage{71} her daring, the Mede was sore afraid at her boldness. --- In Judith, militant
Judaism lays bare its very harshest contradiction. If the chosen people is to
endure in the world, then every means of salvation must surely be well-pleasing
to Jehovah; and yet cunning and deceit cannot possibly further the holy purpose.
A woman's highest passion is needed to master this contradiction. That is why
Judith's resolve is a secret she cannot confide even to her own people. None of
the men of Bethulia suspects how she wrestles in prayer, and adjures `the God of
hosts, that he smite the servant of destruction by the deceit of her lips.' No
one grasps why the grieving widow has so suddenly chosen to put off her mourning
garments, and to renew her beauty with fragrant ointment and costly ornaments.
The young men regarded her with silent admiration as she went out of the city
gate; and they looked after her `until she came down the mountain, until she
passed through the valley,' and they could see her no more. No woman dares more
for her people than Judith. Her sacrifice springs from the invincible faith in
Jehovah. She believes that the God who saves her people will also save her
womanly honor; and in the happy outcome she finds her faith confirmed. With the
fullest assurance that the will of the Almighty has revealed itself in the
carrying out of her deed, she climbs the mountain of Bethulia again, and cries
to those who keep watch: \emph{Open, open now the gate! God, our God, is with
us, and still shows strength in Israel!} And when she lifted up the head of the
captain of the host, and cried: \emph{Praise God! Praise, praise, praise God!}
then all the people understood that she was blessed above all women that are on
the earth. And men and women blessed Judith, and they went out to meet her with
garlands and with songs; and she went before all the people, and led the women
% --- p. 72 ---
\opage{72} %
in the dance\footnote{That Judith is apocryphal cannot be unknown to the
Believer.}. Yes, all the world can understand that the people of Israel must
sing the praise of Deborah, must extol an Esther and a Judith.''

``But the virgin in Nazareth! Who knows her? Who speaks a word in her praise?
The eye of history cannot discover her: she hides modestly in the shadow of
family life. No prophetic call has gone out to her: she does not seat herself on
the judgment seat to judge in Israel. No great and mighty thought of
deliverance surges in her mind: her voice is not heard in the streets. The
betrothed one is occupied with the quiet thoughts of her own heart, quietly,
unnoticed, insignificant to the world, like life in Nazareth. And yet! yet the
Gospel testifies: Here is more than Judith! Here is more than Deborah! --- When
woman blazes up in the heat of battle and intervenes actively in history, she is
for a moment stronger than man, because she is placed in a higher tension.
Without wavering, without stumbling, she rushes forward toward her goal, for she
acts on inspiration; or rather, she does not act herself, but is an instrument
for an irresistible power that acts through her. The all-sacrificing passion
that has transformed her being, the all-powerful resolve that has concentrated
her strength, in short: the consecration that endows her with supernatural
strength, is conditioned by the presupposition of faith, and faith in Jehovah has
sufficiently shown what the Spirit can do through a woman. But in all such
scenes woman is really carried away from herself --- beyond her essential natural
limit, instead of being preserved for her original destiny and, in soulful
individuality, opening the hidden riches of her heart. If the flower of
womanliness is thereby hindered in its unfolding, then the essence of faith is
no less so. Religious faith, namely, expresses itself one-sidedly insofar as it
expresses itself only as heroic power of action; for such
% --- p. 73 ---
\opage{73} power of action can, after all, also spring from quite another source. The
opposition between the supernatural and the natural, the holy and the worldly, the
purest self-denial and the most forceful self-will, still lacks its definite
expression. The life of faith therefore demands a quite different and purer
manifestation: that the soul in true virginity may feel itself pervaded by the holy
Word, and see itself in God's own light. In this sense the Gospel testifies of Mary:
here is more than Judith, here is more than Deborah!''

``Only he who understands the consummation of faith in Israel understands Mary; for
what she keeps in her heart is --- faith's own secret. Ingenious theologians awaken a
dim intimation of the natural determination of the Favored One when they let her
disposition express woman's pure receptivity, and see in her the noblest flower on the
tree of life of the Hebrew people. Against this there is, so far, nothing to object:
Mary is certainly a flower of the people, the fairest flower in the garden of
Palestine, the pure heart-shoot of the root of womanhood. But with this mystical
figurative language one still gets no further than the hint that the mother of the
Son of God must, above all women, have been a pure, a lovely, a gracious nature. But
what a riddle! If it is a holy nature that is spoken of, then Mary is of course
sinless; then we confuse the mother with her Son, and fall back into the childish
thought-play of legend. If she is to be a sinful, yet nevertheless pure and guiltless
nature, then we run up against a contradiction that no human understanding seems able
to endure. If, on the other hand, we abandon the dubious game with natural
determinations, and say it in plain words: Mary is neither a sinless nature nor a
being bound in sinfulness, but she is --- \emph{the believing soul}; then the task
changes. For this does not merely mean that the Favored One is more believing than
pious women in
general,
% --- p. 74 ---
\opage{74} or that she has an unusually strong faith at all; but it is expressly
emphasized that she has faith itself as faith, that her true being is one with the
being of faith, that consequently each of her determinations of soul coincides
immediately with a determination of faith. Only now does it become a religious task to
understand who the Favored One is; and for this understanding the believer will strive
to the uttermost, for faith wants to understand itself. `But,' objects the reasoning
understanding, `if it is nothing more than that, we shall soon be done with the
matter. If Mary is neither more nor less than the believing soul in pure simplicity
and simple purity, then anyone who has any faith at all should surely be able to
understand her; the believer is believing, that goes without saying: like is
understood by like.' --- But what a riddle! In the first case understanding became so
difficult that it exceeded human powers; in this latter case it comes so easily that
here, strictly speaking, nothing is understood, since here there is nothing at all to
understand. Perhaps a psychological investigation may help us to a deeper conception;
let us at least try it.''

``In womanly immediacy faith is a feeling, an inwardly strong, but inexplicable,
unutterable feeling. The indeterminate feeling must, however, be more closely
determined by that which is felt; and, since \emph{the believing soul} knows herself
as a member of a sinful world, she must surely, in relation to the God of holiness,
also feel her guilt. Let the psychologist try, then, whether he can describe to us this
indescribable feeling. `Yes, Mary feels her guilt. A child who in all her innocence
heeds no danger, is weighed down by no transgression, she nonetheless perceives an
anguish of soul, a fear and pain that can be fully known only by the penitent sinful
woman. It is the threats of the Law, it is the judgments of the prophets, it is
David's cry of lament that terrify her; she feels it in her heart that
% --- p. 75 ---
\opage{75} \emph{all are gone astray, that there is none that doeth good, not even
one;} yes, she feels it in uncomprehended dread that the God of Israel is a jealous
God, and his wrath a consuming fire; she feels it deeply (for her being is, after all,
feeling); she feels the need of the suffering for the mercy of the Eternal; she feels
the longing of the pious for the deliverance of Israel; a sigh passes through her
soul, a quiet sigh, but so deep and heavy that he who understood this uncomprehended
sigh would have to say to himself: O the melancholy one, she is unhappy!' --- But now
`the believing soul' stands not only in an immediate relation to punishing justice,
but in an equally immediate relation to compassionate love. This feeling too the
psychologist must be able to describe.
`No, Mary is not melancholy, though she sighs in secret over the sin of men. Her sigh
breathes itself away like a child's sigh; the shadow of sorrow passes the innocent one
by: she breathes again as soundly and lightly as one who possesses a blessed hope, an
unnamable joy. The word of promise stands written in the Law, in the Prophets, in the
comforting Psalms; so she writes this word deep in her heart, and the heart
understands the revealed word: here, then, is the unfailing pledge of saving grace.
Yes, she feels it (her being is, after all, feeling) that salvation comes from
\emph{the God of Abraham}, from the God who \emph{showeth strength with his arm, to put
down the mighty from their thrones and exalt the lowly.} She feels it (for her whole
being is, after all, feeling) that \emph{God's mercy endureth from generation to
generation on them that fear him.} For this reason too she is called the Favored One,
because she, above all souls, has felt the divine grace so purely and deeply, so
strongly and inwardly.' --- Thus the psychologist has indicated the two opposite
feelings into which the fundamental religious feeling must necessarily divide; and
Mary must surely know them both from the ground up, since she is --- \emph{the believing
% --- p. 76 ---
\opage{76} soul.} He who cannot understand her dread and pain will never understand
her hope and joy either. Take away the pain, and the trustful security is only a
youthful boldness, only a carnal security without religious depth and holy
earnestness: all her grounds of comfort then vanish like powerless phrases. Take away
the joy, and all her sorrow over the sin of the world and the misery of the people is
only a melancholy reverie; a morbid brooding darkens her being. --- But what a riddle!
Should these opposite feelings perhaps weaken one another mutually, so that the soul
only partly perceived the pain of fear because it now also partly perceived a relief
in hope; then the Favored One would surely languish in faith, and live on in the
enervating indeterminacy of half-feelings? Or should it be possible that `the
believing soul' could in one and the same instant feel such a pain and yet such a
joy, such a fear and yet such a hope? Would not the one of these strong feelings
stifle the other, and the Favored One find herself lost in an uncanny mood, a dull,
oppressive brooding that makes it impossible to breathe? --- Let the opposite feelings,
then, part from one another in momentary breakthroughs; let `the believing soul' be
moved in passionate unrest between both extremes: the indescribable pain of fear, the
unnamable joy of hope; let her sink abruptly into the darkness of sorrow, rise
abruptly into the rapture of joy: what is she then with her strong, overwhelming
feelings? She is a fermenting possibility, capable of venturing everything according
to the conditions of chance and the prompting of the moment; but then she is also in
the power of selfish passions, and cannot be --- the handmaid of the Lord. She is a
heroic possibility that only awaits a sign from fate in order to feel herself capable
of everything, capable of offering herself on the pyre for the salvation of Israel,
capable of a Judith's cunning for the exaltation of Jerusalem; but then she is also a
vehement,
% --- p. 77 ---
\opage{77} restless nature, and then cannot be the one transfigured in faith, not the
one favored by heaven, not the one blessed among women.''

``Religious pain and the inward joy of pious confidence must, since each of these
moods rules immediately, doubtless express themselves in utterly conflicting states:
pain casts down, joy lifts up; pain proclaims the punishing wrath, joy the mild,
compassionate grace. `The believing soul' is thus at strife with herself. This strife
too can be presented with the help of descriptive psychology. The first immediate
breakthrough of pain can very easily bring a youthful heart to lose courage and
despair of God's mercy, precisely because the impression is so strong, and the
inexperienced disposition is overwhelmed by the feeling. (That the soul cast down in
repentance cannot accuse itself of any single transgression does not come into
consideration here at all; what is decisive is that the not-sinless nature feels
itself in its sinfulness over against the Holy One.) This feeling acts crushingly,
since it acts with the whole unweakened force of immediacy. But as despair draws near,
self-collection awakens; the will strives to win a hope, and the bleeding heart
appropriates mercy. Then the light gleams, a ray of God's love scatters the darkness,
and the cast-down soul lifts itself up --- in blessed confidence. The lifting up,
however, is also immediate, and the peace of confidence a feeling hitherto unknown;
therefore the trust is so unbounded, and the joy so excessive. According to youthful
expectation, the hardest strife should thus now be over, and the danger past. But this
hope will disappoint, and this disappointment is all the more painful as it seems to
betray a falling away from blessedness itself. `What can this mean, that the peace of
heaven visits a man, only to depart again and leave impotence and emptiness and doubt
and unrest
behind?
% --- p. 78 ---
\opage{78} Is the heavenly, then, as fickle as the earthly? Was it not God himself who
breathed upon this heart and proclaimed his forgiveness? Has a new guilt perhaps
angered him, so that he now punishes in the coldness of wrath, or is it a presumptuous
imagining that flesh and blood would live together with the Holy? Can the frail one
satisfy justice? Can God's love will to kill a soul? Impossible, impossible!' --- Here
we have pain again. It is difficult to decide how far the following transition from
the renewed pain to the regained joy presupposes a still harder strife than the one
that began with the first breakthrough; but it is certain that faith lives only in
continued strife. It is certain that faith presupposes and brings about the strife in
which the pain of repentance and the joy of reconciliation are continually separated
as conflicting moods, since being cast down always conditions being lifted up; but it
is also certain that faith in turn strives to end the strife.''

``Must the believing soul, then, suffer many a temptation and pass through many
dangers before the right understanding can come; must Mary, then, above all others,
have been tried in hard trials and have fought the long, toilsome fight of inner
experiences with God and herself before she could win understanding, and reach the
stage of consummation, and be greeted by the angel as the Favored One? It would thus,
in this sense, be the fair privilege of old age to possess the true transfiguration
of faith. But what a riddle! It is not Mary who is presented to us as the aged one; it
is Elisabeth. Of Elisabeth it may hold that through strife she works her way upward
step by step toward consummation. Descriptive psychology, which cannot reach Mary, may
for that reason well reach Elisabeth. Now since Elisabeth understands Mary, let us
turn again to the aged one, to see whether under her guidance we can perhaps gain a
new insight into
% --- p. 79 ---
\opage{79} the essence of faith. `What, then, is this wondrous power of soul that
reconciles the pain of repentance with the joy of confidence, so that the two grow
together, so that joy itself is nourished by pain, and pain transforms itself into the
quiet joy of an inward sadness?' Ask the aged one who has fought the good fight, and
the answer will sound short and definite: `it is true humility.' --- The immediate
passing of strife into reconciliation, the liberation of the mind from the selfish in
unconditional submission to the law of holiness, is given with humility; and humility
is, like virginity itself, at one and the same time an ideal determination both of
faith and of womanhood. Elisabeth, who has learned by experience what humility is,
also knows that the many finite humiliations men are wont to pass through in order to
win true humility cannot be excused by the general sinfulness of our nature, but are a
consequence of the hardness of hearts. Therefore she does not boast either of her many
trials, or of her long strife, or of her great age; but she looks up with humble joy
to the Favored One, and testifies for her, and exclaims: \emph{blessed is she that
believed.} If it now stands firm that humility itself can be just as immediate as
faith, and that faith thereby receives the wondrous gift of being able to be both the
strife itself and the reconciliation of the strife, then Elisabeth has solved the
riddle for us, in testifying of Mary: \emph{blessed is she that believed.} --- It is not
the extraordinary gifts of nature, not the rich sum of all womanly perfections, that
we are to admire here; Mary is not called to shine in history, but to be the handmaid
of the Lord. She is not destined to make her fortune in the world, but to find grace
with God. It is, then, only the gift of the one perfection, the gift by which a
woman's soul finds grace with God, the gracious gift of being pure before the Holy One
in humility of heart,
% --- p. 80 ---
\opage{80} of which Elisabeth testifies when she greets the Favored One: \emph{blessed
is she that believed.} Yes, blessed is she that believed! --- Great deeds has Israel's
strong faith accomplished in the hour of need and deliverance to the glory of Jehovah;
but only in Mary's heart does faith work the perfect deed with the wondrous power of
inwardness. For this is the perfect deed: that a virgin woman, whose heart perceives a
world's distress, nonetheless lives in the world simply secure, as one who knows of no
danger; that she, whose soul can feel a world's pain, nonetheless lives in the world
simply glad, as one who does not know what pain means. Bold, glorious words has Israel
spoken in the mighty language of faith; it sounds to us in the cry of pain, in the
praise of blessedness, in the wail of terror, in the persuading strength of
confidence; in the sighs of the suffering, in the jubilation of the victorious, in the
lament of the oppressed, in the loud thanksgiving of the saved, as no other language
in the world does. But the true word of faith, the word that is born deep in the heart
of humanity and rises up out of the quiet ground of the soul, the pure word of faith in
the inward language of devotion to God, is after all Mary's word to the angel:
\textbf{Behold! I am the handmaid of the Lord, be it unto me as thou hast said.''}

\begin{center}
  *\quad*\\[-4pt]
  *
\end{center}

\begingroup\large
``In the beginning was the Word, and the Word was with God, and the Word was God. It
was in the beginning with God. All things were made by it, and without it was not one
thing made of that which is. In it was life, and the life was the light of men, and
the light shone in the darkness, and the darkness comprehended it not (Joh. 1, 1--5).''
\par\endgroup

\begin{center}
  \rule{0.38\textwidth}{0.4pt}
\end{center}
% --- p. 81 ---
\opage{81} %
\section*{III.\\[4pt] The Holy Family.}
\addcontentsline{toc}{section}{III. The Holy Family}

\begin{center}
  \rule{0.08\textwidth}{0.4pt}
\end{center}

It is a law of actuality that everything which lays claim to independence must have
its place. That is why we say of any utterance which on closer testing proves
unjustified that it belongs \emph{nowhere}. The eternal Spirit is indeed omnipresent,
and consequently not bound to any single place; but in his actual difference from the
world God must nevertheless be said to have his own particular place; which is why
religion too, and rightly so, teaches us that God has his place in heaven, that he
dwells up there in the home of eternity and blessedness. The human spirit, although in
its free striving and infinite yearning it likewise transcends every outward
limitation of place, naturally also has its place; for it dwells with men down here on
earth. So one can speak not merely of a sensible relation of place between the visible
heaven and the visible earth; but one can, in a deeper sense, speak of a supersensible
relation of place between the heavenly Spirit and the spirits that live on earth.
Likewise, as regards every nation, one can ask not merely about the place where it
dwells or has dwelt in the outward world; but one can and must also ask about its
spiritual place, i.e.\ about the definite position it occupies in the series of
national spirits. Where the Israelite people in the sensible sense had its place and
its home, anyone knows who has some knowledge of the situation of the lands on the
globe; but when this people's spiritual place is to be more closely determined, faith
begins to speak a language in which worldly understanding cannot rightly find any
meaning.
% --- p. 82 ---
\opage{82} Revelation-faith, namely, asserts ``that the Hebrew national spirit dwelt a
little higher than all the others in the realm of spirits; that it dwelt precisely at
that wondrous place where temporality borders so closely upon eternity that a mortal
ear can distinctly hear God's own voice, and perceive a rushing of the seraph's wing;
that it was at home in that region of the world where the sun stops short, and the
moon stands still at the word of one who prays; in the land where the clouds bear the
cherub's chariot, and a Jacob's ladder reaches from earth to heaven.'' It is now this
spiritual place that has given the sensible one its name, and therefore there is a
land on earth that we call --- \emph{the Holy Land}.

It is a law of actuality that every independent working must have its time. When it
is nevertheless said that God is exalted above all time, that a thousand years are to
him as the day of yesterday that is past; then an essential difference between eternal
and temporal time is thereby expressed; but by no means that the determination of
time itself should fall outside God's whole being and working. A deity that is
indifferent to time can be nothing other than the timeless ground of existence; but
such a ground must, precisely because it is in itself timeless, borrow all
determinations of time from finitude. From this it follows in turn that a timeless God
sinks immediately down into finite time, and becomes a servant of temporality, instead
of mastering all determinations of time and being the Lord of the times for ever. If
now everything, even the God of eternity, must have its time, then the same holds of
actual revelation. ``There was a time when God came down from heaven and took on
visible form; there was a time when the Eternal himself spoke the language of
temporality, and God's angels visited the children of men.'' This beginning is
sufficient to raise a far-reaching dispute between faith and unbelief. --- What
fortunes and
% --- p. 83 ---
\opage{83} vicissitudes the Israelite people, worldly speaking, has passed through in
the course of the times; what strength and weakness it has displayed in happy and
unhappy times, this question can very well be settled before the general tribunal of
science; and anyone who knows history thoroughly is entitled to have a say. But when,
on the other hand, it is a matter of an inquiry into whether these great changes of
the times are to be ascribed to the temporal conditions of the world alone, or are
essentially to be referred to the foreordained counsel of Providence in eternal time,
then human knowledge becomes divided against itself, and hands out weapons to both
sides. It is possible, of course, that the visions of which religious tradition tells
were only such images as are thrown off by a pious, but lively and strongly stirred
imagination; but it is also possible that those divine visions at certain times
presented themselves to the soul's eye with such unconditional necessity and so
decisive a stamp of reality that, by virtue of their supernatural cause, they lay
claim to the same objectivity as that which we ascribe to outward things because these
involuntarily force themselves upon our senses. It is possible, of course, that the
divine word that was spoken in Israel had at bottom no other source than the one that
wells up in every national disposition when the idea of religion stirs mightily within
it. But it is also possible that the all-working God himself once in days of old
shaped his thoughts of eternity in the temporal language of men; that this human word
of God at its times sounded so clear and pure, so distinct and definite to the soul's
ear, that the hearer would as little dare to call it his own word as the pupil who has
heard and understood his master may truthfully say of himself that he is the inventor
of the teaching. Now if it holds of the pupil that he only shames himself when he
denies the master the honor due him, when, instead of pointing
% --- p. 84 ---
\opage{84} back to the origin of the fundamental thought, saying: ``thus speaks the
wise master,'' he only distorts the original word by wishing to set it forth in his
own name, and confuses the primally powerful thoughts by paraphrasing them in his own
feeble, colorless language: then this holds in a far other and higher sense when God
himself has undertaken to be the teacher of men. It is then also possible that Moses
and the prophets could not, and in the holy name of truth dared not either, open their
mouths to the people without beginning with the words: ``Thus saith Jehovah, the God
of Israel.'' Thus possibility stands against possibility on both sides. The doubter
halts and languishes; but the denier of faith and the believer know, each in his own
way, how to make a choice. Revelation has its actuality and its definite times; this
fundamental thought of faith can burst the bounds of physics, as the spark in powder
can burst a rock. ``Revelation has its own times. There was a time when Jehovah was
the guest of Abraham in the groves of Mamre; there was a time when the heaven of the
Spirit opened, and the heavenly tabernacle became the pattern for the earthly; there
was a time when the prophets in Israel were driven by God's Holy Spirit, and could, in
the midst of the manifold misdirections of finitude, rely so surely on its infallible
promptings; there was a time when the Hebrew national spirit, in spite of all sins and
transgressions, was consecrated to carry the eternal Word through the vicissitudes of
the times until the fullness of time.'' It is now this supernatural Word that has
given the natural national character its name, and therefore it was said to the chosen
people that it should be a kingdom of priests, --- \emph{a holy
people}.

It is a law of actuality that everything that would come forth into existence must
have the force to break its way through. When the seed-corn breaks and the foliage
bursts out, nature has
% --- p. 85 ---
\opage{85} spring; when feelings break forth in the ground of the mind, and strong
thoughts make ready for a breakthrough, a spiritual springtime is heralded. The
elements of nature, the various powers of the realm of spirit, wrestle with one
another everywhere, as far as actuality reaches. The new life begins with a break, and
death in turn wrestles with life. And what is revelation with its election but a
break? It is a breakthrough of the holy light that brings Abraham to break the natural
bonds of kinship, and leave the land of his fathers. And as it goes with the
patriarch, so also with the race. Jehovah wrestles with his people. Israel is in more
than one respect the man who --- wrestles with God. Through continual breaks the
election is repeated within the circle of the elect, so that this circle, despite the
multiplication of the race, becomes narrower and narrower. First Ishmael falls away,
as a broken piece, as a branch that is broken off a tree in order to be planted in
another soil; later a new break occurs, by which Esau is pushed aside, so that the
blessing may rest all the more richly upon Jacob. But within the house of Jacob too
the break of separation continues; for in the stem of Jesse the Spirit aims at a
higher breakthrough, and at last idolatrous Israel breaks with Judah, which under
manifold interruptions continues to hope in the promise. What thus happens with the
whole people repeats itself in certain ways with the individual elect: they are not
transfigured in moral perfection; but the divine rays break through the natural soul,
as the rays of the sun through the dark cloud. It is this refraction of revelation in
the dim medium of the natural spirit that gives the people with its elect its peculiar
character. ``The light shineth in the darkness, and the darkness comprehendeth it
not.'' Jehovah's light rises over the children of Israel as the light of the sun over
the mountains of Judea: only a few prominent points sparkle in the brilliance of day;
but
% --- p. 86 ---
\opage{86} between them appear --- deep chasms, nocturnal shadows. --- When it is
seen that the election takes place through a break, the Believer understands that this
is precisely the condition for the Word to become flesh, and appear as an independent
member of the race. ``Had the aim been to assert the general supremacy of the holy Law
over the drives of human nature and its rebellious self-will, then no break in the
life of the race itself would have been needed; then the Almighty could, after all,
have revealed himself at once in the clouds with power and glory, and as the God of
the human race proclaimed his threats and promises to the whole of humanity, instead
of condescending to become a family god, and beginning by calling Abraham aside in all
quietness, and choosing him alone out of the mass. If the aim had been to mix the holy
together with the worldly in an ambiguous unity, then every break of election would
have been superfluous; then the creating Spirit could, after all, have scattered his
sparks among all the national spirits, strewn his seed evenly in all souls, and
thereby brought about a fermentation without goal or limit. But if, on the contrary,
it is the final end of election that God's own Son is to be born in the race as the
son of a man, then a place must be prepared for him in the very midst of the race. All
separations and breakthroughs therefore aim at a last separation, so that the small
chosen circle may form which encloses the purest in the race, the Favored One, in whom
the holy Word itself comes to breakthrough.'' The chosen people thus has the essential
determination of forming the foundation for --- \emph{the Holy Family}.

The Holy Family! This concept is composed of two heterogeneous components, and
therefore requires more careful consideration. If one lays the stress on the adjective
``holy'' alone, then the determination of spirit thereby acquires so one-sided a
preponderance that the natural moment lying in the main word
entirely
% --- p. 87 ---
\opage{87} loses its significance. The expression ``the Holy Family'' then comes
essentially to designate a community of holy members. That these members, moreover,
stand in family relations to one another can be regarded as accidental and
indifferent, since it is not the bond of kinship that here hallows the individuals,
but conversely the holiness of the individuals that gives the family its epithet. But
if the concept of family is volatilized in this way, then it becomes at the same time
impossible to connect any clear concept with the designation ``holy.'' The members of
the family of which we are speaking here, namely Zacharias and Elisabeth, Mary and
Joseph, are, after all, neither \emph{holy} in the lower sense in which, in everyday
use, one carelessly and half mockingly says of an awakened person that he is one of
the holy; nor in the special sense in which the Catholic uses the word when, for
example, he speaks of ``the intercession of the saints''; still less in the very
highest sense, as when we think of the holy God. If, on the other hand, one lays the
stress solely on the main word ``family,'' then the natural determination becomes so
predominant that the epithet ``holy'' must be regarded as referring to a property
derived immediately from the relation of kinship. In such a presupposition, then, lies
the view that the members of the Holy Family should each be able to lay claim to be
honored as holy on account of their supernatural descent. One could thus speak of a
holy family in roughly the same sense as when mythology speaks of lineages descended
immediately from the high gods. But that such a conception of the concept of holiness
stands in decisive contradiction to the idea of revelation needs no further proof. ---
Holy Scripture has not, indeed, expressly given the family circle in which John and
Jesus are born and brought up the name of holy (the designation has passed from the
Catholic Church into the more general
% --- p. 88 ---
\opage{88} language of art); but the Gospel nonetheless places this family in such a
light that on faith's behalf we may very well call it by this significant name, since
we learn from revelation itself how the designation must be explained. The strict
opposition between the natural determination and the determination of spirit has,
namely, taken up into itself the supernatural as an essential middle term. The gift of
nature, which points to the relation of kinship, is at the same time a gift of grace
by virtue of election. Not with every one of Israel's pious women does the wondrous
thing happen as with an Elisabeth; and how could the divine Word become flesh in
Mary's heart, if the supernatural principle had not by virtue of the choosing hallowed
and consecrated the natural ground? In the Holy Family, then, the individuals are of
the noblest race, not merely with respect to the natural relation of kinship, but
because they stand in the very closest kinship with the Holy. Likewise the
supernatural has pervaded the determination of spirit; for hope, faith and self-denial
are indeed inner works of freedom in man himself; but such a hope as the one that
lives in this family circle, such a faith and faithfulness that these individuals,
though frail human beings, can yet be entrusted with the Most Holy, is thinkable only
by virtue of a quite special supernatural assistance. They are hallowed, then, by
God's Spirit, in that each of them hallows and consecrates himself to receive and
guard the Holy. The hidden wonder thus goes before the manifest one, the preparation
before the actual breakthrough. --- The ancients imagined that the giant Atlas bore
heaven on his broad shoulders. This is only a fable. The sensible heaven weighs upon
the earth at no point; indeed it has, strictly speaking, no weight. On the contrary,
one could, understood figuratively, rather say that it is heaven that embraces and
bears the earth, since the globe, after all, floats in the wide space of heaven.
% --- p. 89 ---
\opage{89} But still more wondrous than the old fable of the giant Atlas is the
evangelical narrative of the Virgin Mary, who bears the Savior of the world, the Son of
the Highest, under her heart. For this presupposes, after all, that the spiritual, true
heaven of eternity has at its place and in its time lowered itself to earth and
rested upon a frail woman; and yet the faith of the Gospel holds fast to this truth,
because it knows that salvation lies in the wonder of omnipotence. The ancients
believed that this earth was the sensible center of the universe; scientific
observations have since taught that this is not so. But the faith of the Gospel goes
much further still. It assumes that a lowly, womanly creature was, in the fullness of
time, the center of the whole universe, that all the rays of the eternal light
gathered in her soul as in their natural focus; and yet no ``scientific observation
shall be able to wrest this conviction from faith, because it knows that salvation is
--- the wonder of love. --- By these and similar `strange' assertions the faith of
the Gospel has of course now gradually come again to stand on a very strained footing,
not only with the modern consciousness and the so-called sound human understanding,
but with all the scientific great powers without exception.'' Only one science has
been found willing to take up faith's cause, namely speculative philosophy; but this
coalition too has dissolved, as speculation gradually became unfaithful to faith, and
faith, which soon noticed the deception, has renounced every ambiguous alliance, and
resolved to wage the war with its own weapons.

How the misunderstanding between philosophy and faith arose can, as far as the sacred
history is concerned, be briefly elucidated in this connection. Everything depends,
namely, on how one conceives the relation between the essential
% --- p. 90 ---
\opage{90} preparation of revelation and its actual breakthrough. Speculative
philosophy accentuates the preparation, because it keeps one-sidedly to the idea;
faith, on the other hand, accentuates the break, since it keeps exclusively to the
fact. To accentuate the preparation is, essentially understood, the same as to press
for mediation, to dissolve the wonder as a necessary member in the inner
self-development of the idea, and to hide the supernatural in the higher forms of
natural life. When, then, the Gospel in all simplicity reports the great,
incomprehensible wonder that a woman in Nazareth has become mother to God's
only-begotten Son, who is himself God, and the world bursts out in the thousand-voiced
Impossible! of offense, then speculation steps forward to convince the unbelievers by
initiating them into its divine-human knowledge. --- The extreme terms (extremes) whose
mediation philosophy takes upon itself are, on the one side: ``the eternal God,
exalted above time,'' on the other side: ``finite man in time.'' ``The natural
understanding must doubtless take it for a settled impossibility that God's eternal
Word, which is itself God, should be able to let itself be born of a human being, and
appear on earth in human limitation; but the speculative doctrine of ideas nonetheless
finds itself able to dissolve the hard stumbling-stone into pure rationality. The
infinite and the finite, the eternal and the temporal, the divine and the worldly
prove, namely, to be sides and moments in one and the same fundamental concept: man
here on earth is precisely the essential middle term through which those extreme terms
join together into true and complete unity. When, therefore, Holy Scripture speaks of
God's becoming man in the fullness of time, this is so far from being superstitious,
arbitrary and unreasonable talk that in a deeper sense it is the highest wisdom. The
idea of the God-man is the idea around which the whole universe turns, the idea for
whose
% --- p. 91 ---
\opage{91} actualization the whole of existence has been laid out from the very first
beginning of things. The Son of God must therefore necessarily reveal himself, for his
coming into the world is the very completion of creation; he must necessarily be born
of a woman, for only under this condition can he shoot forth as the flower on the tree
of the life of the race.''

Had philosophy now been able to solve its task, to translate all evangelical
representations into the necessary forms of pure knowledge, to halt all calculations
of probability, and to bring all the contradictions of criticism to silence: then
--- yes, then the faith of the Gospel itself would indeed not have gained a farthing by
the whole translation, but speculation, on the other hand, would have won for itself a
superhuman victory, and thrust religion from the throne. The solution of the task was,
however, impossible. Philosophy, which began by casting away all presuppositions in
order, by the dialectical self-development of thought, to compel the cognition of the
absolute idea, gradually saw itself forced to abate something of its scientific
pretensions. The supernatural facts of the Gospel were all too numerous, and the
individual miracle narratives all too objectionable, to be justified before pure
thought. One had therefore to leave such ``lower components'' to the treatment of
criticism, and apply all one's strength to making comprehensible the great fundamental wonder of Christianity, namely the Son's becoming man. But since the evangelical idea
of incarnation continued to press and weigh upon pure knowledge with the whole weight
of the supernatural, speculation began little by little to look about for positive
points of support, and now wanted to call the general presuppositions of faith to its
aid. But this was directly contrary to the agreement. The task, after all, was to
mediate and justify everything by virtue of pure knowledge; how, then, could absolute
science decently avail itself of a makeshift? A little of faith and a little of pure
knowledge gives an unclear
% --- p. 92 ---
\opage{92} mixture that makes a poor showing, in life as well as in science. Under such
circumstances it was therefore quite natural that the masked speculation took off its
mask, and showed the world what sort of task it could and would solve. This open
declaration had an astonishing effect. The problem of miracle was solved; but the
solution turned out a little differently from what the credulous had expected. The
revealed secret of mediation was, namely, this: that the Christ of the Gospel can be
called the true Messiah only in a very improper sense. ``All the perfections with which
faith has decked out this single individual are in manifest contradiction with the
limited concept of a historical personality, and must, in the name of the idea, be
distributed over the whole race. It is the human race itself that is the true,
sinless God-man; it is the human race's own spirit that, in the steady progress of the
centuries, itself performs the great wonder-works; it is the human race that is its own
Redeemer.''

Abandoned by every other help, the faith of the Gospel must therefore now try to pull
itself together, to see whether by its own powers, or rather by the supernatural
assistance on which it always counts in the hour of decision, it can perhaps defend
itself and its Gospel. This defense can be carried through only by accentuating the
breakthrough, inserting the wonder as middle term between the extremes, and deriving
mediation from the supernatural facts. --- In order, however, to test more closely
what powers stand at faith's disposal in this respect, we shall for a moment make use
of the form of dialogue, and let the Believer \textit{(A)} meet various objections that
can be put forward by a Non-Believer \textit{(B)}. Thus: \textit{B.} When from a
natural-scientific standpoint we consider the starry heaven, and reflect that our
globe is only as a point, a speck in the universe, must it not
% --- p. 93 ---
\opage{93} seem to us highly improbable that the infinite God of reason, who otherwise
never interrupts his laws, should have made precisely this earth the stage for such
special and miraculous manifestations as those described in the Old and the New
Testament? \textit{A.} Yes, certainly! \textit{B.} When, furthermore, from the
general-scientific point of view of criticism, we institute a comparison between the
Bible's miracle narratives and the heathen myths, do we not discover so many striking
analogies that, in order to pass a calm and impartial judgment, we must find it in the
highest degree improbable that the biblical myths should have a greater claim to
historical credibility than the heathen ones? \textit{A.} This too I suppose I must
grant. \textit{B.} Well then! We go a step further. Let us assume, for example, that
what is said in the Annunciation really happened to Mary; I shall then show you what
lies in this representation. It is told (Luk.\ 1, 39) that ``Mary arose in those
days, and went with haste into the mountains to \emph{Juda's
city}.'' Do you not sense what horror lies hidden in these few words?
\textit{A.} What do you mean? \textit{B.} Now pay attention! Perhaps I may succeed in
opening your eyes. --- Immediately before the passage cited (namely v.\ 38) it says
expressly, after all, that ``the angel departed from her.'' Mary thus walks quite
alone on the road over the mountains. Now put yourself vividly into the event. Imagine
that you yourself, together with a companion, met her on the road. When she has passed
you, the companion stops, points at her, and says to you in a low but solemn voice:
``Just look at this woman; on her life depends the salvation of the whole world!''
Would you not be struck dumb with astonishment, not at the woman hurrying up the
mountain path (for she is like many another woman), but at your companion's madness?
The more closely the picture is viewed, the more clearly one discovers the absurdity.
Mary walks alone in the mountains! Suppose
% --- p. 94 ---
\opage{94} %
that a robber attacked and murdered her; then this robber would at the same stroke
murder the salvation of the whole world; suppose that a ravening beast tore her to
pieces; then the ravening beast would bring the salvation of the world to nothing;
suppose that her foot stumbled, and she plunged down the mountainside into the abyss;
then this fall would make the raising up of the world impossible. It is not my
intention to confuse your imagination with a fantastic conjuring show; I am speaking
in earnest the plain truth. If Mary is an actual human being, and the world in which
she moves an actual world, then she too, even in her favored state, is exposed to
accidental misfortunes; for accidents and misfortunes are, once and for all,
inseparable from actual existence. If, then, when she went over the mountains, she
bore Him on whom the salvation of the whole world depends, then there has also been a
time when God gave the salvation of the whole world over as a prey to chance. But is
such a representation not now to be regarded as an utterly irrational and absurd
representation? \textit{A.} Yes, truly, a highly irrational and absurd representation!
\textit{B.} Then you have hereby passed the sentence of condemnation on the old faith
of the Gospel, and I hold you captive in your own words. \textit{A.} Do you really
think so? \textit{B.} The matter is as clear as day. Have you not yourself granted me that the wonders of revelation must, before the general-scientific tribunal of physics
as well as of the criticism of myth, fall away as --- highly improbable, absurd and
irrational? \textit{A.} Yes! \textit{B.} Are you then not now also forced to admit
that the faith whose content, according to the foregoing admissions, is in the very
highest degree improbable, absurd and irrational, thereby necessarily stamps itself as
a faith in the very highest degree absurd, irrational and reprehensible? \textit{A.}
No! That I have no intention of granting you. \textit{B.} That is a caprice, I
suppose? \textit{A.} It is my dead earnest. \textit{B.} Then may we ask for an
explanation.
% --- p. 95 ---
\opage{95} %
\textit{A.} The faith of the Gospel rests in the wonder. \textit{B.} A poor
explanation, at the very moment we have agreed to grant the impossibility of the
wonder-work. \textit{A.} The impossibility of the wonder-work? God preserve me from such an
admission! I have only, in order to show you that I could follow your scientific train
of thought, conceded that the wonders of revelation must seem to the common
consciousness of the world highly improbable, absurd and objectionable. More than that
I have not conceded at all. \textit{B.} But what is the point of this hair-splitting?
Since you can so easily ``follow my train of thought,'' you will surely also be able
to see that the wonder, as a historical fact, must be judged from a
historical-critical standpoint. But now everyone knows that the boundary line between
the most improbable and the untrue, between the most absurd and the impossible, is so
fine that the historical eye cannot discover it, and therefore rightly pays no heed to
it either. Whether, then, I say of the historical wonder-work that it is in the very
highest degree improbable and absurd, or I straightforwardly assert its impossibility,
surely comes to roughly one and the same thing? \textit{A.} Therein you express a
fundamental error which, from the beginning of the modern criticism of the Bible in
the last century down to this day, has brought about the most lamentable confusions
in faith as well as in science. \textit{B.} A bold assertion, which might well stand in
need of a proof. \textit{A.} Let us consider the matter with all deliberation. It lies
in the concept of wonder that, in comparison with the usual phenomena of nature, it
must appear as the most improbable thing one can imagine; were it lost in
probability, it could, after all, be explained in a quite natural way, and
consequently would be no actual wonder. But does this now mean, in other words, that
the wonder is in itself an absurdity, and is therefore to be classed with the kind of
meaningless contradictions that in logic one is accustomed to
% --- p. 96 ---
\opage{96} compare with a round square? \textit{B.} The contradiction is admittedly not
quite so striking, for the doctrine of revelation, after all, calls the supernatural to
its aid; but nonetheless the miracle is still an actual impossibility, since it
interrupts the law of nature, and this, once and for all, cannot be interrupted.
\textit{A.} Surely you will not pass off this assertion as a scientific proof?
\textit{B.} It is, however, to be regarded as the irrefutable principle on which the
consciousness of culture and the whole of modern science rests. \textit{A.}
Why, then, is the wonder an impossibility? \textit{B.} Because the law of nature
cannot be interrupted. \textit{A.} Why can the law of nature not be interrupted?
\textit{B.} Because it would be unnatural. \textit{A.} Now you are forgetting what you
yourself said just now. \textit{B.} No, not exactly because it is plainly unnatural;
but because ... because the supernatural cannot break through the natural. \textit{A.}
So because the wonder is impossible? \textit{B.} Yes! \textit{A.} Summa Summarum:
the law of nature cannot be interrupted, and why not? Because the wonder is
impossible. The wonder is impossible, why? Because the law of nature cannot be
interrupted. \textit{Ergo:} the wonder is impossible because it is impossible. Do you
not now find that this is an exceedingly trustworthy principle on which the
consciousness of culture and the whole of modern science rests?
\textit{B.} No matter! I believe in its infallibility, and no one ... \textit{A.}
Hush! Let us not begin on the confession of faith; otherwise I could easily resort to
the same evasion, and assert that the wonder is possible because it is possible.
\textit{B.} But one can, after all, so vividly sense in oneself and in the whole world
that the wonder must be an absolute impossibility! \textit{A.} Just remember our
``general-scientific standpoint,'' and then have the goodness to answer me quite
definitely on the following questions: Can one prove the impossibility of the wonder,
or can one not? \textit{B.} One cannot. \textit{A.} Can one, from the so-called
improbability of the wonder, prove its
% --- p. 97 ---
\opage{97} unreality? \textit{B.} We did, alas, agree that the wonder must be
improbable. \textit{A.} If this must now necessarily be granted, then it is
surely not so very difficult either to prove that the modern biblical
criticism, which by its ``general-scientific'' calculations of probability
partly attacks, partly defends the Gospel, is in its innermost essence
meaningless. \textit{B.} How so? \textit{A.} The negative criticism, which aims
at annihilating the credibility of the sacred history by demonstrating the
very highest improbability of the wonder-works, constantly ends up, without
itself suspecting it, emphasizing one of their most unmistakable criteria; and
the apologetic criticism of the Middle Way, which exerts itself to establish
their probability, constantly ends up, contrary to its intention, robbing them
of their most utterly indispensable mark. Is this not a perversity that seeks
its like in vain in any other science? \textit{B.} Perhaps! But --- we will not
quarrel about that.

The faith of the Gospel is grounded in the wonder of revelation. It has
nothing to fear from philosophy, since philosophy is just as little able to
prove the impossibility of the wonder as its necessity. Nor has the sacred
history anything to fear from the general criticism of probabilities, whether
this criticism, for the rest, girds itself for defense or for attack. If
natural science asks what moved God to reveal himself on this globe, then the
Believer points confidently to --- the wonder of love. If criticism advances
with the smug assertion that one may not judge the miracle narratives of the
Old and New Testaments by any other rule than the pagan myths, then the
Believer holds fast to --- the wonder of love. If worldly cleverness wants to
make itself important with the accidents and incidents that in the real world
must have threatened the Savior's life, then the Believer merely reminds it
of --- the wonder of love. Such a monotonous repetition
% --- p. 98 ---
\opage{98} of one and the same ``unprovable'' presupposition (and that in a world that
dares to pledge its salvation on God's not being able to perform a single
wonder) sounds, to be sure, so feeble and toneless that many a defender of
the Gospel would doubtless be ashamed of it, and
hasten at once to seek out some scientific proofs for himself. But for just
that reason the method may be quite serviceable for practical exercise; with
the scientific mediation there is, after all, no hurry at all. If, on the
other hand, the Believer is ashamed to hold so simple a discourse in the face
of the strong thinkers and acute critics who know so well how to master Holy
Scripture; if he is afraid to come forward openly with his ``unscientific
presupposition'': then he may just as well surrender sooner as later; for his
opponents already have the victory in hand. If, on the other hand, the
possibility of the wonder is once fully secured against every
general-scientific attack; if its historical form of reality is sufficiently
safeguarded against the insinuations of the criticism of probabilities; if the
faith of the Gospel at last awakens again from its long stupor, and learns to
see that it commits suicide if it does not openly and unreservedly confess
before all the world that it is a faith in wonders, a faith in the
incomprehensible wonder of redeeming love: then the critical justification
comes as of itself, then there is no trouble at all with the scientific
mediation.

And now back to the Holy Family! --- An awakened soul who has taken the faith
in wonders in vain and intoxicated himself on visions and manifestations would
perhaps expect here an altogether superhuman glorification of the elect who
gather about the Savior's cradle. That the plain simplicity of the Gospel in
this respect too forms a laudable contrast to the bombastic exaggerations of
the legends has, however, often
enough been pointed out and acknowledged by sober critics. But, in order that
this verdict may be more than a religious
% --- p. 99 ---
\opage{99} judgment of taste, the supernatural must on the one hand be sharply emphasized,
and on the other hand its real, though sparser, appearance in the canonical
presentation must show itself to be grounded in the nature of the matter.

The Holy Family lives as if on the border between two worlds: eternity casts a
transfiguring glow over its individual figures; and yet it belongs so wholly to
the present that not a single one of the bonds by which it is united with civil
society can be said to break. The powers of heaven stir in its midst; and yet
all that is wonderful takes place so quietly, so hiddenly, that the busy world
around really notices nothing at all. Yes, hidden and secret is the wonder of
love by which the holy is born in a believing soul; hidden and secret too was
the wonder of love by which the holy Word became flesh in the race. --- If we
now consider
the interaction between the two worlds a little more closely, the relation
looks at first glance as if the supernatural world were really the weaker,
because its figures are so ``airy'' and its effects so hidden; whereas the
natural world, on the contrary, seems to have the decisive preponderance of
reality on its side. But if the wonder is not an illusion, then the matter
is of course the reverse; then the supernatural kingdom of course has the
essential supremacy; then it is of course this present world that would have
to dissolve and vanish like a mirage if those powers of eternity broke forth
in might. What is it, then, that holds the mighty ones back, so that, since
they are after all so near us in their invisibility, they do not shake, do not
unsettle this natural
universe, do not derange this order of things? The Believer has but one
answer: God's infinite love! It is this love that gives the wonder its true
evangelical basic character. Yes, that the heavenly visions can show
themselves to a human
% --- p. 100 ---
\opage{100}eye without the natural sight being at once extinguished forever; that
heavenly voices can sound in a human ear without this natural ear being at
once closed forever; that a breakthrough of the supernatural power of the
creating Spirit can take place in a human soul without disturbing its nature
and confounding all its feelings and thoughts: is this not a wonder of love?

\subsection*{1.\\[4pt] Mary's Visit to Elisabeth. --- The Birth of John the Baptist.}
\addcontentsline{toc}{subsection}{1. Mary's Visit to Elisabeth --- The Birth of John the Baptist}

\begin{center}
  {\footnotesize (Luk.\ 1, 39--80.)}
\end{center}

Revelation is no play of shadows that mocks the laws of existence; it is no
game in which the heavenly ones frivolously mingle with the earthly: in all its
utterances the same severity, on all sides the same basic character of divine
earnestness! They are palpable blows with which the supernatural strikes the
elect. \emph{Zacharias has become speechless, the angel has bound his tongue;
Elisabeth keeps herself hidden in her own house, to avoid the eyes of men; even
Mary is so silent and still in her home that she cannot say so much as a word
to Joseph.} There is in this falling silent, in this hiddenness and
reserve, something that involuntarily fills us with dread and, when we dwell
on it, almost makes the blood run cold. It is as if the human had all at once
been overwhelmed by the divine, as if the individuals were so carried away
from this order of things and from one another that they could no longer find
coherence and meaning in the present. But this lostness in the oppressive
darkness points, on the other hand, to a natural underlying working under which
the appropriation takes place. The bound condition passes over into a higher
liberation: a blessed life floods
% --- p. 101 ---
\opage{101}the souls with light and joy. It is this jubilation of liberation and
understanding that resounds in Mary's hymn, in Zacharias's prophetic song of
praise. Both songs of praise are, moreover, so distinctively determined by the
actual surroundings, and each of them answers so entirely to its occasion,
that the interaction between the supernatural and the natural becomes all the
clearer thereby to a believing eye. --- How naturally, how purely humanly it
can be explained that Mary, under the given circumstances, visits her kinswoman
Elisabeth, and on entering her house is received with open arms, with a loving
greeting; how naturally and purely humanly it comes about that the two
friends, each of whom brings with her a hope, a joyful blessing on which the
whole meaning of her life depends, now meet, the young with the old, pour out
their hearts to one another, rejoice, and thank God. But in their greeting and
mutual communication there at once occur such outbursts and utterances as
admit of no natural explanation at all.
``And it came to pass, when Elisabeth heard Mary's salutation, the babe leaped
in her womb; and Elisabeth was filled with the Holy Spirit.'' If this
introduction is to be understood as a later fabrication because it touches on
``the improbable,'' then the expressions of joy that follow are likewise either
enthusiastic outbursts of the women themselves or ingenious inventions of a
poetizing memory. If, on the other hand, we discover how the supernatural still
continues to stir as in a dim background; if we regard its momentary influence
on Elisabeth as a significant after-effect of the stronger breakthroughs that
took place in both annunciations: then Elisabeth's greeting to Mary, ``Blessed
art thou among women, and blessed is the fruit of thy womb,'' has its good and
sufficient ground; the joyful surprise in which she makes known that the
visitor stands high in dignity above the one who
% --- p. 102 ---
\opage{102} receives the visit, ``Whence is this to me, that the mother of my Lord cometh to
me?'', then points to an immediate influence of the Holy Spirit. To Elisabeth's
mood Mary's rapture answers in turn; the favored women understand one another:
their souls melt together in the one blessed fundamental feeling that the time
of salvation has now at last come.

It is the same with the following event. Elisabeth's time was fulfilled, and
she brought forth her son. ``And neighbors and kinsfolk heard that the Lord had
made his mercy great toward her; and they rejoiced with her. And it came to
pass on the eighth day, that they came to circumcise the child; and they called
him after the name of his father, Zacharias.'' A more lifelike and
warm-hearted presentation of a Jewish family scene than this one will surely
not be demanded under the circumstances at hand. But just this kindly and
natural sympathy of the relatives gives occasion for the supernatural to come
into view all the more definitely. The good relatives are taken aback when they
learn that Elisabeth will not let the child be called by the father's name.
``Why, there is no one at all in the family who is called John.'' The choice of
this name therefore seems to them so strange that they nod to the dumb man to
learn his opinion; and when the father, for further confirmation, now writes on
the tablet:
``His name \emph{is} John,'' they marvel greatly. Scarcely are those taking part
prepared for the extraordinary by their astonishment and the presentiments it
must have called forth, before a new breakthrough also takes place. Zacharias is
freed all at once from his dumbness: his mouth is opened, the band of his tongue
is loosed, he speaks and praises God. Yet the supernatural is so absorbed into
the natural that the wonder itself is thereby almost transformed into a
psychological fact. With the wonder's
% --- p. 103 ---
\opage{103}entrance the mood is raised to the point at which the song of praise can begin.
Its content may be taken as a solemn answer to the expectant utterance of those
present: ``What, I wonder, shall become of this child?'' For this too belongs to
the purpose of revelation: that not only does the wonderful take place, but the
understanding of what has taken place is at once introduced and prepared. ``And
fear came upon all that dwelt round about them, and all these things were
noised abroad throughout all the hill country of Judaea. And all they that
heard it laid it up in their hearts.''

Before we go further, however, in the religious development of the narrative's
content, we will stop here for a few moments to hear the objections of
criticism.

\emph{First Objection.} ``Even if one is willing to grant the religious
significance of the wonder to a certain degree, it still by no means follows
from this that the Evangelists have rendered the individual wonderful
occurrences, with all their attendant circumstances, in complete agreement with
the actual truth. A cursory examination will convince us at once of the
contrary. In order that Joseph shall not leave Mary, the angel must, of course,
reveal the secret to him. But from whose mouth can Joseph, before the dream
occurred, have received that ominous information about Mary's pregnancy?
Matthew acts as if it were nothing, and thereby betrays that he confuses his
own notion of Joseph with Joseph's actual position. Further: Elisabeth's
greeting presupposes, of course, that she already knows beforehand how matters
stand with the entering Mary. But where does she have this knowledge from? The
leaping babe can hardly have told her anything; a special revelation is not
mentioned, and just as little is there any trace of the news being
communicated to her by the natural way. Luke pays no heed at all to this point,
and thereby shows that he
% --- p. 104 ---
\opage{104}confuses his own notion with the actual event. But the greatest difficulty
still remains. From what sources have the Evangelists drawn these accounts? Had
Zacharias perhaps written down what happened to him and his wife? Had Mary
perhaps left something in writing, or had she perhaps later told what she knew
to one or another among the Apostles? To whom, in that case, would Mary have
confided such things sooner than to John? But how, then, is it to be explained
that John in his Gospel pays no heed at all to these birth miracles? These
difficulties could be multiplied with many similar ones. Everyone who will not
harden himself against the truth, and deliberately shut his eyes to the most
utterly plain hints of criticism, will at least see himself compelled to grant
this: that the whole cycle of family scenes preceding the Baptist's appearance
presents itself as pictures in the \textit{Laterna magica} of memory: fleeting,
indefinite family legends which, for the glorification of Jesus' birth and
childhood, were no doubt reworked at a later time, but still not sufficiently
worked together!''

\emph{Reply.} It cannot have been the intention of the Evangelists to enumerate
everything that took place in the Holy Family;
only the most remarkable scenes and their most characteristic main features are
brought out. As for the missing intermediate links, it is just as easy for the
Believer, who dares to take the supernatural proclamations into the reckoning,
to invent a connection for himself according to his way of seeing, as it is
easy for the denier of faith, for whom the supernatural stands in the way
everywhere, to invent for himself as many absurdities as he himself sees fit.
Such subjective hypotheses are like soap bubbles: they glitter and burst.
Instead of engaging in questions that cannot be answered, the Believer had
best let the hidden be hidden, the indefinite
indefinite,
% --- p. 105 ---
\opage{105} and say quite simply: ``No one knows!'' Thus: Who gave Joseph the first news
of Mary's pregnancy? ``No one knows!'' In what language did the Holy Spirit
speak to Elisabeth of Mary's election? ``No one knows!'' From whose mouth did
the Evangelists hear the family histories they tell us in this Division? ``No
one truly knows!'' This may then serve still further as proof that the faith of
the Gospel is and must be a matter of freedom, which cannot possibly be forced
out by historical proofs. If I am disposed toward the denial of faith, I shall
not lack the means to dissolve the sacred history into sheer contradictions;
if, on the other hand, I am fully gripped by the spirit of the Gospel, and have
learned to see what the wonder means, then I shall also find in the hints of
Scripture means enough to demonstrate the inner agreement of the sacred history
with itself, both in the whole and in all its essential members and parts. It
is by no means criticism that turns the scale when it is a matter of faith or
unbelief; on the contrary, it is the difference between faith and unbelief that
determines and guides criticism. The apologist's exertion to refute the
historical-critical objections of the denier of faith is wasted labor. Only when
the negative criticism makes a show of entering into the believing conception,
and then with all sorts of distortions seeks to bring this conception into
contradiction with itself, should the faith of the Gospel make reply; for the
impossible and what is in itself absurd I cannot believe.

\emph{Second Objection.} ``It is the unmistakable task of faith to make the
sacred history as real, as historical, as possible. If unbelieving criticism
did not keep on with its teasing contradictions, the Believer would certainly
never think of subjecting the biblical narratives to any strict scientific
examination. Enamored of the supernatural facts, he naturally has
% --- p. 106 ---
\opage{106} so much to do with the heavenly visions, the inexhaustible prophecies and the
manifold mysteries of divine wisdom that there is scarcely time left to engage
in the laborious investigations of doubt. When, nevertheless, the more recent
biblical criticism has at length
brought matters to the point that even the most ecclesiastical theologian dares
not deny that the narratives in this first Division in particular rest on a
floating tradition, and that the events consequently are one and all wrapped
in the veil of legend, then this concession must quite simply be regarded as
extorted. That faith nonetheless
makes a virtue of necessity and appeals to free appropriation shows itself
plainly enough as an attempt to hide its embarrassment. It lies, of course, in
the nature of revelation that the faith of the Gospel must stand and fall with
the sacred facts: only he who accepts the fact can acknowledge the idea, only
he who holds fast to the letter can be sure of the spirit. If it is now granted
that the narratives in question all lie in the ambiguity of legend, then faith
has lost its proper foothold. If one harbors even the very faintest suspicion
that the individual features are possibly not in perfect word-for-word
agreement with the actual events, then one has at the same time given up what
dare not be given up, the evangelical form of reality, and must then also
resign oneself to deriving the significance of the fact from the idea. But with
this turn the Believer, before he himself knows it, is in the midst of the
labyrinths of philosophy and criticism. As long as the faith of the Gospel can
refrain from thinking of any discrepancy whatsoever between the reality of the
event and the form of the presentation, it can keep doubt out well enough; but
as soon as it has dealings with legend, the possibility of discrepancy is
posited. Let one make this possibility as narrow as one likes;
% --- p. 107 ---
\opage{107} it is always a crack through which doubt steals into the believing
consciousness. In short: if the faith of the Gospel would avoid the reef of
philosophy by clinging to the letter, it runs aground on historical criticism;
if it tries to escape criticism by taking refuge in the idea and the spirit, it
drowns in philosophical skepsis. He who in this strait thinks he can help faith
by recommending the Middle Way to it must certainly be a clever man, so clever
that he does not even heed what it is that is at issue here. To follow the
Middle Way is namely the same as to presuppose a merely \emph{approximate}
agreement between the real and the form; but then of course philosophy comes
at once with its ideas, and criticism with its legends, and then --- yes, then
the faith of the Gospel meets with a sudden dissolution. The faith of the Gospel
must surely have an inkling of this hidden self-contradiction; for it certainly
philosophizes to a certain extent, and tolerates criticism up to a certain
point, but soon becomes passionate and defiant when science presses in upon it
to help it to a higher self-consciousness.''

\emph{Reply.} On the preserve of legend criticism plays its freest game; as the
boy shakes his kaleidoscope, it need only shake its hypotheses, and at once the
most utterly peculiar figures form. The ordinary criticism of the
understanding blurs the supernatural features of the events, and shows us the
individual fragments in the ordinary woodcut of reality,
without the luster of the idea. The philosophical criticism of the idea
transforms the natural as well as the supernatural components into poetic
moments in the vividness of the presentation, and exhibits a series of mythical
aerial pictures that have nothing to do with reality. Both methods naturally
lay claim to the very strictest scientific rigor; but ``strict science'' is
nonetheless not able to reconcile
% --- p. 108 ---
\opage{108} them with each other. The criticism of the understanding, which tries its hand
at the ``natural explanation,'' sees quite well that one cannot without further
ado give up a historical tradition because it contains legendary features, and
therefore calls the thoroughgoing criticism of the idea \emph{fantastic}. This
charge of unscientific procedure the philosophical critic now pays back from his
standpoint with ample surplus. For he proves how uncritically the natural
interpretation goes to work in imagining ``that it has separated the mythical
and poetic from an interpretation when it plucks off some twigs and blossoms of
this shoot, but leaves the mythical root
untouched as the purely historical.'' Therefore he calls such a
half-criticism outright \emph{void of ideas and narrow-minded}. This disagreement
is by no means accidental or merely provisional; it has its deeper ground in
the matter itself. The thoroughgoing criticism of the idea is really fantastic,
since in principle it undermines all historical faith: were an event to be
regarded as unhistorical merely because it illustrates a great idea to a
striking degree, how many real events would then have to be erased from
history! The ordinary criticism of the understanding is in truth narrow-minded,
inasmuch as it cannot even notice in itself that it must kill the idea in
order to save the fact. With the disappearance of the supernatural, revelation
itself disappears. If, then, the supernatural components are subtracted as
supposed additions or fabrications, the natural skeleton must fall away from
the sacred history; for histories of revelation that reveal nothing are among
the manifest contradictions. Before strict science undertakes to help faith
to ``self-consciousness,'' it would surely be quite expedient if it tried to
help itself to a somewhat clearer consciousness of the significance of the
sacred legends. --- This consciousness faith has in itself with true and
original certainty. The Believer knows that the
% --- p. 109 ---
\opage{109} spirit of legend in the Gospel is not a spirit of semblance and illusion, but
the holy and infallible Spirit of truth; herewith every doubt as to the
credibility of the accounts is absolutely excluded. Many conjectures may well
arise about the course the tradition has taken before it at last came to rest
in our present account; but if the Spirit of truth has followed it on all its
detours, then we may also rely on its having been safeguarded against every
aberration. Many opinions may well form about the Spirit's influence on its
chosen instruments; but if the wonder is not a fancy, then the Spirit of truth
has fully understood how to express the divine content in the human form, so
that every falsification of the revealed facts is and remains
unthinkable. Many disputes may well arise about the precise relation between
the immediate events and their evangelical presentation; but if it stands
unshakably firm that the same Spirit that worked in the event has also been
at work in the presentation, then the reader of course has the very same
assurance as an eyewitness; for no eyewitness would have been able to find
meaning in those events had he not seen everything in the light of the Spirit
and heard everything with a spiritual ear. The Believer may therefore, in this
respect, undertake whatever critical reflection he pleases: he always comes
nonetheless to the result that the form of the Gospel is the only true and
absolutely satisfying form, the form in which the Word has transfigured the
wonder, and the Spirit penetrated the letter.

\emph{Third Objection.} ``When a man has once taken it into his head that he
will see, hear, feel, think or speak nothing except in the name of faith, it is,
to be sure, a quite natural consequence that no criticism can embarrass him.
One cannot, after all, be both in faith and outside faith at once, cannot
% --- p. 110 ---
\opage{110} in one and the same breath affirm and deny faith. Hence it is also of little
use to engage with the kind of people whose whole acuteness has entered so
exclusively into the service of faith. Since they always lie in conflict with
sound reason, they are at times even quite inventive, and know how to turn and
twist their assertions until these at last acquire even a certain semblance of
scientific rigor. The objection now to be put forward is therefore not
calculated to bring faith to reason either, but has only the purpose of
enlightening the unbiased as to what sort of trivialities the Believers foist
upon the Spirit of truth. According to the presupposition of the faith of the
Gospel, one must of course assume that the hymns put into the mouths of Mary
and Zacharias are immediate outpourings of the Holy Spirit. But under such a
presupposition one will then also excuse the critic when he finds it striking
that an inspiration which flowed so immediately from a divine source did not
turn out more original but, to recall a word of Strauß, shows itself so
bedecked with reminiscences of the Old Testament.''

\emph{Reply.} So: the hymns lack poetic originality! Quite right. Mary is no
bel esprit, and Zacharias no poet. If it is perhaps the standard of theater
criticism that is to be applied to the evangelical scenes, then we must at once
grant the critic that he is right, and only help him to express himself with a
little more definiteness. ``The two scenes in Zacharias's house, namely Mary's
visit and the neighbors' congratulations on the occasion of the son's birth, are
naturally designed to motivate a couple of religious-lyrical impromptus. But in
this respect the Holy Spirit can hardly be said to have been equal to its task.
Elisabeth gets
only as far as some rather far-fetched outbursts of joy, and Mary's whole
discourse is second-hand. The
% --- p. 111 ---
\opage{111} in itself quite beautiful beginning, `My soul doth magnify the Lord, etc.,'
betrays plainly enough that the idea is borrowed from Hannah's song of praise
(1 Sam.\ 2, 1). Insignificant deviations from the original, which occur here
and there, are compilations from prophets and psalms. Indeed, even the one
strophe in which the favored one depicts her own exaltation (V. 48), `Behold,
from henceforth all generations shall call me blessed,' is a pure echo of
Elisabeth's exclamation (V. 45), `And blessed is she that believed.' The whole
attempt thus turns out rather poorly: not an original thought, not a single turn
of phrase that betrays any real talent! The same holds of Zacharias's part.
With the exception of some very obscure and tasteless allusions to the names of
the persons concerned (namely, by the word `mercy' (V. 72) there is an allusion
to the name `John'; by the expression `to think on,' that is, to remember
(V. 72), to --- `Zacharias'; and by the words `the oath which he sware' (V. 73),
to --- `Elisabeth'), the whole long general introduction (V. 68--75) is kept in
so loose a connection with the situation that it might just as well suit any
other edifying occasion. When at last the song of praise reaches the turning
point where the father's joy should break forth in all its strength, `And thou,
little child, etc.' (V. 76), the author gets no further than a repetition of
what we have long since heard from the prophet and the angel.'' --- Let this,
then, be the sorry outcome of the negative criticism.

If we have thus, to the best of our ability, tried to enter into the modern
consciousness's way of thinking, then ``strict science'' will surely not object
either to the faith of the Gospel now also applying its standard and attempting
a conception of the hymns at hand that answers to its concept of the Holy
Spirit. Good heavens! The world is always so very fond of the
% --- p. 112 ---
\opage{112}talents and the original ideas, and thinks so little of the thousands who,
without being talents and without being able to grasp the original ideas, must
nevertheless also live a human life by the daily bread of the spirit. We do not
at all doubt that the modern consciousness would surely become believing if only
it could soon get a quite new, original revelation; but what we do very much
doubt is that the originality or non-originality of these hymns could in this
respect make the very least difference either way. What would it have helped,
had the Holy Spirit taken thought to insert a couple more original songs of
praise into the New Testament, when the New Testament is itself now already so
old that its ideas, for that reason alone, can no longer surprise by their
novelty or win favor by their originality? But is it not a peculiar thing,
after all, with these many original revelations in which the wisdom of the age
is continually born only to die anew, and the word blossoms only to fade again?
The game of surprise that the World-Spirit, under all its brilliant disguises,
plays with its awakened children is evidently grounded in a vain use of the
word of life and truth. The world demands that the new spirit must speak with
new tongues, and in this it is not wrong either: the world is just as clever
and shrewd when it is a matter of demanding and requiring as it is
unreasonable and foolish when it comes to receiving and using. In the frivolous
use of the spiritual word lies the vanity. The genius who brings the new wisdom
to light is spiritually lost in his original production: instead of presenting
what he has himself actually lived through, he lives through, properly speaking,
only the ideal pain and joy of the presentation. The awakened of the age, who
feel themselves called to acknowledge and spread the original wisdom, then set
themselves in motion and now begin in their fashion to work and labor; to that
extent there is, to be sure, striving
% --- p. 113 ---
\opage{113} and laboring enough in the service of the idea. But it goes with the spreading
and appropriating activity as with the producing: all would rejoice in the
rising idea, all are intent on world-transformations on the grand scale; but no
individual finds time to make experiments on the small scale, as for example by
first ordering his own life according to the rule of wisdom. That is why the
original forms of revelation, whose watchword is surprising genius and whose
essential aim is the enjoyment of ideas, change so quickly. The more
surprising the form, the greater the momentary enjoyment; and the more
momentary the enjoyment, the
more quickly the slackness sets in again that expects a new surprise. In the end
the brilliant production doubtless comes into so enormous a disproportion to
the enjoyment that the world-consciousness devours more original ideas in a day
than all the geniuses of the world can produce in a generation. If this is now
how matters stand, then it is surely very good that the Holy Spirit has not
spent its inexhaustible creative power on outbidding the forming and poetizing
genius, but has gone the opposite way, and has laid down its divine fullness of
wisdom in such great, simple fundamental words as can, with healthy frugality,
be heard again and again without surprising changes of form, since in their
first original form they are destined to be --- the daily bread of all souls.

\emph{My soul doth magnify the Lord, and my spirit rejoiceth in God my Savior.}
As many as have ears for the favored one's song hear in these words the pure,
clear echo of revelation's own unchangeable fundamental word, of the word of
eternity of which the promise confirms that it shall endure when heaven and
earth pass away. Or should these words perhaps also be perishable, like a
man's words, or should they perhaps
% --- p. 114 ---
\opage{114}ever need an essential change? What change, then, would one propose? Are they
not intelligible enough in themselves? Why, any child can understand them, so
clear and plain is their meaning! Or are they perhaps not unintelligible enough
in themselves? If anyone will really set about exploring their secret, behold,
then it is as if a living spirit stirred in them. They then keep an eye, as it
were, on the investigator, from which side he is coming, so that he shall not
take them by surprise; they watch closely what presuppositions and what
disposition he brings with him, so that he shall not take them in vain; they
close ranks, so to speak, the nearer the interpreter advances, and at last form
an insurmountable dividing wall between faith and unbelief. The unbeliever
therefore, with all his cleverness and shrewdness, cannot coax a single
explanation out of them, still less press into their innermost depths. In vain
he summons all his power of thought: the immediate intelligibility of the words
bars him from every deeper understanding. Here there is at bottom nothing at
all to think, and therefore thought must stand still. Here no comparison is of
use; for even the bitterest pain, in a natural heart, is not so different from
the highest joy as this joy in turn is different from the favored one's joy in
God. Here there is no other way out: either these words express an empty fancy,
and then all further investigation is superfluous; or they presuppose a
wonder (a wonder of the soul), and then every explanation is impossible. With
faith, on the other hand, understanding is given. This is so incontestably
certain that the one can even be measured directly by the other: faith by joy,
and joy by faith. If faith is false, then joy in God is false too; if faith is
weak, then joy is weak too; if faith is strong and living, then joy in God is no
less so. As long as a man lives in faith, he also experiences joy; even if he is
chastised by God
% --- p. 115 ---
\opage{115} in the pain of repentance, the pain nonetheless testifies that the fire of the
Holy One burns in him, that God himself stands quite near him and holds him
fast, and this certainty is surely a higher joy. So then the words of the song
of praise rest unchangeable in their fundamental word, until this world itself
perishes in the last change. --- But is it now also certain that these words
can endure when heaven and earth pass away? Behold, eternity presses on; it
gives portents in the word, and reminds us that all that is finite comes to an
end. But eternity is not, for all that, untimely importunate; it knows how to
give time, to bide its time; it gives change time, so that the natural heart
may at last change after all and grasp the unchangeable. As long as
changeability lasts, everything is veiled and mixed: sorrow mingles with joy,
and joy with sorrow; the sorrow of the world mingles with godly sorrow, and
with joy in God there is so easily, oh so very easily, mixed the insidious
self-love's own vain joy. But eternity presses on, that the mixed may clear
itself and the heterogeneous separate. Only the soul's purest joy, joy in God,
contains, of course, no mixture. How, then, should the words of the song of
praise be able to dissolve and pass away, when they contain no mixture at all?
--- The temporal time of change is for choosing freedom at the same time a
time of deliberation. Eternity gives its choosers time for deliberation;
therefore its words fall differently according to the occasion of the times.
There is a word, ``He that is not with me is against me''; it
alternates with another word, ``He that is not against me is with me.'' The
former, the severe word, points toward the eternal; the latter, the milder word,
takes the temporal into account; for temporal time is the time of deliberation.
Deliberation belongs to finitude, and must have an end (a time for deliberation
that never ceases, because with all one's deliberating one can never get
through deliberating, an eternal time for deliberation, is a contradiction).
But when
% --- p. 116 ---
\opage{116} the time of deliberation ceases, then surely the mixed feeling, the mixed will,
the mixed sorrow and the mixed joy must cease too. That sparing word, ``He that
is not against me is with me,'' then vanishes in the last severe word, which
sounds into eternity like a word of doomsday: ``He that is not with me is against
me.'' But now it is surely impossible that any soul should be able to be with
God without at the same time magnifying God and rejoicing in God. The last
condition is therefore this: He who is not with God in the perfect joy of
eternity must wage an eternal strife against God, and have joy for his enemy.
--- It is a horror to think of an eternal damnation and an eternal torment; but
it is an infinite comfort that the judging righteousness reveals the divine
grace. The great question that on the Day of Judgment shall pass through all
consciences may then well be interpreted thus: ``Do you now grasp, O human soul,
that you were originally created for joy; do you now understand that it was to
this goal that all the unrest of finitude, all wants and sorrows, all pains and
all sufferings and all the tears of longing pointed, as to the last and final
end: have you now at last learned what it means to rejoice in God? Here no
halfness counts, no evasive ambiguity; here the question is Yes or No!
Blessedness lies in a Yes, unblessedness in a No!'' --- If this is now how
matters stand, and if we rightly consider that it really must be so, then there
can be no talk whatever of changing or improving the words of the song of
praise, since these words contain, of course, the clearest answer to the highest
question. At the entrance to the kingdom of blessedness, then, he ``who finished
the course and kept the faith'' shall have no need either to ask the Spirit for
a new word, as though the old one were already antiquated and would not quite
suit so solemn an occasion; but it will be enough, yes, to perfection perfectly
enough, when
% --- p. 117 ---
\opage{117} the strife is ended, when the judgment is ended, and heaven opens its gates, to
strike up Mary's song:

\emph{My soul doth magnify the Lord, and my spirit rejoiceth in God my Savior.}

We dwell on the beginning of the song of praise; for this beginning is not
merely a part of a whole, but itself a whole; is not merely a single song of
praise among several others, but the song of praise itself in all evangelical
songs of praise, the song of praise of the redeemed in time and eternity. True
enough! it is short, very short, this song of praise, it has but few words; but
in these few words the Gospel of blessedness is enclosed as in its fundamental
word. He who understands this word understands all the songs of praise of
faith; he who misunderstands it misunderstands them all. But in order that
anyone may be able to understand it, the fundamental word must prove its power
in a wonder of the soul; for true joy in God is a true, an indescribable wonder. He who misjudges the wonder misjudges the word as well, and when the
fundamental word disappears, then it is surely no wonder that the individual
words collapse into dead sentences and empty phrases. The charge of criticism
about lack of originality therefore falls as of itself. One clings to sentences
and phrases. If one finds similar sentences in the Old Testament, then it must
naturally be from there that they are borrowed. Now the verdict is pronounced
in a hurry. ``If the phrases of the song of praise are borrowed from the Old
Testament, then of course they are not original; if they are not original, then
of course they are not primordial either; if they are not primordial, well,
then of course they are spiritless, and the whole thing comes down to a feeble
imitation.'' --- Whether, however, the formal criticism does not blunder in
thus confusing the concepts of originality and primordiality, we will try to
elucidate from another point of view. Originality is a concept of art,
primordiality a concept of nature. Let us for once cast a glance at nature.
Without
% --- p. 118 ---
\opage{118} primordiality no nature, be it divine or worldly: every production of nature
must necessarily bear the stamp of primordiality; but it does not yet follow
from this that it is original. The original tolerates no repetition, but
nature engages in repetitions to infinity.

We remember very well what may be said, for that matter, about the unlike in
all that is like, about the difference within the homogeneous, how the one leaf
is not like the other, the one drop of water not quite like the other. But this
difference has properly nothing at all to do with originality. Every leaf on
the same tree, indeed on all trees of the same kind, repeats the same basic
form, without this uniformity in the least degree impairing the impression of
primordiality. Even if the foliage in a following spring took on exactly the
selfsame forms as in a preceding one, the spring would not lose the very least
of its primordial beauty. Indeed, one would hardly even become aware of the
uniformity (for it undeniably takes a strong memory to be able to remember how
the leaves of the trees sit and form themselves one year after another). With
productions of art it is another matter. Here one shuns uniform repetition,
because it always comes at the expense of primordiality. --- If criticism now
takes it into its head to judge the works of nature by the same presuppositions
as those of art, one notices the error at once. For as long as the instinct of
nature works equally immediately in all productions, one has primordiality in
every repetition, and does not ask about the original at all. Or would one not
smile if a critic scorned the nightingale's song because it now did not at all
``turn out original, but showed itself so bedecked with reminiscences of that
old, long-familiar nightingale song, of which one must surely be tired, since
one had heard it so often''? Only when
% --- p. 119 ---
\opage{119}repetition rests on a free imitation must the determination of originality
step in to establish primordiality. Originality is derived primordiality,
primordiality in creative imitation, but not the creative faculty's own
primordial basic determination. The genius must work his way through to
original presentation; true originality is his goal, the form of his highest
manifestation. But the Creator's Spirit, the Unoriginated, in whom and out of
whom all things originate, is highly exalted above the limitation of
originality.

To make a quite considerable contribution to the general confusion of language
of criticism, one would only need to deliver a eulogy on the Creator's
originality: God the Father would then, in so witty a form of speech, become
a pure genius, indeed an enormous genius, the genius \textit{par excellence},
the Son his genial-original self-revelation, and the God-man Christ equally
genial and equally original in word and deed. ---
The confusion that proceeds from this is many-sided enough. One confuses the
holy with the profane, the supernatural with the natural, production of nature
with production of art, the matter itself with the idea, reality with
possibility, what has been lived through with its presentation. When all
boundaries are thus made fluid, then of course every definite presupposition
becomes fluid as well. A confusion of concepts is a supple presupposition.
Modern criticism thus has the great advantage of being able to prepare for
itself the most utterly convenient presuppositions, and yet boast of mastering
all presuppositions and of being ``presuppositionless.'' To be
presuppositionless is here, though, pretty much the same as to begin with loose
presuppositions. Should this assertion perhaps seem somewhat exaggerated to one
or another? Nothing is gained by exaggeration! We will therefore test the
confusion of concepts complained of in the example that lies nearest to our
evangelical task.

When there is talk of rejoicing, the joy may well be of
% --- p. 120 ---
\opage{120} very different kinds, sensuous, spiritual, worldly, religious, etc.; but
whatever joy one may name, it lies in its nature partly to be primordial,
partly to express itself in a distinctive way. If we now consider in the first
place the purely natural or worldly joy, we discover at once a difference
between joy itself and the ideal presentation of joy. It is one thing to be
glad and to express one's joy involuntarily in words and gestures, another to
present the expressions of joy in ideal imitation: it is one thing to sing for
joy, another to sing of joy. The glad man who involuntarily seeks an expression
for his joy is not always a poet; and even if he were, he is now so glad, so
glad at heart, that he does not think of composing at all, but only of living
through the joy. On the other hand, the poet is not always glad either as he
strives to conjure up the image of the glad man (how many a poet has not
moistened his images of joy with the tears of want). If they are now, each for
himself, the glad man and the poet, to be subjected to a criticism, how
different must not the presupposition then be? The happy man who lives through
joy gives his heart vent and seeks an expression for his feelings. Natural
feeling demands a natural expression; that goes without saying. But the stress
falls, however, on the mood, not on the expression; on joy itself, not on the
particular presentation of the idea of joy; on primordiality, not on
originality. The word streams from the overflowing heart, and no one thinks of
holding fast the fleeting word as long as he is listening to the glad man.
Whether he who is filled with joy cannot at once find the right word, whether he
is awkward and repeats the same outbursts, makes no difference at all. The more
repetitions, the more heartfelt the joy. The glad soul is far too glad to think
of the original; therefore he shall also be allowed to choose his words as they
may best fall. Are
% --- p. 121 ---
\opage{121} %
they merely ordinary turns of phrase? Why, they become expressive by the mere
intonation! Are they merely dead sentences? Why, he makes them living in the
inwardness of his feeling! Are they perhaps long-familiar expressions borrowed from
others? What does he care about that --- they become his true property in the
originality of joy! --- It is otherwise with the poet. He does not give us joy
itself; but he renders the idea of joy: the idea is everything through its form,
and therefore original presentation is the highest thing. The poet is not bound to
actuality, and precisely for that reason he must make artistic use of possibility.
When feeling surges, when the waters of language are stirred, images and thoughts
stream in upon him; then enthusiasm must prove its power in sobriety. In the midst
of the tumult of moods the poet's spirit must sit calm and keep its overview; for an
account must be rendered for every word, even the most insignificant. If the poet
loses the ideal equilibrium, the presentation fails too: one can hardly, after all,
at one and the same time both be drunk oneself and yet also \emph{play} the drunkard.
--- Now if the presuppositions that suit artistic production are confused with
those that hold for natural utterance, how unclear and meaningless the judgment must
then become!

But the confusion of concepts goes further still. For so long as the question
concerns only worldly joy, it may in many cases be doubtful whether the immediate
enjoyment of life is higher than the enjoyment of the idea, the thing itself higher
than its presentation, or the reverse; but when the talk is of religious joy, the
doubt is removed. No one will be able to deny that the soul's purest joy, joy in God
(provided that such a joy is really given), must be worth infinitely more than all
its ideal presentations, however full of genius and original these may otherwise be.
The attack of criticism is therefore altogether unwarranted. That Mary knew Hannah's
hymn,
% --- p. 122 ---
\opage{122} and perhaps many times edified herself by it, that Zacharias searched in
the Law and the Prophets, that both thus have the most essential forms and
expressions for their pious feelings and conceptions from the Old Testament --- who
will doubt this? The opposite would indeed be altogether inexplicable to us. But the
point is that here neither can nor should any original presentation be thought of,
any more than any arbitrary imitation. The whole emphasis rests, on the contrary, on
immediate originality. The Spirit breathes upon the heart, and the right expression
comes of itself. --- Only a befuddled criticism can wish to see here a mechanical
composition of borrowed sentences. The soul that ``rejoices in God its Savior''
truly has no need to strain its memory in order to search out a phrase: it has in
itself what it seeks, or rather, it does not seek at all, it finds, finds what the
heart gives, finds the simple words of its joy, as the bird finds its song when the
sun rises.

If we now consider the hymns before us, each by itself, in its wholeness, it
becomes clear that their inner distinctiveness really proceeds from the particular
position and state of soul of those who speak.

Mary, who has heard the angel's message, experiences its power, grasps its meaning.
The outward manifestation is thereby, so to speak, translated into an inward
revelation. All that the Favored One inwardly experiences she utters in the hymn.
Only such a form of communication suits such a content. The hymn itself thereby
acquires a two-sided significance: an unconditional one, namely, insofar as Mary
speaks from the prompting of the Spirit, as for herself alone; and a conditional
one, insofar as she at the same time answers Elisabeth's greeting and communicates
herself to her. As regards the secret of the Annunciation, there can be no thought
here of any private conversation
% --- p. 123 ---
\opage{123} between the two friends. Only in the highest religious mood could they
utter the unutterable. Elisabeth is no curious woman, that she should permit
herself to question ``the mother of her Lord'' closely, in order to learn more
precisely how the angel looked when he came in, how distinctly he spoke, how
strongly his voice rang, and in what manner or by what way the heavenly one vanished
again. Had Zacharias's wife begun with such questions, scarcely any answer would
have followed. But now, when Aaron's tried daughter, ``herself filled with the Holy
Spirit,'' strikes the fundamental word in her expressive greeting: ``Blessed is she
that believed; for it shall be fulfilled which was told her from the Lord,'' ---
Mary can answer in the same language and pour out her heart in the song of praise.
--- But from yet another side this hymn stands in a remarkable relation to the
manifestation.

The angel had indeed expressly spoken to Mary of the Son's royal dignity, saying:
``The Lord God shall give unto him the throne of his father David. And he shall
reign over the house of Jacob for ever, and of his kingdom there shall be no end
(Luk.\ 1, 32--33).'' To become mother to such a son must surely, in the worldly
sense, be regarded as an indescribable good fortune. ``Yes, the woman who receives
such an election truly has something to rejoice over!'' And yet, how remarkable:
these words of the angel do not re-echo with any particular emphasis at all in the
song of praise. No! Mary does not think at all of making her fortune in the world.
Whether she is to see her son on David's throne, and herself sit at his right hand
or his left, or whether she is perhaps to stand beneath his cross, and a sword of
sorrow is to pierce through her soul --- of that the Favored One as yet knows
nothing, nothing at all. --- No sidelong glance disturbs the pious thoughts, no
worldly hopes are allowed to mingle in her pure joy. This alone she knows
% --- p. 124 ---
\opage{124} with the certainty of faith: that Israel through her now receives a
Savior by God's mercy, that the Lord has regarded the lowliness of his handmaiden,
and that therefore all generations shall call her blessed. If it were of any use to
bring proofs against those who always boast of bringing proofs against the Gospel,
we would need only point to these wonderfully moving words: \emph{Behold, from
henceforth all generations shall call me blessed.} --- Or perhaps these words too
have no originality? Then tell us, pray, where they are borrowed from. What other
woman has ever permitted herself to use such language? True enough! The presumption
of fanaticism is unfortunately not unknown in the world; but who is so presumptuous
that he dares to think he hears here the effervescence of a fanatical woman? In the
outburst of rapture the Favored One is as one transfigured, and her voice as a
spirit-voice from the home of the blessed; but she is pure of all overstrained
notions, free of all infatuating fancies. In the humility of her heart she has not
forgotten the lowliness of the handmaiden, not forgotten that she is still in the
corruptible, not forgotten what the actual world looks like. The picture of the
world indeed shows itself to her with all its well-known features: ``Here are the
proud, who mock God in the thoughts of their heart; here the mighty sit on their
thrones and oppress the lowly; here the rich gorge themselves in excess, while the
naked freeze and the hungry sigh for bread.'' Thus, just thus, does the world look;
and Mary is not mistaken. But whence shall help now come? Does she perhaps call down
divine vengeance upon the ungodly, does she in an inflamed state of mind threaten
with the Son's crushing wrath, does she already behold in the spirit how the mighty
ruler shall go forth to battle, to shed the blood of his enemies and tread upon the
necks of the vanquished? No, far from it! Not by a single word are we reminded of
those expiatory sacrifices of retribution with which self-seeking
% --- p. 125 ---
\opage{125} fanaticism is only too eager to sate its greedy passion. --- In all
meekness Mary utters the assurance that salvation comes from him ``that is mighty,
whose name is holy, and whose mercy is from generation to generation on them that
fear him.'' It is the saving and redeeming God who penetrates the world with the
word of power, and himself overturns all worldly relations, without letting himself
be hindered by any finite barrier. The defiant cannot withstand; for he \emph{worketh
might with his arm, and scattereth them that are proud in the thought of their
hearts.} The tyrants cannot withstand; in vain do they arm themselves; for the Holy
One tolerates no injustice: he \emph{putteth down the mighty from their seats, and
exalteth the lowly.} But the needy and helpless may confidently hope in him whose
righteousness is pure mercy: \emph{the hungry he filleth with good gifts, and the
rich he sendeth empty away.} This, when we consider the matter in general, is no
new discovery at all: every believer knows how it is to be understood, and Mary knew
it beforehand in faith. But because the light of faith rises in joy, it all becomes
new. In understanding the wonder of salvation Mary understands everything; for she
understands with the heart. Therefore her spirit rejoices in God, therefore she is
now blessed in her joy, therefore she knows with the undeceivable certainty of
humility that all generations shall call her blessed.

When we find the tone of the following hymn somewhat subdued, we do not mean by this
that its content, in and for itself, suffers from any essential deficiency, but only
that the impression is particularly conditioned by the presuppositions under which
Zacharias utters his joy. Here is not that soulful inwardness, that bliss of
rapture, found in the preceding hymn. The beginning itself: ``Blessed be the Lord,
the God of Israel, that he hath visited and
% --- p. 126 ---
\opage{126} redeemed his people,'' already indicates the difference. Zacharias has
``been filled with the Holy Spirit,'' so that he must bear witness to the
fulfillment of the promises; but he is not so filled with God that he can live, move
and breathe in God, as Mary can. His position is, in relation to hers, the position
of the observer and spectator. His hymn really begins where Mary's ends. Whereas she,
the Favored One, immediately grasps the present and finds the fulfillment in her own
heart, the old man's gaze, on the contrary, seeks back to the historical past of the
people, and rests upon the house of David, upon the prophets, which have been since
of old. Zacharias is not in such a way a center of salvation that he should think
first of himself; but it is his inward joy that he is saved by the salvation of his
people. ``Security against the violence of enemies, the worship of God in holiness
and righteousness'' are the benefits that fill the aged man with gratitude and make
him forget the shortness of his life and the bitterness of death. For his own house
he asks nothing; he does not think that the newborn son shall one day bring honor to
the family and preserve his name in the remembrance of later generations. The little
child is a delight to the old man's heart, his soul --- clings to the boy; but he
does not at all think of keeping him for himself; he gives him away, as, absorbed in
beholding him, he gives him his blessing: \emph{And thou, little child, shalt be
called the prophet of the Highest; for thou shalt go before the face of the Lord to
prepare his ways.} The child is thus not chosen to be the Lord himself, but the
servant who shall give the Lord's people ``knowledge of salvation by the remission
of their sins.'' In the conclusion of the hymn, as well as in the beginning
(V.\ 66), we hear an echo of Mary's song of praise. ``The light from on high'' has
already visited Israel (V.\ 78), ``to shine upon them that sit in
% --- p. 127 ---
\opage{127} darkness and in the shadow of death, to guide our feet into the way of
peace'' (V.\ 79).

Thus Zacharias stands in the evening of his life at the boundary where the old ends
and the new begins; thus the family spirit that receives into its bosom the prophet
of the Highest is itself transfigured in the Holy Spirit; and thus we also
understand the words of the conclusion (V.\ 80): ``The child grew, and waxed strong
in spirit, and was in the deserts till the day he presented himself to Israel.''

If we now look back on these two scenes in the holy family, we can convince
ourselves from yet another side that the facts of revelation exist only for
believers. Their whole significance rests, as we have seen, exclusively on this:
that the holy joy was not only actually experienced but also has inner truth in
itself, indeed infinite validity. But this cannot be established by any historical
criticism whatsoever. Suppose even that a critic could take Mary and Zacharias
themselves into examination and find them willing to give every possible
explanation, we would still be no nearer, since the critic would not merely ask them
``what they believe they have experienced,'' but --- ``whether they have now also
really experienced it''; and such a separation is in this case a sheer impossibility
for the persons concerned. No critical inquisition, even if the events took place
right before its eyes, would be able to obtain any information other than that
contained in the Gospel narratives. Only the Spirit itself can give fully valid
testimony to their inner truth. If criticism will not believe these narratives for
the sake of the Spirit that speaks through them, then it would be of no use either
if Mary and Zacharias rose from the dead to attest their historical actuality. Mary
rejoiced because she believed; Zacharias rejoiced because he believed: whoever
% --- p. 128 ---
\opage{128} would share their inward joy in God, let him give up all the doubts of
criticism, and come hither to share their faith.

\subsection*{2.\\[4pt] The Birth of Jesus. --- The Shepherds in Bethlehem.}
\addcontentsline{toc}{subsection}{2. The Birth of Jesus --- The Shepherds in Bethlehem}

\begin{center}
  {\footnotesize Gospel, Luk.\ 2 K.\ 1--20 V}
\end{center}

\begingroup\large
And it came to pass in those same days, that there went out a decree from Caesar
Augustus, that all the world should be registered for taxation. This registration
was the very first, which took place when Cyrenius was governor of Syria; and they
all went to be registered for taxation, every one to his own city. Then Joseph went
up out of Galilee, out of the city of Nazareth, into Judea, unto the city of David,
which is called Bethlehem (because he was of the house and lineage of David), that
he should be registered for taxation with Mary his espoused wife, being great with
child. And it came to pass, while they were there, that the time was fulfilled that
she should be delivered. And she brought forth her firstborn son, and wrapped him in
swaddling clothes, and laid him in a manger; for they had no room in the inn. And
there were shepherds in the same country, lying in the field and keeping watch by
night over their flock. And lo, the angel of the Lord stood before them, and the
glory of the Lord shone over them, and they were sore afraid. But the angel said
unto them: fear not! for behold, I proclaim unto you a great joy, which shall come
to all the people; for unto you is born this day a Savior, which is the Lord Christ,
in the city of David. And this shall be a sign unto you: ye shall find a child
wrapped in swaddling clothes, lying in a manger. And suddenly there was with the
angel a multitude of the heavenly host, praising God and saying: Glory
% --- p. 129 ---
\opage{129} to God in the highest! and peace on earth! and in men good pleasure!

And it came to pass, when the angels were gone away from the shepherds into heaven,
the shepherds said one to another: let us now go even unto Bethlehem, that we too
may see this thing which is come to pass, which the Lord hath now revealed unto us.
And they came with haste, and found both Mary and Joseph, and the child lying in the
manger. And when they had seen it, they proclaimed everywhere the word which was
told them concerning this child, and all they that heard it wondered at those things
which the shepherds told them. But Mary kept all these words, and pondered them in
her heart. And the shepherds returned again, glorifying and praising God for all the
things that they had heard and seen, as it was told unto them.
\par\endgroup

\begingroup\small
\begin{verse}
``Every one whom the world has forsaken,\\
We bring comfort in the lonely night.\\
Every one who must suffer all the day,\\
We may console at the midnight hour.\\
The hour the Savior was born on earth,\\
A comfort was born for the bitterest pain.\\
A comfort for him who hopes and believes,\\
An echo of the choir of the blessed,\\
Descends then into the human heart.''
\end{verse}
\endgroup

In these beautiful, touching strains the poet has interpreted for us ``the angels'
song of praise,'' and has thereby attuned us to an inward appropriation of the
gospel of the blessed that is proclaimed to the shepherds in the still night. But
the whole Gospel the poet has not explained. Such an attempt he surely has not dared
to venture. The poet is indeed an artist, and the true artist has a deep reverence
for every one of the
% --- p. 130 ---
\opage{130} manifestations of the divine, and therefore never takes more than what is
given him, lest by self-willed presumption he violate the holy and be punished by the
offended spirit. The various arts must therefore, in relation to the many-sided
manifestation of the idea, hold out their hands to one another: one and the same
divine bond encompasses and connects them among themselves, as many as they are.
--- So now the painter too comes along, and hears the Gospel, and receives the
inspiration, and paints us a picture, a lovely picture. It is the picture with the
little child in the manger, on whom all eyes rest. ``The Virgin Mary, the poor and
yet so rich mother, sits near the manger and nourishes her soul --- on beholding the
child. Beside her stands Joseph, the faithful friend, so inwardly still: his eye
watches over the child. And the shepherds draw near, one with another, one after
another, and their eyes seek toward the manger to find the child. And there is a
fading darkness that withdraws in the picture, and there is a victorious light that
transfigures the picture. It is the victory of light over darkness; yes, the light
must conquer: see, it shines around the child in the manger!'' --- Thus the
beautiful art of colors, too, gets its share in the presentation of the Gospel; but
it does not get the whole Gospel. Here there is still something the artist dare not
touch; it is Caesar Augustus, who registers all the world for taxation. Not that the
artist would not indeed dare to paint a Caesar Augustus; but --- he dare not paint
the great emperor in the same picture in which he paints for us the little child.
Yes, even if he painted the emperor in all his glory in a picture by itself, he
would still not be able to paint for us all the bustle that arose in towns and
cities when the decree went out to register all the world for taxation. All the
world! What artist is capable of painting all the world? Therefore, if anyone
demands to know more about Caesar Augustus, how great his power has
% --- p. 131 ---
\opage{131} been, how far his empire extended, and how many peoples he could register
for taxation, we must refer him to another artist and bid him consult world history.
--- Historical presentation, too, is indeed an art. But world history, which knows
so much to tell about the emperor in Rome, has but very little, almost nothing, to
tell us about the child in Bethlehem. The long dispute about the census itself is so
intricate and learned that the simple cannot follow its course; but the Gospel is
indeed for the simple in faith. With few and simple words the Evangelist touches on
the historical turning point, and at once we receive an impression of the greatness
and power of the world empire. For this a child can understand: that this Augustus
must be a great and mighty emperor, since he can let such a decree go out, that all
the world be registered for taxation. With few and simple words the Evangelist
relates that Joseph, together with Mary his betrothed wife, went from Nazareth to
Bethlehem because he was a humble man who had to obey the emperor's strict decree.
But this the simple one can understand: that besides Caesar Augustus there must yet
be another Lord who has command over all the world: for Christ, the son of David,
had to be born in Bethlehem, the city of David. In plain words the Evangelist tells
us that Mary's firstborn was laid in a manger because there was no room in the inn;
but --- the heavenly hosts indeed bear witness to the shepherds in the field that the
pomp of the earth and the glory of the world are only emptiness and vanity, a nothing
before Him, the King of glory, who rests in the manger. --- As all these main
features now gather for the beholder into one picture that unites the earthly with
the heavenly, as all these strong contrasts of existence break upon one another and
shine through one another in the inner infinity of the picture, how must not all
human art of
% --- p. 132 ---
\opage{132} presentation feel itself a poor, vain piecework over against this depth of
the riches of revealed truth and grace! ---

``Truth and grace?'' whispers a spirit of doubt, and negative criticism begins again
to raise its voice. ``Of historical truth there is but little here in this
narrative; but, it shall not be denied, all the more of edifying illusion. --- In
order that the Jesus child may be born in Bethlehem, nothing less is required than
that Caesar Augustus must register the whole world for taxation! Should anyone now
think that for the attainment of so reasonable an end altogether unreasonably much
was nevertheless demanded both of Caesar Augustus and of the whole world, he need
only reflect a moment and see that the Gospel census really first takes place at the
time when Cyrenius (Quirinius) was governor of Syria. This discovery will at once
reconcile the critical connoisseur to the Evangelist's `credible account'; for from
this it is surely evident enough that in such credible accounts it does not,
strictly speaking, matter at all what actually happens here on earth, so long as one
gets all the more precise a description of the remarkable things that take place in
heaven. For according to the testimony of history, it is not Quirinius who was
praeses of Syria in the last years of the reign of Herod the Great, but Sentius
Saturninus. After him followed Quintilius Varus, and only a long time after Herod's
death did Quirinius take up the administration of Syria.''

--- So great, so enormously great is this critical stumbling-block when the eye of
unbelief really gets to see it: the longer one looks at it, the larger it becomes!
--- So small, so insignificant and small is this critical stumbling-block when it is
seen with a believing eye: the longer one looks at it, the smaller it becomes! --- In
the Gospel's original text it is not written, after all, that the whole world was
actually taxed when Christ was born; but it stands thus: \emph{a decree went out
% --- p. 133 ---
\opage{133} from Caesar Augustus that all the world should be entered on the rolls}
(in order, that is, thereafter to carry out a census and, if that were then possible!
to collect the tax). Had it now been Luke's intention to acquaint us in detail with
the many tax histories of that time, he would certainly have expressed himself more
precisely, thus: ``This enumeration of the people was the first step toward the first
Roman levy of taxes in Judaea, a levy of taxes which, however, only came about later,
namely ten years afterward, when Cyrenius was governor of Syria.'' But as it is,
since the Evangelist does not wish to lead us to study the system of taxation in
Rome, but rather to proclaim to all the world the glad tidings of the Savior's birth
in Bethlehem, he is, as regards the levy of taxes itself, well entitled to choose a
briefer mode of expression; and this all the more since the historical emphasis must
here necessarily fall on the imperial \emph{decree}, not on the more particular
execution of this decree.

``But,'' --- continues modern criticism --- ``even if one is willing to let the
taxation go as best it may, what then becomes of the angels and the heavenly hosts?
Let us for once hear a serious word from David Strauß: `First it may fairly be asked,
what purpose the angelic manifestation is to serve? The nearest answer is: to make
the birth of Jesus known. But it is made so little known thereby that the Magi
afterward bring the first news of the newly born King of the Jews to Jerusalem, which
lies so near, and that in the later history there is no trace at all of such an
event at Jesus' birth. If, then, the purpose of this extraordinary event cannot have
been to make Jesus' birth known in wider circles, one must presumably assume with
Schleiermacher that its aim was to act immediately upon the shepherds themselves.
But in that case
% --- p. 134 ---
\opage{134} one would have to presuppose, with the same author, that the shepherds
had, like that Simeon, been especially filled with messianic expectations, and that
God wished to reward and strengthen this pious faith of theirs by that vision. But
neither does the narrative report anything of such a disposition in the shepherds,
nor is any lasting effect upon them remarked; rather, according to the whole
presentation, nothing concerning the shepherds seems to have been the aim of the
vision, but merely the glorification and making known of Jesus' birth as Messiah. But
since the latter, as we have already remarked, was not attained, and the former, as
an empty piece of jugglery, was no aim worthy of God, this circumstance, even
considered by itself, places no insignificant obstacle in the way of the
supernaturalist interpretation of this story.'\,''

Without at once invoking the faith of the Gospel, we may in the present case freely
turn to any person who has some sense and feeling for the appearances of the divine
in this world of the senses, and ask him whether this Gospel, with Caesar Augustus
and the child in the manger, with the shepherds in the field and the angels' choir
--- whether this Gospel, we say, must not so appeal to his better nature, and move
the heart so wondrously and blissfully, that he might involuntarily sigh: ``God
grant that it were truth;'' but --- ``alas, this vision is surely all too blissful to
be true!'' --- The blissful has no truth, and the true no bliss: here is a feeling of
the secret pain of existence. Therefore man has a natural dread of the eternal riddle
of cold truth; therefore one thinks with horror of that audacious youth who by the
light of the moon crept into the temple to lift the veil from the goddess's face.
``The blissful has no truth, and what has no truth in itself can
% --- p. 135 ---
\opage{135} %
never become actual.'' This is one of the comfortless principles on which the modern
consciousness relies; therefore negative criticism can hardly conceal its
unscientific indignation, although with apparent impartiality it puts the
``scientific'' question: ``what purpose, then, should the whole singing angel choir
serve?'' The actual has no bliss, and the blissful no actuality: is it not so? this
thought is like a gnawing worm that gives the soul no rest until a person resolves
either wholly to deny the deeper demands of the spirit, or else to bow in humility
under the wisdom of the Almighty and accept the wonder. Before it comes that far,
however, it is natural enough that he whose better nature does not permit him to
defy the spirit of bliss, but whose worldly culture nonetheless in so many ways sets
itself against the immediate actuality of the wonder, should make many attempts to
come to terms with the law of worldliness and reconcile himself with the ``idea'' of
revelation. --- As an example of such an attempt we may here adduce the aesthetic
interpretation of the Gospel accounts.

``What a barren and fruitless dispute, after all, is the heated dispute over the
credibility of revelation! There stand these `orthodox,' zealous for the literal
expression, and unable to see the simple truth for sheer wonders. Everything is
supposed to have the same actuality, both the natural and the supernatural; as if a
certain sum of miraculous phenomena, for all their fleetingness, could really
constitute the essential and abiding foundation of faith. This willful and
narrow-minded clinging to the conspicuous and accidental in turn provokes the
one-sided critics, who with all their sharpness of sight still cannot catch sight of
the thing itself, because their understanding is so occupied with discovering the
exegetical arbitrariness of the supernaturalists, and
% --- p. 136 ---
\opage{136} chasing after historical and metaphysical contradictions in every
direction. This dispute is as intolerable to sound understanding as it is revolting
to pious feeling. The error lies in this, that both parties dispute about an
actuality whose proper meaning they have neglected to make clear to themselves. ---
The angels who appeared to the shepherds were hardly such external figures as anyone
else could also have seen with his natural eyes, had he at the right time been
walking over the fields outside Bethlehem; the brightness that shone round the
shepherds would surely have been just as little visible as the tones they heard would
have been audible to any outsider, even if he had happened to be in their vicinity.
On the contrary, the whole manifestation must be assumed to have taken place within
the pious men themselves; and when this is assumed and understood as it ought to be
understood, negative criticism can no longer take offense at this fact.''

``This fact? But then there is no fact here at all!'' --- ``Objections of this
nature prove precisely on how low a level the dispute keeps itself. Everything that
has no external, tangible actuality is regarded as unactual; just as if there were no
inward actuality at all. Is not the essence of truth invisible, then? Must not the
figures of imagination be seen with the inward eye? Would anyone dare to assert that
the visions which in the happy moment of inspiration rise before the poet's inner
sense are only empty fancies, because we do not find such figures in the external
world? Whence come the higher inspirations, if not from the unfathomable Creator
Spirit, who himself gives to all things life and breath and all? Truly, blissful
dreams too are from God! What else, indeed, should it mean that dreams and visions in
the Old Testament (as regards the New Testament, think of Joseph's dream,
Matth.\ 1, 20, of Peter's vision,
% --- p. 137 ---
\opage{137} %
Ap.\ G.\ 10, 10---11, of Joh.\ Aabenb.\ 1, 10 and other places) are set forth as
divine manifestations? How in the contrary case could one understand the prophetic
visions of an Ezekiel or a Daniel? Or perhaps such visions do not belong at all among
the revealed facts? --- The objection that all visions are at bottom only offspring
of man's own imagination falls away when we consider the matter more closely. How am
I to know that an external object has independence, except by being able to perceive
the resistance it offers, and to feel the necessity that compels me to represent it
to myself as it is once given? But should inner actuality not have a similar mark?
--- When someone says that he loves, or that he hates: how should he then be able to
know within himself that he speaks the truth, if he did not immediately feel the
involuntary promptings either of love or of hatred? The same holds of the knowledge
of truth: living conviction is always grounded in an immediate feeling. The mark of
inner actuality is therefore the necessity with which the powers of existence affect
the inner sense. Thus we can understand the visions of revelation: when the divine
gains shape, and presents itself to the soul's eye with such determinacy and
necessity that the seer cannot possibly confuse the vision itself with the other,
arbitrary productions of the imagination, then we must ascribe to the revealed figure
both essential independence and inner actuality by virtue of a supernatural
influence. --- The objection that in this way it becomes very difficult to separate
the true from the false, the divine from the merely human, must likewise fall away on
closer examination. There can, for example, be much dispute about original and
derivative, true and false poetry;
% --- p. 138 ---
\opage{138} but it is nonetheless evident that this dispute cannot be decided by
historical criticism. If that were the case, one would have to know beforehand who
was the true poet, and the whole investigation would turn on whether the poem in
question was by its alleged author or not. This dispute is meaningless. Who is the
true poet? Answer: he to whom the Muse brings the poetic manifestation. Who, then,
has received the poetic manifestation? Answer: he who produces the true poetry. By the
inner truth of the poetry, then, I am to know the true poet, the fortunate one whom
the Muse visits; and not the reverse.''

``It is the same, too, with the extensive dispute conducted about the separation
between the canonical and apocryphal Gospels. This dispute may perhaps be useful
exercise for the learned among themselves; but it yields only a meager return for
religious knowledge. Historical criticism can never come to an end; the longer one
disputes, the greater the confusion becomes, and why? Because the dispute is in
itself meaningless. Who is the true Evangelist? He who proclaims to us the true
Gospel! But how then am I to know the true Gospel, except by its satisfying my
deepest religious need? The Spirit of God bears witness with my spirit that this must
be the true, unadulterated Gospel. --- What has here been set forth in general finds
a special application to the account before us. What a difference between this
canonical Gospel and the apocryphal fabrications with which a pious but
superstitious imagination has sought to adorn the birth of Jesus! When Joseph, it
says in the Protevangelium of James, had brought Mary on an ass to Bethlehem to be
registered for taxation, she began near the city to behave now sorrowfully, now
merrily. When she is asked the reason, she gives the answer that she (after
1 Moseb.\ 25, 23) saw two
% --- p. 139 ---
\opage{139} %
peoples before her, the one weeping, the other laughing; that is to say, according
to the one explanation, the two parts of Israel, for which Jesus' appearance was to
be (Luc.\ 2, 34) the fall of the one and the rising again of the other; but,
according to the other interpretation, the people of the Jews, which afterward
rejected Jesus, and that of the Gentiles, which accepted him. When Mary is
immediately thereupon seized by the pains of birth, Joseph brings her to a cave
lying by the road, and now describes the solemnity of the hour of birth in the
following terms: I (Joseph) saw the pole of heaven stand still, the air stand
astonished, and the birds of heaven tremble. I looked out over the earth: the sheep
were scattered about; they neither went forward nor stood still. The shepherd lifted
his hand to drive them on with his staff; but his hand remained raised. I looked
toward the stream of the river, saw the kids stretch their heads over the water, but
without drinking, etc.''

``The objection that, according to this direction, the separation between the
canonical and the apocryphal presentation still at most becomes only a separation
between the genuine and the failed myth rests on a sheer misunderstanding. ---
Churchly-minded theologians have in a certain way reason to take offense at the
frequent use of the word myth when the talk is of narratives in the Old and New
Testaments. But in order to counteract the mythographic reinterpretation, it is of
little use to try to translate the inner manifestations into outer actuality, or to
appeal to dry facts and historical testimonies. Instead of simply setting the actual
against the mythical, one would surely do better to demonstrate the inner opposition
between myth and revelation. These two concepts are not related to each other simply
as the false to the true (myth too expresses in its way an essential idea, and to
that extent a truth); but they are related
% --- p. 140 ---
\opage{140} %
as the profane to the holy, as the natural to the supernatural. The idea that breaks
through in the form of myth points back to the world's own god-inspired natural
ground; the idea, on the other hand, that gains shape in revelation is permeated by
the Godhead's own supernatural being, and destined to proclaim God's holy will. On
this, surely, rests the real difference between the gods of mythology and the angels
of the Gospel. With this difference it is then also given that myth must be related
to revelation as the vanishing to the abiding. For since the natural and merely human
has a one-sided preponderance in the mythical representation, since the non-divine is
thereby in an unclear way confused with the divine, myth to that extent contains an
untruth; it lies, then, in the nature of the case that maturing thought must
gradually discover this untruth and lead beyond the mythical standpoint. It is
otherwise with revelation. In the holy idea the divine content has so worked through
the form that there is no more of the finite and human left here than precisely what
is required in order that the pure, supernatural truth may be made intuitable for man
and affect the whole man. Now as culture advances, reflection may indeed contribute
its part to making the infinite content of revelation inward; but never will human
thought at any stage of culture be able to dissolve the Gospel representations
without falling into an inhuman contradiction with the divine truth.''

``A striking proof of the correctness of this proposition we have here in the Gospel
of the birth of Christ. This Gospel has been handed down to us in so noble, pure and
perfect a form that, without altering its inner character, it accompanies the
religious mind through the various life-stages of culture, of experience
% --- p. 141 ---
\opage{141} and of self-knowledge, just like a star which, although it stands
motionless in the heavens, yet seems to move in order to accompany the wanderer. This
Gospel, so comprehensible, simple and lovable, is indeed as though it were written
only for children and childlike souls. The childlike imagination becomes as easily
familiar with the angels as with the shepherds. The separation between outer and
inner actuality does not hold at all for the child; here, then, there can be no
question of any myth: `it is all both certain and true in every way.' Then the child
grows up into a youth; but the Gospel grows with him. A painful struggle begins
between the hard, many-edged actuality and the longings and demands that rise in the
young man's breast. `The present will not satisfy; yet satisfaction must surely come,
hope cannot deceive: yonder in the distance the ideal beckons.' If the youthful
imagination now has strength enough to hold fast its ideal (and it is indeed the
glorious privilege of youth to have such strength); if the religious mind has both
strength and courage to lift its eye toward the prototypes of bliss, and to apply
these prototypes to the laboring aspiration of actual life (and it is indeed the
beautiful prerogative of youth to have such strength and such courage): must not the
future then shine for the inspired youth with the brightness of the Gospel, and the
angels themselves be summoned as witnesses that there really is \emph{a great joy,
which shall come to all the people.} Then the struggle continues: the youth grows up
into a man; but the Gospel grows with him, and transfigures its promises in the
struggle. `On this earth bliss is indeed not to be found; but the earth too is
nonetheless blessed by God. Yes, the blessing has taken away the curse; honest
striving bears its fruit, labor has its reward, the whole moves forward toward the
consummation yonder.' Certainly the sharp separation
% --- p. 142 ---
\opage{142} between the actual and the ideal now enters here; but then reflection has
also reached the maturity to make due distinction between inner and outer actuality.
The power of making inward grows amid the many outward-directed endeavors. Instead of
evaporating into a myth, revelation now first opens its inner and true kernel: the
conviction of manhood confirms the faith of childhood. --- With a good conscience,
then, the older man can read this Gospel again to a child, and answer every question
of guilelessness: `Is it really true?' You should believe it, my child! `But you
believe it too, don't you?' `Yes, this we should all believe, for it is God's word!'
--- Then the man ages and grows old, but the Gospel follows the old man, though it
cannot age. Life is a cycle, it is said; this saying too has a deeper meaning. The
old man becomes a child again, and comes nearer to the eternal. Outer actuality with
its noise and strife withdraws further away: it grows still around him. But as it
grows still, as in the night all around, the visions rise up from the ground of the
soul: the inner world opens before the eye. `Is it hope that beckons from a better
home? Is it a memory of childhood that greets in farewell? Why, it is the angels of
the Gospel; they praise God in the highest, they proclaim to man his good pleasure;
they shine peace over the earth.'\,''

The aesthetic interpretation can in many ways vindicate its justification over
against merely historical criticism of the understanding; it can in many turns
contribute to asserting the right of feeling and imagination over against
prosaic-finite reflection; it can, especially when it passes through speculation,
with great superiority defend the main thesis that the God of revelation must be just
as much the God of the ideal as of actuality, just as much the God of imagination as
of thought; in other words: the aesthetic
% --- p. 143 ---
\opage{143} interpretation asserts the demand that, as surely as there is a revealed
religion, there must also be a religious intuition of the imagination, for whose
content and validity the God of truth himself will vouch. --- How aestheticizing
speculation otherwise goes about such a proof does not concern this investigation at
all. The interpretation indicated above is only to serve to show whether the faith of
the Gospel can be satisfied with such an explanation or not. To this question we
must at once, without further ado, give a negative answer. The aesthetic-speculative
procedure suffers from a fundamental weakness that is all the more dangerous since,
with its soothing circumlocutions, it can otherwise easily beguile religious feeling.
Natural piety is only too ready to use the likeness it bears to faith to pass itself
off as faith itself and establish itself in its place. --- To transform every
angelic manifestation without further ado into a magical spirit-vision will not do
at all; for then, for example, the angel who rolled the stone from Christ's tomb
would also have to be such a spirit-vision. Matthew (K.\ 28, 2---6 V.) relates the
event itself thus: \textbf{And behold, there was a great earthquake; for the angel of
the Lord descended from heaven, came, and rolled back the stone from the entrance,
and sat down there. But his form was like lightning, and his raiment white as snow.
But the keepers trembled for fear of him, and became as dead men. But the angel
addressed the women, and said: fear ye not! for I know that ye seek Jesus, the
crucified. He is not here; for he is risen, as he said. Come hither, see the place
where the Lord lay.} --- That this narrative concerns something more than a mere
spirit-vision is self-evident. --- Now if the aestheticizing interpretation would
separate the vision both of the ``keepers'' and of the ``women'' from the natural forces
% --- p. 144 ---
\opage{144} at work in the earthquake and the lightning, then such a separation not
only conflicts with the plain words of the text, but the whole interpretation at the
same time surrenders itself to the very criticism it is intent on combating. The
proposed separation between myth and revelation is, when looked at a little more
closely, itself mythical: the mythical lies in the ambiguity. For if the supernatural
breakthrough of revelation in the believer's conscience could be explained by the
inspired eruption of the idea in the poet's imagination, then the latter would have
to be just as supernatural as the former. The difference between the evangelical and
the poetic manifestations is thereby reduced to a mere difference of degree. Let this
difference of degree be otherwise as great as it will: the concept of the
supernatural is nonetheless lumped together with the concept of the natural; and so
we have the mythical interpretation once again. If, on the other hand, the meaning is
to be this, that the poetic manifestations occur according to natural influences,
while the evangelical ones take place through a wonder wholly inexplicable to us:
what then is gained by the appeal to the separation between inner and outer
actuality? --- Against the separation itself, insofar as it keeps to pure
generality, we have nothing to object. The objection concerns only the purpose for
which it is put forward; for it is evidently designed to keep the manifestation in as
hovering a form as possible, so that the miracle shall not arouse offense. Against
that kind of attempt the faith of the Gospel must resist with all its might. Faith
is in the very highest sense a vital matter; but precisely because it is a vital
matter, precisely because he who would live in faith must set his hope on something
wholly other than he who has nothing to do with faith, precisely therefore it must be
possible to know with indisputable determinacy what it is we have to hold to here.
--- If the Highest, by which a person determines the whole direction of his life, is
something that both can
% --- p. 145 ---
\opage{145} and should be comprehended: then this Highest can in any case not be God (for no
human being can, of course, comprehend God himself); but if, on the contrary, this
Highest is something that neither can nor should be comprehended, then it surely betrays
a want of sound sense if anyone wastes his time trying to comprehend that by which he
is to determine his life in every moment. --- Just as it is in general a great
recommendation for a well-bred man that he knows when he should be silent and when he
should speak, so everyone, in relation to revelation, ought above all to make it his
business to know before whom and in what place he should be silent, and with whom and
in what place he is permitted to reason. --- Now since faith, with the utmost
strictness, bids the understanding keep silent before the wonder of the Almighty, that
ingenious separation falls away of itself. About the manner in which the angels became
able to manifest themselves to the shepherds, the most highly developed man knows no
more than the innocent child. Nevertheless there is a great difference between the
knowledge of the two. The child, namely, does not know at all that there is anything
here standing in the way of comprehension. The man, on the other hand, knows this very
well; but he also knows why he neither can nor should comprehend. Once this is given,
the manifestation of the angels is to be taken quite straightforwardly as an immediate
fact.

The faith of the Gospel is infinitely interested in the critical treatment of the
supernatural events. Only insofar as it has understood this interest of its own has it
understood itself. On this self-understanding it concentrates all its comprehending.
Faith --- be it noted --- has no wish whatever to comprehend the interaction of
revelation with the laws of nature and sensible actuality; but it does wish to
understand its own existential relation to the wonder of the Gospel. --- It is quite
true, what Lessing has said, that ``the transition by which one would build upon a
historical report
% --- p. 146 ---
\opage{146} an eternal truth is a leap.'' But this objection cannot touch the faith of
the Gospel at all, which of course presupposes precisely that everyone who calls himself
a Believer must have made this ``leap.'' Let us here examine a little more carefully
what the meaning of this talk may be. --- No finite experience and no science of reason
is able to wring from the thinking human being the admission that there must
necessarily be a supernatural revelation. If the human being nonetheless says: ``I know
with unshakable certainty that God has revealed himself in Christ for the salvation of
the world, I know it just as surely as I know that there is a God, and that I myself
live in this now'': then he surely makes a leap, since he at one stroke leaps away from
all long-winded reasonings and grasps the unconditional certainty in faith. --- No
historical research and no proof of reason is able to force upon the thinking human
being the conviction that God's revelation is to be found precisely in these gospel
narratives. If the human being nonetheless feels convinced that these narratives, with
all the facts belonging to them, are the right, unadulterated expression of the positive
content of revelation, then he cannot, of course, have attained this recognition by the
historical-critical way, but must have come to it by an immediate leap --- in faith. ---

Finally, it holds here that the transition by which the Believer breaks off all
calculations of probability and, in defiance of the contrary analogies of worldliness,
finds the entrance into the supernatural world of the evangelical facts, cannot be
separated from the immediate transition by which he comes at all to certainty that
there is a real revelation: both recognitions are attained by one and the same ``leap.''

Nevertheless there is a criticism of the understanding that reasons thus: ``positive
revelation and subjective
% --- p. 147 ---
\opage{147} faith presuppose each other, to be sure; but in order that they may rightly
unite for mutual support, each of them must, on its own, be subjected to a sober
deliberation, a careful testing. Only from such a sober and scientific testing will the
complete and many-sidedly grounded faith of the Gospel then be able to result.'' ---
This soberness, however, can easily bring about a good deal of confusion. How is it,
now, that one usually goes about it? --- With solemn, dragging half-heartedness, that is
to say, with a certain decorous acknowledgment of the historical testimonies of
antiquity in general, and an unctuous reverence for our evangelical traditions in
particular, one begins the historical-critical course. But what is the result? Does one
really, by applying the higher and the lower, the internal and the external criticism,
at last arrive at a decisive judgment on the authenticity of the sacred books? ``Not
quite!'' --- Does one perhaps, by adding the general introduction to the special one,
finally get the long investigation of the canon of the New Testament concluded? ``Not
quite!'' Does one then find oneself able to overcome all doubts about the genuineness
of the four Gospels and their agreement with one another? ``Not quite!'' --- But the
supernatural facts, those one must surely see to securing against every scientific
attack, as one secures the apple of one's eye, since it is, after all, chiefly on them
that the subjective faith in revelation is to lean when it suffers trial? ``No, not
quite! These facts, on the contrary, have the very weakest testimonies of all. Of
course, something must naturally have been seen, something extraordinary; but what
this extraordinary thing is, in what it actually consisted, on that the sober-minded
criticism, to put it plainly, dare not yet have any definite opinion.''

If it has now, by these sober concessions, become manifest
% --- p. 148 ---
\opage{148} to all that the credibility of the Gospels stands on such weak feet, then we
must look about us with all the greater expectation to the other side, to see whether
by the ``sober deliberation'' we might perhaps discover the firm, unshakable point of
support of faith in the religious mind itself. --- How then do matters stand with
subjective faith? Does the Believer perhaps now and then receive such inner
manifestations and immediate visions that he can with certainty infer from them the
reality of the evangelical facts? ``By no means! The age of miracles is long past; let us
above all see to keeping enthusiasm out.'' --- Does subjective faith in revelation, then,
have such clarity and reliability in itself that it could on its own defend itself
against every doubt about the reality of the supernatural, even in the case that
negative criticism should really succeed in dissolving the content of the Gospels into
fantastic myths and arbitrary legends? ``No, by no means! Our religious consciousness
will always need a firm, clarifying word; the believing view, which stands in need of
examples and patterns, must continually strengthen itself by looking to his pattern who
is the Way and the Truth and the Life.'' --- A very edifying answer! But if this is now
to say, in other words, that the Believer needs these examples merely for the sake of
vividness and clarity, then one could surely, without troubling oneself much about
historical reality, just as well transform the whole gospel book into a popular picture
book, and religious conviction would still be equally secure? ``By no means! That is
precisely why it is expressly said in the \emph{Introduction} that our weak faith needs
a historical point of support in positive revelation.'' --- If we now sum up what is
brought out on both sides by the sober deliberation of this half-criticism, then we must
indeed admit that the fidelity with which subjective
% --- p. 149 ---
\opage{149} faith and positive revelation support each other in this case is not exactly
all that reliable.

An example: A man seeks certainty in a matter which is naturally of the very highest
importance to him, because it is a matter of life. He is advised to turn to A. A
comforts him as best he can and ventures the supposition that he will surely obtain
complete certainty if only he will consult B. B likewise gives him the best of hopes, if
he will now just go back to A, and press him, and compel him to declare himself with a
definite Yes. But A solemnly assures him that he can by no means venture such a
declaration before B has said Yes. Then B begins to swear by all that is holy that for
conscience' sake he cannot say Yes before A comes out with it. Meanwhile the man is in
much anxiety and great perplexity; he has only these two to hold to, and the matter is
of the utmost importance, because it is a matter of life. Then, by great good fortune,
he meets on the way a clever man, a critical head, who gives him the following advice:
``It is quite true, to be sure, that neither of the witnesses dares to swear to the
matter with a complete Yes; but each of them nonetheless has the best supposition. A good
supposition is, when we rightly consider it, surely much the same as half a Yes. Two
halves make a whole, that is a settled matter. So be quite at ease: add the two half
Yeses together, and you have the complete confirmation.'' Must the man not now be highly
obliged to so sober-minded a man for his many-sided deliberation? --- Truly, if the
sober-mindedness of half-criticism really knew what it itself is doing, one might almost
be tempted to suppose that it was quietly setting out to make fools of both the Gospel
and faith.

What a distance between the merely probable and the true!
% --- p. 150 ---
\opage{150} Since, then, it is impossible to produce a whole truth by adding together two
probabilities, it lies in the nature of the case that the faith of the Gospel can attain
full assurance only by a decisive leap. --- In order, as far as possible, to forestall
any misinterpretation, we will consider the matter from yet another side. The meaning of
the ``leap'' mentioned is, namely, not that a Believer should be able all at once and,
so to speak, in one bound, to cross over from uncertainty to an assurance complete in
every respect and a perfect faith. On the contrary, the true and living faith lives
only in a continued appropriation. In this appropriation the Believer must continually
work against doubt, in that, under the constant influence of the Spirit that stirs and
works in the immediate testimony of the congregation (the Apostles' Creed), he partly
``searches the Scriptures,'' partly tests himself and his inner experiences as before the
face of God; insofar he finds himself, then, in an incessant searching and striving.
``Ye shall, says Christ (John 8, 32), know the truth, and the truth shall make you
free.'' But in order that this striving after knowledge of the truth may lead to the
goal, the Believer must know with determinacy what it is that causes doubt and weakens
conviction. --- According to the assertion of criticism, it is naturally the
incomplete and scarcely credible revelation that must bear all the blame; but according
to the believing view the whole blame falls on the human being, who does not know how to
appropriate the revealed truth. The evangelical facts and the testimony of the Spirit in
the awakened conscience cannot possibly be subject to doubt, unless the Believer is to be
permitted to take God to task and to contradict himself.

``The faith of the Gospel does not shun the light, it will not creep into hiding from
science, it will hold out well enough against half-criticism, it is not at all afraid of
the modern consciousness. If it is
% --- p. 151 ---
\opage{151} challenged according to the laws of honor and chivalry, it is at once ready to
put on its armor and enter the lists with its adversary. No rich man can fight with
greater zeal to keep his possessions from thieves and robbers than the Believer fights
to maintain the evangelical facts against the attacks of criticism. But once he has
become so familiar with the content of revelation that he lives and breathes in it,
then the critic too will find it hard enough to wrest from him even a single tittle of
the evangelical narratives. --- Negative criticism can point out particular features of
profane history that do not seem to agree with the Evangelist's statement; but what
does that prove? Is profane history perhaps infallible? Is it settled that the
contradiction complained of is more than apparent? Can an Evangelist not have had reason
to bring out a particular event that has been passed over by the historian, and
conversely? --- Negative criticism can further point to a particular apparent
discrepancy among the Gospels themselves. But all the narratives nonetheless bear the
same fundamental stamp. The diversity can be explained by the fact that each of the
witnesses has seen the matter with his own eye and made his selection according to his
own particular purpose. The deviation, where such a one really is found, is moreover so
inessential that in any case it only reminds us of the purely human side of the
presentation. --- Negative criticism can further raise difficulties with regard to the
oldest testimonies about our canonical Gospels, found in the Fathers, in the heretics
and the enemies of Christianity. But when these testimonies are weighed in the scales of
faith, the confirming ones are so far from being outweighed by the denying ones that the
latter even shrink to a pure nothing in comparison with the great, unanimous testimony of
the Church. --- Negative criticism can finally raise the objection that the Believer's
% --- p. 152 ---
\opage{152} %
judgment is not impartial. This objection we must expressly bring out; for it is almost
the only one that has any weight. The Believer is, we must admit, in these matters by no
means what is otherwise called impartial in the historical-critical sense. But before we
condemn the partiality, we must consider the matter a little more closely. It is not
with the facts of revelation as with the facts of world history. These latter can and
must be judged from the standpoint of impartiality, that is to say, from the standpoint
of disinterestedness; for, since their content is of a finite nature, it can in a higher
sense be a matter of indifference to the researcher whether they really took place thus
or otherwise, if only he can find out how. With the facts of revelation it is another
matter; their content has infinite significance: they lay claim to be believed unto
salvation. If one would now call him who believes them partial, then one must also call
him who does not believe them partial. Unbelief is thus on this point just as partial as
faith. But where impartiality is an impossibility, partiality ceases to be a vice; this
the resolute Free Thinker knows just as well as the sincere Believer. Only half-criticism,
which would be at once both within faith and outside faith, does not know it. Thus it
comes to stand not only on the higher standpoint of sober-mindedness and many-sidedness
and impartiality, but also on the higher standpoint of indifference and mediocrity and
thoughtlessness. Mediocre half-criticism is in certain respects invincible. It is in
vain here to recommend simple, childlike faith: mediocrity puts on a scientific air; it
demands insight into divine things; it wants now to think for itself. Then speculation
comes and lays its heavy thoughts on its round head; but hardly does mediocrity feel how
it aches
% --- p. 153 ---
\opage{153} in the brain before it puts on a simple face: ``no, no human being is after all
able to comprehend the mysteries of God''; now it would much rather keep to ``positive
revelation,'' and then to --- ``the critical introduction.'' --- Against two powers,
however, mediocrity will find it hard enough to hold its ground, and these two powers
are: the world-spirit with the titanic forces of nature, and the Gospel of the Savior
with the blessed angels.

How far the doctrine of the Middle Way can, on the whole, come to terms with the general
world-spirit remains a question by itself. Here we shall especially bring out this, that
half-criticism is not able to appropriate the Gospel. Above all, then, we must take good
care not to confuse reasoning mediocrity with immediate simplicity. The difference shows
itself at once when the relation between the ideal and the imperfect, the blessed and the
actual, is to be more closely determined. Mediocrity, which clings fast to the finite,
limits the ideal according to natural conditions, and gives up the blessed insofar as it
comes into conflict with the acknowledged actual. Simplicity, on the other hand, looks up
to the ideal and forgets finitude; appropriates the blessed and thanks God. Mediocrity is
conceited, simplicity is humble; mediocrity is refractory, simplicity is obedient;
mediocrity doubts and reasons, simplicity is silent and believes. Because of this
difference the one cannot even understand the other. That this assertion suffers from no
arbitrary exaggeration the Gospel will be able to prove. --- Joseph and Mary are simple
souls; but the mode of thought of mediocrity they do not know. Had they, in the days when
the imperial command went out, struck into the critical Middle Way, then the journey from
Nazareth to Bethlehem would certainly have become a most toilsome journey for them both.
Indeed, if their finite understanding
% --- p. 154 ---
\opage{154} had at that moment come to calculate the contradiction between the blessed and
the actual, the conflict within them would surely have become so desperate that they
would certainly never have reached --- ``the manger,'' but would have perished on the way.
Meanwhile, the criticism of mediocrity may perhaps find it hard enough to discover in what
this difficulty consists. One will rather suppose that the conflict whose possibility is
indicated here can far better be derived from a dialectical whim than from the
evangelical narrative.

As half-criticism reflects, and compares the various scenes in the history of the holy
family with one another, it meets such a superabundance of manifestations that it must
find it very doubtful that the blessed can really have expressed itself in so many and
such immediate communications. ``It is not enough that Joseph's dream-vision stands so
little motivated over against the annunciation to Mary; but there occurs hardly a single
turning point without a wonder at once taking place. The birth of Jesus is proclaimed to
the shepherds by a wonder; Simeon appears in the temple as by a wonder; by a wonder Joseph
receives the admonition that he is to flee with mother and child into Egypt, to escape
the persecution of Herod; by a wonder the family's return from Egypt is again determined
(Matth. 2, 19); by a new wonder Joseph is warned against Archelaus, and so finally finds
the way to Nazareth (Matth. 2, 23). Indeed, when one finds oneself everywhere surrounded
and guided by so many wonders, it cannot be so very dangerous to undertake even a longer
journey. All that was lacking was that an angel in a dream should also have explained to
Joseph why the imperial command absolutely had to be obeyed. Instead of making a fuss
about the short journey from Nazareth to Bethlehem, one must rather marvel that Joseph and
Mary on their various journeys are made partakers of so many manifestations that they
even seem to forget the one over
% --- p. 155 ---
\opage{155} %
the other. When (Luke 2, 18) all marveled at what the shepherds said to them, the
Evangelist adds (V. 19): \emph{But Mary kept all these words, and pondered them in her
heart.} Nonetheless it is said afterwards, on the occasion of Simeon's words (Luke 2,
33): \emph{And Joseph and his mother marveled at the things which were spoken of him}
(Jesus). If this marveling, as some think, is to be referred to how the holy mystery could
have been revealed to Simeon, then the words ought not, as they now do, to point to
Simeon's address as the object of marveling, but to the person of the one addressing them.
Finally, and this is the most remarkable of all, we learn (Luke 2, 50) that the parents
did not even understand the words the twelve-year-old Jesus spoke to them, when they had
found him in the temple sitting in the midst of the teachers: \emph{Why did ye seek me?
knew ye not that I must be about my Father's business?}''

This view is now not at all so strange, if one will only place oneself on the standpoint
of mediocrity and regard the scenes cited at a proper distance. The observer knows
beforehand that Mary, according to the narrative, has now really become the mother of the
Son of God; from this he concludes that Joseph must just as easily have been able to know
it. The observer knows in advance that, even if the child must be laid in the manger
because there was no room in the inn, the shepherds' report of the heavenly hosts will
soon console both Mary and Joseph over their abasement; from this he then concludes that
Joseph and Mary, even without the shepherds' message, could easily have consoled
themselves. In short, the observer is in his thoughts always a step ahead of the event:
from this he then concludes that Mary and Joseph must likewise be so. The observer is so
farsighted because he knows the whole story by heart; from this he supposes that Joseph
% --- p. 156 ---
\opage{156} and Mary must have been just as farsighted; and forgets that they had to live
through ``the incredible'' precisely step by step. --- That this is a backward view is
surely self-evident; but here is yet another contradiction of which the criticism of
mediocrity makes itself guilty. It assumes, namely, without further ado, that a
manifestation is a manifestation, and that the matter can thereby be settled. ``Once one
has received a real manifestation, and with it a clear instruction about God's express
will, then one also knows what one has to go by in the main: all the rest God may then
well leave to the human being's own conduct.'' --- Yes, once one has received a real
manifestation! But how does it come, then, that mediocrity itself has such difficulty in
convincing itself that the manifestations discussed here are real?

The criticism of mediocrity now thinks something like this: ``Had I only been in Joseph's
or Mary's place, I would surely have convinced myself whether the first alleged
manifestation was true or false; and when I had then, by a sober deliberation, come to
agreement with myself, I would also have known how to take my measures accordingly.''
--- Such a line of reasoning proves precisely how little mediocrity is able to put
itself in either Joseph's or Mary's place. The difficulty lies precisely in this,
\emph{that the supernatural is heterogeneous with the natural.} This difficulty cannot
possibly be removed by a human being's himself experiencing the manifestation. On the
contrary, such a human being must feel precisely how heavily the opposition between the
blessed and the actual weighs and presses with all its harshness. Let us just for a
moment try to put the reflection of mediocrity itself into Joseph and Mary.

Some months before Emperor Augustus's command was to be
% --- p. 157 ---
\opage{157} carried out, the betrothed have thus, each on their own, learned the divine
counsel in a particular manifestation. The condition in which Mary now actually finds
herself seems to confirm the truth of the manifestation; then doubt rises in her soul.
--- ``It was a blessed moment: the angel stood before me and greeted me as the one
highly favored. But am I then really the one highly favored? Can Joseph convince himself
of it? If the world suspected my secret, it would rise up against me, to murder me for my
blasphemy, or to mock me for my madness. It was a blessed vision; but the vision has
vanished: this condition is unblessed! Is this what it is, in actuality, to be the one
highly favored? I am conscious of no guilt; but is the whole thing not like a sorcery?
Can not the spirit of darkness also put on the form of an angel of light?'' --- We need
not go further to see what the reflection of mediocrity would work under such
conditions: a single doubt, and Mary is lost --- in madness. The same holds of Joseph.
In his inner being too the blessed must fight a terrible fight with the actual, if
reasoning criticism finds entrance. --- ``A blessed dream, it solved the dark riddle for
me! The angel breathed his peace deep into my soul: so I could understand Mary again.
--- If I now dared to trust a dream, then my lot would be the fairest on earth. As one
who is of the house and lineage of David, I should then stand by the cradle of the
Messiah. Under my fatherly care the Son of God, Immanuel, should then grow up and
increase in stature and in wisdom. But that is surely impossible. --- Only a dream! What
am I in the actual world? A humble craftsman in Nazareth; so low has David's line then
sunk. --- Absurd! What am I in my relation to Mary? As a deceived man. What would people
judge, if I stood forth and spoke to them of my
% --- p. 158 ---
\opage{158} riddle? They would laugh, and be right. --- So I am a fool, then! The actual
world looks at me and mocks me with its clever eyes: it scoffs at me, it laughs at my
dream.'' ---

It does not go so easily with the ``sober deliberation'' as half-criticism seems to
think; the longer one deliberates, the heavier understanding becomes. --- What is the
reception in Bethlehem like? Mary needs a lodging; her hour has come. But there is no room
in the inn! she finds only a little place by the manger. --- ``Is this to be the
fulfillment of the prophet's prediction (Micah 5, 1)? Can this be called the fullness of
time? It seems in actuality to be a most inconvenient time. To be cast aside like this,
like one who is given over to the power of chance, see, that is to be the sign for the
one highly favored!'' --- No, with the reflection of mediocrity it is here altogether
impossible to make any headway. The contradiction between the blessed and the actual is
all too glaring. --- But in pure simplicity the contradiction is overcome; for simple
faith keeps doubt out. Joseph the carpenter is certainly no critical head. Half-criticism
may well smile at his insignificance. For --- is it not so? --- he is after all a most
simple man, this Joseph? simple, that he believes in a dream; simple, that he takes the
forsaken virgin to himself again; simple, that he so unconditionally obeys the emperor's
command; simple, that at so critical a moment he takes Mary to Bethlehem without having
ordered a lodging beforehand. --- But if it is now true that the simplicity of faith alone
is able to unite the blessed with the actual, then we must call Joseph and Mary blessed,
that they were simple souls: happy for them that they did not know the mode of thought of
mediocrity.

If this is how matters stand with the appropriation of the supernatural facts, then a
Believer will also understand
% --- p. 159 ---
\opage{159} %
what the many different manifestations in the holy family circle may well mean. --- Under
the general presuppositions of worldliness and of merely natural piety, human life is in a
certain sense laid out far more simply than it can be by the measure of faith. The natural
man knows what he has to hold to, insofar, namely, as he sees that in this sensible world
he must use his senses and his understanding, and come to terms with circumstances. Here
there is only one fundamental law, the law of the world: here there is only one world.
Otherwise with the life in faith. By virtue of the personal relation to God the human
being comes into relation to a supernatural order of things. The Believer must thus live
at once in two worlds, the ideal and the actual, the blessed and the perishable. For such
a dweller on the border, existence must necessarily be far more complicated than for him
who has set up his dwelling in the midst of the well-built, well-secured regions of
naturalness. The double world of faith is everywhere crisscrossed by opposite ways,
opposite movements, opposite laws. In that the Believer would subordinate himself to the
eternal, he must, as it were, challenge finitude; for the law of eternity is to be
fulfilled in the finite. In that he would obey God, he must at the same time watch over
his understanding; for folly and madness do not please God. On such a thread of life there
then forms a multiplicity of knots that cannot be untied without God's particular aid.
The more decisively the human being acts, the more unconditionally must he commit the
matter to God, and the more unconditionally he gives his cause into God's hand, the more
strongly is he in turn called upon to exert his abilities and powers. The Believer is one
highly entrusted, who must render account. --- If this can hold of the life of faith in
general, how much more must it not be able to hold of the life of the holy family in
particular. See, Joseph and Mary
% --- p. 160 ---
\opage{160} %
are highly entrusted: God's own Son is entrusted to them. The criticism of mediocrity need
not fear that the many manifestations should have a numbing effect on their minds: at
every decisive step they need God's wondrous aid. The divine child is, after all, itself a
wonder, an extraordinary gift from heaven: what wonder, then, that heaven keeps watch over
it, that the angels watch over it, that the Spirit of God draws the pious and gathers them
around this child, that they may rejoice in worship and bear testimony. --- Insofar as the
marveling of Joseph and Mary presupposes such a general understanding, this marveling is
to be taken first of all as an inwardly blessed admiration, which testifies that these
simple souls receive with every new manifestation a new, original impression of the
divine. But on the other hand every breakthrough nonetheless comes at the same time as
something unexpected and unforeseen. The calculation of probability is broken off: no one
knows what God will do in the next moment; no one comprehends what is now to happen with
this child in this actual world, how it is to develop, what stages it is to pass through.
For precisely that reason each single scene, with its various attendant circumstances,
must acquire an infinite significance for Joseph and Mary. In this respect their marveling
can be regarded as a religious startle, a holy surprise that awakens the deep,
half-understood presentiment. --- It surely cannot possibly go with the inward
understanding of the Son of God's entrance into the world as with the finite understanding
of an ordinary event. The criticism of mediocrity does Mary wrong when it pictures the
matter thus: ``Had it been I who had received beforehand from the angel of the
annunciation the explanation that the child was really God's only-begotten Son, I would
certainly not have fallen into any particular marveling over the shepherds' report, but
would have thought in my quiet
% --- p. 161 ---
\opage{161} mind: yes, of course, the angels must naturally praise and extol my son, for
this son of mine is, after all, the Messiah himself, the Son of the Highest. But let it be
what it may with the first surprise and the first marveling: over Simeon's words in the
temple there was now absolutely no reason to marvel. All that the pious old man speaks in
prophetic inspiration as he holds the child in his arms, Mary and Joseph must surely have
known as good as by heart beforehand. Instead of letting herself be blessed by Simeon, Mary
ought rather herself to have blessed the old man, because he spoke so clearly and so
correctly. The further the story advanced, the surer one must, of course, become of one's
case. When therefore Jesus in the temple answered his anxious mother: `I must be about my
Father's business,' Mary ought evidently, according to what had happened so far, to have
exchanged a meaningful glance with Joseph, instead of their both now standing there like a
pair of --- natural men (1 Cor. 2, 14), who do not grasp the things that belong to the
kingdom of God.''

If mediocrity, however, did not take upon itself with quite so much assurance to judge of
things it does not understand, such a misunderstanding could never arise. Joseph and Mary
could, in all their simplicity, well understand that the divine child must be the Messiah,
God's own Son; but they could also understand that this child was now entrusted to their
care, that they were therefore not only to rejoice over this wonderful gift of God, but
also to learn to grasp God's will and counsel. Therefore there can be no question here of
superficial assurance and habitual security. Mary, who looks upon her son, looks into an
unfathomable depth of wisdom and grace; therefore she is so watchful, and yet so inwardly
still: ``she keeps all these words, and ponders them in her heart.''

% --- p. 162 ---
\opage{162} But, continues the doctrine of the Middle Way, even if Joseph and Mary, together
with the others in the holy family circle, were in this way able to join the blessed
with the actual in their pious simplicity: still, that is not the side from which the
matter must be taken if the Gospel is to penetrate into the present age. The blessed will
no longer approach consciousness in such naive appearances. Just as the individual human
being cannot possibly become a child again, however much he might wish to, once he has
reached the riper age, so it is also impossible for humanity, at its present stage of
culture, to screw itself back to that childlike standpoint of antiquity. Whence does it
come that the educated man, even with the best will, cannot believe these childhood
miracles, if not from this, that the mighty spirit of the age votes against them? The
present age is, unfortunately, unbelieving rather than believing. To offer an unbelieving
age so pointed a presentation of miraculous scenes is only to increase unbelief and make
bad worse. If, then, one would take the matter from the right practical side, let one
allow the evangelical images to make their aesthetic impression, and for the rest see to
bringing out the religious-moral substance. The human race cannot save itself; the Savior
was born into the world, that it might become manifest to all that God does not forsake
human beings, even if in their straying they sometimes seem to forsake him: ``\emph{Glory
to God on high!}'' True, there is strife in the world, and the more freely the forces act,
the greater the exertion becomes, the more vehement the strife. But if only we strive to
the best of our ability for truth and right, then we can calmly look up to him who was
born to strive the great strife, and hear in our own conscience the heavenly voice:
``\emph{Peace on earth!}'' True, evil is in the world: ungodliness cannot please God. But
neither the secret guilt nor the open crime can be rooted out
% --- p. 163 ---
\opage{163} by miracles and faith in miracles. Let us therefore with manly courage hold out
in the honorable struggle of freedom and morality; we can then console ourselves that God
has revealed himself in Christ, not as the God of wrath, but as the God of love:
``\emph{Toward men good will?}''

Whether the doctrine of the Middle Way, moreover, finds itself able by such edifying
grounds of consolation and moral reinterpretations to satisfy the restless world, is here
left undecided; we remark only this, that such grounds of consolation can also be had
without particular regard to the Gospel. We readily grant that mediocrity is a power in
the world; nor are we unappreciative of the economical good sense that so nicely
understands how to give the temporal a suitable flavor of the eternal, the profane a
seemly admixture of the holy, and thereby finds itself able to hand out to the needy, if
not exactly the bread of life, then at least some serviceable rules for living. Yes,
mediocrity is a power in the world: when the horizon darkens and the waters of the age
are stirred, then it is usually mediocrity itself that in its own person grasps the helm
and takes the word of command. That will do in unsettled weather, as long as the storm is
not too strong. But if the hurricane breaks loose, if humanity becomes incensed because it
feels itself disappointed and cheated of the Highest, then the passions can easily surge
up, and knock mediocrity over, and turn into despair.

--- It is, after all, a strange saying, when it is said that the educated world cannot
possibly return to that stage of childlikeness and simplicity on which souls may rightly
rejoice in the Gospel. As if for us human beings there could be any other relation to God
than the childlike relation of simplicity. --- Reasoning half-criticism seems here to
forget what it nonetheless so often says itself, and what thoroughgoing speculation
sufficiently
% --- p. 164 ---
\opage{164} confirms, that the depths of wisdom are unfathomable. But even if the wise had
now begun to grow proud, and the scribes to harden themselves: should these signs of the
times then really be a proof that humanity can no longer return to that stage of
childlikeness and simplicity where it can find its consolation and joy in the
manifestation of the blessed? Has every other person nowadays perhaps become a thinker, a
critic or a scribe? Were there perhaps in the time of Emperor Augustus and King
Herod no conceited wise men and no hardened scribes?

But if it is so, that the noblest and most gifted human beings must always rejoice in
understanding the blessed with faith in all simplicity: should the Gospel then not still,
as before, be able to be preached to the poor? The angel who testifies to the shepherds of
the Savior's birth does not, after all, say: I proclaim to you a joy which can yet befall
only the wise and the scribes; but he says expressly: \emph{behold, I proclaim to you a
great joy, which shall befall all the people.} --- It is, after all, a curious saying, when
it is said that the faith of the Gospel demands too much of the people of the present age,
that the appropriation of the revealed facts should now exceed human powers. How much,
then, was required and demanded of the pious of those times? Revelation came as a gift
from heaven! Nothing is reported of the shepherds' having worked for any awakening, or
taken part in any grand struggle of liberation: they tended to their humble work, they lay
in the field and kept night watch over their flock. But when the angel of the Lord stood
before them, and the glory of the Lord shone round about them, when the glad tidings were
proclaimed to them; then they did not set about criticizing and reasoning, but they
believed the manifestation, and said to one another: \emph{let us now}
% --- p. 165 ---
\opage{165} \emph{go to Bethlehem, and see this thing which the Lord hath made known to us.}
And they were not disappointed either; it was a truthful angel that had brought them the
Gospel; \emph{and they came with haste, and found both Mary and Joseph, and the child lying
in the manger.} --- Is it then demanding too much of human beings, when it is only
required that they should believe the word of revelation, and seek salvation where it is
to be found? Can the clever criticism of the understanding then not understand that true
simplicity must be the completion of human development just as well as its beginning? ---
If the wise and the scribes will not see this, but misuse spiritual culture, and put their
own arbitrary notions in the place of the eternal truths: no matter! existence itself will
then, by the power of circumstances, surely teach the people to see where the danger is,
and with whom salvation is to be found. True enough, one does not exactly become anything
great in the world by accepting the Gospel of Christ. The shepherds, too, surely got no
promotion in their outward station; but --- they did see the child in the manger, and they
made known what had been spoken to them about this child. And it had been spoken to them
that this Gospel was to bring peace on earth. That is a truthful saying: if once all the
people could grasp and appropriate the great joy that shall befall them through the Gospel,
then brother would be reconciled with brother, then human beings would forget their
quarrels and strife, then there would be ``peace on earth.'' And it was spoken to the
shepherds that God, for the Savior's sake, found good pleasure in human beings. That is a
truthful saying: if only all the people could grasp and appropriate the great joy that
shall befall them through the Gospel: then hearts would bow in obedience under the holy
law; then the times of freedom would dawn, for no one who obeys God's law can be the slave
of arbitrary notions: then the Spirit of truth would
% --- p. 166 ---
\opage{166} testify that there was ``toward men good will.'' --- That it is now so hard
before this Gospel can find entrance with all the people, that it must pass through
tribulation and distress, that childlikeness is first regained in chastisement, and
simplicity bought with pain, of this the angel of the proclamation does not think at all.
The blessed do not heed the refractoriness of carnal hearts; to them it is as if it went
without saying that a fallen race must surely be quite willing to let itself be saved. The
heavenly hosts that shine for the shepherds in the peace of the night do not look at all
at the world's strife, straying and hardening; but they look at the manifestation of
salvation, reconciliation and blessedness in the world; they look at the eternal essence
of humanity, at the image of God in which we are all created, at the only-begotten of God
who rests in the manger, and by whom this essential image is again to gain form. Therefore
they rejoice on behalf of human beings, therefore their song of praise sounds in blessed
choir: {\large Glory to God in the highest, and on earth peace, and toward men good will.}

\begingroup\small
\begin{verse}
``The hour when the Savior was born on earth,\\
There was born a comfort for the bitterest pain,\\
A comfort for him who hopes and believes,\\
An echo of the choir of the blessed\\
Descends then into the human heart.''
\end{verse}
\endgroup

\subsection*{3.\\[4pt] The Gospel Forewarnings.}
\addcontentsline{toc}{subsection}{3. The Gospel Forewarnings}

In order to be able to grasp any one, be it whichever it may, of the upheavals that take
place partly in the world of nature, partly in the world of spirit, one must above all see
to directing one's attention to the introductory movements, to the forewarnings and the
first outbreaks, i.e., one must above all see to living through the beginning.

% --- p. 167 ---
\opage{167} If, on the other hand, an upheaval occurs so unforeseen that the former order
of things is already overthrown and all bound elements set loose before anyone suspects
it, then one is easily overwhelmed and is unable to grasp anything, because one oneself
loses all composure. The whole scene then stands before one as an inexplicable and
terrible sight, until one at last gets so far as to be able to go back again and
bethink oneself of the beginning.

That the question of the beginning must always exert a peculiar attraction on the
contemplative mind is grounded directly in the nature of the case. The immediate beginning
is like a springing point that forms the boundary between that which passes away and that
which comes into being; a strange, ambiguous something that stirs in the transition from
the possible to the actual.

Before we make any closer application of this to the evangelical history, however, we will
first assure ourselves of the many-sided significance of the beginning by choosing a
couple of examples, from nature and from world history.

What scene of nature is better suited to awaken the spiritual in a human being than the
storm with lightning and thunder? And yet, if such a storm suddenly (as, e.g., in a
mountain region) falls upon the wanderer with all its violence, there is hardly time for
contemplative enjoyment. The blinding lightning, the deafening din, and then especially
the drenching rain surely make one altogether unfit for lofty reflections and
spirit-filled observations. Should the traveler perhaps take courage and defy the four
elements, in order properly to enjoy the great spectacle and lose himself in admiration
of the majesty of the spirit of nature? That will certainly not do. Majesty is in revolt:
he tolerates no contemplation, and will hear nothing of any admiration. Contemplate Him,
the Terrible One? Another time! Now he is angry, his brow dark: he threatens.
Admire Him?
% --- p. 168 ---
\opage{168} He is so hot, so wild, his eyes burn: how it flames in those black brows! ---
And these heavy rumblings, this wild laughter! ---
Flee, wanderer, if you can, to the mountain's cave! away into shelter! for if he notices
you, then you are a dead man.

If, on the other hand, we choose the moment when the storm is gathering and approaching its
beginning, then the event of nature takes on quite another appearance, then there is
something here to contemplate. A strange calm, an almost oppressive stillness is spread
everywhere. The air is so heavy, the sunbeam stings, all forces slacken: it is as if nature
felt faint and were about to swoon. This condition is in actuality most unpleasant: one
suffers, and waits for a change. Then a little cloud shows itself at the rim of the
horizon; it grows, it multiplies little by little, it towers up, even straight against the
wind: that is a forewarning. Soon the heavens are overcast with masses of cloud in
manifold layers of the most varied shapes and colors. An inexperienced person, who did not
know what all this was to mean, might easily hit upon the thought that the whole thing was
perhaps a kind of jugglery, that nature in its indisposition had fallen into slumber and
was now dreaming of --- dragons and giants, patriarchs and angels, floods and
conflagrations, and displaying its dream-images in the air. But we know, of course, that
it is nonetheless no juggler's show; we know that these strange masses of cloud are
workshops for sound and strong forces. There is work going on in there, although
everything outward points to arbitrary whims; there is work going on in there: glimmering
flares, swiftly vanishing streaks are glimpsed through the dark layers. There is work
going on in there: the spirits of the storm, who dwell in the airy tents, bustle back and
forth with fire and light, as if they were to kindle the lightning. And yet all is still
in ambiguous calm, in anxious stillness: the whole is as yet only a possibility. The
forewarnings have shown
% --- p. 169 ---
\opage{169} themselves: the preparations are at an end. Then it happens in an instant; as at a
sign from One up there who leads the way, it happens in an instant: the lightning
flashes, and the air quivers, and the wind stirs; and the rain cloud stoops, and
sinks low, and its fragile floor bursts, and great, heavy drops fall on the leaves of
the trees, on the dusty earth. It is a refreshing moment; we experience the beginning.

And now in history, what upheaval could be more magnificent here, and more
attractive to the observer, than the upheaval that we among us (significantly
enough!) designate by the foreign word: a revolution? But here too it holds that he
who would grasp the spiritual impression must not plunge heedlessly into the
maelstrom of events. The forces of nature are terrible when they rise in revolt;
but --- be it said both to the honor and to the shame of mankind --- when the
passions flare up and set the edifice of the state on fire, then all the strife of
the elements is tame as an innocent jest compared with the frenzied raging of an
inflamed people. --- As regards history, it goes without saying that the observer,
in order to find the right point of departure, must go back to the beginning. But
here too the point is to \emph{experience} the beginning; for the strife that
possibility here, at the first steps of transition, wages with the given actuality
is infinitely significant. In order to illustrate this more closely, let us give
the two contending powers, possibility and actuality, a more definite form, so that
we may the better grasp the impression.

The given actuality, then, is in the first place an antiquated body politic that
rests on its historical ground, and for that reason lays claim to be regarded as the
natural creation of the idea of right, as the express image of justice. But this is
unfortunately a sad self-contradiction; actual right has at every point been turned
into essential wrong. The prince upon his throne calls himself the head of the state
and
% --- p. 170 ---
\opage{170} the father of the people. He is actually in the right; he upholds what is
inherited, he does the same as his fathers did before him. And yet he is
nevertheless essentially in the wrong. As by a sevenfold ring wall of privileges
and arbitrary favors, the head of the state is separated from its body, and the
father of the people, who is kept as if shut up in his castle, surrounded by the
favored, who call themselves the pillars of the throne and guard the entrances of
the palace, cannot, however gladly he might wish to, hear the voice of the people
calling as from afar. Every time the great now make use of their prerogative, they
do, under these conditions, essential wrong; and when the oppressed try to take the
law into their own hands, they commit an actual wrong. So oppressive, so
intolerable can the situation be. And yet the whole still is as though it were in
order: the calm of uniformity and the stillness of habit are spread everywhere. But
it is an uncanny calm, an alarming stillness, for it smolders in the depths, and
the contradiction works in secret. Then come the forewarnings, one after another,
in all ambiguity. --- The law is defied. Is it desperate self-defense or the
presumption of brutality? ``It means nothing. Only an effervescence of restless
heads, loose, powerless attacks on the established order.'' --- But --- new
phrases are forming, new slogans: strange flashes of thought gleam in the masses.
See, the people stirs like one who is tormented by uneasy dreams and struggles to
wake. --- ``Should there really be danger afoot? Then let the authorities show
themselves with threats and promises, and remind the subjects of their oath of
allegiance and their duties.'' It helps! Where the authorities step forth, the
people draws back: the individuals do after all feel a secret awe, an involuntary,
obscure reverence for the higher power. --- ``Yes, did I not say so: the whole
thing is at bottom a farce; the well-disposed people has not yet forgotten
obedience.
% --- p. 171 ---
\opage{171} One need only step forward boldly and overawe it with the law; then it helps at
once: now it will soon come to rest.'' --- But no, it does not help! It burns in
the depths, the contradiction works in secret. Restless crowds move to and fro. The
enormous mass of the people is like a surging sea; one hears the distant roar: it
portends storm and tempest. ``If it should come to earnest, then the Most High
himself must bestir himself: something extraordinary must happen!'' --- So the
prince then descends from his throne, steps out of the palace door, and speaks
comforting, warning and exhorting words to his dear people, as a father to his
children. --- ``See, it helps. The swarming multitude breaks out in jubilation. What
a touching scene! The whole thing is surely a good-natured enthusiasm. The
trusting people has longed to see the father of the land. Now it has seen his face
and heard his promises: now there is no danger, now it will surely pass.'' --- But
it burns in the depths, the contradiction works in secret. The pillars of society
shake. Every shock brings new complications, new public calamities. It is like the
eruption of a subterranean fire that no one can quench. ``If the powers of the
underworld mean to stir, then it is time to consult the oracle: after all, one must
hear the public voice!'' --- So then the spokesmen of the people step forward. In
long rows they take their places opposite the mighty: the ambiguity mounts to the
highest pitch. Do not wonder that these men of the people are so thoughtful, so
earnest, so sparing of words: they stand, after all, on the border of two worlds,
the world of possibility and of actuality, of the future and of the past. Do not
wonder that these strange figures are as stiff and cold as if they were cast in
bronze: the sword of justice is over their heads, and the sword is unquiet, for ---
possibility contends against actuality. Do not wonder that these new spokesmen are
so
% --- p. 172 ---
\opage{172} moved, that their tongues glow, that they whisper among themselves like
conspirators who have consecrated themselves to an unforgettable deed. You will
understand them if you only remember whence and why they have come; if you see what
it is that is going on here round about. --- ``A countless swarm covers the
surroundings in wild groups: fantastic beings, uncanny shapes, distorted faces such
as no one has ever seen before by the light of day! Why, it looks like a carnival,
a gruesome popular entertainment, a masquerade.'' --- No, it is no masquerade and
no carnival; but it is possibility at work, now approaching its actual beginning.
--- One lingers, and waits, and breathes so heavily. --- ``What is one waiting
for?'' --- A decisive word (for the word is and remains, after all, always the
watchword of mankind.) Hear how they contend in there about right and wrong! ---
One lingers, and waits, and breathes so heavily.
Then it happens in an instant, as at a sign from One down there who leads the way,
it happens in an instant: a word! a cry! a shrieking cry! it is \emph{freedom and
salvation!} And the spell vanishes, and the stream of the people swells, and breaks
all dams, and sweeps away the throne, and buries a world of the past with all its
misery and all its glory in its gulf: here the observer has found the right moment,
now we experience the beginning.

But, if we consider rightly, did not the coming of Jesus into the world bring about
a change, a spiritual upheaval, such as was never seen before or since? If there is
any turning point in history of which it can truly be said that a hitherto unknown
possibility broke through and actualized itself in a new beginning, then it is
surely the turning point the Apostle calls \emph{the fullness of time.}

Let one judge the work of redemption otherwise as one will: it must surely have
been a great moment when these powers of Christ stirred for the first time, and this
work took its
% --- p. 173 ---
\opage{173} actual beginning. It must surely have been a great moment when the voice began
to cry in the wilderness: \emph{Prepare ye the way of the Lord, make his paths
level. Every valley shall be filled, and every mountain and hill shall be brought
low, and the crooked shall be made straight, and the rough ways shall be made
smooth; and all flesh shall see the salvation of God.} It must truly have been a
great moment when He who knew himself as the only begotten Son of God came forth as
\emph{the Son of man} in the assembly of men, and repeated the Forerunner's stern
word in the mild exhortation of love: \emph{Repent ye, for the kingdom of heaven is
at hand.}

Are we now in any way able to recall to ourselves the original impression of these
events, to experience this beginning? --- There is an admonition we must hear
first, and it is the admonition of the Risen One to the doubting disciple:
\textbf{Because thou sawest me, Thomas, thou believedst; but blessed are they that
have not seen, and yet believe.} ``Mysterious as His departure from the world is
His coming into the world as well,'' says criticism, and it is right; but the
mystery is --- revealed to the Believers.

If anyone is bent on \emph{comprehending} the mystery of this beginning, he may
freely accuse the Evangelists of spiritual narrowness, mutual disagreement and
contradiction; they do not answer him, they do not defend themselves, they sit as
if scattered, each by himself, as mute as images of stone. If, on the other hand,
we say: We do not demand to \emph{comprehend} (that is here both too much and too
little for us); but our desire is to \emph{experience} the beginning (as indeed the
Believers in the inwardness of appropriation gladly wish to experience the whole
life of the Savior); then there is a stirring in the four witnesses, then it is as
though the Evangelists drew near to us and thereby to one another,
% --- p. 174 ---
\opage{174} each with his Gospel, and pointed to the first things, and said: \emph{Come
and see!}

If there is any point in the sacred history of which one can truly say that it has
been grasped from all sides, it is surely the beginning.

If we turn to John, John the Evangelist --- the Church sees him transfigured in an
image: the eagle at his foot, heaven in his eye ---, then he interprets for us the
eternal, the original beginning, the beginning from above. \emph{In the beginning
was the Word, and the Word was with God, and the Word was God.} But when he then
adds, as though he soared away over the intermediate links of finitude: \emph{And
the Word was made flesh and dwelt among us (and we saw his glory, the glory as of
the only begotten of the Father), full of grace and truth;} then the eternal melts
into the temporal, and the light dazzles our eye. If we now feel that we are on
earth, and lack ``the wings of the eagle'' and the clarity of vision, and need the
finite representation, then Matthew and Luke offer their guidance. They then relate
to us how it came about that the Only-begotten, who was in the bosom of the Father,
descended into the bosom of mankind and was held to be a son of Joseph and Mary: we
meet in them the Holy Family in the most remarkable scenes of the Fore-Gospel.
While John, then, dwells especially on the \emph{supra-historical} beginning,
Matthew and Luke on the other hand communicate the \emph{pre-historical} events,
and now we understand Mark as well, who begins directly with the public life, and
thereby at once produces the strong \emph{historical} impression.

Nonetheless theology and criticism have prettily agreed to misunderstand this
beginning, and that in a twofold respect. The many unsuccessful attempts to find
out the connection
% --- p. 175 ---
\opage{175} between the Fore-Gospel and the actual history show this plainly enough. First it
is assumed that the pre-historical accounts must in every way be placed on the same
level with the historical ones. But when one then becomes aware (which for that
matter is self-evident) that it does make a difference whether it is a mysterious
annunciation of an extraordinary man's birth, or this man's own public appearance,
that is being described --- that one cannot, for example, speak of the angel Gabriel
in the same language as of a John the Baptist --- then one thinks one has made a
great critical discovery. Instead of making the ideal facts of revelation
intelligible to oneself, and considering what it actually is that the faith of the
Gospel requires, one now rushes straight over into the opposite extreme, draws a
line through the divine, and so reduces the whole Fore-Gospel to a mythological
prelude. --- This goes fairly easily now, not exactly because it is ``science
itself'' that does it, and everything that ``science'' establishes must of course
be well established, but chiefly because the whole matter about which the
proceedings turn is, after all, nothing but a simple matter of faith. If it were
only some other and more considerable matter, for example an ordinary matter of
honor, a matter of nationality, or a most serious matter of money, whose settlement
depended on the right interpretation of these events, we should soon see sound
common sense rise up and seize ``science'' by the collar, and shake the long
``critical introduction'' until theology at last learned to see what an
extraordinary beginning has to signify. But --- ``a matter of faith? --- Nothing
more!''

For the critical misunderstandings of theology, however, the Evangelists bear no
blame. The Evangelists are men of the Spirit, and men of the Spirit must surely
know that the extraordinary beginning is experienced and presented only by
beginning to
% --- p. 176 ---
\opage{176} see the forewarnings of possibility behind the covering of actuality, as behind
a half-rent curtain.

Let us therefore for a few moments bid farewell to the modern consciousness, and
wander in the footsteps of the Apostles back to Palestine, to the Holy Land, to the
point of time when the public life of Jesus took its actual beginning. See, the
Evangelist Luke will himself lead us on the way, as he relates the following:

\begingroup\large
Now in the fifteenth year of the reign of the Emperor Tiberius, Pontius Pilate
being governor of Judaea, and Herod being tetrarch of Galilee, and his brother
Philip tetrarch of Ituraea and of the region of Trachonitis, and Lysanias tetrarch
of Abilene, Annas and Caiaphas being the high priests, the word of God came unto
John the son of Zacharias in the wilderness. And he came into all the country about
Jordan, preaching the baptism of repentance for the remission of sins
(Luk. 3, 1--3).
\par\endgroup

Just before the beginning sets in, we are thus called upon to cast a glance at the
condition of the Jewish people at that time, a sad, a most lamentable condition.
Divided and parceled out between the tetrarchs and the Roman governor, the little
kingdom that through so many centuries has striven for its independence, and
through so many dangers and devastations has in wonderful ways been saved and
raised up again by the hand of Providence, is now nevertheless extremely near its
abyss. Everything points to an impending dissolution. --- The old Law does indeed
still stand: Herod has adorned the temple of Jehovah. But the spirit has departed,
and the strength has vanished. Proud Pharisees, with long prayers and broad
borders, sit in Moses' seat, and are zealous for the faith of the fathers, and teach
the people the many commandments. It is a shadow play: the many commandments are
commandments of men, and the empty words do not call
% --- p. 177 ---
\opage{177} the dead to life. The cold Sadducee despises the comfort of religion and the
faith of the multitude; he has enough in his own virtue: no superstition, no
prejudice shall shut him out from the enjoyment of the pleasures of culture. But the
people itself, how does it fare? It suffers violence and does violence: its heart
bleeds. Without teachers and without leaders, the unstable multitude is \emph{as
sheep upon the mountains that have no shepherd.} It is a sad, an almost intolerable
condition. But what then is to become of this people? Yes --- ``what will happen,
then?'' This question wanders like an uneasy thought about Jerusalem and the whole
land, alarms the earnest and entertains the frivolous. What will happen, then; for
it is a matter of freedom and salvation?

If only the God of the fathers were willing as the sons of the fathers are! ---
``then the hero of freedom would come like the lightning, for a revolt smolders in
the depths.''

If only the God of the fathers were willing as the sons of the fathers are! ---
``then Zion would rise above all mountains, and a kingdom would arise such as was
never before seen on earth.''

If only the God of the fathers were willing as the sons of the fathers are! ---
``then the Messiah himself would come before cockcrow, in the next night watch: it
is indeed the last hour, it is a matter of freedom and salvation.'' --- He must
surely come. ``Yes, he shall come! Whether he comes from above, as \emph{the Son of
man in the clouds of heaven}, or springs up from \emph{the root of Jesse}, as a
\emph{shoot on the old stem}: it is all the same, if only He comes. It is all the
same, if only He can \emph{speak unto the kings of the earth in his wrath, and
terrify the heathen in his fury.} It is all the same, if only He \emph{smites}
these enemies of ours \emph{with his rod of iron, and dashes them in pieces like a
potter's vessel!}''

And see! He has come. --- The God of the fathers has kept his
% --- p. 178 ---
\opage{178} promises; but the sons do not understand it. The deliverer has come: but they do
not know it. \emph{He stands in the midst of them; but --- they know him not.}

But what will happen, then, now that it is a matter of freedom and salvation?

The terrible will happen. This people will crucify the Messiah. This people will, in
the uproar of the passions, attempt deliverance by its own strength; it will end in
despair: the destruction of Jerusalem will become for history a memory of horror, to
which men attach the notion of unspeakable misery.

``But --- before this happens --- the wonder will happen! Before the destruction
happens, the Lord of glory will reveal himself in the midst of his people. Before
this old temple, which is made with hands, is torn down, He raises up the new
temple, which is not made with hands. But --- the old is torn down: the veil is
rent, the altar of sacrifice burns, the walls of Zion fall: it is the fulfilling of
the Law, it is to the world a portent that the kingdom of God has come with freedom
and salvation.''

If we now stand firm here on the ground of actuality, and look back with the eye of
posterity upon the first beginning, upon the wondrous portents by which it is
signified to the Believers that the kingdom of heaven draws near, then we too gain
an immediate impression of the significance of the Fore-Gospel in the sacred
history.

The condition of the Jewish people at that time we have considered. This sad
actuality is now, so to speak, the curtain through whose rents the Evangelists Luke
and Matthew let us see single manifestations of the quiet kingdom of blessedness;
just as at night one sometimes sees --- a single star through the rents of the
drifting clouds.

% --- p. 179 ---
\opage{179} %
\textbf{The portents come and vanish; for the kingdom of heaven is at hand.}

\emph{In Jerusalem, in the sanctuary of the temple, an old priest stands before the
altar of incense. He feels so strange, so deeply moved --- alone with God in the
Holy Place. Then comes the vision, a divine vision! It terrifies him, fear falls
upon him: see, the angel of the Lord stands on the right side of the altar of
incense! It is the forewarning of the prophet of the Highest, the second Elijah.}
--- But the multitude of the people, waiting outside at the hour of incense, grows
uneasy, and begins to ask itself why the priest is tarrying in there. ``Surely no
mishap can have befallen him? --- Surely the old man has not fallen into a faint?
--- No, now he is coming out of the sanctuary. But what a peculiar look! He makes a
sign: he is dumb! --- Something extraordinary must have happened to him; he must
have seen a vision in the temple.'' So the multitude speaks, and disperses, and ---
forgets the occurrence.

\textbf{The portents come and vanish; for the kingdom of heaven is at hand.}

\emph{Is there yet a place on earth to which the curse has not reached? See, it is
She, the pure virgin with the pure heart! And the angel comes in unto her and
greets her; and they speak together, the two of them, as soul to soul, inspired by
God. --- Why, it looks as in the regions of Paradise, in the Garden of Eden, before
the cherub drew his sword!} --- But it must surely be a fancy; we are not, after
all, in Paradise. This is Nazareth, the little town of Nazareth by the Sea of
Tiberias.
% --- p. 180 ---
\opage{180} Surely nothing extraordinary can have happened here; not a single person in the
town has learned the least thing. --- ``Could it perhaps be the rumor that is going
about concerning Her who is betrothed to the carpenter Joseph? Quite right! There
is a suspicion. Still, it will no doubt smooth itself over.'' And see! it really
does smooth itself over. Joseph takes to him Mary, his betrothed wife; civil
morality is satisfied; why, just see how it smoothed itself over! ``Such incidents
one should not reckon so strictly.''

\textbf{The portents come and vanish; for the kingdom of heaven is at hand.}

It is a great joy that befalls the house of Zacharias in the city of Juda.
Neighbors and kinsfolk hear that the Lord has shown great mercy upon Elisabeth; and
they gather, and rejoice with her. --- ``But it is strange, though, that the child
is not to be called by his father's name!'' --- Zacharias writes on his tablet:
\emph{His name is John;} --- and they all marvel. See, then comes the portent:
\emph{The mouth of the dumb man is opened, his tongue's bond is loosed, he speaks
and praises God. So at last it becomes known to all that a real wonder has
happened here. And it is noised abroad on all the hills of Judaea, and all that
hear it lay it up in their hearts.} --- Yes, it is noised abroad! --- ``The old
man was as if transfigured with joy: he looked at his little son and was so glad.
Could not joy loose the bound tongue? Might it not after all have come about
naturally that the joy of a father's love opened the mouth of the dumb?'' ---
However, the world has other things to think about.

A decree has gone out from Caesar Augustus to the governor in Syria that the whole
land shall be taxed, and every man appear in his own city to be enrolled. --- In
Bethlehem,
% --- p. 181 ---
\opage{181} the city of David, the influx is so great that poor folk find no room in the inn.
--- ``A young mother has been delivered of a son; she wraps him in swaddling clothes
and lays him in a manger.'' --- How strange! --- ``Is it then anything to wonder at,
that the poor mother must lay the child in a manger when there is no room in the
inn?'' --- O what hope, O what joy! --- ``Is it now anything to rejoice over, that
a mother of the common people has been delivered of a son: the poor man's heir is
not exactly wont to awaken any great general joy.'' --- But behind that busy day
the quiet night came up. \emph{It is a wonder after all, it is a joy after all.
See, the heavenly powers marvel: to them it is a wonder of blessedness. God's
angels in heaven rejoice; they will not hide their joy. --- ``It must be proclaimed
to the fallen race that the Savior of the world is born.'' But the children of men
heed it not: they sleep and dream --- of the day's toil. --- See, these shepherds
are awake, though; they keep watch by night over their flock. ``It must then be
proclaimed for the shepherds' comfort that the Savior was born this night; it must
be proclaimed to the poor men that the joy shall come to all the people.''
--- And the shepherds in the field hear the angel's message, and --- the vision
shines forth in the deep night: a multitude of the heavenly host! Holy strains in
the silent night: Glory to God in the highest!}

\textbf{The portents come and vanish; for the kingdom of heaven is at hand.}

\emph{``And on earth peace?''} --- On the actual earth there is no peace. What,
pray, is the peace of the earth? A gleam of sun on a stormy day! And there is no
\emph{``good will in men''} at all. Who can now take pleasure in a Caesar Augustus
who
% --- p. 182 ---
\opage{182} would tax all the world? And \emph{``a multitude of the heavenly host?''} Yes,
a blessedly fruitful multitude of superstition. Edifying stories from some
shepherds who do not know what they themselves are saying! --- ``A shining
apparition in the air; the cries of some birds of passage --- see, that is no doubt
the whole of it.''

\textbf{But the portents come and vanish; for the kingdom of heaven is at hand.}

Twelve years pass; they soon vanish (and yet in that time much can be heard, and
much forgotten).

It is the feast of the Passover that is being kept in Jerusalem. The temple is
surrounded by halls. In these halls the scribes have schools, which are called
temple schools. In one of these temple schools it now happens that a half-grown
child steps in among the teachers, and hears them, and begins to ask them
questions. --- ``What questions and what answers from a child of twelve!''
\emph{All that hear Him are astonished at his understanding and answers.} ---
``But who can have prompted this child to seek teachers here in the house of
God?'' --- No one knows. It is probably his own impulse; he has no doubt gone off
from his parents. ``This too will no doubt be cleared up in the end.'' --- Yes,
certainly it is cleared up. See, there come his parents. ``How anxious the mother
has been for her dear son! She fetches him with a gentle reproach.'' --- It is a
carpenter's son from Nazareth. --- Splendid gifts! ``If he went to the school of
the learned, he might become a great rabbi.''

But he does not go to the school of the learned. \emph{Now he goes down with his
parents, and comes to Nazareth, and is subject unto them. So he grows up, and
increases in wisdom and stature, and in favor with God and man.} --- ``A fine
testimonial,'' say the learned to all
% --- p. 183 ---
\opage{183} whom they meet from Nazareth, ``a very fine testimonial. But if one would be the
people's teacher, one must first let oneself be taught; if one would act upon the
world, one must know its way of thinking and make use of the guidance of the
competent.'' --- Now He is nearly thirty years old. --- ``The best time for
learning is thus past: he will never become a true rabbi after all.'' --- ``No,
mankind now knows a little more than this self-taught rabbi!'' criticism repeats, to
the honor of the modern consciousness\footnote{%
Die geistig weltliche Bildung Jesu überschreitet dasjenige nicht, was bei
reichen Anlageu in den Volksschulen Galiläas und auf den Volksversammlungen zu
Jerusalem erworben werden konnte; andre, unter gebildeten Völkern und in
glücklicheren Zeiten geboren, waren ihm darin überlegen. \textit{Dr.} Karl Hase.
Kirchengeschichte. Leipz. 1834. S. 25 \S\ 83.%
}. But the Evangelists lift the curtain, point to the portents, and speak to us in
the language of the Spirit: \emph{Come and see!}

Yes, come and see why this teacher, this One, can teach all, without himself being
``taught.'' Come and see why this master, this One, can move the thoughts of a
world without conforming to the world's way of thinking. Hear it, then, and
understand! When this One lifts up his voice and begins to speak, when He stretches
out his hand and begins to act: then the kingdom of heaven is at hand, then there
takes place a true, an actual beginning.

\begin{center}
  \rule{0.22\textwidth}{0.4pt}
\end{center}

\chapter*{Second Division\\[4pt]
  \large The Sacred History: The Life of Jesus in the World}
\addcontentsline{toc}{chapter}{Second Division. The Sacred History: The Life of Jesus in the World}
\markboth{Second Division}{Second Division}
\label{ch:afsnit2}

% --- p. 184 ---
\opage{184} %
\begin{center}
  {\footnotesize While I am in the world, I am the light of the world (Joh.\ 9, 5).}\\[10pt]
  \rule{0.22\textwidth}{0.4pt}
\end{center}

\section*{\textit{A.}\\[4pt] The Redeemer.}
\addcontentsline{toc}{section}{A. The Redeemer}

\begin{center}
  \rule{0.08\textwidth}{0.4pt}
\end{center}

\subsection*{\textit{I.}\\[4pt] The Forerunner.}
\addcontentsline{toc}{subsection}{I. The Forerunner}

It goes without saying that the modern consciousness, which cannot at any point in
the gospel history reconcile itself with the evangelical conception, cannot approve
either of the use of the Fore-Gospel by which the Believer seeks to make
intelligible to himself the relation between the public appearance of the Baptist
and that of the Savior.

``It is not even so difficult to get something put in place that resembles a
transition from the private to the public life, especially if the investigator, in
place of simple proof, will make do with a half-fantastic piece of declamation. It
is, after all, a fairly easy matter to point out a sort of agreement between the
mythical legends of the pre-history and the proper, so-called historical division,
since the gospel narratives from beginning to end undeniably preserve the same
half- or wholly mythical stamp, and
% --- p. 185 ---
\opage{185} everywhere display the same interesting mixture of --- Dichtung und Wahrheit.
Only he who is able to rise to the scientific summit of the present age will be in a
position to separate these components. When the process of separation is
completed, it becomes the great task of science to prove how the composition
originally came about. This much stands firm: that the supernatural facts are all
without exception invented. This must, however, be understood scientifically: they
are of course not invented in the way one invents in the present age. The modern
poets know very well that a poetic production is only a fiction. But this has not
been the case with the evangelical poets. Those authors, unknown to us (the names
Matthew, Mark, Luke, John are only names!), themselves believed in their own poem
as in a revealed truth. Further it must be noted that the narratives were not
invented by single individuals, but by whole congregations. These congregations
formed themselves, no one knows how, little by little in the various regions,
partly of Jewish, partly of heathen extraction. That the Christian Jewish
congregations must have invented somewhat differently from the congregations of
heathen origin is easy to prove.\footnote{%
  Cf.\ e.g.\ C.\ H.\ Weisse. Die evangelische Geschichte. Vol.\ 1. Pp.\
  140--232. Die Sagen von der Kindheit des Herrn.%
} The many different contributions that have come forth in this ingenious manner
were no doubt from the beginning in some obscurity and confusion, but were
afterwards worked over by individual talents, and thus received their present form,
and rounded themselves off into a cycle of clever, baroque, childlike myths rich in
ideas. The Church's four canonical Gospels are for the most part
% --- p. 186 ---
\opage{186} composed of nothing but such myths. Yet this is not all. The inventive urge of
the oldest congregations, or whatever one may wish to call it, has of course also
caused a quite noticeable distortion of the actual facts and historical
personalities. The oldest congregation, namely, was exclusively enthusiastic for
Jesus, and adorned his person with every possible perfection. From this it then
followed in turn, as is understandable, that John the Baptist was put in the shade,
and his merits disparaged even to such a degree that he came to play the
subordinate role of the Forerunner; whereas it is hardly altogether improbable that
Jesus himself was for a time the Baptist's disciple. Once the old congregation had
hit upon the idea that John the Baptist --- to whom, for that matter, the theologian
Hase in his `Leben Jesu' gives the honorable testimony that he was `without a
Messiah's perilous greatness, and without a miracle-monger's dubious presumption'
--- was to count only as a forerunner, then the conjecture lies very close at hand
that those stories of birth and childhood were invented in order to show in
particular how far this forerunner, in respect of the gifts of the Spirit, stood
below the true Messiah.''

These ingenious assertions of modern knowledge are now certainly of such a nature
that simple people who, without exactly being initiated into the most learned
proceedings of criticism, have even only an obscure sense of the fundamental
difference between heathen mythology and Christian revelation, will probably have
trouble seeing how it comes about that such assertions can pass current as
emanations of a higher wisdom. The matter can, however, be explained: every person
who feels within himself that God neither can nor will reveal himself naturally does
everything to guard himself against the opposite view. This holds just as well of
the learned as of the unlearned.
% --- p. 187 ---
\opage{187} %
Only the manner in which the rejection takes place is somewhat different. If one
hands the Gospel to the unlearned denier of faith, he says at once, quite candidly
and straightforwardly: ``No, thank you; I have no use at all for the old
superstition. Whoever it was that could make up such stories! what is it to me? I
can well feel in myself that the whole thing is nothing but superstition.'' --- The
learned denier of faith, on the other hand, cannot so easily dispose of the ancient
gospel writings; that would, after all, be both impolite and unscientific. He must
therefore first explain how it all came about naturally with all these supernatural
things. But to explain how it came about naturally with everything supernatural is
certainly a very difficult and laborious task. One ought therefore to hold lofty
science excused when now and then the explanation goes as best it can.

The Believers, who are not at all burdened by such difficulties, hold themselves
assured, on the contrary, that the Gospel is proclaimed to them for comfort and
joy. Living appropriation makes them never tire of contemplating the holy types and
their mutual relation from every side; \emph{for, they say, God, who said that the
light should shine forth out of darkness, hath caused it to shine in our hearts, to
enlighten the knowledge of the glory of God in Jesus Christ} (2 Cor.\ 4, 6).
--- Does not the divine Law say that a servant shall be beneath his master, and
that the Forerunner ought to step back for him who works the consummation! Were it
only in words and phrases that the Fore-Gospel exalted Jesus above John, then
suspicion would have some ground for hinting at human invention and arbitrary
slighting. But from the very beginning we see in the gospel accounts two mutually
different manifestations of the wondrous, the one as an afterglow of the Law in the
days of its old age,
% --- p. 188 ---
\opage{188} the other as a dawn of grace and truth in the fullness of time. The two
annunciations and the two hymns of praise are related to each other as the vanishing
to the abiding. And if we compare the two birth narratives, the joy in the city of
Juda with the joy in Bethlehem, are not these too related to each other as the
vanishing to the abiding? --- The Forerunner must step forth in order to step back
again: he is prepared for it from his youth. In the wilderness one attracts no
notice, and see, he seeks in all stillness the lonely places, and remains in the
wilderness until the day he shows himself to Israel. The Consummator too must be
hidden from the world and grow a quiet growth before he can come forth as the
Redeemer of the people. But already from the cradle his precedence is visibly
marked: the Lord of glory who rests in the manger lets himself be greeted before he
withdraws. --- It is commanded in the Law (2 Msb.\ 13, 2; 3 Msb.\ 12, 3; 4 Msb.\
18, 15) that all the firstborn which in Israel were called \emph{holy to the Lord}
should be redeemed by a sacrifice, and the firstborn presented in the temple. If
this firstborn is now the promised Son of David, then Israel must surely also know
him and come to meet him. And see, it happens (Luk.\ 2, 25--38). \emph{There was a
man in Jerusalem whose name was Simeon; and the same man was just and devout,
waiting for the consolation of Israel; and the Holy Ghost was upon him.} --- This
man came up into the temple just at the right moment; he came up there by the
prompting of the Spirit, he came up there on Israel's behalf, he came up there, you
see, according to the promise that he \emph{should not see death before he had seen
the Lord's Anointed.} --- Now when the parents \emph{brought in the child Jesus, to
do for him after the custom of the Law,} see, the aged Simeon was already in the
temple, and
% --- p. 189 ---
\opage{189} came to meet them, and took the child in his arms, and praised God and said:
{\large Lord, now lettest thou thy servant depart in peace, as thou hast said. For
mine eyes have seen thy salvation.}

So Simeon stands in the temple with the child in his arms, and speaks the wondrous
word of ``the light of the Gentiles'' and ``the glory of Israel''; and the
prophetess Anna, the daughter of Phanuel, the aged widow\footnote{%
  On \emph{Anna's patience in expectation} cf.\ S.\ Kierkegaard: \textit{MDCCCXLIV}.
  Two Upbuilding Discourses (Copenhagen. Available from the bookseller P.\ G.\ Philipsen).%
}, who never departed from the temple, stands by, and mingles her praise with
Simeon's; and both bear witness before the Lord that the Israel of the Law has long
grown gray in expectation, and is full of days, and would gladly depart in peace,
since his eye has seen the salvation that \emph{is prepared before the face of all
people}.

Yes, before the face of all people, before the face of the Gentiles too. See, the
royal star is kindled far off in the East: the Messiah in the manger lets himself
--- be greeted before he withdraws. --- ``Where is the star, then, where? I cannot
find it at all!'' --- If you would discover the sign of the star, you must see with
the eye of faith. --- ``Self-luminous! Such a light has never before lit up the
plains of Chaldaea. --- A wandering star! It beckons and points toward the land of
Canaan; it leads and shines in the footsteps of Abraham! --- So the old prophecy is
true, now the King of the Jews is born.''

\begingroup\large
Matth.\ 2, 1--2. Now when Jesus was born in Bethlehem of Judaea, in the days of
Herod the king, behold, there came wise men from the east to Jerusalem, saying,
Where is this King of the Jews that is born? for we have seen his star in the east,
and are come to worship him.
\par\endgroup

But the wise men from the East come from afar, and are such
% --- p. 190 ---
\opage{190} strangers in Judaea; therefore they cannot now fully interpret the star's sign,
but go somewhat astray, and must first go to Jerusalem, to King Herod, to hear the
explanation. And see, King Herod is indeed willing to furnish the strangers with all
the information they desire.

\emph{And when he had gathered all the chief priests and scribes of the people
together, he demanded of them where Christ should be born. And they said unto him,
In Bethlehem of Judaea: for thus it is written by the prophet} (Mich.\ 5, 1).

Yes, even Herod believes the prophet on the word of the chief priests, and would now
so dearly like to know where the child is, that he too may come and --- show it his
worship. --- And the wise men from the East follow King Herod's counsel and the
scribes' direction, and depart, \emph{and see, the star, which they had seen in the
east, goes before them, and stands over where the child is.}

So the wise men of the East came from afar, and yet came into the house over which
the star stood. And see, they know the newborn King by the shining of the star;
\emph{and they fall down and worship the child, and open their treasures, and offer
it gifts, gold and frankincense and myrrh.}

--- So now there are two kings in the land of the Jews; old Herod will not suffer
it: the Idumaean will hear nothing of the son of David. ``The magi steal away; they
are making sport with the old folk legend! --- Then we must murder the danger in its
swaddling clothes, so that it does not grow up and become great!''

--- \emph{In Rama was there a voice heard, weeping and crying and great
lamentation}\ldots{} ``When a Herod strikes, he hits!'' --- But no! He came a moment
too late: the
% --- p. 191 ---
\opage{191} angel came first! God's angel stood on guard for the son of David, and the
Idumaean --- he came too late.

\emph{But Joseph arose, took the child and his mother by night, and departed into
Egypt. And he was there until Herod was dead, that it might be fulfilled which was
spoken of the Lord by the prophet, saying: Out of Egypt have I called my son
(Hos.\ 11, 1).}

This, then, is the evangelical conception of the Gospel. Scientific criticism, which
prefers another conception, may now continue its researches, to see whether it can
thus comprehend the mystery of the Spirit, whence it comes and whither it goes.
We, meanwhile, will turn with the Believers to the public life, and convince
ourselves that the same relation repeats itself here too between the serving
Forerunner and the royal Consummator. How different, after all, is their coming
forth! --- John himself testifies before the council's envoys (Joh.\ 1, 20):
\emph{I am not the Christ.} And we must believe his testimony. For he who baptizes
with water unto repentance has, after all, only received the Spirit's call to
preach the Law that was given by Moses; and must therefore rebuke all his own:
\emph{The axe lieth at the root of the trees; every tree which bringeth not forth
good fruit shall be hewn down, and cast into the fire} (Matth.\ 3, 10).
But He who himself fulfills the Law, and has all grace and truth, speaks another
word to his disciples after the first greeting (Joh.\ 1, 51): \emph{verily, verily,
I say unto you, Hereafter ye shall see heaven open, and the angels of God ascending
and descending upon the Son of man.} --- He who baptizes with water unto repentance
is himself under the strictest discipline of the Law, a Nazarite who may not drink
of that which grows on the vine, lest he
% --- p. 192 ---
\opage{192} fall into temptation, a Nazarite who must live on the food of the wilderness,
lest he have need of the world. But He who has himself conquered the tempter in the
wilderness, He who baptizes with the Holy Ghost and with fire, see, he may sit down
at table with the wedding guests and --- turn water into wine. --- He who baptizes
with water unto repentance himself admonishes his own disciples not to take him for
the true Master (Joh.\ 3, 26--31): {\large a man can receive nothing at all, except
it be given him from heaven. Ye yourselves are my witnesses, that I said, I am not
the Christ, but that I am sent before him. He that hath the bride is the bridegroom:
but the friend of the bridegroom, which standeth and heareth him, rejoiceth greatly
because of the bridegroom's voice: this my joy is fulfilled. He must increase, but I
must decrease. He that cometh from above is above all: he that is of the earth is
of the earth, and speaketh of the earth: he that cometh from heaven is above all.}

And He who has come from heaven and is above all bears witness in turn to the
Forerunner (when the latter had already ended his work and exchanged his wilderness
for a prison), saying to the people concerning John (Matth.\ 11, 7--10): {\large
What went ye out into the wilderness to behold? A reed shaken to and fro with the
wind? Or what went ye out for to see? A man clothed in soft raiment? Behold, they
that wear soft raiment are in kings' houses. Or what went ye out for to see? A
prophet? Yea, I say unto you, and much more than a prophet. For this is he, of whom
it is written, Behold, I send my angel before thy face, which shall prepare thy way
before thee.}

So then the Baptist is hereby recognized as the true Forerunner,
% --- p. 193 ---
\opage{193} then he also receives his testimony; and in truth, the man who receives
such a testimony from such a mouth is not arbitrarily pushed aside, is not unduly
placed in the shade. But John was not --- the Christ. He neither could nor would
give himself out as such. He worked in the calling to which he was called: there
was a barrier of finitude that he could not break, an allotted limit that he could
not overstep. --- Thus speaks the divine voice that judges with righteousness
(Matth.\ 11, 11): {\large Verily I say unto you: among them that are born of
women there hath not risen a greater than John the Baptist; but the least in the
kingdom of heaven is greater than he.}

Here, then, is his barrier. He is a true Forerunner, but only a Forerunner.

``But'' --- the modern consciousness continues --- ``things still do not really
hang together with this Forerunner. Criticism proves it on various grounds.
\emph{In the first place:} the span of time that the Evangelists allow their
Baptist is far too short, if he is to be assumed to produce so great an effect.
An older critic, Cludius, remarks in a treatise `über die Zeit und Lebensdauer
Johannis und Jesu' that so young a preacher of repentance could hardly have made
an impression, and been taken for a prophet of old, and above all for an Elijah.
\textit{Dr.} Strauß adds: `To this rule there are indeed not a few exceptions;
but to move us to such an exception in the present case, Luke's report 1, 26,
that John was only 6 months older than Jesus, is far too weak, since it is made
entirely in the interest of legendary poetry, and must therefore be given up as
soon as the slightest improbability appears.'
\emph{In the second place:} a young man who wants in any way to break through
cannot possibly permit himself such a
% --- p. 194 ---
\opage{194} language as this Forerunner uses. Heaven preserve us! How is it,
really, that the Evangelists have that man speak? \emph{You brood of vipers!}
What an objectionable expression! it is sheer abuse. Is that any way to behave?
Surely no reasonable man really addresses a whole people or a public like that
when he wants to advance his cause and win public opinion. Such terms of abuse
the estimable John surely never used either; and it is really one of Bruno
Bauer's merits to have pointed this out. \emph{In the third place:} this
Forerunner, after all, stands on an entirely different standpoint in the first
three Gospels than in the fourth. Matthew, Mark and Luke have him appear simply
as the preacher of the Law who expects a Jewish Messiah; but the Evangelist John,
by contrast, initiates him without further ado into the mysteries of his Gospel,
and thus contradicts the others. \emph{Finally and in the fourth place:} the
relation between this John and Jesus is screwed up so high above the usual
conditions of humanity that it has thereby become an altogether ideal and
consequently unreal relation. How improbable, indeed how highly unreasonable is
it not, that John, who after all already had a name before Jesus appeared, should
at once have been so willing to yield his place to the younger beginner that he
even referred his own disciples to him. No, that is really asking too much. There
are indeed exceptions to the rule; but criticism will not grant such an exception
in the present case. Imagine a John the Baptist, a man in his best years, an
excellent man, a man with rare gifts and strong passions! Should such a man in
reality be able to be without a feeling of his own ability, without any
consciousness of his own merit, without the higher, upward-striving ambition that
is the condition of true energy of action? No, that is incredible. A John does
not yield to a Jesus in this way. The actual law of the development of
% --- p. 195 ---
\opage{195} character runs thus: first a tension between two such men; next a
polemic; finally the weaker yields to the stronger. --- In this respect too we
must thank David Strauß and Bruno Bauer for their very significant
elucidations.''

For the judgment of the faith of the Gospel, such misgivings must now fall away
as an empty figment of the brain of a paltry cleverness. --- ``The young preacher
of repentance cannot get time to break through?'' Of course, there are even
``not a few exceptions to this rule!'' But modern criticism ``is nevertheless not
to be moved to grant such an exception in the present case.'' It is not to be
moved to it, it will not! Yes, that is another matter: when it will not, one
certainly ought not to force it either.

The Believers, on the other hand, would gladly make an exception with John the
Baptist. He does not exactly seem to them to be an ordinary preacher of
repentance. No, such a man surely does not need many months to do his errand: he
is certainly well prepared; he will break through, all right. Just trust him:
when he raises his voice in the wilderness, it is heard round about in the
cities. It is, after all, a well-known voice, it is the voice of the Law to the
people of the Law: the people know the voice at once. {\large Behold! Jerusalem
goeth out to him, and all Judæa, and all the region round about Jordan. And they
are baptized of him in Jordan, and confess their sins (Matth.\ 3, 5--6).} ---
Yes, so he does break through after all! Just trust him, the man in the shirt of
camel's hair, when he girds the leathern belt about his loins, and cries out into
the crowds to those who feel so secure because, after all, they have Abraham to
their father: {\large You brood of vipers! who warned you to flee from the wrath
to come? (Matth.\ 3, 7)} --- then one understands him at once, then he truly
breaks through.
% --- p. 196 ---
\opage{196} %
Pharisees and Sadducees understand him at once: it is they, after all, whom he
strikes. And the word sinks into the people, and they understand at once that he
is a great prophet who preaches the Law to life. For, is it not so? --- the words
of the Law are, after all, so comprehensible that every human being can
understand them without long explanation.

When, then (Luk.\ 3, 12--13), publicans come to John to be baptized, and say to
him: \emph{Master, what shall we do?} and the answer runs: {\large exact no more
than that which is appointed you;} should the conscience of publicans then need
many months before it learned to understand this? --- Or when (Luk.\ 3, 14) the
soldiers too, who come to be baptized, ask him: \emph{what then shall we do?} and
he says to them: {\large do no violence, use no deceit; and be content with your
wages;} should the conscience of the soldiers then need more than ``half a
year'' to grasp these commandments! --- No, --- \emph{the axe is laid unto the
root of the trees!} and when the axe is laid to the root of the trees, one
easily understands the baptism of repentance. Then the waters of the Jordan are
stirred, then John breaks through as in an instant, when conscience, awakened by
the stern word, proves to every human being that the prophet is right.

The old Elijah was a zealous prophet; his spirit was fire, the fire burned in
him, the flames refined him: he went up to heaven in the blaze of fire. The
second Elijah is a zealous prophet; his spirit is fire, the fire burns in him,
but --- the flame does not blaze up toward heaven: it scatters its sparks on the
bank of the Jordan. And consciences are kindled, and the people are in
expectation, and --- all men muse in their hearts of John, \emph{whether he were
not the Christ.}

See, so this Forerunner can break through after all. But, while he breaks through
so powerfully, he still does
% --- p. 197 ---
\opage{197} not push himself forward, does not overshoot the mark, does not
forget his errand. --- To push forward like a rushing river, not to halt but to
plunge, and then to end in the ambiguity that to halt is, after all, to plunge,
see, that the world can understand: it happens so according to the order of
nature. But to soar as in an instant toward Elijah's heaven, and then to sink and
descend, with a servant's humiliation, that the world cannot understand, that is
a little above the order of nature. And the man in the shirt of camel's hair,
with the leathern belt about his loins, how tall and strong he is! see, he stands
as if in the sand of the wilderness, and yet a little above --- the order of
nature! --- And John answered, and said unto all: {\large I indeed baptize you
with water, but one cometh that is mightier than I, the latchet of whose shoes I
am not worthy to unloose; he shall baptize you with the Holy Spirit and with fire
(Luk.\ 3, 16).}

The second Elijah is a zealous prophet, His spirit is fire; the fire burns in
him, so that he must prepare the way for his Lord, yet without standing in his
way. It is with fire and power that the Baptist threatens the \emph{brood of
vipers that thinks it can escape the wrath to come.} It is with fire and power
that before the assembly of the people he casts forth the image of the judging
Messiah: {\large His fan is in his hand, and he will thoroughly purge his floor,
and gather the wheat into his garner; but the chaff he will burn with fire
unquenchable (Luk.\ 3, 17).} It is with fire and power that in the humility of
the heart he appropriates the manifestation at Jesus' baptism: {\large And I knew
him not; but he that sent me to baptize with water, the same said unto me: upon
whom thou shalt see the Spirit descending, and remaining on him, the same is he
which baptizeth with the Holy Spirit. (Joh.\ 1, 33).}

John is absorbed with the whole strength of his soul in each of the various
moments of his
% --- p. 198 ---
\opage{198} short life. If he looks at the recalcitrant Pharisees, at the brood
of vipers, then it is as though he would be consumed in the zeal of wrath; but if
he sees Jesus coming, \emph{the Lamb of God, who bears the sin of the world,} is
it not then as though the hard man would dissolve in worshipping admiration? ---
{\large This is he of whom I said: after me cometh a man which is preferred
before me; for he was before me (Joh.\ 1, 30). He must increase, but I must
decrease (Joh.\ 3, 30).} This the modern consciousness cannot understand, and
therefore it is said of the Evangelists' John that he changes standpoint. But ---
the man in the shirt of camel's hair, with the leathern belt about his loins,
does not know at all what it means to change standpoint. This man is as of one
casting, with one will and one aim: it is the Forerunner's will, the Forerunner's
aim. There is an urge, a haste, a hurrying zeal in this man, which assures us
that in a short time he will end his day's work, and step back, and receive his
reward, and come to rest. This happens not only in such a way that the one thing
follows after the other, but also in such a way that the one follows from the
other.

He ended his day's work. He called the word of the Law to life in the conscience
of the people, and the people became assured that the man who had baptized Jesus
of Nazareth in the river Jordan, and borne witness to him, was a truthful prophet
whose testimony could not be shaken (Matth.\ 21, 25---26). He ended his day's
work. His sharp word penetrated to high and low, it forced its way into the
prince's ear, it fixed its sting in Herod's conscience. It was no word of
rebellion that the prophet of the wilderness had spoken against the tetrarch in
Galilee. John had not said to Herod Antipas: ``come down from thy high seat, for
the children of Abraham will suffer no princes.'' No, thus a Forerunner of
Christ does
% --- p. 199 ---
\opage{199} not speak. It was not John who would break the human laws; but it was
Herod who had committed a breach of the divine Law. The men who wore ``the soft
clothing'' kept silent about it, and no doubt had to keep silent. But the man in
the shirt of camel's hair could not and dared not keep silent; for he said it
right out: {\large It is not lawful for thee to have thy brother's wife (Mark.\
6, 18).} --- And the word bore its fruit, it did not fall powerless to the
ground. Herod could now understand that John was a prophet, and he appreciated it,
and rewarded him for it: he had him cast into prison. There is earnestness in this
recognition; it is in truth a royal reward: a Herod Antipas could not, after all,
reward a John the Baptist otherwise. --- ``The prophet has worked enough, now he
must go to rest.'' --- But the rest of prison and a restless spirit that burns
with zeal for the coming of the kingdom of heaven do not agree well. Even the
greatest among them that are born of women must surely have his dark hour. Then
John receives another reward, that he may come to rest. For thus runs the
Savior's greeting: {\large go and shew John those things which ye hear and see.
The blind see and the lame walk, the lepers are cleansed, and the deaf hear, the
dead rise up, and the gospel is preached to the poor (Matth.\ 11, 4---5).}

When John in his prison hears this greeting, then he surely comes to rest. Now
his day's work is ended, now he has nothing more to fear in this world. ---
``Death!'' --- The prophet of the Most High does not fear death. --- Yes, death,
when might it come? In a lonely prison time grows long, especially for one who
--- waits. --- Should there then be no one so sympathetic as to think of this?
Should the world in Israel now really have forgotten to give a prophet ---
death? --- In the great world one forgets so many things, and seldom thinks of
the
% --- p. 200 ---
\opage{200} imprisoned man. Herod no doubt has his own thoughts at times; but ---
he drowns them in wine, he gives them away for a dance.

Gospel, Mark.\ 6, 21---28.

\begingroup\large
And when a convenient day was come, that Herod on his birthday made a feast to
his lords, and the high captains, and the chief men of Galilee; and the daughter
of Herodias came in, and danced, and pleased Herod and the guests, the king said
unto the damsel: ask of me whatsoever thou wilt, and I will give it thee. And he
sware unto her: whatsoever thou shalt ask, I will give it thee, unto the half of
my kingdom.
\par\endgroup

But then there is no one here at all who thinks of the prophet. The daughter of
Herodias? She does not think, after all: she dances! See, ``when one dances, one
has no time to think of an imprisoned man''. No, the girl does not think for
herself, but --- she goes in to her mother. Then Herodias wakes up; she was
sitting so alone by herself --- in thought. She was thinking just then of the
prophet's word: \emph{thy brother's wife!} --- And the daughter asks her mother:
\emph{what shall I ask?} --- {\large Not the half of the kingdom --- that is too
little --- the head of John the Baptist on a charger!}
\par\noindent{\small (cf.\ Matth.)}

\begingroup\large
And she came in straightway with haste unto the king, and asked, saying: I will
that thou give me straightway in a charger the head of John the Baptist. And the
king was sorry; yet for his oaths' sake, and for the guests' sake, he would not
reject her. And straightway the king sent one of the guard, and commanded his
head to be brought. And he went and beheaded him in the prison; and he brought his
head in a charger, and gave it to the damsel, and the damsel gave it to her
mother.
\par\endgroup

--- So even this head came after all into \emph{kings' houses!} --- ``A little
late, if you will, for the man did not exactly wear \emph{the soft
% --- p. 201 ---
\opage{201} clothing.}'' --- He was God's prophet, Herodias! --- ``Now he is like
an image of death.'' --- Herodias, he was a stern prophet! --- ``Now he is mild as
an angel, and keeps silent.'' --- But, Israel, He was your prophet!

\emph{Yea, I say unto you, and much more than a prophet. For this is he, of whom
it is written: behold, I send my angel before thy face, which shall prepare thy
way before thee.}

So then He received his reward. So then He came to rest. He heard, after all, the
Bridegroom's voice. Now His joy is fulfilled.

\subsection*{\textit{II.}\\[4pt] Christ among His Own.}
\addcontentsline{toc}{subsection}{II. Christ among His Own}

\begingroup\large
He came unto his own, and his own received him not. But as many as received him,
to them gave he power to become the children of God, even to them that believe on
his name\dots\ (Joh.\ 1, 11---12).
\par\endgroup

``But'' --- the modern consciousness continues --- ``if Christ could not even
break through with his own, and win a reasonably complete recognition among his
countrymen and contemporaries: how can one then expect that, after so many
centuries' many hundreds of misunderstandings, he should be in a position to
procure himself any victorious recognition in the present? --- To be sure, it is
said: `he that believeth and is baptized shall be saved, and he that believeth
not shall be damned'; likewise it is said `that there is given to us men no other
name for salvation than the name of Jesus, and that God hath highly exalted him,
and given him a name which is above every name, that at the name of Jesus every
knee should bow, of those in heaven and on earth and under the earth' etc.\
etc.; but if this is now really
% --- p. 202 ---
\opage{202} so, then --- good night, poor reason: for you there is no rescue! ---
Whether it is conceivable that Christ's appearance among the Jews alone can be
enough to save all rational creatures in the whole universe; and if not, whether
it is then conceivable that on every inhabited planet (for that all the globes up
there are like empty rooms the theologians cannot, after all, make us believe) a
particular manifestation of the same Christ is to take place; whether it is
conceivable that no fall into sin has occurred on any of the other globes, so
that Christ need merely pay a fleeting visit to each one of them in particular in
order at once to be known as by `his own'; or, if all the planet-dwellers are
sinners, like us, whether then the same Christ is to be scourged, mocked,
crucified and die the death of atonement on each single globe in all the solar
systems --- that would become an enormous history of suffering, --- however, we
will not dispute about that now, since a matter that lies so far from us can,
after all, easily be dismissed with a \textit{non liquet}. --- Nor shall we take
issue either with the theologians or with our Lord as to why so many generations
on earth had to die without salvation and blessedness before the Redeemer could
finally come and his name be spread abroad. We gladly grant the teachers of the
Church all the apologetic advantages that can be drawn from that fleeting
utterance (1 Pet.\ 3, 19) that the dead Christ `preached unto the spirits in
prison'. We concede that in less than three days he may have converted all the
heathen souls that had passed away in death from `the flood in the days of Noah'
until `the fifteenth year of the reign of Tiberius Caesar', if only the
theologians will in return give a little proper account of the one main point of
redemption that concerns ourselves directly. --- Is it really true that all
salvation and perdition turns exclusively on a single name in heaven and on
earth: how then are we to understand that this great
% --- p. 203 ---
\opage{203} signature of redemption has been so mystified that no one, not even
the most honest inquirer, can with any certainty decipher its hieroglyph? Let us
not dishonor the divine reason by setting up so evangelical an idol. `Were my
salvation and damnation solely and entirely dependent on my relation to the
historical-personal Christ, must I not then expect to see this image of Christ
stamped on history with such pure, indelible features that I could not fail to
recognize it without failing to recognize all history? must I not expect to see
the essence of the Godhead shine forth from this image so unmistakably and
clearly that I could not deny it without becoming conscious that I was
deliberately denying the truth?' But this is, as everyone knows, not the case at
all. The Gospels are neither complete nor exact. On the contrary! the historical
image of Christ is a colorful sketch, certainly, but still a fleeting one,
treated, as it seems, even with a certain inspired, simple-minded carelessness.
The divine that individualizes itself in this image is stamped so hard that it
would in truth be altogether indigestible, if the power of reason were not able
to dissolve the composition and separate the substance from the contingency of
the phenomenon. --- And yet one dares to assert that a human being's salvation
and damnation shall be decided by the person of Christ! Fortunately Scripture
contradicts itself; fortunately Christ himself (Luk.\ 23, 34) forgave those who
mistreated him, because they did not know him; fortunately the first among the
Apostles (Ap.\ Gj.\ 3, 17) himself acknowledged that the Jews crucified his Lord
out of ignorance; fortunately, then, one can, despite the threats of the Church,
be ignorant of this Lord without being damned; fortunately one can, then, live
rationally and die calmly, and (if there is such a thing as salvation) live again
and be saved without exactly troubling oneself about the name of Jesus.''

Even though we must admit that the number of those who with such
% --- p. 204 ---
\opage{204} self-assurance could resolve to reject or mock the name of Christ has
not yet gained the decisive preponderance in Christendom, the Evangelist's
testimony: \emph{He came unto his own, and his own received him not,}
nevertheless retains its original significance for the present time. --- There
is a fairly widespread prejudice which, in all brevity, is set forth thus: ``The
Jews, who had Christ himself in their midst and were eyewitnesses to everything
that happened through him, must surely have had every possible opportunity to
assure themselves of what he was really like, whether he really could perform
miracles, whether he really was also pure and without sin, in short: whether he
really was in every way the one he said he was. It is otherwise with us who live
in the present. Show us Christ as he is! we shall then soon declare whether we
will believe him or not. But this is now a sheer impossibility. We do not, after
all, have the actual Christ among us, and are therefore legitimately excused if
we do not place ourselves in any altogether definite relation to him.'' This
prejudice is very reassuring, and consequently very widespread. By keeping the
matter suspended in a certain indefiniteness, one is prepared for more than one
eventuality. Suppose that the Gospel is a deception; then one does not lose much
by it, since after all one does not take that matter particularly to heart.
Suppose, conversely, that the Gospel is a truth, that Christ really will come one
day, as he himself has said, to judge the living and the dead: well, then one can
after all make him an excuse, and say it straight out, that one was not quite so
well prepared because the Gospels were so inexact and theology harbored so many
doubts in the critical introduction. When a prejudice is so convenient and at the
same time so soothing, what wonder, then, if it is both deeply rooted and very
widespread? --- That it is nonetheless a paltry prejudice, the Gospel itself
proves.
% --- p. 205 ---

\opage{205} It is thought that the Jews, Christ's own, found it so easy to assure
themselves of the truth in his person because he himself spoke with them and was
bodily present among them. This advantage is certainly very ambiguous. To be sure,
the Apostles were eyewitnesses, and could appeal to --- \emph{that which they had
heard, that which they had seen with their eyes, that which they had looked upon,
and their hands had handled} (1 Joh.\ 1, 1); but Caiaphas and Pilate, who examined
him, the servants who scourged him and struck him in the face, were also
eyewitnesses, could also see him with their eyes and handle him with their hands,
and still they did not become Apostles. --- To be sure, it was a great advantage
for the disciples that they could question him personally as their Lord and
Master, and hear the word from his own mouth; but the Pharisees and the scribes
had the same advantage, after all, and still they did not become his disciples.

That some time had to pass before his own properly became aware of him and came
in earnest to ask themselves the question: \emph{who then is this?} may readily be
granted. But it did not last very long, after all, before he became an object of
everyone's attention. Indeed, the chief men themselves in the high council even
meddled in the matter at last, to be on the safe side, and took him into
examination. Why does the authority take a human being into examination, except
to inform itself of who he is, where he is from, and what it is he is up to?

So Jesus at last came under examination, and was brought before the high priest
Caiaphas, where all the scribes and elders were assembled; so Christ himself had
to appear in his own person before the council. And the whole council stared at
him, all eyes shone toward him, and yet there was no one in the council who could
recognize him. And witnesses were heard, that they should
% --- p. 206 ---
\opage{206} explain who this man was; but they could not explain it: their
testimony did not agree. And the high priest stood up and said: \emph{I adjure
thee by the living God, that thou tell us whether thou be the Christ, the Son of
the living God.} And Jesus did not keep silent; but answered (as was fitting,
since he was, after all, under examination): \emph{thou hast said it; nevertheless
I say unto you: hereafter shall ye see the Son of man sitting on the right hand of
power, and coming in the clouds of heaven.} So the assembled council heard his own
answer; but this answer was not recognized as valid in court. --- \emph{What think
ye?} But they answered and said: \emph{He is guilty of death.} So near was Christ
now to these elders, and yet --- so far away! Yes, we hope and believe that the
Son of man who now sits at the right hand of power is not so far away from any of
us as he was far away from Caiaphas when this high priest rent his clothes and
said: \emph{He hath blasphemed God; what further need have we of witnesses?
Behold, now ye have heard his blasphemy.}

``But'' --- objects prudent sobriety --- ``one cannot reckon so exactly with the
embittered Caiaphas. Everyone knows, after all, how hatred can blind. Caiaphas was
partial, and naturally had to be so by his character and position.'' --- Let us
then follow the course of the narrative, and accompany Jesus to the judge who did
not share the high priests' passionate hatred, since he was no Jew; to the judge
who would so gladly let the strange prisoner go, since he finds no fault in this
man; to the Roman governor Pontius Pilate.

So the Savior of the world stands in Pilate's judgment hall. The examination
takes a calmer course: might Pilate not now recognize
% --- p. 207 ---
\opage{207} him? An important question is put, a remarkable answer follows: might
Pilate not, after all, discover who the prisoner is who now stands directly before
him. --- \emph{Art thou not a king then?} --- \emph{Thou sayest that I am a king.
To this end was I born, and for this cause came I into the world, that I should
bear witness unto the truth. Every one that is of the truth heareth my voice.} ---
And Pilate hears; but --- he does not recognize the voice. It is so foreign a
voice to him; it is as though this voice came from a dark depth into which no
light can yet penetrate. --- \emph{What is truth?} --- So Christ was so near his
judge, and yet --- so far away! Yes, we hope and believe that the King of truth,
who was born to this end and came into the world for this cause, that he should
bear witness to the truth, is not so far away from any of us as he was far away
from Pilate when this doubting Roman asked: \emph{what is truth?}

``But'' --- objects prudent sobriety --- ``one dare not reckon so exactly with a
Pilate. He was, after all, a foreigner, who did not heed and perhaps hardly even
knew the Jewish expectations. It is really asking too much that a Roman governor
should see a Savior of the world in a Jewish prisoner. Pilate did for him what he
could under the circumstances.'' --- Let us then follow the narrative a step
further. Pilate, after all, leads the prisoner out before the people. Should
there not be someone in the crowd who could recognize him, and say a word for his
justification? --- \emph{He had, after all, spoken openly to the world; he had
always taught in the synagogue and in the temple, where the Jews came together,}
--- should there then be no one present among the Jews who could recognize him?
should the people not acknowledge the Prince of Peace whom only a few days before
it had greeted with palm branches and cries of
% --- p. 208 ---
\opage{208} Hosanna? --- But there is no one at all who will acknowledge him: he
is as though transformed before the eyes of the crowd. Before, they thought they
saw a great prophet in him; now they do not recognize him at all. Pilate asks the
people: \emph{Will ye that I release unto you the King of the Jews?} But all cried
out as with one mouth: \emph{Not this man, but Barabbas!} But Barabbas was --- a
robber. --- What help is it, then, that the Messiah, the King of the Jews, is
present to the senses among his own, when his own nevertheless see in him a
criminal who is worse than a robber?

``But'' --- objects shrewd sobriety --- ``one cannot reckon so exactly with the
crowd. It is, after all, not enough that the Messiah is the King of truth in word
and deed. If he is to be recognized by the great crowd, then still other,
outwardly visible marks are required. See, that is why actual kings have both
scepter and crown, and wear a purple robe when they show themselves to the
people, so that by such conspicuous and significant signs they can work upon the
great multitude.'' --- Let us then follow the narrative; for, in truth, it goes
to the utmost!

The Roman governor and his soldiers ``do for the King of truth what they can
under the circumstances.'' \emph{Behold! they strip him and put on him a purple
robe. And they plait a crown of thorns and put it upon his head, and a reed in his
right hand; and --- they fall on their knees before him and worship him: hail to
Thee, Thou King of the Jews!}

``But'' --- objects prudent sobriety --- ``one surely dare not reckon so exactly
with these soldiers. They were, after all, coarse men: who pays any heed to the
very lowest mockery of coarseness?'' --- But --- sobriety against sobriety! ---
it is in truth an ingenious mockery: there really does seem to be meaning in the
% --- p. 209 ---
\opage{209} homage these soldiers pay their prisoner. If Christ is a king, then it
also befits him, after all, to wear the royal insignia: the crown and the scepter
and the purple robe! Who dares say of all this that it is only an invention of
coarseness? Is the purple robe perhaps not pure and clean? why, he dyes it with
his blood! Is it then no genuine crown that he wears? see, it is a crown of
thorns! alas, the King of truth can wear no other crown in this world. --- So the
actually present Christ has all the outward marks that befit him. --- \emph{And
Jesus came forth, and wore the crown of thorns and the purple garment; and Pilate
set him forth before all the people and said: behold, what a man!} --- But the
people could not see at all now; the Jews could not recognize their own King,
and therefore they all cried with a loud voice: \emph{crucify! crucify!} --- He
fulfilled the Law, and still those learned in the Law could not recognize him,
but justify themselves before Pilate and say: \emph{we have a law, and by this law
he ought to die, because he made himself the Son of God.}

\begingroup\large
--- Blessed are the eyes which see the things that ye see. For I tell you that
many prophets and kings would have seen the things which ye see, and have not seen
them, and heard the things which ye hear, and have not heard them (Luk.\ 10,
23---24).
\par\endgroup

That these words of Jesus to the disciples must refer to a different seeing and a
different hearing than the merely outward and sensible is now, according to our
preceding consideration, surely easy to see. Instead of wondering at the
blindness of those Jews, we must much rather ask ourselves how it was possible
that even only some among his own could have recognized him in the form of a
servant, received him as the Lord of glory, and believed on his name. To this
question the Evangelist (Joh.\ 1, 12---13) gives us
% --- p. 210 ---
\opage{210} %
{\large the answer that he gave to them that received him in faith power to become
the children of God, which were born, not of blood, nor of the will of the flesh,
nor of the will of man, but of God.} --- When, then, someone stands up and says:
``show us Christ as he is, and it is enough for us!'' we will remind him of the
Evangelist's words, and ask: have you also the eye with which Christ wills to be
seen, and the ear with which he wills to be heard? If not, then your demand is in
vain.

\emph{Show us Christ as he is!} This demand the Gospel satisfies, in that it
presents the Redeemer to us in relation both to the believers and to the
unbelievers among his own, and thus explains to us the two-sided nature and
essence of this fundamental relation. We convince ourselves of this at once by
considering the first scenes in the Life of Jesus, and seeing how the Forerunner
prepared the way for his Lord.

Gospel, Joh.\ 1, 35---52.

\begingroup\large
The next day after John stood again, and two of his disciples. And looking upon
Jesus as he walked, he said: behold the Lamb of God! And the two disciples heard
him speak, and they followed Jesus.
\par\endgroup

On the occasion of this testimony, ``scientific'' criticism unfortunately has a
multitude of difficulties to contend with. --- In the first place it is not easy
to find out where the Baptist actually got the idea that Jesus should be the Lamb
of God that (V.\ 29) bears the sin of the world. ``What kind of lamb is being
alluded to here? The pair of lambs that (2 Mosb.\ 29, 38) were offered morning
and evening one naturally cannot think of, since these offerings were not
offerings of atonement. It must therefore be another lamb; but which? Next, it is
difficult to see what the Baptist can actually have had in mind with the
expression: to bear the sin of the
% --- p. 211 ---
\opage{211} world, since the word `to bear' can naturally be taken in more than
one sense.'' This is what the scientific interpreters nowadays call --- eine
Schwierigkeit. But suppose that the art of interpretation even got this
``Schwierigkeit'' removed (which, however, can hardly be hoped for so soon); then
there still remains here another, and a far greater one\footnote{%
  Those who have not yet properly acquainted themselves with German theology's very
  newest critical editions of ``Leben Jesu'' can hardly form any conception of what
  ``Schwierigkeiten'' the modern consciousness has already discovered in the sacred
  history. As an example we cite here ``eine große Schwierigkeit'', which precisely
  with regard to the Baptist's testimony about Jesus (Joh.\ 1, 29---34) has been found
  by Bruno Bauer, a discovery about which Ebrard in turn has expressed himself with
  much indignation: Er (the Baptist) soll alle fünf Verse,
  im Ganzen 102 Wörter, darunter 22 Artikel und einsylbige Präpositionen, und 7
  \textit{enclyticæ} und \textit{atona}, und 2 \textit{καί}, also im aramäischen nur 70
  Wörter, d.\ i.; 11 Oktavzeilen --- er soll sie alle gesprochen haben! --- Dennoch sagt
  Bauer (Pag.\ 22): ``Der Herr muß also noch weit entfernt gewesen seyn'', (damit diese 70
  Wörter gesprochen werden konnten). ``War das der Fall, so ist das Reden und Hinzeigen
  haltlos und sieht gezwungen und ungeschicht aus. War aber der Herr so nahe, daß der
  Täufer beqvem auf ihn zeigen konnte, so bleibt für die vielen Worte keine Zeit, er
  müßte sie denn den Umstehenden auf das schleunigste und hastigste in das Ohr gezischelt
  hoben.'' --- Das ist Bubengeschwatz. J.\ H.\ A.\ Ebrard. Wissenschaftliche Kritik der
  evangelischen Geschichte. Frankfurt a.\ M. 1842. 1ster Theil. Pag.\ 309.%
}. --- ``If John the Baptist really recognized Jesus as the Messiah, and referred
his disciples to him: how then did it come about that he did not at once
% --- p. 212 ---
\opage{212} give up his own imperfect activity, and have himself received with all
his followers into Jesus' company?''

Yes, asks the Believer, how may it come about that the scribes will no longer
understand\footnote{%
  Naturally this cannot hold of all scribes. Cf.\ \textit{Dr.} A.\
  Neander. Das Leben Jesu Christi. Hamburg 1837. Pag.\ 73---92.%
} that the Forerunner best serves his Lord by performing his errand as Forerunner?

\begingroup\large
But Jesus turned, and saw them following, and said unto them: What seek ye? And
they said unto him: Rabbi (which is to say, being interpreted, Master), where
abidest thou? He said unto them: come and see. They came and saw where he abode,
and abode with him that day; for it was about the tenth hour.
\par\endgroup

The first meeting, the first glance, the first word, the first impression: how
inestimable is not all this that is first in the memory of a precious
personality! The disciple who loved the Lord, and whom the Lord loved, recalls
here the moment when he first saw his Savior and Friend face to face. It was an
unforgettable moment. He himself and another disciple (it was Andrew, Simon
Peter's brother) had just heard John's testimony; they then both began in all
timidity to follow after the Lord, but did not dare to speak to him or to stop
him on his way. Then the Master turned round to them, spoke to them, invited them
to himself, and let them stay with him \emph{that same day.} It was, then, the
first time they saw this countenance, as he turned round; it was the first time
they heard this voice: \emph{what seek ye?} It was his first invitation:
\emph{come and see!} It was a great, an unforgettable moment: \emph{it was about
the
% --- p. 213 ---
\opage{213} tenth hour.} When John wrote his Gospel, he had, according to the
testimony of the Church, become an old man, an apostolic elder. He had by then
seen much, and heard much, lived through much, and no doubt also forgotten much of
what he had heard and seen in the world; but this first meeting with the Lord and
Master, this first invitation, this first --- blessed moment of expectation,
fulfillment and joy he could never forget: \emph{it was about the tenth hour.}

A similar impression this first meeting must surely also have made on the other
disciple. For Andrew, after all, hastens to seek out his brother and proclaim to
him the glad tidings.

\begingroup\large
And he brought him to Jesus. And when Jesus beheld him, he said: thou art Simon
the son of Jona, thou shalt be called Cephas (which is by interpretation Peter).
\par\endgroup

``But'' --- the modern consciousness continues --- But no! we will not listen to
the modern consciousness now, but rejoice with the first disciples, and believe
the Gospel, and hasten like Philip to find Nathanael.

\begingroup\large
Philip found Nathanael and said unto him: we have found him, of whom Moses in the
Law, and the Prophets, did write, Jesus, Joseph's son, of Nazareth.
\par\endgroup

``Among the names of Jesus' disciples, Nathanael's certainly does not belong to
the world-famous; it does not shine in the history of the Church beside the names
of such disciples as Peter, Paul and John; we know but very little of the
fortunes of his life, and as good as nothing of what he wrought and accomplished
in the Lord's service; but to the name of Nathanael there nonetheless attaches the
idea of the most lovable and precious thing that can be said of a human being,
the idea of an upright soul whom the heavenly Father himself
% --- p. 214 ---
\opage{214} had stamped with a holy simplicity and an unfeigned love of
truth.''\footnote{%
  \textit{Dr.} H.\ Martensen's Sermons. Second Collection. p.\ 5.%
}

\begingroup\large
And Nathanael said unto him: can there any good thing come out of Nazareth? Philip
said unto him: come and see!
\par\endgroup

On Nathanael, then, this message makes a surprising impression. --- ``Joseph's
son of Nazareth! Can despised Nazareth then come to such honor and dignity?'' ---
The honest man begins as though he would doubt Philip's report and open an
independent investigation; but --- here there is no time for evasive
deliberations: \emph{come and see!} And Nathanael comes in the tension of
expectation, as one who wants to put the matter to the test. Then he hears the
Lord's voice: {\large behold, this is truly an Israelite, in whom there is no
guile;} and now he gets --- other thoughts. --- {\large Whence knowest thou me?
--- Before that Philip called thee, when thou wast under the fig tree, I saw
thee.} --- No more is needed.

But is this trusting Israelite not, after all, just like a child? At first he
wanted to be an independent thinker, and consider how far the promises in the Law
and the Prophets could be made to refer to Nazareth. Now, when he stands before
Joseph's son of Nazareth and hears his first words, he is at once convinced, and
why? Because Joseph's son of Nazareth has recognized him, Nathanael, to be an
Israelite without guile, because this son of Joseph of
Nazareth has seen him, Nathanael, under the fig tree before Philip called him:
see, that is why the honest Israelite is now fully convinced that Joseph's son of
Nazareth must be the Son of God, the King of Israel. --- Such is the first
impression! Let us not dwell critically on the fact that the independent thinker
% --- p. 215 ---
\opage{215} Nathanael is so suddenly transformed into pure trustfulness; let us
not reproach the honest Israelite for losing his equilibrium in the abrupt
transition. The abrupt transition is, after all, precisely faith's own immediate
beginning. But it ought nonetheless to be known that it is --- a beginning.
Therefore the Lord not only gives Nathanael a gentle and loving admonition (V.\
51), that the honest Israelite should not be overhasty; but at the same time he
builds on the first impression, he adds the promise, as he begins the journey
with the expectant disciples:

\begingroup\large
Verily, verily, I say unto you: hereafter ye shall see heaven open, and the angels
of God ascending and descending upon the Son of man.
\par\endgroup

``But'' --- the modern consciousness continues --- ``among a few people not of age
a gifted man can easily play the superior one. A hint, an obscure allusion, a
friendly invitation is enough. Indeed, if men in general were as credulous as
these disciples, it would not even be so very hazardous an undertaking to give
oneself out for a Messiah. It is perfectly clear from Philip's words that this
limited person had not the slightest notion of what criticism is. Had Philip
known only a little bit about criticism, he would surely have thought it over a
while before he took Jesus, Joseph's son of Nazareth, for the one `of whom Moses
in the Law and the Prophets did write'. Of course, one can hardly reproach Philip
for not knowing criticism: he was, after all, a Galilean man, a man of the
people, a man without higher education. Besides, in those days the critical
understanding had not yet properly become conscious of itself, least of all among
the Jews.''

It is possible, after all, it may well be, that these Galilean men came to faith
so easily because they lacked natural understanding and
% --- p. 216 ---
\opage{216} sound judgment. Yes, it is possible, after all! Meanwhile, before we
sign so harsh a verdict, we will nonetheless inquire a little more closely of the
Evangelists whether by their guidance we might perhaps find out how things
actually stood with criticism in Galilee. We may, for example, hear the Evangelist
Luke (4, 14---30).

\begingroup\large
And Jesus returned in the power of the Spirit into Galilee, and there went out a
fame of him through all the region round about. And he taught in their
synagogues, and was glorified of all. And he came to Nazareth, where he had been
brought up, and went into the synagogue, as his custom was, on the sabbath day,
and stood up to read.
\par\endgroup

The Galileans in Nazareth, then, do not attempt to put any obstacle in Jesus'
way. They are no doubt even rather glad to see that the young teacher of the
people, a man of their own town, a gifted man, who had already begun to win a
name, now also stands up to read in their synagogue.

Then they gave him the book of the prophet Esaias, and --- their expectation is
not disappointed, for the reader knows at once how to find his place, and --- has
a clear voice, and reads to them with force and power.

\begingroup\large
The Spirit of the Lord is upon me, therefore he anointed me; he hath sent me to
proclaim the gospel to the poor, to heal them that have a broken heart, to preach
to the captives that they shall be set free, and that the blind shall receive
sight, to set at liberty them that are bruised; to proclaim a year of grace from
the Lord.
\par\endgroup

The Galileans in Nazareth are not so uneducated after all: all the hearers sink
into attention, there is a listening silence over the whole assembly.

\begingroup\large
And he closed the book, and gave it to the attendant
% --- p. 217 ---
\opage{217} again, and sat down; and the eyes of all in the synagogue were fastened on him.
\par\endgroup

The Galileans in Nazareth are by no means unappreciative of the gifts of the Spirit: they
gladly give themselves up in the hour of devotion; their minds can be lifted up by a godly
discourse.

\begingroup\large
But he began to say unto them: This day is this scripture fulfilled in your ears.
\par\endgroup

And the Galileans listen, and hear his speech, and are seized by the Spirit. {\large And they
all gave him praise, and wondered at the gracious words which proceeded out of his mouth . . . .}

--- So the Galileans in Nazareth are, after all, very near to grasping the fulfillment of the
prophecy, very near to forgetting all irrelevant considerations, very near to giving
themselves up to the gracious words. Then a disturbance arises in the assembly, the silence is
broken, the attention disturbed: what kind of being is this, then, that all at once forces
its way into this devout synagogue and causes such confusion? --- Why, it is --- criticism!

{\large Is not this Joseph's son?} --- See, here lies the problem, the critical problem, in
its first natural nakedness. The assembly in the synagogue was not content to hear the words
of Jesus; but --- the eyes of all were fastened on him. (This straining of outward sight has
no doubt contributed to weakening the inward sense.) If only these hearers had been willing to
lower their eyes a little, and in addition to open the inner ear still more: perhaps they
would have appropriated the gracious words, perhaps they would have known him by his voice.
But since they did not cease to stare at him, neither could they cease to ask themselves:
\emph{is not this Joseph's son?} As it is expressly said in Matthew (13, 54---55):
\emph{Whence hath this man such wisdom and the mighty
% --- p. 218 ---
\opage{218} works?} --- Now if one cannot be content to appropriate the wisdom and believe the
mighty works, but must, however it may go, absolutely have it fathomed whence he has this
wisdom and how he has become able to do such works, why, then there is no end to the
objections of offense. --- Modern science prides itself on having discovered the truth ``that
the absolute Idea cannot concentrate itself in a single limited individual.'' That is a vain
conceit in which modern wisdom is pleased to please itself. The discovery, essentially
understood, is not new: it was already made in the synagogue in Nazareth. The faith of the
Gospel and negative criticism have the same fatherland: both are descended from Galilee.
``But'' --- continues the modern consciousness --- ``criticism is now developed quite
differently, and has quite another power than back then, when it was in its first Galilean
childhood.'' This claim too is only half true, and rests moreover on a vain conceit.

When modern wisdom begins in its language: ``So absolute, infinite an Idea, and yet nothing
but this finite individual? An impossibility!'' --- then this reads, in good Galilean:
``Such wisdom, such gracious words, and yet nothing but Joseph's son? --- An absurdity!''
--- So the men of Nazareth were then just as eager to explain away this absurdity as the
reasoning critics are to this very day. The difference is only that the men of the synagogue
could not take the task quite so lightly as the talents of the newer school. Modern criticism
can, when it comes to the pinch, dissolve the historical Christ pretty much as it pleases;
the Galilean criticism, on the other hand, did not have that advantage. The historical Christ
was himself present in the synagogue, the eyes of all were fastened on him: just listen how
they ask and ask,
% --- p. 219 ---
\opage{219} %
in order to get over the real stumbling-block! ``Is not this the carpenter's son?'' If only
they could have got free of this presupposition, it would have helped a good deal. But that
they could not do. They knew his history far too well for that; they could go into the finest
details, why, they had his whole genealogy at their fingertips: \emph{Is not his mother called
Mary? and his brethren James, and Joses, and Simon, and Judas? And are not all his sisters
with us? Whence then hath this man all these things? And they were offended in him}
(Matth.\ 13, 55---57). --- Modern criticism can, when it finally comes to the pinch, deny not
merely the idea of a God-man but the absolute self-validity of the Idea altogether. This
evasion, to be sure, the Galilean criticism did not know either. The men of the synagogue
still stood in a certain way under the influence of the Spirit; it struck them when they
heard him begin with the prophet: \emph{The Spirit of the Lord is upon me,\ldots\ he hath
sent me to proclaim the Gospel to the poor\ldots} they had to admit themselves that he spoke
gracious words. --- ``But'' --- one continues --- ``then this is precisely a proof that
criticism at that time lacked the requisite power.'' --- You are mistaken, you learned
gentlemen; it is precisely a proof that the Galilean criticism had fully as great a power as
the modern. For although those Galileans in Nazareth could not (what is naturally an easy
matter for the modern consciousness) avoid the first immediate impression of the Savior's
personality, their criticizing understanding nevertheless had strength enough to strike faith
down. More, surely, one will not demand. When therefore \emph{the Lord's Anointed} would not
comply with their defiant demand that he do a sign, but began to speak to them of ``the
widows in Israel in the days of Elias,'' and of ``the many lepers in Israel in the time of
Eliseus the prophet'':
% --- p. 220 ---

\opage{220} \begingroup\large
Then all they that were in the synagogue, when they heard these things, were filled with
wrath. And they rose up, and thrust him out of the city, and led him unto the brow of the hill
whereon their city was built, that they might cast him down. But he went through the midst of
them, and departed.
\par\endgroup

Criticism in Galilee, then, did not after all stand on so low a level as modern science would
now so gladly have us believe.

``Who then was in the right: those Galileans who became his first disciples, and believed him
at his word that they should now see heaven open and the angels of God ascending and
descending upon the Son of man? or these Galileans who in the synagogue attacked him as a
false prophet, and thrust him out of the city in order to cast him down from the brow of the
hill: which of the parties was then in the right?'' --- why, that the science of the present
had better refrain from deciding. The matter was left undecided at the time. Christ himself,
in his own person, could not convince the men of the synagogue of his messianic calling; how
then should any teacher in the present be able to convince the non-believers that the Jesus of
the Evangelists is the real Christ. But does not the converse now hold as well? Are not all
the tumults of unbelief in the synagogues of criticism still too powerless to rob him of a
single one of the disciples who are his own in faith?

\begingroup\large
I know them, and they follow me. And I give unto them eternal life, and they shall by no means
perish, and no one shall pluck them out of my hand (Joh.\ 10, 27---28).
\par\endgroup

It is of course possible that this word might one day fail; it is of course possible that
this teacher, critically speaking, misled his first disciples, and that these misled ones then
in turn misled the many others in the course of time; we make the modern consciousness this
accommodating concession: \emph{it is of course possible!} But now the modern
% --- p. 221 ---
\opage{221} %
consciousness ought in return to show us the goodwill to admit that it has nevertheless not
been quite so easy to beguile the Galileans, and that on the whole things hang together
strangely with Jesus' first appearance. --- Should his messianic dignity really be found to be
grounded in a lamentable mistake, then it would undeniably have been of great benefit to
mankind if the strong men of the synagogue had succeeded in casting him down from their
hill-brow and at the same time freeing the world from so seductive a deception. But what shall
we say now? If it did not succeed then, when he was after all so near the brow, how should it
succeed in our days. If he went unharmed through the midst of the incensed crowd that wanted
to hold him fast in the bodily sense, then there is surely very little hope for the strong men
who now, from the various synagogues of criticism, exert themselves to thrust him out of the
evangelical representation in order to cast him into --- pure consciousness. If he escaped
those, who were after all so near to seizing him with their hands, why then should he not also
escape these, who can only seize him in thought and set him down in the concept. --- ``Often
enough we hear from the critical synagogues that this or that strong man has now at last
succeeded in leading the Nazarene to the brow of the hill and casting him down. But --- that
is a pure mistake! There was perhaps something that fell down. But this something is by no
means the Christ of the Gospel; no, it was only an old doctrinal edifice that fell, because it
could no longer stand. At times, to be sure, it seems as if the light of the Gospel lost its
shine, as if the salt had lost its savor and its strength, and the word died away on our lips.
In such interim times the negative spirit exults, and the cry goes up in the synagogues of
criticism: `now Christ is dead; it is we who have now at last got him killed and laid under
guard in the tomb.' A pure mistake! There was perhaps
% --- p. 222 ---
\opage{222} %
something that went into the tomb. But this something is not the Word of life which was from
the beginning; no, it was faith's human word, which like a Lazarus was laid in the tomb, and
waited for Him who was to come again to wake it to life.'' See, therefore Christ is and
remains an eternal riddle for pure self-consciousness. The science that otherwise comprehends
everything cannot possibly comprehend how it happens with this Gospel of poverty that it is so
often crucified and buried, and yet again and again rises as on the third day, and comes to us
in the hour of sorrow, and comforts every troubled soul, as the Savior comforted the deeply
grieved sister: \emph{said I not unto thee, that if thou wouldest believe, thou shouldest see
the glory of God.}

``Show us Christ as he is, and it sufficeth us!'' --- Only the spirit of the Gospel can
satisfy this demand. There is no scientific proof that Jesus of Nazareth is the only begotten
Son of God; and there is no scientific proof that he is not the Son of God.

``But'' --- so comes prudent circumspection --- ``if it must now be admitted that this matter
cannot be decided by irrefutable grounds, then it is, as we remarked above, still after all
the best course to leave the whole thing undecided, and keep plainly to custom and usage. Our
Lord surely cannot demand that we human beings should now all at once start up and want to
decide what cannot possibly be decided as long as this world stands.''

A decision can and shall take place! When Christ came to his own, he brought with him all the
conditions for the possibility of the decision. But the decision itself was not to take place
outside man; it took place within man; and so it takes place still to this very day. Whether
Philip was in the right when he testified to
% --- p. 223 ---
\opage{223} Nathanael that Joseph's son must be the Messiah, or the men in the synagogue were
in the right that this son of Joseph ought to be cast down from the hill-brow, that was
decided at the time not by that which is outside man (for all saw and heard the same man), but
by that which is within man, and so it can be decided in this now.

``But'' --- continues prudent circumspection --- ``one can very well have a certain confidence
in Christ, and go to church and honor religion, without therefore needing to squeeze oneself
into a painful `Either --- Or' and put thumbscrews on faith. Such decisions are best left to
those who fish, not like St. Peter to catch men, but, as clever heads, to catch ideas and
surprise with the novelty of their presentation. While these decisive authors make an end of
everything and turn all things upside down, they set especially great store by the dexterity
with which they know how to turn Christ's coat. How many times and in how many ways this holy
coat has now been resewn and turned in the last 40 years is well enough known. As soon as a
critic or philosopher has got the coat turned, or seen to a little resewing, the word goes out
at once in the literature: `a decision, a quite new decision!' This decision is naturally
determined solely and entirely by the side of the coat that is turned out. If it is, e.g., the
philosophical-speculative side that is turned out, then Christ appears as the personified
Concept: `He that comprehendeth the Concept shall be saved; he that comprehendeth not is
condemned already' (which is why a Hegelian in his day had also to raise the question whether
women and immediate poets could really be saved, since they were not organized to comprehend
the Concept). If speculation now becomes threadbare (and that happens as a rule as soon as it
is taken along for everyday use), the philosophical
% --- p. 224 ---
\opage{224} side must again be regarded as the wrong side, and it then becomes once more a
necessity to get the coat turned. Meanwhile one cannot without further ado turn the old side
of faith out again; for it too is, frankly speaking, somewhat threadbare. Here, naturally, a
change must be made. What does one do then? One accentuates incomprehensibility to the most
incomprehensible degree; and only then does one turn the right side, that is to say, the side
of faith, out again. The old, then? No, by no means! Something entirely new, or at any rate as
good as new. That first, plain immediacy, which we know well enough, is now worked up into a
quite peculiar, piquantly surprising immediacy, a, so to speak, decatized immediacy, which
gleams and sparkles when one strokes it in the direction of --- \emph{the Paradox}. As it goes
with the comprehensible and the incomprehensible, so too with Christ's mildness and severity.
If he steps forth in the philosophical coat of mildness, he is so calm, so contemplative, so
ethically indifferent, that one need only rest in his metaphysical embrace in order at once to
perceive the eternal harmony of existence, the blessed unity that pulses in all the manifold.
But if the conception now changes, if the coat has got the mark of the Paradox: then woe unto
the contemplative and to them that suck at speculation in those days, for then shall
tribulations and anguish and terrors come upon men. Now he is sheer will, sheer fire and sheer
spirit: the scourge of zeal is in his hand, and even if he does not drive the whole
congregation out of the sanctuary, his holy love is still on the point of bursting for us the
noblest bond of humanity. --- See, this alternation of extremes undeniably gives material for
very instructive reflections; but a well-organized head will take good care not to stake too
much on one extreme. One knows how it goes: it lasts only a little while, then the extreme so
commended passes over into its opposite, then one regrets what one has done, and has
% --- p. 225 ---
\opage{225} besides the shame of losing one's composure. Altogether, extremes are a quite
peculiar matter: one dare not condemn them; on the contrary, when one speaks in pure
generality, one ought even to acknowledge warmly their necessity for the maintenance of forces.
But as regards one's own person, let one take good care not to go to any extreme: that honor
is always best left to others. While the opposite extremes now provoke each other and level
each other, the doctrine of the Middle Way stands behind in the still luster of a silent
superiority, draws nourishment from all extremes, and then at last comes forward with the
final, well-considered, clarifying and reassuring word.

``This now holds of all extremes, but especially of the Christian ones. Christ did not come
into the world to split men up into extremes (this splitting is unfortunately only too well
grounded in our own natural passions); but he came into the world to gather what was
scattered, and thereby to breathe a new life into the world. Yes, he is the bread of life, not
for some few, but for all: he nourishes our natural life too, he creates bread in the
wilderness. --- Without, naturally, being able to vouch for the miracle of the feeding as a
wonder-work, or to lay any special weight on this admittedly very disputable sign, we shall
merely use this hint to draw attention to the universally human in Christianity: men receive
bread from heaven so that they may then learn to work the work of God on earth.''

\emph{Show us Christ as he is!} --- We have now heard the wisdom of circumspection; let us then
turn again to the Gospel.

The decision of self-wise circumspection in mediocrity was not unknown in Galilee either.
There was a multitude which, after having enjoyed the wondrous feeding (Joh.\ 6 Kap.),
% --- p. 226 ---
\opage{226} crossed over the Sea of Tiberias to seek Jesus in Capernaum. This multitude, which
likewise did not want to lay any special weight on the single sign, but all the more emphasis
on the edible bread, had nevertheless conceived a certain confidence in Jesus, wanted also to
honor religion to a certain degree, and therefore began to inquire what it might mean to work
the work of God. So long as the Savior's answer kept to generalities, did not demand too much
of the hearers, or go to any striking extreme, this multitude was, as it seemed, quite well
enough satisfied with the discourse. But when the testimony concerning \emph{the bread of life}
now became more forceful and more severe; when the Lord touched on the Paradox, and repeated:
\emph{my flesh is meat indeed, and my blood is drink indeed}; what did the multitude do then?
It became extremely dissatisfied, and murmured in the synagogue at what he said. They found
that it was a hard saying; they found that he went to an extreme; they found that it was
intolerable, and then --- they went their way.

\begingroup\large
From that time many of his disciples went back, and walked no more with him
(Joh.\ 6, 66).
\par\endgroup

See, that was the multitude's decision.

\begingroup\large
Then said Jesus unto the twelve: Will ye also go away.
\par\endgroup

Christ would gladly draw all to himself; but he will not tone down his demand; he will not at
all employ any ingratiating persuasion of accommodation in order to keep the many hearers and
the many disciples.

\begingroup\large
Then Simon Peter answered him: Lord, to whom shall we go? thou hast the words of eternal
life.
\par\endgroup

See, that was Peter's decision.

Who then made the best choice: the multitude that went away,
% --- p. 227 ---
\opage{227} and left to Peter ``the honor'' of following his Lord to the uttermost, or Peter,
who, without heeding the offense of the multitude, received the words of eternal life, even if
he had not yet grasped the secret of the discourse; which of the parties was in the right: the
self-wise multitude or the believing disciple; that was not decided by what is outside man
(for all saw and heard the same man); but it was decided by what is within man; and so it is
decided still to this very day.

``Show us Christ as he is!'' --- \emph{He came unto his own, and his own received him not. But
as many as received him, to them gave he power to become the children of God, even to them
that believe on his name: which were born, not of blood, nor of the will of the flesh, nor of
the will of man, but of God.}

``Show us Christ as he is!'' --- It is a great responsibility that was laid upon mankind when
Christ came into the world, as unto his own. Therefore every human being in the world shall
give account of how it has seen him, whether it has been willing to receive him and believe on
his name.

``Show us Christ as he is, and it sufficeth us!''

Will you then receive him, and believe on his name? --- The promise is comforting, but ---
the threat terrible.

\begingroup\large
Whosoever will confess me before men, him will I also confess before my Father which is in
heaven; but whosoever shall be ashamed of me and of my words in this perverse and sinful
generation, of him also shall the Son of man be ashamed, when he cometh in the glory of his
Father with the holy angels (Matth.\ 10, 32. Mrk.\ 8, 38).
\par\endgroup

% --- p. 228 ---
\opage{228} %
\subsection*{\textit{III.}\\[4pt] Christ Is a Sign That Is Spoken Against}
\addcontentsline{toc}{subsection}{III. Christ Is a Sign That Is Spoken Against}

If we are to point out in one word the main point in the evangelical fortification against
which recent science aims and fires from all sides, as against the weakest and yet most
important point, we need only utter the word: \emph{contradiction}. ``Everything that exists,''
it is said, ``is a contradiction. The positive religions make no exception:
Alles was entsteht, ist werth, daß es zu Grunde geht.'' Now comes the explanation. ``The idea
of revelation has its root in the race. From the innermost ground of human nature this seed
shot up at the beginning of history in a mighty shoot. The first trunk split into various
branches; each branch, through the lush vegetative power of the primeval world, became in turn
an independent trunk with smaller but leafy branches, with peculiar flowers and fruits.
Afterwards the power weakened, and the vegetation declined. The side trunks rotted, the
branches dried up, the foliage withered, the fruits fell to the ground without sap or kernel.
--- This image has no doubt often been used; but let us carry it out a little further. That
the side trunks, which spent themselves in wild shoots, fabulous ramifications and fantastic
splendor of color, gradually withered away and fell off, was naturally no harm to the main
trunk. On the contrary: the actual principal trunk thereby gained precisely a heightened
vitality, since it could thus draw all the nourishment of the vegetation to itself alone. All
at once there then took place a transformation in its fluids and saps; it was as if a sigh
went through its marrow; it shot a new shoot; it was the true root-shoot. Small, insignificant
and lowly in its beginning, this root-shoot gradually transformed itself into the trunk
itself, and
grew then into a great tree, so that
% --- p. 229 ---
\opage{229} the birds of heaven came to build nests in its branches. --- If one now asks which
religions are designated in the image by the wild shoots, everyone knows that it is the
heathen ones. If one should further ask how it came about that those side trunks withered so
early, the answer would surely not be so difficult either. As the heathen religions lost the
luxuriance of vegetation with which they had originally
grown up, and had flourished in fantasy and feeling, it was a matter of course that they could
no longer give souls any healthy, satisfying nourishment. The peoples felt a lack, an
emptiness, but did not themselves know what ailed them. Then came
science with
criticism, and criticism discovered
at once where the fault lay: \emph{the whole thing was namely a contradiction}.''

``It is a profound feature of the Nordic myth of the ash Yggdrasil that four harts stood and
nibbled its leaves, that it became rotten with time, that all the serpents of Helheim gnawed at
its root. What are these serpents of Helheim but the manifold gnawing, devouring contradictions
of time and existence! Heathendom shed its bloom, and fell by contradiction. With Judaism it
went likewise. This trunk too withered and laid bare its contradictions when the young,
life-lush Christianity had sucked out its marrow and strength. These general truths are now
granted by most. `But,' it is said, `with Christianity it is a quite different matter. Here is
a trunk that never rots; this tree of life can never wither.' --- Without wishing to dispute
the superiority of Christianity over the other religions, we shall merely point out one
fundamental condition which must hold just as fully for Christianity as for any other religion
whatsoever. Every religion must grow a natural growth, and express itself in the race as a
higher, irresistible drive of life. (One may
% --- p. 230 ---
\opage{230} readily substitute for `the natural' the concept `the supernatural'; this change of
concept makes no difference here, provided only one will grant that even the supernatural must
be nature in its kind). Thus: a religion lives a healthy and true life in a people only so
long as it is able to preserve its spiritual vegetative power, and to work in the form of
\textit{natura naturans}. --- If one denies this fundamental law, one not only falls out with
psychology but directly contradicts Jesus' own teaching. --- It says in the parable
(Mr.\ 4, 26---28):
\emph{So is the kingdom of God, as if a man should cast seed into the ground, and should sleep,
and rise night and day; and the seed groweth and becometh high, he knoweth not himself how. For
the earth bringeth forth fruit of herself; first the blade, after that the ear, after that the
full corn in the ear.} What a difference between Christianity in its first centuries and
Christianity in the present. One might almost be tempted to believe that original Christianity
had been exchanged on the long journeys, or that it was not the same thing that is designated
here by the same name. In those days Christianity was so stirring and living, so lush and
alluring, that no noble human being and no people could resist it. Heathen priests hated it as
a poisonous weed, mighty statesmen wanted to root it out as a dangerous parasitic plant; it was
no use; it
grew almost wild; it shot up out of the ground of the heart as grass and herbs shoot up out of
the earth. In our days it goes otherwise. Now one would gladly have Christianity thrive and
grow everywhere; but there is, as it seems, no real vigor of growth in it. Wise
statesmen plant, pious priests water the noble
growth as best they can, and still it will not thrive. Is it mankind's fault? Let us not be
unfair. Experience certainly
% --- p. 231 ---
\opage{231} shows that we human beings nowadays are not as good as we could and should be; but
neither, after all, were the people of those times. It is no use at all trying to put the
blame on the race; the one generation has by nature no advantage over the other: Christianity
itself teaches, moreover, that we are all sinners. No one who has any overview of the turning
points of history will be able to deny that there is a disquieting likeness between the
Christianity of the present and heathendom in its last days. Now too the peoples feel an inner
emptiness, a lack, a spiritual dissatisfaction they cannot properly explain to themselves. Now
too scientific theology has transformed itself into criticism; now too criticism has
discovered where the damage is: \emph{the serpents of contradiction have long been gnawing at
the root of Christianity}.''

``Existence is in much pain when the holy dissolves in contradiction; but --- it is no use
getting angry at science. The priests of faith who curse criticism because it lays bare
contradiction act like those tyrants who would murder their own scouts when these bring them
unwelcome tidings. What childish folly! As if it were science that cast contradiction into
Christianity; whereas, conversely, it is Christianity itself that bears contradiction within
it, and precisely for that reason suffers from an incurable disease. If we examine the
doctrine of the Church: a system of contradictions! every dogma is a peculiar contradiction.
If we compare the evangelical accounts: contradiction follows upon contradiction. If we
consider the historical Christ: an existing contradiction! If we look into the inner life of
the Believer: O sorrowful contradiction! --- In vain will the Protestant Church defend itself
against criticism, and veil the contradiction. At the Reformation criticizing thought gained
free access to the mysteries of revelation. To be sure, science at
% --- p. 232 ---
\opage{232} that time played a different role from now. At the time of the Reformation the
critical process was like a vinous fermentation; in our time it has become like an acetous
fermentation; but the one follows from the other.''

According to the view set out here, it might now almost seem as if the faith of the Gospel
were incompatible with true criticism, as if we had only the choice between two possibilities:
either to renounce all stricter thinking and believe blindly, or to give up the faith of the
Gospel in order to serve science.

Thus the matter is presented in modern literature, and from there the rumor spreads ever
further in educated circles that
\emph{``negative criticism has now proved to be the last and highest standpoint of
cognition.''} But this rumor is nevertheless an untruth, and this untruth would soon come to
public knowledge, if only the people of the present could get a little more time to think over
such things for themselves.

Instead of contradicting the modern theory of the principle of contradiction, we shall, for a
change, begin by granting its truth: greater compliance ``science'' can surely not demand.
Thus: the whole of existence lies in contradiction. Each time this or that finite thing
perishes, a contradiction is resolved; but it is only a finite contradiction. Each time this or
that finite thing arises, a new contradiction forms, but only a finite contradiction.
Countless contradictions can in this way arise and vanish; but contradiction itself, the inner
moving unrest of existence, the rhythmic pulse of life that beats in all nerves, can never
cease. Every human life bears contradiction within it. Nevertheless an ordinary human being can
very well die away without existence suffering from it, precisely because the contradiction in
the single individual is so weak and slight. But if we now put in place of the everyday human
being an extraordinary
% --- p. 233 ---
\opage{233} world-genius, a hero: what a change! The hero too bears his contradiction within
him, and must at last succumb to the general law of existence. But when the turning point
arrives, when he stands at the boundary and is about to fall, then there is indeed a fire in
the world and an uproar in men, as if the whole of existence were to go to
ruin with this one man.

The contradiction in the heroic genius penetrates so deeply into the fundamental contradiction
of existence itself that the hero's fall is on the point of setting the crown on the
corruptibility of life and making history meaningless. Meanwhile this shock too passes, like
another earthquake. Existence orders itself again. A criticism befalls the hero; this
criticism now proves that the contradiction was more accidental than necessary. ``The great
genius could, after all, have done the whole thing otherwise: posterity at least will know how
to do it quite otherwise.'' --- We go a step further. Instead of the single individual we
choose, e.g., the Idea itself in one of the positive religions.

Religion has to take up the finite into the infinite, the temporal into the eternal, to check
evil and promote good. In order to solve this task, religion must necessarily engage not
merely with these or those finite contradictions, but with the contradiction in all
contradictions, i.e., with the fundamental contradiction of existence itself. So long as the
positive religion masters contradiction, the pious people has a comfort in all its sorrow, a
refuge in all its need. It is almost incredible what a people can suffer and go through when
religion only continues to afford it the requisite influx of strength and courage. But it is
also almost incredible how a people can be weakened and confused when contradiction gains power
over religion: no finite means, no outward gain replaces the dwindling vital force. Seldom or
never is the dissolution of a folk religion accomplished without great
% --- p. 234 ---
\opage{234} anguish of soul in the people. But even if the dissolution took place almost
imperceptibly; even if the multitude in all its immaturity did not properly understand what
loss it suffered, but rejoiced in the novelty of the change, just as a child who has lost its
mother rejoices in the time of mourning over --- ``the new dress'': experience would
nevertheless soon make the people feel what it had lost; for with the fall of religion
contradiction breaks loose at every point of existence.

In order best to see how we now come to terms with our modern theory of contradiction, we set
up the following main theses:

1) \emph{Without religion the human race cannot live.}

2) \emph{Every religion is, according to its idea, designed to master the contradiction of
existence.}

3) \emph{The difference between the imperfect and the perfect religion is not that the former
takes up contradiction and the latter excludes it. But the difference consists in this, that
the imperfect religion, which conceives its contradiction superficially, also resolves it only
superficially, and must therefore at last succumb to contradiction; whereas the perfect
religion grasps the fundamental contradiction in its whole depth, and masters it with infinite
superiority.}

4) \emph{When the imperfect religion is dissolved and replaced by a more perfect one, this
transition is to be regarded as a progress to the better. But the perfect religion can neither
be dissolved nor replaced by anything else without the existence of the human race having at
the same time to prepare itself for its own dissolution.}

% --- p. 235 ---
\opage{235} Which of the positive religions is, for the rest, the perfect one, whether such a
one is already given us or we are still to expect it, about that we shall in this connection
pick no quarrel with modern science. On the other hand, with regard to the critical doctrine of
contradiction, we shall merely indicate a few distinguishing marks by which the perfect
religion is always to be conceived and judged. --- If every religion is to be regarded as a
link between the eternal and the temporal, then perfection must be determined as much from the
standpoint of eternity as from that of temporality. The imperfect religion has a side, a
moment of the eternal, but lacks the substance of eternity. It therefore holds only for a
certain lower stage of cultivation and for a certain time. When development advances, a
turning point comes at which the deficiency shows itself. The religion is then subjected to a
criticism, and dissolved into dialectical contradictions. This dissolution does not, however,
mean that the human spirit will give up the eternal (such a supposition is in any case only a
provisional fancy); but the critical dissolution means that mankind now feels the need for a
deeper and truer knowledge of God, that it wants to exchange the old fantastic pictorial forms
for purer and truer forms of eternity. Insofar as criticism carries through the dissolution,
the contradiction can rightly be said to fall on the religion's own imperfection; the
scientific standpoint is then, in this case, higher than the religious. But with the perfect
religion the matter stands the other way round. The content of the eternal is here expressed in
such forms and prefigurations that these cannot possibly be dissolved and disrupted without
mankind at the same time having to bid the eternal farewell and give up the spirit in despair.
If science nevertheless makes a critical attack on the religious forms of eternity, then it
makes, with or against its better knowledge, a desperate attempt to
% --- p. 236 ---
\opage{236} dissolve the eternal itself. In this case the contradiction does not fall on
religion, but it falls on science, and makes criticism insane. Let us then consider the member
of temporality in the religious relation. Since the imperfect religion mistakes the
determination of eternity, it at the same time misses the temporal moment. This happens in
various ways. On the one hand the temporal can be held fast in such a way that the divine is
made finite, and thereby at last perishes in finitude. On the other hand the temporal can be
volatilized into the eternal in such a way that actual existence is transformed into semblance
and shadow. Only the perfect religion knows how to grasp finitude at the point where it lets
itself be penetrated by the eternal without losing its own stamp of actuality. This fundamental
point is not to be sought in the general regions of nature or of history; but it is to be
sought in man. It is not the sun and not the planets, not the wild beasts of the field and not
the birds of heaven, that feel the contradiction of existence; but it is man. It is man as man
who is the wondrous middle term between the temporal and the eternal. It is man in his
determinate actuality, man as this single individual, divided between temporality and
eternity, who with full right can say of himself: I am the contradiction, it is I who must
suffer the pain of existence. --- If now religion alone can master the contradiction, and if
human existence is itself the contradiction, then in the strictest sense it is not enough
either that religion come to man's aid; but the perfect religion must itself become man, and
the perfect man exist as religion itself.

As the founder of the religion is, so too is the religion. At the stage of imperfection the
founder of a religion can be separated from the religion he founds. The contradiction in his
character and personal relations can be subjected to a criticism; can
% --- p. 237 ---
\opage{237} to that extent be condemned by science. The perfect founder of a religion, by
contrast, cannot be condemned by any criticism, since he is indeed religion itself, and the
absolute religion masters the deepest contradiction. By attacking the religion existing in the
person of the founder, criticism attacks the innermost kernel of human existence. For could it
be proved that the worm of contradiction gnawed at this kernel too, then it would at the same
time be proved that the fruits of life were all, without exception, worm-eaten. But a criticism
that casts everything, all aims and all endeavors, down into a bottomless contradiction must
surely be consistent enough to include itself: this attempt is thereby stamped as meaningless
and self-contradictory.

Having thus shed light on the relation of contradiction to religion in general, we shall now
find it easier to determine the relation of scientific criticism to the faith of the Gospel.
--- What is criticism? The thing itself must surely correspond to the name. The word criticism
is of Greek origin. To judge is in the Greek language \textit{κρινειν}
\textit{(krinein)}, a word that also means to distinguish; because no one is able to judge
without at the same time distinguishing and separating, joining the homogeneous and separating
the incompatible. So that our inquiry shall not take a learned turn, however, we shall at once
leave the verb and keep to the Greek noun, which, remarkably enough, is known to all, namely
the word: crisis. Everyone knows what it means when it is said: a crisis is impending. In all
spheres of life (even in the world of money) considerable crises do take place from time to
time. The turning point at which the crisis sets in is called the critical moment. In such a
moment the indeterminate must be determined, the obscure clear itself, the undecided be
decided. Now it is --- a matter of life and death; now it must either break or bear. The
critical moment is in this respect at the same time a moment of danger, a
% --- p. 238 ---
\opage{238} moment of anxiety and suspense. In every critical turning point there is always
something that must be ventured, something that is at stake. When many interests are set in
motion, the passions are aroused: the critical point in time is usually accompanied by strong
shocks. The crisis as a rule brings with it that something remains standing while other things
fall, because they have not strength enough to hold themselves up: the critical moment is to
that extent to be regarded as a moment of trial. The crisis is a purgatory that consumes all
that cannot stand its test. But sometimes the crisis takes place in such a way that the
heterogeneous constituents, even if they are clarified out of their ambiguous mixture,
nevertheless continue to hold themselves over against one another, now as indifferent, now as
hostile factors.

When a decisive crisis begins in mankind, the even intermediate state must at once vanish.
Where the Middle Way formerly went, there now goes the rushing stream that sweeps everything
along with it, and by new mixtures produces the fermentation from which the final separations
at last result. The movements in the opposite directions then sometimes take place with such
speed that people who formerly shared one fate suddenly see themselves separated as by a
yawning abyss. --- Memorable crises can thus take place in the merely temporal world. The
tension can be so great, and the swing to both sides so strong, that only the extreme points
stand out. In such critical moments it is: wealth or poverty, freedom or slavery, eminence or
lowliness, power and authority or oppression and wretchedness. And what passions, and what
interests, and what convulsive exertions can such crises not bring about in the human world,
although in all these cases the matter really concerns only temporal and finite goods.

One can indeed almost never grow tired of contemplating
% --- p. 239 ---
\opage{239} the heroes. It is an infinite satisfaction to see with what assurance the hero
appears, and then the irresistible urge that draws men to him. In hosts the many flock about
this one, and, it goes without saying, he does not send them away either. It looks almost as
if he had sent for them all; he greets them, these many thousands, with so heartfelt a
satisfaction and so inviting a superiority as if he had himself created them all --- for his
own use. It is precisely the most beautiful side of spiritual superiority that it not only
presses down but also lifts up the multitude of men. The practical genius not only raises
himself above all his contemporaries, but he also raises all these contemporaries a
considerable step up above themselves. This holds just as fully for those who are against him
as for those who are with him. The former he teaches to hate, the latter to love, with a
passion hitherto unknown; in the former he awakens a power of resistance they would otherwise
never have known, in the latter a self-sacrifice that must surprise themselves. The hero's
greatness is known by the admiration and adoration of his followers; but it is known perhaps
still better by the hatred and curses of his enemies. The practical genius has the attracting
and the repelling power in equal degree; when the heroic character shows itself, then it is in
truth a critical moment; the hero produces crisis.

Yet this crisis is not the highest. Let the world-idea be as great as it will, the idea of
blessedness is greater. Let a man sacrifice everything, his life, his soul, for a worldly idea,
he still does not bring the greatest sacrifice: his spirit is intoxicated, his soul stupefied,
he does not rightly know what life is worth. But he who sacrifices everything to win, not
temporal honor, but eternal blessedness; he who fully knows what it is that costs him his life,
and what life costs: he brings the greatest sacrifice. The relation in which the
% --- p. 240 ---
\opage{240} subordinate feels himself consecrated, on behalf of the idea, to serve the superior
is not yet the most inward relation of personality. Even if the servant is ready to sacrifice
everything to his master, there is still something he ought not to sacrifice to him, and that
is his own self, yes, it is conscience. In the name of conscience, then, all serving
individuals are inwardly and essentially freed from the superiority of genius. In the religious
relation, by contrast, conscience goes along; the whole man must give himself up entirely, and
keeps nothing back, not even the least remnant of his own self. For that reason alone the
religious crisis stands, by its nature, high above the merely heroic.

Let one take good care that alien considerations do not creep in and displace the point of
view. The imperfect religion readily takes up one or another world-idea into its distinctive
content, and thereby gains all the greater power to break its historical path. The founder of
the religion here appears not merely as the prophet of the deity and teacher of mankind, but
also as the lawgiver of the state, as the hero of the people. Divine cooperation now heightens
the natural passions, the power of action takes an outward-going direction, fanaticism glows in
enthusiastic hosts, and the religious crisis takes shape as an epoch in world history.
Nevertheless such an epoch-making awakening may by no means be confused with the true crisis
that reaches into the very depths of conscience. This can first set in with the perfect
revelation, thus on the presupposition that the founder of the religion exists as religion
itself.

Since here the coming of the kingdom of God becomes earnest, every finite change in the world
must be reduced to something inessential. Everything outward must provisionally remain as it
is, so that the true transformation may be accomplished in the inward. He who himself brings
% --- p. 241 ---
\opage{241}the kingdom of God with him into the world has no other glory to offer the world.
While so many other great men make themselves known by having so strong a feeling and
practical sense for everything that profits society and benefits the state, it almost
seems that ``this One'' is so absorbed in the invisible, supersensible kingdom of God
that he thereby forgets all sensible and visible necessities. While so many other great
men make themselves known by being, in an eminent sense, both politicians and patriots,
it almost seems that ``this One'' must stand far behind as both politician and patriot.
Were he a politician and patriot, he would surely hasten to issue a well-considered
appeal to his countrymen and fellow citizens for a speedy abolition of the many abuses
and for the most vigorous means of guaranteeing the many rights. If he could only issue
such an appeal and come forward with some practical proposals: it would always raise him
in the eyes of the crowd. Even if the first appeal were later lost in other cries; even
if the practical proposals should later prove --- impractical: it would nevertheless be
a support for his influence and a help to --- the holy cause. For the time being it
could, after all, always win him a name, a party; and this party would then at least
have the satisfaction that its divine leader stood behind no one as politician and
patriot. --- But now he issues no such appeal. He does not cry out into the world, but
he cries into the soul; his whole appeal is lost in the strong, penetrating call:
\emph{What does it profit a man, if he gained the whole world but lost himself, or
suffered harm to his soul.} But this call is, after all (we admit it), no exclusively
patriotic appeal.

The Divine One, then, does not rightly know how to begin with the
% --- p. 242 ---
\opage{242} social, the expedient, the inspiring. The only contribution with which he
enriches the many sources of livelihood, the only proposal by which he undertakes to
protect the many rights and to promote the many trades, he lays down in the one
penetrating admonition: \emph{Seek first the kingdom of God and his righteousness, and
all these things shall be added unto you.} But this proposal, which rests on so dubious
a presupposition, and this ``Seek first,'' which points out so troublesome a detour,
surely lack the practical momentum of the moment, and can certainly make no great claim
on political attention. The divine Teacher and Leader must therefore, for the sake of
his kingdom of God, renounce every outward advantage, every distinction by which the
other great men otherwise know how to gain a hearing for their word with the crowd.
Instead of appearing as the great man who comes bearing the welfare of the whole, he
must even submit to appearing as the lowly man who can offer his own no worldly
advantage. If one asks about power: in this world he has no power whatever. In order to
use power, he would, after all, have to show himself as master; in order to show
himself as master, he would, after all, have to have servants and let himself be
served; but he has no servants, and cannot let himself be served: he must himself
serve. If one asks about riches and honor: it is certainly in vain to turn to him.
Were he to reward his friends with worldly honor, he would, after all, himself have to
be esteemed and honored in the world; were he to be esteemed and honored in the world,
he would surely also have to seek his own honor and not disdain the honor of men
either; but for the sake of his kingdom of God he dare not seek his own honor, nor take
honor from men; no, that honor he must count as nothing. And if he were to reward his
followers with gold and gifts, he would, after all, have to be rich, and very rich at
that. But he is, after all, no rich
% --- p. 243 ---
\opage{243} man; on the contrary, he is a poor man. In order to make men rich in God, he
must submit to being poor in the world, so poor that he has not even where to lay his
head.

The Divine One, then, on these terms cannot possibly appear before the eyes of the
crowd as any eminent man: how, then, must he appear? As a man, a sound and true and
natural man, but only as a man, to judge by the outward. Nonetheless this man must, in
his own person, still make demands on men as no other: his station and dignity require
it, for he is of high birth: God's own Son! indeed, he has a high rank: God revealed on
earth! To exist as the perfect religion is, in other words, to exist as God-man. To
speak to men in God's name is then, for the God-man, the same as to speak in his own
name; he alone knows within himself that he stands to God as the only-begotten Son to
the Father, and therefore he bears witness of himself: \emph{He who sees me sees the
Father: I and the Father are one.} What a contrast between the outward and the inward!
And what, then, has this man by which he can ground the truth of his own testimony? The
word of truth must be proved in the deed: \emph{The Son can do nothing at all of
himself, but what he sees the Father do; for whatsoever things he does, these the Son
also does likewise. For as the Father raises up the dead and makes them alive, even so
the Son also makes alive whom he will.} But what a contrast between the outward and the
inward! The God-man bears witness of himself, and proclaims the mysteries of the kingdom
of God; but other men of God have, after all, also proclaimed the mysteries of the
kingdom of God, and yet no one has
% --- p. 244 ---
\opage{244} ventured to ``call God his own Father, and make himself equal with God.'' The
God-man bears witness of himself, and does great, strange works, heals the sick, and
raises the dead. But other chosen instruments of God have, after all, also done great,
strange works, have also healed the sick and raised the dead; yet none of them has laid
claim to be honored as God himself. --- In the outward, finite sense the God-man demands
nothing for himself; for he is, after all, in the form of a servant: \emph{if I honor
myself, my honor is nothing;} but in the spiritual and eternal sense he demands
everything for himself, for he is, after all, the Mediator between God and men:
``\emph{No one comes to the Father but by me.}'' --- Such loftiness in such lowliness,
such authority in such obedience, such self-exaltation in such self-humiliation has
never, before or since, been seen in the world. Now, then, the choice is to be made,
not between a greater and a lesser degree of the true, but between the truth itself
and the infinite deception; now it is no longer a matter here of finite gain and finite
loss; but here it is a matter of an infinite gain and an infinite loss, of an eternal
life and an eternal death, an eternal blessedness and an eternal perdition. He who
believes in the Son of God has eternal life. But the Son of God says: \emph{He who will
save his life shall lose it, but he who will lay down his life for my sake shall save
it.} In order to follow the Son of God, a man must be ready to forsake everything; for
the Son of God says: \emph{He who loves father and mother more than me is not worthy of
me.} In the relation to the Son of God, then, conscience comes along: he who loves him
must love him with infinite love.

But now the converse also holds, that he who takes offense at the Son of Man must take
conscience along
% --- p. 245 ---
\opage{245} into the offense, and look upon ``this deceiver'' as upon a man who ``has a
devil.'' It is terrible not to be able to love the Son of God in the Son of Man; for he
who cannot love the Son of God must hate him, and he who hates him in the name of the
spirit must hate him unto death. Any other enthusiast and deceiver I can surely also
hate, when I discover his deception. But as there is a limit to the finite deception,
there must also be a limit to my hatred. I can have pity on the enthusiast and feel my
superiority over against the deceiver (for it is, after all, at bottom not I but the
deceiver himself who is deceived). The further the deception is carried, the more the
wretchedness of the deceiver will shine through. But if, on the other hand, I were
mistaken about ``the Son of Man,'' and thought to see in him a deceiver, where then
should my hatred find its limit? Can anything more revolting well be thought than that
one who is only a man should demand of me that I honor him as a God? Or can a greater
encroachment well be made on my personal freedom than this, that another man would place
himself between me and my God? The offense is unlimited, since it is, after all, not the
Son of Man who would deceive me, but, on the contrary, I who am caught in a fearful
self-deception. The offended one's imagined superiority over the supposed ``deceiver''
is lost here in a secret feeling of real impotence, and with the feeling of this
impotence the hatred grows in turn. It is a dreadful thought --- who dares think it
through! --- that any man can hate the truth itself as if it were an infinite
deception, and yet sense the power of the truth in the infinite passion of hatred;
indeed, he who hates the Son of Man must hate him unto death.

Under these presuppositions we may now rightly say of Christ that he not only came into
the world at a critical point of time,
% --- p. 246 ---
\opage{246} but that he also came in order to bring about the decisive crisis. When we
speak thus, we do not speak of ourselves; but the Gospel bears witness to this from the
beginning to the end. In John 12, 31 it says: \emph{Now is the judgment of this world;}
and in the same Gospel, 3, 19---21: \emph{But this is the judgment.} In both places the
word ``\emph{judgment}'' is a translation of the word of the original text:
``\emph{Krisis}.'' --- ``Now is the Krisis of this world'' --- ``But this is the
Krisis'' --- The crisis, then, is this, that the light has come into the world, and has
made a separation between the men who love darkness rather than light, because their
deeds are evil, and the men who do truth and come to the light, that their deeds may be
made manifest. The crisis is this, that Christ stands in the world as the sign of
contradiction, and divides mankind into two parties, which contend with one another
over truth and lie. Not only the history of Jesus himself, but the whole of world
history from Christ onward, may be regarded as an essential continuation and further
carrying-through of this crisis. --- A modern thinker has somewhere put forward the
assertion that the course of world history would have to have been concluded with the
ascension of Christ, if the evangelical presentation of this event were credible. Not
all men are so full of genius that they can with ease enter into such an idea, but
happily that may also be all the same; for the evangelical presentation proves just the
opposite. It says expressly in John's Gospel (3, 17): \emph{For God has not sent his Son
into the world that he should judge the world, but that the world should be saved
through him.} That is to say: Christ did not come into the world in order at once to
condemn the non-believers; but he came in order to bring about the crisis whereby the
unbelieving fall into contradiction with the believers, if perhaps by this
contradiction the light might at last drive out the darkness, and all thereby be saved,
and come to the knowledge of the truth.

% --- p. 247 ---
\opage{247} As the crisis is, so must the criticism be also. Both are governed by the
contradiction. The faith of the Gospel is so far from evading the modern theory of
contradiction that it rather strives to impress it, and that not from an untimely
desire for dialectical exaggerations, but from an irresistible need for a decisive
knowledge. When already in the Fore-Gospel we hear Simeon say to Mary about the child:
behold, this one is set for the fall and rising again of many in Israel, \emph{and for a
sign which is spoken against,} how could we then doubt that the crisis and the
contradiction condition one another, and that Christ himself must be the critical
principle, since he is set for a sign of contradiction. To misunderstand the
contradiction is thus at the same time to misunderstand the crisis; but to misunderstand
the crisis is, in other words --- to miss the critical task.

The sign of contradiction has in this way been misjudged, not only by such as those to
whom the Gospel is foolishness, but also --- by such as, on behalf of pure knowledge,
think they see the wisdom of God in the face of Christ. The first misunderstanding is
characterized in negative criticism; the latter in dogmatizing speculation. Both
misunderstandings spring from one and the same presupposition. For both agree that
faith itself must be dissolved and perish if the contradiction under which the concept
of the God-man suffers cannot find its solution in science. The speculation that means
well by faith labors to dissolve the contradiction in the idea. Will this endeavor
succeed? Speculative dogmatics, after all, unties the knot for us in the christological
idea as a maiden unties a bow. A pity that this idea, when it has been properly worked
through to the clarity of the concept, proves to be in so striking a contradiction with
existence. This is the fatal turning point where speculation goes out of itself and is
transformed into negative criticism.
% --- p. 248 ---
\opage{248} The criticism that means ill by faith now labors to dissolve the concept into
sheer contradictions by demonstrating the impossibility of the idea's existence. The
speculating Hegel teaches that the truth of the God-man confirms itself in the idea; the
criticizing Strauß objects that the idea of the individual God-man negates itself in
existence. This result, sufficiently demonstrated in the history of science, would
undoubtedly already have struck positive theology a stunning, almost deadly blow, had
not this science, by great good fortune, been hovered over by a protecting spirit, an
ever-watchful genius that could avert the danger and absorb the blow, namely ---
mediocrity.

The genuine doctrine of the Middle Way of course discovered at once that Hegel could not
be used. For Hegel held to the orthodox concept of the Son's unity of essence with the
Father. What an enormous blunder on the part of speculation! ``Both sound understanding
and well-explained Christianity teach us, after all, amply enough that such words of
Christ as \emph{I and the Father are one} (John 10, 30) cannot possibly be meant so
seriously that they should actually allude to any true unity of essence of the Father
and the Son. On the contrary! these words in Holy Scripture belong precisely among
those that stand longing for a rational interpreter who can come and give them the
requisite reinterpretation and bring their hidden meaning to light. When, therefore,
Christ says: \emph{I and the Father are one}, this should preferably be taken as a
rhetorical expression for the general truth that there really is a likeness of essence
between the holy God and the moral man. The sensible interpretation of this rhetorical
expression Hegel (remarkably enough, since it is otherwise found in so many
commentaries) unfortunately did not know; and this exegetical ignorance of his has in
turn been the cause
% --- p. 249 ---
\opage{249} of his wasting his time and diligence on fathoming the Son's unity of
essence with the Father.'' --- That the doctrine of the Middle Way was thus able from
the very beginning to overlook the whole of speculation, its exit as well as its
entrance, and thereby avoid a reef on which so many able thinkers have unfortunately
been stranded, is, among the many remarkable occurrences of recent times, by no means
the least remarkable. Meanwhile we believe that the matter can be explained without any
miracle, if only we note that the many thinkers who ran aground in the waters of
speculation ran aground precisely because they looked too deep into logic, and that in
roughly the same sense as when one says of a man that he has looked too deep into ---
the bottle. Sobriety is a noble virtue; and this must be granted to mediocrity: it
knows how to keep within measure. No one, however unbelieving he may be, will so easily
permit himself to mock the sober representatives of the half-criticism as those
unbelievers at the first Pentecost permitted themselves to mock the Lord's apostles,
when they said: ``they are full of new wine.'' --- Thus in the dangerous moment when most
thinkers looked too deep into logic, mediocrity sat erect and sober at the helm, and
observed the precaution of not looking at logic at all, so as not, like so many others,
to become confused, get out of the right track, and miss the course. Had the theology of
the Middle Way been given to see its own most reverend face in logic --- indeed, it
would have been a dreadful moment! The Hegelian Medusa would certainly have turned
mediocrity into a stone, and so after all have forced it to sink to the bottom in the
depths of speculation. --- For the enlightenment of the uninitiated, merely the following
remark. When the cautious expression ``likeness of essence'' is so urgently recommended
to the teachers of the Church, instead of the Church's own bolder ``unity of essence,''
one sees at once what it actually is that the doctrine of the Middle Way aims at. It is
the difference and
% --- p. 250 ---
\opage{250} opposition between the essence of God and that of man that it wishes to
impress with due emphasis. And it certainly deserves all appreciation that this
difference be properly impressed. But had the theology of the Middle Way only felt
disposed to make even a fleeting acquaintance with speculative logic, it would
undoubtedly have acquired some distrust of the effect of its much-recommended
expression: ``likeness of essence.'' True likeness of essence must, after all, be
grounded in a unity of essence. Only a superficial, indifferent and accidental likeness
can fall outside the essential unity. This essential unity logic calls
\emph{identity}. Of this identity it is expressly taught that it cannot be thought
without \emph{difference} and cannot subsist without \emph{opposition}. The exertion of
the half-criticism in impressing the difference is to that extent rather superfluous.
--- It was well, however, that mediocrity paid no heed to logic, and nevertheless
overcame Hegel and made an end of speculation; if it should now only succeed with just
as much good fortune in overcoming Strauß and making an end of Feuerbach (and surely
it must succeed some day): then Christianity would, after all, have in this theology a
truly solid and defensible bulwark, and be secured for many years against all manner of
hateful attacks.

Meanwhile it might well be a long way off before all these matters are settled. Perhaps
the struggle drags on, and it is a while before the half-criticism sees itself in a
position to make use of the heavy artillery and deliver a decisive blow; perhaps, when
the half-criticism after such a decisive blow then returns home again to faith, it
returns home with the usual honors, i.e., with half a victory; perhaps, when the enemy
(for enemies always have their own interpretation) then interprets the half victory for
us as a whole defeat; perhaps, we say, a new notion will then again occasion new
complications, and the struggle begin all over. If we now take all these circumstances
into serious
% --- p. 251 ---
\opage{251} consideration, and consider besides that the faith of the Gospel always is
and remains a vital matter, whose decision the person concerned cannot well postpone to
a time convenient for science, then it must become quite evident to us that the
Believer ought at once to look about for another way out, and expect another solution
to the riddle of the contradiction.

It is the presupposition itself that must be changed. Were the believing appropriation
of Christ's Gospel to suffer any detriment because science now and then fell into
embarrassment with the concept of the God-man, then Christianity would long ago have
perished; but the matter stands the other way round: \emph{the incomprehensibility of
the God-man is the religious presupposition that must be regarded as the right,
unshakable foundation of Gospel faith.}

When, now, science, as is fitting, continues to search after the contradiction; when by
its unceasing searching it constantly enriches itself with new, more cutting, still
deeper-penetrating contradictions; when it finally melts and works together all these
contradictions into one single great fundamental contradiction, in order to crush ``the
prophet of Nazareth'' under the hundredweight of this fundamental contradiction: what,
then, does this mean? It does not mean at all that the perfect religion shall now fall
before the contradiction, but it means that ``the prophet of Nazareth'' is the true sign
of contradiction, in truth both God and man. And when, then, the secret accomplices of
pure knowledge whisper to one another: contradiction, contradiction, what a
contradiction! and when with this whispering they steal forth from the ambush of
reflection in order to dissolve ``the actual God-man'' into a misty generality: what,
then, does this prove? It does not prove at all that ``the God-man'' must now dissolve
into a generality; but it proves that pure knowledge is on the point of
% --- p. 252 ---
\opage{252} dissolving itself in the ``indissoluble contradiction'' of the God-man. And
when, then, this secret becomes known; and when, then, all the loud-voiced critics begin
to shout: contradiction! contradiction! what an enormous contradiction! and with these
shouts set the world in motion to storm the ramparts of the Gospel and thrust ``the
historical Christ'' out of history: what, then, does this bring about? It does not
bring about at all that Christ passes out of history; but it brings about that the light
of the Gospel rises in faith; indeed, it brings about that the aged Simeon steps forth
with the child in his arms, and testifies before all the storming host what he once
testified before Mary: {\large This one is set for the fall and rising again of many in
Israel, and for a sign which is spoken against.}

The world contradicts Christ; it does so with great proficiency (proficiency comes with
practice, and it has practiced at this from the beginning). A guileless soul, who does
not know the secret of contradiction, might perhaps easily hit upon making the following
reflection: ``If Jesus really did all the miracles ascribed to him in the four canonical
Gospels, and accomplished each of these wonder-works exactly as the Evangelists tell it
to us: must not even the most hardened Pharisees at last have convinced themselves that
he was, if not God's own Son, then at least an extraordinary messenger of God? When
they nonetheless not only persisted in denying his divine mission, but even sought his
life, then this fact furnishes, as it seems, a quite conspicuous proof that the
evangelical miracle narratives must be, if not downright invented, then at least
embellished, or in some other way distorted.'' --- Such misgivings arise only from a
misunderstanding of the essence of the supernatural and of the immediate effect of the
wonder on the eyewitnesses, a misunderstanding that proves only all too
% --- p. 253 ---
\opage{253} emphatically how little one is able to enter into the actual circumstances of
that time and judge the evangelical facts, when consciousness has become a stranger to
the original impression of revelation. Certainly it is a contradiction to want to kill
a man who himself raises the dead. The chief priests and the Pharisees acted (John 11,
47) quite certainly according to the principle of contradiction; that shall never be
denied. But this contradiction is so far from disproving the reality of the wonder that
it rather proves its peculiar power, in that it shows us the repellent effect that
every wonder-work must necessarily exert on the non-believers. If we see the matter
from this standpoint, the evangelical accounts will in all respects demonstrate their
inner truth. --- In order that the true wonder may take place, the condition of faith is
presupposed. In Matth. 13, 58 it says expressly that Jesus did not do many mighty works
in his own country and in his own house, \emph{because of their unbelief.} The wonder is
not offered at all to the unbelieving. The clearest proof of this we have in the
following account (Matth. 16, 1---4): \emph{And the Pharisees and Sadducees came forth,
tempted him, and desired that he would show them a sign from heaven. But he answered and
said unto them: in the evening you say, it will be a fine day; for the sky is red; and
in the morning: there will be a storm today, for the sky is red and dark. You
hypocrites! the face of the sky you know how to judge, but the signs of the times you
cannot. This evil and adulterous generation seeks a sign, and no sign shall be given
it, but the sign of the prophet Jonah. And he left them, and went away.} --- What, now,
can have moved the Pharisees and Sadducees here to forget their own
% --- p. 254 ---
\opage{254} disputes and unite in so solemn a procession in order to request of the
wonder-worker a special sign, when he otherwise did so many (indeed, in their opinion,
certainly far too many) signs among the people? The answer lies close at hand. The
scribes, who despised the people, naturally also despised the signs that awakened faith
in the people. The immediate relation of unbelief to the wonder was thus the same then
as it is today. The same contempt with which a David Strauß in the present speaks of
those healings of the sick in Galilee, the same contempt and mistrust the Pharisees and
Sadducees already showed in that time toward the same healings of the sick. The scribes,
who felt themselves so exalted above the people, doubtless thought that for their part
they could well lay claim to a more distinguished sign, a sign whose nature they
themselves were now pleased to determine, a sign that was able to convince them. This
sign was naturally not to be any ordinary sign, nothing so insignificant as, for
example, that a poor blind man got his sight, or a madman his reason; their sign was
not to take place for the benefit of any single human being at all; no, it was to be
precisely a sign exalted above human existence, a magnificent sign, surprising all the
world, a sign that could shake nature itself, a shining, splendidly radiant sign --- a
sign from heaven!

``That Jesus did not, after all, seize this opportunity to show his opponents in an
instant who he was must certainly arouse astonishment. It was, after all, undeniably
both highly placed and respectable men who here condescended to request a sign of him.
Had he now seized the opportunity, these scribes would doubtless have become his most
zealous followers, instead of which, disappointed in their expectation, they naturally
had to become his bitterest enemies. Had he only made use of this occasion and given
the sign demanded, how much would then not have changed in his outward situation;
% --- p. 255 ---
\opage{255} the many misjudgments, the many rebukes, indeed the heavy suffering and the
bitter death might then perhaps have been avoided. But he did not seize the opportunity,
he did not make use of the occasion, he left them, and went away --- in anger. How
criticism must marvel at such conduct!'' Then let it marvel! But should a critic here
wish to try to make use of this narrative as a proof that ``the wonder-worker presumably
did not dare venture an attempt at a miracle in the presence of so many attentive
witnesses'': then we must remind criticism of what happened with Lazarus in Bethany
(John 11), and still earlier with the widow's son at Nain (Luke 7, 12---17). How public
and open is this wonder, with what vividness is this scene presented! ``The Lord saw
the funeral procession as he came near to the gate of the city; he saw the widow and had
compassion on her from the heart, and said to her: weep not. And he came up and touched
the bier (and they that bore it stood still), and he said: young man, I bid thee,
arise! And the dead man sat up, and began to speak, and he gave him to his mother
. . .''

If anyone thinks that such a wonder-work must necessarily awaken faith in all those
present, he is remarkably mistaken. To convince ourselves of this, let us place, for
example, some hardened Sadducee or other as eyewitness. How, now, would a Sadducean
critic judge this scene? Without doubt thus: ``It was indeed a true piece of luck for
the young man that he was not buried alive; but it was a fatal piece of bad luck that
this alleged prophet, this wandering miracle-monger, should come to the gate of the
city just at the moment when the apparently dead man was being carried out, and at once
be so fortunate as to make the discovery and exploit the occasion. One must grant him,
however, that he knows how to make it both solemn and touching. See, fear stands
% --- p. 256 ---
\opage{256} painted on every face: it is at bottom a shattering scene. The poor people
praise God that the prophet has saved a dead man from the grave; they had better ask
the authorities to see to it that honest folk get time to die before their neighbor has
them buried.'' --- A modern critic need not at all be ashamed of the Sadducee's view; on
the contrary, he can even justify it by virtue of science's own logical principle. If we
think a little more closely, we at once find the proof; here it is: ``If the young man
had been really dead, he would not have risen again. Now he has, after all, risen again.
\textit{Ergo}: he has not been really dead.'' This inference is stringent and pure, as
pure logic and pure knowledge.

We can go a step further still. Suppose that a wonder-work were so conspicuous, so
incontrovertible, that no eyewitness ventured to deny its wondrousness: would such a
wonder-work then be able to force faith forth in the non-believers?\footnote{%
In a treatise excellent from its standpoint (Ueber das Wunder. 1839.) Feuerbach has set
forth this truth, among other things, in the following terms: ``Der Glaube an das
wunderbare Faktum ist immer zugleich der Glaube an die \emph{Wahrheit der Vorstellung}
z. B. der Auferstehung, welche durch dieses Wunder bewährt werden soll. \emph{Es ist
eins, ob ich das Faktum} oder die dem Faktum zu Grunde liegende \emph{Vorstellung} (oder
Lehre) hinstelle: die Vorstellung wird mir dadurch richt gewisser, nicht klarer, nicht
näher gebracht, sie mag mir nur \textit{in abstracto} als Vorstellung oder \textit{in
concreto} als sinnliches Faktum dargestellt werden. Wer mir eine Vorstellung als Faktum
hinstellt, der will nur kurzen Proceß mit mir machen. `Es ist Faktum! \textit{punctum
satis}. Was willst du dagegen machen? Du mußt es glauben'. Aber eben dadurch werden
meine Zweifel nicht aufgeklärt, sondern nur niedergeschlagen --- das Faktum ist ein
Gewaltstreich --- keine innern Gründe angegeben, die die Sache wahrscheinlicher
machten. Das dogmatische Faktum ist illiberal, anmaßend, tyrannisch, es macht mir eine
Vorstellung zum Gesetz, es schiebt mir sie ins Gewissen hinein, indem es dieselbe zu
einem Faktum, d. i. etwas nicht mehr zu Bezweifelndem macht, ohne doch die Zweifel und
Schwierigkeiten aufzulösen''. Ludwig Feuerbach sämmtliche Werke 1ster Band Pag. 20.}
On this point too the Gospel can give us
% --- p. 257 ---
\opage{257} the necessary enlightenment. In Matth. 12, 22---24 it says: \emph{then there
was brought unto him one possessed, who was blind and dumb; and he healed him, so that
the blind and dumb both spoke and saw. And all the people were terrified, and said: is
not this the son of David? But when the Pharisees heard it, they said: this one does
not cast out devils but by Beelzebul, the prince of the devils.} What further need have
we of testimony! The Pharisees could not and dared not deny the fact; they could not and
would not believe the wonder: how, then, did they help themselves out of the
embarrassment? Answer: by an interpretation. Indeed, the art of interpretation is truly
a wonderful art. If the raw mass takes offense at Christ, it reaches for stones; but
when the scribes take offense at Christ, they reach for an interpretation. This means
has never yet failed. What defense is so powerful as a supple interpretation! The divine
authority of the Son of Man would sooner subdue the whole enraged crowd, which can defend
itself only with stones, than a single scribe who knows how to defend himself with an
interpretation. And what weapon of attack is so dangerous as a cunning interpretation!
Let
% --- p. 258 ---
\opage{258} the crowd band together against the Son of Man as against a robber with
swords and staves: it is of no avail, he still gets through, he conquers; but when the
scribes begin to entangle him in their net, and bare the dagger of interpretation, then
he must give himself up for lost, and cannot get through. All other demons the Son of
Man can subdue; the demon of interpretation alone he has hitherto not been able to
subdue. If modern criticism reflects a little, it can very well sign the Pharisees'
interpretation: \emph{He does not cast out devils but by Beelzebul.} Who, then, is
Beelzebul? The spirit of illusion, which is itself an illusion. Yes, the interpretation
holds good: Christ's wonder-work is a demonic illusion.

What has been said here of the wonder holds also of the word. The relation to God, the
righteousness and freedom that express themselves in Christ's word cannot be understood
without an unconditional faith in his wonder-working power. The Pharisees contradicted
his word because they contradicted his wonder-works; the modern critics do likewise.
However great the difference may be between the culture and way of thinking of that time
and of the present, the contradiction of unbelief is nevertheless at bottom the same.
Whether we mingle among those old Pharisees who were zealous against the uncalled rabbi,
explained his miracles as effects of demonic sorcery, and took offense at him because he
violated the Law of Moses, disturbed the good order of things, and made himself the Son
of God; or whether we place ourselves in the ranks of the enlightened Free Thinkers and
mock ``this equivocal Jesus, whose invented miracles prove superstition's attempt to
break the law of nature, whose fantastic doctrine of sin and the forgiveness of sins
confuses the moral concepts of true duty and sound virtue'': the result of the
contradiction is in both cases the same:
% --- p. 259 ---
\opage{259} \emph{Jesus of Nazareth has, by giving himself out as the only-begotten Son of
God, disturbed the divine laws in the world of nature as in the world of spirit,
confused many souls, and stamped himself as a deceiver of the people.}

It is one of the great labors of the world-spirit to promote and multiply this
contradiction in all directions and carry it through all crises. But while the
world-spirit employs the many contradictions only to confirm the negation, the Spirit
of Christ also now goes through the ages in order to take up the same contradictions,
turn them round, and dissolve them in faith. Reflection then fights in turn for a
believing appropriation.

\emph{How unjust and false, then, is the assertion that Jesus' wonder-working power
should in any of its manifestations whatsoever have brought about a disturbance of the
forces of nature, and wrought confusion in the world of souls.}

True it is that the Lord once, when he had sat down among the wedding guests at Cana in
Galilee, changed water into wine. But what forces were here disturbed? Could it wrong
the pure water that a thought of the Godhead changed it into wine? Or could it perhaps
disturb the order of nature that the Son of God would hallow one of the beautiful but
fleeting moments of human life with a reflection of eternal glory? The world has, after
all, so many teachers who prefer to begin by changing wine into water; one could then
easily forgive ``this One'' that he --- for a change! --- made the beginning by changing
water into wine.

True it is that the Lord once, when he was crossing the Sea of Gennesaret with his
disciples, rose from his sleep, and rebuked
% --- p. 260 ---
\opage{260} the weather, and stilled the storm. This, now, according to the world's
judgment, is supposed to be a revolting encroachment on the rights of nature. But is it
not, after all, according to the world's judgment, a natural right that men defend
themselves against the blind forces: why, then, will one accuse ``this One'' because he
heard the disciples' cry of fear, and in the moment of danger made use of his divine
right in order to put to shame those ``of little faith.'' The world has, after all, so
many strong spirits who rush along with storm and roar with tempest; one could then
easily forgive ``this One'' that he --- for a change! --- will stop the storm and make
the roaring tempest --- dead calm.

True it is that the Lord, when he walked as a man among men, needed only to stretch out
his hand in order to heal the sick, or speak a word in order to raise the dead. He had
power, he worked with might; but did he then disturb the human forces? Everywhere that
the mighty of the world rightly use their power and work with might, it can never fail
that the number of the lame and crippled is greatly increased; why, then, will one so
harshly accuse ``this One'' because he, for a change, uses his power only so that the
number of the lame and crippled is somewhat diminished. The world has, after all,
rulers enough who can kill the living, and are found willing to sacrifice the many in
order themselves to live; why, then, will one take offense at ``this One'' because he
--- for a change! --- now prefers to raise the dead, to sacrifice no one but himself,
and at last to go to death for all. --- Nonetheless he was, after all, a ruler: the word
that can raise the dead could surely also kill the living. The world may well dread
such a ruler, for --- ``how easily he might misuse his power.''

``Let one imagine a noble autocrat, a philanthropic despot (despots too can, after all,
sometimes be philanthropic)
% --- p. 261 ---
\opage{261} equipped with such a power; let one imagine that this despot had enemies; let
one imagine that these enemies misunderstood his long-suffering, and interpreted the
lenient use he made of his mysterious power as a sign of real impotence; let one imagine
that these enemies were ungrateful and unjust and bloodthirsty alike, so that they would
not only kill the ruler himself, but at the same time bring all his philanthropic
arrangements to nothing: how must the noble autocrat then not be inwardly indignant at
such enemies; how must he in the moment of danger, when he stood alone surrounded by
these ungrateful, unjust, bloodthirsty enemies, not be tempted to rise up with
supernatural might, make use of his power, and annihilate them all with a single
word!'' --- The Lord Christ had such a power; but --- he did not use it. He too had
enemies, unjust, ungrateful enemies: he was angered at them, he was inwardly grieved
over them; but the power he did not use. The enemies surrounded him and took him
captive: he stood alone in the midst of them; they thirsted for his blood; he needed
only to speak a word to cast them all to the ground; but --- the power he did not use.
Had he at any moment of his life forgotten himself and misused his wonder-working power,
it would certainly not have been forgotten by his many accusers. False witnesses brought
all manner of charges against him; but there was, after all, no one who could come
forward and testify to the misuse of power. No Pharisee could come forward and complain
that ``the captive deceiver of the people had once by a word of sorcery robbed him of
the light of his eyes.'' No Sadducee could come forward and complain ``that the captive
enthusiast had once struck him in anger and robbed him of the use of his limbs.'' No,
he did not disturb the natural order, he did not crush human
% --- p. 262 ---
\opage{262} independence. He would, after all, not tear down but build up, not wound the
healthy but heal the sick.

There were many sick who had need of the physician. And the great Physician began by
going about in all Galilee, and healing all sickness and all infirmity among the people.
He could heal, for the Creator of nature laid the power of healing in his holy word; he
would heal --- all who called upon him in unconditional trust and faith. And there were
so many who had need of the physician, because they suffered ill. Now in the throng, and
now on the lonely roads, our eye meets these sufferers. Their cry reaches him at every
step he takes; and he stops his course, and does not grow weary of hearing their
prayers. There it is a leper who falls down before him: ``\emph{Lord, if thou wilt,
thou canst make me clean!}'' Yonder comes the centurion from Capernaum: ``\emph{Lord, my
servant lies at home sick of the palsy, grievously tormented!}'' Here it is the sick
woman, she who for twelve years has hoped in vain in the art of the physicians; see, she
seeks him in the press of the crowd, she follows after him so quietly, if only she might
come so near to him that she touched the hem of his garment; for, she says in her heart,
\emph{if I but touch his garment, I shall be healed.} --- There are so many sick who
have need of the physician, and here is a physician who can help, and who never grows
weary. No, he never grows weary of speaking the word and stretching out the hand. To the
one it sounds: \emph{I will, be thou clean!} to the other: \emph{I will come!} to the
woman: \emph{Daughter, be of good cheer, thy faith hath saved thee!} To the one sick of
the palsy: \emph{Son, be of good cheer, thy sins are forgiven thee!} and to the blind by
the wayside: \emph{Be it unto you according to your faith.}

% --- p. 263 ---
\opage{263} %
\emph{So his fame went out throughout all Syria, and they brought unto him all that
suffered ill, that were afflicted with divers diseases and torments, those possessed,
and lunatic, and sick of the palsy, and he healed them. And there followed him much
people from Galilee and Decapolis and Jerusalem and Judea and from beyond Jordan. But
when he saw the people, he went up on a mountain, and when he had sat down, his
disciples came unto him, and he opened his mouth, and} --- to teach the people from
the mountain. --- Who can speak here of disturbance and of leading the people astray,
here, where everything round about breathes as in the peace of eternity, while the
Master lifts up his clear voice? See, the mountain is his teacher's chair, the disciples
sit at his feet, silent multitudes listen to the fundamental words of the Gospel, souls
drink from the well-spring of life! Who dares call this mountain a mountain of
disturbances? Pious legend has called it the Mount of Beatitudes: it is a good name, it
echoes the Savior's words; it is a blessed name, it opens so bright a prospect. If
anyone will now look down from this Mount of Beatitudes toward the desert where Sinai
looms gray in dark shadows as the mountain of terrors: then he shall find true what John
testifies, that \emph{the Law was given by Moses, but grace and truth came by Jesus
Christ.}

Gospel of Matth. Chap. 5.

\begingroup\large
Think not that I am come to dissolve the Law or the Prophets: I am not come to dissolve,
but to fulfil. For verily I say unto you: till heaven and earth pass away, not even the
least letter or one tittle of the Law shall pass away, till all be fulfilled.
\par\endgroup

% --- p. 264 ---
\opage{264} Has a deceiver of the people ever spoken to the crowd in such language? It is,
after all, well known that he who would win the favor of the crowd must conform to the
crowd's wishes. But it can hardly be the crowd's wish that the old Law should be
impressed down to tittle and letter; it can hardly be the crowd's wish that what is
burdensome should be made still more burdensome. The crowd wishes no exaggerated
impressing of what exists at all. What the crowd wishes is, on the contrary, a more
reasonable change in the Law, answering to the demands of the time, a change that
naturally aims at a mitigation. This the popular orator knows very well. When, then, it
is really of consequence to him to win the people and work upon the crowd, he makes sharp
speeches against the Law and what exists. The more sharply he speaks against the Law,
the more he rails against what exists, the better he is understood by the crowd; the
better he is understood by the crowd, the more easily he leads the people astray: the
restless crowd will not merely be led, ah no, it will --- be led astray. It is, after
all, well known that he who would move the mass, and thereby bring about a great
upheaval in the people's conditions, must conform to the mass's power of comprehension.
But he who conforms only to the mass's power of comprehension dare not place the change
he intends to bring about in the inward man. This change can, after all, only take
place through each individual man himself being changed: how difficult to understand
such a change, how slow, how laborious to carry it through! Never yet has any mass of
people been enthusiastic for such an inward, invisible and essential change. What the
mass desires is a change in the outward. This change can then show itself especially
in a striking contrast between the old
% --- p. 265 ---
\opage{265} and the new. The new must naturally be an outward improvement, and this outward
improvement must then presumably also let itself be carried through without the individual
persons exactly needing to exert themselves all too much with their own inward
self-improvement. This the popular orator also knows very well. If, therefore, it is his most
important task to set the masses in motion, he at once begins to attack the old and to extol
the new. The better he now succeeds in throwing the corruption onto the old and coaxing
blessedness out of the new, the better he shall succeed in being understood by the mass; and
the better he succeeds in coming to an understanding with the mass, the more easily will he
also succeed in leading the people astray: the mass will not merely be led, ah no, it
will --- be led astray.

\emph{Think not} --- so speaks the divine Teacher from the mount --- \emph{think not that I
am come to dissolve the Law or the Prophets.} Christ must surely have known men's impatient
longing for outward change, since he speaks thus. From the mount, from his high seat, he was
able not merely to survey the numerous multitude of the assembly, but in truth also to see
through the mass. As his gaze then saw through the mass, and saw the restless drives
fermenting in the crowd, he laid great emphasis on the words: Think not that I am come to
\emph{dissolve the Law or the Prophets}; I am not come to \emph{dissolve}, but to
\emph{fulfil}. It is as if by this warning he would lay down his whole authority in the
speech of the Prophets, in the words of the Law: \emph{For verily I say unto you, till heaven
and earth pass away, not even the least letter or one tittle of the Law shall pass away, till
all be fulfilled.} --- It is as if he could not enjoin this admonition strictly enough; he
even lays a punishment on the misunderstanding; he goes into
% --- p. 266 ---
\opage{266} particulars, hear him! \emph{Whosoever therefore shall loose one of these least
commandments, and shall teach men so, he shall be called the least in the kingdom of
heaven.} And yet he loves the people, and yet he speaks to the crowd: how is this now to be
understood? Is he perhaps himself in every way secure in the old, and satisfied with what
exists? Has he, then, nothing new at all to offer? Surely it can never be his meaning that
everything in actuality shall remain as it is. Hear him! \emph{For I say unto you: except
your righteousness exceed the righteousness of the scribes and Pharisees, ye shall in no case
enter into the kingdom of heaven.} --- So he too is not, after all, content with what exists,
and by no means secure in the old. He too will bring about a change. But what change? Yes,
it is precisely about this that the contradiction turns!

When Jesus' speech had sounded among the crowd, there were surely many who afterwards must
have felt themselves in contradiction with the speaker. The restless, who felt the deep,
well-enough-known need for ``an outward change,'' must surely have found it highly
contradictory that ``the great Prophet should thus begin by securing what exists and holding
to the old.'' The chief priests and the rulers of the people must surely have felt themselves
highly alarmed that ``the bold popular orator was casting so dangerous a leaven into the mass,
and --- what presumption! --- dared to speak to the crowd of a new righteousness, a
righteousness that should even exceed that of the scribes and Pharisees.'' --- The
contradiction was thus many-sided enough; for, as both parties thus contradicted Christ, they
at the same time, of course, contradicted each other.

Yet it is now centuries since these contradictions first broke out in the vehement strife of
those parties. The world's
% --- p. 267 ---
\opage{267} circumstances have changed since that time; those antiquated party quarrels over
the old and the new can no longer disturb the calm contemplation of the Messiah's appearance
among his own.

Theological science, then, now sees itself to that extent also able to dissolve those
contradictions in the many-sidedness of the idea itself. How much that is fitting, apt, yes,
even edifying, can now be brought forward about the relation between the lawgiving on Sinai
and this Sermon on the Mount, between Moses and Christ, between Jehovah and the heavenly
Father, between the Law and the Gospel. But nevertheless the power of science does not reach
so far that it should be able to remove the original ground of the misunderstanding or
abolish the cause of the strife: each time the existential conditions come into force anew,
the contradiction too breaks out in new, surprising forms. An unmistakable proof of this is
furnished by the negative criticism. Let the offended Free Thinker himself give us a little
lecture.

``There is, after all, no one in the present age who with greater cunning and more talent has
understood how to undermine Christianity than theologians and pastors. In order to lure the
crowd to themselves and --- to Christ, they have themselves let themselves be lured out onto
the scientific glare ice of culture, where one stumbles at every moment and loses the foothold
of faith. --- In order, however, to keep one's balance under such precarious circumstances,
one must now see to it that one turns and twists with the greatest possible suppleness. How
does it go, then? Quite well! One merely pretends that there is the most intimate and
beautiful harmony between drive, reason, and Christianity; then one keeps one's balance, then
it goes quite well. Is it, then, really possible to hold together such heterogeneous
components? Yes! by neutralizing the hostile forces, lulling the drive to sleep, flattering
reason, and drawing the sting out of Christ's words, it is to a certain degree even a little
more than possible. The present age, which has at last
% --- p. 268 ---
\opage{268} %
given up those medieval-unheimlich attempts at making gold, has instead gained all the
greater practical skill in the manufacture of all sorts of fancy goods and interesting
compositions. If one cannot make gold, well, one can at any rate make silver, namely new
silver: that is always something. It is roughly the same with Christianity. The genuine ore
of the Bible's word, that pure gold of the Gospel which once blazed in the `shield of faith'
and the `helmet of salvation,' has quite vanished. No one knows what has become of this gold;
one digs and gnaws, searches and seeks everywhere, but --- it is no longer to be found. What
is to be done about this now? For want of the best, one chooses the next best. In hard times
one helps oneself as one can --- by an ingenious composition. But New Christianity and New
Silver are nevertheless not \textit{al pari}. To deceive one's neighbor with the new silver is
dangerous for name and reputation, for the assay says incorruptibly that `the old is better';
but to beguile the crowd with a new Christianity, yes, that is an honest business, for the
falsified assay assures us, after all, that `the new is better.'\,''

``What is brought forward here is no empty accusation, but a notorious fact: a few examples
will suffice to set the matter in its proper light. The task is to examine whether the
theologians and pastors of the present age now also really teach and preach as Christ himself
preached.''

``In the so-called Sermon on the Mount Christ begins by calling those blessed who are poor in
spirit. To gain citizenship in the kingdom of heaven, then, one need only be poor in spirit.
Here at once is something that appeals to natural feeling. A truly edifying sermon on this
text will surely be able to procure, not only for the pastor himself but for the whole
congregation (the parish clerk and the pew-opener included), a devout relief from all doubting
% --- p. 269 ---
\opage{269} thoughts, a joyous victory over all proud scribes, a refreshing triumph over the
vain wisdom of the world. That what is hidden from the wise and the scribes can be revealed
to the babes (likewise to the parish clerk and the pew-opener) is undeniably, and beyond all
contradiction, a very edifying thought. But should anyone among the babes of the congregation
then want to try going further with the simple principle that all Christians must therefore,
in religious respect, be equally wise, since they are all without exception poor in spirit:
what must he then not experience? Why, then he must surely experience that there is, after
all, a remarkable difference, that the one in poverty of spirit is, after all, a good deal
wiser than the other. See, the parish clerk, for example, is at once wiser than the
pew-opener; the pastor, again, is far wiser than the parish clerk, the dean wiser than the
pastor, the bishop wiser than the dean, and the \textit{Summus Theologus}, since the abolition
of the papacy, the wisest man in Christendom. It goes as best it can with the evangelical
word: `blessed are the poor in spirit.' A \textit{Summus Theologus} cannot possibly be poor in
spirit; but notwithstanding, he must, after all, be assumed to be able to become blessed,
inasmuch as he is the wisest man in Christendom; on the other hand, the pew-opener can never
hope to become \textit{Summus Theologus}, but notwithstanding, she must surely be assumed to
be able to become blessed, inasmuch as she is --- poor in spirit.''

``But how does it stand, then, with `becoming blessed' and `being poor in spirit'? If the
simplest Christian knows just as fully as the wisest what essentially belongs to faith and
blessedness: what, then, does it mean to be \textit{Summus Theologus} and the wisest man in
Christendom? --- It means that the Christian faith needs thorough learning and scientific
% --- p. 270 ---
\opage{270} %
culture in order to be rightly understood, ably defended, and duly preached in the Church.''
--- No, is that really true? Was Christianity, then, not rightly understood and ably defended
and duly preached in the time of the Apostles? Was St.\ Peter, when we look more closely,
perhaps \textit{Summus Theologus}? Or was it perhaps the theological doctor's diploma that
overtook Paul on the road to Damascus? --- ``With the Apostles it was quite another matter;
they had, after all, the supernatural gift of grace of the Holy Spirit.'' --- An entirely apt
answer! equally illuminating, equally weighty, and in every way equally welcome, whether it be
given us by the pew-opener or by \textit{Summus Theologus}. So: when the congregation had the
gift of grace of the Holy Spirit, Christianity needed no \textit{Doctores Theologiæ}; now, on
the other hand, when the Church no longer has the supernatural gift of the Spirit,
Christianity cannot subsist without a \textit{Summus Theologus}. Yes, now we understand it:
according to the evangelical fundamental law all believers are simply and plainly poor in
spirit, and Christ in that respect did not think of making any ecclesiastical distinction
between being a pew-opener and being \textit{Summus Theologus}; the distinction must rest
solely on the supernatural gift of the Holy Spirit; but since this gift of grace has gradually
diminished and in the very latest time has for the most part quite gone out, ecclesiastically
minded men have been intent on making up the loss by more natural means. On this occasion,
then, the fundamental law has been revised anew, and, in consideration of the needs of the
time, one has even found oneself entitled to add a little editorial note, reading as follows:
``According to the Founder's divine charter (cf.\ the first article of the Sermon on the
Mount), all true Christians without distinction are poor in spirit, and no one should in that
respect be granted any other precedence than that which the Holy Spirit himself might be
pleased to assign him in the congregation; but
% --- p. 271 ---
\opage{271} since the times are worldly, and the worldly times advance in a steady increase of
all manner of learned knowledge and rational insight, it has been found necessary, for the
needs of the Church, to enrich faith with sundry goods of learned knowledge, and for that
reason the poor of the congregation should be kept under constant supervision by the greater
landowners or scientific proprietors, who, under the presidency of \textit{Summus Theologus},
naturally rank according to their literary fortune.'' God knows whether the pew-opener, too,
can just so readily see the deeper meaning that underlies this editorial note; but it is
certain that \textit{Summus Theologus} can see it very well, and we venture to think that we
too can see it. ``Faith has at bottom nothing to do with learned science, and yet faith ---
\textit{mirabile dictu} --- cannot subsist without learned science.'' A superb ambiguity! It
would take threefold witchcraft now, should the highly learned patrons of the Church ever lose
their balance. If science now and then becomes a bit too intrusive, if one gets into a bit of
a fix with the negative criticism: well, then one can crawl together in evangelical
simplicity, take the fundamental law literally, and fold one's hands: `Blessed are the poor in
spirit'! If, on the other hand, the layman steps forward too boldly, if the babes, likewise the
parish clerk, likewise the pew-opener, now also want to put in a word in the discussion of
ecclesiastical questions: what does one do then? One naturally appeals to the editorial note
and proves the necessity of interpretation from the analogy of Scripture. Scientific
interpretation always aims, after all, at finding the light, and there are parallel passages
enough in holy Scripture that bear witness to the light. With the help of the light the whole
becomes clear: the light of Scripture shines only in interpretation, the light of
interpretation shines only in hermeneutics, but the light of hermeneutics first dawns in
theology: should a \textit{Summus Theologus}, then, not
% --- p. 272 ---
\opage{272} know how to convince the congregation that Christ's word about the poor in spirit
may by no means be taken according to the letter.''

``But, to end this ambiguous game, one will surely in the end be forced to strike out on one
of two roads. Either one admits outright that Christian conviction necessarily depends on
rational insight and rational grounding: and then the doctrine of poverty of spirit has lost
its evangelical significance: the knowledge of reason conquers Christianity. Or one takes the
doctrine of the spiritually poor quite literally. In this case one preserves Christianity, but
comes into contradiction with reason. The inner light of which the Gospel speaks is, mind you,
by no means the light of reason; no, it is a quite peculiar, mysterious, supernatural light, a
light that by a miracle can dawn just as well for the uncultivated as for the cultivated. So
much, then, is settled: he who seeks this inner light in poverty of spirit should not turn to
our scientific theologians and modern pastors; he must far rather choose the poor mystics and
strict monks of the Middle Ages as his teachers.''

``Christ further calls those blessed who have a pure heart, and promises them that they shall
see God. These words too, `blessed are the pure in heart,' commend themselves to pious
feeling. It is, after all, so noble, so lovable to have a pure heart. Yes! in pure words on
pure paper we would all very gladly be pure in heart; but alas, unfortunately! in real life,
in flesh and blood, no heart is pure: we are all, after all, `sinners.' It helps but little
that the pastor explains to us that we should, as far as possible, strive to purify our
hearts; and that the theologian promises us that, if only we faithfully persevere in this
endeavor, together with the grammatical explanation of holy Scripture's Hebrew
% --- p. 273 ---
\opage{273} \textit{bare-lebab}, we shall then also, as far as possible, come into an ever
closer and closer communion with God. This doctrine of an infinite approximation is the
downfall of Christianity. The promise of seeing God sounds, according to the doctrine of
approximation, like a mocking irony over human frailty: `it is as easy for a man to see God as
to be pure in heart; both are equally impossible.' In the spirit of Christianity, on the other
hand, Christ's words are so expressive and so distinctive that a certain artfulness is
positively required to get them duly misunderstood. By `the pure in heart' one cannot here
think of the gradual, step-by-step improvement of natural men; for the natural man does not
grasp at all `the things that belong to the Spirit of God; for they are foolishness unto
him.' No, here one must think of something quite other, here one must think of that
mysterious supernatural metamorphosis, that miraculous, incomprehensible transformation which
`God the Holy Spirit' alone can bring about in his elect, when at his pleasure he creates in
man a pure heart --- in faith. The difference, then, is clearly enough expressed. If one
chooses to preach about the natural man's moral approximation to God, then one has thereby
recognized the unavoidable necessity of the development of reason, and at bottom abolished
Christianity. The natural heart may improve itself as much and as long as it will; it will
nevertheless never come to see God, and surely does not demand to either. If, on the other
hand, one chooses the miracle of regeneration as one's starting point; well, then one upholds
Christianity, but comes into a hard contradiction with reason. It is naturally left to each to
choose which he will; only this we must repeat most emphatically, that he who would defy
nature and reason with the miracle of regeneration must not expect any
% --- p. 274 ---
\opage{274} support from the pastors and theologians of the present age, but should rather
seek it from the mystics and monks of the past.''

``What has been said here of faith holds also of works: revealed religion and the
corresponding morality are related as cause and effect.''

``Christ, who preaches under the open sky, forbids his followers all use of the right of
retaliation; let us just see whether any of the much-favored preachers of the present age dare
do the same when they preach in church. It is a very rewarding task to have to soften the
hard commandment: `An eye for an eye, a tooth for a tooth.' Were the church ever so overfilled
with people, there would hardly be one to be found who in our humane times would take it amiss
of the pastor that he warmly enjoined conciliation: good Lord, we are all brothers, after all.
But now the positive commandment, the commandment by which the Law passes over into becoming
Gospel, the strangely mild-severe commandment: `if any man will go to law with thee, and take
away thy coat, let him have thy cloak also': what pliancy in language, what delicacy in
expression, what softness in the voice is not required in order rightly to apply this
commandment in a truly edifying discourse. Suppose the pastor, in the midst of the discourse,
came into the spirit; suppose that zealous spirit which was upon Christ when in his Sermon on
the Mount he burst out: `till heaven and earth pass away, not even the least letter or one
tittle of the Law shall pass away, till all be fulfilled'; suppose, we say, that this spirit
of zeal gained such power over the pastor that he could not gain power over Christ's words,
but in the holy ardor of self-forgetfulness began in full earnest to enjoin upon the
congregation the terrible commandment: `if any man will go to law with thee, and take away thy
coat, let him have thy cloak also': what unrest would not then arise among the `devout
listeners,' that
% --- p. 275 ---
\opage{275} %
is, if it could only occur to the devout `that the man, by God, really means it.' `What a
dangerous exaggeration, to let religion abolish the right of property! Every decent man must,
after all, feel at once that such fanaticism must naturally lead to --- pure communism.' ---
It is of little use now that a theologian meddles in the matter and tries to calm the tempers
by mumbling out an interpretation about the ideal. However the ideal may otherwise be laid
out, it must always in one way or another be reconcilable with the actual; but now every
sincere Christian knows, after all, `that pure communism is --- in open conflict with the
actual.' --- Need we cite still more examples?''

``In his Sermon on the Mount Christ further calls the meek and the peaceable blessed. To the
former he promises that they shall `inherit the earth'; of the latter he says that they shall
be called `the children of God.' --- How attractive, yes, how ingratiating can such gnomic
sayings not become when they are treated with taste by an eloquent preacher. In the first
promise, however, one hears at once that its fulfillment can hardly come about naturally. One
need not exactly see so dreadfully deep into the course of things to be able to tell oneself
that meekness is by no means the virtue by which nowadays one takes kingdoms and lands. The
pastors know it very well, if not exactly from their own experience, then at any rate from
the example of their predecessors. With all the greater confidence, then, can they commend
peaceableness. If only they would now act accordingly! we would in all sincerity bid them
welcome, we would with heartfelt joy hail them, as it is written: `how beautiful are the feet
of them that preach peace, that bring glad tidings.' But (however it may go with peace) one
remark we can for truth's sake in no way
% --- p. 276 ---
\opage{276} hold back: the peaceableness that Christ enjoins stands in open conflict with
natural drive and reason. It says expressly (Matth.\ 5, 39): `If any man give thee a blow on
the right cheek, turn to him the left also.' How, now, does one bring this command into
harmony with the natural drive of self-preservation? Imagine that an insolent person (there
are, unfortunately, still insolent people) flew up and gave a peaceable man a box on the ear
`on the right cheek': what would one then expect of the offended party? If he struck back,
one would surely find it quite natural; if, instead of returning the blow, he offered his
adversary his hand in reconciliation, one would surely praise his peaceableness and admire his
superiority; but if the peaceable man, on the other hand, without further ado turned the other
cheek to his adversary: how would the modern Christians (we presuppose that those present were
nothing but modern Christians) then judge such conduct? Some would perhaps take it for a
somewhat exaggerated joke; others for a whim of defiant humor; others would even suppose that
the man, dazed by the first blow, probably did it in a kind of delirium; but hardly anyone
would, in the moment of surprise, think of making a serious application of Christ's
commandment and calling the peaceable man `a child of God.' But is it not ridiculous to set up
an ideal that no one can realize without at once being misjudged, even by those who otherwise
in writing and speech prove abundantly that they do, after all, pay homage to the ideal.''

``To remove every shadow of arbitrary exaggeration, we shall, in conclusion, further point out
that the point of Christian virtue is by no means that natural and solid soundness of soul
which shows itself in genuine moral-rational energy of action, but much rather a certain
inexplicable sickness of soul that feeds on pain and
% --- p. 277 ---
\opage{277} lives on suffering. The more a true Christian can come to suffer `for
righteousness' sake,' the more he can be hated, reviled, cursed, and persecuted by his fellow
men, the more inwardly he feels his blessedness, the surer he is of the reward in heaven
(Matth.\ 5, 10---12). With such a view of life Christ's commandment of love (Matth.\ 5,
44): `Love your enemies, bless them that curse you, do good to them that hate you, and pray
for them which do you harm and persecute you,' agrees most exactly. If suffering is the
highest end on earth, then the Christian surely has every reason to love and bless his
enemies. If there were no one who would hate him, he could not be persecuted and reviled; if
he could not succeed in being persecuted and reviled, he would not be sure of coming to heaven
and becoming blessed. Whom, then, has he on earth to thank that he becomes blessed in heaven?
Properly speaking, his enemies. Yes, if unnature could ever become nature, then Christ's
commandments would be most natural; if unreason could only transform itself into reason, faith
would become the most reasonable thing of all. --- `There is, after all, no one in the present
age who understands how to undermine Christianity with so dazzling a shrewdness and so
calculating an ingenuity as theologians and
pastors.'\,'' ---

What, then, does the Free Thinker accomplish with all these objections? Should we become
afraid of him, that he will now go and make us believers unbelieving, and take the life out of
Christianity? No, there is no danger of that! On the other hand, if he is otherwise duly
equipped, he will surely be man enough to make an end of half-criticism and confound
mediocrity. An able Free Thinker is undeniably one of the most useful members of the literary
community: he benefits Christianity, and procures the world a relief. --- That he benefits
Christianity is easy to prove. For, as he
% --- p. 278 ---
\opage{278} %
protests on behalf of drive and reason against the validity of revelation, he contributes, in
spite of any \textit{Summus Theologus}, to setting the supernatural origin of the Gospel in
its proper light for us. That he at the same time procures the world a refreshment, of which
under certain circumstances it may surely stand in great need, is self-evident. What the
Believer worships as the supernatural, unbelief must, after all, abhor as something
unnatural? The world feels itself, to put it plainly, encumbered by Christianity, and needs
air. When the ecclesiastical edifice of doctrine rises higher, and makes itself so broad that
it either shuts in or at any rate overshadows the other dwellings of the spirit, when theology
is just about to wall up the last opening, and it begins to get stuffy in the congregation:
then the Free Thinker starts up, knocks out a windowpane, and lets in criticism's sharp but
nevertheless refreshing draft. The more the Free Thinker, moreover, by his wanton attacks,
sarcastic turns, and teasing exaggerations, is able to shake the theologians' school-edifice
and put the Church's pastors in a comic light, the greater the applause he naturally reaps
from applauding unbelief. It is really not so much the weight of his reasons as the boldness
of his contradictions and the daring of his assertions that carries away the beau monde; for
the thing is that the world needs air. Let the Free Thinker by all means get air, then; let
him just go off, followed by the chosen host of enlightenment: it does not, after all, harm the
health of faith that a fresher breeze passes through the Church. We then take our place again
among the Believers, and leave it to the religious mood whether we might perhaps find, if not
the keynote, then at least the fundamental thought in the evangelical discourse.

``We cannot preach like Christ; oh no! we are to preach Christ. But that we truly preach
Christ shall be known by this, that the wisdom in our speech turns into
% --- p. 279 ---
\opage{279} foolishness for everyone who will not confess our Lord and Savior. Should the
wise of the world come together and want to explain to us that saving truth cannot possibly be
in Christ alone, but that there must necessarily be a purer, a more universal truth for the
comfort and salvation of all: we would nevertheless have to stand before these wise ones as
ignorant fools, as children under age; we would then have to answer simply and briefly that
Jesus Christ alone is the author and finisher of our faith. And if the wise then spoke up
again, and wanted to take us to task and test our reasons, \emph{why} we have adopted
precisely this faith and no other, what it is that has guided our choice: then we would, after
all, have to be to these wise ones as ignorant fools, as children under age, and answer simply
and briefly: we believe \emph{because} we believe in Christ, we believe \emph{because} --- we
believe; here we stand and can do no other. When the wise of the world hear this answer, they
will surely not tempt us to the utmost, but smile in farewell, and rejoice `that we can still
remember our conviction so well, although we have surely forgotten our reasons; then they will
surely congratulate us that we truly belong to Christ and worthily bear our name: the poor in
spirit.' --- And we will not blush at this salute of honor, as though it betrayed a mockery;
we will not be angry with the wise, but bow deeply and greet them kindly in return; for if
they have already gathered so many reasons that, on account of the superabundance of their
reasons, they cannot find any whole conviction, then they must, after all, be themselves ---
the rich in spirit. --- But, lest a doubting soul be confused by this greeting, and dispute
with itself whether it were best to stay with us or to go off at once with the wise, this
discourse shall venture an attempt, that it might yet explain to us how it can come about that
a man who goes furthest in the acquisition of spiritual riches can end by feeling himself ---
spiritually poor. --- If human
% --- p. 280 ---
\opage{280} wisdom is the highest goal, then we must say it without circumlocution to everyone
who is so fortunate as to be able to count himself among the gifted: turn away from us! we can
give you no help; you have, after all, now heard from the mouth of the wise that we are poor
in spirit. Not that we would send you away; oh no, we would rather keep you with us, and
gladly share our poverty with you. But if it is the wisdom of thoughts that is to save your
soul, then just take your leave, the sooner the better, for alas! you truly have no time to
waste. The wisdom of thoughts will lead you by the ways of doubts: there are many ways, if all
the ways are to be tried; yes, believe us, there are many doubts, if all the doubts are to be
tested. So we part: God knows whether we shall meet again. We now turn toward Galilee, your
face is turned toward Hellas. Our way is level and short, we found in faith a shortcut to
Cana; your way is both long and narrow: you must now wander by day, by night, in the dream of
your thoughts, over hill and dale, until you find the city of the wise. We rejoice over a
little, your spirit desires the great. --- You know that poor wedding feast: it was begun with
joy and merriment. Alas, the poor man's joy lasts but a short while! The will was so heartily
good, but the means only very slight. They had not reckoned it closely enough; one guest too
many had been invited. \emph{They have no wine.} But --- he is, after all, a high-born guest,
he speaks a word of authority: see, the servants draw the good wine, want is turned into
abundance: it is poured for poor guests. It is still to this day an equally singular riddle:
perhaps a bridegroom can grasp what no governor of the feast comprehends. --- Yet, we know it,
after all: it is too little a thing for you to ponder this riddle. You demand another and
richer feast: your spirit desires the great. So go, then, in peace. Far off on your way lies a
splendid city. Wisdom lifts her shield high above the high castle. Here the gods of taste
% --- p. 281 ---
\opage{281} hold their high feast. At the rich banquet sit the wise of antiquity: they weave
for themselves --- an immortal wreath. In the humor of the mood, in the wit of the speeches,
in the clarity of the thoughts, a spring of wine bubbles; perhaps it can gladden and refresh
him who thirsts for the wisdom of the world. So go, then, in peace! our best wishes accompany
you on the way. But when you have then come to the longed-for place; when you open the door to
the banquet hall; when the sight surprises you, and your heart beats with admiration: then take
good care that you also find the wisest among the wise. He too is a strangely odd guest: you
know him easily, he does not resemble the others at all. If you hear him, his voice captivates
your ear: the thought is an intoxicating wine, but the speech resembles pure water. Keep close
to him; for he knows the riddle of wisdom. And if you will now rightly attend to his speech,
and to the sting that is in his words, and to the smile that plays about his lip, then he
tells you in the language of wit what it is that the knower knows, then he tells you a veiled
word; but the word reads --- nothing. --- If you can now sate your soul with this secret, and
refresh your heart with this hungry nothing; if you can learn the fine but difficult art of
transforming want into abundance and using hunger as satiety, why then you may boast of being
a spiritually rich man, even if you know within yourself that your riches are want and
poverty. But if you cannot do this, you too must confess that this wisdom of hunger is, after
all, like an ingenious folly: when the feast fades, when the intoxication ends, and the soul
secretly languishes: then think of Cana's humble house, think of Mary's words: they have no
wine; yes, think of us and of our humble wedding feast and of our poor wisdom. Turn around,
and cast your glance on Christ; turn around, and hear his voice from the mount: {\large
Blessed are the poor in spirit, for theirs is the kingdom of heaven.}''

% --- p. 282 ---
\opage{282} ``We cannot preach like Christ; we are to preach Christ. But that we truly preach
Christ shall be known by this, that our preaching always sounds like a dark saying, a
bewildering word, for everyone who will not believe the wonder of regeneration. How often is
it not said that faith must be a matter of the heart. This is understood by those who love
Christ with their whole heart, and it is conceded by those who hate him with their whole
heart. That is why the strife of faith and unbelief is so hard and terrible a strife; for
there heart strives against heart, and the true man dwells in the heart. But harder than the
strife that is thus striven between man and man is the divided heart's own strife with
itself. --- There was perhaps one who tried his hand at love of the world, but had to turn
away because his heart was sickened by its vanity; there was perhaps one who lived high with
the wise of the world, until the many thoughts drained the springs of life from him, and he
felt himself --- poor in spirit. If now such a one, like a traveler, turns to us and asks the
way to the kingdom of heaven, then we must not lead him wrong, as though he could walk straight
in at once if only he turns his back on the world; but we must in all sincerity set before him
what we daily set before ourselves, --- the Redeemer's word: {\large Except a man be born
again, he cannot see the kingdom of God.} What does it help to want to leave the world, when
one cannot after all reach up to God? what does it help that a man doubts all the wisdom of
the world, when he cannot after all recognize God's wisdom in Jesus Christ? If you would in
spirit and truth recognize God's wisdom, then you must believe in the Lord Christ; if you would
believe in Christ, then you must love him with your whole heart; if you would love him with
your whole heart, then you must also be pure in heart, for
% --- p. 283 ---
\opage{283} %
only the pure can love the Pure: \emph{Blessed are the pure in heart, for they shall see
God.}''

``When the stranger who asks the way to the kingdom of heaven now sees himself halted by this
testimony, he will surely by no means at once find himself uplifted in the involuntary joy of
expectation; but, since he is an experienced man, he will more likely feel the weight of care,
and begin to take himself to task. And when he then begins to take himself to task, and the
practiced understanding discovers the difficulties, it may easily happen that he sinks --- into
doubt and despondency; and when the doubting understanding then rises in proportion as he
himself sinks, it may easily happen that the despondent one at last takes us, too, to task,
`whether we, then, are already in the kingdom of heaven; whether our faith, then, knows
nothing at all of doubt; whether we, then, love Christ above all things; whether we are
perhaps the pure in heart, the elect, who can see God and live?' --- We are not angered at the
stranger that he thus stands here and doubts and asks: he is, after all, asking the way to the
kingdom of heaven; we do not reproach him that he is so despondent. Alas, despondency comes of
weariness. This pilgrim is weary and tired. He knocked at so many a sanctuary, and went in,
and found it --- godforsaken; he searched so far about in the world; but never did he find what
he sought; now he sits here with a troubled mind, asking himself whether with us too he should seek
without finding.''

``So will we speak as to the stranger, while in all sincerity we speak to ourselves; so will
we exhort him, while we exhort ourselves; so will we bear witness to the wonder of faith and
of regeneration.''

``Our faith begins in need. --- `Then I shall surely soon come to faith,' you think, `for I
truly need a divine help.' --- Beware, though, lest you misunderstand the very first thing. It
is not with the need for
% --- p. 284 ---
\opage{284} God as with the need for daily bread. A merely natural need, were it the heart's
strongest need, can never lead to God. No one but God can draw our heart to God; but God is
above nature. It is only the world's confused talk when it is sometimes proclaimed that the
need for God is natural, like the need for food and clothing, or like the soul's natural need
to love another human being. Certainly our nature is a creation of God, and the drive toward
God to that extent also natural. But our nature is, after all, also fallen away from God, and
the natural drive ambiguous. When, then, the need for God awakens in the heart, it depends on
how it is apprehended. If it is felt as a merely natural drive, then it stands, according to
its origin, on an equal footing with the other drives. But the natural drives strive against
one another: when the one leads to God, the other leads away from God. It is therefore
doubtful which one is right; it is doubtful which will in the end become the strongest; it is
doubtful whether the drive toward God will not at last run out into a pious fancy.''

``What, then, gives this drive its first wonderful stamp as a supernatural drive, a need for
God that is brought about by God? --- Faith! This is the supernatural beginning of our life,
that need works by means of faith, and faith by means of need, because God's own Holy Spirit
works in them both. Consider now Mary, the believing soul, as you see her at the wedding feast
in Cana. Her attentive eye discovers a lack; she at once interprets it as a deeper need. The
guests, if they were made aware of the same need, could surely all explain how it had arisen
in an entirely natural way, and how it could in turn be remedied in an equally natural way, if
only one had the natural means. But in Mary the natural want worked a supernatural need, a
need to
% --- p. 285 ---
\opage{285} see God's glory. So her soul was simple in the simple expectation, so her heart
was pure through the pure faith, so she lifted up her eyes and looked at her Son with a
questioning, pleading glance: \emph{They have no wine.} --- The world calls this a peculiar
behavior. But you, who are finished with the world's glories, its festive customs and
conventions, just go and try the Gospel. Learn to understand Mary's need and Mary's faith;
turn in your prayer to Christ, and you shall find that it purifies your heart and clarifies
your need, that the Lord will for your sake too make --- a supernatural beginning.''

``Our faith begins with the wonder of regeneration. So it stands in our catechism: `I believe
that I cannot by my own strength or reason believe in Christ or come to Christ, my Lord, but
that it is the work of the Holy Spirit!' --- `But,' you say, `if I can do nothing by my own
power, but must await everything from the work of the Spirit, when will the wonder then
happen, how long must I then wait?' --- How long you must wait! An idle question! In the
washing of regeneration in the name of the Triune God you must offer your Lord everything,
your natural understanding, your natural strength, and give yourself over to his grace. For
this is the wonderful beginning, that the sacrament works through faith, and faith through the
sacrament, because God's Holy Spirit works in both. --- `But,' you say, `how long must I then
wait before this working of grace proves its power in inner experience?' --- How long you must
wait! --- Until the Lord's hour comes. Hear what may befall you (it has already befallen so
many): when you strive with flesh and blood, when your zeal grows as you vainly summon all your
strength, perhaps you will then one day press in upon Christ and cry to him that he might
transform you as by a wonder, that he might at once
% --- p. 286 ---
\opage{286} stretch out his hand and speak the word, and hasten with help before it is too
late; perhaps it will then also befall you that you do not perceive his holy presence at all,
not the word of power that should raise up your courage, not the creating breath that should
renew your heart. It is as if you were turned away with your importunate prayers; it is as if
he turned his face away; it is as if you heard the stern voice: \emph{What have I to do with
thee? Mine hour is not yet come.} Do not then despair, do not be disheartened, and, we beg you
for heaven's sake, never think of calling your God to account; never think of being angry, as
though you were the offended one. Look at Mary, the believing soul! She too had to hear the
stern word: \emph{Woman, what have I to do with thee? Mine hour is not yet come.} But she was
not angered at her Son because he rebuffed her reminder. She believed nonetheless that he would
hear her prayer. She felt, precisely in the hard word, that his hour would soon come; see, that
is why she held herself ready and placed the servants at his command. Go, then, and do
likewise. If you stumble at the hard word, then understand that his hour will soon come; if it
is as though he would send you away, then keep your mind ready in expectation, place your
serving thoughts at his command, and you shall surely find that for you too he will make a new,
a wonderful beginning.''

``And when the transformation then happens (it happens before you yourself know of it), then
your heart is strangely moved, then your soul is purified in the pure faith, then the Gospel
tastes like the good wine, then you see your Savior's glory, then you have an inner
experience. --- Come again then, and hear our humble discourse; yes, come
% --- p. 287 ---
\opage{287} %
again, and hear His voice from the mount: {\large Blessed are the pure in heart, for they
shall see God}.''

``We cannot preach like Christ; but always we are to preach Christ. But that we truly preach
Christ shall be known by this, that we confidently obey his law and do not fear the world. ---
Not fear the world! But the world is fearful, after all: it hates Christ; its hatred of us is
to be a testimony that we love Christ from the heart. It can never come about naturally that
we, the weak, who can do nothing, should now be able to awaken the world's hatred. It can
never come about naturally that the hatred of men and the persecution of the world can be
transformed in us into a pledge that we follow after Christ. Let it then be our rich comfort
that what is impossible for the powers of nature is possible for God's free grace; let it then
be our rich comfort that, though we must through much tribulation enter into the kingdom of
heaven, we are also to hear his voice who leads the hosts of the suffering: {\large Rejoice and
be exceeding glad, for great shall be your reward in heaven; for so persecuted they the
prophets which were before you}.''

``Not fear the world! But the world is fearful, after all, when it rises up within us, when it
stirs in drives and passions. It can never come about naturally that we, who would much rather
have only --- \emph{an eye for an eye and a tooth for a tooth}, should love our enemies, and do
good to them that hate us, and --- overcome evil with good. Let it then be our rich comfort,
when we obey Christ's law, that we can do nothing of ourselves; let it be our rich comfort that
He who stilled the storm on the raging sea can still the storm in the raging mind. ---
Therefore, when someone gives you a blow on your right cheek, and it surges in your blood that
you might yet slake your revenge, and pay
% --- p. 288 ---
\opage{288} him evil for evil, then remember Christ's word, then hear his
voice from the mount: {\large Blessed are the meek, for they shall inherit the earth.}
--- When you hear this voice, then it becomes still in you; and when it becomes
still, dead calm in you, how strangely does it not come about: you become so
inwardly strong, yes, far stronger than your adversary, for you become, after all,
stronger than yourself. And when you have first become so strong that you are
stronger than yourself, can you not then freely let the adversary strike --- even
your left cheek?''

``Not fear the world! No, faith knows that the world, alas, is mortally sick,
and does not fear the sick one. These proud plans on which the millions build; they
are, after all, not a work of health, no, they are the fevered man's strong dreams;
and these outbursts of powers that follow blow upon blow, alas, they come not of the
heart's health, they come of the nerves' weakness: it is the sufferer's violent
convulsion. Faith knows that the race of men is a suffering race, faith knows that
all souls are sick. The difference is only this, that some sick ones are proud, and
therefore with great care and much diligence hide their sickness, and count
themselves among the healthy, so that they shall not seem to have need of the
physician. Other sick ones are faint-hearted. These do indeed feel their inward
weakness, but still have not the courage to consult the great Physician, because he
takes it so strictly and speaks so harshly:
{\large If thy right eye offend thee, pluck it out, and cast it from thee; for it is
profitable for thee that one of thy members perish, and not that thy whole body
should be cast into hell.}
This now seems to them a hard saying; therefore they do not dare to submit to the
strict Physician's treatment, but seek counsel from other, more lenient, but
unfortunately incompetent physicians, and spend all their substance without being
healed. But all the sick,
% --- p. 289 ---
\opage{289} who in trust and faith take refuge with the great Physician still find him, as before,
unchanged, the same. Still, as before, he stretches out his hand; still, as before,
he speaks the word, to the one: \emph{I will! be thou clean!} to the other:
\emph{I will come!} \emph{but} to the soul that suffers most of all, and is in
unspeakable pain, and, long given up by all other physicians, now seeks him in
the press, only to touch his garment, his gentle voice sounds:
``\emph{Arise, be of good cheer, thy faith hath saved thee!}''

``We cannot preach like Christ; but always we shall preach Christ.''

--- Who, then, is this One, who speaks the word and stretches out his hand? \emph{The great
Physician!} --- ``But --- with little practice.''
\emph{The true Prophet!} --- ``But --- \textit{post eventum}.''
\emph{The Savior of the world!} --- ``But --- without a world.'' \emph{God's only-begotten Son!}
--- ``Born of --- romantic longing.''

\begingroup\large
Think not that I am come to bring peace on earth. I am not
come to bring peace, but a sword. For I am come to set a
man at variance against his father, and the daughter against her mother, and the son's wife against her
husband's mother. And a man's household shall be his enemies (Matth.\ 10, 34---36).
\par\endgroup

Christ is a sign that is spoken against.

\subsection*{\textit{IV.}\\[4pt] Christ Reveals the Thoughts of the Hearts.}
\addcontentsline{toc}{subsection}{IV. Christ Reveals the Thoughts of the Hearts}

In rejecting the Fore-Gospel with its many-sided hints
and allusions, the reasoning criticism must of course go its own ways in order to investigate
how Jesus came upon the idea at all that he, and no one else,
should be the Messiah; and also how, from the beginning, he understood
% --- p. 290 ---
\opage{290} this his Messianic calling. ``What, then, was Jesus' real intention? What plan
had he laid when he came forward? Did he hold to the same plan from beginning
to end, or was his provisional plan perhaps in large part reshaped and
altered, as circumstances and the obstacles he met made it
necessary?'' --- That the answering of such important, such serious, such highly
interesting and difficult questions must naturally cause criticism a good deal of
trouble and much racking of brains, anyone who has any sense for scientific
thoroughness will readily see. One is still far from being clear about
this matter: opinions are divided. Some (especially among the older critics) think that
Jesus from the beginning intended to establish a Jewish
Messianic kingdom; they appeal in particular to his ``political'' entry into
Jerusalem (Matth.\ 21, 1---11). Only after this attempt had foundered on
the pharisaic obstinacy of the chief priests and the fickleness of the crowd was the plan
altered to this effect: that his kingdom should not be of this world at all, but should
be spread as a kingdom of spirit and of truth among all the peoples of the earth.
Now although, as was said, many learned and acute men have shared the opinion that
some change or other must have taken place in Jesus' first plan, they
have nevertheless not been able to agree on the manner in which the change must be supposed to have
taken place. Opinions are divided: some, namely, think that it is Jesus himself who
altered the plan when he saw that he would not succeed in
restoring the Davidic-theocratic throne; others, on the contrary, think that the
ingenious reinterpretation by which the Jewish-political Messianic kingdom is transformed into a
spiritual kingdom of God is owed to the apostles, who, when Jesus had to succumb, had
had conduct enough to rescue a cause that was as good as lost.
% --- p. 291 ---

\opage{291} In the modern School one is chiefly occupied with investigating how
Jesus' consciousness was really constructed. The theologian Hase makes the very
important remark ``that his Messianic consciousness cannot have been clear when he
lay in the manger.'' ``It may be,'' it says, ``that the first germ of this
consciousness lay in the mother's hope; yet by laying too much weight
on this one would encroach upon Jesus' independence.'' ``According to a rabbinic tradition,
every Israelite, and in particular every son of David, was to wish to be the Messiah; it is the
custom of great souls to turn the wish that governs their life into deed.''
``Both,'' Hase continues, ``can be conceived: either that the consciousness of his
Messianic vocation developed immediately, side by side with his
self-consciousness and his recognition of the Messianic hope of the people, or that it rather
issued, as a definite resolve, out of inner doubt and hopes.'' ---
The anti-theologian Bruno Bauer likewise speaks much, both of Christ's
God-consciousness and of the apostles' religious consciousness and the whole age's
world-consciousness. All that is now lacking is that a new
critical-philosophical genius should arise who could speculate out thoroughly how at
bottom all these consciousnesses hang together. If one could only find out
how Jesus' God-consciousness became one with his self-consciousness, and this his
self-consciousness in turn with his higher world-consciousness: then the knot would be untied, and the
secret out; then one would at last, for once, get to know what lay in
Jesus' original plan, what it is that he really intended.

The Believer will gladly leave it to modern criticism to juggle the many
opinions and arrange the many consciousnesses; he contents himself meanwhile with
Simeon's direction in the Fore-Gospel: \emph{This one is set for a sign which is spoken
% --- p. 292 ---
\opage{292} against,} {\large that the thoughts of many hearts may be revealed.}
Here, then, we have a hint which, at least in one principal direction, puts us on the
right track.

The Savior in the Gospel does not stand alone. Only once in a while does he withdraw from
the multitude, to rest in prayer to the heavenly Father. Some hours
he spends in a narrower circle of pious souls devoted to him. At times
he walks with his closest disciples through the cornfields of Galilee; but most
often we find him in the midst of the manifold throng of the people, where he then works with all
his repelling and attracting power: the crowd is startled, the disciples listen,
the Pharisees lie in wait, and the eyes of all watch him when he speaks the word and stretches out
his hand. --- If he is now really the true light that lighteth every man, then
every human figure that moves in his vicinity must also become
transparent in relation to him, and show itself to us in the right light. Here
are souls enough whose inner being he can move; here are hearts enough whose thoughts he
can reveal. Yes, he reveals the thoughts of the hearts, not only of those who are
with him, but also of those who are against him; he lets the light rise over the evil and
the good. --- First, then, we turn to his disciples, to see how the thoughts of their
hearts are revealed.

As regards the calling of the disciples, and in particular of the twelve apostles, modern
criticism has from its standpoint raised a number of objections which the doctrine of mediocrity
will find troublesome enough to refute. Strauß brings out, among
other things, the following difficulty. ``According to the concordant
narrative of the first two Gospels (Matth.\ 4, 18---22. Marc.\ 1, 16---20), Jesus, as he walked by
the Galilean sea, called the two brothers, Peter and Andrew, and immediately afterward
James and John, away from their fishing nets, to follow him. The fourth
% --- p. 293 ---
\opage{293} Gospel too relates right at the beginning (Joh.\ 1, 35---52) how the first
disciples attached themselves to Jesus; among these are likewise Peter and
Andrew, and probably also John, since Andrew's unnamed companion is
usually taken to be him. Now since these narratives concern different
events, the question arises which is the earlier and which the
later. The interpreters agree in placing the Johannine calling as the
earlier. The fourth Gospel, it is said, relates only Jesus' first
acquaintance with these men, after which they did not at once become his
constant companions. Only on the occasion preserved by the Synoptics did
Jesus call them to be his constant companions, real disciples.'' --- Upon
this preparation there now follows the critical blow. ``If the words \emph{follow me}
in the Johannine narrative (Joh.\ 1, 44) are to be taken as a merely provisional
summons, then the words \emph{follow after me} (Matth.\ 4, 19. Mark.\ 1, 17)
must necessarily also be understood of a merely provisional invitation. But this
interpretation is impossible. How else could Peter afterward (Matth.\ 19, 27), in the name of his
fellow disciples, have reminded Jesus of such a merely provisional calling with the
words: \emph{behold, we have forsaken all things, and followed thee; what shall we
have therefore?}''

``As a special advantage,'' Strauß continues, ``of this arrangement of the two narratives,
the rationalist interpreters point out that only in this way does it become
comprehensible, what one would otherwise have to be astonished at in the highest degree, namely,
how Jesus, merely as if in passing, at the first glance, could choose four
fishermen as apostles, and among them at once hit upon the two most distinguished apostles;
and also how the four busy men could, at the enigmatic call of a man not more closely
known to them, at once leave their trade and set about
accompanying him. For by a
% --- p. 294 ---
\opage{294} comparison with the fourth Gospel one sees that Jesus had long
beforehand come to know these men, and likewise shown himself to them in his
excellence; from which both the felicity of his choice and their
readiness to follow him can now be explained. But precisely this
apparent advantage is what utterly breaks the staff over the arrangement described
of the two narratives. For nothing can be more decidedly against the intention of the first two
Gospels than the presupposition that a relation had already existed between
Jesus and both of the pairs of brothers called. For if the narrative in
both lays such great weight on their leaving their nets and attaching themselves to Jesus, then
in the narrator's opinion this must have been something extraordinary, which it
would not have been if the men had already been in Jesus' following before. And likewise with
regard to Jesus, the point of the narrative lies in this, that with prophetic spirit
he at once, at the first glance, found the right men; that he, according to Joh.\ 2, 25,
\emph{needed not that any should testify of man, for he himself
knew what was in man;} whereby he fulfilled a demand which the Jews
made of their Messiah.''

These difficulties can easily multiply when the negative spirit lays its
blessing on the critical observations. Well, that does no harm either. Let
criticism just drag mediocrity about a little: the faith of the Gospel meanwhile
remains standing on the fundamental word that the \emph{thoughts of many hearts shall be
revealed.} Thus the modern consciousness is, after all, revealed at once as soon as it
sets about settling the relation of the disciples to their Lord and Master.
Strauß makes no secret of what it is that arouses his offense. ``A
seeing-through of man at the very first glance, such as Jesus showed here,
far surpassed everything possible for the happiest and most
% --- p. 295 ---
\opage{295} practiced knowledge of men according to the laws of nature. A man's inner being can only with
certainty be known from a series of utterances and actions: a gift of penetrating
into the dispositions of others without their help already verges on the visionary,
and so onto a territory for which the expression the rabbis used for this
gift in the Messiah, \textit{orando judicare}, is by no means too monstrous. Hardly
less improbable, on the disciples' side, is the unhesitating following,
inasmuch as Jesus was not at that time as famous in Galilee as later on, and one
would almost have to infer from this readiness a power in Jesus' voice and will,
independent of preparation and reasons, to penetrate immediately into
people's hearts; which would be yet another magical trait, besides the visionary one mentioned above,
to complete the incredibility of the narrative's historical
character.'' --- Now it might indeed seem as if negative criticism made
itself guilty of a questionable contradiction, in so summarily denying
Christ the gift of being able to see through a man with his glance, while it
itself arrogates to itself the ability to see through with \emph{its} glance both
Christ himself and the Galilean men who met him on the way; but that is, after all,
another matter: the modern consciousness has, you see, seen through --- ``the laws of
nature.''

Upon the believing heart, however, it makes a wondrously deep impression to see the Son of God
in so immediate a contact with men who are, like the rest of us, frail
beings of flesh and blood: how alike, and yet how unlike; how near, and yet ---
what a distance! Criticism thinks that the various callings contradict and exclude
one another; but what holds here of these callings in public life is what has
already been said of the many manifestations in the Fore-Gospel: the understanding is
conditioned by the psychological presupposition. --- Criticism thinks that the number of the twelve
apostles and of the
% --- p. 296 ---
\opage{296} seventy disciples (Luc.\ 10, 1) is highly suspicious, because this
number is at the same time symbolic. On this occasion we must take up yet another Straussian
remark: ``For just as the twelve apostles, with regard to the twelve tribes of Israel,
indicated Jesus' destination for the Jewish people, so, it is
said, the seventy, or, as some authorities have it, the seventy-two
disciples, were representatives of the 70 or 72 peoples which, like
so many tongues, according to Jewish and early Christian view were to be found on
earth, and thus pointed to the universal destination of Jesus and his kingdom.
For the Jewish nation in and for itself, too, this number had a significance as a
holy number: 70 elders Moses chose as helpers (4 Moseb.\ 11, 16. 25);
the Sanhedrin had 70 members, and the Old Testament as many Greek
translators.'' --- Should anything be able to lead us to go along with this criticism,
it would have to be the question that now immediately follows: ``Had Jesus, who
was set in the very midst of the press of circumstances, nothing more important to do than to
gather together all possible numbers, and according to them surround himself with various circles of
disciples?'' In the conclusion, however, it runs thus: ``So much seems
indeed to be clear from Ap.G.\ 1, 21, that Jesus, besides the Twelve, had still
other constant companions: but that these formed precisely a corps of
Seventy, or that so many of them were chosen, does not seem
to be duly proved.'' Since the critic himself yields ground in this conclusion, as
if he were not after all so entirely sure of his case, we too will in this point
prefer to hold to the Gospel and content ourselves with the symbolic number.

Criticism now thinks quite surely that an apostle's calling is an apostle's calling, roughly
in the same sense as when one says: a word is a word, a man a man, a
contract a contract.
% --- p. 297 ---
\opage{297} But faith recognizes that there is a great difference between a word and --- a word,
between a man and --- a man, yes, a great difference between letting oneself be called
to be an apostle and entering into a contract. How far the faith of the Gospel is right in
this assumption will appear when we consider more closely the personal relation
between the Lord and his disciples.

It is a remarkable struggle that faith and doubt, understanding and misunderstanding,
devotion and self-love must have waged in the hearts of these disciples.
The preference which the Lord himself (Joh.\ 1, 43; Matth.\ 16, 17---19) grants Peter
seems to have aroused a violent jealousy in the others. They come to Jesus
(Matth.\ 18, 1) and ask: \emph{Who is the greatest in the kingdom of heaven?} This
question admits of a double interpretation. On the one hand it bears witness, of course, to their inward
striving after perfection. In true understanding, in living faith, in loving
self-sacrifice, the one will not stand behind the other. But the question
at the same time betrays their natural self-love, in that the one would not gladly
give place to the other; but each thinks to himself that he might possibly
in the end become the greatest. The incident with the rich young man, who
could not leave all his goods to follow Jesus (Matth.\ 19, 16---22),
prompts Peter to the question: {\large Behold, we have forsaken all things,
and followed thee; what shall we have therefore? But Jesus said unto them: Verily I say
unto you, that ye which have followed me, in the regeneration, when the Son of man shall sit
on the throne of his glory, ye also shall sit upon twelve thrones, judging the twelve
tribes of Israel (Matth.\ 19, 28).} That is a word that falls into the disciples' ears.
``What a reward, what glory, when the kingdom is one day established, to have such
places of honor in store!'' But what Jesus promises them in
% --- p. 298 ---
\opage{298} the language of eternity and blessedness they now translate to themselves in the sense of temporality and
worldly happiness: thus arises the dispute about the first place of honor.
Christ has already several times foretold to them what lay ahead.
Hardly has he (Matth.\ 17, 12) come down with the three chosen ones from the mount of the Transfiguration
before he reminds these three that the \emph{Son of man also must suffer.} Afterward
(Matth.\ 17, 22---23) it is said expressly: While they abode in Galilee,
Jesus said unto them: \emph{The Son of man shall be betrayed into the hands of men; and
they shall kill him, and the third day he shall be raised again.} And they
(the disciples) \emph{were exceeding sorry.} The same reminder is repeated once
more (Matth.\ 20, 18---19): \emph{Behold, we go up to Jerusalem, and
the Son of man shall be betrayed unto the chief priests and scribes; and they
shall condemn him to death, and shall deliver him to the Gentiles to mock,
and to scourge, and to crucify him.}

These repeated predictions would surely seem bound to displace the sensuous
conception of the Messianic kingdom. ``Nevertheless the disciples continue to hold
fast to their Jewish expectation: here, then, is a contradiction.'' Undeniably! But
this contradiction stems from the disciples' misunderstanding, and cannot at all be
used as proof against the credibility of the account. The misunderstanding is explained
naturally by the disciples' situation at the time. On the one hand they heard the Master's
foreboding sayings about the impending suffering and death; but on the other hand they
also heard, after all, his express promise of the resurrection and the glory. The first
of these predictions is canceled in the last. The very first portent is
a riddle to the disciples. Jesus' word to the Jews
% --- p. 299 ---
\opage{299} (Joh.\ 2, 19): \emph{destroy this temple, and in three days I will raise it up,}
they (v.\ 22) first came to understand when he was risen from the dead. How
Peter, James and John must have felt themselves affected by his supernatural
glory as they followed him down from the mountain where they had seen him with Moses and
Elias! --- As if to damp the all-too-strong awakening, there now follows the reminder
that it behooved the Son of man to suffer like the old prophets and the second Elias.
But the talk of Elias is, after all (Matth.\ 17, 12), a figurative speech: the blessed vision
displaces the painful impression again; what wonder, then, that the disciples take the
literal all too figuratively, and the prefigurative all too literally. It is
indeed said expressly (Matth.\ 17, 22) that the disciples were \emph{exceeding
sorry} when they heard the crushing prediction; but the sorrow vanishes
again, for wonder-work, after all, follows upon wonder-work (Matth.\ 17, 27. 19, 2):
``should not he who can do such works also be able to establish the kingdom?''
The provisional battle and strife which the Master designated by his crucifixion and
death would, in their opinion, soon be over. When the brief strife was
ended, then followed the resurrection, and with it the regeneration. And when then
the great hour of the regeneration came, yes, then it was a matter of not missing out on
the place of honor. The only one among them all who really comes to consider the danger
is Thomas the doubter.

What a difference between this Teacher and these disciples: the Master forewarns of his
crucifixion and death; the confident disciples look only to his resurrection and
their own glorification. Hardly has Christ (Matth.\ 20, 19) uttered the words
\emph{mock, scourge, crucify,} before we, as a surprising proof of the
misunderstanding, (v.\ 20---22) meet the following scene: {\large Then came to him the mother of Zebedee's
children with her sons, falling
% --- p. 300 ---
\opage{300} down before him, and desiring a certain thing of him. And he said unto her: What wilt thou? She
saith unto him: Say that these my two sons may sit in thy kingdom, the one on
thy right hand, the other on thy left. But Jesus answered and said: Ye know
not what ye ask. Are ye able to drink of the cup that I shall drink of, and to be baptized with the
baptism that I am baptized with? They say unto him: We are able.}

Yes, they are able! of that there is now no doubt at all here: they are quite able to drain the cup; they
are able, if it must be so, to empty it in one draught; if only they can then get
--- the place of honor.

Had John and James not placed so unshakable a trust in Jesus' word about the
kingdom of heaven which he would establish, they would hardly have put forward their request; but
conversely, had they had the true recognition of this kingdom's nature and essence, they
would surely not have permitted themselves such a request either. They loved their
divine Teacher and Friend, and therefore they now so gladly wished for themselves a place at
his side; but their self-sacrificing love at the same time betrays a highly questionable
self-love, for they desired, after all, --- the place of honor. --- So slow to learn is
sensuous man when he is taught by God the secret of love and self-sacrifice.
And yet how high these twelve Galilean men tower above the
whole fickle crowd that gathers and scatters, comes and goes! Behold, they have
given up the temporal to find the eternal; they have forsaken all things to
follow a man who in this world had not himself where he might lay
his head. They have an ear for the spiritual word, they have an eye for the divine
glory. Blessedness is so near to them, and yet the actual is at odds with the
blessed. Sensuous sight discovers nothing, nothing at all: there is, to human
view,
% --- p. 301 ---
\opage{301} %
as yet no sign that any actual kingdom is to take shape; and yet it is to them as if
they already saw heaven open, saw all the supernatural signs of the regeneration; yes, as if
they saw the twelve thrones right before their eyes.

So, then, these twelve disciples walk with their divine Teacher and Master forward, upward, toward
the goal of the regeneration; forward, upward, toward the mountains of the spirit and of eternity. But finitude
weighs down, the foot stumbles; the Master must at every moment stretch out his hand and watch over them, that
none among them grow dizzy and fall from the perilous height.

Simon son of Jonas, the rock-man of faith, who is called Peter, comes first, of course. It is his
firm purpose to remain with his Lord always: for ``this One hath the words of eternal life; with him
it is good to be.'' --- The eager disciple feels within himself that in faith he can do
all things, and will venture all things, and does not know his own frailty. What the Master does, the
disciple will do likewise. So he tries it with the test of strength; but --- whence comes
the strength? --- Jesus walks on the sea; it is in the fourth watch of the night (Matth.\ 14, 25);
why then should Peter not also be able to walk on the sea? {\large Lord, if it be
thou, bid me come unto thee on the water.} Then a voice sounds through the darkness: {\large Come!}
And Peter knows the voice, and steps out of the ship, and walks on the water to come to
Jesus. {\large But when he saw the wind boisterous, he was afraid;} \emph{and beginning
to sink,} he cried, saying: {\large Lord, save me!}
\emph{And immediately Jesus stretched forth his hand, and said unto him:} {\large O thou of little faith, wherefore
didst thou doubt?} --- It was somewhat too early that the disciple wanted to be as his Master. ---
Peter loves his Lord; for this Lord is, after all, \emph{Christ himself, the Son of the living God}
(Matth.\ 16, 16). But hardly
% --- p. 302 ---
\opage{302} %
has he made his great confession, and received the promise of the keys of the kingdom of heaven, before Jesus
again makes known to his disciples that \emph{he must go unto Jerusalem, and suffer many things
of the elders and chief priests and scribes, and be killed.} ``It is dreadful
to hear. Should it go so badly with Christ at the last? It must surely be rashness.
The Master must certainly be overheated against the scribes. It is probably best that the
disciple who can speak for them set him right in time, before misfortune happens.'' And Peter took
him aside, and began to rebuke him, saying: {\large Lord, spare thyself; this shall not be
unto thee.} But it was, after all, a sinful rashness that the disciple would rebuke his
Master. And the Lord turned and said unto Peter: {\large Get thee behind me, Satan! thou art an
offense unto me; for thou savorest not the things that be of God, but those that be of men.} --- So Peter must
surely learn at last that a disciple is not above his master, and that the greatest in the kingdom of heaven
must be the servant of the others. It is humility, then, that it all comes down to: now
he understands it; now he will surely conduct himself accordingly. Thus we find him at the
last supper (Joh.\ 13, 3---10). {\large Lord, dost thou wash my feet?} --- \emph{What I
do thou knowest not now, but thou shalt understand hereafter.} --- {\large Thou shalt never
wash my feet.} --- \emph{If I wash thee not, thou hast no part with me.}
--- {\large Lord, not my feet only, but also my hands and my head.} --- How hard
true humility is to learn! Peter now believed quite surely that he had grasped
it; but this certainty --- came somewhat too early.

At last, in the quiet circle, the melancholy words of farewell sound from the Master's lips: {\large
Little children! yet a little while I am with you. Ye shall seek me, and as I said unto
the Jews:
% --- p. 303 ---
\opage{303} Whither I go, ye cannot yet come, so now I say to you} (Joh.\ 13, 33). Then
Peter's heart will break: {\large Lord, whither goest thou? --- Whither I go, thou canst not
follow me now; but thou shalt follow me afterwards.} Then the disciple surges up: {\large Lord, why
cannot I follow thee now?} ``I am ready for anything; see, I have a sword, I will
draw it against thy enemies; yes, I will lay down my life for thy sake.'' This he now promises from
the bottom of his heart, and yet it is rashness. He feels that he is able, and yet he is not yet
able; he relies on his own strength, and yet has not yet tested himself; he is defiantly confident in his
self-sacrificing devotion, and heeds not the Master's urgent warning: \emph{Wilt thou lay down
thy life for my sake? Verily, verily, I say unto thee:} {\large this day, even in this night, before
the cock crow twice, thou shalt deny me thrice} (Mark.\ 14, 30).

Beside Peter we find in this chosen circle the two sons of Zebedee, James and
John. The Master called them \emph{the sons of thunder} (Marc.\ 3, 17). This name points
to strong passions, and the Gospel has also preserved a couple of traits that remind us of
their vehement temper (cf.\ Mark.\ 9, 38. Luk.\ 9, 49 with Luk.\ 9, 51---56).

And it came to pass, when the days were fulfilled that he should be received up, he stedfastly set
his face to go to Jerusalem. And he sent messengers before him, and they went, and entered into
one of the towns of the Samaritans, to make ready lodging for him. And they did not receive him, because his
face was turned toward Jerusalem. But when his disciples, John and James, saw this, they
said: {\large Lord, wilt thou that we command fire to come down from heaven, and consume them,
even as Elias did?} \emph{But he
% --- p. 304 ---
\opage{304} turned, and rebuked them, and said:} {\large Ye know not what manner of spirit ye are of.}

This trait might now seem to stand in open contradiction with the ecclesiastical picture of
the Evangelist John as the apostle of love and inwardness. But this contradiction is
more apparent than real. True inwardness comes into being precisely as the
feelings that break forth and the strong, outward-striving passions are bent back into
the discipline of self-mastery and the blessed calm of meekness. He who stilled the storm on
the sea of Gennesaret was surely able also to tame the lightning in the breast of the sons of thunder, and to command a
dead calm. It is above all the immediate, soulful element in devotion to the Master's own
person by which the apostle John is distinguished from Peter. Whereas the latter (Philochristus)
apprehends Christ chiefly as the almighty Lord, John (Philojesus), on the contrary, apprehends
Jesus as his divine friend. Therefore his anger is kindled against the unworthy Samaritans
who would not grant this friend a lodging; therefore he prayed that fire might fall down from
heaven and consume the recalcitrant: alas, he had not yet learned to be --- still.
Only afterward, when at the supper he rested his head on the Master's breast; afterward, when
on Golgotha he stood at Mary's side beneath the Reconciler's cross, yes, afterward the
son of thunder learned to be still.

Some way behind the three chosen ones, Peter, James and John, the other disciples now follow. The
simply anxious Philip (Joh.\ 6, 7. 14, 8.) would so gladly see the Father; and
Thomas the twin, despite all his exertions, still cannot work out where the way leads
(Joh.\ 11, 16). If the interesting, that intermediate concept, could be applied at all
to the sacred history, we would doubtless have reason to call Thomas an
interesting character. --- He is surely a calculating man of the
% --- p. 305 ---
\opage{305}%
understanding, this twin. At times he stands quite still and ponders; one does not
rightly know whether he will come along any further or not: one might almost fear that the Master
together with the others will go on without him. But he thinks it over after all, and then follows along for
faith's sake; for however he reckons, and however long he ponders, he still
always gets out that ``faith must be --- the best.'' It is so hard that it should always go
against sense and understanding; but --- Thomas follows along nonetheless, for ``faith must be the best.'' It
goes so narrowly when the way is steep and the prospect dark; but --- Thomas follows along nonetheless, for
``faith is, after all, the best.'' If only he could soon catch sight of his goal and get the proof in hand. Still,
that will surely come; provision will surely be made for him too, if only he holds out and remains
faithful to the end. Once he gets so far that he can see the print of the nails in the Risen One,
and thrust his hand into his side: \emph{yes, my Lord and my God!} then provision will be made
for him too, then he will get the proof in hand.

So, then, all these walk with their divine Teacher \dots{} yet there is one more to
name; but he always keeps in the shadow. This figure prefers to appear somewhat behind
Peter, opposite John. --- ``He has a defiant look, and clutches the purse in his hand:
what, then, does this man want in the circle of the intimates?'' He does not know himself; he is drawn along. It
goes through him, marrow and bone, like heat and frost, when the Lord looks upon his Twelve and
forewarns that there is one among them who is a devil (Joh.\ 6, 70). ``What, then, is this
man doing here in the brotherly fellowship of the chosen?'' He does not know himself; he is drawn along.
Light and darkness strive in him, around him: it passes thus in silence. Only once
does a word escape the compressed lips. --- It was in Bethany, six days before the
% --- p. 306 ---
\opage{306} %
last Passover. Lazarus, who had been raised from the dead, made a supper for Jesus and his
disciples. Martha served them. But Mary took a pound of costly ointment of spikenard, and anointed
Jesus' feet, and wiped his feet with her hair; and the house was filled with the odor of the ointment. Then
said Judas, Simon's son, Iscariot: \emph{Why was not this ointment sold for 300
pence, and given to the poor?} (Joh.\ 12, 5). --- It is no good sign when a
Judas with the purse takes up the cause of the poor so eagerly. --- This disciple too was to be raised
up into the heights of the kingdom of spirit. But when the hour of decision came, he grew dizzy; it went black before
his eyes: he gave his Master a kiss, and --- plunged down, --- how deep? No one knows
it; --- the abyss is dark.

Jesus' enemies too are revealed, as they draw near to strive --- against the light. It is especially
the sanctimony, vain pettiness and hierarchical arrogance of the Pharisees that are struck
and chastised by the word of truth. The image of these scribes is mirrored in the Savior's discourses, and
stands out in sharp lines in his parables. \emph{When thou doest thine alms, do not
sound a trumpet before thee, as the hypocrites do in the synagogues and in the streets, that they may
have glory of men.} --- \emph{And when thou prayest, thou shalt not be as the hypocrites; for they
love to stand and pray in the synagogues and in the corners of the streets, that they may be seen of
men.} --- \emph{And when ye fast, be not of a sad countenance as the hypocrites; for they
disfigure their faces, that men may see them fast; verily I say unto you:
they have their reward.} --- It is these Pharisees and scribes who transgress the commandment of God
for the sake of their traditions;
% --- p. 307 ---
\opage{307} %
behold, it is these Pharisees and scribes who sit in Moses' seat, and bind heavy burdens
and grievous to be borne, and lay them on men's shoulders, but will not themselves move them with
one of their fingers. It is these Pharisees and scribes who so gladly would sit uppermost at table
at feasts, and have the chief seats in the synagogues. It is these Pharisees and
scribes who so gladly would be greeted in the markets and called Rabbi, Rabbi! of
men: it is these Pharisees and scribes the thoughts of whose hearts are now
revealed. The parable of the Pharisee and the publican, however short and simple it may be, illuminates
the innermost corners of a pharisaic heart. That the portrait is a likeness, that every single trait is drawn
with the eye that never errs, the Pharisees themselves prove in their continued zeal against the
hated teacher of the people. No fact is so strong that it can move these hearts;
no man's healing so wonderful that it could for a single moment bring
these obstinate opponents to forget the sabbath law: the warmer love breathes,
the colder these hearts become; the clearer the light shines, the blinder these
eyes become.

Now since the scribal enemies do not at once dare to lay hands on Jesus, for fear of
the people, they are compelled to try the means of the word. Thus there unfolds a series of
disputes which, with regard to the opponents' way of thinking, is of thoroughgoing significance.
--- The law of Moses presents two principal sides. Insofar as it is apprehended in spirit and truth, it points
beyond itself and leads to the Gospel: thus was it preached by the Baptist,
thus is it acknowledged and fulfilled by Christ. But regarded as a foundation, as
preparatory revelation, it is at the same time burdened with such a one-sidedness, and contains
such components, that in its imperfection
% --- p. 308 ---
\opage{308} it can very easily take it into its head to set itself in the way of the perfect. This makes it
explicable that into the interpretation of the law of Moses a spirit can intrude which stands against
Christ. Pharisees and Sadducees formed two parties which, each in its own way, were possessed
by this spirit: the former, because they clung to the petty additions of tradition; the latter,
because it accorded so well with their irreligious morality to hold to the very first and
meagerest foundation. Insofar, now, as one can say that Pharisees and Sadducees were as if
possessed by this spirit of one-sidedness, they were unfree souls, and their strife with the Savior
rests in that respect on a sorrowful misunderstanding. But their blindness is at the same time a
selfish hardening against the truth that is clear as day, and therein lies their guilt. Therefore
Christ says to these Jews (Joh.\ 8, 37. 40):
\begingroup\large
I know that ye are Abraham's seed; but ye seek to kill me, because my word hath no
place in you \dots{} If ye were Abraham's children, ye would do the works of Abraham. But now ye seek
to kill me, a man that hath told you the truth, which I have heard of God:
this did not Abraham.
\par\endgroup

This general strife individualizes itself, and is made visible in a multitude of
separate scenes. In Matthew's Gospel (22, 15; cf.\ Mark.\ 12, 13. Luk.\ 20, 20) it is
related that the Pharisees went and took counsel how they might \emph{entangle him in
his talk.} It is certainly very cleverly devised by the Pharisees to want to entangle Christ in his talk.
The plan will surely succeed, it can never fail: ``this big-talking pseudo-prophet surely
does not know the Scriptures; he is a Nazarene who at bottom has learned nothing at all.'' It is
furthermore a proof of these scribes' cleverness that they have the disciples they send to
him accompanied by Herodians; for the Herodians are, after all, politicians, and
% --- p. 309 ---
\opage{309} %
it is a political question that they intend to put before him. That they put a political
question before him is, however, the cleverest of all; ``for this rabbi of the people doubtless knows
nothing of politics, since his appearance and his whole conduct are, after all, in such open conflict with
all the rules of political prudence.'' The question itself is, moreover, laid out with extraordinary
shrewdness; but, of course, it has also sprung from the shrewd counsel of the assembled Pharisees.
It is now only a matter of hitting upon a suitable introduction; then the whole
will succeed, then he will surely walk into the snare. ``The introduction must naturally be designed to make
him as secure and trusting as possible, so that he suspects no mischief and wriggles out of it.
How shall one best begin? One must begin thus: \emph{Master, we know that thou
art true, and teachest the way of God in truth, neither carest thou for any man; for thou regardest not
the person of men.}'' --- ``Yes, there is absolutely no doubt that he must feel as
trusting as he is pleasantly surprised when he hears this address, especially the last words: neither
carest thou for any man; for thou regardest not the person of men.'' The question itself can therefore
be put: \emph{Tell us therefore, what thinkest thou? Is it lawful to give tribute unto Caesar,
or not?} ``There now, let him think it over, whether he can find an answer. If he answers that it
is not lawful, then he can, of course, be informed against by the Herodians and condemned as a rebel against
Caesar. If he answers that it is lawful, then our pharisaic disciples will surely know how to
present him to the people as a traitor to the fatherland. If he does not answer at all: well, so
much the better, then he can simply be declared defeated.'' Yes, it can never
miss: he must be entangled in his talk. --- ``What does he say?'' {\large Ye hypocrites, why
tempt ye me?} ``Does it begin in that tone!'' --- ``Yes, so there is again someone who has
% --- p. 310 ---
\opage{310} %
been with him and given him a hint beforehand.'' It is no good, after all, for so many to be
in on one secret. ``But what can he want with the tribute money?'' {\large Whose is this image and
superscription?} ``A strange question!'' ``It might almost seem as if he
meant by it that all those who contented themselves with Caesar's coin were also obliged
to pay Caesar tribute.'' --- Hardly have the self-wise fools thought these thoughts in the
wickedness of their hearts before the word of truth strikes down among them as in a flash of lightning, and strikes their
clever design to the ground. Behold, the Master, who was to be caught in his talk, needs only to
cast a glance at the tribute money with Caesar's superscription and the stamped image in order at once
to stamp an absolutely binding truth with his image and his superscription; for as long as
there is religion and state in the world, Christ's commandment: {\large Render unto Caesar the things which are Caesar's,
and unto God the things that are God's,} will be the last, decisive watchword.

Since the attempt now turned out so unfortunately for the Pharisees, the Sadducees must see whether they
can perhaps do a little better. Yes, the Sadducees will surely know how to hit the point: they
stand, after all, high above the Pharisees in unprejudiced culture and critical knowledge. The task they
set before Jesus is even in the most perfect agreement with pure knowledge. One
would almost believe it had sprung out of the modern consciousness's own head. \emph{A
woman has in this world of the natural been married to seven men one after another: to which of
them shall she now belong in the resurrection?} --- Those learned in nature praise nature because by
small means it accomplishes great things; believers praise the Redeemer because in simple words he
utters great truths. The whole wisdom of the Sadducees consists in ridiculing
the folly of the Pharisees. The Pharisees clothed the doctrine of the
% --- p. 311 ---
\opage{311}%
resurrection in such gross representations that eternal life came to look like a straightforward
prolongation of the temporal. Christ now reminds them that the resurrection is a transfiguration
in which everything corruptible vanishes. The individuals will indeed rise (the woman herself
as well as her seven husbands); yet \emph{in the resurrection they neither marry,
nor are given in marriage;} but {\large they are as the angels of God in heaven.} This part of
the answer overthrows all Sadducean objections: now follows the real confirmation of
the hope of the resurrection. --- The mildness with which Christ takes up the Sadducees' question
forms, moreover, a remarkable contrast to the severity with which he received the
pharisaic emissaries. To these it runs at once: \emph{Ye hypocrites, why tempt ye me?}
But to the Sadducees he says merely: {\large Ye do err, not knowing the Scriptures, nor
the power of God.} The Savior has compassion on these who have gone astray; he is ready to explain
the Scriptures to them and to instruct them about the power of God. Therefore in the last
part of the answer, too, he now keeps as close to the Sadducean presupposition as is at all conceivable. ---
``Abraham, Isaac and Jacob are, after all, long since dead. If God, then, continues in the holy
Scriptures (2 Mos.\ 3, 6. 15. 16) to call himself the God of Abraham, Isaac and Jacob, then these
venerable dead must surely live with God an eternal life; for is it not a sheer contradiction that God
should be the God of the dead? Were the Sadducees' confession of the God revealed in the law of Moses
only sincere and true, they would also perceive that the hope of the resurrection is unshakably
grounded in the power of God. The resurrection is a wonder of omnipotence. The true assurance that
the resurrection of the dead really lies ahead is to be sought only in living faith in ---
\emph{the God of the living.}''

But when the Pharisees now heard that Jesus, to the great
% --- p. 312 ---
\opage{312} astonishment of the people, had put the Sadducees to silence, they gathered together again. But it cannot
help them after all that they gather together again; for the Master surely knows the great commandment
in the law: {\large thou shalt love the Lord thy God with all thy heart, and with all thy soul, and
with all thy mind.} Yes, he even seems to indicate a defect in the lawyer's question,
in that, perhaps not without a personal allusion to the craftiness of those asking, he of his
own accord adds the law's second great commandment: {\large thou shalt love thy neighbor as thyself.}

Jesus himself at last puts an end to all these importunate but impotent attempts, in that
he (Matth.\ 22, 41) turns to the assembled Pharisees and puts to them the
question: {\large What think ye of Christ? whose son is he?} They say unto him:
\emph{David's.} That is a very correct answer: the Pharisees are, after all, learned in the Scriptures. But when it
now runs further (v.\ 45): {\large If David then call him Lord in the spirit, how is
he then his son?} it becomes evident that in the pharisaic learning there is, after all, a certain
something that is lacking. ``And this something?'' Yes, it is --- the spirit; for about David's Lord in the spirit
they now know nothing whatever. --- Modern science, which, after all, according to its own
claim, both has the spirit and knows its business well, would so gladly put in a word here, and
make the world aware that Christ ``naturally did not `know' --- the critical
introduction to the O.\ T., since the psalm (110, 1) from which he takes his proof
is supposed to be not \emph{by David, but to David.}'' This learned remark may, for that matter, be
just as good as many of the pharisaic ones; but anyone who knows something of the criticism of the spirit
surely sees that the true Messiah must, according to the given promises, be both David's
son and at the same time David's Lord. There is no subtlety at all
% --- p. 313 ---
\opage{313} in the Master's question: the Pharisees, who sit in Moses' seat, ought surely to know
who Christ is according to the Spirit; and if, notwithstanding, they do not know it, their
ignorance only betrays how little fitted they are to be the leaders of the people. It is thus
proven by the fact itself that the Lord is right when in another place (Matth.\ 15, 14) he
says to the disciples, who draw his attention to the Pharisees' taking offense at his speech:
{\large Let them alone! They are blind leaders of the blind: but if a blind man leads a
blind man, both shall fall into the ditch.}

What holds of Christ's words holds in the same sense of his wonder-works. Each single wonder
has the fundamental determination of revealing the thoughts of the hearts. As a proof of this
we may, for example, adduce the healing of the man born blind (Joh.\ 9, 1---39). Modern
criticism cannot touch this narrative without at once betraying the innermost secret of
``science.'' ``Es folgt hierauf,'' says a philosophical critic (C.\ H.\ Weiße. Die
evangelische Geschichte, \textit{II.}\ B.\ S.\ 249), ``noch eine weitschweifige, für uns,
wenn wir einmal die richtige Ansicht über die Beschaffenheit unsers Evangeliums gefaßt
haben, ziemlich interesselose Erzählung von der Art, wie diese Begebenheit unter den Juden
bekannt wurde, uud von der Bewegung, die sie unter ihnen hervorruf.'' --- Yes, of course,
``wenn wir einmal die richtige Ansicht \dots{} gefaßt haben,'' then the Evangelist will
surely fail to suit criticism. For if he has reproduced the narrative according to the more
detailed reports of other eyewitnesses, then one cannot, of course, believe ``die große
Genauigkeit des Details.'' If, on the other hand, he, who was presumably present himself when
the wonder took place, has on his own account endeavored to grasp each single feature with the
greatest exactness: ``so fällt auf, wie er diese Sorgfalt des Nachforschens
% --- p. 314 ---
\opage{314} auf ein so nichtiges Geschwätz gewendet hat.'' Nevertheless it is precisely this
great ``Genauigkeit des Details'' that we have to thank that not merely the train of thought
in Weiße's criticism --- which might, however, perhaps be of lesser importance --- but a whole
group of evangelical character-portraits steps forth into the light as at a divine sign. The
thoughtless disciples with their Jewish prejudices; the curious neighbors with their
superficial astonishment; the evasive parents with their fear of men; the hardened Pharisees
in their great embarrassment; the man born blind in his wondrous passage from darkness to
light: all these characters are grasped with such sureness and inner truth that for the
beholder they at once transform themselves into individual figures of flesh and blood. Their
peculiarity of soul is reflected from all sides in the wonder itself. It is especially the
man born blind and the Pharisees who draw our attention. The blind man who sits by the way
and begs (V.\ 8) is like a figure wrapped in obscurity. Surrounded on every side by a
spiritual and a sensible darkness, the unfortunate man rests as if in himself alone. The
customary cry of sufferers: Lord, save me! is not heard from his mouth. He has no inkling at
all that it is --- \emph{the Light of the World} that is now drawing near to him. The blind
man cannot see his Savior; but the Savior, who stands beside him, can see him, and see what
is hidden in his darkness. The misfortune is not here a punishment of God for some particular
transgression: he was born blind, \emph{that the works of God should be made manifest in
him.} Thus the wonder is prepared (V.\ 9). That Christ has not seen wrongly and chosen an
unworthy man shows itself at once; the blind man believes the unknown man at his word, obeys
his command unconditionally, goes away to the pool of Siloam, and comes back seeing. --- The
wonder is now a fact, which works differently upon different people. In
% --- p. 315 ---
\opage{315} relation to the busy curiosity of the neighbors the man born blind is rather
sparing of words. Without entering into empty conjectures he answers their question briefly
and definitely: ``it is I; a man who is called Jesus smeared clay on my eyes; at his word I
went to Siloam, washed, and received my sight; more I do not know.'' This explanation is in
the most inward agreement with the man's character. Had he let himself go in lengthy guesses
about the personality of Jesus, then we would surely have to be greatly surprised that he had
not already beforehand, when he came to hear the name Jesus, known Jesus as the prophet of
whom all the people spoke. Now, on the contrary, since he names him to the neighbors only as
\emph{a man who is called Jesus,} he lays his whole peculiarity bare in this statement. Shut
up in himself with the deep darkness that surrounded him, the blind man had hitherto
concerned himself but little with the opinions and judgments of men: this reserve gives him
not only a perhaps melancholy tinge, but also indicates (V.\ 12) the strength of temperament
with which he now comes forward as he is led before the Pharisees. The Pharisees' indignation
that the healing had taken place on a Sabbath, and their evident ill will toward ``the man
Jesus,'' do not keep him from setting forth the event as it happened, and, to their question:
\emph{what do you say of him, since he opened your eyes?} from answering straightforwardly:
\emph{he is a prophet.} --- Nothing, however, can set the Pharisees' embarrassment in a
clearer light than the interrogation that now follows. This much is settled, that Jesus, cost
what it may, must and shall be a sinner: \emph{anyone who confesses him to be Christ shall be
put out of the synagogue.} To justify this resolution one now had two ways out: either to
insist on the violation of the Sabbath law or to deny the fact outright. Now things get
% --- p. 316 ---
\opage{316} tight for the scribes. ``This man does not keep the Sabbath; how then can he be of
God?'' But on the other hand: ``He who does such signs, how can he be a sinful man?'' With the
violation of the Sabbath it will not work: \emph{there is division among them} (V.\ 16). Well
then, one must attack the fact, and begin the interrogation over again. ``The man has not been
blind; the whole thing is only a deception; let us call his parents.'' The parents, who know
the threat and know what follows upon it, do not go into more than what is strictly
necessary; this, however, they cannot deny, that it really is their son, and that he was born
blind; for the rest, they will have nothing further to do with the matter: \emph{he is old
enough, ask him yourselves.} So the turn comes round again to the man born blind: ``For God's
sake he must confess outright what the Pharisees after all already know very well
beforehand, that the man is a sinner.'' This, however, exceeds the blind man's power of
comprehension; his sound, solid answer (V.\ 25): \emph{whether he is a sinner, I do not know;
one thing I know, that I, who was blind, now see,} is so striking that the Pharisees'
following question (V.\ 26) now indeed passes over into ``ein nichtiges Geschwatz''; but this
``Geschwatz'' has nonetheless revealed the thoughts of their hearts.

It is not merely the thoughts of individual hearts that are revealed through Christ; but it is
also made manifest through him how the different positions that individuals occupy in society
partly hinder and partly ease their entrance into the kingdom of truth. In this regard there
now holds the reversal that what in a worldly sense may be esteemed an advantage is, from the
standpoint of revelation, rather to be regarded as a hindrance. --- Neither Christ nor his
disciples ever inveighed
% --- p. 317 ---
\opage{317} against the power of princes or the authority of the great, but in relation to the
ideal of humility and self-denial it is a very doubtful advantage to be mighty and esteemed in
the world. {\large You know that the princes of the nations rule over them, and the great
exercise power over them. But so shall it not be among you; but whoever will be great among
you, let him be your servant (Matth.\ 20, 25---26).}

It was, after all, rank and position that made it so hard for the chief priests and the
Pharisees to understand Jesus. Those among them who did understand him nevertheless dared not
confess him. \emph{Nevertheless,} says John (12, 42), \emph{many even of the rulers believed
in him; but because of the Pharisees they did not confess it, lest they should be put out of
the synagogue.} {\large For they loved the honor of men more than the honor of God.} How
fearfully cautious is Nicodemus (Joh.\ 3, 1. 2), and yet he is one of the boldest (Joh.\ 7,
52). Very characteristic is the following account (Joh.\ 7, 32). The Pharisees heard that the
people ``murmured such things about him'' (Jesus), and the Pharisees and the chief priests
sent officers out to seize him. These officers nevertheless came back with their errand
undone (Joh.\ 7, 45), and to the question: \emph{why did you not bring him here?} they
answered (V.\ 46): \emph{never has a man spoken as this man.} Without exactly coming to faith,
these officers have nevertheless, in an involuntary impression, perceived the power of the
divine word, to which their high superiors, on the other hand, show themselves altogether
inaccessible. The chief priests and the Pharisees feel wounded in their dignity; their wounded
sense of rank vents itself in the following outburst (V.\ 47---49): \emph{are you (officers)
also deceived?} {\large Has any of the rulers believed in him, or of
% --- p. 318 ---
\opage{318} the Pharisees? But this mob, which does not know the law, is accursed.}

Wealth is a gift of fortune: Christ does not banish wealth, and does not preach against the
right of property. But he does make manifest what these goods of temporality are in relation
to the one thing needful. {\large Lay not up for yourselves treasures on earth, where moth and
rust corrupt, and where thieves can break through and steal; but lay up for yourselves
treasures in heaven, where neither moth nor rust corrupt, and where thieves do not break
through nor steal (Matth.\ 6, 19---20).} This admonition is not set forth without grounds.
{\large For where your treasure is, there will your hearts be also (V.\ 21).} Should anyone,
however, feel himself able both to lay up treasures on earth and yet be with his heart in
heaven, then he makes a beautiful exception to the rule cited; then he is, though the rich
man, yet as one who owns nothing; then he makes himself \emph{friends by the unrighteous
mammon} (Luc.\ 16, 9); then his wealth will become a blessing both to himself and to others,
and no disciple of Christ shall presume to demand of the rich man his goods or to speak a word
against the fortunate one.

But how hard it is to make such an exception to the general rule; how hard it is to have
one's treasures on earth and yet be with one's heart in heaven: this we learn from the story
of the rich young man (Matth.\ 16, 16---23; Mr.\ 10, 17; Luk.\ 18, 18). The young man who ran
up and fell on his knees before Jesus on the way, and asked him: \emph{good Master, what shall
I do that I may inherit eternal life?} by no means belonged to those who misuse the gifts of
wealth in the service of unrighteousness or of base desires; on the contrary, he knew all the
commandments, he had kept
% --- p. 319 ---
\opage{319} them from his youth. Yes, youth, a good upbringing, and on top of that very great
possessions! he felt so fortunate! Now he wished to add to the many goods of fortune the
highest good as well; he had inherited so many treasures on earth, now he wished, if it could
be done, to inherit eternal life as well. The manner in which he puts forward this wish of his
certainly betrays a youthful rashness; but it also indicates that a longing for the perfect
has awakened in him; he must surely be prepared to make, if not the greatest, then at least a
quite acceptable sacrifice, since he so boldly asks the ``good Master'': \emph{what shall I
do?} --- The young man was without doubt inspired at the moment he ran up and fell on his
knees on the way. Thus we understand what Mark reports (10, 21): But Jesus looked at him, and
loved him, and said to him: \emph{one thing you lack; go away, sell what you have, and give it
to the poor, and you shall have treasure in heaven; and come, take up your cross, and follow
me.}

The ``good Master'' now wished at once to ease the fortunate young man's way to the coveted
perfection by freeing him from a burdensome load: alas, the fortunate one did not grasp the
liberation. \emph{But he became displeased at this saying, and went away sorrowful; for he had
great possessions.}

To be rich in earthly goods and yet lay up treasures for oneself in heaven is, if not
impossible, then hard, indeed very hard: {\large It is easier for a camel to go through the
eye of a needle than for a rich man to enter into the kingdom of God (Mrk.\ 10, 25).}

Knowledge, insight and cultivation are precious goods, well worth striving after, both for him
who would work in the temporal and for him who would appropriate the eternal. Christ, who is
himself \emph{the Light of
% --- p. 320 ---
\opage{320} the World}, likewise casts no shadow on thorough knowledge, any more than on true
enlightenment (Matth.\ 5, 13---16); but it has nevertheless been made manifest through him
that all the advantages which in the world of spirit are acquired through learning,
enlightenment and higher cultivation are ambiguous advantages when it is a matter of
appropriating the divine truth in childlike faith and simplicity of heart (Matth.\ 5, 3). How
near was that lawyer who (Matth.\ 22, 35) tempted him with questions about the great
commandment, how near was this learned man to coming to the knowledge of the truth (Mr.\ 12,
34): Jesus saw that he answered discreetly, and said to him: \emph{you are not far from the
kingdom of God.} What was it, then, that held the scribe back, so that he did not at once
enter into the kingdom of God? It was without doubt his learning in the law! --- Indeed, when
even in our days we see the learned search and search in the holy books without, for all
their learning and searching, coming to the knowledge of the truth, how often must not the
words of Jesus (Joh.\ 5, 39---40) come to our mind: {\large You search the Scriptures,
because you think that in them you have eternal life; and it is they that bear witness of me.
And you will not come to me, that you might have life.}

If, now, rank and position and outward advantages lay such great hindrances in the way of a
believing appropriation of the Gospel of salvation, then it follows from the nature of the
case that the Savior must especially have appealed to those who had no such hindrances to
overcome. On this presupposition, then, Christ can be said in a special sense to be the man of
the people. Though a son of David, he nevertheless himself came forth from the midst of the
people. To the people he addressed his word, and the people heard him gladly, for he taught,
after all, as one who had authority, and not as the scribes (Matth.\ 7, 29). To the poor he
preached the Gospel (Matth.\ 11, 5), and in
% --- p. 321 ---
\opage{321} the people he found --- the poor, the mourning, the suffering. From the people he
chose his disciples, and in the temper of the people he found a natural protection against his
embittered enemies. But he was not a man of the people in the sense that he would make a pact
with the dark passions of the masses in order to overthrow the mighty. Not in order to use the
people as it is, but in order to reveal the thoughts of the hearts, and thereby make manifest
what the people is, the Son of Man steps forth, lifts up his voice, and speaks the word of
truth.

What, now, is the popular spirit, what is a people? It might perhaps seem strange, indeed in
our days, when everyone knows just as well what the popular is as what the people is, almost
indecent to ask in this way. But for a slow head there is nevertheless still some reason to
put forward this question. From every side there is crying to the people and from the people
and for the people; but that is not to say that the many cries explain what the people is, or
wherein the popular really consists. Little is gained by a formal delimitation of the general
concept. If one says: the people is in a whole country what the family is in a single house,
then one points to the national. A people is a nation. But every nation in turn contains many
different elements of the people: the great family has many adopted sons and many
illegitimate children. If one says: the people is a common designation in which all the
difference that is grounded in rank and position, cultivation and ability, is absorbed, so
that man thereby comes to count as individual, and the individual as man: then this universal
concept can in turn be regarded from several different sides. For insofar as the difference
of individuals vanishes in the natural unity of the life of the race, the people is a
\emph{mass}; insofar, on the other hand, as the unity acts and expresses itself through the
flock of individuals,
% --- p. 322 ---
\opage{322} the mass is a \emph{crowd}. The concepts of mass and crowd can in turn be grasped
from a double standpoint. The mass has as yet no cultivation; it is only about to be
cultivated: the mass is raw, the mass is ``rabble.'' As cultivation begins to penetrate the
mass, it separates into a multiplicity of reasoning and criticizing individuals: through a
many-sided, superficial cultivation the crowd becomes a ``public.'' Both concepts, ``rabble''
and ``public,'' are one-sided concepts of reflection. If it would be a bloody insult to
humanity were someone to say of the people that it was only rabble (\textit{le peuple}!), it
would be an equally harsh, if not quite so palpable, affront were someone to assert that the
people was only a public. --- The ground of the popular spirit is one with the original natural
ground of our race. No one deserves the name of a man of the people merely because he stirs up
the mass; and no one ought to be called a character of the people because he knows how to lead
the crowd and ingratiate himself with a public. No, only he is in truth a man of the people
who is able to address man as man. The word of the people is an authoritative word that
penetrates deep; a creative word that makes the soul fruitful; a blessed word that gladdens
the heart and rejuvenates the being of man. --- O, what a difference between becoming a
darling of the crowd and being --- a man of the people!

Only in relation to Christ is it made manifest in truth where the people is. --- ``Where, then,
is the people?'' In the mass! but the mass, as mass, is not the people itself. On the
battlefield of finite interests the strife over the right and significance of the masses will
never be able to cease. Anxious rulers fear every word that can set the crowd in motion and
bring the mass into ferment. This selfish anxiety is alien to the popular spirit: the mass
shall ferment; for only out of the fermenting mass can the people clarify itself and come
% --- p. 323 ---
\opage{323} to view. But the heated hunters of the people, who hunt up the passions of the
masses so that the raw elements must ferment and seethe without cease, lay the people waste:
in the wild ferment the popular spirit, after all, sinks clean to the bottom.

\begingroup\large
He told them another parable: The kingdom of heaven is like leaven, which a woman took, and
hid in three measures of meal, till the whole was leavened (Matth.\ 13, 33).
\par\endgroup

As long as this leaven works in the mass, the Redeemer works with it, and helps the people out
of the ferment. Woe to him who dares to knead his own unwholesome ferment into this leaven and
corrupt it! Against such a one the mass shall one day rise in its ferocity, and avenge the
wrong the people suffered.

``Where, then, is the people?'' --- In the crowd! but the crowd, as crowd, is not the people
itself. When the dark mass draws near to the light, the crowd can glimpse; but only when the
light penetrates into the heart of the crowd can the people see. There is still complaint over
the darkness of the world: the crowd cries for light. So the obliging leaders go out, each
with his lantern, to light the way for the people. Alas, how the many lantern-men do confuse
the mass\footnote{%
If the kingdom of heaven is, in relation to the mass, to be likened to leaven, then the word
of life is, in relation to the reasoning crowd, rather to be likened to salt. As the crowd
articulates itself (\textit{sit venia verbo}) into a public, and loosened reflection begins,
little by little, to shake the first simple truths, the human temperament undergoes a most
precarious change. The mass, which lives in God's free nature and breathes the rawer but
sharper air, will always preserve a certain natural freshness. But the reflecting, reasoning,
all-criticizing crowd, which drifts about in the sultry atmosphere of half-culture, can very
easily take harm to its health. It goes with the public in the great cities as with fat meat
when there is thunder in the air: it takes on a thought, it has a thought, that is to say, it
turns rancid. If, then, the reasoning crowd, which surely has some salt, but alas all too
easily takes on the many rank thoughts, is to keep itself fresh in the first natural
fundamental thought, then there must be vigorous salting. Therefore Christ says to his own
(Matth.\ 5, 13): \emph{You are the salt of the earth. But if the salt loses its strength,
with what shall it be salted? It is good for nothing else than to be cast out and trodden
under foot by men.}%
}, and dazzle
% --- p. 324 ---
\opage{324} the crowd! --- Only the holy, simple light is able to shine for all, and only when
it in truth shines for all is the people enlightened.

\begingroup\large
You are the light of the world; a city that lies on a mountain cannot be hidden. Nor does one
light a candle and put it under a bushel, but on a candlestick; and it shines for all who are
in the house. Let your light so shine before men, that they may see your good works, and
glorify your Father who is in heaven (Matth.\ 5, 14---16).
\par\endgroup

``But now the instinct of the people?'' Yes, the instinct of the people will always seek the
good. But the temper of the people, which in womanly modesty conceives the good, is corrupted
when it gives itself over to the vain imagination that by its own power it can create and
bring forth the good itself. The modern impudence which fawns upon the crowd, sanctifying
natural drive in its sentimental language of flowers, must praise instinct in lofty tones. But
when the ``friend of the people'' has then himself come to the helm, he all at once becomes
% --- p. 325 ---
\opage{325} so silent, and is in wondrous fashion transformed. It is as though he now suddenly
had an inkling that the instinct of the crowd could all too easily err. What the reason for
this may be we will not examine in this connection; but much rather see how the very ground of
the popular spirit is illuminated by the revealed word.

{\small Gospel, Matth.\ 13, 3---9.}

\begingroup\large
Behold, a sower went out to sow. And as he sowed, some fell by the way, and the birds came and
devoured it. But some fell on stony ground, where it did not have much earth, and it soon
sprang up, because it had no depth of earth. But some fell among thorns, and the thorns grew
up and choked it. But some fell on good ground and bore fruit: some a hundredfold, some
sixtyfold, some thirtyfold. He who has ears to hear, let him hear!
\par\endgroup

If it is now true that the soil of the people has so varied a ground: how then could the
instinct of the people be infallible? But even if the fertility were equally great
everywhere, the fruits would nevertheless not be equally good. In the soil of the people falls
the good seed as well as the evil: on this field the tares grow together with the wheat
(Matth.\ 13 24---30).

``But who, then, shall lead the people, and how, then, shall it be brought up?'' Two questions
in one breath! Yes, it becomes an endless strife when one will not go back to the
unconditioned first. The people, after all, is in the mass, it tumbles about in the crowd: if
the mass is not organized and the crowd not governed, the people will be tossed about until at
last it is altogether dazed. So it holds unconditionally that there must be ruling: the more
vigorously there is ruling, the stronger the people becomes, when it has --- mark well ---
found the right ruler. The people is oppressed in the mass and sighs in the crowd:
% --- p. 326 ---
\opage{326} So it holds unconditionally that the mass must ferment and the crowd stir, that the
bound may be loosed and obtain his freedom. The faster the liberation can proceed, the louder
the people rejoices, when it has --- mark well --- seen and recognized the true liberator. The
true ruler and the true liberator are not two, but one person, i.e., the Redeemer\footnote{%
There is surely no one at this time who can hear the word ``freedom of the people'' without at
once thinking of politics. ``Since, as the saying goes, one cannot learn politics from
Christianity (unless it be the well-known priestly politics), it is only an empty phrase when
it is said that Christ is the Redeemer of the peoples.'' --- Let us consider this matter in
all brevity.

\textit{a)} All the world knows that the freedom of the people must be grounded on truth and
right; but all the world also knows that truth and right have their great difficulties.
Enterprising leaders of the people undeniably want truth to a certain degree, and justice to a
certain degree. But why only to a certain degree? Naturally because it is only politic to
serve truth and right when one has a little profit from the service oneself. To serve gratis
is impolitic. --- In the Gospel the relation is reversed. Truth says to everyone who goes
about with great plans: ``If you would benefit the state and be a man of the people, then for
God's sake never think of profiting from the mass or of seeking comrades in the crowd. But if
you are the unconditional servant of justice, then only step forth with all boldness. Even if
it seems that you, despite all your exertion, labor in vain, only be of good cheer; even if
you draw upon yourself the hatred of the crowd, only be of good cheer! I, the almighty Truth,
shall surely vouch to you that you nevertheless benefit the people and receive your
reward.''%
}; o, let us say it freely to all peoples: there is only one Redeemer.
% --- p. 327 ---
\opage{327} Modern wisdom asserts that \emph{the voice of the people is the voice of God}; the
Gospel teaches us in what sense.\footnote{\textit{b)} All the world knows that truth and right
cannot be realized in civil society without entering into the conditions of finitude. Quite
right! What, then, does politics do? It adjusts its conviction to circumstances, and so loses
the unconditioned in endless conditions. In itself justice is, after all, nothing, or as good
as nothing: everything depends on circumstances. The prince who yesterday cried with a loud
voice that he was called to rule by the grace of God, today has a weaker voice, and bows
beneath the grace of the crowd --- owing to circumstances. The teacher of the people who
yesterday laid down his firm conviction of many years in the well-written article in which he
judged the living and the dead, today became a minister by entering into an entirely new
conviction --- owing to circumstances. The mass, which yesterday was meek and quiet, then
became today, so as not to lag behind the prince and the minister, untamed and unbridled ---
likewise owing to circumstances. Thus everything is thrown into confusion, and in the end the
freedom of the people is probably clean forgotten --- owing to circumstances. --- In the
Gospel the relation is reversed: the conditioned is made dependent on the unconditioned.
Divine justice asks everyone who intends to lead the people: ``Have you now really so firmly
grounded a conviction that you dare risk everything to carry it through: then only step forth
boldly, and I shall be your advocate. But if your conviction is of the kind that decides right
and wrong according to circumstances, then woe to you, you corrupt the people; woe to you, you
shall receive the harsher judgment; woe to you, even if you won thousands and tens of
thousands of comrades in the crowd, you shall one day be judged by the people!''} The entry
% --- p. 328 ---
\opage{328} into Jerusalem (Matth. 21, 1---16. cf. Mark. 11. Luk. 19. Joh. 12) and the walk to
Golgotha (Matth. 27. cf. Mark. Luk. Joh.) are the two scenes in which the voice of the people
rises so high that its cry reaches up to heaven.\footnote{\textit{c)} All the world knows that
great strength and much self-sacrifice are needed to help the people to its right. But all the
world also knows that every leader of the people has his ambition. Nothing will ever come of
the people's right, however, so long as ambition is to be remunerated \textit{ad usus
publicos}. He who has not the strength and courage to serve justice even in unconditional
self-denial can never be the friend of the people; but he who, out of true love for the
people, would be the incorruptible servant of justice can hardly become the friend of the
crowd. So it happens that he who loves the people from his heart is misjudged by the mass and
hated by the crowd as an unbending despot. --- See, that is why the liberation of the people in
the Gospel is not calculated on ambition either, but on self-denial. But the high road of
self-denial leads to Golgotha. --- This, then, is the politics we can draw from the Gospel.}
At the entry the cry is: Hosanna! it is the people's jubilation, it is the crowd's cry. How
Christ interprets this cry is clear from his answer to the embittered Pharisees: \emph{``Master,
rebuke your disciples.''} {\large If these should hold their peace, the stones would cry out
(Luk. 19, 39---40).} \emph{``Do you hear what these} (the children who cried out in the temple.
Matth. 21, 15---16) \emph{are saying?''} {\large Have you never read: out of the mouth of babes
and sucklings you shall prepare praise for yourself.} So the Hosanna-cry sounded as a voice of
the people through the rejoicing crowd. By this voice of the people it was (Matth 21, 5)
{\large proclaimed to the daughter of Zion: behold, your King comes to you meek, and riding on
an ass, i.e., on the she-ass's
% --- p. 329 ---
\opage{329} foal.} --- But now the other cry, the cry that sounded in the streets of Jerusalem
when the crowd raged: \emph{Crucify! Crucify!} the cry of vengeance that delighted the hearts of
the Pharisees and filled Pilate with loathing and horror: is this cry, then, also the voice of
the people, or has the spirit of the people in this wild scene departed from the crowd? --- If
we conceded such a separation between the people and the crowd, the concept of the popular
spirit would lose its reality. The cry \emph{Crucify!} must therefore be of the people just as
much as the cry \emph{Hosanna!} In the one case it is a good spirit that is upon the people, in
the other it is an evil spirit; but in both ambiguity reigns. In the cry Crucify! too Christ
must hear a voice of God (it behooved him, after all, to suffer this before he entered into his
glory); but if the raging of the crowd is nevertheless not God's work, then neither can the
jubilation of the crowd be God's work. If the people's cry Crucify! cannot immediately pass
itself off as the voice of God, then neither can the people's cry Hosanna! The one here
presupposes the other: had Hosanna! not first been cried, then Crucify! would presumably not
have been cried afterward either. The one can only be understood through the other: had the
people not lived in the crowd, there would never have been the cry: \emph{Hosanna!} Had the
crowd not raged in the mass of the people, there would never have been the cry:
\emph{Crucify!} The misunderstanding of the crowd, which hid itself in \emph{``Hosanna''}! is
loudly cried out in \emph{``Crucify''}!

Modern wisdom teaches that \emph{world history is the world's judgment.} --- If this wisdom in
earnest believes what it itself says, then the conduct of the present is altogether
incomprehensible. If history is really so strict and earnest a judge, how then can it come
about that this swarm of insignificant individuals of every sort, who throng together from all
the corners of private life, gets the wondrous
% --- p. 330 ---
\opage{330} courage to come before the face of the earnest judge and clamor that it now
absolutely must get into world history and have its names recorded in the annals of the
nations? What, then, is one to think of modern wisdom? --- Is getting into world history only a
strenuous amusement, like joining in the noisy game that is played at emperor and king? is
seating oneself on the high seat of power an entertaining bit of fun, like sitting for a
little while on the chair of astonishment? has immortality now become so liberal that every
petty half-hour hero can at once have his fame entered in the annals of the nations? does one
now really become renowned when one can merely manage to get one's honored name scratched on
some pane or other in some public building or other (be it church, theater, exchange or town
hall): well, that is another matter, then the conduct of the present is quite comprehensible.
If it is really true, what the modern consciousness asserts as quite certain, that there is no
immortality in heaven, then it is naturally in order that everyone who wishes to outlive
himself must address himself to history. With Herostratic presumption individuals break loose
against the barriers of existence, bent on setting the whole traditional world on fire, merely
in order to become immortal. One ought, however, to think a little more closely about
historical immortality. Without mentioning that when the traditional world burns down all the
annals could easily burn along with it, and the coveted immortality go up in smoke, one ought
nevertheless, in case the annals should come safely through so great a conflagration, to think
a little more closely about historical immortality. --- There was once a great and mighty
emperor; he had made himself so immortal that no one could come near him without becoming
immortal through him. But this immortality worked differently upon different people. If the
emperor was very gracious and wished to reward someone,
% --- p. 331 ---
\opage{331} he bestowed on him --- immortality; if now and then he became angry and really
wished to punish someone, he likewise bestowed on him immortality. All his faithful servants
he favored with an immortal name; but to the jailer he said at last: I condemn you to
immortality. This, then, is one of the features of world history; let us now turn back again to
the Gospel, to see what sacred history teaches about judgment and immortality.

That history judges means, in other words, that it reveals what is hidden, clarifies the
eternal out of the temporal, makes separation between right and wrong, truth and lie. So long
as the verdict of eternity is lacking, the judgment of history is unreliable. If it is now true
that the eternal only through Christ becomes flesh and blood in the temporal, then it must
surely also be true that history only through him receives its proper authority as judge, that
the judging principle only in his person enters history with undimmed clarity. He came to save
all, but in actuality he saves only the believers. --- But as the Revealer Christ is in truth
the judge of all the world. This authority as judge he did not first receive in his exaltation,
when he was taken up into heaven and sat down at the right hand of the Father; he had it
already in his humiliation, when he went to Golgotha and sank beneath the cross. He who cannot
recognize this judge in his humiliation will surely not yet recognize him in his exaltation
either, but will nevertheless learn to know him when the day comes on which we shall all be
made manifest before his judgment seat. --- In a worldly sense it is the enemies who forced
Christ to bear his cross, but in a divine sense it is Christ himself who induces his enemies to
lay the cross upon him. In the language of finitude it is the overpowering enemies who say to
the Son of Man: ``See, now we lead you to Golgotha, that you may suffer your deserved
punishment as blasphemer
% --- p. 332 ---
\opage{332} and rebel against Caesar''; but in the language of eternity it is the Son of God
who says to the blinded children of men: ``See, this is now your hour, but it is also my hour:
this is your hour and the power of darkness, but this is also my hour, which the Father has in
his power. Therefore it now behooves you all to go to Golgotha with me, that what is hidden in
darkness may be made manifest in the light, that a righteous judgment may go forth over this
world.''

Thus has the Son of God and of Man judged humanity in his own race; in order to make the
judgments of history clear, he has condemned a whole generation to immortality. What mildness,
and yet --- what severity in this judgment! So many familiar figures, so many kindred souls
meet us here in the mild light of the Gospel. Here are names and deeds to which we may well
apply the words the Savior himself spoke over the pious woman who anointed him in Bethany
(Matth. 26, 13): \emph{Verily I say unto you: wheresoever this Gospel shall be preached in the
whole world, there shall also this, that she has done, be told for a memorial of her.} But here
is also another judgment, a judgment passed on names and deeds of which one might well say
that, with all the guilt they bear within them, they have nevertheless not raised themselves to
the extraordinary by their own power, or earned their own way up to the immortality of crime.
We are thinking of the judgment that for Christ's sake has come upon a Herod and a Pilate, an
Annas and a Caiaphas. In these examples we see how dangerous it can sometimes be to come too
near the light. How many highly trusted men have not found themselves in the same difficult
position as that Roman governor! How many a procurator has not, both before and since, with the
same bearing understood how to be the friend both of Caesar and of justice, and with the same
shrewdness
% --- p. 333 ---
\opage{333} known how to give up the innocent and wash his hands before the people. Many such
deeds are overlooked, many names are forgotten; but there is nevertheless one name that can
never be blotted out of the memory of history, and that is the name Pontius Pilate. --- Many a
clever statesman has, both before and since, taken his seat in the council of the mighty and
spoken of the necessity of sacrificing the single individual to the claim of the many, and has
perhaps even taken pleasure in his own wisdom. But in the great multitude of politically wise
sages single names could easily be forgotten. See, that is why the Gospel, after all, has also
wisely provided that there should at least be one who did not miss the well-deserved
immortality. No priestly statesman has probably ever spoken with greater self-confidence than
Caiaphas when he said in the council (Joh. 11, 50): \emph{You know nothing at all; nor do you
consider that it is expedient for us that one man die for the people, and that the whole
nation perish not.} A profound word, which fell at the right time. But neither has any
prophetic statesman ever received such recognition of his political insight as Caiaphas, for
it is expressly conceded that the wisdom in his words must be referred to a higher inspiration
(V. 51): \emph{this he said not of himself; but, being high priest that year, he prophesied
that Jesus should die for the people.} Yes, Caiaphas --- can be sure enough of his
immortality.

Seen from this standpoint, ``the procession to Golgotha'' is to be regarded as a triumphal
procession, whereby the earthly life of Christ completes its revealing activity. And what
triumphal procession\footnote{Seen in this light, the procession to Golgotha must presumably be
taken to begin from the Garden of Gethsemane.} offers
% --- p. 334 ---
\opage{334} a richer spectacle too! See, Jews and Romans, people and Pharisees, mocking
soldiers, weeping women, vengeful priests, disheartened disciples: all of them meet on the same
way, all of them move toward Golgotha. Here now goes Judas, muttering in his darkness; there
stands Peter, weeping in his remorse. See this youth, aged in the earnestness of sorrow: it is
the beloved disciple; and she, the courageous woman who walks at his side, though she can
scarcely move her foot: it is the Mother of Sorrows. These two now form a solitary pair in the
midst of the throng; these two now walk together the narrow way --- up to Golgotha. Along this
way the King of sufferings will draw them all; on this way he will gather them all, on Golgotha
he will speak to them all --- the last word, on Golgotha he will condemn them all to ---
immortality.

\subsection*{\textit{V.}\\[4pt] Christ and the Spirit of Darkness.}
\addcontentsline{toc}{subsection}{V. Christ and the Spirit of Darkness}

The faith of the Gospel and the wisdom of the world are incompatible. The ambiguous union in
which they have hitherto so often worked for one another and yet against one another is to be
compared with an unhappy marriage: every approach, however mutual and well-meant, ends in
painful misunderstandings; every clash strikes a new spark for new disputes and domestic
quarrels. Of so unhappy a couple prudence is wont to cite its well-known rule: ``if they
cannot agree to live together, then let them above all agree to part.'' But here too divorce
has its great difficulties. Modern knowledge is no merely external tormenting spirit that one
can be rid of by an energetic step; the negative consciousness is no ordinary
% --- p. 335 ---
\opage{335} domestic tyrant whom one can escape by giving him notice in the name of faith and
moving to another quarter of existence. When the wisdom of the world rules over the whole
world, where can the faith of the Gospel then expect to find a quiet, secure refuge? Even if
faith were once more to take the veil and go into a convent, it would not help. Even in the
most pious soul it must put up with living together with doubt and criticism. We are, after
all, all by nature critics in relation to the supernatural facts, all doubters in relation to
the Gospel of faith, all of us together children of time in relation to the manifestations of
eternity. Should this assertion seem to anyone too harsh, we can easily put it to the test.
You who boast of your perfect faith, hear a word: can you move mountains? --- We are, after
all, entitled to ask thus when we rightly remember what Christ himself promised his disciples
(Matth. 17, 20): \emph{verily I say unto you: if you have faith as a grain of mustard seed, you
can say to this mountain: move hence to yonder place, and it shall move; and nothing shall be
impossible to you.} The most pious among the pious will now probably have to reflect a little
before he answers. If he answers us: ``these words are spoken as in a parable; such a promise
cannot be understood literally; no, such wonderful things faith cannot in reality
accomplish'': then we hear at once the doubter in this critical reinterpretation. If, on the
other hand, he answers us: ``The promise stands fast. For the believer nothing shall be
impossible. I believe in truth that if I only dared, in unconditional trust in Christ's
omnipotence, to utter the word: \emph{move!} the mountain would move and --- vanish. But, I do
not know how it is, I nevertheless dare not say to the mountain: move! I am so afraid of
tempting God. No, when I rightly consider the great heavy mountain, it is
% --- p. 336 ---
\opage{336} impossible for me to say: move! --- I fear the madness of fanaticism,'' --- if the
Believer answers in so candid a language, then we shall by no means be dissatisfied with his
answer, or tempt him to a presumption; but we must nevertheless quietly point out to him that a
most precarious doubt hides itself in his ``perfect faith.'' The most pious is thus, after all,
in his way a doubter and critic, just as all the rest of us are. If he cannot move for us the
``great, heavy mountain'': then he will surely also have trouble enough moving any one
whatsoever of the --- great heavy mountains of cares or of sufferings or of the agonies of
death.

If Christ's word has now convinced us that deep within every man there sits a doubter and free
thinker, then it is also not conceivable that the faith of the Gospel should be able to lie
down to rest in any heart without heeding criticism. The Believer dare not expect that the
inner strife shall be ended because at one moment or another he has been lucky enough to win a
victory over unbelief. The Believer knows that a great danger is always bound up with a great
victory. When faith rejoices and sings \textit{Te Deum}, the enemy, as often as not, gathers
new reinforcements and steals forth from a secret hiding place to take the victor by surprise.
Therefore faith dare never celebrate any victory feast without first doubling the watch,
having all the outposts out, and keeping the strictest military discipline, so that none of
those on watch shall grow drunk with joy and sleep at his post. Faith itself dare never sleep;
it is literally --- sleepless. But it is not for that reason comfortless. We all need, after
all, the rest of sleep insofar as we live a temporal life: the natural man has a horror of ---
sleepless nights. But he who in faith would live an eternal life dare not think of sleeping;
for in eternity one never sleeps. To want to hush doubt, in order to
% --- p. 337 ---
\opage{337} settle down to rest in faith and hope for eternity, is of all self-contradictions the
most uncanny. Against hardened criminals the torture has sometimes been used of not
letting them sleep until they had confessed. That is in truth a terrible punishment.
But for one who means to go into eternity without having learned to keep awake, eternal
life is an even more terrible punishment. Unbelief says that the dead sleep an eternal
sleep. This proves that unbelief has its life-root in the temporal. Defiant unbelief
challenges death, because it relies on being able, without further ado, to fling itself,
crime and all, into a grave; unbelief goes with bold steps to meet its eternity, because
it relies on this eternity being --- a sleep. But suppose now that eternity is sleepless;
then it is surely a serious matter if a person who was fully and firmly prepared to find
in the grave an eternal sleep now came down there and found --- an eternal sleeplessness;
or would it not be ghastly if some soul had to grope its way through the long eternity
like a melancholy sleepwalker, like a Lady Macbeth, sleeping, and yet --- sleepless?

For flesh and blood it is a dangerous thing to have dealings with the spirits. --- Name
whatever spirit you will; they all bring you into danger. If only we knew that unbelief
was right; yes, if only we knew that man's life was concluded and ended with this
existence: then we would at once have to award the wisdom of mediocrity the highest
prize, and call out to every person, be he learned or unlearned: ``take heed of the
spirits, lest you fall into temptation. Remember that you are flesh and blood! Never come
too near any spirit, never nearer than the bounds of prudence and sobriety and
cultivation allow. One can, after all, in truth be both cultivated and learned without
exactly coming too near the spirit. If you must, once in a while, so as not always to
put your light under
% --- p. 338 ---
\opage{338} the bushel either, venture a little further; then by no means carry it beyond
fleeting spiritedness. One can at a pinch be spirited without on that account having the
spirit in the serious and perilous sense of the word. Humane spiritedness shows itself
chiefly in the ease with which one makes the transition from faith to doubt, and
conversely from doubt to faith. If it is faith that is in fashion, then one lets one's wit
(that is, if one has a little wit) play against the spiritless doubters; if on the other
hand unbelief is uppermost, well, then one plays and mocks with faith. If one only has
the courage in this way always to take the side of the victors, one can always be sure
of advancing with the times from victory to victory: spiritedness is mobile, like the
mobile spirit of the age. Insofar as you now follow these admonitions of ours, you do
what is in your power. But should you nevertheless notice that a real spirit --- oh
horror! --- was approaching and wanted in earnest to get hold of you, then you must in
good time do everything in your power to get away. You cannot, to be sure, simply run off
at once (for the man must be quick on his feet who is to enter a race with a spirit); but
you ought at least to seize every means of self-defense at the disposal of a reasonable
person. Only be careful, but do not lose heart! There are, after all, means enough, if
one only knows how to use them. Besides the usual household remedies: social
stimulations, trips abroad, the latest news, and the like, one might well also count a
learned occupation with old things, especially a certain kind of biblical criticism, such
as: how many siblings Christ had? how many jugs of wine the guests drank at the wedding
in Cana? how it came about that the tower in Siloam fell down? how many Jameses one finds
in the N. T.? etc., among the innocent means by which even very learned men have known
how to secure themselves, with enviable success, against the designs of the spirit. Never
trust any
% --- p. 339 ---
\opage{339} %
spirit! even spiritedness tempts. Every spirit is dangerous for flesh and blood; but of
all spirits under heaven the spirit of religion is the most dangerous of all.''

When the enemies of revelation really want to attack Christianity, they usually like to
dwell with a certain complacency on the many aberrations of spirit that blind religious
zeal has caused in the course of time. ``Terrible,'' they say, ``is world history when it
opens its exhibition with the bloody scenes of crime and displays to us what is dark and
savage in our nature; but all this is only child's play in comparison with the crimes and
demonic scenes of horror that the history of the Church has to show. In secular revolts,
however raw and savage they may otherwise be, we still always recognize some trace or
other of the original majesty of the human being; but the misdeeds of church history rise
to the absolutely inhuman. What madness do we not meet in such self-tormentors as the
pillar saints and flagellants! What works of darkness have the annals of the monasteries
and the Inquisitions not to bring to light! There is a society, --- humanity shudders at
its mere name, although the name itself is the holiest in Christendom, --- it is the
\emph{Society of Jesus}: what if history could gather into one picture and lay bare to us
the main features of the spiritual exploits this sinister society has now already been
carrying out for 300 years --- `for the greater glory of God'; could any human eye bear
the sight of such a picture? --- By the fruits shall the tree be known. If you would see
the ripe fruit of orthodoxy in a truly glowing light, then above all do not forget to
think of a St.\ Bartholomew's Night.''

Without circumlocution we admit at once that the Free Thinker's complaint is in this
respect only too well founded. It does no good to try to weaken the impression of the
dark, uncanny stains of church history
% --- p. 340 ---
\opage{340} by emphasizing the bright side of revelation; nor can it help to appeal to the religious
sobriety of the present: the fine equilibrium in the religious way of thinking of the
present could easily turn into a troubling testimony to the weakening of Christianity as
well as of religiousness in general. We avow the conviction that new religious
breakthroughs could hardly occur without calling forth new aberrations; we readily admit
that there can be no spiritual power on earth capable of occasioning such sorry
misunderstandings, such revolting misdeeds, and such terrible errors as the religion of
Jesus Christ. But this, then, is also our contention: that among the many historical
proofs that the theologians know how to record in defense of Christianity, the proof
that can be derived from this ought to be counted as one of the very strongest.

There is a saying that he who would climb high may fall deep. Now surely no one in the
world of spirit can aim to climb higher than he who resolves to sacrifice everything in
order to follow Christ. Let us here consider the peculiar position of such a person a
little more closely. The first step is an irrevocable resolution, by which the awakened
one casts from him the usual means of support, breaks with the surrounding world, and
ventures out into the ultimate decision. Here a dangerous turning point now sets in.
Temporality is not so easy to dismiss; it follows after, though at a distance; it waves
farewell, as if it would once more call the fugitive back. The world? Yes, it is dark,
angered at the presumptuous one who dares to turn his back on it. But eternity? It is
cold in all its loftiness; it looks down so thoughtfully on the bold one, as if it would
ask him: do you yourself really know what you are doing? --- Once a person has worked his
way out of the guardianship of finitude and is on the point
% --- p. 341 ---
\opage{341} of becoming spirit, then at the entrance to the realm of spirits he must surely know
himself able to give an account of himself: ``What do you want here? Why did you not stay
with the others? Why did you give away all your goods? You could, after all, have kept
them yourself and still kept the commandments. Why did you leave father and mother, kin
and friends, as if your true destiny were to be found only out here in the regions of
emptiness, in the desert of renunciation?'' These questions press forward like surprising
promptings, and voice themselves the more vehemently, the greater the break the
resolution has produced. But whence do such promptings come? It may be that it is only
flesh and blood that tempts the one renouncing; but it may also be that it is the spirit
itself that now wants to test its swing. It may be that it is a seducing spirit that
only wants to plunge the person into ruin; but it may also be that it is a saving spirit
that wants to secure the weak one against overwhelming dangers. These questions are
altogether ambiguous; and yet they demand, each for itself, a definite, an instantly
decisive answer. Blessed is he who finds the right answer! he saves his soul. Woe to him
who mistakes the answer! he perishes. If the question comes as a prompting, the answer
must surely also come as a prompting. In determining himself spiritually, a person must
give himself entirely into the power of the spirit. Here, then, is the turning point at
which the spirit itself will come; but --- which spirit? --- For flesh and blood it is a
dangerous thing to have dealings with the spirits!

Before we seek the final enlightenment on this obscure point in the Gospel, however, it
would surely be expedient if, for a preliminary discussion of the matter, we first
consulted a little more closely with the modern consciousness. A consciousness that has
already come so far that it can rightly be said to have the light in itself must surely
be able to find its way in the dark. This
% --- p. 342 ---
\opage{342} is in any case a great advantage, since one then need not delay oneself at all by
waiting for any light from heaven. Let us listen!

``It is the old strife between light and darkness, between good and evil, between Ormuzd
and Ahriman, which in the mythology of the New Testament is presented to us as a personal
battle between Christ and the devil. Which of these two rulers finally keeps the upper
hand cannot be in doubt. The worse the devil rages, the more Christ's strength grows. We
see it in the fate of the earliest Church under the persecutions; we see it again in the
history of the Reformation. Never, surely, has the Antichrist been nearer to laying the
Church waste than when, in Luther's day, he had encompassed the whole of Christendom with
his nets, stirred up the Turk, made the Pope his creature, and bribed the Roman clergy.
It is uncertain how far the corruption would have gone had the devil not at the right
time found his match in the stout-hearted Doctor Martin; but it is entirely certain that
Luther would never have become the man he became had he not seen himself compelled to
take on the devil in his own person. Many a midnight hour did these two fight the hard
fight with each other; and Luther would hardly have come off well had Christ not sent him
extraordinary reinforcement. The devil was a sharp dialectician, but every time he
advanced with doubt, Luther struck at him with faith; the Evil One, he was a cunning
sophist; but when he wanted to choke the anguished conscience with the severity of the
law and the punishment of sin, Luther seized on the merit of Christ, and then the devil
had to yield. Had the great Reformer not been tried in such a struggle, where then would
he have found the courage to defy both the Pope and the world and all its hosts? Now, by
contrast, since the bold man of faith had seen the dark enemy himself face to face, he
feared no human
% --- p. 343 ---
\opage{343} power: he knew the author of evil. `The heathen,' he says, `do not know whence
misfortune comes so suddenly; but we know that it is sheer work of the devil; he has such
halberds, lead bullets and guns, such spears and swords, with which he shoots, throws and
stabs among us, when God gives him leave.' --- Never did this hero of faith lose sight of
his hereditary enemy. Not in a single place could the bird of ill omen stir without
Luther, who had his eyes about him everywhere, at once noticing who was afoot. `Therefore
let no one doubt that where a fire breaks out, so that a village or a house burns down,
there always sits a little devil who keeps blowing into it so that it may grow bigger.
Likewise when someone dies of pestilence, or drowns, or is killed in a fall, the devil
does it.' It always gives an extraordinary security to know whom one personally has to
deal with, in evil as well as in good. Living faith in Christ is always conditioned by an
immediate sense of the devil's dangerous nearness. `A Christian should know that he sits
in the midst of devils, and that the devil is nearer to him than his coat or shirt, yes,
than his own skin, that he is all around us, and that we must therefore constantly be at
his hair and fight with him.' --- With complete warrant, therefore, David Strauß also
reproaches the more recent theologians for how frivolously they act in setting aside the
doctrine of the devil while still wanting to hold on to the historical Christ. `The whole
idea of the Messiah and his kingdom,' it says, `is as little possible without a kingdom
of demons with a personal head as its counterpart, as the north pole on a magnet without
the south pole. If Christ came to destroy the works of the devil, then he would not have
needed to come if there were no devil; if there is a devil, but only as the
personification of the evil principle, --- fine, then
% --- p. 344 ---
\opage{344} a Christ as an impersonal idea is also enough. Here the pious narrowness that fears
losing Christ along with the devil has seen far more rightly than Schleiermacher with his
postulate that belief in the devil ought in no way to be set up as a condition for faith
in Christ.' ''

``With the orthodox blockheads of the old Lutheranism there is naturally nothing to be done
here. It is to you, you priests of cultivation, who are so deft at furbishing faith and
polishing the rust off Christianity; it is to you, you theologians of progress and the
new science, who surely intend at the very first opportunity to bring revelation to
reason for us; it is to you, you lofty, liberal spokesmen of Gospel criticism, that this
speech is addressed. What, honestly speaking, is your opinion of the significance of
thoroughgoing criticism? You fraternize with the Free Thinkers when they attack the
doctrine of the devil; you turn your backs on them scornfully when they attack the
doctrine of redemption. Now what is this haughty air of superiority supposed to mean? If
you have really, in your quiet minds, already given up the cause of Christianity, and
wish only, for certain reasons, to make a quiet, almost imperceptible transition to
something better; why then will you prostitute yourselves by speaking so loudly against
the very consequences of the idea of reason? But if, on the other hand, it is your
intention to cast away the hell of revelation in order to enjoy the remnants of
Christianity in a better condition; then modern science shall have a word with you. You
must not think that negative criticism lacks all sobriety because it sometimes has a
little fire in its head; you must not rely on its always harming its own cause by
youthful rashness, as though it could tear the mask off you with one jerk; no, it can
also go to work soberly and calmly; it will keep picking at all the threads of the
doctrinal system until at last it has unraveled your web. Here straightaway are a few
% --- p. 345 ---
\opage{345} points for closer consideration. 1) You smile at Luther's fantasies; you blush with
embarrassment when one speaks to you of the Church's personal devil. But where does that
come from? It comes, surely, from your sound reason's being unable to accept that the old
story of the fall of the evil spirits from God should count for more than a myth. You
give up the devil and all the evil angels, because their supernatural existence seems
unnatural to you. You give up the good angels, and for the same reason. But why then will
you not go on, since you are on so good a path? If hell must disappear because it is
unnatural, the good angels must disappear because they are supernatural; why then do you
continue to preach about Christ's seat at the right hand of the Father? Where, then, is
the right hand of the Father? `Everywhere!' Quite right, an empty abstraction! Where, then,
does the clarified Christ sit? `Everywhere!' So it goes with the clarified Christ as with
the clarified devil: the clarification dissolves them both into the emptiness of
abstraction. 2) You reject the devil because your sound reason is outraged at the thought
that a supernatural evil will should act upon the will of man. You see very well that the
evil that is in this world can well enough be explained by sufficient grounds without the
devil's cooperation. Fine, praiseworthy progress! But since you have now in so reasonable
a manner said A, why will you not also say B? You assert the autonomy of the human being
with respect to evil, and in that you do right; but why can you not bring yourselves to
assert the same autonomy with respect to the good? The latter is, after all, at bottom
far more reasonable than the former. Evil is unnature (the negative), the good is nature
(the positive). If the human will, merely on account of a certain natural unnature, can
plunge itself into so much corruption; then it must surely also be
% --- p. 346 ---
\opage{346} able, with the help of the positive nature of the idea, to save itself out of corruption.
3) You now speak so readily of the progress of science, and in particular of criticism,
in recent years; you do, after all, also give lectures, and acute and tasteful lectures
at that, on the perfectibility of Christianity. This is in order so far as it goes, and
entitles us all to the fairest hopes. But can you not, or will you not, then see where
all this progress must necessarily lead in the end? The personal devil disappears; you
applaud it. The personality of the Holy Spirit disappears; you almost concede it.
Christ's visible ascension disappears; you pretend not to notice it. But if anyone now
hints that God the Father himself must necessarily bow to the same consequences, and
consequently disappear as well; then he may be sure of your very highest disfavor. When
the talk is of a hell beyond, you will gladly hold to immanence; but when science then
refers everything to the immanent heaven, you at once begin to peer upward after the
beyond and to reason at length about transcendence. What would Luther say if he rose from
his grave and heard you, his most reverend disciples, discoursing and philosophizing in
this way about immanence and transcendence? He would probably say something like this:
`Transcendence there, transcendence here! let me now hear what you have in your heaven,
and where you put your hell!' --- Only posit a God the Father, and you thereby at once
posit a supernatural nature, you show us the glory of the Son in the heaven of the
angels, you open to us a glimpse into hell. Deny the supernatural, plain and simple: you
then deny not only Satan and his angels, but you deny at the same time both the Father
and the Son.''

``Here you now have our scientific letter of defiance, or, if you like, our amicable
message and greeting. Everything
% --- p. 347 ---
\opage{347}%
depends on which party you choose. Will you go back again, or will you go forward? In the
middle of the road you dare not remain standing. You can defy reason and go over to the
supranaturalists; you can pay homage to reason and go over to us. But if you linger on
the Middle Way, if you try to damp the sparking thoughts with the ashes of tradition, if
you perhaps attempt to quench the firebrand of enlightenment by dipping it in the oil of
half-faith: then the negative spirit shall swiftly come upon you, and its criticism
chastise you, and the flames of its wrath illuminate you, and consume you all, together
with your --- Christian remnants.''

For flesh and blood it is a dangerous thing to have dealings with the spirits.

{\small Gospel of Matth. Ch. 4. V. 1--11.}

\begingroup\large
Then was Jesus led away by the Spirit into the wilderness, to be tempted of the devil.
And when he had fasted forty days and forty nights, he was at last an hungred. And the
tempter came to him, and said: If thou be the Son of God, then say that these stones
become bread. But he answered and said: It is written: Man liveth not by bread alone,
but by every word that proceedeth through the mouth of God. Then the devil took him with
him into the holy city, and set him on the pinnacle of the temple, and said unto him: If
thou be the Son of God, then cast thyself down; for it is written: He shall give his
angels charge concerning thee, and they shall bear thee up in their hands, lest thou dash
thy foot against any stone. Then said Jesus unto him: It is written again: Thou shalt not
tempt the Lord thy God. Again the devil took him up with him onto an exceeding high
mountain, and showed him all the kingdoms of the world, and the glory of them, and said
unto him: All this will I give thee, if thou wilt
% --- p. 348 ---
\opage{348} fall down and worship me. Then said Jesus unto him: Get thee hence, Satan! for it is
written: Thou shalt worship the Lord thy God, and serve him only. Then the devil left him,
and behold! the angels came to him, and ministered unto him.
\par\endgroup

Half-criticism, which fears that the objectionable in this event will offend the present
age and have a harmful effect on the health of the cultivated, hastens to give an
enlightening interpretation, lays its hand on its heart, and assures us most urgently
``that the whole story, by God, has really got in here through a pure mistake.'' ---
``Either Christ, when he fasted in the wilderness, had some restless dreams; or else he
did not fast in the wilderness at all, but merely wanted, in a parable, to make vivid to
the disciples what temptations he in particular had to overcome who would found the
messianic kingdom. The disciples then completely misunderstood this parable, and
afterward related it as a historical occurrence. Many other and learned opinions could
still be adduced; but when all these learned opinions have been heard, everyone is, after
all, free to think what he himself likes, and that is precisely what is edifying.'' ---
This way of seeing certainly deserves every appreciation, if judgment must at all costs
be passed in spiritlessness. But world-conquering unbelief, meanwhile, has its spirit
just as well as faith. Both agree that the understanding of this story is of
thoroughgoing significance for the conception of the Gospel history; both agree that the
wilderness of temptation is a fitting battlefield, where they can best meet as thought
against thought, will against will, and try their strength against each other.

\emph{The negative spirit.} ``What, then, is this temptation supposed to serve? Surely one
will not ascribe a redemptive value to it, nor say that God first had need to put Jesus to
the test. The forty days' fast, too, is somewhat
% --- p. 349 ---
\opage{349} strange. One really does not comprehend how Jesus, after 6 weeks' lack of all
nourishment, can still hunger, and has not long since died of hunger. `He fasts at the
prompting of the Spirit.' But the Spirit of God must surely have known best that Satan
would attack him precisely by means of this fast, and would take the hunger it produced
to his aid in his temptation. It is, after all, a most unreasonable way to challenge the
tempter, a foolhardiness that does not become even one who is most sure of himself.''

\emph{The faith of the Gospel.} Do not take offense at Him who wanders there through the
desolate places, as though he sought rest. It is no melancholy brooder who wants to bury
himself in the solitude of the wilderness: see, it is the Son of Man, who is driven by the
Spirit of God. He comes straight from the river Jordan, from the baptism of consecration,
and is preparing to carry out his work. Do not take offense at him because he fasts so
strictly and so long; a person's time of preparation is earnest and strict as a time of
fasting: as the work is, so must the preparation be. To fast for 40 days and 40 nights
certainly exceeds a human being's powers; but to redeem a world from sin and death surely
also exceeds a human being's powers. --- Fear not that hunger will waste away this faster
in the days of preparation: the Spirit, the Spirit of God, which drove him out into the
wilderness, will nourish him in the wilderness until the trial is ended. Fear not that he
will now become entangled in illusion and ensnared by the tempter's conjuring tricks: the
Spirit remains over him, the Word proclaims itself in him, that he may know himself, and
choose himself, and fulfill all things in obedience. Fear not that the Son of Man will
burn up in the fire of trial and lose the sense of nature: \emph{when he has fasted the 40
days and the 40 nights, he does after all hunger at last,} --- a real human being in
flesh and blood.

\emph{The negative spirit.} ``There is no personal
% --- p. 350 ---
\opage{350} devil. But even if there were a personal devil, he still could not show himself
visibly; and even if this were the case, he would hardly have behaved as he is supposed to
have behaved on this occasion. If the devil knew Jesus' higher nature, how could he then
have thought of trying to seduce him; but if he had not already at that time seen a
higher being in Jesus, he would certainly not have gone to the trouble of revealing
himself to him.''

\emph{The faith of the Gospel.} The spirit of negation cannot bear that the most blessed of
all possibilities should be realized here on earth. The prince of this world is an
all-demolishing understanding, which always reckons on the cunning of self-will and the
infinite disproportion of human striving to the highest good. But the fallen spirit, with
all his dazzling shrewdness, is still no all-knowing spirit: he sees the Messiah, but in
his own light; he sees in him a manifestation of the powers of the holy kingdom, he sees
that these supernatural powers are here placed at the disposal of a single human being;
he reckons that such a human being may grow dizzy when he catches sight of the world and
learns to use his powers. The fallen spirit sees here in the Messiah a dreaming
possibility, a sleepwalker who will surely fall if one only knows how to fall upon him
rightly, if one suddenly calls out his name: the fallen spirit sees that the highest and
holiest --- is at stake; so he follows his impulse and joins in the game.

The Gospel knows that the devil is spirit, and that spirit cannot be made sensible in any
outward form: the tempter therefore cannot be seen either; he can only be heard. The
tempting word is each time calculated for the momentary state of soul of the one tempted,
insofar as this can be determined by an involuntary impression of the world. The Savior,
who has fasted so long, awakens as from an
% --- p. 351 ---
\opage{351} infinite self-absorption; finitude asserts its claims, the Son of Man finds himself in
the midst of the wilderness and feels his hunger; then the spirit of darkness draws near,
we hear the tempter's voice: \emph{if thou be the Son of God, then say that these stones
become bread.}

A shrewd word, put in at the right time! This summons has all the marks of a temptation.
It seems, after all, so well founded: how degrading that a Son of God should feel a want,
as though he were the most insignificant of men; how unworthy that he who is preparing to
save a whole world and to bestow on the hungry race the bread of true life should now in
reality lack the bread with which he himself might still his hunger! The summons seems,
moreover, so innocent, and in itself insignificant: ``if thou be the Son of God, whose
mere word can create bread in the wilderness, then make a trial at once, it can do no
harm: you must, after all, make a beginning some time.'' And yet this summons conceals
within it all the consequences of corruption. It cannot, after all, be complied with
without the wonder-working power having to enter the service of natural appetite and
sensuality; but then --- the kingdom of heaven perishes. {\large Man liveth not by bread
alone, but by every word that proceedeth out of the mouth of God.}

\emph{The negative spirit.} ``Granted that the first temptation is quite well contrived,
but what good is that, when the evangelist nevertheless prostitutes himself in trying to
set forth the second? The line he has here put into the devil's mouth betrays clearly
enough the clumsiness of its inventor. When the first temptation, as Strauß has it, did
not succeed, the devil, as a clever tactician, ought to have had an even more alluring
temptation in readiness; but instead of this we now find (in Matth.) the proposal of the
neck-breaking undertaking of throwing oneself down from the pinnacle of the temple.''

% --- p. 352 ---
\opage{352} \emph{The faith of the Gospel.} How the negating spirit of modern criticism must
sometimes wrap itself in spiritlessness so as not to betray itself! It seems deliberately
to overlook the main point of the second temptation. The tempter surely knows what he is
doing. It is not in order that the one tempted should now, like any other person, throw
himself down from such a height and end his life by suicide, that he has, in the power of
spirit, carried Jesus with him to the holy city and set him on the pinnacle of the
temple. The shrewd design aims precisely at the Messiah's having the angels called, that
they may bear him up in their hands, lest he dash his foot against any stone. If
sensuality is denied, fantasy must surely take all the stronger a flight. The second
temptation is, to that extent, also far more fantastic than the first. ``I see that you,
O Son of God, despise bodily necessity; you breathe only for the spirit's lofty goal, and
that I must find praiseworthy. But if you are now really Lord of the spirits, and if your
spirit has the right swing, then you must show it in deed. We spirits heed not the laws
of gravity, any more than the bounds of space: such things apply only to matter. Well
then, make a trial at once! One is so light when fasting; you can fly through the air.
Fear not, cast thyself down, see, it is written: \emph{he shall give his angels charge
concerning thee, that they shall bear thee up in their hands; a Messiah dare heed no
outward hindrance, he dare not dash his foot against any stone.}''

Must we not recognize in this the wild, self-willed flight of the vainly striving spirit,
a flight that only too often causes the chosen ones of humanity, alas, to be seized by
dizziness, and the great ideals to dissolve into air? --- \textbf{Thou shalt not tempt the
Lord thy God.}

\emph{The negative spirit.} ``No, now it gets better and
% --- p. 353 ---
\opage{353} better. The noble criticism of the Middle Way still has a shred of modesty left. It at
least lets the tempter nicely hide himself in the fog of myth, and speaks a little
vaguely into the fog about the second and third stages of the temptation, about the
location of the pinnacle of the temple, and about the high mountain that is found nowhere
in geography. But this blind literalism is priceless: no locomotive can be more
consistent when it comes to going straight at the thing. We have, to be sure, heard often
enough that the wisdom of the Gospel is to be foolishness to the world; but we would never
have believed that this hidden wisdom would in the end announce itself as an open
stupidity. Yet that is how far one goes on this point. It is, after all, by Strauß! an
utterly stupid demand that the Gospel finally has the tempter make. `All this will I give
thee, if thou wilt fall down and worship me.' --- `To worship the devil by falling at his
feet,' says Strauß, `was a request which, however much might be gained by it, every pious
Israelite must have rejected with abhorrence.' ''

\emph{The faith of the Gospel.} We by no means dare deny that the very highest shrewdness of
the God-denying spirit must finally end in stupidity. But the demonic stupidity does not
on that account cease to show itself as shrewdness. The last temptation has in our Gospel
so peculiar a stamp that a link would certainly be missing if it were dropped. For
whereas the second temptation aims at intoxicating the Messiah with false, airy ideals,
the third sets out to catch him in the real snare of worldliness. Yes, the wilderness
itself echoes the tempter's speech: ``I see that you, O Son of God, are no fantastic
enthusiast. Your sound, strong spirit works for real greatness. Yes, that I can
understand, so it ought to be. Let us but hold fast to concrete life. But where do we
find full concrete life except in the world? Yet --- you have no
% --- p. 354 ---
\opage{354} world: until now you have lived so narrowly! So I will lead you on and lift you up onto
the mountain of contemplation. Here you gain an overview, here you descry a goal worthy of
your striving: all the kingdoms of the world and their glory! --- Now you know who I am.
--- \emph{All this will I, the prince of the world, give thee, if thou wilt fall down and
worship me.}''

As the temptation thereby reaches its height, the tempter himself is as if infatuated by
his own jugglery; as he would breathe his world-poison into the pure soul of the Son of
Man, he is struck by the lightning of the Word: \textbf{Get thee hence, Satan, for it is
written: Thou shalt worship the Lord thy God, and serve him only.}

So the spirit of darkness sinks --- down into its darkness; but the heavens clear, God's
angels now come to minister to the Son of Man.

In relation to the temptation story it becomes quite evident to us what it really is that
the faith of the Gospel has to hold to. Were faith to be conditioned by learned researches
and metaphysical investigations, then with respect to the account before us there would
be questions enough to answer. How could Christ fast for 40 days and 40 nights? How could
the devil carry him to Jerusalem and lift him up onto the pinnacle of the temple? Where is
the mountain from whose summit one has a view over all the kingdoms of the world and their
glory? Is the concept of a personal devil not a pure contradiction? etc. --- That such
captious questions must be able to embarrass even the most thorough theologians, everyone
sees. But what does faith now say to all this? Nothing. It lets exegetes and
archaeologists, critics and dogmaticians entertain one another with their learned and
metaphysical problems as much as they like, and meanwhile holds to the Gospel in all
simplicity. Yes, faith is really simple, so thoroughly simple that it does not even quite
dare to
% --- p. 355 ---
\opage{355} speculate on what in all the world Christ can have eaten in the 40 days when he fasted,
whether he perhaps lived all that time on ``roots and potatoes''? And it fears still more
to have dealings with any theosophical devil. Should a profound exegete, however, wish to
stamp this refusal as a regrettable sign of unscientific obstinacy, then we must, on
behalf of politeness, convey our respects to the profound one, and ask him to excuse the
faith of the Gospel for not yet having come so far that it can comprehend the wonder of
revelation. Should it, against all expectation, come upon someone or other who knew
anything proper about the law by which the supernatural acts upon the natural, then it
hereby solemnly promises at once to send word both to Messrs.\ exegetes and archaeologists
and to Messrs.\ critics and dogmaticians, and let them know. But as long as such a one has
not been found, as long as that much-inquired-after fundamental law is still an unknown X,
it begs to be excused, and otherwise remains in all simplicity with the Gospel narrative.

``But,'' object the interpreters, ``if the faith of the Gospel wants in this way to be
exempted from all scientific trouble, and to take it upon itself, without our guidance,
to interpret the sacred text; what does it then become but an arbitrary conception of
certain favorite expressions, or a spiritless clinging to the letter of the narrative?
The temptation story can serve here as an example. For if one wants to hold to the words
and sentences of the account, one at once becomes uncertain whether to follow Matthew or
Luke, since the narrative proceeds in a different order in the former than in the latter.
While Matth. observes a gradual ascent, both in the tempting summons and in the manner in
which the tempter is rebuffed, Luke on the other hand has had regard to the natural
transition from one place to another, from the wilderness to
% --- p. 356 ---
\opage{356} the mountain, and from the mountain to the city. How fruitful such comparisons can
become for the right understanding of the sacred books, thinking Christians will readily
be convinced. A cursory comparison between the two temptation stories shows clearly
enough that we do not here stand on firm historical ground; for both arrangements cannot
possibly be equally true. But when a careful discussion thus leads us with inner
necessity from the historical foreground over into the more remote domain of the
spiritual, is it not then well grounded in the nature of the case that religious faith
must let itself be guided by scientific interpretation?''

Faith's answer is a double one; this double answer sounds with different emphasis. The
faith of the Gospel must most decidedly decline every guardianship, whether it be offered
by philologists and archaeologists, or by exegetes and dogmaticians. Religion is no
learned matter, and can never become one. The only demand faith makes of the learned
interpreters is for a good, vigorous and, as far as possible, literal translation of the
original text of the sacred books. The upbringing, instruction and cultivation that
otherwise condition the right appropriation of the evangelical content fall within
another sphere, and have nothing to do with the learned enterprise. It is necessary to
bring out this principle, not because it is flatly denied in the theological schools; no,
on the contrary, because it is conceded only too readily, namely as --- a figure of
speech. In the theological compendia that treat of the relation between the School and
life, it seems, to be sure, at the beginning as though the stream of life were now to be
allowed to well up from its own springs, as though true faith were now really exempted
from all scholastic burdens and extraordinary impositions; but the fulfillment does not
always answer to the promise: before one knows it, the figures of speech change color,
and the whole looks as though Christian
% --- p. 357 ---
\opage{357} %
knowledge of the truth nevertheless depended, if not exactly on the entire Academy of
Sciences, then at any rate on the prodigiously learned development of all the disciplines
of the School. Against this misunderstanding the faith of the Gospel hereby lodges its
serious protest, and this is its first answer. --- As regards the second answer, the
style and tone can be much softened. Insofar as one has in view a humane and higher
cultivation, the Believer will readily concede the necessity that those who are to govern
the Church or speak the Word in the congregation should as a rule go through a
theological preparatory school. As regards learning and science, the Believer is so far
from wanting to abolish any discipline or suppress any particular tendency that he must
rather take a lively interest in the scientific metamorphoses that the ecclesiastical
forms of doctrine gradually undergo. On the other hand, it is a pure mistake if anyone
thinks that faith itself should meddle in these disputes and be zealous enough to take
sides for or against. Whether theology can at a given time get the better of philosophy,
or vice versa; whether a Free Thinker or an orthodox loses his balance; whether it is
perhaps a skeptic or a dogmatist who now sinks to his knees in the arena of debate,
religion troubles itself little about that. As an attentive spectator the faith of the
Gospel then sits at the barrier and looks with its living, gentle eyes upon all the
mighty strife, without letting itself be provoked; secure in its simplicity, it calmly
follows the course of the strife, and has its own thoughts in silence. Should any of the
great powers approach and make as if to violate it, it shall not be without all
protection either. There is a not inconsiderable force that can take the field every
time scientific arrogance would bring faith into danger. Christian reflection then sends
out a small observation corps for the security of the religious preserves, and posts
everywhere a multitude of
% --- p. 358 ---
\opage{358} small but vigilant reflections as on lookout, to watch how the most learned theologians
and acute critics deal with the Gospel.

With these elucidations we now turn again to the temptation story. There is naturally no
passage in the New Testament that the interpreters cannot explain. The more difficult
the passage, the greater the artistry the explanation presupposes. Of artful explanations
there are in turn two kinds. The lower or simpler we call a \emph{re-explanation}. It is
used when one wants to give the biblical narrative a more elegant and agreeable tinge, to
remove stains, and in general to win from the presentation the side that most appeals to
consciousness. Such a reinterpretation or modification is fairly easy to carry out. Far
more difficult, by contrast, is the other and higher kind of art, namely genuine
\emph{explaining-away}, whose successful execution may be regarded as the highest triumph
of hermeneutics. This art consists in the interpreter's being able, without exactly doing
open violence to the text, to remove an entire event, so that a historical narrative at
last comes to look as though it narrated nothing. As the explanation reaches its highest
clarity, it suddenly becomes quite dark, and one sees nothing at all; that is to say, the
fact goes into the light: pure light is like pure darkness. How many times has such an
explaining-away not been applied with success to the story of the temptation! When the
exegete sets about performing this feat, he is always well prepared: first he feels his
way in some introductory turns of phrase, then he adjusts the layout, then he picks at the
expressions, then he enumerates the difficulties, and at last fixes attention on the
point he himself has chosen as his main point. Thus one remarks: ``on the devil's side,
the flatness in his arranging and introducing
% --- p. 359 ---
\opage{359} %
of the individual assaults; on Jesus' side, the patience in continuing these negotiations
and roaming about with the spirit of darkness, instead of breaking off at once and
withdrawing from contact with the unclean.'' By carrying these hints a little further,
one gradually gets farther and farther away from the first, immediate impression. Now the
different accounts are compared, now the essential is separated from the inessential, now
the text's clothing is transformed into something entirely inessential: then the devil
disappears, then Christ disappears, then the whole story disappears.

Faced with such an artistic performance, the faith of the Gospel must now above all guard
against every passionate outburst. It does little good for indignant orthodoxy to start
slamming the church door and ringing the great bell: a science is, after all, a science,
and an art an art all the same. On the other hand, the Believer does well to take up the
Gospel and see with his own eyes whether there really is so considerable a difference
between Matthew and Luke that any suspicion against the historical character of the
narrative could be based on it.

Had Luke related the temptations in the same order and with the same words as Matthew,
would recent criticism not then have claimed that the story itself had only rather weak
warrant, since the one evangelist had presumably copied the other? Now, on the other
hand, one will surely be compelled to assume two sources, or, if you will, two
independent uses of the same source: the one testimony then supports the other. ``But
what of the mutual conflict between the accounts?'' There is no conflict at all, when we
look more closely. The discrepancy cited could be invoked as an actual contradiction only
in the case that both evangelists had given the definite temporal sequence of the
temptations with equal definiteness.
% --- p. 360 ---
\opage{360} But if we consider the individual transitional expressions in Luke: V. 5. ``And the devil
led him up onto a high mountain.'' V. 9. ``And he led him to Jerusalem.'' V. 13. ``And
when the devil had made an end of the temptation \dots'', then it shows clearly enough that
no particular emphasis at all can fall here on the chronological order. Both accounts,
though mutually independent, agree so well with each other that the art of
interpretation, if it otherwise lay in its interest, would surely know how to use such an
agreement for a thorough defense of the credibility of the story. But the supernatural,
which is here a stumbling-block to the critic, cannot, after all, be a stumbling-block to
the Believer. The faith of the Gospel is therefore within its rights when, without letting
itself be misled by any explaining-away artistic performance, it holds in all immediacy
to the historical fact. --- ``Whence, then, does the Believer know that there is a
personal devil?'' Among other things, from this temptation story. For it is plain that
the tempter could neither act as an arbitrary power nor speak with a voice of his own
were he not a being free and self-conscious in himself. From every other inquiry the
origin of evil will always know how to withdraw itself from view and wrap itself in
obscurity; only He who is himself the Light of the world is able to draw the secret of
wickedness forth into the light and reveal --- the spirit of darkness. We are hereby led
to the following conclusions:

\begin{center}
  \textit{1.}\quad\emph{The temptation story cannot be a\\
  myth.}
\end{center}

It is the distinguishing mark of myth to leave every fundamental principle in ambiguity,
and to confuse genuine personality with mere personification. Myth personifies the good
just as well as evil, light just as well as darkness, Ormuzd just as well as Ahriman.
Revelation, on the other hand, is recognizable by this,
% --- p. 361 ---
\opage{361} that it unmasks the ambiguity in the essence of evil as of good. If the Devil is nothing
other than a Jewish Ahriman, then Christ is consistently to be regarded as a Jewish Ormuzd. ---
To conceive the Tempter as a mythical personification is to take up a fundamental
ambiguity into the sphere of the Gospel, and to occasion a confusion that obliterates
the stamp of revelation.

\begin{center}
  \textit{2.}\quad\emph{The temptation story has nothing to do with the\\
  distorted image of superstition.}
\end{center}

The notions of the popular imagination about the Devil's manifold manifestations and
Mephistophelean adventures find no justification whatever in the sacred history. So far
is it from our being given here any hint that the Devil should have personal dealings
with all sorts of eccentric individuals, that he even keeps himself hidden from all the
persons who act in the Gospel, Christ alone excepted. We can perhaps best illuminate
this point by looking at the two characters on whose fall the Tempter seems to work in
particular, namely Judas and Peter.

Without a demonic middle term the relation between Christ and Judas will remain
wholly incomprehensible. Nor, then, has modern science neglected to direct its
attack against this point. After having stressed the discrepancy that supposedly
exists between the first three Evangelists and John concerning Jesus's foreknowledge of
Judas's betrayal, Strauß continues thus: ``According to this (namely what is reported
Joh. 6, 64. 70), Jesus would thus have known from the beginning of his acquaintance with
Judas that the latter would betray him; and not only had he foreseen this outward event,
but, since he knew what was in man (Joh. 2, 25), he had also seen through Judas's
motives, namely, that he would
% --- p. 362 ---
\opage{362} commit the deed out of avarice and greed for gain. And yet he is supposed to have made him
treasurer, i.e.\ to have put him in a post in which his inclination to procure gain for
himself in every way, even an impermissible one, must find the richest nourishment? He is
supposed to have given him occasion to become a thief, and, as though deliberately, to have
hatched himself a traitor? Viewed from the economic standpoint alone, does anyone entrust a
purse to one who he knows will steal from it? And now, as regards the pedagogical, who puts
the weak man in a place that tempts his weakness so constantly that it is to be foreseen
that sooner or later he must succumb? But in truth Jesus did not thus play with the souls
most closely entrusted to him; he did not thus show them the opposite of what he himself
taught them to pray: \emph{Lead us not into temptation} (Matth. 6.
11), that he could have appointed Judas, of whom he knew beforehand that out of greed for
gain he would become a traitor, as treasurer; or, if he did make him this, he cannot have
known that beforehand.''

For the faith of the Gospel, which nonetheless does not let itself be misled by exegetical
pettifogger's tricks and impudent reasonings, the relation appears in quite another
light. Were Christ's foreknowledge in conflict with the free development of character,
then the same would of course have to hold of God's omniscience. Modern criticism proceeds
from the presupposition that God's providence is incompatible with man's freedom. The
Believer proceeds from the opposite presupposition, and must therefore, as regards Christ,
again and again recall Simeon's words: \emph{This one is set for the fall and rising
again of many in Israel, and for a sign which shall be spoken against, that the thoughts
of many hearts may be revealed.}

As the Revealer, Christ has then also willed to call Judas,
% --- p. 363 ---
\opage{363} and to take him up among the chosen. We are surely entitled to presuppose strong and
great dispositions of character in this disciple. Judas seeks Jesus because he is
spiritually moved, and feels an urge to test the powers of the spirit. Jesus calls Judas
because in this man he sees the unusual natural force that at one and the same time
conditions an extraordinary rising again but also makes possible a terrible fall. Which of
these possibilities is to become actual in time is not predetermined by any compelling
necessity: the crisis is laid in the freedom of the will. On Judas Iscariot a higher
power will and must reveal itself; then the spirit will come, but --- which spirit? --- It
is a dangerous thing for flesh and blood to have dealings with the spirits!

As Judas now resolves to follow Jesus, the secret powers that in all eternity work
against one another begin to work in him. He is to recognize in the Master's person the
Son of God, the true Messiah. So works the Spirit of Christ that he ought now to receive
the word in all humility, and appropriate the promise in simplicity of heart; but --- in
doubt and misunderstanding a tempter breathes: deep in selfish hiddenness lurks the spirit
of darkness. Thus this disciple does not begin, like the others, with glad confidence and
open-hearted adherence to the Lord and Master; but he begins in the anxiety of tension, in
concealed passionate unrest, in sinister ambiguity. For Judas, too, the idea of an
approaching Messianic kingdom dawns. Even a Judas Iscariot has an eye for Jesus's personal
superiority. --- ``When He lifts up his voice: what authority in his words! The multitudes
listen. --- When he stretches out his hand: what mighty works! the palsied man stands up,
the dead man rises on the bier.'' --- ``If only one had such a word, and could do such a
work!'' --- ``But can it then really also be the Messiah
% --- p. 364 ---
\opage{364} himself, this peculiar man who roams about here from town to town, heals cripples
indiscriminately, and wastes his words on every Tom, Dick and Harry? --- No overview, no
calculation, no plan!'' --- ``Yesterday in the wilderness, yes, that was a wonder! We
disciples went about and handed out the loaves; the men regarded us with holy
astonishment; the women looked at us as at the Almighty's angels; yes, that was a sign for
the people, a sign for the hungry: they ate and were filled. But --- then all at once the
whole thing was over. The Incomprehensible One withdrew into deep silence, the crowds
dispersed; not a word about when we should gather again.'' --- ``Today, then, the people
have found him again, though with much trouble, him about whom all are asking; the crowd
gathers, perhaps even ready to take him by force and make him king. --- But what will come
of this after all? Everyone is in the greatest expectation: then he speaks, --- yes, he
speaks, --- like one who wants to parody his own good deed.''

\textbf{Except ye eat the flesh of the Son of man, and drink his blood, ye have no life in
you.}

--- ``Oh yes, such a speech will do its work. The assembly murmurs, and breaks up. See how
they now go away again, one cluster after another. Naturally! Who can bear to listen to
such a speech? But --- the spirit he does have, and --- for the spirit he must, after all,
wait.''

If we presuppose in the Iscariot a mood corresponding to the reflections suggested here,
then the scene mentioned before (Joh. 6, 67---71) will also make a new, original
impression.

The offended crowd has withdrawn: it is so quiet around the Master. Here now no one is
present but the few who cannot leave Him who has himself called them His own, the few who
cannot go away because they do not know to whom they should go. ---
\emph{Lord, to whom shall we go?
% --- p. 365 ---
\opage{365} thou hast the words of eternal life. And we have believed and known that thou art the
Christ, the Son of the living God.} --- So Peter testifies in the name of all, yet there is
One who does not heed his testimony. So the word of confession will spread its clarity over
all the disciples; yet there is One whose countenance becomes --- still darker. So these
chosen ones are to join together more inwardly in faith, in hope, in the true union of
love; yet there is One who remains standing in his own thoughts, present and yet ---
absent, unmoved and yet --- transformed: he does not believe, but he broods --- on a way
out; he does not hope, but he waits --- for the spirit's assistance; --- he does not love,
but he dedicates himself --- to attempting the uttermost. Then the Lord's voice sounds:
\textbf{have I not chosen you twelve, and one of you is a devil?}

In this portent we hear at once a prediction and a warning. He who knows what is in man,
and knows the impression that hard saying must have made on the disciples, sees also what
at this moment is stirring in Judas Iscariot. To that extent, then, the Evangelist is right
when he says (V. 64): ``Jesus knew from the beginning who they were that believed not, and
\emph{who it was that would betray him.}'' But since the prediction here concerns a deed of
freedom, not of necessity, it must at the same time be regarded as a warning. The
prediction is therefore not set forth in the form of finitude; Christ does not turn
expressly to the individual disciple, he does not say outright: ``I know it, Judas, that
you are a traitor; I know too that you must necessarily be one, since you are possessed by
the spirit of darkness.'' But he sets forth the prediction quite generally; he speaks as
to all the disciples, and yet there is only One among them who knows within himself whom
the word concerns, only One who in this
% --- p. 366 ---
\opage{366} prediction must hear his judgment, as a terrifying warning: ``Beware, Judas, for here is
One present who is --- a devil.''

But what, then, should the disciple do who found himself struck by such a warning? Beat
his breast, go aside and --- weep bitterly. No, for that Judas felt himself too strong. He
would not let himself be crushed by a word; so he stayed where he was, and as he was,
outwardly cold and pale, so that he could hide the movement within.

From this point of view the reciprocal relation between guilt and fate is not difficult
to show. Had Judas fully known the voice that warned him, he would surely have come to
himself; had he fully known the Tempter in the dark promptings, he would surely also have
come to himself. But --- the light makes him blind. So he stands in obscurity; and when
one stands in obscurity it is always hard to discern and distinguish: so Judas cannot
rightly distinguish which is the Savior and which is the Tempter; so they run together
before his sight, so the Savior becomes tempter, and the Tempter savior, and alas! then
comes the heavy fate. --- Hardening is no empty act of arbitrariness: we determine as we
\emph{are determined}. But on the other hand it now also holds that no being can take the
free will by force: we are determined only as we \emph{determine ourselves}. Why does
Judas not fear his Lord's wrath? Because he does not \emph{will} to fear. Why does he not
understand the warning word? Because he does not \emph{will} to understand --- himself. It
is certain that the hardened will at bottom loses its freedom; but it does not therefore
lose its responsibility. The responsibility is that the free will has thrown its freedom
away by giving itself over into the power of evil, and thereby making itself the slave of
an alien will. The evil will, too, is spirit; but --- what a spirit!

\textbf{It is a dangerous thing for flesh and blood to have dealings with the spirits.}
% --- p. 367 ---
\opage{367} The negative spirit of criticism cannot, of course, put up with this conception, but will
naturally always keep to the reliable rule that a prediction is a prediction: ``either
Jesus expressly foretold that Judas was a traitor, or he did not foretell it at all.'' The
modern consciousness thinks thus: ``either one has a foreknowledge, or one has none.'' This
fundamental thought at once gives science an easier swing, and spares criticism many
troubles. If one can really get off so easily with the distinction between foreknowledge
and non-foreknowledge, then it is indeed no witchcraft to dissolve the whole evangelical
history. \textit{Dr.} Strauß, who never lets many-sided considerations be lacking, has not
neglected, either, to consider the relation between Christ and Judas from the
ingenious standpoint of private feeling and sociability. On this occasion he raises the
following interesting question: ``How could Jesus bear to have about him throughout his
whole public life a man on whom he knew all instruction was wasted? Must not every hour of
his intimate intercourse with his disciples thereby have been embittered for him?'' --- On
such criticism, however, the faith of the Gospel dare not waste its instruction. Instead
of dwelling on the Straussian question why Jesus did not simply dismiss Judas as a
disagreeable person, we must much rather consider the relation from the standpoint of
revelation, and raise the religious question: ``How is it conceivable that Judas, after
having felt himself struck by that rebuke, could bear to remain in Jesus's company?'' ---
This question breaks forth with irresistible urgency when we rightly try to enter into the
Gospel. That Jesus's personality was strong enough to bear a defiant disciple's hardening,
that no cloud of misunderstanding and no shadow of misjudgment was able to darken the sun
of love, that here
% --- p. 368 ---
\opage{368} %
there can be no talk at all of any human being's being able in a personal sense to embitter
or \emph{vex} the only-begotten Son of God --- this everyone must surely admit who will not
flatly deny the authority of the Revealer. No, on the Redeemer's behalf the question here
is not about pleasant human company; the question here is about the thoughts of hearts,
that what is hidden in darkness may come forth and be recognized in the light. --- But that
Judas could bear it, that this disciple, who after all must have felt within himself that
his state of mind could not possibly be hidden from the Master's gaze, that Judas himself
did not find his position intolerable, and seek the first opportunity that offered to
withdraw from the painful relation and slip away, is a circumstance that surely deserves to
be considered somewhat more closely.

First of all, then, we must notice that outwardly the relation, even after that portent had
been given, remained entirely unchanged. Strauß is certainly right when he maintains that
there is no trace in the Gospel of Judas's having been set behind the other disciples. In
the relation to the Lord, merely human partiality and individual favor are excluded. Judas
can thus maintain his place in the circle of the Twelve without anyone's attempting to set
him aside. Even if the other disciples had had an inkling that it must surely be the
Iscariot whom the Lord had meant as --- a devil: what conclusion could they then draw from
this? Was not Peter himself (Matth.\ 16, 22) punished with an equally harsh address? --- In
outward respect, then, the Twelve occupy the same position; in the circle of these
personalities it is everywhere equally far from periphery to center. Only in the light of
inwardness does the distance vary, and the difference make itself felt: Peter, James and
John come forward with a special precedence. But since this
% --- p. 369 ---
\opage{369} precedence is conditioned solely by the true understanding and personal appropriation of
the Master's words and works, each of the others had, after all, the hope that on the same
condition they could raise themselves to the same level. As regards the outer side of the
relation, then, there is nothing here that could induce Judas to withdraw. All the greater
weight must therefore fall on the inner side of the relation.

From this point of view we at once bring out the main proposition, \emph{that Judas, in
consequence of his ever-growing hardening, must gradually have entered into a demonic
relation to Jesus.} --- How painful, to know oneself in a personal discord with a superior;
how painful, to have to remain in the superior's company though the discord continues; how
painful, indeed how intolerable, to have to waste away under such a relation in all
silence, without its ever coming to an open break or a definite declaration. If we now add
to this that the superior whom the discord concerns is a man like no other, God's
only-begotten Son! then we get an inkling of what pressure must have weighed on Judas when
he really began to feel himself in secret conflict with his Lord and Master.

Judas must thus bear a superhuman burden, since he must immediately and daily bear a
personal discord with Christ himself. To bear such a burden, an extraordinary strength is
presupposed. --- ``What, then, is this Judas?'' --- A strong spirit, and yet --- a most
wretched man. The secret is that his unnatural strength is essentially impotence; the
secret is that the disciple who seems to master the superhuman is not master of himself;
the secret is --- an incomprehensible
contradiction.\footnote{The all-explaining criticism can, of course, easily explain this
secret along with all the rest. Should the contradiction become too strong, one has the
myth close at hand, and in the myth there is naturally a remedy for every predicament. If
one will not flatly dissolve the traitor into pure myth: very well, then rational
psychology is at one's service in every way. The Iscariot plays in the demonic, the modern
consciousness plays with the demonic: kindred spirits meet halfway. The traitor is on the
whole no unrewarding person to deal with. If a critical advocate with some talent will
only take up Judas's cause a little seriously, he can be fairly sure of a decent fee. An
agreement is easily reached. Judas supplies the advocate with material for an
incomparably interesting proceeding; in return the advocate procures his client a
resounding satisfaction: Judas keeps his strong spirit, his defiant pride, his grandiose
despair; what is base and vile in his character falls back, if not on Christ himself, then
at least on the apostles' false, erroneous report, and the evangelical Faust passes with
distinction.}
% --- p. 370 ---
\opage{370} In the absolute sense the Lord and Master does everything to educate Judas and save this
soul from perdition. To doubt this is to deny Christ. But in the sense of finitude nothing
is accomplished: the highest spiritual relation cannot be altered and cannot be adjusted by
half measures and provisional explanations. The divine word, which is sharper than any
two-edged sword, must pierce through, until it divides asunder both soul and spirit, both
joints and marrow, and judges the thoughts and intents of the heart. The highest is
offered: it becomes man's affair to use what has been entrusted. --- Judas is thus a real
apostle; Christ has, after all, called him, and Christ stands by his word. When the Lord
(Matth.\ 10, 5) sends out the Eleven, he makes no exception of the
% --- p. 371 ---
\opage{371} %
Twelfth. The Iscariot is shown the same trust as the others: he receives the same
consecration, the same blessing, the same authority and promise as all the others. But
neither, then, can he dazzle or deceive the divine oversight; the eye that sees into the
heart looks after him when he goes away, to see how he understands his mission; it looks at
him when he comes back again, to see in what spirit he has discharged his calling. It is
thus no mistake over which Christ complains when he utters these words: \emph{Have I not
chosen you twelve, and one of you is a devil?} But it is a terrible contradiction that
weighs on Judas.

--- ``Apostle? Devil? An apostle-devil! --- Quiet! who can know what follows? Quiet, quite
quiet! Might not just such a one be needed here?'' ---

Judas seeks refuge in hiddenness; but what a contradiction, to have to secure oneself in
hiddenness and yet stand face to face with the Revealer himself. This relation the disciple
cannot see through; he does not himself know what it is that at once repels and attracts
him. He cannot even decide with certainty whether his most secret thought is really hidden,
or perhaps already discovered and judged. He is stifled and tormented, but cannot say
definitely what it is: a fleeting thought, a dark prompting, something and yet ---
nothing. The outward is unchanged, and yet what unrest and strife within! What is it, then?
He is like one cast off, and yet in reality he has as yet committed no misdeed. If only one
of his fellow disciples would try to raise an accusation, then with defiant assurance the
Iscariot would step forward and ask: ``Lord, say, what have I done wrong? Have I not for
your sake left all things, just as well as the others? do I not preach the coming of the
kingdom, like the others? Do I not
% --- p. 372 ---
\opage{372} await the manifestation of your glory, as fully as Peter and all the others? what, then, is
my crime?'' --- It would surely strengthen Judas to meet such a summons; it would relieve
him to be able to conduct such a defense. But in truth it would not improve him. And so no
summons comes, no challenge: the silent one must bear his burden alone. Yet quite alone?
--- There is a spirit that works in secret; reserve gives strength, after all: the silent
one grows --- a silent growth. --- He will sacrifice everything; in selfhood he sacrifices
himself: what, then, has he to expect?

\emph{Verily I say unto you, that ye which have followed me, in the regeneration, when the
Son of man shall sit on the throne of his glory, ye also shall sit upon twelve thrones, and
judge the twelve tribes of Israel.}

--- ``Twelve thrones, twelve apostles, Twelve! So it will come out even after all: the
Twelfth is --- among the chosen!'' ---

It is not to be doubted that the Iscariot, as the powers of the spirit work themselves up
in him, will surely continue to hold fast to the expectation and to direct his gaze toward
--- the throne. If staking everything on a relation to eternity and demanding the highest
by virtue of a wonder were without further ado the same as being a Believer, then we could
surely say of Judas with full right that he not only had faith but was even capable of the
extraordinary by means of faith. Judas is no rapidly promoted Free Thinker who in the
greatest haste has disposed of the riddle of existence by stirring spirit and matter
together in a stroke of genius; he is no genteelly blasé wizard-doctor who, after having
despaired in philosophy and jurisprudence, in medicine and theology, now stands poised to
plunge into the intoxication of the moment in order to suck a draught of immediacy. No,
freshly and eagerly and vigorously Judas grasps after the spirit, as
% --- p. 373 ---
\opage{373} %
he feels himself grasped by the supernatural power. Nevertheless this being-grasped is not
faith; on the contrary! it is precisely to be determined as the inversion of faith; his
inverted faith is, so to speak, a demonic faith, for he is grasped, after all, in the
deepest egoism. If Judas could only find a Christ who answered to this faith, then he would
soon become stronger than Peter, then he would surely venture everything for the Messianic
throne; yes, then he would speak with the tongues of angels, and give his body to be
burned, so that he might then rise victorious in the spirit. But now Peter's Christ is,
after all, not --- such a Messiah; now the doubtful expectation is a secret torment; so the
strong spirit must exercise itself in small things, so the low comes forth --- with
cunning, so Judas consoles himself with the purse, and makes himself a thief for a penny.

When the evil principle is to take shape in aesthetic presentation, the imposing character
must be freed from every paltry pettiness and, so to speak, rise up on the cothurnus. In
this case evil may well arouse our hatred, but not our contempt. The poet must with a
certain care protect the evil power, so that it can set minds in the right tension, and
thrill the spectators with secret horror. It is otherwise with the ethical presentation.
Here one pays no heed to the theatrical effect of evil; one keeps simply to the actual. But
in actuality the evil that breaks all divine and human laws will naturally not always
conform to the aesthetician's law. The high and the low, the dreadful and the contemptible,
the grandiose and the wretched mingle with one another in the character that defies the
holy spirit of morality. This holds especially of Judas. --- How in the finite sense it came
about that precisely this disciple attained, as Strauß puts it, to becoming ``treasurer,''
of that the Gospel reports nothing. But if what has already been said
% --- p. 374 ---
\opage{374} %
about his outward position stands firm, then there is nothing to prevent his having well
been able to obtain such an office.

It is only a losing of oneself in the lowest considerations of finitude that enables
criticism to contrive its objections. If we set the matter forth in the light of
revelation, is it then not clear enough what moved Judas to grasp after the purse and take
upon himself the care of the poor? --- The Evangelist John, who keeps an eye especially on
Judas, has, in the words with which the Iscariot rebukes Mary at the feast in Bethany
(Joh.\ 12, 4---5), expressly brought out one of the lowest traits in the traitor's
character: \emph{this he said, not that he cared for the poor, but because he was a thief,
and had the purse, and bare what was put therein.} --- Pride and coldness, sly cunning and
paltry greed for gain can very well be united in the being that has now taken charge of
Judas.

In presenting to us the life of Jesus, the sacred history presents to us at the same time
the terms and conditions under which apostolic faith went through its becoming and
development. In this respect Judas provides a striking counter-image both to John and to
Peter. The evangelical faith of John is in inwardness; but while the Johannine inwardness is
silent because it is in unspeakable love, this silence is in turn one with the lovable
open-heartedness that breathes in Peter's confession. The demonic faith of Judas is also in
inwardness; but this inwardness is the reserve of hardening, is bound in hiddenness and
reflected in guile, and therefore neither can nor will reveal itself. In Peter's relation to
the Lord there appears a continual alternation of strength and frailty, a constant rising
and sinking: the disciple works his way forward like a man swimming in a stormy sea. Judas,
too, rises and sinks, feels his strength, and senses his
% --- p. 375 ---
\opage{375} impotence; but here the high and the low pass demonically over into each other: Judas is
like a heaven-storming Titan and yet a wretched, contemptible man, a strong spirit that
makes itself a thief for a penny.

With ingenious prolixity and rare thoroughness the theologians still keep reasoning up and
down, back and forth, in order at long last to find out what can really have moved Judas
to behave so badly toward his Lord and Master. ``Hier aber entsteht'' --- as \textit{Dr.}
Neander says --- ``die schwierige Frage, durch welche Beweggründe Judas zu einem solchen
Frevel sich bestimmen ließ.'' --- Without exactly daring to dispute with the
\textit{Summus Theologus} ``über diese schwierige
Frage''\footnote{\textit{Dr.} Neander (Leb.\ Jes.\ S.\ 573) does indeed think that the
Evangelist John was not altogether fortunate in finding out the ``motive''; in proof of this
he recommends, among other things, ``die scharfsinnige Schrift eines viel versprechenden
jungen Theologen, Jesus und Judas, von \textit{Dr.} Gustav Schollmeyer''..., and then puts
forward (on the occasion of Joh.\ 12, 6) the following remark: ``Johannes mußte sich zwar
wohl solcher Markmale erinnern, aus welchen er den Grund zu einer solchen Anklage gegen
Judas ableiten zu können glaubte, aber er konnte sich als Mensch irren.'' Very true, --- er
konnte sich als Mensch irren. But to whom, then, should one now rather keep: the old John or
the young \textit{Dr.} Schollmeyer? This too will surely in the end become ``eine
schwierige Frage''! On the one hand John, as an eyewitness cited for many years, undeniably
has a certain advantage over \textit{Dr.} Schollmeyer; but on the other hand ``der junge
scharfsinnige vielversprechende'' \textit{Dr.} Schollmeyer has, after all, made a new
discovery, and this discovery has in turn been ratified and cited by \textit{Dr.} Neander.}
the Believer can nevertheless in all
% --- p. 376 ---
\opage{376} simplicity well understand that the disciple who thus misjudged the call of the spirit,
and hardened himself against the Savior's love, could not become a true apostle, but had to
end as an instrument of the Devil; that the evil which smoldered in Judas had at last,
under the influence of the final catastrophe, to reveal itself in a despairing breakthrough;
that the unhappy man who seems to act with the coldest deliberation, and by the pettiest
calculation (30 pieces of silver!), nevertheless does not rightly know what he himself is
doing. But when the Believer has thus summoned up every psychological presupposition in
order to explain the misdeed, he suddenly becomes aware of the inexplicable, looks down into
the bottomless abyss of contradiction, into the darkness of despair, sees the secret of
wickedness, and --- shudders.

\emph{Now when the supper was held, and the devil had already put into the heart of Judas
Iscariot, Simon's son, that he should betray him (Joh.} 13, 2): then the great hour had
come. Judas too is a guest at the Lord's table. Here he sits at the supper, one of the
Twelve, well prepared in --- silence. Then the solemnity of the feast begins with a
cleansing (V.\ 3). \emph{When Jesus knew that the Father had given all things into his
hands, and that he was come from God, and went to God: he rose from supper, and laid aside
his garments; and he took a towel, and girded himself. After that he poured water into a
basin, and began to wash the disciples' feet, and to wipe them with the towel wherewith he
was girded;} then a word falls to Peter: Ye are clean, but not all. --- ``Who, then, is the
one who cannot become clean?'' The Lord knows it, Judas knows it, the spirit knows it; but
--- what a spirit! --- And Jesus looks at his Twelve, and speaks to them, and gives them a
heartfelt admonition
% --- p. 377 ---
\opage{377} %
(V.\ 16): \emph{The servant is not greater than his lord; neither he that is sent greater
than he that sent him. If ye know these things, blessed are ye if ye do them.} Then a word
falls, a word of Scripture, so stern and heavy: he that eateth bread with me hath lifted up
his heel against me. --- ``Who, then, is the servant who will now rise up against his
lord?'' --- ``Where, then, is the unblessed one who is thus to fulfill the Scripture?'' ---
The Lord knows it, Judas knows it, the spirit knows it; but --- what a spirit! --- And the
Master spoke again to the disciples, and said: \emph{Verily, verily, I say unto you, He that
receiveth whomsoever I send receiveth me; and he that receiveth me receiveth him that sent
me.} But when Jesus had said this, he was vehemently moved in spirit, and testified, and
said: verily, verily, I say unto you, that one of you shall betray me. --- What a moment!
All the disciples are in dismay; they look at one another, doubting of whom he spoke. Yet
no, not all are in dismay. There sits One among them: the glance may well lurk sidelong,
but --- the man himself is undismayed. He is a man who has already taken his resolution,
and when a man has taken his resolution, then he must hold fast and be --- undismayed.

\emph{Now there was one among his disciples, who lay at Jesus' bosom, whom Jesus loved. To
him therefore Simon Peter beckoned, that he should ask who it might be of whom he spoke. He
then leaned up to Jesus' breast, and said unto him: Lord, who is it?} Jesus answered: it is
he to whom I shall give this sop, which I dip. \emph{And he dipped the sop, and gave it to
Judas Iscariot, the son
% --- p. 378 ---
\opage{378} of Simon.} --- When surroundings such as these, in quiet, solemn earnestness, gather mind
and thought upon a single crime; when this crime, made clear in word and symbol, towers up
before the soul's eye, and then sinks down again to weigh as with the weight of mountains on
man's conscience, then it takes strength of spirit still to will the crime, and to hold fast
to the resolution. It is at such a turning point that we learn to shudder at the superhuman
power of evil. --- \emph{And after he had received the sop, then Satan entered into him.
Then said Jesus unto him:} what thou doest, do quickly. --- ``Why does Judas stand up so
suddenly?'' --- ``Where is he going with the purse?'' --- ``A purchase for the feast?'' ---
``To give something to the poor?'' --- ``Why does he go away from the supper?'' --- The Lord
knows it; Judas knows it \dots\ {\large When he then had received the sop, he went
immediately out: and it was night.}

--- If rational psychology, however, is of the opinion that the conception of Judas's
character held by the faith of the Gospel is supposed without further ado to prove the
existence of a personal devil, it may advance whenever it likes and fight the air: it is
then contending here against an imagined opponent. Faith knows just as well as unbelief that
a proof of this kind is wholly insufficient.

What, on the other hand, can be established for psychological consideration is that a
demonic hardening such as we have found in the traitor's character can, in the name of the
spirit, transform a human being into an inhuman being. But an inhuman being is not master
of itself. In demonic hardening the inhuman evil thus takes on the character of a
superhuman spirit. That this spirit should moreover have a personal existence, independent
of human individuals, unbelief must deny just as surely as faith must affirm it. The
Believer, who makes the temptation story
% --- p. 379 ---
\opage{379} the basis of the conception of all temptations, dare not exchange the evangelical
expressions (Joh. 13, 2. 27) for the \textit{rational} phrases of the present; he dare not
engage at all in any ``re-explanation,'' still less in any ingenious ``explaining away.''
Just as impermissible is it to misuse such statements as objections to the
self-determination of the human will. The evangelical history is in the deepest sense the
history of freedom. The conflict of the spirits is here in every respect subject to the law
of freedom. The tempting spirit that gains power over Judas also lies in wait for Peter; the
attempt has a different outcome, but why? Because the personality of the one tempted is
here a different one.

What an opposition there is, after all, between two such as Peter and Judas, and yet what a
likeness! Both move from the same height toward the same abyss; the one stumbles, the other
plunges: how far is it, then, from the denier to the betrayer? --- Few thinking people can
hear the Master's words to the first among his disciples (Matth. 16, 23): \emph{Get thee
behind me, Satan!} without involuntarily having to ask themselves: ``What dreadful crime,
then, has the unhappy Peter committed? Just now (V. 17---20) he has been exalted to heaven,
now he is on the point of being thrust down to hell. That is, after all, a terrible
transition. True enough! the all too self-wise disciple wanted to correct his Teacher; the
presumptuous one forgot the reverence he owed his Lord; but he did it, after all, with good
intentions, he did it, after all, because he could not bear that the divine Lord and Master
should give himself up to the chief priests and scribes. Did such rashness, then, deserve
such chastisement? From the purely humane standpoint it seems, after all, as if the gentle
Teacher ought rather to have helped his disciple out of his error by a calm admonition and
% --- p. 380 ---
\opage{380} an instruction resting on good grounds.'' --- It cannot be denied that the rebuking words
make an impression as though we perceived an involuntary outburst of momentary anxiety and
vehement anger. But what can now be the reason for such severity? We surely dare not assume
that Christ suddenly became afraid that Peter would talk him into yielding somewhat, sparing
himself, and thereby avoiding the suffering; nor can it be mere indignation that the bold
disciple had dared to offend his dignity that called forth this utterance of wrath. On the
contrary, Christ's words explain themselves as of their own accord when we rightly bethink
ourselves and see what spirit it is that here approaches and lies in wait for Peter. The
disciple, intoxicated by the eulogy (V. 17. 18) and the promise (V. 19), is now on the
point of taking the whole kingdom of Christ in vain. It is to be a glory without
self-sacrifice: a disguised worldly ideal is thus in the act of unfolding its glitter
between Peter and Christ. Had this delusion time to steal into the disciple, would it not
then wind itself about his heart, and twine itself through all his thoughts and feelings;
and would not all future instruction then be wasted on one who had once lost the true
impression of the Savior's person? See, Christ punishes as though he would annihilate
Peter; the word flashes as though it would shatter the rock on which the church is yet to
be built; but it is necessary! The punishment is here the immediate expression of love. The
false impression must be driven out by a word of power. To get the dear disciple back
again, the Lord must thrust him away --- in wrath: \emph{Get thee behind me, Satan! thou art
an offense unto me; for thou perceivest not the things that be of God, but those that be of
men.} --- And the rebuke worked, not as with Judas; for Peter turned, and saw his Lord
again --- on the Mount of Transfiguration (Matth. 17, 1). --- But there is, after all, a
remarkable
% --- p. 381 ---
\opage{381} likeness between Peter and Judas! If Judas can revolt us by his betrayal, Peter can
certainly terrify us by his denial.

``Yes,'' sighs the irresolute interpreter, ``if only we knew how things really stand with
this denial.'' --- The modern consciousness, which takes foreknowledge away from Christ and
cultivates an almost divine omnipresence, has then also quite rightly been in play on the
occasion of the story of the denial; it has here again lured half-criticism into the
pettiest investigations, and got mediocrity to stare at something that does not concern the
main matter at all. According to a table of denials drawn up by the famous interpreter of
Scripture, Paulus of Heidelberg, Peter must have denied the Lord at least 8 times. ---
``But,'' adds Strauß, ``by thus holding the Evangelists' accounts apart out of respect for
their credibility, one runs the risk of attacking Jesus's still more important credibility;
for he, after all, had spoken only of a thrice-repeated denial.'' And that is quite right.
To this there is further added the circumstance that the summary table lacks the proper
exactness. For the count is governed not by the denier's answers but by the number of
questioners (the portress, Joh. 18, 17. Several who stood by the fire, Joh. 18, 25. The maid
by the fire, Luk. 22, 56. Mark. 14, 66. Matth. 26, 69. One who is not more closely
designated, Luk. 22, 58. The maid at the entrance, Matth. 26, 71. Malchus's kinsman, Joh.
18, 26. One who would know Peter by his Galilean dialect, Luk. 22, 59. Several others who
join in, Matth. 26, 73. Mark. 14, 70). If now, in order to enter with complete exactness
into the particulars, one stares fixedly at ``the several who stood by the fire,''
\textit{item} ``the several others who joined in,'' then one gets not merely 8 but 8 times
8, that is, in a round number, 64 denials. Quite exact
% --- p. 382 ---
\opage{382} %
the reckoning can probably never become; but this proves precisely how infinitely exalted is
the goal toward which science strives. Were the exegesis of the story of the denial to be
quite exact, one would have to be able to calculate approximately how large the staff of
servants was in whose company Peter found himself. But to determine this, one would first
have to make an estimate of how many people in all could have found room in the high
priest's courtyard; such an estimate in turn presupposes an exact knowledge not only of the
high priest's position at the time, but especially of ancient, and in particular Jewish,
architecture in the time of Christ. This knowledge falls under archaeology. One has here a
really fine example of how one discipline lends the other a hand.

We can gladly wish both exegetics and archaeology every possible success and progress; we
can also concede that science on the whole is called upon to strive after the greatest
possible exactness, without our therefore needing to make the believing conception dependent
on specialized archaeological investigations. For the faith of the Gospel it is of little
interest which and how many the questioners were; the Believer directs his whole attention
to the denier, and can therefore very well understand how it came about that the various
Evangelists have designated the questioners in a somewhat different way, while they yet
agree most closely in this, that Peter denied his Lord three times. --- But is it not,
after all, a remarkable transformation that takes place in this disciple from the moment he
steps into the high priest's courtyard? Could it be the new and strange surroundings that
act on him so strongly? The changed surroundings must certainly have made a strong
impression on the Galilean's temper. As Peter sat among the servants at the coal fire and
warmed himself, he became, after all, a witness to how the
% --- p. 383 ---
\opage{383} chief priests had his Lord and Master interrogated and mistreated. Why, then, does he not
now think of Christ's many predictions? Alas, no, now he has forgotten them all; for ---
``he is in danger.'' ``But Peter sat without in the courtyard, and a damsel came unto him,
and said: thou also wast with Jesus, the Galilean.'' {\large But he denied before them all,
and said: I know not what thou sayest.} Why, then, does he deny it? Because he --- is in
danger. But it cannot possibly stop at the first denial: these unbelieving, intrusive
people will not, after all, leave Peter in peace; see, they crowd together around him,
challenge him and question him, one after the other; so he must set about denying again
and, if it might help, --- \emph{confirm the denial} with an oath: {\large I know not the
man.} If Peter has become a perjurer, then surely the spirit of darkness will stand by him;
yes, now he is in danger. ``But a little after came they that stood by, and said to Peter:
`Surely thou also art one of them; for thy speech betrayeth thee.''' {\large Then began he
to curse and to swear: I know not the man; I know not what thou sayest.} --- The lie itself
is now on the point of becoming truth. Peter really no longer knows his Christ; alas, it has
become quite dark around him. What if the Master now did not know his disciple either!
\emph{But immediately, while he yet spake, the cock crew the second time, and the Lord
turned, and looked upon Peter,} and saw that he was --- in danger\footnote{%
And how did Christ look upon Peter? Was this glance repelling, was it like a glance of
farewell? O no, it was as when the mother sees the child in danger through its own
carelessness, and now, since she cannot get to it to seize hold of the child, overtakes it
with her reproachful, to be sure, but also saving glance. S.\ Kierkegaard. Works of Love.
First Series. Page 187.}. {\large And Peter remembered the word of the Lord,
% --- p. 384 ---
\opage{384} how he had said unto him: Before the cock crow twice, thou shalt deny me thrice.}

There is yet some likeness between a Judas and a Peter. The natural testimony that the evil
which is accomplished by man's will is nevertheless, in its innermost movement, something
inhuman, and to that extent superhuman, we have, after all, in remorse and regret. From this
side, too, the appearance of the two disciples in the Passion story is most remarkable. The
story is not, properly speaking, interrupted by this, but, so to speak, episodically broken
through. In the high priest's palace the scene is doubled. The observer sees, as through the
veil of night, two groups at once: to the one side ``the chief priests and the elders and
all the council, who sought false witness against Jesus''; to the other side the servants
around the coal fire, where the first among the Lord's disciples stands and curses and
swears: I know not the man. All at once Peter finds himself otherwise engaged, and ---
vanishes. --- \emph{And the men that held Jesus mocked him and smote him. . . .} The
disciple's sudden vanishing and the immediate continuation of the Passion story make a
strangely deep impression. While contemplation is held captive by the scenes that follow one
upon another, thought is at the same time occupied with the disciple left to himself, bowed
down in remorse and contrition. What became of the denier now? {\large He went out, and
wept bitterly.} With this the narrative has briefly indicated that the repentant disciple is
torn out of the power of evil, that his remorse will lead him to rising again. Thus the
night passes away. ---
\emph{But when the morning was come (Matth. 27), the chief priests
% --- p. 385 ---
\opage{385} and the elders of the people took counsel against Jesus to put him to death. And they bound
him, and led him away, and delivered him to Pontius Pilate the governor:} --- then the
other, the unhappy one, comes into view. He too has spent a terrible night. \emph{When Judas,
who betrayed him, saw that he was condemned, it repented him, and he brought the thirty
pieces of silver again to the chief priests and elders, and said:} {\large I have sinned in
that I have betrayed innocent blood.} \emph{But they said: What is that to us? see thou to
that. And he cast the pieces of silver in the temple, departed, and went and} --- ended his
life in despair. Can this despairing soul now be torn out of the power of darkness? Has He,
with whom all things are possible, then mercy also for this ---
\emph{son of perdition} (Joh. 17, 12)? We do not know; nor are we given any rest to brood
over it, for the Passion story strides restlessly on; behold, Jesus stands before the
governor (Matth. 27, 11). \emph{And the governor asked him, saying: Art thou the King of the
Jews? . . . .}

\emph{For this purpose,} says the Apostle (1 Johann. 3, 8), \emph{the Son of God was
manifested, that he might destroy the works of the devil.} This does not come about by
freedom being violated; but the works of darkness are destroyed when the darkness vanishes,
when day dawns in the world of spirits. From the standpoint of naturalness every one of the
appearances of evil must necessarily be disguised, since it always is and remains relative.
It is this relativity that weighs upon the human will as a stupefying ambiguity. In one form
of expression evil is ``unavoidable'', in another ``excusable'', in the third ``questionable,
but
% --- p. 386 ---
\opage{386} yet, when it really must be, surmountable'': in this way the will is gradually
spun into so toughly yielding a web of twine that it cannot possibly break through with
original freedom. Entangled in this web of ambiguity, human nature now conducts itself quite
strangely. At one moment the will takes the enormous weight of moral guilt and lays it on
man's own head, in the proud consciousness that he, even with this burden, will be able to
climb the heights of the ideal. But when the burden then presses, and the great postulate
dissolves into a disappointment: what does the man spun in do then? Then he drowns the
freedom of the individual will in the dark necessity of the life of the race, laughs at
responsibility, and declares moral guilt to be a mere nothing. In this way the human will
certainly gets the upper hand, not only over evil, but also over good: the difference
between good and evil is transformed into a difference of degree. With this ambiguous
difference of degree the Believer, for his part, simply cannot and will not be content, and
therefore takes refuge in the Gospel of enlightenment. Here, then, one begins not with an
unclear commingling but with a separation in principle. On the level of revelation evil
stands against good as spirit against spirit, as will against will. The evil will is in its
essence superhuman, and yet human: its dominion is grounded in ambiguity itself. Only by
lifting the veil and bringing to clarity the difference in principle between That which is
not of flesh and blood and That which is of man's own nature can the power of evil be broken
down: ``for this purpose the Son of God was manifested, that he might destroy the works of
the devil.''

If we consider the places in the gospel history where Jesus points to the spirit of darkness
(Luk. 10, 18. 22, 31---32. Joh. 8, 44. 14, 30. 16, 11), we should truly have to set a high
value on the \emph{explainings-away} of the science of the School if we
% --- p. 387 ---
\opage{387} were to be convinced that these utterances are only vague accommodations. The
objection that not all the expressions can be understood literally says nothing, since here,
after all, the talk is not of a sensuous individual but of a restless, hovering spirit. That
one must always speak spiritually of the spiritual is surely self-evident. He who has not
dulled his sense for the immediate effect of the word by an all too habitual intercourse with
ambiguous turns of phrase will surely hear something more than an unclear adaptation in an
address such as the one we find in Luke (22, 31---32): {\large And the Lord said: Simon,
Simon! behold, Satan hath desired you, that he may sift you as wheat. But I have prayed for
thee, that thy faith fail not; and when thou art once converted, strengthen thy brethren.}
--- Christ does not forget the natural conditions of freedom; the Master looks sharply at
each of the personalities that surround him; he knows, after all, human relations and
conditions; but he also sees the superhuman power that knows, in all its cunning, how to get
its promptings introduced into every crisis: he sees the danger, and at the decisive moment
gives warning of the enemy lurking in the darkness. By thus leading relative evil back to its
first origin, namely the demonic will, revelation calls the fundamental human will up out of
its natural slumber, sets before it the choice between the kingdom of darkness and the
kingdom of light, and drives the problem of freedom to its utmost point.

When the sun of the world goes down, and twilight spreads over the realm of freedom; when all
the exemplars in the sanctuary of the good are shrouded as by a veil; when men begin:
``Peace, peace!'' but there is no peace; and: ``Freedom, freedom!'' but there is no freedom;
when only a fool can believe in an eternal truth, and only a madman strive for an unvarnished
justice: behold, then the Revealer comes with the light, and sets it in the
% --- p. 388 ---
\opage{388} sanctuary, and speaks a word: ``beware, soul, of the web of unrighteousness;
behold, the spinner is in the midst of the web!''

``But,'' continues the modern consciousness, ``let the old belief in the `spirit of darkness'
pass for what it is; that might do, if one could then get off with one Satan. But it is well
known to every man that no one can have dealings with the biblical devil without at once
getting a whole swarm of unclean spirits on his neck. A whole legion, i.e.\ 6666 demons in
one possessed man: pooh! that is a frightful story. What sorrows, troubles and worries have
those possessed men in Galilee not caused the mystical theologians of our day, e.g.\ an
Olshausen! --- Concerning the nature and origin of the evil spirits we find, as Strauß also
says, nothing expressly noted in the Gospels, except that they belonged to `Satan's
household' (Matth. 12, 26); which is also why what one of them does is simply ascribed to
Satan (Luk. 13, 16). But even if one stops at this boundary, it still goes equally awry,
whether the theologians now try to modernize the New Testament notions of the demonic or set
about judaizing our present concepts. We shall not now detain the educated reader with
scientific questions such as these: how does it come about that such diseases as melancholy,
madness, epilepsy and the like, which nowadays are treated according to the laws of nature,
could in those days have had a supernatural cause? How are personal demons conceivable? Can
anyone imagine the relation between self-consciousness and the bodily organs to be so loose
that a foreign self-consciousness could insert itself here and take possession of a man's
body? If one would yield a little to the philosophy of nature, and reduce the demons to
impersonal forces that act only according to the laws of the kingdom of darkness: how then
does one explain that the demons so often pass by the
% --- p. 389 ---
\opage{389} wicked and hunt down the less wicked, just as Asmodeus picks out Sarah and her
husband to torment and destroy, not because these persons were wicked, but because Sarah's
beauty enticed him (Tob. 6, 12. 15)? --- We shall not, as said, waste the reader's precious
time with general considerations; only, in order after all to enjoy a little of the naivety
of the evangelists and the insipidity of the exegetes, we must dwell for a moment on the
anecdote of the possessed at Gadara (Matth. 8, 28. Mark. 5, 1. Luk. 8, 26). After a stormy
crossing over the Galilean sea to the eastern shore, Jesus meets a possessed man who dwelt
among the tombs in that region and was wont to rage terribly against himself and others.
(According to Matthew, to be sure, there were two possessed men; but by now we are so used to
such discrepancies.) Let us now consider how unreasonable it is that the spirits in the
possessed man or men should at once have recognized and addressed Jesus, even from far off:
in Matthew the address sounds like a fearful attempt to hold off Jesus, who comes so
inconveniently; in Luke it is a humble approach to one already present; in Mark even a hasty
seeking out of the Messiah, who was still far off. What a discrepancy! Let us further
consider how Messrs.\ Demons answer Jesus when he asks the possessed man: what is thy name?
`Legion!' By all the spirits of the air, one must admire the composure and dignity with
which a Strauß has known how to enter into this gibberish. `For,' it says, `even if one can
imagine how one demon could take possession of a human organism while suppressing the human
consciousness, one cannot at all imagine that many demonic personalities could possess one
human being. For since this possession is nothing other than making oneself the subject of
the consciousness in an individual, but consciousness in reality can only have
% --- p. 390 ---
\opage{390} one apex, one center: it is in any case not conceivable that several demons
should be able at one and the same time to take possession of a whole human being, and the
multiple possession could only ever be present as a successive alternation of possessions by
different demons, and a whole host of them could not, as here, dwell at the same time in one
human being and leave it together.' Let us furthermore consider how improbable it is that
the demons (assuming that they were real demons) could have entertained and obtained the wish
to go into animal creatures. Strauß thinks that they must in any case have been rational
spirits; but when we see the effect they produce on the swine, we might, with all the respect
we otherwise have for the far-famed critic, almost be tempted to suppose that they must have
been raving spirits. For scarcely have they entered into the swine before they at once drive
the whole herd to hurl itself into the lake; whereupon Strauß himself rightly asks `what the
demons had then gained by entering into the animals, when they at once destroyed them, and
so deprived themselves again of the tolerable interim lodging they had so urgently begged
for.' Let us then, in order finally to be done with this caricature of Jewish superstition,
consider in conclusion how utterly bare of any moral significance this miracle stands here,
unless it should be this, that Jesus made an open encroachment on the lawful property of the
Gadarenes. When even an Olshausen places the Gadarene swine of the N. T. beside Balaam's ass
in the Old, it is surely evident enough that there is no rescue for this miracle story.
--- But once we have sawn this one out, we shall doubtless get more.''

The Gospel book and the Gospel are not the same. Should it really prove possible to show that
a fragment had been inserted in the gospel history that did not bear the stamp of the
Gospel, then
% --- p. 391 ---
\opage{391} the Believer would pass it over in all silence, yet not hastening with any
explaining-away, but dwelling on the Savior's word: \emph{blessed is he, whosoever shall not
be offended in me.} That the faith of the Gospel, however, will also know how to appropriate
this narrative, we can perhaps best establish by indicating the following conception:

Christ, having crossed the Sea of Gennesaret with his disciples, has here come to the outer
border of Galilee of the Gentiles (Peræa). To these border-dwellers too he will reveal himself
as the Redeemer. This cannot now come about by his beginning to speak the word and expound
the Scriptures; for these Gadarenes are sunk in coarseness and live in a half heathenism,
which shows itself among other things in their keeping, contrary to the prohibition of the
Mosaic law, great herds of unclean animals. The most powerful means of awakening the
instinct of faith in these coarse minds was thus to begin with a saving act. To such an act
Jesus is occasioned when the possessed man came out from the tombs and cried after him. Once
the supernatural deed had been accomplished and made its impression, the clarifying word
could follow. Christ's act is twofold: \textit{a)} he heals an unhappy human being (perhaps
two); \textit{b)} he lets the unclean spirits go into the unclean animals. By this the act is
at bottom transformed into a speech, and this speech contains a question of conscience, which
we can now expound more closely for the dismayed Gadarenes: ``What do you judge of this
stranger who has so unexpectedly come to visit your region? See, on the one hand he has
healed a man who was unhappy, yes, sorely unhappy: now this man sits at his feet, clothed and
in his right mind. Such things are accomplished only by supernatural powers. But on the other
hand, on the occasion of this healing, he has also caused you great damage: \emph{the whole
herd of swine}
% --- p. 392 ---
\opage{392} %
\emph{has run violently down the bank into the lake and perished in the water.} What do you
think now of this wondrous man? Should he not seem to you to bear some likeness to one of the
old prophets (you have surely heard tell of the prophets in Israel)? Might he perhaps, by
letting the unclean spirits destroy your unclean animals, wish to remind you that you too are
Israelites, not heathens (for this much you have surely heard, that to the Israelite the swine
is an unclean animal)? Should your conscience remind you; should you be able to perceive in
your hearts that only a messenger of God can do such works: then the man stands here before
your eyes, as if he were waiting for you to come nearer and ask him something.'' --- Thus a
choice was set before these Gadarenes; it was therefore their own fault when in their terror
they valued the loss of an unlawful property more highly than the salvation of a human being.
\emph{But the man out of whom the devils were departed besought him that he might be with him;
but Jesus sent him away, and said: Return to thine own house, and tell how great things God
hath done unto thee. And he went his way, and published throughout the whole city how great
things Jesus had done unto him} (Luk. 8, 38---39).

When someone otherwise wants to talk about what he does not understand at all, when, e.g.,
someone not skilled in medicine wants to hold forth on how physicians really ought to treat
the sick, one pays no heed to his talk; and if he nevertheless persists in it in a
presumptuous tone, one is perhaps even close to despising his talk. Must it not then seem
strange to us that science will allow its critics to reason so far and wide about things
that lie entirely outside this human knowledge! A mathematician is to blame if he demands of
his pupils
% --- p. 393 ---
\opage{393} the solution of a problem that is in itself insoluble; ought not the critic then
also be blamed for confusing the understanding of the sacred history with utopian questions?
--- ``How many demons can be accommodated in one possessed man?'' ``What entitles Christ to
let the unclean spirits go into the Gadarenes' herd of swine?'' With such questions the faith
of the Gospel makes short work, and turns away with the same aversion from the defending as
from the attacking interpreters. If, on the other hand, anyone will put to himself the
religious question, what it is, then, that the Redeemer reveals to us in healing the
demoniacs too: then the spirit of the Gospel will surely also help him to a fruitful
understanding. It is not contributions to the further progress of physics or medical science
that we seek in Jesus' wonder-works; it is the comforting assurance that there is no bodily
evil or spiritual misery so great that the Redeemer cannot and will not heal it. The only
evil, then, that remains is man's own moral guilt. To be saved from this evil, the will must
will its own salvation: there are demons that are cast out by the immediate authority of the
apostolic word; there are demons that are cast out by prayer and fasting; there are demons
--- that are not cast out.

\subsection*{\textit{VI.}\\[4pt] Christ is mighty in word and deed.}
\addcontentsline{toc}{subsection}{VI. Christ is mighty in word and deed}

When one has a view one wants to carry through, it is always a great advantage to be able to
appeal to a recognized authority, that is, if the authority can in any way be won over to the
view. Why is it that theologians and anti-theologians of every color still set so high a
value on biblical expertise
% --- p. 394 ---
\opage{394} and learned exegesis of Scripture? It is, of course, because the Bible is an
authority. The Free Thinker is just as busy studying his Hebrew and Greek Testament as the
orthodox man. ``Authority against authority: the authority of the spirit of the age or of the
Bible?'' The combatants know very well that men, with their independent thinking and thinking
independence, will in the end settle down after all with one reliable authority or another.
But to win the authority for oneself requires dexterity and shrewdness. The positive
theologians of the present day, who think that they will surely uphold the authority of
Scripture if only they hold fast to the scriptural principle with both hands, are, so it
seems, a little clumsy. Not to mention that the biblical world-view finds itself in so hard a
conflict with enlightenment, the roving criticism of Introduction keeps on teasing these Bible
theologians by snatching away from them one piece of their canon after another. ``Now,'' it
says, ``one must utterly reject Matthew as an eyewitness''; then there is a great deal to
correct in Mark; then something is wrong with Luke; then comes --- it is almost not to be
borne --- at last the sorrowful tidings ``that the Gospel of John can hold out no longer''.
How the otherwise so self-sufficient scriptural principle will end up under these
circumstances, the Lord knows. But the negative critics, too, act unwisely more often than
not; for partly their attacks on the authenticity and credibility of the sacred books are,
after all, all too rash in many essential points; partly the philosophy of life that modern
knowledge would put in place of biblical doctrine still sounds a little too high-flown for
plain folk. ``No God? no providence? no life after this? why, it is dreadful to think of''.

No, the true word-heroes, the genuine rationalists, have understood quite otherwise how to
make the authority of the Bible fruitful both for themselves and for their fellow men. A very
% --- p. 395 ---
\opage{395} %
welcome helping hand the theologians of reason received already in the last century from the
philosopher Kant. According to Kant, namely, the Christian sources of religion in the O.\ and
N.\ Test.\ must always be referred to a meaning that agrees with the practical laws of the
pure religion of reason. This interpretation, even if it does violence to the text, apparently
or really, should be preferred to the literal one when the latter, as is the case with many
biblical stories, either contains nothing at all for morality or even works against the
moral motives. By borrowing the authority from Scripture but the truth from reason, the
enlightened theologians could now, as regards the gospel history, fairly easily come to terms
with Jesus' words; but to get on with Jesus' deeds was, on the other hand, not so easy a
matter. On the one hand it was important not to impugn the credibility of the evangelists too
far; on the other hand the great multitude of meaningless and offensive miracle narratives
had nevertheless to be disposed of \textit{à tout prix}. On closer examination it had to turn
out that everything had happened naturally. We set down here some samples of the so-called
natural explanation. --- \emph{Jesus turns water into wine.} This miracle is --- no miracle.
``How is it conceivable that Jesus should have used so extraordinary a means for so slight a
purpose? He who in the wilderness would not still his own hunger by a miracle cannot possibly
have performed a miracle to procure wine for a party that had surely already had enough.
Add to this that the quantity of wine Jesus procures for the guests must seem, according to
the text's statement, far too unreasonably excessive. Six jars, each of which held 2 to 3
measures, metretes; this gives (when each metretes is reckoned at $1\frac{1}{2}$ Roman
\textit{amphoræ}) in all 255---378 Württemberg cans.'' ``What a quantity for a party that
% --- p. 396 ---
\opage{396} was already fairly drunk!'' exclaims Strauß. ``What enormous jars!'' cries
\textit{Dr.} Paulus, and then hastens to reduce the measure and explain the event. Instead of
letting Christ turn water into wine, Paulus now undertakes to turn the miracle into a quite
natural occurrence. ``Jesus had secretly seen to having a portion of wine in readiness. In
order to surprise with his gift he is reluctant to let anything be noticed before the time,
and therefore rebuffs his mother, who is close to betraying the secret. Although the wine was
presumably drunk on the spot, one can nevertheless, on \textit{Dr.} P.'s assurance, be certain
that Jesus knew how to maintain decorum just as fully as the Roman Marc.\ Porc.\ Cato when the
latter held small drinking parties at his country house.'' The moral point of the narrative
is that the disciples see Jesus' glory, i.e.\ his good humor: ``Ein so freundlich
wohlthätiger, ein den Frohsinn der Gesellschaft auf eine scherzhaft geheimnißreiche Art
befördernder Mann, nicht ein Freudenstörer, nicht auf menschliche Wünsche hoch herabsehend ist
unser Messias!'' (H.\ E.\ G.\ Paulus. Commentar über das neue Testament). Meanwhile, however,
one obstacle after another is met with here, and the road through the holy land of miracles
would certainly be quite impassable if the natural art of interpretation did not simply take
it upon itself to carry out what in its opinion was not possible for Christ, let alone for
his disciples, namely to move mountains. Let us go on, then, and see how the rationalists move
their mountains about.

\emph{Jesus walks on the sea.} To walk on the water, and at night into the bargain, and in a
strong storm, is flatly against the physical laws of gravity. So that this narrative shall
not bring the Bible into discredit and weaken its authority, the just-mentioned \textit{Dr.}
Paulus offers to move a mountain. This wonder he performs by a quite simple means. Into the
words of the original text, to
% --- p. 397 ---
\opage{397} walk on the sea, he reads the meaning that Jesus walked so near to the sea that he
could speak with the disciples from the shore. For since it was a stormy night, there is
every probability that the disciples sailed as near the coast as possible. ``But now the
moral?'' Well, the moral seems to be this, that the disciples kept closer to their Master than
the interpreter to his text. --- \emph{Jesus feeds 5000 men in the wilderness with 5 barley
loaves and 2 small fishes.} The believing simplicity that sees in this only a supernatural
effect of the blessing that lies in Jesus' word naturally cannot understand Joh.\ 6, 13
otherwise than that, \emph{after} the feeding, 12 baskets were filled with the fragments that
remained over from the 5 loaves. But \textit{Dr.} Paulus and practical reason, which must in
equal degree take offense at so unnatural a multiplication of the natural gifts of God, have
here too joined together to move a mountain and by this operation to drive out the miracle.
This enormous undertaking is carried out, however, not by sorcery or ecstatic incantations,
but in all prosaic sobriety, by inserting a parenthesis. By means of a single little
parenthesis, placed in the right spot (cf.\ Comment.\ üb.\ d.\ Joh.\ S.\ 327 and 337), the
whole event is turned around, and the account then comes to look as if the evangelist really
meant to tell us that the 12 baskets were filled \emph{before} the feeding, and that in the
following way: Among the 5000 men who had gone out here into the wilderness to hear Jesus,
there may with every reason be supposed to have been a number of well-to-do and provident
people who had not neglected to take provisions with them. Christ's true beneficence consists
now, properly speaking, in this, that by offering his disciples' ``5 barley loaves and 2 small
fishes'' he set the better provided a fine and good example. By all the contributions being
collected and
% --- p. 398 ---
\opage{398} %
fairly equally distributed, it came about \textit{miraculeusement} that the whole great
multitude got an ample meal. From this one learns what wonder-working blessing is laid up in
a proper natural interpretation. The moral is that the exegetical parenthesis can --- move
mountains. One more example: \emph{Jesus is transfigured on the mountain} (Matth.\ 17, 1---8).
This transfiguration stands in great need of an explanation. Here several ways out at once
present themselves. One can, e.g., suppose that the disciples at once fell asleep up there on
the mountain and dreamt the whole thing. But since it is not altogether natural that three
people should dream the same dream at once, it is surely better to let \textit{Dr.} Paulus
give his hand to \textit{Dr.} Kuinoel, so that with combined help they can move yet another
mountain for us. The real scene then consists essentially in this, that Jesus had a secret
meeting on the mountain with a couple of men, Paulus does not know which, but Kuinoel
considers it likely that they must surely have been two such as Nicodemus. Even before the
two men arrived, Jesus prayed, and the disciples, who took no part in the prayer, fell asleep.
On hearing the strange voices they wake up. Jesus, who probably stood on a higher point of the
mountain than where his disciples lay, now appeared in an unusual radiance. This radiance came
from the first rays of morning, reflected by a nearby layer of snow. The drowsy Peter, and
after him the others, hit, one knows not for what reason, on the idea that the two men are
--- Moses and Elias. The disciples' consternation at last rises to its height when they see
the two strangers vanish in a bright morning mist that came down just as they were about to
go away, --- and hear one of the ``Nicodemuses'' call out from the cloud of mist the words:
\emph{this is my beloved Son, in whom I am well pleased; hear ye him.} ---
% --- p. 399 ---
\opage{399} The moral is that simplicity smiles, but science becomes --- serious.

The serious interpreter, who will not tolerate the scientific reputation of hermeneutics
being exposed to danger by such dubious samples, can then with decorum and dignity set forth
the following: ``To raise oneself above the one-sidedness of the great predecessors is not
yet sufficient proof that one has oneself penetrated more deeply into the matter. The
excesses of psychological-historical interpretation hermeneutics as a science long ago
repudiated, without therefore being able to disregard the apparatus of learning and the
acumen that those scholars set in motion in order to penetrate into the individual parts of
the text and to order the historical features in their pragmatic causal connection.
Theological interpretation knows with sure tact how to hold fast the mysterious moment of the
wonder in the gospel accounts, and therefore need not heed the untimely admonition that one
might possibly have intended for it by a bantering presentation of individual blunders.''

We bow deeply before hermeneutics as a science; nor should we have presumed to say a word
against ``theological interpretation'', if only we had been able to rely on its now also
being in every respect clear with itself concerning ``the mysterious moment of the wonder''.
But besides the fact that hermeneutics as a science has hitherto disdained to bestow on us a
reasonably thorough treatment of the religious reciprocal relation between the word and the
wonder, there is, to tell the truth, something else that makes us anxious with regard to the
very vague determination of ``the mysterious moment'', and that is an unreserved, but, as it
seems, fairly well-founded objection by Strauß. ``The most modern thing,'' it says,
``nowadays is, while giving up the strict concept of wonder as \textit{miraculum}, still to
recognize
% --- p. 400 ---
\opage{400} \textit{mirabile}, an operation of powers of a higher rank, as it is called, in
the life of Jesus and in biblical history generally; and one can then partly appropriate the
sonorous philosophical formulas, partly make use of the more recent observations in the field
of animal magnetism. But with this one has robbed the wonders --- which have kept only the
name --- of all power to prove the divinity of Christianity. For on this standpoint the
wonders are no longer actions of God himself, who intends by them to procure credence among
mankind for the wonder-worker; rather, the power to work miracles is a natural capacity,
which individual men come by in the same way as by other similar gifts, such as bodily
strength, a peculiar fineness of a particular sense, or else a talent for music, mathematics
and the like, namely through a concatenation of natural causes, which, as is well known,
distributes such natural gifts without regard to the moral character of individuals. But,
one says, the theological world-view, faith in divine providence, demands after all that God
distinguish with such high gifts only those men who deserve such a distinction. --- How so?
Has he then never distinguished a seducer with beauty and eloquence, a cruel man with
strength, a malicious man with understanding? --- Yes, indeed, with such natural gifts; but
the power to work miracles is something else. --- Then tell us plainly: do you regard the
capacity to work miracles as something supernatural? Then you still hold fast to the old
concept of miracle, against which we have said enough: or is it your earnest intent to
comprehend it as a natural gift? (the distinction between the higher and the lower nature is
meaningless), then it too, like all natural gifts, will stand in a contingent relation to
man's moral worth, and the possible possession of it can be no proof for a founder of
religion. This makes it superfluous to test more closely
% --- p. 401 ---
\opage{401} %
how far this notion can be applied to the biblical narratives of
miracles\footnote{D.\ F.\ Strauß. The Christian Doctrine of Faith, translated by H.\ Brochner.
1st Part. Page 222.}.''

That the ``theological interpretation'' of half-criticism must surely feel the lash of the
modern consciousness's whip can also be noticed in this, that with its general reservation of
``the mysterious moment of the wonder'' it always shrinks from an unconditional appropriation
of the individual, particular wonder-work. To cover this weakness, a view has gradually been
introduced according to which the wonder becomes a mere side issue, the word, by contrast, the
real main issue. As a scriptural recommendation for this separation of word and wonder one
has even directly appealed to Christ's authority (Matth.\ 12, 39. Joh.\ 4, 48 and other
places), and for further security supported this divine authority with a declaration of ``the
strong-in-faith'' Luther: ``God has,'' says Luther, ``had to give the heathen outward signs
that one could see with the eyes and grasp with the hands, has had to draw to himself by
outward signs those who before knew nothing of God, just as one must throw apples and pears
to children. But that Christ without ceasing works in Christendom with his divine, almighty
power, that is in truth to cast out the devil, drive away serpents and speak with new tongues
for believers. Those visible works are signs for the foolish, unbelieving crowd; but we who
already know such things and believe the Gospel, what need should we have of them?'' --- By
reducing the wonder to something accidental and bringing out the content of the word as the
essential and abiding, the interpreter now thinks he does so much the better, since in this
way he keeps step with culture and brings theology into a friendly understanding
% --- p. 402 ---
\opage{402} %
with philosophy. ``One is,'' it says in Hegel, ``to believe in wonders, and this is to be a
means of believing in Christ; it is quite possible that it is regarded as a means, but it is
also demanded for its own sake. This faith, which is thus demanded, is faith in a content that
is contingent, i.e.\ it is not the true faith, for the true faith has no contingent content.
This must be especially noted with regard to the so-called Enlightenment: it has got the
better of this faith, and when orthodoxy still continues to demand such a faith, it can no
longer preserve it in certain people, because it is a faith in a content that is not divine,
not the testimony by which God testifies of himself as spirit in spirit. Whether the guests
at the wedding in Cana got more or less wine is quite indifferent, and it is just as
contingent
whether that man's withered hand was healed, for millions of people go about as cripples and
with withered limbs without anyone healing them.'' --- By reducing the wonder right down below
the word, the interpreter further gains the intelligent superiority over all the evangelists
that with regard to each individual wonder-work he can at once decide whether it is
expedient or not. ``Was it necessary, was it fitting for Jesus to perform such a miracle?''
Such questions are just the thing for exercising the understanding and sharpening the
judgment. If one opens the newest commentaries, how many wonder-works do the exegetes not
still see themselves compelled to reject, because they can neither see their use nor find
them consonant with Jesus' dignity. The walking on the sea cannot be accepted; it has
something ghostly about it. The stater in the fish's mouth cannot be accepted; it is without
moral value. Of the wine in Cana and the swine from Gadara we will not speak at all. By thus
rejecting a considerable number of gospel
% --- p. 403 ---
\opage{403} wonder-works one gains the advantage that the few but choice miracles one keeps
(that is, if one keeps any) can all the better satisfy the moral as well as the cultivated
taste. If one furthermore understands, as regards the miracle narrative itself, how to keep
the individual features properly hovering, and cautiously to intimate ``the mysterious
moment'', then it surely cannot fail that the wondrous event will indeed work with edifying,
uplifting power on the religiously cultivated mind. --- Once the theological interpreter has
with due decorum pushed the wonder into the background, why, then it is naturally an easy
matter to win even the very freest disposal over the word. What should stand in the way here?
The content of the word is, after all, spirit, and pure spirit is in truth both pliant and
compliant toward anyone who knows how to rise a little above the letter. On this height the
``theological interpretation'' of hermeneutics can then truly enjoy the pure joy of
developing and presenting Jesus' pure doctrine, of contemplating and admiring Jesus' pure
personality, his awe-inspiring loftiness, his gentle condescension, his earnest morality, his
untiring activity, etc. All the eternal, instructive truths that are otherwise found so
widely scattered, in a Confucius and a Lycurgus, a Pythagoras and a Socrates, a Plato and an
Aristotle, a Cato and a Cicero, a Horace and a Seneca, behold, all these eternal, instructive
truths are now found here compressed and uttered and enjoined by the teacher who taught as he
lived, and lived as he taught, until he went away and left us an example.

--- But what use is it, after all, for the theologian to torment himself with so much
interpretive diplomacy, when in the end it nevertheless falls upon himself? He turns to
Christ in order to obtain from this highest authority a reduction in the value of miracles.
But the wonder-worker cannot possibly slight his own works. Christ speaks against the crowd's
sensuous misuse of
% --- p. 404 ---
\opage{404} the outward signs, but not against the essential validity of the miracle itself.
\emph{If I do not the works of my Father} (Joh.\ 10, 37---38)\emph{, believe me not. But if I
do, though ye believe not me, believe} the works. The interpreter next turns to Luther, and
now really seems to find some support in this very considerable authority; but let it be well
noted that the ``apples and pears'' the open-handed Reformer bestows on the supplicant are by
no means to the taste of the present day. Luther, who inveighs against the Catholic
miracle-mongering, against the monks' holy tales and childish legends, certainly has no
inkling that his declaration should ever be used against the gospel wonder-works themselves.
If he now got to see how his humorous remark is used, there is no doubt that the man, as
strong in faith as he was zealous, would instantly wish his ``apples and pears'' turned into
--- red-hot stones, if that could help ``the theological interpreter'' --- out of theology.
If the interpreter then finally knocks at the philosophers' door, intending by their
reputation and assistance to procure himself a usable definition of the concept of wonder,
one in keeping with the demand of the times; well, then he will certainly soon notice that he
has come into cultivated company. But the awkward thing is that he who has once let himself
slide down the steep slope of consequences cannot all at once stop halfway by merely saying:
stop! since the speed, as is well known, grows the stronger the longer it goes. If one would
begin with Schelling and Hegel, one must put up with ending with Strauß and
Feuerbach\footnote{Wer einmal ein Wunder glaubt, dem ist überhaupt nicht die geistige Macht,
noch die Macht der Natur, sondern allein die Wunderkraft die höchste, wahre, die Welt
bestim\-%
}. ``But
% --- p. 405 ---
\opage{405} %
the theologian ought, after all, also to use his own sound reason!'' --- ``Yes, he ought, and
he does, too, as best he can; but, alas, it is not yet enough.'' --- ``No, in these
enlightened times of unbelief there is, to put it plainly, as good as no theological
enlightenment that will properly carry through. But is that the fault of theology? Take, e.g.,
the expediency of miracles: how often has apologetics not worked over this important point,
and that in the sweat of its brow! but --- it does no good all the same. How many individual
miracles has theological interpretation not already striven to get out of the way, and that,
by God, out of sheer compliance toward the enlightened; but --- the world is ungrateful! ---
it does no good all the same.''\footnote{mende und regierende Macht, dem erscheinen alle
Gränze, alle Unterschiede, alle Gesetze als \emph{rein willkührlich}\dots\ Heute macht das
Wasser naß, morgen trocknet es vielleicht; heute bewegt sich die Sonne, morgen steht sie
vielleicht stille \dots\ die Wunderkraft kan mir absichtlich ein \textit{Quid pro quo}
vormachen, damit ich mich nicht auf mich verlasse, sondern stets in dem Abhangigkeitsgefuhl
von der Almacht des wunderthatigen Willens erhalte. Wenn ich, wie Bileam, einen Esel reden
hore, so weiß ich nicht mehr, ob ich ein Esel bin, oder ob der Esel ein Mensch ist. L.
Feuerbach. U. d. W. S. 6.}

It is certainly an extremely difficult task to have to judge a miracle's necessity,
expediency and moral validity by rules of common reason. What appeals to one repels another;
indeed, the changing of moods goes so far that what at one time seems edifying to a person
may at another time seem offensive to the same person. Seen from one point of view, the
gospel wonder-works are all together highly edifying; but from another point of view
% --- p. 406 ---
\opage{406} %
they are all without exception equally inexpedient, equally superfluous, equally offensive.
God knows we are not exactly much taken with the modern consciousness, but ---
\textit{suum cuique} --- we nevertheless dare not deny that there is something in the remark
with which Strauß introduces his dissolution of the concept of miracle (Doctr.\ of Faith 1st
Part p.\ 199): ``That the doctrine is divine, the wonders are to prove to me; but that these
themselves are divine, and not rather devilish, I am to see from the doctrine. One who stood
in league with demonic powers would not deliver a doctrine so excellent in moral respect, nor
lead such a life; but whence do I know, on the ecclesiastical standpoint, what is moral,
except from this doctrine? The doctrine is as it is, and apart from it I have no standard by
which to measure its excellence; or, if I have such a one, then I presuppose something in me
that stands above revelation.'' --- A step further. Once one has begun to raise doubts about
the authenticity of the wonder-works, how will one then be able to avoid raising strong
doubts about Jesus' own sinlessness and freedom from deception? If the evangelists have
mistaken their apprehension of the outward signs, what certainty is there that they have not
also mistaken their apprehension of Jesus' own inner life? It is surely self-evident that
there can be still less any irrefutable historical proof of a man's perfect sinlessness than
of the authenticity of miraculous facts. Or, if this is not evident enough by itself, we can
of course appeal to an authority. ``The hatred of Jesus' enemies,'' says the theologian Hase
(Life of Jesus. p.\ 39), ``which did not dare to assail the purity of his conduct, Pilate's
acknowledgment, the noble Roman lady's sympathy, the robber chief's confession on the cross,
Judas's despair: all this testifies only to an upright man whose blood was shed guiltless.
That the Gospels relate nothing ignoble about Jesus could, without
% --- p. 407 ---
\opage{407} %
violation of truth, also be the case in the biography of a lesser man, written down by
friends to honor his memory and for the edification of his admirers. The stern Baptist's
homage, the apostles' unconditional reverence and their definite declaration that Jesus was
righteous, holy and without sin, is an irrefragable testimony to his moral loftiness; but they
did not know his earlier life and the secrets of his heart well enough to be able to testify
that sinful desire had never found room in it; it is also quite improbable that this strict
concept of sinlessness, as more recent times have conceived it, was known to the apostles.
Therefore their testimony has not much more significance than when Xenophon (\textit{Memorr.
1, 1. \S\ 11}) praises his teacher for this, that no one ever saw him do anything impious or
heard him say anything unholy.'' --- If one now has no sufficient guarantee of Jesus'
sinlessness, then the theological praise of the infallibility of his doctrine and the divine
authority of his word is also empty talk. Christ must therefore put up with standing in line
with the other distinguished teachers of mankind; the content of his doctrine is to be judged
according to the same laws as any other human doctrine. Even if one would, as Hase proposes,
ascribe to him ``perfect holiness in life and perfect clarity in God-consciousness, so far as
both are granted to human nature''; this last clause could nevertheless permit no suspension
of the rule stated here. But by applying the rule one abolishes Christianity. Let us consider
this matter with all the attention it deserves.

As the truth is, so must its communication be. Without stirring up discord in the kingdom of
truth we can with good reason distinguish several kinds of truths. There are truths of reason
that are wholly impersonal; there are truths of reason that are partly
% --- p. 408 ---
\opage{408} %
personal; and there are revealed truths. Each of these main kinds requires its own particular
communication. As an example of the first class we can cite the mathematical truths; since
these are impersonal, the communication is impersonal too, that is to say: here no closer
question is asked about the personality of the one communicating. Certainly it is required of
the teacher that he have the right insight and also the gift of being able to set forth the
doctrine clearly and intelligibly. But these qualities do not, in the proper sense, work
personally. Once the proof has become evident to the pupil, it can at bottom be indifferent
to him what the teacher is like: he is now himself the holder of the rational truth. Pure
insight here displaces authority. If one could construct a speaking machine that delivered
mathematical theorems by a good method, would not such a machine then in the main accomplish
the same as a mathematical lecturer, and consequently also have the same scientific
authority? --- It is otherwise with the personal truths of reason; they would rather be
communicated personally. If we look, e.g., at moral upbringing, how much does not depend in
this respect on the educator's character and personal relation to the one being brought up.
Nevertheless personality is still only a relative magnitude, and its significance relative.
The guardian must certainly begin by commanding, and the ward by obeying. But the command
that gets its first validity from the authority of the one commanding, and works morally
insofar as it is followed in faith and trust, is nevertheless at the same time designed to be
justified one day by rational grounds: that the minor may himself come of age, that practical
reason may take over from the guardian's authority, is surely the final aim of upbringing.
``But now the educator's own good example''? Certainly a good example is a precious thing,
rich in much blessing; but for all that one must nevertheless not forget that the personal
% --- p. 409 ---
\opage{409} example is a means that in and for itself must never become an end. The guardian who
acts morally merely to give his ward a good example is moral in roughly the same sense as the
gentlefolk who go to church merely to give the servants a good example are religious. As a
moral intermediate determination, the example will always be subject to ambiguity. Of the
bad example one says that it is contagious; this implies that the personal example works
as if by immediate contact. But precisely because the effect of the example always presupposes
the individual appearances of the personality (naturally not always its sensible
presence), it is constantly exposed to misinterpretation. The morally distinguished characters who
are praised by posterity on account of ``the good example'' have almost all been condemned by
their contemporaries on account of ``the bad example.'' Think of a Socrates, a John Huss, and
others. ``How does it come about, then, that the good example first shows itself when the person has vanished,
and thus works differently on those who come after than on those who live alongside?'' It comes from this,
that the eminent character is the bearer of the idea. As long as the person himself exists and works in
the service of the idea, he is commonly a great nuisance and much trouble to his contemporaries. Without
being able to see through his task (for as long as the idea is still grown together with its personal
organ, it leads, so to speak, a cryptic life) they receive all the more palpable an impression of
his sharp, pronounced peculiarity. To their eyes it looks for all the world as though it were
merely this single human being whose presumptuous self-will everywhere made a point of opposing
the established order, everywhere would get in the way, and cause disturbance. If the balance is to be
restored and understanding to come, a transformation must first take place. This
transformation does not come about through the people of the great world around him being converted to recognition
% --- p. 410 ---
\opage{410} of the law of the ideal. Far from it! these contemporaries have, after all, their own laws, once
agreed upon; why then should they be so foolish as to submit to a single man's
``self-willed whims''? Still less can the reconciliation be brought about by ``the
obstinate man,'' that is to say: the servant of the idea, condescending to abate something
and violate the idea merely to humor those who live alongside him. No, with the obstinate man it only
gets worse the longer it goes on. He who should have given his fellow men a fine example of
compliance and pliancy must, out of inner obedience to the command of the idea, resign himself to giving them a
hated example, an example of outrageous self-will. So there is only one power that can
restore the balance and bring about the transformation. This power is death. How the
distinguished personality comes to his end in the finite sense of the word is, with regard to the idea
itself, fairly indifferent (the man could, for that matter, just as well die in his bed); but as a
rule it ends with the catastrophe that the contemporaries put the obstinate man to death,
and thus, without themselves knowing it, pave the way for good understanding. --- In
religious-solemn usage it is of course said of every deceased person that he is one transfigured; as far as
the idea is concerned, however, it comes down to finding out what it really is that becomes
transfigured. The life-content of some people is already so comprehensible and clear, long before they
die, that death was really not needed at all to get them transfigured; others, and among these
belong the ethical characters of the first rank, cannot in their whole lifetime and with all their
exertion bring out the word that should make their true justification understood, and
must therefore wait for death one day to open their letter of authority, once it has
pressed the seal upon their lips. The transformation by which death helps such a
personality to become ---
% --- p. 411 ---
\opage{411} ``one transfigured,'' is two-sided. It ends the discord. Discord presupposes, after all, that there
is drawing in contrary directions, as the two are drawn into contention with each other. With ``the transfigured one''
there is no discord: he is no longer drawn into contention with men; for he is, after all --- drawn
away. Then hatred is quenched, then reflection begins. What the thousand eyes of the passionate age
could not see, because they constantly stared themselves blind at the obstinate human being, the
calm eye of reflection can now clearly see, namely the idea in whose service the obedient one has
been sacrificed\footnote{Without being able to cite any particular passage, the author would merely remind in
general of the many dialectical determinations belonging here which are found in
Mag.\ Kierkegaard's writings.}. Then posterity receives a great, an uplifting example.
But how? At the time when the man by his personal appearance should have given a good
example, he gave a bad example; now, on the other hand, when the bad example has become good, the
man no longer works personally. Thus: as long as the person is, the good example is not;
when the good example is, the person is not; can there be any clearer proof that
the good example must be a little more than personal? --- What is seen from this standpoint
to move through all the shining examples of the past is, without circumlocution, the
great, universal fundamental idea of morality. The ideal does not work immediately by means of the personal
example; but the personal example is raised to an ideal by means of the idea. Whether the actual
person who fell as a sacrifice to the idea really also resembled ``the
transfigured one'' in all respects, posterity cannot decide, and does not need to either: it has the ideal, and
that is enough. All theological efforts on these terms to assert Christ's personal
authority and historical exemplar by ``an inference back from the moral
effects of Christianity'' are powerless. It remains
% --- p. 412 ---
\opage{412} thus: the ideal stands above the historical exemplar, the idea of reason is the only authority.

With the revealed truths the case is quite otherwise; they have their
last and only ground in the divine authority. If a truth of revelation must first be
derived from the universal idea, or made dependent on rational grounding, it is all over
with it. The revealed law has, according to the O.\ T., its principle in Jehovah's will. ``Thou
shalt do it, for I am Jehovah.'' \textit{Sic volo, sic jubeo, stat pro ratione voluntas}:
other grounding is not given, and from the standpoint of revelation cannot be given either.
Insofar as God now uses a human being as his instrument, and chooses him to receive and proclaim
the revelation, everything will at the same time rest on the unconditional authority that the founder of the religion can
procure for himself. Not that the founder of the religion, as a finite individual or as an independent thinker and
sage, should make the least claim to recognition; but on God's behalf, as God's
messenger (thus not on the ground of the idea of reason, but by virtue of a wonder), he must demand of
men unconditional trust. --- He who will not be content with the wonder, but
wishes for a rational proof that Muhammad really is Allah's prophet, cannot be a
Muhammadan; and likewise he who does not dare simply to assume that God has revealed himself to
the patriarchs cannot be an Israelite either. Moses and the prophets have no other grounding
than this: \emph{thus speaks Jehovah, the God of Israel.} The same holds of Christ. To
justify Christ's divine authority by virtue of the idea of reason that runs through
his life and teaching is, in other words, to undermine his authority, and to set him far below
the philosophers. The philosophers at least strive, each in his own way, to
defend their personal conviction by adducing
% --- p. 413 ---
\opage{413} grounds, but Christ knows only one single form of grounding, and it is this:
\emph{verily, verily I say unto you.} --- ``But,'' exclaim the enlightened, ``when one can with
such effrontery attack sound reason in order to defend revealed religion,
one can surely also take on the slight trouble of turning light into darkness and
darkness into light, for the difference surely does not mean much. The orthodox Turk is just as
justified as the orthodox Jew and the orthodox Christian. Reason has thus lost its
right to distinguish between the false and the true religion; \textit{eh bien,} then stake and
fire are once again the only proof of a revelation's truth.'' --- Certainly it is hard
to have to take reason captive under the obedience of faith\footnote{``Ein dogmatisches Faktum ist
daher die tollste Chimare, die je in einen Kopf gekommen, toller als ein holzernes Eisen.
Das Dogma (d. h. als Lehre) hat nur innere Merkmale; es macht auf Wahrheit Anspruch, es
appellirt mit diesem Anspruch an den Gedanken, es giebt sich der prüfenden Vernunft preis;
es macht seine Geltung nur abhangig von der Note, die es beim Examen der Vernunft
erhalt.'' L. Feuerbach. a. St. Pag. 20.}; but precisely for that reason the true revelation
is, after all, designed to be a matter of freedom. He who knows himself so secure under the aegis of natural
reason that he does not consider it necessary to enter into voluntary captivity can, after all,
stay where he is, and arrange himself in the most reasonable way. The enlightened need not
become afraid at all: if revelation is too hard for them because it humbles reason; if negative
knowledge is too hard for them because it deifies reason, well, then they can seek their comfort and
reassurance where it is best to be found, namely in mediocrity. We can assure
the enlightened that the doctrine of the Middle Way shall not wound them with all too sharp consequences.
% --- p. 414 ---
\opage{414} They will find here the noblest mixtures and prettiest compounds of the
Generally-Natural and the Specially-Natural, the Merely-Rational and the A-Little-More-than-Rational, the
Wholly-Comprehensible and the Not-Quite-So-Wholly-Comprehensible, that anyone could wish for. The shell that
surrounds the kernel of mediocrity (if it really has any kernel) has precisely the
excellent property that it is neither quite hard nor quite soft, but tough. From this
it is explained that one can do almost anything one likes with this shell, that one can
squeeze it and press it, and knock it and pound it, as much and as long as one likes, and yet
never get the kernel out of it. Should the enlightened now wish to conclude from this that perhaps it has
no kernel at all, and therefore wish to be informed how matters then stand with
the kernel of the faith of the Gospel: the following remark might perhaps serve as preliminary guidance. ---
The whole present embarrassment and apparent bankruptcy of theology has without doubt been occasioned
by a little carelessness. Once, as it seems in an unguarded moment, one had the
misfortune to confuse a pair of concepts that ought not to be confused, namely \emph{rational
grounding and believing appropriation.} That these two concepts resemble each other in many respects
is evident at first glance; but that they nevertheless belong in quite opposite
quarters of the sky is not difficult to prove either. Believing appropriation divides into
two functions: \emph{submission to the divine authority and understanding of the
content of revelation.} Neither of these functions is conceivable without the cooperation of reason;
but natural reason works here in such a way that it continually cancels itself. --- To the
question: ``Shall I accept a revelation or not?'' reason answers quite straightforwardly:
``If you cannot, dare not, will not live without faith in a living God, then in God's
% --- p. 415 ---
\opage{415} name accept a revelation; but if you have accepted one, then above all do not forget
that you have accepted it on an unconditional authority.'' --- Just as the understanding can in many individual
cases find it comprehensible that it cannot comprehend, so sound reason can
also find it reasonable that it must deny itself for the sake of revelation, and put
authority in the place of ground. But once it has, after mature deliberation, taken such a step,
because it saw its absolute necessity, it will no longer engage in
groping after grounds where grounds are impossible. The ground of reason is thus, with reason's
knowledge, excluded once for all: authority rules with unlimited, divine
power; it does not let itself be bent, not shaken, but is and remains unchangeable, more
unchangeable than the fixed star in the heavens. And to the question: ``Which of the
different revelations then shall I accept, the Muhammadan, the Jewish, or the
Christian?'' reason answers in all simplicity: ``You shall accept the revelation that
most perfectly of all satisfies your soul's innermost need.'' Reason knows how to supply no other
criterion for distinguishing between the false and the true, the imperfect
and the perfect revelation. But this criterion is also sufficient when it is
applied rightly.

The Muhammadan, the Jew, and the Christian can naturally not agree in judging the
perfect revelation; but science with all its grounds and groundings is certainly
not able to bring about agreement either. It is not reasoning
interpreters of Scripture, but well-equipped preachers of repentance, who could come into
consideration here, insofar as the dispute can be settled by human assistance. The disputants must
not, then, begin by raising a swarm of grounds (such things only throw dust in
% --- p. 416 ---
\opage{416} %
the eyes); but they must all the better understand how to draw the bow of authority, and aim at
the conscience, so that the word may strike the right point. That Christianity is the perfect
revelation cannot be proved in any doctrinal system (for the doctrinal system is acknowledged only by those who
are already convinced beforehand); but the proof is conducted every time the Spirit of Christ penetrates into a
human being, and speaks to the heart in the mother tongue of conscience. ``But --- when one has once
for all presupposed the authority of revelation, when one stands on the standpoint of the Church and
speaks as a believer to believers, shall one then not be able to develop the inner rational
coherence of the doctrinal system, and explain the content of faith in the cognition of the idea, according to the old word:
\textit{credam ut intelligam?}'' --- The ecclesiastical standpoint! The Church stands on the
standpoint of the confession of faith. If we consider the many symbols, from the Apostolic down to the
Formula of Concord, it becomes clear to us that neither the individual articles of faith nor their
mutual coherence can, in the stricter sense of the word, become an object of rational
grounding. The individual articles of faith are dogmas, and a dogma is, after all, a proposition of authority.
Whoever will not accept the particular dogma the Church is wont to exclude from its community by the
little addition: \textit{anathema sit;} but \textit{anathema sit} expresses, after all, no rational
grounding. As regards the inner coherence of the dogmas, the one principal dogma cannot really be
derived at all from the other; but each takes its place as an independent authority beside
the other. One could for that matter just as well begin with the 3rd as with
the 1st article of the creed, provided one gets them all in; for the main thing is that
the one has the same originality as the other. --- ``But besides the various greater and
lesser symbols, the Church has, after all, its dogmatics too?'' --- Yes, it has; and besides
dogmatics and
% --- p. 417 ---
\opage{417} %
symbolics it has, moreover, both isagogics and hermeneutics and apologetics, together with
polemics and irenics and methodics; but --- \textit{hinc illæ lacrymæ!}

If we now turn to the other function in believing appropriation, namely understanding, the
ideal transformation, the true inwardizing of the content of revelation, it is easily shown
wherein this differs from rational grounding. In its scientific zeal to
comprehend Christianity's ``objective world of ideas,'' systematizing theology has on its own
authority (at least without any evangelical warrant) permitted itself that ``famous
separation of Christ's word and wonder'' by which both finally collapse in impotence. Since the
many different wonder-works naturally will not go along in the universal
development of the idea, ``rational grounding'' has in return let them stand for what
they are, ``supernatural individual utterances of Christ's person, well suited to yield individual
ideal reflexes, but nevertheless belonging most nearly to a historical-critical treatment.'' With the word,
on the other hand, it has gone the opposite way. ``Christ's word is, after all, eternal truth, and eternal
truth must, after all, be able to justify itself before thorough thinking, without regard to time, place
and person, solely by virtue of the idea. Instead of believing the word in all simplicity for Christ's
sake, one must then, by continued study, at last be able to bring it so far that one believes Christ
solely for the sake of the word, that is to say: for the sake of the idea'': now rational grounding is in full
swing. --- But is it not, after all, evident that by such a procedure one misses the
intended end? The purpose of rational grounding is, after all, to confirm the doctrine of faith; but
how in all the world can one think of confirming the doctrine of faith when one impairs and
weakens faith? According to the method stated here, both the word and
% --- p. 418 ---
\opage{418} the wonder must finally cease to be objects of religious faith: the word, because one lays so
one-sided a stress on its universality that it is absorbed as a moment in the idea;
the wonder-work, because one holds it in so contingent a form of singularity that it falls outside
the idea. Certainly, according to the Church's conviction, Christ's word is an eternal, unchangeable
word of truth; but this assurance is grounded solely on Christ's authority. The Church can in no
way permit that any ground of reason be interposed between Christ and his word. A similar
word, a similar truth, may have been put forward by other sages, and to that extent become an object
both of cognition and of recognition, but, mark well, according to rational grounding, not
according to faith; only when the congregation hears that it is Christ's word, or at least the word that has been
spoken in his name with the apostolic seal of the Spirit, is it transformed into a word of faith. All
long-winded reasoning about whether and how far it really is Christ's word, whether and how far one
is now to recognize the truth in its truth, etc., the congregation dismisses with the short answer: ``I
believe the word because I believe that it is Christ's word; in the word I believe the Lord Christ himself,
for Christ is, after all, the Word.'' Only when the congregation has in this way got peace from all
intrusive apostles of grounding does the next thing follow, namely understanding. \emph{My doctrine},
says Christ (Joh. 7, 16---17), \emph{is not mine, but his that sent me. If any man
will do his will that sent me, he shall know of the doctrine, whether it be of God, or whether I
speak of myself.} In these words we have the principle and the standard for all Christian
instruction, the so-called lower and the so-called higher, the practical and the theoretical.

The consideration of the individual wonder-works leads us to the same result. To every
objection against any of the
% --- p. 419 ---
\opage{419} %
evangelical miracle stories the congregation repeats the same answer: ``I believe this wonder
because I believe that it is Christ's wonder; in the wonder I believe the Lord Christ himself, for
Christ is, after all, the wonder of God's love.'' By thus cutting off the distracting
prolixity of reasoning, the community of believers closes more inwardly together, and the Church, despite the
storms of the age, stands unshakably firm on the foundation of the wonder.

\emph{What think ye of Christ? whose son is he?} This question, which was first
put to particular individuals, on a particular occasion, sounds through all ages with
the same originality and the same weight: the individuals change, but the occasion constantly returns.
Revealed truth cannot be obtained by an abstraction; it must be appropriated by a practical
application. Abstraction wants me to push the occasion aside and hold fast to the
universal, but application wants every subsequent occasion to be compared with the
first, wants the eternal word to work with its unchangeableness in all finite changes of circumstance,
and to shine through existence. It may be granted that the present has a different
circle of notions from the Jews in Christ's time; but this difference is of an all too finite
nature to effect any change in the conception of the Gospel's peculiar expression.
The interpreters' theories of accommodation with regard to the notions of the good angels, of
the Devil, of the demons, etc., are only empty evasions. The present is, as far as these principal
notions are concerned, surely just as much a minor as that time: there is only One who vouches
for their absolute validity, and that is Christ himself. The word of truth he speaks to that time,
he speaks to all times, and as he spoke it to that time, so he speaks it until the
end of days. The finite conditions of appropriation may change, but
% --- p. 420 ---
\opage{420} %
the word itself remains unchanged: the Spirit that works understanding will, in the right application,
each time assure us that the perfect word has been spoken once for all. With the perfect
deed the case is the same. That a human deed is imperfect is seen
not merely from its fitting only certain circumstances of time, and to a certain degree; but, absolutely
understood, the imperfection lies in this, that the agent, who is himself imperfect, is only
able to concentrate an imperfect force of personality in his action. Otherwise with the
divine deed. What Christ does, when he steps forth and acts in the Father's name, he
does not do merely from a momentary mood or out of exclusive regard for his nearest
surroundings; but he does it in the name of redemption, he does it as on behalf of the whole
human race; each time he acts, he acts in spirit and truth absolutely, and
the action then holds once for all. --- \emph{What think ye of Christ? whose son is
he?} Each time this question is repeated, the conditions for answering it must also be
repeated, and these conditions are at all times the same: \emph{the word and the wonder.}

If we institute a comparison between Jesus's various wonder-works, we meet at once the
peculiarity that, while some of them have a more immediately edifying character, others
seem on the contrary to be offensive, or at least almost offensive, to the present's general
sense of piety. Let us, e.g., put side by side the following two: \emph{``The raising of the widow's son''}
and \emph{``The cursing of the fig tree''}. The first of these miracles ``the theological
interpreter,'' however he may otherwise manage with ``the mysterious moment,'' will
doubtless find quite edifying, at least for the common man. But the second --- the cursing of the
fig tree (Mrk. 11, 13---14)
% --- p. 421 ---
\opage{421} --- it is hardly advisable to recommend for edifying use: what would the modern
consciousness, what would criticism, indeed what would Strauß say to this? ``One must,'' says this
critic, ``already find it strange that Jesus should thus have expected figs on the
tree out of season; at least, when he found none, he ought to have recognized the groundlessness of his
expectation, and refrained from so wholly unfair an action as the curse was. The moral
purpose of punishment, namely to bring the one punished to acknowledge his faults and to improve himself,
falls away completely when it is a tree that is in question; to become embittered at a lifeless
object can rightly be interpreted as a lack of culture'' etc. --- If ``the theological
interpreter'' now goes along with this conception, then on closer consideration he must surely also find the
first-named wonder-work (Luk. 7, 11---15) just as offensive: ``Was it perhaps no real
miracle? But then it is highly offensive that it has taken on the appearance of one. Was it a
real miracle; why then should precisely this man be raised from the dead? He was his
mother's only son, and she was a widow. Yes, that is all very well; but surely Omnipotence will not work
a miracle every time a poor widow loses her son? Not a word is even mentioned of this
widow knowing Jesus, or turning to him in faith; only so much is hinted, that she walked by
the bier and wept. \emph{And when the Lord saw her} (V. 13), \emph{he had compassion on her inwardly,
and said unto her, Weep not.} Christ thus lets himself be moved by involuntary compassion;
and, while millions of widows otherwise must weep at their sons' biers without anyone helping them
by a miracle, this one alone is helped as by a divine stroke of luck.'' --- But if now
both wonder-works, when seen in this light, are equally offensive, since in them
both there appears an infinite
% --- p. 422 ---
\opage{422} %
disproportion between means and end: should this not be able to awaken in ``the theological interpreter''
a suspicion that with his endeavor to find out the true end of the wonder he has gone
astray? Instead of seeking the fault in the Evangelists and complaining that ``the holy
men'' did not have culture and criticism enough to distinguish between facts and
parables, one ought, if not exactly from unconditional reverence for ``the holy men'' and ``the
holy books,'' then at least for the honor of ``hermeneutics'' and ``the Scripture principle,'' item ``the mysterious
moment,'' to reflect a little whether one could not after all discover that the fault lay much nearer, e.g.\ in
one's own eye. --- In judging the ``purposiveness'' of the wonder, theological half-criticism has
evidently gone astray in a double respect. On the one hand, one has in all historical dryness
undertaken an inspection of the various wonder-works, as of other outward
individual deeds, and then added up these many particulars as so many proofs of
Christ's divinity. This is now called in apologetics proving the truth of Christianity from
the miracles. On the other hand, out of a predominantly moral interest, one has transformed the
various individual deeds into nothing but symbolic intimations, abstracted their
universal moral significance, and appealed to a half determinate, half indeterminate total impression. This
total impression then works, in morals, a true glorification of Jesus's exemplar. That both
procedures are, however, equally unevangelical is surely self-evident. --- What a
scientific blunder, to wish to set up any single miracle as proof of Jesus's
messianic dignity! To form the premise of the proof, the miracle would in the first place have to stand
unshakably firm as historical fact; but doubt can be raised about every historical fact.
Supposing even that no doubt arose against
% --- p. 423 ---
\opage{423} the fact: it would still be impossible to infer from the nature of the wonder to
the messianic authority of the wonder-worker. Moses and the prophets, after all, also performed miracles. If now
each single proof that rests on the single miracle must be acknowledged insufficient, then
it avails but little to try to help the matter by putting together a plurality of such
miracle-proofs: a sum of insufficient proofs still gives no true conviction.
Just as superficial and one-sided is the symbolizing abstraction, by which the individual deed
is overlooked, and Christ's actual figure is volatilized into a dim general image whose
indeterminate features are lost in the mythical.

If, on the other hand, we understand that it is Christ's own self-manifestation that is the absolute end of the wonder,
then the abstraction is cancelled, and the individual deed coincides with the symbol. The raising of
the widow's son is in reality occasioned by Jesus's coming to the city gate just as they
were carrying out the dead man. The Savior does not place himself in so fantastic-objective an idea-relation to the
great whole, to the whole people and the whole of mankind, that ``a single, insignificant human being's''
grief should be indifferent to him. No, the Son of Man shares the conditions of finitude with the
lowliest human being. It really is the sight of the weeping widow that has moved him to compassion, yes,
an inward compassion. Moved in the mood of the moment, he then steps up to the bier, and speaks a
word, so that even deaf death must hear the voice of pity. --- ``But then Omnipotence is, after all,
set in motion to help this single one.'' Yes, certainly God in heaven will not disdain to
come to the help of the single human being (and then need not first ask either
the laws of nature or the universe for permission); but in spite of this there is here no disproportion whatever
between means and end. As the wonder is accomplished, the individual
% --- p. 424 ---
\opage{424} deed is at the same time raised above the contingent limitation of finitude, and posited as a symbol
of the divine compassion with the whole mortal race. In temporal respect the widow's joy
over her son is surely too insignificant to be set down as the highest aim of the wonder-work.
Enlarge the finite as much as you will: it remains too little. Put the Queen of Sheba
in place of the widow of Nain: it remains too little. Take a thousand heroes of freedom and lay them on
the bier in place of the widow's son: it remains too little. But precisely because any
relative end whatsoever is here too insignificant when separated from the absolute aim, precisely
for that reason what the world calls insignificant is significant enough, if only it can become a means for
the absolute. To do the perfect deed, Christ need not make alliance with
the millions, and set kingdoms and countries in motion; no, he can be content with the lesser: a withered
hand! a weeping widow who follows her son to the grave!

Just as the raising of the widow's son becomes edifying for the believer solely because it is Christ's
deed, so too the cursing of the fig tree receives its religious significance, and that for the
same reason. Although the action is symbolic, we must nevertheless not forget that it is also an
actual individual deed. According to the Evangelists' account (Matth. 21, 18. 19. Mrk. 11, 12. 13),
the cursing of the fig tree follows, after all, almost immediately upon the entry into Jerusalem. This
circumstance may well come into consideration when the occasion is to be determined more closely. The homage
the multitude had just shown Jesus by receiving him with palm branches and cries of Hosanna was,
humanly speaking, the greatest recognition any prophet in Israel could expect.

But, divinely understood, the reception has only a symbolic significance. The people seem to know,
but in
% --- p. 425 ---
\opage{425} reality do not know their Messiah; Israel seems, against all expectation, to grasp God's counsel, to be willing to
bear \emph{fruits worthy of repentance}, and that even at a time when, according to the usual
conditions, one should as yet expect no fruits. But it is only a semblance. --- We may well
suppose that such thoughts stirred in the Lord when he was with his disciples on
the road between Bethany and Jerusalem. With such a fundamental mood the momentary
impression united. He hungered, and as he hungered, he saw a fig tree afar off; and the fig tree
had leaves, and the fullness of the leaves entitled him to the expectation that this tree, notwithstanding the season,
must also have fruit (for a fruitful fig tree sets fruit as it sets
leaves); Christ thus in this moment looked on the outward, approached the tree, and was disappointed in
his expectation: \emph{he found nothing but leaves; for the time of figs was not yet.} Thus
this tree stands as a living image of disappointing Israel. As Christ sees the disappointment,
he is kindled with anger, and the disciples hear the words of the curse: \emph{No man eat fruit
of thee hereafter for ever!} sound over the tree. ``So Christ himself can become angry?'' Yes, in truth, they do not know
the gentle Jesus who have not heard the angry Christ. That pietistic sweetness has united
with modern sentimentality and pictured the Redeemer to us as the mild Lord Jesus who always
looks so gently cannot be imputed to the faith of the Gospel. When we (Matth. 11, 16---24) hear his complaint
against the perverse generation, his cry of woe over Chorazin and Bethsaida, his sentence of punishment over
Capernaum, his thundering speech against the Pharisees (Matth. 23), then we understand that he too, the
Meek One, the Lowly of Heart, has the consuming zeal of the old prophets (cf. Joh. 2,
17), that there is, in relation to men's perversity, indeed a wrath in him, a holy
% --- p. 426 ---
\opage{426}%
wrath that could strike the race with lightning and burn up the earth. But the zeal is mastered in
obedience; holy wrath must serve compassionate love. Instead of punishing
disappointing Israel, the wrath breaks out in a symbolizing action, and the withered fig tree remains
standing in the Gospel as a warning example. --- When the believers then gather to praise the
gentle Redeemer; when the congregation rejoices in its hope of salvation, and greets in the Spirit the King
of Glory with palm branches and Hosanna; see, then it stands there, the withered fig tree, and reminds us that
the King of Glory nevertheless is and remains a King of Truth, who never departs from his law, but seeks
the fruit where he sees --- the leaves.

The difference between rational grounding and believing appropriation can now be expressed
thus: grounding presupposes a ground of truth that is higher than Christ, and for precisely
this reason can never reach Christ. Faith presupposes that Christ is one with the truth itself,
and can appropriate the fullness of the truth only by appropriating Christ.

Grounding passes through two stages: \emph{finite reasoning} and \emph{the
speculative thought-process.} Criticism of the understanding, or finite reasoning, is recognized by this, that
one begins by dividing the indivisible, separating the inseparable. ``The word'' and ``the wonder'' and
``Christ himself'' one separates in such a way that in the end one cannot recognize the unity. In order to
understand the wonder, one wants first to know the word; but in order to know the word, one must first have
understood the wonder. The demand constantly flips over, and with this ceaseless flipping over a
settled knowledge becomes impossible. If now the word and the wonder are put together as the one term, and Christ
himself as the other, then the same difficulty repeats itself. In order to know Christ himself,
one must first know the word and the wonder; but in order to grasp the
% --- p. 427 ---
\opage{427} true significance of the word and the wonder, one must, after all, first have known Christ himself. In this way
the investigation of the understanding always goes around the center; reasoning runs round in a circle. ---
With due contempt for this endless understanding of finitude, speculation then takes a
decisive step, and will now grasp Christ in --- ``the absolute idea.'' But what then is this
Christ-idea? --- a hovering ambiguity! It does not help to try to throw one's weight about with
phrases and say: ``Christ is the idea, and the idea is Christ himself''; for there must, after all, be a
reason here why one still distinguishes the two designations: \emph{``Christ himself''} and \emph{``the idea''}.
Thus the alternatives stated above return again: either I must frankly admit to
myself that, when all is said and done, I really recognize Christ only for the sake of the
idea; or, conversely, confess it outright that in the idea I look solely and alone to Christ's
person. In the first case grounding rules, the emphasis falls on knowledge, and the
theological edifice of doctrine is surrendered to philosophy. In the latter case Christ's authority steps
into the place of the ground, and rational grounding must give way to believing appropriation.

The distinguishing mark by which Christ himself is set apart from all other extraordinary messengers of God
must, however, be sought chiefly in this, that he speaks the word not only in the Father's name, but also in his
own name; that he not only bears witness to the truth, but also bears witness to himself: \emph{I am
the truth} (Joh. 16, 6). Every other prophet must deny his I, and be content to use the word
as a means to reveal God's will; but Christ has from the Father the authority to
posit his I as the center of salvation, and to use the word as a means for his own
self-revelation. Likewise with the wonder. No prophet and no apostle dares step forth and say that
he does the wonder in his own name; but
% --- p. 428 ---
\opage{428} %
he must deny himself, and confess that he does it in the name of the Almighty, or in Jesus's
name (Ap. G. 3, 6); Christ alone has from the Father the authority to do
the wonder-work in his own name (Joh. 5, 20---23). {\large For the Father loveth the Son, and showeth
him all things that himself doeth; and he will show him greater works than these, that ye may
marvel. For as the Father raiseth up the dead and quickeneth them; even so the Son also quickeneth
whom he will\dots\ that all men should honor the Son, even as they honor the Father.} The wonder too, then,
Christ posits as a means for his personal self-manifestation. --- If it now stands firm that
the evangelical Christ is inseparably one with his wonder, as with his word: how then dare
half-believing half-criticism venture to take the Lord himself to task over the fittingness and
purposiveness of his individual wonder-works? If faith in Christ is an absurdity, well then,
let it be rejected together with its Gospel: there is, after all, always some sense in that. But to want to
stand on the standpoint of the Gospel, and yet reason about acceptable and rejectable
miracles, is an unmistakable sign, not exactly of immediate presumption, but certainly
of scientific thoughtlessness. Before the wonder-worker's supreme authority the self-conceited
understanding must bow, and be silent, as surely as Christ is --- the Lord Christ.

\begingroup\large
Matth. 9, 2---7. And behold, they brought to him one sick of the palsy, lying on a bed; and when Jesus
saw their faith, he said unto the sick of the palsy: Son! be of good cheer, thy sins are forgiven thee.
And behold, certain of the scribes said within themselves: this man blasphemeth God. And when Jesus saw
their thoughts, he said: wherefore think ye so evil in your hearts? For which is easier? to say:
thy sins are forgiven thee? or to say: arise and walk?
% --- p. 429 ---
\opage{429} But that ye may know that the Son of Man hath power on earth to forgive sins, then saith he
to the sick of the palsy: arise, and take up thy bed, and go unto thine house. And he arose, and departed
to his house. But when the multitude saw it, they marveled, and glorified God, who had given
such power unto men (cf. Mrk. 2, 1. Luk. 5, 17).
\par\endgroup

The vital nerve of this Gospel is touched palpably enough in the question: ``which is easier? to
say: thy sins are forgiven thee? or to say: arise and walk?'' --- If we listen to
reasoning criticism, then it is naturally easy enough to say the former, but on the other hand very
difficult to acquit oneself well of the latter. --- If Jesus shared the prejudice of his contemporaries and
countrymen that a man's illness must surely be a punishment for some sin or other, and
that the illness therefore could not be healed before this cause of the illness had been removed, then the
words: ``thy sins are forgiven thee'' come as of themselves. What, on the other hand, makes the narrative so
questionable is the manner in which the actual healing is said to take place. In the depiction of
the scene when the palsied man is brought to Jesus, Strauß finds a remarkable ascending ladder
among the three accounts. ``Matthew,'' it says, ``narrates quite simply how, when
Jesus had returned to Capernaum from an excursion to the other shore, they had brought him a
paralytic lying stretched out on a couch. Luke describes precisely how Jesus, surrounded
by a great crowd, especially of Pharisees and scribes, taught and healed in a house; and
how those who carried the paralytic could not, because of the crowd, get through
the door to Jesus, and therefore let the sick man down to him through the roof. If one considers the
structure of Oriental houses, to whose flat roof an opening leads from the uppermost story, and
takes
% --- p. 430 ---
\opage{430} into account the rabbinic usage, in which `the way through the gate' is opposed to `the way through
the roof' as a no less regular way, then with the expressions: `and let him\dots\ down
through the tiling,' one can hardly imagine anything other than that the bearers, having come either by a stair
leading directly from the street, or from the roof of the neighboring house, onto the flat roof of the
house in which Jesus was, let the sick man down together with his bed to Jesus, as
it seems with the help of ropes, through the hole that was already there in the ceiling. Mark, who
agrees with Matthew in setting the scene in Capernaum, and with Luke in depicting the
great throng and the ascent of the roof occasioned by it, goes, moreover, by fixing
the number of bearers at four, still further than Luke in this, that he has them, without regard to the
doors already at hand, break open the roof, and let the sick man down through the opening that
they themselves made.'' --- After thus having given the interpreters something to think about, and set up
so many obstacles at the entrance that pondering mediocrity must necessarily get
stuck --- in ``a hole of its own in the roof,'' the critic now passes on to treat the sick man's
healing. ``One has,'' he says, ``tried to present the outcome as a natural one by
declaring the sick man's condition to be a nervous weakness, in which the imagination that the illness
must continue as a punishment for his sins was the worst part; one has appealed to
similar cases of a rapid physical healing, and assumed a longer continued after-cure;
but even if there should be something true in the alleged analogies, it is nevertheless
incomparably more easily possible that healing stories of `the lame' and `the palsied' formed themselves in
legend according to the messianic expectations (Jes. 35, 6), than that they could actually take place.''

Whatever one may now judge of these critical intrigues:
% --- p. 431 ---
\opage{431} %
so much is evident, that the word cannot work where the wonder is misjudged. As one gets into
embarrassment with the healing, one gets at the same time into embarrassment with the forgiveness of sins.
The difference is only that the invisible word can more easily be explained away than the visible
fact. Tell me honestly what you think of these words of Jesus: \emph{arise and
walk!} and I shall tell you what you think in your heart of these words of Jesus: \emph{thy
sins are forgiven thee.} But precisely this reciprocal relation, in which the visible places itself as
mark and safeguard for the invisible, shows clearly enough, after all, that the wonder happens by means of the word,
so that the word in turn can work by means of the wonder. Instead of reasoning about
the lodging and the roof and the many accidental localities, which no antiquarian can any longer
determine with exactness, the faith of the Gospel fixes its gaze on the wonder-worker himself. And
when we then see the Savior in this scene; see the scribes and doctors of the law, who had come from
every town of Galilee and Judea and Jerusalem, sitting round him, when the power of the Lord was present
to heal; see the press of people in and outside the house; see the sick man let down together with
his bed through the tiling into the midst of the people before Jesus's feet, and hear the Master say,
when he saw their faith: \emph{Man! thy sins are forgiven thee:} then we know the
speaker by his voice. And when he thereupon puts the scribes to shame, and makes it known
to all \emph{that the Son of Man hath power on earth to forgive sins, and bids the
palsied man: I say unto thee, arise, take up thy bed, and go unto thine house,} --- so that the
palsied man immediately arises, and ``takes up that whereon he lay,'' and goes to his house
and glorifies God: then we understand that the Redeemer's visible deed is one with the invisibly
working word; then Christ stands before the inner sight as he was then, and still is, the
Son of the living God, mighty in word and deed.

% --- p. 432 ---
\opage{432} %
By being reflected in Christ himself, the wonder becomes as comprehensible to us as the word. ---
Christ feeds the hungry in the wilderness (Joh. 6, 35), and bears witness to himself: \emph{I am
the bread of life. He that cometh to me shall not hunger; and he that believeth on me shall
never thirst.} Is the latter now not at bottom just as difficult to comprehend as the
former? If one is content with a natural explanation of the miracle of the feeding, such as
\textit{Dr.} Paulus has prepared it for us, then one must also be content with an equally
meager explanation of the word about the bread of life. The same \textit{Dr.} Paulus who comprehends
the feeding as a humane use of provisions brought along also comprehends the word, after all,
as an ingenious use of a by no means unusual allegory, ``so verstund er (Jesus)
natürlich nicht ein Geniessen seines geistigen `Wesens' als eines solchen, sondern das
Aufnehmen seiner Einsichten''. If one asks what these ``Einsichten'' then really consist in,
particular weight must be laid on v. 32: \emph{Moses gave you not that bread from heaven,} which
in turn can be interpreted thus: ``the bread Moses gave you, i.e.\ the manna in the wilderness, did not (as
the fathers in their superstition have hitherto supposed) fall down from heaven.'' In this one then sees (if only one
looks rightly) ``eine Spur, daß Jesus zum Theil Naturkentnisse hatte, wie sie unter seinem
Volke ungewöhnlich waren!'' If one cannot now still one's hunger with the scant nourishment that
is offered in these ``Einsichten,'' but must, in order after all to sustain life, go a step further,
and seek the feeding with Strauß: will not the word then with this interpreter too obtain the same
value as the wonder? As the miracle of the feeding is transformed into a myth, the whole
discourse about the bread that comes down from heaven is likewise transformed into a myth: the nourishing power that
is ``the bread of life'' must be ascribed solely and alone --- to the mythical idea, since the Gospel of John
% --- p. 433 ---
\opage{433} naturally cannot hold its ground. In the eyes of the faith of the Gospel, on the other
hand, the whole takes on an entirely different appearance. With the inward confidence that the
Gospel of John, with God's help, will surely be able to hold its ground, the Believer does not
let himself be confused by the theologians' manifold attempts to get round the matter; but,
while he continues everywhere to look to the Savior, he at the same time keeps an eye on the
negative criticism, insofar as this criticism, by pointing out the exegetical misdirection,
must also \textit{nolens volens} be able to help him with one direction or another. If it is an
exegetical blunder to want to push back the immediate impression of the fact by keeping the
wonder at the greatest possible distance, then it is certainly also a blunder to want to bring
the individual wonder-work nearer to representation than the Gospel permits, by inventing a
multitude of finite intermediate links. The feeding miracle can serve as an example in this
respect. ``According to the number of groups acting in our narrative,'' says Strauß, ``three
views are possible: the multiplication took place either in Jesus's hands, or in those of the
distributing disciples, or in those of the receiving people. The last notion is partly petty to
the point of the fantastic, if one would imagine Jesus and the apostles distributing the pieces
sparingly, so that they might stretch, and these then swelling up in the hands of the
recipients; partly it would surely not even be possible to produce a piece, however small it
might be, from 5 loaves, which, according to Hebrew custom and since they were carried by a
boy, cannot have been very large, and still less from two fishes, for each individual of a mass
of 5000 men. Of the two other modes of representation I find it fitting, with Olshausen, that
the provisions multiplied under Jesus's creative hands, and that he gave the distributing
disciples one piece after another. One can then seek to picture the event in a twofold way, by
either imagining that,
% --- p. 434 ---
\opage{434} as often as the loaves and fishes came to an end, new ones had come from Jesus's
hands; or one can think that the individual loaves and fishes grew, so that, when a piece was
broken off, it was replaced again until by calculation the turn could come to the next. The
first notion seems to be foreign to the text, which, when it speaks of leftovers of the 5
loaves, hardly presupposes an increase in their number, and so only the second remains, by
whose poetic elaboration Lavater has only done the orthodox view a poor service. For this
wonder belongs to those that can seem in any way credible only so long as one knows how to keep
it in the half-darkness of an indeterminate representation: as soon as one would draw it forth
into the light, and examine all its parts closely, it dissolves into a phantom of mist. Bread
that swells up in the hands of the distributors like sponges one has moistened; fried fish
whose broken-off parts grow back by degrees, like the claw one tears from a living crayfish,
obviously belong not in the realm of the actual, but in an entirely different one.'' ---

For the pure cognition of the idea and the rational grounding such a caricature may well seem
outrageous; for the believing appropriation, on the other hand, it is in a high degree both
awakening and instructive. Had Christ, instead of blessing the five loaves and two fishes, at
the command of omnipotence let herbs and trees shoot up as in an instant from the barren soil;
had he created a paradise in the middle of the desert, that the hungry crowds might have eaten
of the fruits of the trees, and drunk from the clear springs (for, and this too the criticism
really ought to have noticed, --- the narrative reports nothing, after all, of the people
getting anything to drink), and rested in the evening hour behind the shelter of the foliage;
yes, had Christ proved his divine power merely by accelerating a natural process: would not
the interpreter then, with the help of a little imagination, easily have been able to find out
the intermediate links, and the rational grounding have
% --- p. 435 ---
\opage{435} enjoyed the pleasure of beholding the idea itself in one of its rich creative
manifestations! Alas, this pleasure is now spoiled. Who can fathom the idea here? Who explains
to us the hidden natural process? --- ``An accelerated natural process,'' says Strauß, ``it
would have been, if a grain had ever borne a hundredfold fruit in Jesus's hand and brought it
to ripeness, and he had poured out the multiplied grains to the people with ever-full hands, to
let them grind, knead and bake them, or let them eat them raw, and as they were, in the desert;
if he had taken a living fish, and suddenly called forth the eggs in its body, and made them
into full-grown fish, which the disciples and the people could then have boiled or fried.''
--- By such blunt objections the negative criticism does the Gospel a double service: it
strikes down the ``grounding'', and gives an impulse to awaken the believing appropriation. By
looking too much for ``the accelerated natural process'' it might easily happen to a spirited
observer, after all, that he got to see so much grain, and so many fishes, and so many lovely
fruits, that, enchanted by the sight, he lost sight of Christ himself altogether. But if the
miracle in this way brought it about that the many glories of the idea of nature overshadowed
and hid the wonder-worker, would it not then lose its evangelical character as a means for
Christ's self-revelation? Now, on the other hand, since the Savior does not at all conform to
``the demands of the idea'', but sets the finite natural conditions aside, lets the desert be
desert, and only out of inward pity for the hungry blesses the five loaves and two fishes, that
the people should not go away fasting and faint on the way; now the faith of the Gospel can ---
not explain the wonder ---, but see Christ himself transfigured in the wonder; and it is by
this, after all, that faith is to live. --- The same holds of the word. Had Jesus in his
discourse gone no further than to say: \emph{I am the bread of life. Your fathers did eat manna
in the wilderness, and}
% --- p. 436 ---
\opage{436} \emph{are dead. This is the bread which cometh down from heaven, that a man may eat
thereof, and not die;} then the scientific interpreter, without in any way violating ``the
mysterious moment'', would surely be man enough to transform all these words for us into fluid
moments in a clear-flowing cognition of the idea; and this so much the more easily as the
speaker at the end himself comes to meet the interpreter with the explanation (Joh. 6, 63):
\emph{It is the spirit that quickeneth, the flesh profiteth nothing; the words that I speak
unto you, they are spirit, and they are life.} Even if the interpreter took the life out of the
spirit by a ``reasonable'' re-explanation, it would hardly even be noticed. Now, however, the
Gospel's own concluding explanation is preceded by a couple of enigmatic words that must
occasion a stumbling-block and bring the interpreter of reason to silence (V. 51): \emph{and the
bread that I will give is my flesh, which I will give for the life of the world.} In these
words there is a sting that wounds reason; and we must surely thank the murmuring Jews, insofar
as their objection brings it about that the words do not flow imperceptibly away in the stream
of the discourse, but stand forth with unmistakable firmness and decisive determinacy, when it
says (V. 53---54): {\large verily, verily, I say unto you: except ye eat the flesh of the Son of
man, and drink his blood, ye have no life in you. Whoso eateth my flesh, and drinketh my blood,
hath eternal life; and I will raise him up at the last day.} At these words the grounding of
reason falls into impotence; but faith comes alive and rises up by means of the wonder, for the
wonder is the power of the word, and the word is life and spirit, and --- Christ is mighty in
word and deed.

``But,'' continues the reasoning criticism, ``all this may now be very good, if only it could
be proved that
% --- p. 437 ---
\opage{437} Christ really did such deeds, and spoke such words; yes, if one only knew whether
the Gospel of John could really hold its ground as well.'' --- We must much regret that we
cannot go into the many hundreds of half and whole reasons with which the learned criticism of
introduction marches up when it would pronounce judgment on the authenticity of the fourth
Gospel. But should anyone have penetrated so deeply into the critical introduction that,
exhausted by emptiness, he at last wished himself well out again, in order after all to find
Christ himself for once (in the criticism of introduction he certainly does not find him); then
we would warmly recommend to him the fourth Gospel, and especially the place where he will meet
the man born blind, just at the moment when the latter has been cast out of the synagogue (Joh.
9, 35---38). \emph{Jesus heard that they had cast him out; and when he had found him, he said
unto him:} {\large dost thou believe on the Son of God?} --- How then did the man born blind come
to believe on the Son of God? He does not yet know the one who has opened his eyes, and
therefore asks: {\large who is he, Lord, that I might believe on him?} The wonder in and for
itself was therefore not what was decisive; but just as little could the word be so. Had the
man born blind explained the healing in the same way as the Pharisees, then Jesus would have
made himself known to him in vain. Now, on the other hand, when he again heard Jesus's words,
he knew him by his voice. The same voice that said: \emph{go away, and wash in the pool of
Siloam,} the same voice now sounded again: \emph{dost thou believe on the Son of God?\ldots{}
thou hast both seen him, and it is he that talketh with thee;} then the man who had {\large
received his sight said: Lord, I believe! and he fell down before him.} --- Thus with the same
wonder by which Christ made the blind man seeing, with the same wonder he made (V. 39) the
seeing blind; and with the same word by which he took the
% --- p. 438 ---
\opage{438} seeing man to himself, with the same word he condemned the blind to remain in their
sin; yes, Christ is mighty in word and deed.

It lies in the nature of the case that no human character under the sun can be so clear, so
pronounced, so unequivocal and transparent, as the Revealer himself. But it also lies in the
nature of the case, after all, that no human being's life in the world can become the object of
such misjudgments and misunderstandings as the life of Jesus Christ. In order to grasp and
present a superior personality, peculiar in its kind, the one who sets himself the task must,
after all, know how to raise himself in inner ideal contemplation, if not above, then at least
to an equal level with his hero; he must, so to speak, make himself one with him, live his way
into his being, think, act, suffer, grieve and rejoice with him: without such an assimilation
the presentation will always remain superficial. But what observer is now able to set himself
on an equal level with Christ? What psychological observer is able to feel and think, act and
suffer with the Savior of the world? --- The reasoning interpretation always talks so much about
how the divine in Christ ought not to annul the concept of the human. This thought may in and
for itself be quite true. But the question is whether one has now also exhausted the concept of
the human? Has one now really gone oneself to the outermost limit of ``perfectibility''? Has one
calculated exactly how much of the divine pure human nature can contain within itself without
bursting? Does the interpreter perhaps know from his own experience how a human being fares who
is about to do a miracle? Can one not, will one not then see that the redeeming wonder-worker is
himself the living center of all the wonder-works, the central wonder that strikes down all
comprehension?

\emph{The Jesus of the Gospel hungers and thirsts like another human being;} I believe it, for I
believe the Gospel. \emph{He}
% --- p. 439 ---
\opage{439} \emph{fasts 40 days and 40 nights;} I believe it, for I believe the Gospel. But is
there now anyone who can explain to me how a human being who otherwise hungers and thirsts like
others can hold out fasting for 40 days and 40 nights? ``The interpreter'' remarks that the
Evangelist's words (Luk. 4, 2): \emph{And in those days he did eat nothing,} --- must be
regarded as a figurative exaggeration that does not exclude simple natural food (cf. Matth. 11,
19). There now! One thus escapes ``the mysterious moment'' by pleading ``a hyperbolic
expression''. If only one could always have ``a hyperbolic expression'' at hand\footnote{No,
rather, then, say it straight out with the Free Thinker: ``Too long a fast! The man-god would
surely have starved to death!''}.

\emph{Jesus sleeps in the ship;} I believe it, for I believe the Gospel. \emph{He stills the
storm on the Galilean sea;} I believe it, for I believe the Gospel. But when it is then asked
how it is conceivable that he who can sleep like a man can subdue the storm like a God, when
that exclamation of amazement: \emph{what manner of man is this, that even the winds and the sea
obey him?} repeats itself among men: what then has the interpreter to answer to this? Well, if
one is determined to answer, an answer will surely come. If one cannot find the answer in the
Evangelist, well, then one can surely find it in Schleiermacher. Schleiermacher is, after all, a
theological genius, and Christ, who speaks --- into the wind, when he would speak to the
disciples, is presumably also a theological genius. ``The mysterious moment'' thus lies in this,
that the Savior paid no heed to the tempest because he knew within himself that he could not
drown before his work was completed. Jesus's exalted calm, which forms so strong a contrast to
the disciples' anxious faintheartedness, is then grounded in his
% --- p. 440 ---
\opage{440} consciousness, indifferent to the storm, a consciousness that naturally
distinguishes every distinguished genius\footnote{No, rather, then, with the Free Thinker:
``Take heart, skipper, you carry Caesar and his fortune!''}.

\emph{Jesus walks on the earth, and as the Son of man has the weight of a real human being;} I
believe it, for I believe the Gospel. \emph{He walks on the sea in the 4th watch of the night;}
I believe it, for I believe the Gospel. But --- what a transition! What becomes of the
explanation now? --- After many theological interpreters have already made themselves
ridiculous by their acute attempts to find out an explanation, the half-criticism has at length
come to the result that Hugo Grotius, who is, after all, also an authority, may on this point
well be placed above the Evangelists, so that without harm to ``the mysterious moment'' one may
readily explain this event away, so as not to awaken doubt about Jesus's human
nature\footnote{No, rather, then, with the Free Thinker: ``This little story surely belongs among
the sailors' yarns.''}.

These parallels could be multiplied with many more; but that is not needed. ``The interpreter''
will nevertheless surely know how to keep up and master the extremes, and each time explain to
us how it really hangs together with the union of the divine and the human in Christ's person.
But we who believe the Gospel, and grow dizzy at the chasm that separates the Savior of the
world from the natural men of reason, we unfortunately cannot comprehend the explanation; alas
no, we cannot fill in the chasm with any cognition of the idea. We respect reason, we honor
acuteness, yes, we admire mediocrity; but --- may hermeneutics forgive us --- we worship Christ.
% --- p. 441 ---
\opage{441} Man's relation to God is the highest expression of the inwardness of
character\footnote{On the peculiar dialectic of the God-relation, cf. the Kierkegaardian
writings.}. The Believer understands that in his finite expressions of life he must in many ways
be misunderstood by others; partly because he has within himself a measure of his own for
existence; partly also because, by applying the ideal measure to himself without himself being
the ideal, he can come into manifold conflicts of frailty with the actual world around him. Now
Christ, who is the ideal itself, can indeed express at every point in finitude the pure,
absolute God-relation without wavering and without misdirection. But what thus from the one side
makes the Redeemer's self-revelation in word and deed intelligible to us, makes it in another
sense quite impossible for us to penetrate comprehendingly into Christ's being, and to explain
his individual words and actions from finite grounds. The God-relation by virtue of which the
Redeemer is the ideal is so wondrous, so incomprehensible, and so impenetrable for human
thought, that it raises him high above the reasoning of observation, and makes him the object of
faith and worship. If only ``the interpreter'' would not disdain the believing appropriation, he
would soon convince himself that the fourth Gospel can surely hold its ground, and be
strengthened in the inner man by giving up the reasoning crowd of the half-criticism, turning
toward Bethany, and going with Christ to the grave of Lazarus.

\emph{Gospel of John, Ch.\ 11. V, 1---46.}

\begingroup\large
Now a certain man was sick, Lazarus of Bethany, of the town where Mary and Martha her sister
were. (But it was that Mary which anointed the Lord with ointment, and wiped
% --- p. 442 ---
\opage{442} his feet with her hair; her brother Lazarus was sick.) Then the sisters sent unto
him, and had him told: Lord, behold, he whom thou lovest is sick. But when Jesus heard that, he
said: this sickness is not unto death, but for the glory of God, that the Son of God might be
glorified thereby. Now Jesus loved Martha, and her sister, and Lazarus. When he now heard that
he was sick, he nevertheless remained two days in the place where he was. Then after that he
said to the disciples: let us go into Judaea again. The disciples said unto him: Master, the
Jews of late sought to stone thee, and goest thou thither again? Jesus answered: are there not
twelve hours in the day? If any man walk in the day, he stumbleth not, because he seeth the
light of this world. But if a man walk in the night, he stumbleth, because the light is not in
him. These things said he; and after that he saith unto them: Lazarus, our friend, has fallen
asleep, but I go, that I may awake him out of sleep. Then said his disciples: Lord, if he sleep,
he shall do well. Howbeit Jesus spake of his death; but they thought that he had spoken of the
natural sleep. Then said Jesus unto them plainly: Lazarus is dead! And I am glad for your sakes
that I was not there, to the intent ye may believe; but let us go unto him. Then said Thomas
(which means Twin) unto his fellow disciples: let us also go, that we may die with him.
\par\endgroup

It would certainly be most unjust toward the theological interpreters if we, confusing the
matter with the persons, everywhere presented the exegetical blunders in such a way that a less
informed reader might construe it as if all recent theologians without exception only
misunderstood and distorted the content of the Gospel. Christian
% --- p. 443 ---
\opage{443} thinkers such as Olshausen, Tholuck, Neander, Lücke and others have in truth worked
faithfully against the rationalists' conceited criticism of the understanding, and vindicated
for many a gospel narrative its edifying character. But this circumstance is so far from
benefiting ``hermeneutics as science'' that on the contrary it only furnishes a new proof of how
often the heart is richer than science, and the master himself better than the art he
practices. --- Let one read the commentaries on the present Gospel, and one will have examples in
abundance.

The interpreters torment themselves to fathom Jesus's foreknowledge, and to find out the reason
why he still remained in Peraea two whole days after he had learned of Lazarus's sickness.
``What does Jesus really mean by the ambiguous declaration: \emph{this sickness is not unto
death?} Is he of the opinion that the sickness is not so very dangerous, and that the help is
therefore not so very urgent? In this case one can well explain to oneself that he still
remained in the place where he was for two days. But since on this assumption he therefore does
after all really begin with an erroneous conception of the nature of the sickness, the next
question arises: how then does he later (V. 14) come to know that Lazarus is really dead? Did
the sisters perhaps send yet another messenger? But of that, after all, the Evangelist reports
nothing. --- It seems far more reasonable, therefore, to hold to Christ's divine foreknowledge.
The words: \emph{this sickness is not unto death,} are then explained by the peculiar addition:
\emph{but for the glory of God, that the Son of God might be glorified thereby.} Jesus thus
foresaw from the beginning that Lazarus, even if his sickness was mortal, should nevertheless
not remain in death, but occasion the glorification of the Son of God by
% --- p. 444 ---
\opage{444} a great miracle. On this assumption one then need not ask about a new messenger:
Jesus learned of Lazarus's death in a supernatural way. But now other, and far greater,
difficulties arise. Why did Jesus (V. 4) give the messenger mentioned such an answer that it
must obviously disappoint the grieving sisters (for the brother's subsequent death seems, after
all, to be a factual refutation of Jesus's words)? Why did he remain two more days in Peraea?
Did he perhaps have far more important, urgent business there? But then this ought surely to
have been indicated in the narrative. Or did he stay away on purpose, waiting for Lazarus first
to die? But so arbitrary a conduct toward the family he (V. 5) loved so dearly seems unworthy,
yes, almost cruel.'' Thus one guesses in the study up and down, back and forth, putters and
potters with loose probabilities, until the negative criticism suddenly bursts in at the door of
the house, dumbfounds the interpreters, and reduces the whole event for us to a poorly told
``anecdote'' (Weisse); and there you have it: it is said that the fourth Gospel cannot hold its
ground.

If we now put the critical reasoning to rest, and let apprehending thought serve appropriation,
then it will also appear that the event, link by link, reveals Christ to the Believer. Jesus's
words (V. 4): \emph{this sickness is not unto death, but for the glory of God, that the Son of
God might be glorified thereby,} indicate the absolute standpoint under which what follows is to
be considered; it is like a verdict of eternity that sounds over the finite dispensation. The
question \emph{how} this verdict is to be understood more precisely admits of a plurality of
different answers. The disciples who stood by could understand these words in their way; the
anxious sisters could, when they received the news,
% --- p. 445 ---
\opage{445} explain them to themselves in their way, and, if we assume that the messenger had
returned before the death occurred\footnote{Which, however, is surely less probable; compare V. 6
and V. 39.}, the dying Lazarus could, after all, hear the same words, and comfort himself with
them in his way; it is possible, then, that these words undergo as many explanations as there
are persons who take part in the event. The different apprehension is conditioned by the
impression each individual soul receives. From this it follows in turn that those who heard the
words from Jesus's own mouth must surely have been just as eager to learn their final meaning as
the most zealous interpreters are to this day. --- On the possibility of such a many-sided
apprehension rests the infinite effect of the divine word; for each one judges himself according
as he judges of the word. But although reflection must move freely in order that understanding
may come, it does not revolve around a loose aggregate of finite reasons and counter-reasons;
rather, appropriation has its fixed point of support, its last and only decisive ground, in
Christ's authority. However you may understand these words in a finite respect, you must always
understand them as Christ's words. That on which everything depends is trust or mistrust toward
Christ. If we now presuppose unconditional trust in Jesus's utterance, then the finite How of
understanding will also gradually unfold from the event itself. With this we then pass on to the
next main point (V. 6). \emph{When he now heard that he was sick, he nevertheless remained two
days in the place where he was.} Concerning this point too, the interpreters must essentially
stand on the same level as the disciples and Lazarus's sisters. The question: why does the Lord
tarry in Peraea, although
% --- p. 446 ---
\opage{446} he knows that his friend lies at death's door in Bethany? no human being can answer.
The dying man's sisters, in all their anguish and sorrow, naturally could not explain it to
themselves; the interpreters, even if they would swear a thousand times by ``hermeneutics'' and
speculate on it till doomsday, cannot explain it either: for since no reason is given, every
attempt at grounding must naturally dissolve into an empty conjecture. On the other hand, the
interpreters, if they could only persuade themselves to look a little more at the \emph{Son of
God} than at ``hermeneutics'', could well explain to themselves why no reason is stated, and
would then not need, for that matter, to give up hope for the fourth Gospel. If the Lord himself
has given no reason (perhaps because his tarrying had its ground in his God-relation,
unfathomable for human beings, and, so to speak, belonged to his own inner self-preparation),
then the Evangelist is lawfully excused for stating no reason at all. To state a merely apparent
reason, as e.g.\ that Jesus was just then engaged in teaching the people and healing the sick in
Peraea, would only cause confusion. We do not doubt that Jesus was, after all, in great activity
when he received the message; for when did the Savior's work ever cease! \emph{my Father worketh
hitherto, and I work} (Joh. 5, 17); but just because the reason that could be drawn from this is
so general, it cannot turn the scale in a particular grounding. The Evangelist, out of
unconditional reverence for the Lord himself, has let him have his own reason; should the
interpreters not do likewise? --- As regards the third main point (V. 11---15), the task is in a
certain respect altered, and yet the conditions of understanding are at bottom the same. The
disciples hear from Jesus himself that Lazarus is dead; they do not brood over how the Master
has come to know it. On the other hand they cannot well accept that he will again come so near
to Jerusalem (in
% --- p. 447 ---
\opage{447} their opinion he should perhaps preferably let his wonder-working word act from
afar). The interpreters, on the other hand, will gladly concede the necessity that he must
himself go to Bethany, but they can now in no way comprehend how in all the world he can have
learned of the death. It is certainly not easy to comprehend, either. The interpreters must
really see to it that they do as the disciples did, bow in obedience under Christ's authority.
--- How far this simple counsel will appeal in any considerable degree to ``hermeneutics as
science'' is rather uncertain; but certain it is that in principle it will appeal to the faith
of the Gospel.

\begingroup\large
When Jesus came, he found Lazarus lying in the grave four days already. (Now Bethany was nigh
unto Jerusalem, about fifteen stadia off.) And many Jews had come to Martha and Mary, to comfort
them concerning their brother. Then Martha, as soon as she heard that Jesus was coming, went out
to meet him; but Mary sat still in the house. Then said Martha unto Jesus: Lord, if thou hadst
been here, my brother had not died. But I know, that even now, whatsoever thou wilt ask of God,
God will give it thee. Jesus said unto her: thy brother shall rise again. Martha saith unto him:
I know that he shall rise again in the resurrection at the last day. Jesus said unto her: I am
the resurrection and the life; he that believeth in me, though he die, yet shall he live. And
whosoever liveth and believeth in me shall never die. Believest thou this? She saith unto him:
yea, Lord, I have believed that thou art the Christ, the Son of God, which cometh into the
world. And when she had so said, she went her way, and called Mary her sister secretly, and
said: the Master is here, and calleth for thee. As soon as she heard that, she arose quickly,
and came unto
% --- p. 448 ---
\opage{448} him. (Now Jesus was not yet come into the town, but was in that place where Martha
met him.) When the Jews then which were with her in the house, and comforted her, saw Mary, that
she rose up hastily and went out, they followed her, saying: she goeth unto the grave to weep
there. Then when Mary was come where Jesus was, she fell down at his feet, saying unto him:
Lord, if thou hadst been here, my brother had not died.
\par\endgroup

It would not surprise us at all if the learned interpreters of Scripture should find that the
believing appropriation, as it has been characterized in the foregoing, makes the task all too
easy for itself, in that it at one stroke withdraws from the many exegetical difficulties. We
answer to this that all the learned contributions of science are received with gratitude,
insofar as they support and guide the evangelical train of thought; but that we cannot otherwise
assume that the Son of God came into the world in order to increase the ``exegetical
difficulties''. Seen from the standpoint of faith, understanding is and must be easy, for if it
were not both simple and easy, the Gospel could not, after all, be accessible to all. --- But in
another sense appropriation is, on the contrary, by no means easy, and it might perhaps be
counted among the most dangerous blunders of scientific interpretation that it makes difficult
what ought to be easy, and on the other hand makes all too easy for itself what is at bottom
difficult. To appeal to Christ's authority is no doubt very easy in a scientific respect, for
about the essence of faith pure knowledge can surely say nothing else than that faith is faith;
but to believe in Christ as God's only-begotten Son, and in this faith to have eternal life, is
in relation to existence and actuality so strenuous a task that it suffices for a whole life.
Every feature in our Gospel reminds us of this. Both the disciples and the sisters related
themselves believingly to Christ, after all,
% --- p. 449 ---
\opage{449} and yet it is as if faith would fail them in the moment of trial. The interpreter De
Wette, of whom it can by no means be said that he evades ``the exegetical difficulties'', seems
on the other hand to overlook the difficulties of faith and of appropriation when, on the
occasion of Jesus's words (V. 15): \emph{And I am glad for your sakes that I was not there, to
the intent ye may believe,} he cannot hold back the question: ``Aber waren die Jünger noch so
glaubensschwach, daß sie dieses sinnlichen Befestigungsmittels bedurften?'' Yes, so it goes. The
disciples are ``glaubensschwach'' when it comes to understanding the Lord; the theologians are
``glaubensschwach'' when it comes to understanding the disciples: alas, we are all
``glaubensschwach'', and our weak faith constantly needs renewed confirmation. --- So long a time
these disciples had now walked with their Lord and Master; so often had they heard his
prediction that he must go up to Jerusalem, and suffer many things; so sharply had Peter been
rebuked for his ``Lord, spare thyself'' (Matth. 16, 22), and still they cannot hold back their
objection (V. 8): \emph{Master, the Jews of late sought to stone thee, and goest thou thither
again?} Even after Jesus had reminded them how safely he walks who walks in the day, and then
declared to them plainly that he was really going to Bethany to awaken Lazarus and confirm their
faith (V. 11---15), they follow him with a heavy heart\footnote{The way in which Lücke and De
Wette take Thomas's utterance (V. 16) we cannot approve. ``Thomas,'' says D. W., ``spricht nicht
in halber Verzweiflung (Thol.): damit thut man dem ohnehin verrufenen Apostel Unrecht; er hegt
zwar, trotz der Beruhigung, V. 9 f., noch ernstliche Besorgniß, drückt aber eine edle
Entschlossenheit aus.'' But \textit{a)} the main question is not whether the disciples could
\emph{resolve} to go along with Jesus, for that surely went without saying. \textit{b)} In this
context the emphasis cannot fall on the threatening danger, for the Lord had, after all, at this
very moment sought to instill confidence in the disciples. \textit{c)} That Th. nevertheless sees
danger and death before his eyes proves that he has not heeded the Master's discourse, and
interprets his words as a doubter. --- ``Es offenbart sich in diesem Zuge des Thomas derselbe
Charakter, der sich in jener Erzählung, C. 20, 25, ausspricht. Er vermag es nicht sich über den
reflectirenden Verstand und die Gefahren, welche dieser sehen läßt, zum festen Vertrauen, zum
kindlichen Glauben zu erheben, der alle Bedenklichkeiten verbannt, sobald er sich auf göttliche
Verheissung stützen kann.'' (Thol.)}.

% --- p. 450 ---
\opage{450} As regards the pious pair of sisters, their concern of soul is here no doubt quite
different from that of the disciples, but with faith itself it goes with them nevertheless much
the same way as with the latter. Before the interpreters pronounce final judgment on the fourth
Gospel, we ask them to cast one more glance at these sisters, whose image has been caught not
only by the Evangelist John, but also, though in a different light, by the Evangelist Luke (10,
38---42). The scene of which Luke gives an account reminds us of an earlier time, when Jesus is
received as a guest in Martha's house. Tired from the journey, the Lord rests here in this lovely
family circle, and prompts each of the sisters to bring her own peculiar cast of soul to light.
By the brother's death everything is now changed in Martha's house, and yet we still recognize
both sisters in their mourning dress. The same Martha who in her domestic sphere was so
restlessly active, and out of loving care for the precious guest ``was cumbered about much
serving'', is here too, in her deep sorrow over the beloved brother's death, still restlessly
active, and seeks her comfort in much communication.
% --- p. 451 ---
\opage{451} The same Mary who ``sat at Jesus's feet and heard his word'', and forgot the outward
manifold over the one thing needful, sits here too, in the gathering of mourners, sunk in
inward, quiet pain, and forgets the outward, and pays no heed to what goes on around her.
\emph{Lord, if thou hadst been here, my brother had not died,} with these words Martha (V. 21)
receives the bitterly missed helper, and in the same outburst Mary (V. 32) voices her grief to
the divine Master. But while Martha, however much she wavers in mood, nevertheless holds herself
upright, and strives in conversation (V. 21---27) with the Lord in order to win a hope, feeling
has so overwhelmed Mary that she must at once fall down at his feet, and give vent to
tears\footnote{If we further compare Joh. 12, 1---3 with Luk. 10, 38---42, the inner agreement
between these character-portraits will be perfectly evident.}.

\begingroup\large
When Jesus then saw her weeping, and saw the Jews weeping which came with her, he was moved in
the spirit, and was violently stirred, and said: Where have ye laid him? They said unto him:
Lord, come and see. Jesus wept. Then said the Jews: behold how he loved him! And some of them
said: could not this man, which opened the eyes of the blind, have caused that even this man
should not have died? Then Jesus, again moved inwardly, came to the grave. It was a cave, and a
stone lay before it. Jesus said: take ye away the stone. Martha, the sister of him that was
dead, saith unto him: Lord, by this time he stinketh, for he hath lain four days. Jesus said
unto her: said I not unto thee, that, if thou wouldest believe, thou shouldest see the glory of
God?
\par\endgroup
% --- p. 452 ---
\opage{452} That Jesus ``was moved in the spirit when he saw Mary weeping, and saw the Jews
weeping'', that he was ``violently stirred'' and --- ``wept'', that he was ``moved inwardly''
even at the moment when he stood at the grave in order to accomplish the wonder of the raising,
this apparent contradiction must naturally weigh upon the interpreters with many ``exegetical
difficulties''. One may now well appeal here to the concept in dogmatics of Christ's twofold
nature, and explain the relation to oneself thus: Jesus weeps with those who weep, because as
man he has sympathy with human beings; he has the stone taken away from the cave of the grave,
because as the Son of God he knows himself to be the Lord of death. But although this
explanation, considered proposition by proposition, is entirely correct, it nevertheless suffers
from a very serious weakness, and cannot possibly furnish any rational grounding. --- As regards
the concept in dogmatics of the God-man, we must remark in the first place that this concept is
so far from standing higher than the evangelical account that it is rather an abstraction from
the confession of faith and the Gospel. If now the fundamental concept thus abstracted has in
itself no independent ground in reason, it further betrays its impotence in that its scientific
content dissolves into general propositions, and to that extent must lack all the immediate and
indispensable intermediate determinations of actuality. That the Son of God could with inward
sadness feel the grief and distress of mortals in a moment like this, when he stood at Lazarus's
grave, can indeed at a pinch be explained from the general concept of a God-man, for the quiet
sadness of sympathy presupposes precisely the infinite superiority; but how Christ could be
moved so strongly that the words of the original text even point to a spiritual shaking cannot
be derived from any dogmatic concept. Olshausen's pious remark, ``daß er über den Tod und dessen
Schrecknisse überhaupt als Sold der
% --- p. 453 ---
\opage{453} Sünde geweint habe'', does not, according to De Wette's critical remark, lie in the
text, and even if this ``überhaupt'' were actually put into the text, it would nevertheless, as
a merely general determination of reflection, weaken the immediate impression.

Only when we exchange the rational grounding for the believing appropriation can understanding
come. --- Who bears the greater burden: the person who is helped by another through a wonder's
taking place, or the person who can help another only through its being granted him to do a
wonder? To a cursory view it might well seem as if only the first is to be called a sufferer,
since he is one in distress. But if we look more closely, the relation will turn out to be the
reverse. He who is to be helped by another through a wonder-work at once casts part of his burden
over on the latter, and yet is still as though he alone were the burdened one, because he cannot
in himself overcome the selfish. But he who is called to render another divine help by means of
the wonder must bear the whole burden, and yet be as though he were not burdened at all, because
he must in himself annihilate the selfish. Insofar as the wonder-worker is a human being, he
cannot, after all, do the wonder by his own power (for by himself he can do nothing); but if the
wonder is to take place by God's power alone, with what unconditional trust must this human being
then not rely on the Almighty's aid? Unconditional trust in turn presupposes an unconditional
obedience: the slightest doubt can bring God's severe punishment upon the
wonder-worker.\footnote{Thus Moses was punished by Jehovah because he doubted when he was to
strike the rock with his staff (4 Mos. 20, 9---12).} --- ``But,'' someone might say, ``even if
this explanation could be applied to Moses and the prophets, it cannot be applied to Christ.
% --- p. 454 ---
\opage{454} Moses and the prophets were frail human beings like the rest of us, and to that
extent had to bear heavily the burden of obedience and self-denial; but Christ is, after all,
the perfect man, in whose sinless nature the selfish principle never rules; for him the burden
must be light. Moses and the prophets surely stood in an inward relation to the Almighty,
insofar as signs and wonder-works were done through them; but this relation was nonetheless a
relation of faith. Now since faith in every human being presupposes the possibility of doubt, it
is to that extent intelligible enough that those chosen men must have felt themselves in great
tension and secret anguish of soul each time they found themselves called upon to attempt a
wonder-work. But with Christ it is an entirely different matter. His relation to God is, after
all, not a mere relation of faith, but a relation of knowledge and certainty high exalted above
faith and doubt, grounded in the Son's unity of essence with the Father.'' --- It is chiefly the
dogmatizing speculation which, in letting the impression of Jesus's human existence vanish into
immanent, metaphysical divinity, occasions such objections. If it is true that he who helps by a
wonder bears a greater burden than he who is helped; if it is true that the burden of actuality
and finitude, humanly speaking, does not diminish but increases with the God-relation itself:
then it is surely also true that the Savior of the world bore a world's burden each time he
stretched out his hand to help the suffering by a wonder. It is such a conception that our
present Gospel presupposes and confirms from all sides. Through the individual human beings
Christ everywhere enters into an immediate relation to the whole distressed human race. What
feelings must the news of the sick Lazarus not have awakened in his soul! Could the concern of
those who took part, could the sisters' anguish and grief even come into consideration, when we
truly
% --- p. 455 ---
\opage{455} want to think of the Redeemer's inward compassion for the dying friend? Jesus sees
that Lazarus's sickness must serve the glory of God, that help must come by a wonder; and yet he
does not follow the immediate impulse of sympathy, does not at once do the wonder, does not
hasten at once to Bethany, but keeps still, and remains two more days in the place where he was.
This lingering slowness affects painfully the natural heart, which longs impatiently for help.
All the sympathizing friends in the town ``where Martha and her sister were'' would surely, if it
were possible, cry out to him: come, O come, before it is too late! the grieving sisters would
surely cry out to him: come, O come! we beg thee for thy mercy's sake, come, O come, we implore
thee by thy friendship for our brother Lazarus, come, O come, before it is too late! --- So the
children of men seem, then, to be more compassionate than God's own Son. This compassion is the
weakness of the distressed, by virtue of the selfish. But the Redeemer dares not follow any
selfish impulse; what is to happen must happen to the glory of God; so he must subdue the
impatient longing of human sympathy, so he must deny himself, and bear the burden of
self-denial. Slowly and calmly he prepares the wonder, slowly and calmly he proceeds to its
execution. And while the helper approaches so slowly, the grieving hearts are near to despair in
grief and pain; but in whose soul the pain has after all penetrated most deeply, only One knows
within himself. And while the divine Comforter approaches the grave step by step, the many who
take part are united in lamentation and tears, as though they were all suffering; but who among
these sufferers suffers most, only One knows within himself. Behold, the Savior of the world
bears the heavy burden of compassion as he walks the slow way to the grave of Lazarus. So heavy
was this burden that the Son of man was near to sinking under its weight, so
% --- p. 456 ---
\opage{456} inward was this pain that the Lord \emph{was moved in the spirit and was violently
stirred,} when he saw the deeply grieving sister weeping, and saw the Jews weeping which came
with her. Thus we can understand, then, that for a moment he had to bow under the burden of
sympathy, and sink into the pain of finitude, and weep with those who weep.

\begingroup\large
Then they took away the stone from the place where the dead was laid. And Jesus lifted up his
eyes and said: Father, I thank Thee that Thou hast heard me. And I knew that thou hearest me
always; but because of the people which stand by I said it, that they may believe that thou hast
sent me.
\par\endgroup

The Believer who in this scene hears the words of the prayer cannot be disturbed by the modern
buzzing of criticism. Let the half-Hegelian Weisse call this prayer ``ein Schaugebet'', it does
no harm; but that the theological interpretation joins in, and would criticize ``Jesus's
consciousness'' for us and school the Evangelist's presentation, cannot in a scene like the
present one but have a disturbing effect. If hermeneutics has really fathomed the secret of the
God-man, then we ask it in the name of science to make its new discovery known as soon as
possible\footnote{``\emph{Aber um des dastehenden Volkes Willen sagte ich es, damit sie
glauben}'' u. s. w. Das Unpassende dieser Wendung muß man anerkennen, ohne mit Diefenbach
(Berth. kr. Journ. \textit{V}, \textit{1}. S. \textit{8}). V. 42. für unächt, oder mit Strauß
\textit{II}. \textit{162}. die ganze Erzählung für unhistorisch erklären zu müssen. Man hat
nur (wozu Lücke \textit{I}. \textit{A}. geneigt war) zuzugeben, daß der Evang. Jesu diese
Worte geliehen habe, und zwar dies Mal nicht aus seiner subjectiven Ansicht, sondern indem er
sich richtig in}; but if it is only the usual
% --- p. 457 ---
\opage{457} interpretive reasoning that will make itself heard again, then we bid the theologians,
in the name of the faith of the Gospel, with seemly reverence for the holy, either to take
themselves away at once from Lazarus's grave, or at least to keep still and quiet while the
Lord prays.

It is a deceptive fancy that the criticism which places the observer in a certain
spectator's calm outside the event should be better suited to assess the facts of
revelation than the criticism which leads the observer as a participant into the event, and
strives for appropriation. Not by laying faith aside, throwing himself back into the
armchair of contemplation and surrendering to the mode of thought of the age can the
interpreter learn to expound the wonder of the raising; but by bowing in faith before
Christ, and inquiring of the mode of thought of eternity, can a simple person appropriate
both the word and the wonder. Would you grasp Jesus's words to Martha: \emph{I am the
resurrection and the life}, then your mood must be like Martha's, consecrated to death.
Would you rightly consider why the Savior weeps, then let your heart first weep with Mary;
would you understand the prayer of thanksgiving that sounds over the grave, then first
perceive in Martha's anguish --- a smell of the dead from the grave! --- How could it be
expressed more clearly that the word and the wonder in their supernatural interaction must
serve the Redeemer's self-manifestation than when he himself bears witness in prayer:
\emph{because of the people which stand by I said it, that they may believe that thou hast
sent me.}\footnote{Jesu
Bewußtseyn versezte, nur aber den Fehler beging das, was in Jesu Seele lag, laut
aussprechen zu lassen. Denn in der That dürfen und müssen wir uns Jesu Bewußtseyn von Gott
und seiner Stellung zu ihm als ein ununterbrochenes, nicht an einzelne Gebetsmomente
geknusptes denken.“ De Wette.}

% --- p. 458 ---
\opage{458} %
\begingroup\large
And when he had said this, he cried with a loud voice: Lazarus, come forth! And the dead man
came forth, bound with graveclothes about feet and hands, and his face was bound about with
a napkin. Jesus said unto them: loose him, and let him go.
\par\endgroup

If the word of life is to release us from the power of death, it must be preached in power;
but if it is to be preached in power, it must be preached in deed. The Lord preached the word
to Martha before he went to the grave; but it was as though the word had not the right
power. So then he goes to the grave and preaches the same word, and behold, now it works in
power. He preaches it so loudly that the world must after all hear: Lazarus hears it, and
rises up in the cave within. He preaches it so visibly that the eyes of all must after all
see, yes, see the dead man, bound about with graveclothes on hands and feet, step out of the
grave and greet --- the light of the world. Thus is the word preached in power, the word of
the resurrection: \emph{I am the resurrection and the life. He that believeth in me, though
he die, yet shall he live. And whosoever liveth and believeth in me shall not die
eternally.}

Christ is mighty in word and deed.

\subsection*{\textit{VII.}\\[4pt] Jesus Christ is the Savior of the World.}
\addcontentsline{toc}{subsection}{VII. Jesus Christ is the Savior of the World}

Savior of the world! Will the world now really also admit that it needs a Savior? ---
``That surely depends on how one understands it. A savior is, after all, always a help; but
the help must suit the need.'' --- What is it, then, that the world needs? --- ``That is not
easy to say. In the world everything is parceled out and divided: men's need is divided
too. This much is
% --- p. 459 ---
\opage{459} certain, however, that each feels his own need best, that no one needs what he
has, and no one has what he needs.'' --- What, then, is the world's drive? --- ``Well, that
is not easy to say. The drive is just what instinct is; instinct is its own master. This much
is certain, however, that the drive drives wherever it will, that each must follow his own
drive, and see to it that he drifts with the times.'' --- What, then, is the world's highest
law? --- ``Well, that is not easy to say. The freedom of the will is the ground of the law;
but each has his own will. When, then, the will turns (which often happens), the ground must
of course turn along with it. This much is certain, however, that the law must bend to the
will if the will is to comply with the law.'' --- What, then, is holy to the world? ---
``Well, that is not easy to say. Holy is everything that I worship and cultivate; no, the
world is not ungodly, but who can count its gods? This much is certain, however, that there
is a difference between private and public worship. In silence each worships his own god;
in public one now worships the god who has won the most votes.'' --- But what in all the
world is this world, then, this real world of men? --- ``Well, that is really not easy to
say. The world is as one takes it; different people take it differently. People who go in
for the practical like to take it as a mixture of evil and good that falls somewhat outside
the concept; but the philosophers, who rest in the theoretical, prefer to take it such that
it coincides precisely with the concept. Some, then, hold that the world by virtue of the
concept is One and All, others that it is pure nothing; but the poets and the fine spirits
are all agreed that this world is motley.''

--- To a superficial glance the glory of the world is beyond all survey. The eye is dazzled
by the show of phenomena, and memory despairs over the multitude of events. The comparing
understanding would hold fast to the most important, the most important of all. In vain!
What
% --- p. 460 ---
\opage{460} comes is always more important than what vanishes: the latest time, the latest
year, the latest day is precisely the most important, the most important of all. --- ``The
lengthy comparison thus brings only a waste of time; one can straightaway lay down the rule:
\emph{the present is always the most important.} By thus going somewhat thoroughly to work
with the inquiry, one will also find that there is reason in this world, that it has, in
other words, awakened to consciousness. That the consciousness of the world is one with the
modern self-consciousness goes without saying (if the world did not reveal itself with life
and soul in all that is modern, how else should it reveal itself?)''. --- Then let this
consciousness speak itself out, let the World-Spirit itself instruct its disciple!

\begin{center}
  *\quad*\\[-4pt]
  *
\end{center}

``The principle on which all existence builds, and has been built from eternity, i.e.\ from
the beginning, is laid down in the few but significant words: \emph{everything must happen
naturally.} That this proposition needs no proof follows directly from its being a
principle. The observations of natural philosophy, daily experience, sound understanding,
and above all immediate sense teach us that the proposition is infallible. But what good is
it to have made such a discovery if one does not know how to apply it properly? If the
principle is to be made fruitful for knowledge and for life, it must rule in its whole extent
over religion just as well as over poetry and art. --- Everything must happen naturally: how
fruitful this proposition is, when it comes to a criticism of revealed religion! As Minerva
once sprang in full armor from Jupiter's head, so there now springs from this principle one
well-armed chief truth after another. Each of these truths at once casts off the wrappings
of reasoning, and stands before clear consciousness in full
% --- p. 461 ---
\opage{461} force, and as at one stroke; for the world has become conscious of itself. Here
they are, then, these short but comprehensible and pithy truths of self-consciousness, as
they follow one another in logical order:

\begin{center}
  \textit{1)}\quad\emph{There is no God.}
\end{center}

Religion speaks to us of a living, omnipresent God, an almighty Creator who sustains and
governs all things. Revealed religion has in its way a right to speak thus, insofar as it
aims at flattering feeling and elevating the spirit. But real truth there nonetheless is
not in this talk. Imagine the God of religion however one will, he must always remain a
supramundane God; but a supramundane God is, of course, a supernatural God. --- Religion
teaches concerning the creation that everything from the beginning happened supernaturally;
but our firm principle overturns this doctrine. Everything from the beginning has happened
quite naturally! motherly nature brings forth all things, and makes the almighty Creator
impossible.

\begin{center}
  \textit{2)}\quad\emph{There is no eternity.}
\end{center}

Religion teaches that after the death of the body the soul passes over into a higher
existence: revelation promises us the resurrection of the dead. It is an age-old wisdom, a
most comforting teaching; one should not disdain it, still less despise it. But real truth
the teaching does not have. Natural life is not eternal; how, then, should eternal life be
possible, when everything happens only naturally? That the body can do nothing without the
soul is granted by all; but let us now also consider this: what, pray, can a soul do without
a body? Without a heart I cannot feel: without a brain I cannot think. When the body falls
into a faint, the soul swoons: death is the last swoon.
% --- p. 462 ---
\opage{462} Should the soul now awaken again from this swoon, it could happen only under
this one presupposition, that it got a new heart, thus a heart that is not of this nature,
and a new brain, thus a supernatural brain; but this assumption is impossible, for
everything must happen naturally.

\begin{center}
  \textit{3)}\quad\emph{There is no conscience.}
\end{center}

When religion enjoins upon us the difference between good and evil, right and wrong, it
proceeds from the presupposition that God, the Holy One, the All-knowing, knows every human
heart. Man's knowing with himself is then inseparable from God's knowing with man: man knows
with God and with himself what is right and wrong. This immediate knowing-with is the right
and proper expression for conscience. By thus paraphrasing man's natural consciousness of
good and evil, religion can undeniably do morality a great service; for most people are
the better for living on in secret dread of an All-knowing One. But truth this doctrine
nonetheless does not have; for --- everything must happen naturally. --- When religion
really wants to stir feeling, or, as it is called, `awaken the sleeping conscience', it
speaks in lofty tones of man's inner, infinite self-worth, of retribution in the beyond, of
the last judgment. On this point too we must acknowledge the beneficial influence of
religion; for without such an apparatus of supramundane motives few people would do their
duty in the world. But real truth ought not to be sought in such figurative talk. Mediocrity,
to be sure, now holds that a really thorough reasoning from natural grounds of reason will
surely be able to vindicate the right of conscience and the inviolability of the law of
duty, without one's precisely needing to touch too hard on the supernatural. But the only
result that comes
% --- p. 463 ---
\opage{463} out of the whole demonstration, when one takes it really stringently and
really scientifically, is without doubt --- mediocrity itself. Should the law of duty really
have true and eternal validity by virtue of conscience, then conscience would of course at
the same time have to have a supernatural validity by virtue of duty; but, as surely as
everything happens only naturally, no conscience shall survive death.''

``The great, simple truths of nature that we have here considered all bear the same stamp,
cold as ice, unmixed and pure as the newly fallen snow. `What would the many busy people
judge, if they once had time to stand still and really consider these stern Norns of
existence?' O you who do know a little of the way of the world, do not wonder that the
children of men cannot press the ice-cold truths of the mystery of nature to the warm heart
without shuddering; but only take courage, hold the thought fast, and learn, then, from the
spirit of this world what wisdom gives. In this respect too it holds that everything must
happen naturally. The spirit of the world is one with the spirit of humanity, and becomes
conscious of itself only in individual gifted men. The great multitude must at all times
live on in a spiritual half-darkness. Few in number are those who feel the strength and
courage to rise up into the cool heights where the eye not only surveys but also sees
through existence. Even these few are granted only at brief intervals to breathe the ether
of knowledge; they must then go down again into the warmer regions of feeling, into the
throng of men, must feel their limitation, and resign themselves to working under the
pressure of finitude. This explains how the idea of the world can become conscious of
itself in the individual man without this man being therefore able to make himself one with
the World-Spirit itself. At this point you will now also understand that pure wisdom can be
both affirmed and
% --- p. 464 ---
\opage{464} denied for the whole world without yet doing the world harm. --- Were
revelation to gain title to an eternal truth, it would look sad for the present world. Stern
is the holy law, terribly stern the judgment of righteousness: how many brilliant passions,
how many proud exploits, how many historical works would it not stunt? Human nature is
already weakened enough under the influence of religion; were the ideals of eternity to
weigh with still greater weight, the race would have to languish and sink under the burden.
See, therefore it is wisely arranged that the number of the pious never becomes
disproportionately large; that even faith in the believers never wholly annihilates, but
only alters and rather heightens, the selfishly veiled worldly drive. How many people, pray,
express in their lives the truth that there really is a God! The few who try it cannot, of
course, get through in their lifetime (unless it should be by a kind of misunderstanding).
The understanding of the world will always know how to construe such religious characters
so that their example does not work too strongly, but so that they, as rare exceptions to
the general rule, fall under a certain species of the so-called holy madness. --- But now
look also to the other side. If the unconditional natural truth of the denial of God
suddenly got through with everyone, without a corresponding improvement of culture being
attained at the same time, would that not likewise look dangerous? Human nature is, once and
for all, so strangely constituted that, although there really is no God, it nonetheless
cannot do without the comfort of religion. Could the multitude convince itself that there is
no God, then lawlessness, willfulness and insolence would rise to a height hitherto unknown;
then the works of culture and the ideas of history would soon go under in the coarseness of
sensual lusts, in the barbarism of material pleasures. The multitude is already rebellious
enough: it cannot bear the pure
% --- p. 465 ---
\opage{465} truth. Therefore it is now also wisely ordained that the denial of God never
exceeds its proper measure. Few are they who by the power of thought can reach the deeper
grounds of atheism; fewer still they who are able so to master all the prejudices of blood
and habit that their life expresses the truth: there is no God. Good-natured people of
feeling lay their hand on their heart and assure us that there can be no real denier of God.
O, how edifying! The World-Spirit must smile, for these fools are in a way wise. People are
as a rule just as soft when it comes to a decisive denial as when a decisive affirmation is
demanded. People have on the whole more than enough in their practical affairs to settle and
decide among themselves: absolute decisions they can confidently leave to history and fate.
--- See, this is now in sum what the world-consciousness knows of the highest wisdom; if you
would go further, then see for yourself, think for yourself, test for yourself whether you
can perhaps also grant what follows.''

``If religion can never be wholly dismissed, and yet cannot rule unconditionally either, then
its content must naturally continue to keep midway between revelation and myth, and the
religious consciousness constantly hover between affirmation and negation. This dialectical
hovering is conditioned by the fact that the human spirit works both as feeling and as
understanding. Insofar as feeling has a one-sided preponderance, the individual will become
only a believer; but if understanding comes to dominion, man must give up faith and become a
knower. For the believer myth is a revelation; for the knower revelation is myth. Both forms
of consciousness are so natural to man that the one of them could not even properly develop
without presupposing the other: they are related to each other as the male nature to the
female. Woman's
% --- p. 466 ---
\opage{466} consciousness is essentially mythical; hence it is that religion must always be
the absolute for her. An irreligious woman and a corrupted woman are at bottom synonymous
concepts. In the man, on the other hand, understanding ought to rule; therefore he stakes
everything on the real, translates revelation into myth, and works as a knower. An
immediately religious man and a weakened man are, in the deeper sense of self-consciousness,
synonyms. Speak to a noble woman about the fundamental relations of existence, about the
holy mysteries of life, then tell her that everything must happen naturally; she will not
understand you, nor does she dare to. But now turn with the same speech to the man: his
smile will at once tell you that he has grasped your meaning and understood the secret of
the principle. The difference in principle between the two sexes must not, however, be held
fast too one-sidedly. Were there really so deep a chasm between the ideal natural
determination of man and of woman, they would of course have to repel and misjudge each
other forever; but daily experience teaches, after all, that they by no means repel, but
much rather attract each other, and rejoice in mutual understanding. If the man is not
without feeling, then neither is the woman without understanding. If He cannot do without
the flower of religion, although he is a knower, then neither can She do without the fruit
of knowledge, although she is a believer. In the beautiful union of man and woman real human
life thus has its complete expression; here is the right equilibrium between feeling and
understanding, revelation and myth. If you ask me whether this inner equilibrium could not
also be expressed in a single person's life, then you now see for yourself what lies in this
question. To exist at once both as knower and as believer is, in other words, to exist in
the man's form of consciousness as well as in the woman's, to be oneself both man and
woman (androgyne). Incline
% --- p. 467 ---
\opage{467} now your ear to my wisdom, and let your mind be still and quiet, that you may
rightly mark my speech, for `wisdom entereth not into a malicious soul, nor dwelleth in the
body that is subject unto sin' (Viisd.\ 1. 4). You easily see that the androgynous existence
of which we speak must naturally not be confused with what is otherwise called a bachelor's
existence, be it with or without asceticism. The male-female existence whose concept we are
here to treat requires, namely, that the solitary person essentially express the whole,
complete human being, so that he, although a man in spirit, has a woman's soul, and,
although a knower, is nonetheless at the same time a believer. Such an individual must then
be not merely a real but also a quite extraordinary person, a human individual of a higher
order. As such human individuals of a higher order mysticism names for us two: \emph{Adam and
Christ.} In order to explain the myth of paradise, and make it really evident how it
actually came about that the first man could be the center of all creation, and unite the
opposite forces of the universe in his unity, so that the fall of the one must necessarily
draw after it the fall of the whole natural world, mystical theologians have represented
Adam as the original image of God, the complete human individual raised above the
difference of the sexes, i.e.\ as androgyne. With the fall of the primal man all the
oppositions of finitude were loosed; the elements stirred, the monstrous world-animal came
apart, and fell from itself into races and species and perishable individuals, when Adam
fell from himself --- into man and woman. As the fall is, so too must the restoration be. If
the frailty of the first Adam consisted in his getting an Eve, becoming a husband, eating of
the tree of knowledge, and becoming a knower, then the perfection of the second Adam,
mystically speaking, must of course show itself in this, that he, although the true, genuine
man, nevertheless does not become a married
% --- p. 468 ---
\opage{468} man; although the knower, nevertheless regenerates faith, and exists as the
restored unity of the forces of nature, i.e.\ as androgyne. To a mythical consciousness all
this now looks very profound, and may well be passed off as a revelation. A woman at least
ought to believe it; not merely because the natural consequence of the fall fell as a pain
upon her, but chiefly because she ought always preferably to believe that the man can do
without her, while she nonetheless feels within herself that her desire shall be to the man.
But if a manly understanding is also to believe it all: then you surely see that the matter
must have its great difficulties. Look about you in the real world; where, pray, do you find
the human being complete in himself, the extraordinary human individual who can endure being
at once both a knower and yet a believer? It would have to be a truly venerable figure, a
priest of the Church; or can you name anyone else? The ordinary faculty theologian you are
not to name. The theologian is a man, indeed a man of science, but no apostle, still less a
woman. `The scientific theologian is much knowing, but little believing.' A true theologian
really does not have much time to occupy himself with the actual matters of faith; he must,
after all, study `hermeneutics as a science', and above all `the critical introduction', and
most of all at a critical juncture, when myth overshadows revelation, and the faith of the
Gospel lies on its sickbed. --- You grow uneasy. Do not be hasty, do not become angry
(wisdom does not enter into a malicious mind); but consider our venerable figure once more.
Do look properly at the ecclesiastical Reverence of orthodoxy, as he enters company in the
dress of his order, grotesque as a man, modest as a maiden. Do you not now find in this
figure, in dress and bearing, in look and mien, in word and voice, something of that
wondrous, mystical doubling which points to such a
% --- p. 469 ---
\opage{469} being's having, as it were, a double soul? Does it not seem to you that this
reverend representative of piety moves precisely in the equilibrium of the male and the
female, that among all real people he most resembles an androgyne? --- There is hardly any
feature in the history of the Middle Ages that proves the religious earnestness of that time
more forcefully than the strict prohibition by which the Pope annihilated the marriage of
priests. It was a tyrannical command and contrary to nature, but from the Church's standpoint
entirely consistent. Priestly celibacy is an indispensable condition if revelation is to
succeed in overcoming myth. The witness of faith whose citizenship is in heaven cannot set
up house and become the father of a family on earth. The cleric lives only to spread the
kingdom of God in the world; but to spread the kingdom of God in the world, he must overcome
the world; to overcome the world, he must strive with the world; to strive rightly with the
world, he must know the world; to know the world rightly, he must be a knower. Not even the
shrewdest, most experienced man of the world shall be able to surpass the true priest in
thorough knowledge of men and many-sided experience. The true priest must know everything,
and yet be in the world as though he knew nothing. He knows desire, but he may not desire;
he understands better than anyone else what enjoyment is, but he may not enjoy. Everywhere
he must find himself touched by the glories of this world; but he himself may touch nothing.
On these terms our priest then becomes a knower; but you easily grasp that there is a world
of difference between priestly and profane knowledge. Profane, selfish knowledge rests in
the affirmation of the world, and leads to unbelief. Priestly knowledge is built on
self-denial, and finds rest only in faith. Only as the cleric in every way kills the
principle of nature in himself can he succeed in transforming the present into a myth; and,
as his consciousness unceasingly
% --- p. 470 ---
\opage{470} transforms the world of reality into a mythical realm, it at the same time
transforms the myth of religion into a real revelation. By this transformation the priest
becomes knowing as a man and believing as a woman, gets the man's spirit and the woman's
soul, and raises himself above them both. Into the hand of such a man heaven can entrust its
loosing and binding power; for he knows how to use it. See him, then, amid the graves, a
witness of the resurrection! whether he consecrates the dead to the peace of the kingdom of
heaven, or first condemns them to the torments of purgatory: one must, after all, believe
him, for he is a believer; deeply one must bow before him, for he is a knower. He stands at
the altar, and weds man and woman to indissoluble union, but, himself consecrated to the
Queen of Heaven, he cannot be wed to any earthly woman: for so it behooves, after all, him
who blesses to be greater than those who are blessed. --- What do you now think of the ideal
that will rise on this basis? Does it not resemble the idea of revelation to a T: so
supernatural, that is to say, so unnatural? That things in the Church went supernaturally
was the firm dogma of papal faith; but that just for that reason things in the world went
only all too unnaturally, the Middle Ages had at last to confess. So the miracle stopped at
last with the Reformation. It was a new fall. As the whole world fell with Adam, so now the
one saving Church fell with the priest. --- The ecclesiastical universe, which had hitherto
subdued the warring elements by concentrating its forces in high-priestly unity, came apart
and fell from itself into faith-races and species and perishable individuals, when the first
priest of the Reformation gave up his androgynous existence
and fell from himself --- into man and woman\footnote{%
``Whether the modern consciousness, not merely in its judgment of the theology of the
School, but especially on the delicate point of the marriage of priests, should betray any
striking one-sidedness, the many-sided reader will, of course, best be able to tell himself.
The moral standpoint of the Middle Ages lies, as is known, a degree lower than the modern.
According to the ascetic consciousness of the Middle Ages, the unmarried state of priests was
accordingly demanded as the most reliable guarantee of the dominion of the spirit over the
flesh. The humane consciousness of the present, on the other hand, sees in marriage a higher
expression of the fulfillment of duty. The Middle Ages, mistrustful of the natural drive,
proceeded from the naive presupposition: `it is easier for a man to marry than not to marry;
\textit{ergo}: the priest may not marry'. The modern consciousness proceeds from the reverse
presupposition: `it is easier to skip marriage than to enter into it; \textit{ergo}: the
priest ought to fulfill his duty and enter into marriage, as a bold and useful citizen'. The
only difficulty that remains, when married life is presented as the indispensable demand of
the ideal, then lies merely in the certainly peculiar circumstance that Christ himself, as far
as is known, did not enter into any marriage. The sting of this problem the biblical
theologian, church historian, dogmatician and symbolician \textit{Dr.} Karl Hase has
certainly also felt, when (Leben Jesu § 71) he raises on behalf of the idea the moral
question of what in all the world can have kept Christ from marrying. One now easily sees
that for a medieval imagination it would be a horror to enter upon this inquiry, and to
picture the Lord Christ settled in Capernaum, paying taxes and dues, with wife and children.
But herein precisely the deep opposition between the Middle Ages and the present shows
itself. Without letting himself be dumbfounded by any narrow-minded objection,
\textit{Dr.} Hase (Streitschrift. S. 105) therefore also raises, in sound confidence as much
in Lutheran humanity as in the modern consciousness, the far-famed
christological-anthropological question: `Mann, wo ist dein Weib?' Yes, --- `Mann, wo ist
dein Weib?' If this christological question were to be properly exhausted and treated,
according to the strictest demands of biblical theology as well as of church history,
dogmatics and symbolics, then the inquiry would surely be of enormous scope, indeed so
comprehensive that a theologically aspiring head would without doubt have enough in it for
his whole scientific career. (Such a work might then perhaps best be divided into the
following 3 main parts: 1st part for the licentiate degree, 2nd part for the doctorate, 3rd
part, a treatise in the Society of Sciences.) --- Should anyone, however, find himself
repelled and disquieted by the view here set forth, he ought surely to consider that the
modern consciousness in this as in all other points is and will be consistent, and must then
surely also have the right to express itself in all its consequences, without regard to the
respective author's own subjective opinion.''
\par\hfill The World-Consciousness's private note.%
}. The monks raged, the Pope
% --- p. 471 ---
\opage{471} thundered, the Protestants rejoiced; but the World-Spirit smiled, for ---
everything must happen naturally.''

% --- p. 472 ---
\opage{472} ``The surest guarantee that religion, although it has no real truth, will yet
never cease to work by the semblance of truth, is the circumstance that its general main
ideas have already managed to penetrate into the popular consciousness and fix themselves in
certain definite, ineradicable phrases. Consider, then, what wondrous magic can dwell in so
seemingly insignificant a thing as a phrase is. When the powerful thought finds a living
expression, it falls in love with it, grows together with it, and settles fast in the
phrase. The powerful thought is then, in its fullness, idea. As the idea bursts forth with
moving
% --- p. 473 ---
\opage{473} force, the phrase sounds as principle; but when the movement is ended and the
result comes, the principle sounds as phrase. Thus the word --- `freedom' is at one time a
mighty principle that moves everything, at another time an empty and powerless phrase. This
difference by no means has its ground in the idea itself; for the pure idea, the idea in and
for itself, is really nothing at all, or rather, it is --- a phrase. `Freedom', `equality',
`brotherly love', `the rights of man', `justice', nothing but ideas, nothing but phrases! The
difference is determined solely by imagination and passion. The phrase in which men at a
certain time concentrate all their passion, imagination transforms into a mighty deity; to
this deity all must sacrifice life and blood: that is the idea. But when passion has spent
its rage, and imagination is exhausted, the idol loses its gilding: now it turns out that the
idea for which one fought and bled was only --- a phrase. So long as the phrase is idea, one
cannot touch it without burning oneself; but when the idea has cooled
to a mere phrase, one can do with it whatever one likes. It goes with the phrase as with a
firearm. Once the pistol has been fired, a child can play with it without coming to harm;
but a loaded pistol is a dangerous toy. So with the phrase. Once it has discharged all its
electricity, one can use it for a pleasant pastime; if, on the other hand, it is charged with
the whole electrifying passion of an age, one must take care, and handle it with caution. It
goes with the phrase as with the fair-wind ship in the old myth. When Skibladner was not in
use, it shrank so small and fine that the god Frey could roll it up and carry it in his
pocket; but when it was drawn forth and swelled before the favoring winds, it became again
so roomy, large and strong that it bore all the gods of Asgaard
% --- p. 474 ---
\opage{474} through the air. So too with the phrase. So long as it rests in itself, in
language, it is only a small, insignificant word; one carries it with one everywhere without
even noticing it. But when the winds of the age stir, and breathe upon the phrase, and puff
it up into pathos; see, then the idea flaunts all the colors of the rainbow, enchants the
multitude, and lifts a whole people from the earth upward --- into the air. Inquire of
history, look at the past, consider the present, and you will learn to understand how it
comes about that the world, without contradicting itself, can at once despise and deify a
phrase. This holds not only of the many ideas and phrases by which real life must most
immediately direct itself; but it follows from the nature of the case that it must above all
hold of the super-real, and then in particular of the ideas and phrases of religion. You find
in world history not a single people that does not in its way assume a governing Providence;
but you nonetheless find in reality not a single people that really believes in a
Providence. This, however, you must not understand to mean that there are naturally always in
every people some who believe, others who do not believe; this catechism explanation you may
gladly leave to the priests. But the matter is essentially to be understood thus: every
people must, with regard to the governance of Providence, necessarily be both believing and
not believing. Herein is no contradiction, for faith and unbelief conform to circumstances.
In extraordinary cases every people is always believing; but by the rule of reality always
unbelieving. With this it agrees that governing Providence is both an idea and yet a phrase.
For extraordinary circumstances the phrase is a powerful idea (poets and priests vie in
blowing it up into pathos); in the ordinary relations of everyday life the idea is again
phrase, and reality
% --- p. 475 ---
\opage{475} goes its way. Listen to colloquial speech in the various circles of society;
everywhere you hear of `Our Lord', of `the Governance', of `the will of Heaven'. The words
`fate' and `providence', `dispensation' and `chance', mingle kaleidoscopically with one
another. The speakers surely think just as much by `the Governance' as by `chance', and just
as much by `the will of Heaven' as by a
`God bless you' and by the greeting: `Adieu'.
Meanwhile one has, once and for all, all these phrases, and that is the most important thing.
If one then feels some need of Heaven's assistance, this assistance can easily be procured.
The Governance is not far off; indeed it can even with good reason be said to be
omnipresent, since it is, after all, always present in the phrase. The religious popular
consciousness as a rule needs to be reminded at least once a year of divine providence. Every
New Year's Day there is, for a couple of hours during divine service, properly a Governance
over the whole land. It happens naturally; one gathers for serious reflections, audits the
budget of dispensations for the year gone by; looks with an uncertain glance at the uncertain
future; then one is seized by presentiments, then one falls into thought, and behold, then
one has the Governance. This, however, is not yet the highest Governance, or rather, it is
not yet the Governance itself in its highest and most powerful manifestation. That occurs
only when violent encroachments, wild passions, restless spirits have put everything at stake
and brought the common welfare into danger. But during such a crisis the Governance then also
comes over the people, and in dead earnest. God's providence naturally cannot concern itself
with every trifle. But where a whole nation is overwhelmed, and a whole people is near to
perishing, there the matter is of importance, there the Governance accordingly gets something
to look after. That the people's cause is always the just cause goes
% --- p. 476 ---
\opage{476} without saying; that Heaven always helps the just cause likewise goes without
saying. If now these two presuppositions go without saying, it undeniably follows that Heaven
both will and must and ought to help the wronged people to its right. If the multitude would
now warm itself at Heaven's justice, if patriotism would all at once be God-fearing, if the
spokesmen of enthusiasm would pledge the Governance as security for enthusiasm's victory;
then you must not exactly push yourself forward and say all too loudly that the governance of
Providence is only a phrase: at such a moment it might go ill with you. But only give it
time, and you shall see that the matter soon takes another turn. The holy law is inhumanly
strict; it demands of society not merely a certain general righteousness, what one might so
to speak call a righteousness \textit{en gros}; but it goes into the particular, and demands
of every individual a quite special righteousness, a righteousness \textit{en detail.}
To satisfy this demand is extremely troublesome. Good Lord! one is, after all, in this world
once and for all; and in this world every man of capacity who would serve the public always
has a dram of personal egoism, or at least his bit of ambition, which he simply must have
got invested along with it in the social catastrophe. If one earnestly wills the real end,
one may not reckon too closely with the means. That justice ought certainly on the whole to
be done, of that all right-minded people are convinced; but that one nonetheless ought not to
go too much into detail, and above all guard against the all too religious consideration
weakening the power to act, on that sensible people will, of course, also soon be in tacit
understanding. The all too lively idea of God's inflexible governance can, as you see, easily
become burdensome under such circumstances. Only for a few days will sovereign
power be able, as by a religious dictatorship, to be laid in the hand of Providence: little
by little it must
% --- p. 477 ---
\opage{477} go back again to the people, where, after all, it really belongs as well. This
gradual transition, conditioned by the gradual sinking of the religious mood, you find
exactly indicated in the use of the phrase, like the degree of heat on a thermometer. First
it is said: `the whole rests in the hand of the Governance' (--- boiling point ---); then it
is said: `one can surely leave something to the Governance, but one is nonetheless also
called upon to do a good deal on one's own' (--- several degrees below boiling point ---);
finally comes the witty ambiguity: `help yourself, and God will help you' (--- very near
freezing point ---). Here, then, you can again quite calmly touch the phrase without danger
of burning yourself on the idea. How wisely everything in this world is arranged, after all!
Yes, is it not admirable? There is, then, after all really a Providence, although in reality
there can be no God and no Providence. The Governance is enthroned in heaven, is
omnipresent, works everywhere, and extends its scepter over states and nations, over races
and individuals. No man and no people can prevail against Providence: the Governance governs
the world. But --- what a wonder! --- the world in turn governs the Governance. Men, after
all, get it with God and his providence exactly as they themselves will. If it is needed,
then Providence is --- almighty; if it is a nuisance, then Providence is --- half-mighty; if
it must be, then Providence is --- powerless. By virtue of Providence's almightiness the free
wills sink down into the dependence common to all; by virtue of Providence's powerlessness
they rise with renewed force to the highest independence. It is the spirit of humanity that
kneels before the almighty Providence, and worships --- itself. It is the spirit of humanity
that abolishes the powerless Providence, and rises up to cultivate --- itself. See, that is
the dialectic of freedom in pure self-consciousness \dots\ You are taken aback, friend! The
whole seems to you a miserable riddle, a desperate
% --- p. 478 ---
\opage{478} contradiction; see for yourself, think for yourself, test for yourself: you will
at last acknowledge that everything does happen naturally.''

``With the continued, humanizing activity of religion, and especially of Christianity, there
is, however, a still deeper connection. This humanizing activity is grounded
in the relation between the common spirit and the personal spirit. But in order that you may
rightly comprehend this relation, a fleeting comparison between the imagining Middle Ages and
the reflecting present will be of importance to you. In the Middle Ages the personal spirit
was far stronger than the common spirit. Rooted in the obscurity of feeling, shaped in the
medium of imagination, the idea of personality needed only to put forth a single bold shoot
of thought in order to reach up from the earth and unfold its crown in heaven. As the bearer
of the idea, individuality was at once loosened and inwardly freed from the all-dominating
nexus of reality. If the idea of religion broke through in an individual, the individual
broke with the world. Even such individuals as, without precisely taking the way straight
toward heaven, kept to finitude and worked for earthly ends, were in a way closed to the
many-sided influence of the common spirit. Every estate, every corporation, every pronounced
character bore, so to speak, its world within itself. Hence the abrupt contrast in light and
shadow, in groups, in physiognomies, that makes the Middle Ages so poetically attractive.
These strongly marked individualities could be held together only by a supramundane bond of
union. The form of the common spirit had necessarily to be theocratic-ecclesiastical; for only
the ecclesiastical common spirit is a spirit of personality. --- What a counter-image do you
not now find in the modern age? Rising civilization grinds off the sharp edges of
individuality, unlocks the holy of holies of inwardness, and polishes every man's
% --- p. 479 ---
\opage{479} soul into a pure and clear mirror of reflection for the whole surrounding world.
The Middle Ages kept to the magical play of colors of individualities; the changing
multiplicity of the images of imagination was like the tales in the `Thousand and One
Nights'. The modern consciousness levels the individual; the whole present wears uniform: the
industrial multiplicity of reality is put on display as in a bazaar. Here is the basic
condition for the freer exchange of ideas, for the equal distribution of spiritual goods.
In the past, indeed almost right down to the most recent time, men had now and then to feel
themselves oppressed by individual eminent geniuses, proud individualities who by their
personal superiority always kept the multitude at a certain humiliating distance. It was an
aristocracy of spirit that hampered freedom. Instead of ruling equally according to the
concept, and distributing its gifts to all, the idea was so capricious that it every moment
showed respect of persons, and as by a kind of divine nepotism sought only to puff up its own
favorites, and favored these few in a quite extraordinary degree. When such a favored one
had then, for a shorter or longer time, like a usurer, sucked the interest and the compound
interest out of the great common capital of the human spirit, and at last enriched himself
with so enormous revenues that he was no longer able to hold them within himself, he lapsed
into genial productivity, and now threw his sparkling ideas out among the people, as a
half-reckless rich man throws gold coins out of a window. That was beautiful, of course,
indeed it was brilliant; but you surely see that it was unjust. The multitude received as a
gift of grace what was due to all as a right. What a fuss was not made when such a product
of genius first saw the light! People gathered and read aloud; or they sought solitude in
order to immerse themselves in the richness of the content. In short, people regarded the
newborn masterpiece with an almost religious
% --- p. 480 ---
\opage{480} reverence, and felt an almost equally childlike joy as the innocents who on the
King's birthday were allowed to go to the square and see the golden apples dance. --- Only
with rising civilization does the idea become just. Instead of following its own whims, and
distributing its gifts at random, the idea is now placed under the control of publicity, and
must resign itself to doing what is right and fair. Instead of choosing a few to be
geniuses, the idea is then forced to make
society itself a genius, so that no one shall have anything to complain of. Instead of
plunging from the lofty peak of personality down over the breast of the multitude, as when a
roaring mountain torrent plunges down over the low plains and causes a flood, the idea must
now nicely distribute itself into many small running waters, must seep through the mass as
through marshy ground, so that the herb of equality may really thrive, and the humming
society grow green far and wide, like a fertile meadow. The highly objectionable inequality
between those who give and those who enjoy then disappears little by little in the world of
spirit. Everyone who has a little to give to the common good must hurry to come forward with
it before it is too late. `But he who has nothing at all to give?' Well, he has
\textit{eo ipso} the right to enjoy.
You see, since there are many who give, even if the individual contributions are only small,
the recipient need not make compliments at all for any extraordinary gift; he can with
unimpaired self-esteem quite calmly make a grab into the public common fund of spiritual
contributions, and, if he is not all too clumsy, even grab so much that he himself becomes a
giver. This is the equity of humanity in the true common spirit. This impartial spirit makes
no distinction whatever between individual and individual. On the contrary! it goes, after
all, with the common spirit in individuals as with the liquid in communicating tubes: it
stands at exactly the same height in them all. The difference will at most be only as a
drop, and even this drop is, after all,
% --- p. 481 ---
\opage{481} at every moment on the point of dripping away. If the modern spirit, according to its
concept, is, as you see, an altogether impersonal spirit, then you will now also understand
from this that the prevailing character of the present age must necessarily be both
unchurchly and irreligious.''

``But even if, relying on the above-mentioned, irrefutable presuppositions, one can now assert
with certainty that the Church will never again regain its medieval dominion, that the period
of Christianity's splendor is in that respect over and will never return, you will
nevertheless easily see from the foregoing that there still remains room enough for `the
Crucified.' --- You do know, surely, the meaning of the old saying, `to crawl to the Cross'?
It is a very telling expression. When someone, for example, has made himself rebellious but
meets his master, and is forced to beg pardon, then it is said, you know, that he must crawl
to the Cross. But who is it now who, in spite of the pure self- and universal consciousness,
nevertheless constantly sees himself exposed to this humiliation? Of course, it is not the
states, not the nations, still less the public or public opinion, that crawl to the Cross;
but it is? --- `The individuals'! Quite right. Where the impersonal universal spirit rises and
rules with might, the personal livelihood of the many single ones, of the individuals, will
always be hard pressed. Hence, even in the most enlightened age, it can never fail that there
must always be found a great many individuals who sigh to heaven and --- `crawl to the
Cross.' Note well that the strong self-feeling of the present age, which alone is able to
bear the ideas of `freedom' and `equality' under the whole enormous weight of modern
culture, --- that this colossal self-feeling was really first able to rise on the basis of
the Christianity of the Middle Ages. It comes about naturally. Classical antiquity made
individuals the serfs of the state; in the immediate natural
% --- p. 482 ---
\opage{482} state of the Nordic peoples the force of individuality was so exuberant that the idea
of the state hardly came into view. Before a proper legal relation could form, and a deeper
reciprocity between the individuals and the state be founded, the national elements had to be
mixed, and a higher power had to reveal itself, to sanction the inviolability of the
individual. It happened --- through the migration of the peoples and Christianity. In order
to come into their right, original right, the individuals, the victors as well as the
vanquished, first had to `crawl to the Cross.' Faith in the Cross crowned the Pope, and
celebrated its triumph in theocratic despotism, that is true; but beneath all this it sowed
the seed of freedom. The Reformation freed consciences; the Revolution proclaimed the `rights
of man': the latter would have been impossible without the former. Faith in the Cross and
self-consciousness are related as myth to reality. First religious faith must objectify the
idea, and cast the liberated self up into heaven; then critical thought must comprehend the
idea, and draw the acting self down to earth; and behold, there he stands, the strong,
self-conscious man, able to fight the fight of freedom and to assert his rights. The ancient
world, however, contented itself with one unconditioned end: either the individual had to
become a means for the state, or the state a means for the individual. To the modern world,
on the other hand, it is reserved to combine two heterogeneous but equally absolute ends: the
idea of the state and the inviolable right of the individual. That the movements which arise
from this must necessarily run directly against each other, that the loud prattle of the
present age about order and freedom in all directions is consequently a piece of nonsense,
you will readily grasp. The combined task is a pure contradiction; this contradiction can be
resolved only in the course of the world, i.e.\ in the ever-accelerating dialectic of
revolution and reaction. In revolution the individuals prove
% --- p. 483 ---
\opage{483} %
their natural right by overthrowing the state. The immediate consequence of revolution is
anarchy; in anarchy the superiority of the individuals culminates: the single ones are their
own masters. With reaction civil order returns again. The energy of reaction gathers itself
in dictatorship; and with dictatorship the state holds court-martial over the individuals.
Here, then, you see a fine example of the all-reconciling, all-mitigating power that lies in
time. The two sovereigns cannot, after all, reign at the same time; so they take turns: one
day anarchy, the next day dictatorship. Equality lies in the alternation: the faster this
alternation proceeds, the fuller is freedom. If it is asked which of the parties can hold out
longest in this conflict, the state or the individuals, then you must expressly note that
individuals and individuals are two different things. Individuals in general hold out well
enough: they are at bottom far stronger than the state; but individuals in their singleness,
these actual, determinate individuals, well, they hold out only poorly: in the long run they
are far weaker than the state. --- Now what these oppressed individuals, who after all also
have the absolute self-feeling together with the inviolable rights, without being able to
obtain their right, --- what these many unfortunates are then to resort to? On this point the
personal spirit will never cease to react against the universal spirit. Here there are
several ways out; which of these ways out the single one will choose depends chiefly on what
is called his natural disposition. Some (and these are the noblest natures) then cool their
resentment in the cold resignation of an ingrained defiance; the mob lives on in dim
despondency, with cowed mind; but there are also many who in such adversity bow beneath
Providence, sigh to heaven, and learn from the priest to crawl to the Cross. --- If the
social conflicts, considered by themselves alone, will at all times cast out so
% --- p. 484 ---
\opage{484} considerable a number of wronged and ruined individuals that the Crucified will be able
to gather from this host alone a not inconsiderable congregation, then the religious domain
expands still far more when we take family life into account and look about us in private
relations. Here, in the still, deep grounds of the heart, the strife of the individuals is
indeed not so visible and noisy as out on the wide world-ocean of history; but precisely for
that reason the pain is all the more inward, and the need for the aid of religion all the
more keenly felt. The Thorn-crowned One was right when he said to his judge: my kingdom is not
of this world; that is to say, he was in a way right, for his words must be understood
\textit{cum grano salis}. Faith in Christ really has nothing to do with the kingdom of the
world. Since it has no reality itself, it cannot possibly hold its ground against the powers
of reality; hence the Church too is now, with the rising universal spirit, pushed back ever
more and more by the state. The universal consciousness demands religious freedom: true
religious freedom is, of course, freedom from all religion. But religion shall not, all the
same, be homeless. Banished from the state, the Crucified will find a refuge in the house of
the private man, and raise his altar at the family hearth. Bearing his cross, he will then go
about from door to door: whoever needs him opens the door; but where there is no need, the
cross-bearer is not let in. As the need is, so must the help be too: the religious need gives
birth to its own helper. In this sense Christianity, however much it may preach and protest,
cannot now conceal that it is nevertheless a product of this world. It comes about
naturally. Had the suffering human heart not needed a gospel of resurrection, then the
Crucified would never have risen from the dead either; that is, in other words: then the myth
of the resurrection would never have been believed in the world. --- It is repeated often
enough, after all, that man's
% --- p. 485 ---
\opage{485} %
being is a contradiction, as the poet says so pathetically: `Was ist der Mensch? Halb Thier,
halb Engel.' Now each time `the half angel' cannot find its footing in temporality, `the half
animal' must of course set to grappling with eternity. If the individual isolates himself, and
wants to live alone, then melancholy comes and sits down beside him, and keeps him company.
All that melancholy brings along for private entertainment consists solely of a dark glass to
look through. So melancholy sits with the dark glass and looks, and looks at the solitary one;
so the solitary one sits with melancholy, and looks, and looks --- through the dark glass ---
into eternity. If the individual remains in the natural relations, and attaches himself
closely to `his dear ones,' so inwardly close that he cannot possibly tear himself loose,
because he feels within himself that he cannot live without them --- what then? Then death
comes and dictates separation. With the dead man there is no trouble, he got the better part:
`he is, after all, well provided for.' However inconsolable he may have been before he went
away, when death came, then he was consoled. But the mourners, whose hearts bleed, the
survivors, who now cannot live and yet must live --- yes, they need consolation. When the
father of the house dies, the family must crawl to the Cross; that is an old story. That is
why, you see, in all houses of mourning across the land you also see the man of the black
gown before the black coffin. There really is no sermon that works so powerfully as a good
funeral sermon. Consider now how many thousands of funeral sermons are preached, even in a
small country, in the course of a year, and you will get an impression of what religion is
for the suffering, of what the Crucified still daily accomplishes as `the Savior of the
World.' --- In the bridal chamber revelation is --- a myth; in the chamber of death the myth
is a revelation. So, you see, everything comes about naturally: people get through the
world.''

% --- p. 486 ---
\opage{486} %
``What do you say? The whole thing an illusion! --- Well, look at that, what wisdom can after
all dwell in a well-ordered human brain: everything, not this or that, but everything, the
whole, is an illusion. For the rest, this says very little, or infinitely much, all according
to how one takes it. It is an easy matter to declare the highest an illusion; but you see, my
friend, it depends on the depth and the reliability with which the recognition is possessed.
The softness of mind that cannot bear the pressure of existence, the excess of feeling that
will not put up with --- `unworthy limits,' the self-indulgent, pretentious recklessness of
the imagination that nevertheless cannot `carry the ideal through,' often enough wring,
especially from youthful hearts, the dreary confession that life is a game, virtue a
hypocrisy, duty an invention, and truth a miserable illusion. Against this sort of effusion
rational reality will always remain in the right: such pathetic exclamations cannot be
taken into account here at all. The task is to grasp in what the illusion really consists;
but, to grasp this, one must give an account to oneself of what the so-called absolute, the
unconditioned, is, when it is more closely determined. --- To imagine the absolute as a
personal God betrays, according to what we have said, an enormously stupid and clumsy
illusion. This illusion, which is to consciousness as a film is to the eye, has a numbing
effect among other things in that the will bows beneath a law from beyond, instead of having
the law for its doings and dealings in itself alone. Likewise, to attribute to the eternal, in
whatever form it may be, any positive validity whatever is the same as, put bluntly, to lead
one's own freedom by the nose and let oneself be caught in a most troublesome and to that
extent regrettable illusion. You will easily see, however, that the eternal cannot be
attributed unconditioned validity without such ideas as, e.g.: truth, justice, holiness, love
% --- p. 487 ---
\opage{487} etc., likewise having to lay claim to unconditioned validity. Now since these ideas
cannot possibly be realized in the present world, the human will, if it were to be
unconditionally bound and obligated to such unconditioned ideas, would \textit{eo ipso} be
obligated and bound to something supramundane and otherworldly, hence to something that does
not exist at all. No, instead of going about in mystical obscurity and weakening its power of
action through supraterrestrial considerations, the manly consciousness ought as soon as
possible to have done with all manner of illusory notions and recognize the absolute, the
eternal, the unconditioned, for what it is, an empty thought, an abstraction, a nothing.
This step, when it is, mind you, taken thoroughly, will instantly change the prospect and
break the spell. Man now not only becomes as God, knowing good and evil (that first advance
was still only a slight beginning), but the consummation is that man becomes wholly freed
from God and precisely thereby authorized to determine in his own person what is good and
evil. It would only be a new folly and a still more miserable illusion if we wanted here to
switch over into the much-favored figures of speech, and say that God is immanent in humanity,
and each man his own God. Such figures of speech are used only when one wants to ensnare
stupidity and flatter the imagination. In the pure, practical-prosaic style it runs quite
plainly: there is no God at all; the unconditioned, the absolute, is as little on earth as
in heaven. --- It is really a very ambiguous compliment to call humanity the immanent deity;
it is even more absurd than if someone were to assert of his coat that the coat in its
entirety was a genuine velvet coat, although he would have to admit that every thread and
fiber was quite simply woolen yarn. Humanity exists, after all, only in the multiplicity of
human individuals; but tell us, in the blessed name of the unconditioned: what is a
% --- p. 488 ---
\opage{488} %
human individual? When it is said of the bodily man: `from earth thou art come, to earth
shalt thou return,' must it not then hold of the psychical man: `from nothing thou art come,
to nothing shalt thou return'? --- Now take care that you do not misjudge the spirit! The
doctrine of the nothingness of the unconditioned is by no means propounded with the intention
that you should now instantly go and lay hands on yourself, in order to become what you
already are --- a nothing. You can always get to a suicide: it is not worth hurrying. The
doctrine of nothingness is propounded, on the contrary, with the intention that, rightly
understood and well digested, it should inspire you with heightened courage, renewed strength
and rejuvenated desire truly to live and to try life in every direction. Just look about you,
where you stand! You do not now stand, as the priests say so edifyingly, on the border
between two worlds, and so after all at odds with yourself whether you should prefer the
heavenly uncertainty to the earthly certainty; but you stand, as befits the free,
self-conscious man, before the countenance of the spirit --- right in the middle of pure
nothing. The absolute vanished, a place became vacant; this place you have yourself taken.
Not that you, dust and ashes, should regard yourself as the absolute; but the spirit sets you
in the place of the unconditioned: what a prospect for the freedom of your will!''

``Let us therefore, in this solemn hour of your initiation, repeat it, and say with calm, firm
conviction: the secret of existence, the marrow and kernel of human life, is, in the name of
modern knowledge, pure, sheer nothing. Your I, your will, your innermost self, be hereby then,
in holy, inviolable silence, released, redeemed and set free from all the bonds of religion
and all the prejudices of morality. Behold, truth and justice, holiness and love, judgment and
conscience, God and eternity now dissolve before your clarified sight into the pure, empty
nothing of the unconditioned. There is then,
% --- p. 489 ---
\opage{489} absolutely understood, no, no limit at all to the freedom of your will. Fidelity and
honor, promises and oaths, duty and right, the confidence of friends and the trust of men, all
serve you for your good, all you may posit as means, and nothing shall stop you on your way;
or, to let this rich truth play for you as in a pun: that which stops you is only that which
stops everything, insofar as in the end it makes all things equal, the beast and the man,
virtue and vice, the hero and the wretch, light and darkness, that is: the pure, absolute
nothing. --- Be still, do not run wild, do not misjudge the impulse of the spirit! The meaning
is surely not that you are to be initiated into vices and crimes. To become a highway robber
or a pickpocket one by no means needs the initiation of the spirit. But what you are
initiating yourself into is to live as a human being with human beings, and to try what you
can make of yourself when you really want to use your powers and trade with the talent
entrusted to you. So: in your outward relation, in the relation to your fellow men, you must
know well how to limit yourself and always bear in mind what the human signifies, and what
men are like. Out there in the throng you shall of course not scoff at virtue and conscience;
in men's ears you shall not be sparing in enjoining either the sanctity of the oath or the
strictness of duty or the inviolability of right; before the eyes of the crowd you shall
not, when it matters, be the last to bend the knee to the God of eternity either. In this
there is no falseness at all, no Jesuitism, no hypocrisy; for you know within yourself that
the holy powers of existence, though for you they are --- a pure nothing, are nevertheless in
reality the highest, the noblest and best that men know: it befits a man to honor the human.
--- Remember that you, whatever you do, always see to it that you have honor and the opinion
of the majority on your side everywhere and in every way. Honor is almost the only
% --- p. 490 ---
\opage{490} positive expression of the aim of freedom. Honor is the crown of life, the flower of
deeds, the garland of women and the idol of men; the world's honor is the world's all: glory
be to honor, the proud deity that holds the beautiful whole together.''

``The instruction is at an end; you now know the inner structure of the modern consciousness,
and that is enough; all the rest comes with practice. --- Now do not fall to musing, do not
stand here like a pillar of salt, but pull yourself together and use the moment while you
have it. With the whole speed of negative infinity steer boldly out onto the stream of time,
and you shall see that the winds of the world will surely bend to you. Just say: \textit{homo
sum, nihil humanum a me alienum puto}, that is: set about everything that is to come forward
in the present; dispatch a great deal in haste, so that you may, if possible, snatch the best
for yourself. Do not neglect yourself, but try what you will; not bashful, not timid, not
dejected! all ways, after all, stand open to you; become what you can, all that you will:
journalist, jurist, sophist, clubbist, priest, dean, minister in the state, bishop in the
Church, talent in science, genius in the folk school, what you can, and what you will! It is a
rich, a glorious course that now lies before you; hurry, then, run where you will! run ---
according to circumstances, but take honor with you! And when you have then shown the world
what you will and what you can, when enemies hate you, and adversaries fear you, and friends
envy you, and the crowd gapes at you, and the esteem of your fellow citizens proclaims your
capacity; then the modern consciousness shall smile contentedly, and bless you in secret.''

``Only one word more! Fear not the end, never ask about the last! As the first is, so must the
last be too:

\begin{center}
  {\small Alles, was entsteht,\\
  Ist werth, daß es zu Grunde geht.}
\end{center}

% --- p. 491 ---
\opage{491} %
\noindent That is the course of the world, that is the ring of wisdom, and herewith \textit{sat
sapienti}.''

\begin{center}
  \rule{0.38\textwidth}{0.4pt}
\end{center}

\begingroup\large
For God so loved the world, that he gave his Son, the only begotten, that whosoever believeth
in him should not perish, but have an eternal life (Joh.\ 3, 16).
\par\endgroup

As a fence, a living hedge, encloses a garden where the gardener plants and weeds, so let
these words be our bulwark and hedge, that they may guard the precincts of faith, and secure
its fruitful regions against wild incursions from the desert of illusions. --- So we will
then once again take counsel with the Gospel, and hear the testimony of the Believer.

``\emph{He hath blinded their eyes, and hardened their hearts, that they should not see with
their eyes, and understand with their heart, and be converted, and I should heal them. These
things said Esaias, when he saw his glory, and spake of him} (Joh.\ 12, 40--41).''

``The present age boasts of its pure light and its clear eyes. With these clear eyes man now
sees so precisely both the laws of things and the forces of nature, where they labor and work
in the least as in the greatest. With these clear eyes man sees precisely what is on the
earth, what is in the air, and what is in the heavens, what glints in the cloud, what teems
in the drop of water, what swarms in the mist of the ether. With these clear eyes man now sees
himself as the center of the forces, the masterpiece of nature, the lord of all creation: man
is spirit, spirit is man! --- But the spirit is bound in this nature; and what
% --- p. 492 ---
\opage{492} then is man according to nature? The highest, the lowest, the most excellent, the most
wretched, mighty and proud, a ruling god, and yet --- like the grass and the worm. What, then,
are we now to say of this pure light and these clear eyes? Come hither, you priests of
enlightenment, you who are so proud because you have seen and recognized that everything must
come about naturally. Tell us: what is the greatest thing in man, the height or the lowliness,
the frailty or the glory, when everything is to come about naturally? Can you really
comprehend how the world came to be without God? Then teach us, and do not keep the wisdom
for yourselves! Explain it, then, quite clearly, how it came about from the beginning that
spirit was born of nature, that the thoughtless gave itself thought, that the lifeless made
itself living, the blind made itself seeing, the deaf made itself hearing. --- `The thoughts
were in the forces, and the forces in matter.' --- Why, what light, what clear eyes! So the
forces themselves could think, then, although, as forces of matter, they were still blind,
thoughtless forces. So wisdom has sprung from thoughtlessly thinking forces, and this wisdom
of the forces has in turn taught you to think those --- thoughtless thoughts! --- Ah, what
kind of spirit is it, then, that can so beguile man that he now, with the freedom of thought
and the clarity of self-knowledge, will deny the God of freedom, fall down before blind
nature, and bow beneath the old fate? --- It cannot, surely, be of the heart alone, for how
should the human heart will its own bondage! nor can it be of thought alone, for thought
despairs when it must in earnest think such a contradiction. It must then be as the prophet
says: he hath blinded their eyes and hardened their hearts, that they should not see with
their eyes and understand with their heart.''

``And the present age boasts of its pure light and its clear eyes. With these clear eyes each
man can now see himself in the whole
% --- p. 493 ---
\opage{493} race, and the whole race in himself, can speak for all in his own name and for his own
in the name of all, and feel himself strong in the great whole. So the whole plumes itself on
its parts; so one counts the numbers of the peoples: no single one is to be forgotten,
neither he who has too little, that he may obtain a more --- for the good of the whole; nor he
who has too much, that he must be content with a less --- for the good of the whole; so one
rouses the lowly in their corners, so one judges the mighty on their thrones, so one opens
the books of law, that there may be dividing and sharing, and dealing and judging in all
justice. But --- kingdom rises against kingdom, nation against nation, brother against
brother. The stronger does violence to the weaker; the mighty scorns the lowly, and the lowly
rebels against the mighty; the poor man envies the rich, and the rich man misjudges the poor;
and still there is to be dividing and sharing in all justice, and everyone repaid according
to his works.''

\begingroup\large
Marvel not at this: for the hour is coming, in the which all that are in the graves shall hear
his voice; and they shall come forth, they that have done good, unto the resurrection of life;
and they that have done evil, unto the resurrection of judgment (Joh.\ 5, 28---29).
\par\endgroup

``And the present age boasts of its clear eyes. But tell us, then, you men with the pure
light: what is a world's judgment without a world's judge? What is a world's judgment when
those who are judged can hide in their graves?\footnote{Alas! terrible deeds have been done
under this heaven, because men blindly let themselves be led in the way of unrighteousness,
and terrible judgments of punishment have come upon the hardened blindness. But let us not
dwell only on what is foreign, on what is far away, but let us also here, before the face of
God, bring forth the thoughts that in these times follow us everywhere. When we see the
blinded hosts, whom we all call servants of unrighteousness, who now offer their arm to
rebellion and faithlessness, and rush forward against our land and people, how the thought
then comes to our mind: Is there, then, among the mighty who sent them out, among the lowly
who freely offer themselves, no one who puts to himself the question that ought after all to
be the first and nearest in all that we undertake: Is it also right, what you are doing? But
if we turn this question to others, then let us not forget that under all the trials of life,
in every temptation, the same question sounds to each one of us; it sounds from above,
downward, from him who judges the earth, it sounds from within, upward, from the depths of our
conscience, of our innermost consciousness: Is it also right, what you are doing? ---
\textit{Dr.} J.\ P.\ Mynster. \emph{Sermons Preached in the Year 1848}. p.\ 105.} What is a
world's
% --- p. 494 ---
\opage{494} %
judgment, if the righteous God will not judge men by a man, by the man whom he hath ordained
thereto, and hath given assurance thereof unto all, in that he hath raised him from the
dead?''

--- {\large ``Whom do men say that I, the Son of man, am?''} --- Our Lord and Savior had
already long walked with his twelve disciples, and waited for the deep presentiment that
stirred in them to rise into the clarity of recognition, when at last (it was in the region
of Caesarea Philippi) he asked them and said: ``whom do men say that I, the Son of man, am?''
{\large But they said: Some say that thou art John the Baptist; but others Elias; but others
Jeremias, or one of the prophets. He saith unto them: but ye, whom say ye that I am? Then
Simon Peter answered and said: Thou art the Christ, the Son of the living God}
(Evangel.\ Matth.\ 16, 13---16).

% --- p. 495 ---
\opage{495} %
``If we now collect our thoughts, and consider this utterance a little more closely, then it
is, after all, only a single, simple confession, by a single man, a lowly man among the
people, a fisherman from Bethsaida. And yet the Lord and Master knows how to be content with
this confession. He does not look questioningly at Peter, as if he first wanted to search
out in what sense such a word might have been spoken; nor does he receive the testimony
coldly and fleetingly, as one who only receives what is rightly due to him. But he rejoices,
and makes no secret of his joy. He who knew within himself that he would one day sit at the
right hand of power and call the dead up out of their graves --- see, nevertheless, how he is
content with little, how he praises the disciple, and rejoices over a single confession of a
single human being, so that not even he who received a world-kingdom with all its glory can
rejoice as he does. {\large Blessed art thou, Simon, son of Jonas; for flesh and blood hath
not revealed it unto thee, but my Father which is in heaven.} --- So Christ rejoices for this
single one's sake that a ray of the Father's light broke through his soul. And he who himself
knows what dwells in man and knows how frail the disciple is, he does not fear this single
one's weakness; but he builds on the confession. Yes, the Lord builds on Peter's confession,
not as a man may build an uncertain hope on another man, because he himself needs help; but
as the Lord who lays the burdens of a kingdom on a single servant, because he knows how
strong, how infinitely strong such a single one can be when the Father who is in heaven
transforms his weakness into strength. --- {\large And I say also unto thee, that thou art
Peter, and upon this rock I will build my congregation; and the gates of hell shall not
prevail against it.''}

When the ruler is anointed, and the crown is to be set upon his
% --- p. 496 ---
\opage{496} %
head, it is a grave, a solemn moment; it is as if the great multitude of the assembly
vanished away into silence, as if the many became one, and the one heard his heart beat. In
this solemn stillness everyone perceives that the crown is heavy; the deep silence testifies
that in the great multitude of spectators there is not one who will envy the mighty one, not
one who dares say to himself: `would that you were now in his place, who sits there before
the altar with the crown on his head and the scepter in his hand'; but each single one goes
into himself, and considers a human being's weakness and the greatness of the burden and the
weight of the responsibility. So much must indeed be endured, and much ventured, that a
temporal kingdom may take shape, and each in particular find his place and his room in the
whole, because in the whole itself there is a room of its own, an exalted place, where power
is enthroned. So the spirit of history now labors through the whole for all the single ones,
and through the single ones for the whole; and still the spirit cannot do full justice ---
cannot master its actual limits. Now it gathers power on a single head, and all complain
because there is only a single one who has the whole power, while the many single ones are so
powerless that the unity presses them to the ground. Now it divides power among all the
single ones, and all complain because each single one now has too much power, and the whole
is threatened with dissolution.''

``How entirely otherwise, then, it goes with the division of power of which the Gospel speaks,
as it shows us the Master's joy over the disciple's confession. See, the Lord and Master
begins to share with his disciple, he gives this single one power, a firm, divine power.
Strange! It is as if the Lord were now venturing everything in sharing his power; and yet he
knows within himself that he loses nothing at all by the sharing. --- Never before was
anything built so firmly on so frail a ground; never before did a single one
% --- p. 497 ---
\opage{497} on earth receive so divine a power; never before was a human being in the world raised
so high as this lowly man, who in the moment of confession surely did not think of exalting
himself, but only of honoring the Lord and Master to whom the power and the honor belong. It
goes wondrously, though, with this single human being: he is still as unsettled as one who
has not tested himself; he is still as wavering as a beginner can be. It might almost seem
as if the Master in his joy were overhasty in thus letting the disciple hear what was in time
to become of him; it might almost seem as if this single one, when raised too high out of
season, could only be in danger of falling all the deeper. And yet --- he is raised high; so
high a human thought cannot soar without growing dizzy. For the Master does not tarry, and
does not deliberate; but see, he enfolds his disciple in the embrace of the Spirit, and raises
him on his almighty word --- up into the heaven of heavens, where there is loosing and
binding.''

\begingroup\large
And I will give unto thee the keys of the kingdom of heaven: and whatsoever thou shalt bind on
earth shall be bound in the heavens; and whatsoever thou shalt loose on earth shall be loosed
in the heavens.
\par\endgroup

So much must indeed, humanly speaking, be ventured, that the single one may win a room, a
place, an abiding place, an infinite freedom in the heavens. But the Lord who ventures
everything knows best within himself how he is also to win. So much must indeed be endured,
that the single one may have a share in the power, the true spiritual power. But the Lord who
knows within himself that he can endure the disciple's weakness, overhaste and
misunderstanding, can begin at once from the foundation; for the confession is not of flesh
and blood, but of the Father in heaven, and the Master builds on the confession. The disciple
can at once have a share in the power: nothing is
% --- p. 498 ---
\opage{498} lost by the sharing, for by virtue of the confession the Lord himself still has the
whole power, and the Master builds on the confession.''

``When the present age wants to speak its modern language, it may well call itself an age of
representation. In order to have a share in power, the many single ones must let themselves
be represented. Thus there are many and many kinds of representatives for the various estates
and the various interests in a whole kingdom. If someone now wants to ask what right and
significance each single believer can then obtain in Christ's kingdom? very well, then we too
put up a representative, the fisherman from Bethsaida. Do not despise him, this
representative! he too submits his mandate for examination: see, he has the keys of the
kingdom of heaven, and whatsoever he binds on earth shall be bound in the heavens, and
whatsoever he looses on earth shall be loosed in the heavens. Do not fear him, this
representative! do not be afraid that such a single one will misuse his power to oppress the
many other single ones! He cannot misuse the power; for the least misuse would be its
forfeiture. He will not misuse the power; for before he received the extraordinary power, he
had to learn fully that he had no power whatever of himself, but that it was given him from
above by virtue of his confession. Do not envy him his precedence, this representative! He
has not come to exalt himself above his brothers, but to speak for the freedom of the single
ones in the kingdom of the spirit and of freedom. If anyone in the assembly wishes to have a
share in his power, then he is willing to hear him, and offers fair terms: `if you can confess
Christ, the Son of the living God, as I confess him, then he will surely give you too a key
to the kingdom of heaven. Come hither, then, and we shall share the power, and lose nothing by
the sharing; for our whole power is in our confession, and by virtue of the confession the
Lord himself has all power in heaven and on earth.' ''
% --- p. 499 ---
\opage{499} ``So strangely, then, do things go in this kingdom of Christ: that the spirit works in the
infinitely great, and yet condescends to the lesser and works in the small; that this spirit
prepares all things for what is to come, and yet will penetrate the present; that right is
to be established in the heavens, and yet there is to be loosing and binding on earth; that
the many judges of the world are at last to fall silent before the one Judge who sits at the
right hand of power, and yet each single one perceive the binding and loosing power in his
own breast: with all this it must surely go wondrously, but then it has also gone wondrously
that the world has received a Savior.''

Evangel.\ Mark.\ 10, 13---16.

\begingroup\large
And they brought little children to him, that he should touch them; but the disciples rebuked
those that brought them. But when Jesus saw it, he was angry, and said unto them: suffer the
little children to come unto me, and forbid them not; for of such is the kingdom of God.
Verily I say unto you: whosoever shall not receive the kingdom of God as a little child, he
shall in no wise enter therein. And he took them up in his arms, and laid his hands upon them,
and blessed them.
\par\endgroup

When in this picture we see the Savior of the world in the circle of the children, we feel
immediately what must have stirred in the pious mothers, that they wanted to bring the little
ones to the Lord. On the other hand, it might perhaps at first glance arouse some
astonishment that the disciples showed such a repelling conduct toward those who brought the
children. They surely did it with the best intentions; but
% --- p. 500 ---
\opage{500} their well-meant zeal was grounded in a misunderstanding: they wanted to free their
Master from an untimely importunity. ``What are these children to do with Jesus now? Had they
only been possessed, from whom the Lord could have bidden the unclean spirits to go out; but
these are nothing but healthy, fresh, well-formed children. Had they only been penitents,
drawing near with a broken heart; but these small innocent creatures as yet know of nothing
evil, and cannot grasp the word about the forgiveness of sins. What, then, is the point of
this intrusiveness?'' Thus the disciples thought. Their practical understanding told them
that it must surely be, if not unworthy, then at least most inconvenient for the Lord to
concern himself with these children. Without troubling themselves with any question to the
Master himself, they then at once set about proving their zeal in service by removing the
throng, keeping the way open and making room. The women, on the other hand, or whoever they
were who brought the children, took the matter in quite a different way. Without any notion
of being a bother to anyone, they only followed the prompting of the heart, and came with
their little ones to ask the Holy One for a blessing. --- So the self-wise disciples now
contend with the simple among the people; but the Savior ends the contention and lifts up his
voice: \emph{suffer the little children to come unto me, and forbid them not; for of such is
the kingdom of God.}

Much contention has been waged in the course of the ages about the moral disposition of human
nature and the relation of the world to the ideal of the good. That the race is fallen, that
every human being consequently participates from birth in the general corruption, is a
presupposition that cannot be abolished without at the same time abolishing the Gospel of
redemption. It is an evident misunderstanding if anyone interprets the Savior's words about
the little children as though these words were to prove to us that the childlike nature was
still wholly pure,
% --- p. 501 ---
\opage{501} sinless and undefiled. But it also betrays a misjudgment of the present scene if anyone
presents human nature to us as so fundamentally corrupt that sin, through Adam's fall, has now
become our proper substance, as if even the child in its innocence were an unconscious
criminal, an unblessed being, possessed by the spirit of darkness. The Savior's words about
the little children contain a testimony about the relation of the Creator Spirit to the
fallen race, a testimony that affords the right guidance when we are to speak rightly of the
conditions for the salvation of the world.

\emph{Verily I say unto you: whosoever shall not receive the kingdom of God as a little child,
he shall in no wise enter therein.} That is the condition of redemption: these are fair terms.
To receive the kingdom of God as a little child is, in other words, to receive it as a
fatherly gift. And yet the demand is in another respect at the same time both serious and
strict. --- To begin as a little child is, in other words, to begin over again. But the older
person who has already ``become something'' always has trouble beginning over again; for in
order to begin over again he must as it were at one stroke give up this something and resign
himself to its becoming nothing in the sense of eternity; this way back to the first
beginning can then be troublesome enough, especially for those who have come further.
Imagine a sage, a philosopher: he has already founded an original world-view, and established
a school. Disciples and adherents spread wisdom everywhere in his name and his name in wisdom:
must it not now come hard enough to such a one to have to bring his original world-view as a
sacrifice to the truth of the Gospel? must it not come heavy enough to the famous teacher to
close the school, and tell the many disciples that the whole of human wisdom is at bottom
vanity, that one must begin over again, that whoever will enter into the kingdom of God must
% --- p. 502 ---
\opage{502} receive the Savior as a little child? Imagine a theologian, a scribe, who has already
grown gray in hermeneutics, a teacher in Israel, who long ago became renowned in isagogics:
must it not come hard enough to literary renown to put the many half-finished investigations
on hold, to write the many interesting reinterpretations into the book of oblivion, to let
the many ingenious, highly remarkable hypotheses fall, and to begin over again to receive the
Gospel as a little child? --- Let us then consider more closely what lies in the demand: to
begin over again as a little child. If it is a true evangelical fundamental demand, then it
must surely be recognizable by this, that everyone who has truly tested its content, and
weighed it on the gold scale of conscience, must say to himself: ``so heavy and yet so
light!''

Become a child again, return to your original beginning, that is:

\begin{center}
\emph{1) Let your being be as God has created you, be\\
natural, be truthful.}
\end{center}

For the fulfillment of this demand the child is presented as a model. Not that it should be a
moral advantage in the child, as though we were all born with essential virtues; no, it is no
special advantage, but an original determination of nature that is in question. The child
cannot deny and cannot conceal the peculiarity of its nature: it cannot become untrue to
itself. In order to deny one's own nature and conjure oneself into a foreign being, one must,
after all, know the evil, cunning artifices of comparison and of imitation and of
dissimulation; but the child has not yet come so far. The child must begin by being as it is,
by being in all naturalness as our Lord has created it; that is why we also say of a little
child that it is --- one of God's creatures. --- True,
% --- p. 503 ---
\opage{503} it soon appears that there are conflicting forces in the child's nature; but these
conflicting forces express themselves without reserve: expression and being correspond to each
other. To be sure, a bright child will early try its hand at manifold imitations; but these
innocent imitations constitute its play: it comes, after all, so naturally to the child to
play. If the play begins to degenerate, if the imitation becomes a little too unnatural, if
the child wants to ``put on airs'' and assume a foreign being, then fatherly seriousness comes
and teaches it to be what it after all at bottom is, a truly natural and guileless child. To
be sure, a clever child will now and then try to help itself with a little untruth, and see
whether it might not also get by with a little concealment and dissimulation. But face to
face with fatherly authority the dissimulation must fall, and the secret come out, and the
sinner go nicely to confession. When the child has then told the plain truth, and received
its punishment, and asked for forgiveness, it can breathe healthily and lightly again, and
involuntarily feel well; for the father's chastisement has, after all, torn it out of the
power of the lie. Were it conceivable that a child could by nature be so ambiguous, so
deceitful, so initiated into all the sorry arts of aping and dissimulation, that at every
moment it deceived both itself and others, would it not then be a sheer impossibility to get
such a child brought up? --- So with man in his relation to God. Do not believe that the
Gospel will abate anything of its ideal demand. Face to face with the Creator, the Almighty,
the All-knowing, a ``grown-up person'' must not merely be something childlike or almost like a
child; no, the relation of childlikeness holds unconditionally: the one most fully of age is as
a child, a little child, face to face with the heavenly Father. This relation must surely be
easy, since it is so natural. And yet the demand is strict; for what the child is by nature
without its own cooperation, the developed human being is to be in spirit and in truth. To give
oneself
% --- p. 504 ---
\opage{504} over to the softness of half feelings, to go into second childhood like one who is old
and enfeebled, is not the ideal. But, full of vitality and health, to summon up all one's
strength in order to attain the weakness whose secret is strength; to subdue the selfish, so
as not in a self-deception to be too strong for God: that is in truth to satisfy the ideal.
One grows uneasy at heart when one sees a warped child. One thinks to oneself: Lord God, what
a sorrow it must be for the father who has such a child. Should the God of truth, the heavenly
Father, then be able to take pleasure in half-natural, untruthful and warped children! ---

Become a child again, return to your original beginning, that is:

\begin{center}
\emph{2) Bow beneath God's fatherly hand, remain true to the\\
ideal, be humble.}
\end{center}

For the fulfillment of this demand the child is the model (cf.\ Matth.\ 18, 3---4 Parall.). On
a preliminary view it might perhaps seem striking that so essential a determination of
inwardness as true humility is should be capable of being represented by a child. Even if one
imagines that the child whom Christ in the Gospel chose as an example for the disciples was
the best-behaved, most lovable and most undemanding of all children in all Galilee, one gets
no further in understanding by that. The task is to find the basic features of humility
prefiguratively designated in the childlike nature. The selfish germ of haughtiness and pride
is naturally already active in the child's soul. Self-will, involuntary self-complacency, the
desire to raise oneself above one's equals, express themselves early: a child can, in the
sense of innocence, be downright both presumptuous and despotic. One should beware of becoming
sentimental when one wants to grasp childlike innocence. A healthy
% --- p. 505 ---
\opage{505} child is too full of natural vigor and freshness of heart to fall under the sentimental. But
the Savior is right when he sets the child before us as a model. However strongly childish
self-will may express itself in particular moments, this is nonetheless decisive, that a
child must --- let itself be told: a little obstinate innocence is a great cross in the house. Yet a
word of authority, spoken at the right time, with the right emphasis, can also as in an instant
annihilate childish defiance; face to face with the father's gaze the child can even
feel itself to be a nothing, can bow beneath the admonition, can give itself over in unconditional trust.
But this, on the level of spirit, is the essential mark of humility. It is often said that
a child can take pleasure in little: this already is a good sign; for he who is himself content
with the lesser is not easily tempted to belittle the greater. However, this is
still not the essential thing; for the child can also take it into its head to covet the greater: a
child too can become envious. But the essential thing is that the childlike disposition, which so openly and
willingly receives every impression, is above all able to receive a deep and living
impression of the greatness of the divine. If you believe in the ideal, if you can speak ingenuously of the
exalted, the glorious, the blessed, then do not turn to the full-grown, those of age, who are
too clever for you; but turn with your speech to the little ones, those not of age, and you shall find a
grateful listener in the child. See, its eye is clear, its ear open; now take care that
you do not misspeak, for the little listener can notice it at once: himself captivated by your words,
he catches you in your word. --- The modern consciousness, which explains everything, of course also explains
to us that the child's nature is only a product of obscure natural forces. This explanation will not
suffice. It does not at all look as though the child, lately come forth from the darkness of the ground of nature,
had to spell its way painfully through the sensuous,
% --- p. 506 ---
\opage{506} in order at last to glimpse a shadow of the supersensuous; no, the little listener, transfigured
in ideal vision, looks rather as though he had lately come down here from the regions of light, and
now stood here, hovered about by blessed recollections, listening for --- the tones of home. A child
believes in the ideal, and bows beneath it. This is the condition of true humility: whoever
will only believe in the divine ideal, and measure himself by its measure, shall surely
learn what humility is.

Become a child again, return to your original beginning, that is:

\begin{center}
\emph{3) Bow beneath God's fatherly wisdom, obey\\
the holy word, be simple.}
\end{center}

For the fulfillment of this demand the child is a model. With a wonderful instinct
the childlike understanding learns to distinguish between good and evil. It is true, to be sure, that this
distinguishing does not exactly penetrate into the depths of ethics; but its originality shows itself in this, that
the opposition is kept so pure. Good and evil are related in the child's conception as light and
darkness. The secret of worldly wisdom, that the good is after all never really good, and the evil
not so entirely evil either, the childlike understanding cannot yet grasp. That is just how it is
with simplicity. There is much earnestness in a childlike soul; from the beginning the child sees
no disproportion between word and deed, intention and execution, promise and fulfillment. If a
child has not far to go from smile to tears, neither has it far to go from word to deed;
lengthy preparations, long-drawn deliberations are against its nature. If you have promised a
child something, then by no means forget it. If you forget it, you will be dunned. No
creditor, be he ever so great, can on this point have a
% --- p. 507 ---
\opage{507} surer memory than a little trusting child. Nor does it help if you begin
by trying to explain your promise away; the child will hear nothing of any interpretation: it
has as yet no taste for ``hermeneutics''. Such is simplicity. --- People talk now and
then a good deal about what are called their childhood recollections. One then learns that
most have had a very happy, some a fairly happy, a few an unhappy childhood.
Insofar as this mainly comes down to mere talk, the matter probably has no great
significance (the point is that every recollection must --- be recollected); on the other hand it should be expressly
emphasized that everyone who really recollects his childhood (be it otherwise happy
or unhappy) involuntarily sinks into melancholy. The turning point about which the
quiet melancholy of recollection moves is essentially to be determined as the transition from the simple to
the manifold, from the trusting security of ingenuousness to the painful emptiness
of offense. ``The child was so near to God, it was as if it could see his angels in heaven. But
then people always talked so much and so confidently about everything earthly, and only
so little and so unreliably about the heavenly: so the good angels vanished. The tempter came,
thoughts came, doubts came: God hid himself yonder --- behind the stars. And yet the child was so near
to the good; the holy law stood written in its heart. Back then! yes, back then deed was to
be as word, back then every promise given was to be kept, back then the
good intentions were to be carried out, then --- evil was to be stifled. But then people began to
interpret the law, and to come out with the many exceptions to the rule; then the
innocent one could well see in its elders that deed must surely be very different from
word, execution very different from intention, fulfillment very different from
promise; then the one not of age came
% --- p. 508 ---
\opage{508} little by little to an understanding with evil, then the simple one at last hit upon
interpretation, then innocence finally became worldly-wise, and alas! then came offense.''
--- \emph{Woe unto the world because of offenses! It is needful indeed that offenses must come;
but woe to that man by whom the offense cometh} (Matth.\ 18, 6). --- The Savior knows
our original nature. He knows that we all have from the Creator a need, an essential
need, for the kingdom of God; but he knows also that the world is offended at God, and gives offense
upon offense. Therefore he watches over the little ones, and takes up the cause of those not of age, whom the world
despises. \emph{Take heed that ye despise not one of these little ones; for I say unto you: their
angels in heaven do always behold the face of my Father which is in heaven.} --- A little child is only an
insignificance in the great world, and can so easily be forgotten in the throng; but for the child in
the Savior's arms he has given his word: it will not be forgotten in heaven.

\emph{Verily I say unto you: whosoever shall not receive the kingdom of God as a little child, he shall
not enter therein.} How, then, is this word understood and applied in Christendom?
Christ's first command, ``Suffer the little children to come unto me, and forbid them not'', has through
the introduction of infant baptism long since been complied with in all the lands where
Christianity has found entrance into the family and the state; does present-day Christendom, then,
with the same care comply also with the attached condition, to appropriate the Gospel of the kingdom of God
as a little child? --- To regard Christianity as a simple
children's doctrine is certainly fairly common; but with this the enlightened at the same time connect
the accessory notion that one cannot remain standing at the first simple thing, but
must gradually
% --- p. 509 ---
\opage{509} %
grow out of one's children's doctrine much as out of one's children's shoes, and look about more closely among the manifold
ideas of the age. The meaning is not exactly that an educated man is simply to give up
religion. The guiding fundamental thought can perhaps best be expressed thus: as long as the individual is still
of minor age, he stands, as regards his world-view as a whole, on about the same
level as the old prophets and apostles: unconfirmed youth should therefore preferably take both
catechism and textbook and Bible history as literally as possible. But once one has been
confirmed, then thinking for oneself begins; and through thinking for oneself, together with the culture of the age, one
gradually gets away from the simple, and soon reaches the level of majority from which one looks
back on evangelists and apostles with about the same feeling as on unconfirmed youth
still attending the pastor's classes. To be a simple person is, in the language of the present, about the
same as to be a narrow mind or, as it is called, one of the saints. People who
are otherwise not exactly bashful can turn quite red at the mere thought
that the world might one day consider them to belong to --- ``the saints''. This prejudice (for that it
really is a prejudice, rigorous thinking should surely be able to prove) is so deeply rooted in the present
that it might almost seem an impossibility to get it eradicated. All who occupy themselves with
speculating on the salvation of the world are agreed that it cannot possibly go on any longer with simple
Christianity and the ever-growing reflection of culture as hitherto. But what, then, is Christendom now to
do? Many and various are the proposals that recent times have put forward for the good of the Church; many and
various are the endeavors by which one seeks to prop up Christianity and keep the dying
Gospel alive. These proposals and these endeavors can, however, all
% --- p. 510 ---
\opage{510} be traced back to two points of departure: a theoretical-intellectual and a popular-practical.

The representatives of ecclesiastical intelligence must, as is natural, be sought chiefly among
the theologians of the School. The sober-minded theologians of the School are all convinced that an open break
with the intelligent world of culture would be not only an irreparable loss for science, namely
Christian science, but might even bring about the ruin of the Church. The task thus becomes to
set ``proof against proof, intelligence against intelligence''. --- Yes, if only that could help!
Still, let us try it; perhaps the world of culture is not so unreasonable after all, if one will only rightly
attend to what it is that it really respects.

``The cultivated world, which on the whole has so much esteem for historical-critical information, must
surely at last be won over to the hermeneutical and critical investigations which aim at
making the genuineness and credibility of the O. and N. T. really evident to sound reason. If
one could only get the synoptic tables to agree with the Gospel of John; but there is that
enormous difficulty with the Passover reckoning! It is indeed a real nuisance, that
Passover reckoning! If the cultivated world gets wind that the Passover reckoning will not work out,
the Gospel of John certainly cannot hold its ground.'' --- Then let us, in God's name, abandon
the Passover reckoning, and try to find a surer foothold.

``The world of culture has respect for the absolute Idea. On this presupposition the Church grounds its
apologetics, its symbolics, its dogmatics, its ethics. It must indeed be said that the true and essential
relation to God is a relation of existence, that religion is thus not merely a task of cognition
but just as much a task of feeling and of will. But it must also be said that
% --- p. 511 ---
\opage{511} %
the Idea of religion is the Idea of truth. Since the absolute Idea of truth must necessarily be determined
as the Idea of science, it goes without saying that the Church cannot do without its Christian science.
Science does not, to be sure, do everything; but it does clear many an eye, and exercises a quiet working.
The cultivated world will certainly never search out the ultimate grounds or in full earnest penetrate into
the depths of the Idea; but for all that a profound, clear and systematically contemplative presentation of
the great chief truths of Christianity will make its impression all the same. It would truly be of
great significance for the ecclesiastical general consciousness if one could now lay before the world of culture a good
and sound Christology.'' --- So much for science.

As for the most popular spokesmen of the folk-church endeavors, the language and the tone are
naturally far more popular. With due disdain both for the foreign newfangled culture
and especially for everything the German calls Wissen and Wissenschaft, they may well, in the name of wit
and witcraft, even feel like thumbing their noses at the ``Passover reckoning'' and ``Christology'' and ``the
absolute Idea'', and to that extent need not rack their brains either over the certainly very
difficult question what it is, after all, that ``the world of culture really has respect for''. Without
delaying over too many broodings they can straightway get the fundamental view, and go straight to the point, and
say it straight out, ``that the course of Christianity's life in the world must everywhere and in every way conform
to the course of the peoples' life in history; yes, say it straight out, that the whole Christian-ecclesiastical
muddle, under which the present suffers in particular, obviously has its ground in the different
peoples having mixed and got into a muddle with one another. It is as clear as day that if one would
become a Christian, one must first be a human being, but if one would be a human being, one must surely be able to
remember
% --- p. 512 ---
\opage{512} %
what one is called oneself, and to which people one belongs. If we would think of interpreting
Christianity, we must surely first think of interpreting the signs of the times and making clear what is in the people;
if we would think of calling the faith of childhood to life so that it is good for something, we must first and foremost
think of ending the battle of the giants and binding the Fenris wolf and cleansing the heart of the people of the foreign
mischief. There is now no doubt that on the whole it must come about thus; it simply cannot be
otherwise. It is so in the Church, and it is so in the world, it is so in the spirit as in nature: with the leafing of the woods in
springtime comes birdsong over the forest, with the folk-spirits at the turn of the year the Savior rises
from the dead.''

See, these are now the most important endeavors of Christendom in the new age. But what are we then to
judge, or rather: what may we suppose that the Head of Christendom would judge of these many
proposals and these many endeavors, if we could rightly appropriate his Gospel?

--- ``But is it not then every man's duty to do what he can in his station and with
his powers for the good of the whole?'' --- Yes, assuredly! It is every man's duty. --- ``Have
hermeneutics and the critical introduction not also their significance, then?'' --- Yes, \emph{their}
significance! --- ``Can it then be denied that the eternal ideas are the guiding stars of the human spirit, and
Christ himself the holy light that would shine through the realm of ideas?'' --- No, that can certainly not
be denied. --- ``Has it now perhaps come so far that no man can speak a living word to living
people before he has brooded away half his lifetime, or at least first asked a brooder's leave?'' ---
No, it has not yet come so far. We acknowledge all able and good endeavors, in whatever
direction they may work; we honor, we hold in high esteem the best men of the Church and of the people who in these
times of darkness still hold fast to mankind's
% --- p. 513 ---
\opage{513} %
highest concern. But it is another question whether the many Christian plans of the present now also
coincide in all respects with Christ's own plan; whether it should now really be found to
be either ``ecclesiastical science'' or ``the national awakenings'' through which the Savior of the
world has resolved in these days to save the world. Before we venture, however, to answer so
decisive a question, we will first try to strengthen our weak insight by considering a couple of
peculiar scenes in the sacred history.

Gospel of John, Ch.\ 3. V.\ 1---21.

\begingroup\large
There was a man of the Pharisees, named Nicodemus, a ruler of the Jews. He came to
Jesus by night, and said unto him: Master, we know that thou art a teacher come from God; for no man
can do these signs that thou doest, except God be with him.
\par\endgroup

The man who in this narrative approaches the Savior of the world is, to judge by human
reckoning, assuredly a man of influence and ``intelligence'', a religiously cultivated man, a man who must surely
also have done a good deal of theological study. (Whether Nicodemus had exactly studied ``the critical
introduction'' we cannot say with certainty; but that he must have studied his
``hermeneutics'' duly becomes more than probable to us when we consider, on the one hand, that he was a
scribe, a teacher in Israel, and on the other take note of his utterances in the following conversation.)
This man, then, came to Jesus by night. And why just by night? Doubtless because he
counted on the quiet working, and feared an open break with the other rulers; for he was
himself one of the rulers. Although it is, as said, somewhat late when he pays his visit
(already well into the night), he nevertheless finds the ``Master'' awake, and is let in, and has now surely
also prepared himself for a thoroughly religious-scientific conversation.
% --- p. 514 ---
\opage{514}%
That the stranger's greeting is both solemn and very appreciative can naturally not
surprise us; but the mind is roused as from a slumber when we hear the Savior's voice:

\begingroup\large
Verily, verily, I say unto thee: except a man be born anew, he cannot see the kingdom of God.
\par\endgroup

This answer comes so suddenly, yes, so surprisingly. Nicodemus is taken aback, and cannot rightly grasp
it: hermeneutics ponders, and has not yet fathomed it.

\begingroup\large
How can a man be born when he is old? Can he enter the second time into his mother's
womb, and be born?
\par\endgroup

We should truly do Nicodemus an injustice if, with certain interpreters,\footnote{Strauß and De Wette are agreed that Nicodemus is to understand the expression \emph{to be born anew} directly of
bodily birth. D. W. admits that ``eine solche, zumal hartnäckige Unwissenheit'' is
striking, indeed improbable ``an einem israelitischen Gelehrten'', and therefore he lays the blame on
``the Johannine type''. --- Yes, that is what comes of hermeneutics. First one misunderstands the turn of
speech, and accuses Nicodemus of an error he does not commit at all. Then one finds that the error
is too gross --- ``an einem Gelehrten'', and lays the blame on the evangelist John, in order to acquit
the Pharisee Nicodemus. Then one finds again that John after all cannot deliberately have meant to
insult Nicodemus, and now at last the whole blame falls, not on John himself, but on ---
``the Johannine type''. Good Lord! if it is to go on like that, then it is indeed no wonder
that the fourth Gospel cannot hold its ground.} we maintained that he understood the Savior's words
quite materially. The Pharisee is not so uncultivated. On the contrary! he is precisely
% --- p. 515 ---
\opage{515} developed enough to reason. He wants, as they say, to feel his way as to how matters stand with
the task: the objection that lies in the evasive question is exceedingly well suited to a
reasoning conversation. But Christ will not dispute with Nicodemus; instead he repeats
the condition, and impresses upon him the necessity of regeneration.

\begingroup\large
Verily I say unto thee: except a man be born of water and of the Spirit, he cannot enter into the kingdom of God.
That which is born of the flesh is flesh, and that which is born of the Spirit is spirit.
\par\endgroup

Nicodemus is taken aback, and cannot grasp it: hermeneutics ponders, and has not yet fathomed it.

\begingroup\large
The wind bloweth whither it will, and thou hearest the sound thereof, but thou knowest not whence it cometh,
or whither it goeth; so is every one that is born of the Spirit.
\par\endgroup

As a teacher in Israel Nicodemus must surely have known his passages of Scripture, that is to say,
the passages of Scripture in the Old Testament. He must then undoubtedly also have gone through such
passages as Psalm\ 51, 7---12. \emph{God, create in me a clean heart, and renew a steadfast spirit
within me.} That he also understood how to arrange his passages of Scripture into ``a well-connected
whole'' according to ``the analogy of Scripture'' we dare not doubt either; he was, after all, ``ein Gelehrter''.
Nicodemus did not, to be sure, know ``the Pauline doctrinal system'' (with such a ``doctrinal system''
literature had not yet been enriched at that time); but he would certainly have been able to give us a
lecture, at once clear and tasteful and correct, on the Davidic ``doctrinal system''. And yet --- how
strange --- the teacher of Israel has now, as it were, all at once gone
% --- p. 516 ---
\opage{516} %
from the ``doctrinal system'', has forgotten the parallel passages on regeneration, and now stands here groping:

{\large How can these things be?}

Yes, how can these things be? Nicodemus ponders, and has not yet understood it: hermeneutics
ponders, and has not yet fathomed it.

Then perhaps it is ``Christology'' that Nicodemus lacks. If Nicodemus could now gain an
insight into ``Christology'', would he then comprehend --- the wonder of regeneration?

\begingroup\large
Verily, verily, I say unto thee: we speak that we do know, and testify that we have seen; and ye receive
not our witness. If I tell you earthly things, and ye believe not, how shall ye
believe, if I tell you of heavenly things? ---
\par\endgroup

And yet the Savior will not disdain to speak to Nicodemus of heavenly things. --- If there is anyone
who can furnish a good ``Christology'' in the Church, it must surely be the Lord Christ himself.

\begingroup\large
And --- ``Christology'' runs --- no man ascendeth up to heaven but he that came down from heaven,
the Son of man which is in heaven.
\par\endgroup

Nicodemus ponders.

\begingroup\large
And --- ``Christology'' runs --- as Moses lifted up the serpent in the wilderness, even so must the Son of
man be lifted up; that whosoever believeth in him should not perish, but have eternal
life.
\par\endgroup

Nicodemus ponders.

\begingroup\large
For --- ``Christology'' runs --- God so loved the world, that he gave his only begotten
Son, that whosoever believeth in him should not perish, but have everlasting life.
\par\endgroup

% --- p. 517 ---
\opage{517} %
Nicodemus ponders.

So in this ``Christology'' there is talk both of the present and of the future, of the
earthly and of the heavenly, of the serpent in the wilderness and of the Son of man on the cross; of the light of the
world, of the judgment of the world, of the darkness in men's souls. And ``the Idea'' dawns, and the discourse swells in
the evangelist's rich, inexhaustible recollection, and ``the holy light shines through the Idea''; but --- it
is late at night:

Nicodemus vanishes.

--- We must not, however, now hold ourselves too one-sidedly to theory either and expect everything from
``ecclesiastical science''; we should, after all, also attend to practical immediacy, to the signs of the
times, ``the living word'' and the great Christian popular movement.

Gospel of John, Ch.\ 4, 1---35.

\begingroup\large
When therefore the Lord knew how the Pharisees had heard that Jesus made more disciples, and baptized more
than John (though Jesus himself baptized not, but his disciples), he left Judaea, and departed again
into Galilee.
\par\endgroup

Without heeding the ingrown national hatred that separated Jews and Samaritans, so that a true
patriotic Jew would rather make a detour through Peraea than defile his feet with the dust of Samaria,
Christ with his disciples went the straight way through Samaria.

\begingroup\large
Then cometh he to a city of Samaria, which is called Sychar, near to the parcel of ground that Jacob gave to Joseph
his son. Now Jacob's well was there. Jesus therefore, being wearied with his journey, sat thus on the well.
\par\endgroup

It was at the hour of noon. The disciples were gone away into the city to buy food. So the
Savior now sat there alone by Jacob's well.

% --- p. 518 ---
\opage{518} %
\begingroup\large
There cometh a woman of Samaria to draw water: Jesus saith unto her: Give me to drink.
\par\endgroup

This woman can, as it seems, present herself as a thoroughly popular representative of practical
immediacy. She has not, to be sure, made hermeneutical studies like a Nicodemus, but she has
nonetheless a very good ear for language: she can hear at once from the stranger's pronunciation that he must
be a Jew.

\begingroup\large
How is it that thou, being a Jew, askest drink of me, which am a woman of Samaria?
\par\endgroup

For her own part the jovial woman of the people would certainly very gladly have given the
stranger who sits here at the hour of noon a refreshing drink from Jacob's well. But --- may she
also do so on behalf of her people, her nation? ``How is it that thou, being a Jew, askest drink of me,
which am a woman of Samaria?'' At this question the water jar nearly falls from her
hand. The national concern occupies the talkative woman throughout the whole conversation to such a degree that
the narrative does not even know whether the Savior in the end got anything to drink.

\begingroup\large
If thou knewest the gift of God, and who it is that saith to thee: Give me to drink; thou wouldest have asked of him,
and he would have given thee living water.
\par\endgroup

It must not now surprise us that the woman with ``the heart of the people'' and ``the living word'' understands
the Savior's words about \emph{the living water} in her own way. That lies plainly in her being so
naive, so popular, so full of the national.

\begingroup\large
Lord, thou hast nothing to draw with, and the well is deep; from whence then hast thou that living water?
Art thou greater than our father Jacob, which gave us the well \dots
\par\endgroup

% --- p. 519 ---
\opage{519} See, antiquity, the great ancestral antiquity lights up for her; therefore she now cannot understand at all
the stranger's language, although its tones sound so akin. On the other hand she can say it straight
out to whom father Jacob gave this well, say it straight out that this well is a well of Samaria,
yes, say it straight out that father Jacob himself drank of this well with all his children and
all his cattle.

\begingroup\large
Whosoever drinketh of this water shall thirst again. But whosoever drinketh of the water that I
shall give him shall never thirst; but the water that I shall give him shall be in him
a well of water springing up into everlasting life.
\par\endgroup

So at last it comes to an explanation. But no, it does not yet come to an explanation; for
the woman will hear no explanation. For her part she cannot, to be sure, deny that the spring the
stranger would open for her must certainly be of a deeper source than all the wells hitherto
known in Samaria. But for certain reasons she will nevertheless not ponder this further. What does
she do, then? She brushes the whole thing off, as ``a jest'': she is so naive, so popular, so full of
the national.

{\large Go, call thy husband, and come hither!}

It takes a serious turn; for this woman's domestic affairs are not exactly in the best
order: she has evidently gone astray in the social. That she has already had five husbands, and that the
sixth she now has is not even her own husband, is admittedly no concern of the popular, and
cannot to that extent be put down to the account of nationality. On this point she is surely somewhat ahead
of the other women in the city of Sychar. Had she, on the other hand, lived in the present, and in particular in the
land of liberty where consciousness enjoys its highest freedom, she could, on account
% --- p. 520 ---
\opage{520} %
of her rich many-sidedness, perhaps best have put herself at the head of an association that intended to
work for --- ``the emancipation of woman''.\footnote{That this view is not the edifying one
the author hardly needs to remark. The historical is not exactly always the edifying.} ---
How matters otherwise stand in her conscience we cannot see. She does not yet
let anything show. But --- heart she has, spirit she has, spirited she is, and that is, after all,
the main thing in the national.

\begingroup\large
Lord, I perceive that thou art a prophet. Our fathers worshipped in this mountain; and ye say that
Jerusalem is the place where men ought to worship.
\par\endgroup

The woman now takes a step forward. Not that she should go back into herself; although the Savior
has even (v.\ 17---18) made use of his divine knowledge in order to help her out of
the enchantment. How should she ever come to see her own soul, as long as she keeps looking
so fixedly at Gerizim and Jerusalem? That the chief ecclesiastical question of the time is not foreign to her
goes without saying. Religion is, after all, precisely the great popular concern. With her
many-sidedness it naturally cannot fail that she must also be enthusiastic about --- a
folk church, and take up the cause of what is distinctive in popular worship; for it is, after all, precisely
the popular distinctiveness of worship that is to set the crown on the national.

\begingroup\large
V.\ 23. But the hour cometh, when the true worshippers shall worship in spirit and in truth; for
the Father also seeketh such to worship him.
\par\endgroup

The Savior has patience with the poor woman. He too touches on the national opposition
(v.\ 22), although in
% --- p. 521 ---
\opage{521} %
a way that can by no means flatter her popular vanity. If there must be talk of the
national oppositions and popular preparations, then it is, after all, not to the Samaritan but to the Jew
that the preference belongs. Yet this preference is a vanishing one: the true religion is raised above
all national cleavages, highly exalted above the petty jealousy of the peoples. The true God is not a
God of Gerizim, nor a God of Jerusalem, that he should bind himself to a piece of ground in order to
become a people's god.

\begingroup\large
God is a Spirit, and they that worship him must worship him in spirit and in truth.
\par\endgroup

Since the woman can now no longer hold her own in the national by patriotic exploitation of the past,
she turns with her patriotic hope toward the future; whereby it must be noted that the only thing she
here skips over is her own present time.

\begingroup\large
I know that Messias cometh (which is called Christ;) when he is come, he will tell us
all things.
\par\endgroup

If only she might now perceive the presence of the Eternal!

{\large I that speak unto thee am he.}

The woman's behavior at this moment is very remarkable. The stranger reveals himself as he is;
the Messiah stands before her. But what impression? One would think that now at last she felt
the admonition of conscience; one would now expect that the sinful woman would involuntarily cast herself down
before the Messiah, and pray: ``Lord, forgive me my sin! Alas! I have had five husbands, and the man I now
have is not my husband.''\footnote{Such a thing would surely not have been left out in an invented
narrative; but the historical is not always the edifying.} What is it, then, that keeps her
upright? Why, it is precisely the social consciousness. Instead of an anguish of conscience: ``The Messiah
% --- p. 522 ---
\opage{522} %
judges me, he knows all things,'' she suddenly gets an inspiration of the sociable:
``I wonder what people in the town will say when they hear all this?''

{\large The woman then left her waterpot, and went her way into the city, and}\par\noindent
told it all straight out, and occasioned --- a great popular awakening.

Here we now stop before these pictures of life, in which the present will find some features of its own
physiognomy indicated as in dim silhouettes. We stop, and examine ourselves, whether we have not
perhaps miscalculated in precisely the same way as that scribe Nicodemus who wrapped himself
in the cloak of obscurity and counted on the quiet working. We stop, and examine ourselves, whether we might
not at this moment go astray, as that naive (but perhaps unassuming) woman
went astray, in the social. We stop; for it is well known that one can misjudge the quiet
working in a twofold way: partly by rating it too high, as if it could exempt a
man from the decisive step, and dress up considerate caution so well that it
positively looked like a cardinal virtue; partly by rating it too low, as if the
inconsiderate appearance in public, the improvising presumption, the practical, pushing immediacy should without
further ado be recommended to public recognition. We stop; for it is well known
that one can go astray in the social in two directions: not only in such a way that, like the woman
at the well, one forgets oneself in wanting to work for ``the others''; but also in such a way that,
painfully affected by the social heedlessness of the awakenings, one shrinks into a selfish exception. The
latter byway is even more dangerous than the former. The popular awakenings of feeling (even if for
the rest they take a one-sided and crooked direction) do
% --- p. 523 ---
\opage{523} %
always have this about them, that they are at bottom so naive: a few knocks in the school of experience will soon
bring the abler individualities to themselves again, and show them where the boundary is. If, on the other hand, a
personality has broken with all sympathies of feeling in such a way that, like a Cynic, he even
denies the purpose of society, then he has thereby --- not freed himself from the idea of society
(for the general consciousness, once awakened, cannot be annihilated), but gone astray in the
social. This is easy to see. He who cannot sympathize at all with his fellow men must
naturally find himself in a negative and repellent relation to his surroundings. But the negative
relation is still a relation. Even a Diogenes needs people at every moment, in order
to get it said that he does not need them.

Under these presuppositions it will now perhaps be easier to understand how, from the standpoint of
the faith of the Gospel, we should grasp and apply those examples.

The preliminary step by which Nicodemus intends to prepare a quiet, calm transition from the
old to the new is neither blamed nor praised by the Savior. Christ does not begin
by calling upon Nicodemus to lay down his office and step out of the Sanhedrin. --- That rich
young man who came running on the road, and fell on his knees before ``the good Master'', and by this kneeling
proclaimed himself one awakened who was surely not afraid of attracting attention, at once received an injunction that
might well serve to make the transition striking even in externals; but the quiet
Nicodemus surely has leave to remain in his present position. The absolute
condition, on the other hand, is just as decisive for Nicodemus as for the rich young man. As it goes with
position, so also with insight. Christ does not speak otherwise to the ruler of the
Jews than to the lowliest among the Samaritans: the teacher
% --- p. 524 ---
\opage{524}%
in Israel and the woman in Samaria, without distinction, find it equally easy and equally hard to understand
him. --- ``But,'' the theologians of the School now object, ``a Nicodemus can by no means
be said to represent ecclesiastical science for us: the Pharisee was, after all, not yet born again, had
not yet properly said: \textit{credam}, and consequently could not begin on:
\textit{intelligam} either''. --- This is indeed very true. But that Nicodemus may nonetheless well
be said to represent for us the complete indifference of the faith of the Gospel toward ``ecclesiastical science''
the theologians themselves will be able to convince themselves by entering vividly into the relation. Had
Christ acknowledged ``the learned man's'' exertions to comprehend the economy of regeneration
(\textit{ordo salutis}), had he seen that a scientific doctrinal system would in the end
become necessary for the preservation of faith, then he would surely have comforted Nicodemus with
the prospects of the future.

``You seek the concept; that is very good; for my Church cannot in any case subsist without
the concept. But before you can comprehend the concept, you must first be born again. Go now and let yourself be born again
of water and of the Spirit, then come again: we will then, in a quiet evening hour, talk further together about the right
\textit{credam, ut intelligam}.''

It was, even if to the disciples' surprise (Joh.\ 4, 27), not beneath the Master's dignity to
speak with the woman at Jacob's well; but it is, even if perhaps to the theologians' surprise, in
truth deeply beneath Christ's dignity to enter with Nicodemus into any learned, scientific
discussion. It is of no use to appeal to the changed conditions either of the spirit of the age or of
science. A thorough researcher may be as convinced as he likes of ``the necessity of the critical introduction
for the preservation of the Church''; if it should befall him that, like that researcher in
Israel, he once found himself alone, quite alone with
% --- p. 525 ---
\opage{525} %
the Savior of the world: he would certainly not dare to raise doubts about ``the fourth Gospel'', and would surely
not think of making a great fuss about ``the Passover reckoning'' either. No, that he
certainly would not. But --- why not? We ask him in his conscience: why not, then? --- A
dogmatic genius may otherwise be as convinced of the ecclesiastical significance of ``Christology'' as
he likes; if it should befall him that he once, like that teacher of Israel, found himself alone,
quite alone with Christ himself: he would surely hardly lay before him the problems of ``the two
natures'' (\textit{communicatio idiomatum}), would surely hardly presume to question the Redeemer
himself concerning the metaphysical relation between the miracle of his power (\textit{miraculum
potentiæ}) and the miracle of his knowledge (\textit{miraculum scientiæ}). No, that he surely would
not. But --- why not? We ask him in his conscience: why not, then?\footnote{Whether
a question of faith is conceived in such a way that its answer must be of religious or of
scientific significance can best of all be tested by imagining the question addressed
directly to Christ himself. If the question is religiously framed, the Lord will surely
answer it in his Gospel, whether it otherwise turns on a spiritual or a worldly
concern. If, on the other hand, the question is scientifically framed, its answer, however ecclesiastical the subject may
otherwise be, must without further ado be excluded from the sphere of faith and referred to science;
whereby it is then to be observed that science does not undertake to comprehend the
incomprehensible.}

As regards the Samaritan woman's popular character, we must still consider the picture from another side.
Jesus' words, to be sure, did not at once make any very deep impression on this woman:
% --- p. 526 ---
\opage{526} %
she has, however, been brought to the supposition that the stranger might possibly be the Messiah.
(Quite sure of it she is not exactly; she must naturally first learn what people in the city will
say.) On the Lord himself the conversation has had a strong effect. (V. 32--34.)

\begingroup\large
My meat is to do the will of him that sent me, and to finish his work.
\par\endgroup

The Savior himself thus expects a popular awakening, and prepares himself to receive the Samaritans
coming from the city, when, doubtless at the sight of the approaching crowd of people,
he addresses these words to the disciples:

\begingroup\large
Say not ye: there are yet four months, and then cometh harvest? Behold, I say unto you: lift up your
eyes, and look on the fields; they are white already to harvest. (V. 35.)
\par\endgroup

So the woman, for her part, has after all also done something for the spreading of Christ's name. By her
``living word'' she has as in a moment gathered a whole congregation, which can now hear Christ himself, and
thereby come to believe. Yes, to that extent she has undeniably deserved well of the town of Sychar. She
has, as the saying goes, been an instrument in the Lord's hand, and has accordingly in this respect received
her historical memorial. But it is another question whether anything has now also really
come of her own faith. We will by no means doubt it, but only add this, that
she surely did not save her soul by intoxicating herself in folk-church-national awakenings, but
by herself fleeing to Christ, collecting herself from all distraction, drinking from the pure source, the
inner spring, which springs up into everlasting life.

That Christendom needs a regeneration can certainly not be denied; but that Christianity is
heterogeneous with any worldly purpose whatsoever can surely not be denied either. If the Gospel
% --- p. 527 ---
\opage{527}%
is God's revealed word, then it must surely be able to help itself by itself, and does not need, absolutely
understood, first to be helped along by anything else: it makes everything a means and instrument for itself,
but, be it noted, in such a way that it does not itself become means and instrument for anything else. ---
The teachers of the Church and the leaders of the people gather and deliberate on the good of Christianity. The former would
span the terrain of culture and help revelation along by ``ecclesiastical science''; but when
one in this way makes science a means for Christianity, Christianity in turn must put up
with becoming a means for science: so it comes out in the end that science is science,
and at bottom has nothing to do with Christianity. The latter, on the other hand, think they work most powerfully
for the cause of the Gospel when they can really stir the popular imagination, blow wind into the flowing word and
breathe the breath of life into the nostrils of nationality. But when by this procedure one would make the popular
a means for the Christian, it happens conversely that the Christian all at once becomes a means for
the popular: so it comes out in the end that the great popular movements go their own worldly
way and have nothing to do with the Christian. Everyone would gladly contribute his share to the good of
the whole; everyone would so gladly save, if not exactly the whole world, then at least his part of
the world; everyone would then in the Lord's name be a little savior of the world. But then come the many ideas
and the many proposals and the many plans and the many deliberations, and alas! then comes ---
the delusion. It almost looks as though one had forgotten who the Savior is; it is as if one would
not merely be helped by him, but also help him along a little; it is as if one would not merely do
something for him, but also, so that he might really come to the fore, do something to him. One would refine
him in ``Christology'', another would lend him a modern cloak, a third would adorn
% --- p. 528 ---
\opage{528} %
him with the national colors: all would do something to him, so that he might now really break
through, and take possession of his kingdom among us. See, thus there is striving and working in all directions:
for and against the state church, for and against science, for and against religious freedom, for and against
standstill, for and against movement; but --- it bears no fruit. Let the vineyard be full of
laborers, let all the laborers be in full swing: both ``the folk-spirits'' and ``hermeneutics'' and ``the
living word'' and ``Christology'' and ``the birdsong'' and ``the critical introduction'', let all these
laborers and labors only be in full swing; but ---

\begingroup\large
Verily I say unto you: whosoever shall not receive the kingdom of God as a little child, he shall in no wise
enter therein.
\par\endgroup

How this fundamental word changes the matter! For it is quite certain that a little child must not
burden itself with all too great tasks. A little child cannot on its own take part in the father's
whole household. --- In relation to the heavenly Father every human being must --- we repeat it ---
feel himself a little child. This is for the believer the absolute, the unconditional, the
first. But this simple relation to God is now at the same time to be expressed and, so to speak,
justified in a multiplicity of relative relations to the world. In these relations to the world lies
difference. Here there is then a great difference between children and grown-ups, those not of age and those of age; and
this difference is not to be understood as if believers were to be minors face to face with
unbelievers, or at any rate helpless children face to face with the worldly-wise. In a deeper
sense the relation will turn around. The more childlike a man is in his relation to God, the more powerfully and
emphatically will he be able to work in his relations to the world; for the simple man everywhere finds a
point of support in the divine will. If, on the other hand, anyone
% --- p. 529 ---
\opage{529} %
would labor for the religious tasks, and yet disdains to relate himself rightly simply and childlikely
to God, it can easily happen to him that with all his cleverness about the world he sometimes behaves
a little childishly in the world. Unbelievers imagine that religiousness weakens a man's
power of action, and prevents him from appearing with authority; but the matter stands the other way round. The
greater the obedience in the relation to God, the greater the authority in the relation to the world: Christ's absolute
authority is explained, after all, by this, that he \emph{became obedient unto death, even the death of the cross} (Phil.\ 2, 8).
When thus the absolute, the unconditional, is for every man \emph{the first}, then the
relative endeavors follow in turn from it as \emph{the second}. --- \emph{Seek ye first the kingdom of God and his
righteousness, and all these things shall be added unto you} (Matth.\ 6, 33). This ``\emph{first}''
no hermeneutics may presume to garble for us. But if only the antecedent is allowed to
keep its place, believers will surely also in the end be able to agree on the consequent
that goes with it. Let everyone, then, who works for ecclesiastical science, and knows in himself that he does not
misjudge its relation to the Gospel, work with blessing in his calling, and contribute his share to
spreading a true knowledge of God. And if anyone has a living and powerful word for the
awakening of the peoples, and knows in himself that he does not confuse the worldly with the holy, let him speak
with all boldness according to the gift of the Spirit that is given him: such a one is a minister of the word, a
steward of the Gospel, for the help and comfort and strengthening of many. In this way the kingdom of God
is spread, and \emph{He who sees that the harvest is great, but the laborers few, will himself pray the Lord of the harvest
that he send forth laborers into his harvest} (Matth.\ 9, 37--38).

But when the relation is turned around; when that which unconditionally is and is to be --- \emph{the first} is only
in the manner of speaking allowed to be
% --- p. 530 ---
\opage{530} %
the first, while in deed and actuality it sinks down to become \emph{the second}; when
that which in the nature of the case is and ought to be --- \emph{the second} only in the manner of speaking becomes the second,
while yet in deed and actuality it steals into being \emph{the first}: then there comes not
the kingdom of God, but an altogether different kingdom. But ---

\begin{center}
Jesus Christ is the Savior of the world.
\end{center}

\begin{center}
  \rule{0.55\textwidth}{0.4pt}
\end{center}

\end{document}
